\documentclass[12pt]{book}%

%\usepackage{fontspec}[greek]
%\usepackage{fourier-otf} % pour les arcs (double brève qui compte comme une longue)
%\usepackage{arcs} % pour les arcs (double brève qui compte comme une longue)
\usepackage{graphicx} % pour l'image de la coupe d'Onésime
\usepackage{fontspec}%
\usepackage{calc} % \ratio
\usepackage{fancyhdr}% en-têtes personnalisés
\usepackage[strict]{changepage}%pour modifier les marges des footnotes
\pagestyle{fancy}
\fancyfoot[LE,RO]{\thepage}
\fancyfoot[C]{}
\fancyhead[LE,RO]{}
%\fancyhead[RE]{\leftmark}
%\fancyhead[LO]{\rightmark}
\newcommand{\enTeteG}[1]{\fancyhead[LE]{#1}}% En-tête principale (pages paires de gauche)
\newcommand{\enTeteD}[1]{\fancyhead[RO]{#1}}% En-tête secondaire (pages impaires de droite)
\renewcommand{\headrulewidth}[0]{0pt}%filet de en-tête haut supprimé
\renewcommand{\footrulewidth}[0]{0.05mm}%filet de bas page
%\usepackage{totpages}
%\cfoot{{\thepage} / \ref{TotPages}}
%\usepackage{xunicode}%
%\usepackage{yfonts}%
\usepackage{hyperref}% pour le courriel
\usepackage{amsmath}% pour \overset
\usepackage{cancel}% pour \inRas
\usepackage{quoting}%
%\usepackage{csquotes} % Pour remplacer les guillemets anglais par des français
%\MakeOuterQuote{"}
\usepackage{ulem}% pour souligner en pointillé
\usepackage{comment}%
\usepackage{xcolor}%
\usepackage{fourier}% pour \widearc de la synizèze
%\usepackage{ragged2e}% pour la justification
%scale = {hscale, vscale}
%\usepackage[scale={0.75,0.8}]{geometry}%\usepackage[total={12cm,20cm}%Harmattan
%\usepackage[total={12cm,22cm}, top=3.5cm, right=4.5cm, footnotesep=1cm, includefoot]{geometry}% pour une page impaire
\usepackage[twoside, bindingoffset=1cm,total={12cm,22.5cm}, top=3.5cm, right=4.5cm, footnotesep=1cm, includefoot]{geometry}
%\usepackage{xpatch}
%\setlength{\hoffset}{-18pt}        
%\setlength{\oddsidemargin}{20pt} % Marge gauche sur pages impaires ATTENTION écrit 2 fois
%\setlength{\evensidemargin}{9pt} % Marge gauche sur pages paires
%\setlength{\marginparwidth}{54pt} % Largeur de note dans la marge
%\setlength{\textwidth}{481pt} % Largeur de la zone de texte (17cm)
%\setlength{\voffset}{-18pt} % Bon pour DOS
%\setlength{\marginparsep}{7pt} % Séparation de la marge
%\setlength{\topmargin}{0pt} % Pas de marge en haut
%\setlength{\headheight}{13pt} % Haut de page
%\setlength{\headsep}{10pt} % Entre le haut de page et le texte
%\setlength{\footskip}{27pt} % Bas de page + séparation
%\setlength{\textheight}{708pt} % Hauteur de la zone de texte (25cm)


%\newfontfamily\policeGr{Liberation Serif}[
%  Script=Greek,
%  NFSSFamily=liberationserif,
%
%]
%\newfontfamily\policeGr{Liberation Serif}%

%\newfontfamily{\policeFr}{Garamond Premier Pro}%
%\newfontfamily{\policeLat}{Book Antiqua}%
%\newfontfamily{\policeAn}{Comic Sans MS}%
%\newfontfamily{\policeAll}{Fraktur}%
%\newfontfamily{\policeIt}{BodoniFLF}%
%\newfontfamily{\policeEsp}{Great Moravia}%
% ...
%\setmainfont{Liberation Serif}%
%\setmainfont{Liberation Serif}%
%\newfontfamily{\policeFr}{Garamond Premier Pro}%
%\newfontfamily{\policeFr}{EB Garamond}%
\IfFontExistsTF{GFS Bodoni}{\newfontfamily{\policeFr}{GFS Bodoni}\newfontfamily{\policeGr}{GFS Bodoni}\setmainfont{GFS Bodoni}}{\newfontfamily{\policeFr}{GFS Didot}\newfontfamily{\policeGr}{GFS Didot}\setmainfont{GFS Didot}}%
%\newfontfamily{\policeLat}{Book Antiqua}%
\newfontfamily{\policeLat}{GFS Neohellenic}%
%\newfontfamily{\policeAn}{Comic Sans MS}%
\IfFontExistsTF{Courier 10 Pitch}{\newfontfamily{\policeAn}{Courier 10 Pitch}}{\newfontfamily{\policeAn}{FreeMono}}%
\newfontfamily{\policeEsp}{Abyssinica SIL}%
%\newfontfamily{\policeAll}{Fraktur}%
%\newfontfamily{\policeAll}{fraktur-bt.ttf}%
\IfFontExistsTF{Bertholdr Mainzer Fraktur}{\newfontfamily{\policeAll}{Bertholdr Mainzer Fraktur}}{\usepackage{yfonts}}%
%\newfontfamily{\policeIt}{BodoniFLF}%
\newfontfamily{\policeIt}{GFS Artemisia}%
%\newfontfamily{\policeEsp}{Great Moravia}%
%\newfontfamily\policeGr{GFSBodoni.otf}[BoldFont = GFSBodoniBold.otf, ItalicFont = GFSBodoniIt.otf, BoldItalicFont = GFSBodoniBoldIt.otf]%
%\setmainfont{GFSBodoni.otf}[BoldFont = GFSBodoniBold.otf, ItalicFont = GFSBodoniIt.otf, BoldItalicFont = GFSBodoniBoldIt.otf]

%\setdefaultlanguage[variant=ancient]{greek}%
\usepackage{polyglossia}%
\setdefaultlanguage{french}%
%\setotherlanguage[variant=ancient]{greek}%
%\setmainlanguage[variant=ancient]{greek} 
%\setotherlanguage[variant=ancient]{greek} 
%\setmainfont{Libertinus Serif} 
%\newfontfamily{\greekfont}{Liberation Serif}[
%  Script=Greek,
%  NFSSFamily=liberationserif,
%]

\usepackage[maxlevel=3]{csquotes}%
%\MakeOuterQuote{"}% pour les guillemets français
%\usepackage{verse}%
%\usepackage}[squarebr]{gmverse}% pour le rejet avec un crochet droit
\usepackage[parapparatus]{reledmac}%
\usepackage[parledgroup, nomaxlines]{reledpar}% 
%\usepackage[parledgroup, nomaxlines, nopbinverse]{reledpar}% nopbinverse empêche de casser un groupe
%\usepackage[parledgroup]{reledpar}% nomaxlines permettait d'avoir un nombre de lignes différents sur des pages qui se font face.
%\parledgrouptrue
%\setgoalfraction{0.9}% sept 2022
%\makeatletter
%\xpatchcmd{\mp@append@Xnotes}{\@parboxrestore}{\@parboxrestore\leftskip-1in}{}{}
%\makeatother
% Pour que les commentaires apparaissent sur les deux pages
\AtBeginDocument{%

%\Xbeforeinserting[B]{\espace{-2cm}}% code exécuté avant l'appel à la note
%\Xafternote[B]{\espace{20cm}}
%\Xbhookgroup[A,C]{\espace{-2cm}}
%\Xprenotes{2cm}
%\setlength{\skip\footins}{-2cm}
%\Xhangindent[B]{3cm}
%\sethangingsymbol{\protect\hfill}
%\versehangrightsquare% le rejet d'un vers se fait avec un crochet droit
\rightnoteupfalse 
\Xmaxhnotes[A,B]{0.50\textheight}%
\Xmaxhnotes[C]{0.22\textheight}%
%\Xmaxhnotes{0.9\textheight}%
%\maxhnotesX{0.45\textheight}%
\setRlineflag{}%
%\newcommand{\decoupeVers}[1]{\linenumannotation{#1}}%
\newcommand{\decoupeVers}[1]{}%
%\setlength{\oddsidemargin}{100pt} % Marge gauche sur pages impaires ATTENTION écrit 2 fois
% La Largeur du commentaire
%\Xwidth{\columnwidth+\marginparsep+\ledrsnotewidth}
% \Xgroupbyline[A,C]
\Xgroupbyline[A,C]
% \Xonlyside[A]{R}% imprime les testimonia page de droite
%\Xafterrule[A]{-1em}% supprime un peu d'esp vertical
%\Xwidth{\columnwidth+\marginparsep+\ledrsnotewidth}
\Xwidth{14cm}
%\Xsethangindent{20pt}%décale les footnotes
%\Xhangindent{-3cm}
%\Xparindent{-10cm}
%\sethangindentX{20pt}
%\Xhangindent = 20pt%
}%


\newcommand\Petit{\fontsize{9pt}{11pt}\selectfont} % x 1,2 pour e second argument qui est l'espacement interligne
\newcommand{\Taille}[1]{\fontsize{#1pt}{12pt}\selectfont}
\newcommand{\TaillePt}[1]{\fontsize{#1}{15pt}\selectfont}
\newlength{\LargAuto}%
\newlength{\LargTheor}%
\newlength{\Vrai}%
\newlength{\TP}%
\setlength{\TP}{12pt}%
\newcommand{\Automatique}[1]{%
\setlength{\LargTheor}{-72pt}%
\addtolength{\LargTheor}{\textwidth}%
\settowidth{\LargAuto}{#1}%
%\divide \LargTheor by \LargAuto%  \real{1.} *
\setlength{\Vrai}{\TP * \ratio{\LargTheor}{\LargAuto}}%
%\edef\TP{12* \ratio{\LargTheor}{\LargAuto}}%
{\TaillePt{\Vrai}#1}%
}%
\newlength{\largtext}%
\newcommand{\Strophe}[1]{\ledleftnote{%
\settowidth{\textelargeur}{\footnotesize{\francais{#1}}}%
%\espace{-\textelargeur}
\begin{minipage}{\textelargeur}\footnotesize{\francais{#1}}\end{minipage}}}%
\newcommand{\renverse}[1]{\rotatebox[origin=c]{180}{#1}}
\newcommand{\italique}[1]{\textit{#1}}%
\newcommand{\citComm}[1]{\begin{quoting}[vskip=0pt] #1 \end{quoting}}% vskip = \topsep + \partopsep par défaut
\newcommand{\passeLigne}[0]{\skipnumbering}%
\newcommand{\fin}[0]{\linebreak\passeLigne\protect\hfill $[$ }% Pour le rejet 
\newcommand{\ligne}[1]{\setline{#1}\hidenumbering}%
\newcommand{\espace}[1]{\hspace{#1}}%
\newcommand{\cesure}[0]{ [...] }%
\newcommand{\nieme}[1]{#1\italique{e}}%
\newcommand{\lettreBreve}[1]{$\breve{\mbox{#1}}$}
\newcommand{\synizese}[1]{$\widearc{\mbox{#1}}$}
\newcommand{\inRas}[2]{\if\relax\detokenize{#2}\relax\text{\sout{#1}}\else$\overset{\text{#2}}{\text{\sout{#1}}}$\fi}
\newcommand{\auDessus}[2]{$\overset{\text{#2}}{\text{#1}}$}
%\newcommand{\etc}[0]{κτλ}%
%\newcommand{\cit}[0]{\espace{1cm}} % indent. pour une citation
\newcommand{\expo}[1]{\textsuperscript{#1}}%
\newcommand{\alinea}[0]{ \ \ }%
\newcommand{\Fr}[0]{ Ar. \italique{fragm.} (Kassel-Austin) }%
\newcommand{\Syn}[0]{ Συναγωγή }%
\newcommand{\SynB}[0]{ Συναγωγή\expo{B} }% la deuxième version pour la lettre A
\newcommand{\ClemAlex}{ Clemens Alexandrinus }
\newcommand{\Eustathe}{ Eustathius }
\newcommand{\AelDion}{ Aelius Dionysius }
\newcommand{\EtymMagnum}{ Etymologicum magnum }
\newcommand{\EtymGen}{ Etymologicum genuinum }
\newcommand{\EtymSym}{ Etymologicum Symeonis }
\newcommand{\Apostol}{ Michael Apostolius }
\newcommand{\Macaire}{ Macarius } 
\newcommand{\Phryn}{ Phrynichus \italique{Praep. soph.} }
\newcommand{\Stobee}{ Joannes Stobaeus }
\newcommand{\EpimAlphab}{ Epimerismi alphabetici in Homerum (ed. Cramer) }
\newcommand{\GregCor}{ Gregorius Corinthius }
\newcommand{\Antiatt}{ Antiatticista (ed. Bekker) }
\newcommand{\Harp}{ Harpocration }% Autre édition de référence : Dindorf
\newcommand{\Rpc}[0]{R\expo{pc}}%
\newcommand{\Rac}[0]{R\expo{ac}}%
\newcommand{\Rdeux}[0]{R\expo{2}}%
\newcommand{\Mon}[0]{ M }%
\newcommand{\PsiUn}[0]{PSI-1194}
\newcommand{\PsiDeux}[0]{POxy-3839}
\newcommand{\PsiTrois}[0]{POxy-3840}
\newcommand{\PsiQuatre}[0]{POxy-4935}
\newcommand{\PsiCinq}[0]{POxy-1176}%P50

\newcommand{\francais}[1]{{\policeFr #1}}%
\newcommand{\latin}[1]{{\policeLat #1}}%
\IfFontExistsTF{Bertholdr Mainzer Fraktur}{\newcommand{\allemand}[1]{{\policeAll #1}}}{\newcommand{\allemand}[1]{{\textfrak{#1}}}}%
%\newcommand{\allemand}[1]{{\textfrak{#1}}}%
%\newcommand{\grec}[1]{{\policeGr #1}}%
\newcommand{\grec}[1]{#1}%
\newcommand{\anglais}[1]{{\policeAn #1}}%
\newcommand{\italien}[1]{{\policeIt #1}}%
\newcommand{\espagnol}[1]{{\policeEsp #1}}%

\newcommand{\Nauck}[0]{N\expo{2}}%
\newcommand*{\auteur}[1]{\textsc{#1}*}%
\newcommand{\cf}[0]{\italique{cf.}\xspace}%
\newcommand{\chap}[1]{\^{#1}}%
%\newcommand{\chap}[1]{$\hat{#1}$}


%\marginparwidth = 1.2\marginparwidth
%\ledrsnotewidth = 1.2\marginparwidth
\newlength{\larg}%
\newcommand{\decalage}[1]{%
\settowidth{\larg}{#1}%
\espace{\larg}%
}
\newcommand{\decalageN}[1]{%
\settowidth{\larg}{#1}%
\espace{-\larg}%
}
%\newcommand{\apparat}[2]{\ledrightnote{\espace{-1.2\marginparwidth}\begin{minipage}{1.4\marginparwidth}\footnotesize{\textbf{#1} #2}\end{minipage}}}%
\newcommand{\apparat}[2]{\ledrightnote{\espace{-2cm}\begin{minipage}{4.5cm}\footnotesize{\textbf{#1} #2}\end{minipage}}}%
%\newcommand{\lecon}[1]{\dotuline{#1}} % la leçon retenue de l'apparat critiqe
\newcommand{\lecon}[1]{\textbf{#1}}
%\newcommand{\douteux}[1]{\dotuline{#1}} % pour les lettres douteuses : un point en-dessous
\newcommand{\douteux}[1]{$\underset{\cdot}{\text{#1}}$}
\newcommand{\difference}[1]{\dotuline{#1}}


%\setsidenotesep{ sep }%
\newcommand{\espApp}[0]{ || }%\espace{10pt}%
\newcommand{\espTest}[0]{ || }%\espace{10pt}%
\newcommand{\ligneApp}[1]{ \ \ \textbf{#1} }%
%% 3 arguments : les numéros de ligne, le lemme et la note elle-même



%\newcommand{\commentaire}[3]{\edtext{}{\Bfootnote[nonum, nosep]{\espace{-2cm}\begin{minipage}{16cm}\espace{-1cm}\textbf{#1~~#2} #3\end{minipage}}}}%
\newcommand{\commentaire}[3]{\edtext{}{\Bfootnote[nonum,nosep]{\begin{adjustwidth}{-1.5cm}{-2cm}\espace{-1.5cm}\textbf{#1}\decalageN{\textbf{#1}}\espace{0pt}\espace{1.5cm}\textbf{#2} #3\end{adjustwidth}}}}% Avec une * : adjustwidth* pour la symétrie des pages
%\renewcommand{\footfudgefiddle}{68}%Note that you must call it before 
% \Xarrangement[C]{paragraph}% entrâine trop d'erreurs, cf. documentation
\newcommand{\debutNote}[1]{\edtext{}{\Cfootnote[nonum, nosep]{\espace{-1cm}\espace{-15pt}\textbf{\textbf{#1}}}}}%
%\newcommand{\debutNote}[1]{\edtext{}{\Cfootnote[nonum, nosep]{\espace{-1cm}\espace{-15pt}\textbf{\dotuline{\textbf{#1}}}}}}%
\newcommand{\note}[2]{\edtext{}{\Cfootnote[nonum, nosep]{\espace{-15pt}\textbf{#1} #2}}}%
\newcommand{\testimonia}[2]{\edtext{}{\Afootnote[nonum, nosep]{\espace{-15pt}\textbf{#1} #2}}}%
%\newcommand{\testimonia}[2]{}
\newcommand{\marqueur}[1]{\textsuperscript{#1}}%
%\newcommand{\cat}[0]{$""_"^"$}%
\newcommand{\cat}[0]{\texorpdfstring{\textsubscript{\textasciicircum}}{cat}}%
%\fnpos{critical-familiar}%
%\mpfnpos{familiar-critical}%

\newcommand{\eff}[1]{%
{\color{orange}#1}%
}%

\newlength{\textlarg}%
\newcommand{\typeVers}[1]{%
\settowidth{\textlarg}{\footnotesize{#1}}%
\hspace{-0.8cm}{\footnotesize{#1}}\hspace{-\textlarg}\hspace{0.8cm}}%
%\newlength{\largtext}%
%\newcommand{\strophe}[1]{%
%\settowidth{\largtext}{~~\footnotesize{\francais{#1}}}%
%~~\footnotesize{\francais{#1}}%\hspace{-\largtext}%
%}%
\newcommand{\strophe}[1]{\footnotesize{~\francais{#1}}}%
%\newcommand{\strophe}[1]{%
%\settowidth{\largtext}{\footnotesize{\francais{#1}}}%
%	\ledleftnote{%
		%\hspace{-0.5cm}\footnotesize{\francais{#1}}\hspace{\largtext}\hspace{0.5cm%}%
%		\hspace{-0.5cm}\footnotesize{\francais{#1}}\hspace{\largtext}\hspace{0.5cm}%
%	}%
%}%
%\newcommand{\personnageGrec}[1]{\ledleftnote{\textbf{\large{#1}}\espace{-1.5cm}}}%
\newlength{\textelargeur}%
\newcommand{\personnageGrec}[1]{\ledleftnote{%
%\settowidth{\textelargeur}{\large{\textbf{#1}}}%
\settowidth{\textelargeur}{\large{\textbf{Στεφανόπωλις }}}%
%\espace{-\textelargeur}
\begin{minipage}{\textelargeur}\begin{flushright}\large{\textbf{#1}}\end{flushright}\end{minipage}}}%
%\newcommand{\Strophe}[1]{\footnotesize{\francais{#1}}}%
%\newcommand{\Strophe}[1]{\ledleftnote{%
%\settowidth{\textelargeur}{\footnotesize{\francais{#1}}}%
%\espace{-\textelargeur}
%\begin{minipage}{\textelargeur}\footnotesize{\francais{#1}}\end{minipage}}}%
%\newcommand{\personnageFr}[1]{\ledleftnote{\textbf{\large{\francais{#1}}}\espace{-1.5cm}}}%
\newlength{\texto}%
\newcommand{\personnageFr}[1]{\ledleftnote{%
%\settowidth{\texto}{\large{\textbf{#1}}}%
\settowidth{\texto}{\large{\textbf{Marchande d}}}%
%\espace{-\textelargeur}
\begin{minipage}{\texto}\begin{flushright}\large{\textbf{#1}}\end{flushright}\end{minipage}}}%
%% Les différents types de vers
\newcommand{\trimIamb}[1]{%
\hangindent=20pt%
\espace{0pt}% la valeur précédente -20
#1%
}%
\newcommand{\trimIambFr}[1]{%
\hangindent=20pt%
\espace{0pt}% la valeur précédente -20
\francais{#1%
}}%
\newcommand{\trimIambCat}[1]{%
\hangindent=53pt%
\espace{33pt}% la valeur précédente -20
#1%
}%
\newcommand{\trimIambCatFr}[1]{%
\hangindent=53pt%
\espace{33pt}% la valeur précédente -20
\francais{#1%
}}%
\newcommand{\tetramIamb}[1]{%
\hangindent=10pt%
\espace{-10pt}% la valeur précédente -20
#1%
}%
\newcommand{\tetramIambFr}[1]{%
\hangindent=20pt%
\espace{0pt}% la valeur précédente -20
\francais{#1%
}}%
\newcommand{\tetramTroch}[1]{%
\hangindent=20pt%
\espace{0pt}% la valeur précédente -20
#1%
}%
\newcommand{\tetramTrochFr}[1]{%
\hangindent=20pt%
\espace{0pt}% la valeur précédente -20
\francais{#1%
}}%
\newcommand{\tetramIambCat}[1]{%
\hangindent=10pt%
\espace{-10pt}% la valeur précédente -20
#1%
}%
\newcommand{\tetramIambCatFr}[1]{%
\hangindent=20pt%
\espace{0pt}% la valeur précédente -20
\francais{#1%
}}%
\newcommand{\dimAnap}[1]{%
\hangindent=70pt%
\espace{50pt}% la valeur précédente -20
#1%
}%
\newcommand{\dimAnapFr}[1]{%
\hangindent=70pt%
\espace{50pt}% la valeur précédente -20
\francais{#1%
}}%
\newcommand{\anap}[1]{%
\hangindent=70pt%
\espace{50pt}% la valeur précédente -20
#1%
}%
\newcommand{\anapFr}[1]{%
\hangindent=70pt%
\espace{50pt}% la valeur précédente -20
\francais{#1%
}}%
% 2 dochmiaques
\newcommand{\diDoch}[1]{%
\hangindent=60pt%
\espace{40pt}% la valeur précédente -20
#1%
}%
\newcommand{\diDochFr}[1]{%
\hangindent=60pt%
\espace{40pt}% la valeur précédente -20
\francais{#1%
}}%
\newcommand{\tetramAnap}[1]{%
\hangindent=20pt%
\espace{0pt}% la valeur précédente -20
#1%
}%
\newcommand{\tetramAnapFr}[1]{%
\hangindent=20pt%
\espace{0pt}% la valeur précédente -20
\francais{#1%
}}%
\newcommand{\chantAgath}[1]{%
\hangindent=90pt%
\espace{70pt}% la valeur précédente -20
#1%
}%
\newcommand{\chantAgathFr}[1]{%
\hangindent=90pt%
\espace{70pt}% la valeur précédente -20
\francais{#1%
}}%
%% Pour les "vers" 295 à 311
\newcommand{\prose}[1]{%
\hangindent=30pt%
\espace{10pt}% la valeur précédente -20
#1%
}%
\newcommand{\proseFr}[1]{%
\hangindent=30pt%
\espace{10pt}% la valeur précédente -20
\francais{#1%
}}%
%% pour les vers courts de structure changeante dans un choeur
\newcommand{\chantChoeur}[1]{%
%{\centering
%	#1\par
%}%
\hangindent=50pt%
\espace{30pt}% la valeur précédente -20
#1%
}%
\newcommand{\chantChoeurFr}[1]{%
%{\centering
%	\francais{#1}\par
%}%
\hangindent=50pt%
\espace{30pt}% la valeur précédente -20
\francais{#1%
}}%
%% pour les vers 1015 à 1055
\newcommand{\paratragiques}[1]{%
%{\centering
%	#1\par
%}%
\hangindent=50pt%
\espace{30pt}% la valeur précédente -20
#1%
}%
\newcommand{\paratragiquesFr}[1]{%
%{\centering
%	\francais{#1}\par
%}%
\hangindent=30pt%
\espace{10pt}% la valeur précédente -20
\francais{#1%
}}%
% pour indiquer que les notes marginales s'alignent au niveau de la première ligne
%\leftnoteupfalse
%\rightnoteupfalse



\newenvironment{groupe}{%
\pstart%
\ledgroup%
}{%
\pend %
}%

\newcommand{\LSJ}[0]{LSJ}%

%% pour voir tous les numéros de ligne (version brouillonne)
\linenumincrement{1}%
\firstlinenum{1}%
%% Lignes à imprimer
\gdef\linenumberlist{2,6,10,15,20,24,30,36,40,45,50,55,60,65,69,74,80,87,91,94,102,105,110,115,120,125,131,135,140,145,150,155,160,165,170,174,180,185,190,196,201,205,215,220,225,229,234,239,245,250,261,265,270,276,280,285,290,295,300,305,310,315,320,325,330,335,340,345,350,355,360,365,370,375,379,385,390,395,400,405,410,415,420,425,430,435,440,445,450,455,460,465,470,475,480,485,490,495,499,501,505,510,515,521,526,530,535,541,545,550,556,561,565,572,575,580,585,591,596,600,609,614,620,624,629,635,640,650,656,660,665,670,675,680,685,690,695,700,708,713,720,725,729,735,741,745,750,755,761,766,770,775,780,786,790,795,800,805,810,815,820,825,831,835,840,845,850,856,860,865,870,875,880,884,890,894,900,905,915,921,930,935,940,946,950,955,960,965,970,975,980,985,990,996,1000,1010,1016,1020,1025,1030,1035,1040,1045,1050,1055,1060,1066,1071,1076,1100,1106,1116,1120,1124,1130,1141,1147,1150,1156,1161,1166,1171,1175,1180,1185,1190,1195,1200,1205,1211,1215,1221,1224,1230}



\usepackage{adjustbox}
\usepackage{array}
\setmonofont{FreeMono}
\setlength{\headheight}{16pt}
\renewcommand{\enTeteG}[1]{\fancyhead[LE]{\small #1}}
\renewcommand{\enTeteD}[1]{\fancyhead[RO]{\small #1}}

% In Spread 2, critical apparatus and testimonia are already displayed on Spread 1 Left Page:
\renewcommand{\apparat}[2]{}
\renewcommand{\testimonia}[2]{}
\makeatletter
\renewcommand*{\led@err@LeftOnRightPage}{}
\renewcommand*{\led@err@RightOnLeftPage}{}
\makeatother
\AtBeginDocument{%
  \Xmaxhnotes[A,B]{0.82\textheight}%
  \Xmaxhnotes[C]{0.55\textheight}%
}
\renewcommand{\commentaire}[3]{\edtext{}{\Bfootnote[nonum,nosep]{\begin{adjustwidth}{-1.5cm}{-2cm}\fontsize{8.0pt}{9.4pt}\selectfont\espace{-1.5cm}\textbf{#1}\decalageN{\textbf{#1}}\espace{0pt}\espace{1.5cm}\textbf{#2} #3\end{adjustwidth}}}}%
\renewcommand{\debutNote}[1]{\edtext{}{\Cfootnote[nonum, nosep]{\fontsize{7.8pt}{9.2pt}\selectfont\espace{-1cm}\espace{-15pt}\textbf{#1}}}}%
\renewcommand{\note}[2]{\edtext{}{\Cfootnote[nonum, nosep]{\fontsize{7.8pt}{9.2pt}\selectfont\espace{-15pt}\textbf{#1} #2}}}%

% Row macro for Spread 1 Left Page (Diplomatic-yet-corrected Ravennas + inline apparatus)
\newcommand{\ligneRav}[4]{%
  \noindent\hbox to 16.5cm{%
    \hbox to 0.55cm{\hfill\tiny\texttt{#1}\hspace{0.1cm}}%
    \hbox to 1.25cm{\hfill\scriptsize #2\hspace{0.15cm}}%
    \hbox to 8.1cm{\adjustbox{max width=8.0cm}{#3}\hfill}%
    \hspace{0.2cm}%
    \parbox[t]{6.4cm}{\raggedright\fontsize{6.7pt}{7.9pt}\selectfont #4}%
  }\par\vspace{0.35pt}%
}

% Row macro for Spread 2 in the Strict 4-Page-Per-Folio Edition (Established Greek / French)
\newcommand{\ligneEst}[3]{%
  \noindent\hbox to 16.5cm{%
    \hbox to 2.5cm{\hfill\scriptsize #1\hspace{0.2cm}}%
    \hbox to 13.2cm{\adjustbox{max width=13.1cm}{#2}\hfill}%
    \hbox to 0.8cm{\hfill\scriptsize\texttt{#3}}%
  }\par\vspace{0.3pt}%
}

\begin{document}
\thispagestyle{empty}
\noindent\begin{adjustbox}{max totalheight=19.5cm, max width=\textwidth, keepaspectratio}\begin{minipage}{\textwidth}\begin{center}
\vspace*{0.4cm}
{\Huge\textbf{ΑΡΙΣΤΟΦΑΝΟΥΣ ΘΕΣΜΟΦΟΡΙΑΖΟΥΣΑΙ}}\par\vspace{0.5cm}
{\LARGE\textbf{ARISTOPHANE --- LES THESMOPHORIES}}\par\vspace{0.8cm}
{\Large\textit{Nouvelle édition philologique et paléographique sur quatre pages}}\par\vspace{0.4cm}
{\large Tome II : Folios 16 à 30 (\texttt{ff. 170r–177r}, vv. 635–1231)}\par\vspace{0.5cm}
\hrule height 0.8pt\vspace{0.6cm}
{\normalsize
\textbf{Architecture des doubles pages (30 folios du \textit{Codex Ravennas} 429, \texttt{ff. 162v–177r}, vv. 1–1231) :}\par\vspace{0.3cm}
\begin{itemize}
  \item \textbf{Première double page — Page de gauche (pair) :} Transcription diplomatique corrigée du \textit{Codex Ravennas} 429 (\textbf{R}), avec \textit{notae personarum} tachygraphiques, \textit{paragraphoi}, \textit{dikola}, grattages (\textit{in rasura}), ajouts marginaux, les quatre \textit{parepigraphai} scéniques, soulignement pointillé des émendations (\dotuline{texte}), \textbf{apparat critique dans le prolongement de la ligne}, et en bas de page : (1) \textbf{\textit{loci paralleli}} et (2) \textbf{suite (en commentaire) de l'apparat critique détaillé}.
  \item \textbf{Première double page — Page de droite (impair) :} Fac-similé haute résolution en pleine page du folio correspondant du \textit{Codex Ravennas} 429 (Biblioteca Classense, Ravenne).
  \item \textbf{Seconde double page — Page de gauche (pair) :} Texte grec établi proprement, avec colométrie et attribution complète des répliques.
  \item \textbf{Seconde double page — Page de droite (impair) :} Traduction française en regard.
  \item \textbf{Bas de la seconde double page (gauche et droite) :} Long commentaire philologique, historique, métrique et scénique.
\end{itemize}
}\vspace{0.6cm}
\hrule height 0.8pt\vspace{0.5cm}
{\Large\textbf{Édition, traduction et commentaire par Louis HAUSEUX (\textit{Bibi})}}\par
\end{center}\end{minipage}\end{adjustbox}
\newpage

\clearpage\ifodd\value{page}\hbox{}\newpage\fi
\enTeteG{ΕΠΕΙΣΟΔΙΟΝ Β' --- Codex Ravennas 429, f. 170r (Classense 172r) --- vv. 635–678}
\begin{adjustwidth}{-1.5cm}{-3.0cm}
\begin{adjustbox}{max totalheight=21.0cm, max width=16.5cm, keepaspectratio}
\begin{minipage}{16.5cm}
{\fontsize{7.2pt}{8.5pt}\selectfont
\noindent\textbf{Folio 16 --- Codex Ravennas 429, \texttt{f. 170r} (\texttt{Classense 172r}) : ΕΠΕΙΣΟΔΙΟΝ Β', vv. 635–678}\hfill \textit{Témoins :} $\text{R}^1, \Rpc, \Rdeux$ (\Mon)\par\vspace{2pt}
\hrule height 0.4pt\vspace{2.5pt}
}
{\fontsize{8.1pt}{9.5pt}\selectfont
\ligneRav{635}{\auDessus{γυ}{ν}}{ὅδ᾽ ἐστὶν \difference{ἁνὴρ} ὃν λέγεις. : τί οὖν ποιῶ;}{\textbf{635} ἀνὴρ R : \lecon{ἁνὴρ} Bentley}
\ligneRav{}{\auDessus{γυ}{ν}}{ἀπόδυσον αὐτόν· οὐδὲν ὑγιὲς γὰρ λέγει·}{}
\ligneRav{637}{μ\auDessus{ν}{Η}}{κἄπειτ᾽ ἀποδύσετ᾽ ἐννέα παίδων μητέρα;}{\textbf{637-40} (Μνη.) R : \lecon{(Μνη.) (637-8) et (Μίκα) (639-40)} Bibi}
\ligneRav{638}{}{\difference{χαλᾷ} ταχέως τὸ στρόφιον ὦναίσχυντε σύ·}{\textbf{638} χάλᾳ R : χάλα Grynaeus : \lecon{χαλᾷ} Bibi}
\ligneRav{639}{\auDessus{γυ}{ν}}{ὡς καὶ στιβαρά τις φαίνεται καὶ καρτερά·}{\textbf{639-40} (Μνη.) 639 et (Γυνή) 640 R : \lecon{(Μίκα)} Bibi}
\ligneRav{640}{}{καὶ νὴ Δία \difference{τιτθούς} γ᾽ ὥσπερ ἡμεῖς οὐκ ἔχει·}{\textbf{640} τίτθος R : \lecon{τιτθούς} ed. Iuntina}
\ligneRav{641}{μ\auDessus{ν}{Η}}{\difference{στερίφη} γάρ εἰμι κοὐκ ἐκύησα πώποτε·}{\textbf{641} στεριφῇ R : \lecon{στερίφη} (sed στεριφὴ) \Mon}
\ligneRav{642}{κλεισθ/}{νῦν· τότε \difference{δὲ} μήτηρ ἦσθα παίδων ἐννέα·}{\textbf{642-3} (Μνησ.) (641-3) R : \lecon{(Κλεισθ.) (642-3)} Bibi\espApp δὴ R : \lecon{δὲ} \Mon}
\ligneRav{}{}{ἀνίστασ᾽ ὀρθός. ποῖ τὸ πέος ὠθεῖς κάτω;}{}
\ligneRav{644}{\auDessus{γυ}{ν}}{\difference{τοδὶ διέκυψε} καὶ μάλ᾽ εὔχρων ὦ τάλαν·}{\textbf{644} τὸ δὶ δὴ ἔκυψε R : \lecon{τοδὶ διέκυψε} Dobree, Dindorf}
\ligneRav{645}{κλεισθ/}{καὶ ποῦ 'στιν; : αὖθις εἰς τὸ πρόσθεν οἴχεται·}{}
\ligneRav{646}{κλεισθ/}{οὐκ ἐνγεταυθί. : ἀλλὰ δεῦρ᾽ ἥκει πάλιν·}{\textbf{646} \lecon{ἐνγεταυθί} R : ἐγγεταυθί Thiersch\espApp \lecon{ἀλλὰ} R : μἀλλὰ Bentley}
\ligneRav{647}{κλεισθ/}{ἰσθμόν τιν᾽ ἔχεις \difference{ὤνθρωπ᾽}· ἄνω τε καὶ κάτω}{\textbf{647} ῶνθρωπε R : ἄνθρωπε \Mon : ἄνθρωπ᾽ Suid. ι 640 : \lecon{ὤνθρωπ᾽} Ellebodius}
\ligneRav{}{}{τὸ πέος διέλκεις πυκνότερον Κορινθίων·}{}
\ligneRav{}{\auDessus{γυ}{ν}}{ὤ μιαρὸς οὗτος· ταῦτ᾽ ἄρ᾽ ὑπὲρ Εὐριπίδου}{}
\ligneRav{650}{—}{ἡμῖν ἐλοιδορεῖτο. : κακοδαίμων ἐγώ,}{}
\ligneRav{651}{}{εἰς οἷ᾽ ἐμαυτὸν \difference{εἰσεκύλισα} πράγματα·}{\textbf{651} εἰσεκύλησα R : \lecon{εἰσεκύλισα} Scaliger}
\ligneRav{}{\auDessus{γυ}{ν}}{ἄγε δὴ τί δρῶμεν; : τουτονὶ φυλάττετε}{}
\ligneRav{653}{}{καλῶς, ὅπως μὴ διαφυγὼν \difference{οἰχήσεται}·}{\textbf{653} οἴχεται R : \lecon{οἰχήσεται} Faber p.788}
\ligneRav{654}{}{ἐγὼ δὲ ταῦτα τοῖς \difference{πρυτάνεσιν} ἀγγελῶ·}{\textbf{654} πρυτάνεσσιν R : \lecon{-εσιν} \Mon\espApp πρυτάνεσσιν R : \lecon{πρυτάνεσιν} \Mon}
\ligneRav{655}{χορ/}{ἡμᾶς τοίνυν μετὰ τοῦτ᾽ ἤδη τὰς λαμπάδας ἁψαμένας χρὴ}{}
\ligneRav{}{}{ξυζωσαμένας εὖ κἀνδρείως τῶν θ᾽ ἱματίων ἀποδύσας}{}
\ligneRav{657}{}{ζητεῖν, εἴ που κἄλλος τις ἀνὴρ \difference{ἐπελήλυθε}, καὶ περιθρέξαι}{\textbf{657} εἰσελήλυθεν R : \lecon{ἐπελήλυθε} Handley \textbf{658} ἀθρῆσαι R : \lecon{διαθρῆσαι} Daubuz}
\ligneRav{658}{}{τὴν Πύκνα πᾶσαν καὶ τὰς σκηνὰς καὶ τὰς διόδους \difference{διαθρῆσαι}·}{\textbf{658} ἀθρῆσαι R : \lecon{διαθρῆσαι} Daubuz}
\ligneRav{659}{ἡμιχ/}{εἶα δὴ πρώτιστα μὲν χρὴ κοῦφον ἐξορμᾶν πόδα}{\textbf{659} (Ἡμιχ. γυν.) \Rpc}
\ligneRav{660}{}{καὶ διασκοπεῖν σιωπῇ πανταχῇ· μόνον δὲ χρὴ}{\textbf{660} πανταχῇ post χρὴ R : \lecon{ante μόνον} Küster}
\ligneRav{}{}{μὴ βραδύνειν, ὡς ὁ καιρός ἐστι μὴ μέλλειν ἔτι,}{}
\ligneRav{662}{}{ἀλλὰ τὴν πρώτην τρέχειν \difference{χρῆν} ὡς τάχιστ᾽ ἤδη κύκλῳ·}{\textbf{662} χρὴ R : \lecon{χρῆν} Bentley : χρῆ μ᾽ Austin}
\ligneRav{}{χορ/}{\hspace{1.1em}εἶά νυν ἴχνευε καὶ μά-}{}
\ligneRav{}{}{\hspace{1.1em} τευε ταχὺ πάντ᾽, εἴ τις ἐν τόποις ἑδραῖος}{}
\ligneRav{}{}{\hspace{1.1em}ἄλλος αὖ λέληθεν ὤν·}{}
\ligneRav{665}{}{\hspace{1.1em}πανταχῇ \difference{δὲ ῥῖψον} ὄμμα,}{\textbf{665} διάρριψον R : \lecon{δὲ ῥῖψον} Hermann}
\ligneRav{666}{}{\hspace{1.1em}καὶ τὰ τῇδε καὶ τὰ δεῦρο}{\textbf{666} post καὶ τὰ τῇδε add. καὶ τὰ κεῖσε Kaibel\espApp Post καὶ τὰ τῇδε add. καὶ τὰ κεῖσε Kaibel}
\ligneRav{666}{}{\hspace{1.1em}πάντ᾽ ἀνασκόπει καλῶς·}{\textbf{666} post καὶ τὰ τῇδε add. καὶ τὰ κεῖσε Kaibel\espApp Post καὶ τὰ τῇδε add. καὶ τὰ κεῖσε Kaibel}
\ligneRav{667}{}{\hspace{1.1em}ἢν γάρ ληφθῇ δράσας ἀνόσια,}{\textbf{667} μὴ λάθῃ R : \lecon{ληφθῇ} Reisig : με λάθῃ Bergk}
\ligneRav{}{}{\hspace{1.1em}δώσει τε δίκην καὶ πρὸς τούτῳ}{}
\ligneRav{669}{}{\hspace{1.1em}ἦν δὲ μὴ λάθῃ, δώσει δίκην, παράδειγμά τ᾽ \difference{ἀνδράσιν} (ἅπασιν \inRas{ἅπασιν}{}) ἔσται}{\textbf{669} ἅπασιν \inRas{ἅπασιν}{} ἔσται \Rpc : ἅπασιν ἔσται \Mon : \lecon{ἀνδράσιν ἔσται} Beer ap. Hermann (1845) : ἔσται ἅπασιν Bergk : alii alia}
\ligneRav{670}{}{\hspace{1.1em}παράδειγμ᾽ ὕβρεως}{}
\ligneRav{}{}{\hspace{1.1em}ἀδίκων τ᾽ ἔργων ἀθέων τε τρόπων·}{}
\ligneRav{}{}{\hspace{1.1em}φήσει δ᾽ εἶναί τε θεοὺς φανερῶς,}{}
\ligneRav{}{}{\hspace{1.1em}δείξει τ᾽ ἤδη}{}
\ligneRav{}{}{\hspace{1.1em}πᾶσιν ἀνθρώποις σεβίζειν δαίμονας}{}
\ligneRav{675}{}{\hspace{1.1em}†δικαίως τ᾽ ἐφέποντας † ὅσια καὶ νόμιμα}{\textbf{675} \lecon{ἐφέποντας} R : ἐφέπειν Hermann}
\ligneRav{}{}{\hspace{1.1em}μηδομένους ποιεῖν ὅ τι καλῶς ἔχει·}{}
\ligneRav{}{}{\hspace{1.1em}κἂν μὴ ποιῶσι ταῦτα τοιάδ᾽ ἔσται·}{}
}
\vspace{3pt}\noindent\rule{4.5cm}{0.4pt}\par\vspace{1.5pt}
{\fontsize{7.0pt}{8.3pt}\selectfont
\noindent\textbf{Loci paralleli (vv. 635–678) ---} \textbf{641} (στερίφη) Hesychius σ 1769 στερίφη· στεῖρα, μὴ τεκοῦσα, μηδὲ τίκτουσα $=$ Photius σ 532 ; cf. Moeris σ 46 (στερίφη $=$ στερρά) \espTest \textbf{641} (στερίφη) Hesychius σ 1769 στερίφη $=$ Photius σ 532 ; cf. Moeris σ 46 (στερίφη $=$ στερρά) \espTest \textbf{647-8} Suidas ι 640 \espTest \textbf{659} (εἶα δὴ) Hesychius ε 8 ἔα δή· ἄγε δή $=$ Photius ε 1 $=$ \Syn ε 1 $=$ Suidas ε 3\par
}
\vspace{2.5pt}\noindent\rule{4.5cm}{0.4pt}\par\vspace{1.5pt}
{\fontsize{6.9pt}{8.2pt}\selectfont
\noindent\textbf{Suite (en commentaire) de l'apparat critique \& notes paléographiques (\texttt{f. 170r}) ---} \textbf{Sigle tachygraphique marginal μν\auDessus{η}{σ} au v. 638 et redistribution des répliques par \textbf{Bibi} (vv. 637–643) :} Comme le précise Louis Hauseux (\texttt{vv 1-641.tex}, l. 4618), R comporte en marge du v. 638 le sigle tachygraphique \lecon{μν\auDessus{η}{σ}} et aucune indication de changement d'interlocuteur jusqu'au v. 640 (Γυνή). Aux vv. 637–640, Louis Hauseux donne 637–638 à (Μνη.) (en lisant l'indicatif 3\ieme{} pers. \lecon{χαλᾷ} au v. 638) et 639–640 à \lecon{(Μίκα)} ; aux vv. 642–643, il restitue la réplique à \lecon{(Κλεισθ.)}. --- \textbf{Confusion onciale ΙϹ $\leftrightarrow$ Π au v. 657 :} εἰσελήλυθεν R corrigé par \textbf{Handley} en \lecon{ἐπελήλυθε}. --- \textbf{Grattage d'un mot entier répété par dittographie au v. 669 :} \Rac{} avait écrit deux fois de suite ἅπασιν ; \Rpc{} a gratté le second (\lecon{ἅπασιν \inRas{ἅπασιν}{} ἔσται}), laissant sur le parchemin un espace blanc très net entre ἅπασιν et ἔσται (Beer ap. Hermann restitue ἀνδράσιν ἔσται). --- \textbf{vv. 637–640 :} μν\auDessus{η}{σ} (ad v. 638, nulla mut. pers. usque ad 640 Γυνή) R : \lecon{(Μνη.) (637–638) et (Μίκα) (639–640)} \lecon{Bibi}. --- \textbf{v. 638 :} χάλᾳ R (χάλα Grynaeus) : \lecon{χαλᾷ} \lecon{Bibi}. --- \textbf{vv. 642–643 :} (Μνησ.) (641–643) R : \lecon{(Κλεισθ.) (642–643)} \lecon{Bibi} \espApp \textbf{v. 642 :} δὴ R : \lecon{δὲ} \Mon{}. --- \textbf{v. 644 :} τὸ δὶ δὴ ἔκυψε R (médivision et iotacisme) : \lecon{τοδὶ διέκυψε} Dobree, Dindorf. --- \textbf{v. 651 :} εἰσεκύλησα R (iotacisme) : \lecon{εἰσεκύλισα} Scaliger. --- \textbf{v. 657 :} εἰσελήλυθεν R : \lecon{ἐπελήλυθε} Handley (confusion onciale ΙϹ/Π).\par
}
\end{minipage}
\end{adjustbox}
\end{adjustwidth}
\newpage
\enTeteD{Codex Ravennas 429 (Biblioteca Classense) --- Fac-similé du f. 170r (vv. 635–678)}
\begin{adjustwidth}{-1.5cm}{-3.0cm}
\begin{center}
\includegraphics[width=16.5cm,height=20.5cm,keepaspectratio]{/workspaces/Thesmophories/dossier_ravennas_thesmophories/folios_individuels_ff162v-177r/16_f170r_classense172r_vv0635-0678.jpg}\par
\vspace{2pt}
{\scriptsize\textbf{Codex Ravennas 429, f. 170r (Classense 172r)} --- ΕΠΕΙΣΟΔΙΟΝ Β', \textbf{vv. 635–678} ($\text{R}^1, \Rpc, \Rdeux$ (\Mon)).}
\end{center}
\end{adjustwidth}
\newpage
\clearpage\ifodd\value{page}\hbox{}\newpage\fi
\enTeteG{ΘΕΣΜΟΦΟΡΙΑΖΟΥΣΑΙ --- Texte grec établi (f. 170r, vv. 635–678)}
\enTeteD{LES THESMOPHORIES --- Traduction française (f. 170r, vv. 635–678)}
\begin{pages}
\begin{Leftside}
\beginnumbering
%% v. 635+ (part 2)
\begin{groupe}
\trimIamb{\ligne{635}ὅδ᾽ ἐστὶν ἁνὴρ ὃν λέγεις.%\apparat{635}{ἀνὴρ R : \lecon{ἁνὴρ} Bentley}
\\\passeLigne}%
\end{groupe}%
%% v. 635b
\begin{groupe}%
\personnageGrec{Κλεισθένης}%
\trimIamb{\ligne{635}\decalage{ὅδ᾽ ἐστὶν ἁνὴρ ὃν λέγεις.}τί οὖν ποιῶ;\\\passeLigne}%
\end{groupe}%
%% v. 636
\begin{groupe}%
\personnageGrec{Μίκα}%
\trimIamb{ἀπόδυσον αὐτόν· οὐδὲν ὑγιὲς γὰρ λέγει.\\\passeLigne}%
\end{groupe}%
%% v. 637-8
\begin{groupe}%
\personnageGrec{Μνησίλοχος}%
\trimIamb{κἄπειτ᾽ ἀποδύσετ᾽ ἐννέα παίδων μητέρα;\apparat{637-40}{(Μνη.) R : %(Μνη.) (637), (Κλεισθ.) (638) et (Γυν.) (639-40) Brunck, edd. post. : 
\lecon{(Μνη.) (637-8) et (Μίκα) (639-40)} Bibi \ligneApp{638} χάλᾳ R : χάλα Grynaeus : \lecon{χαλᾷ} Bibi%\espApp χάλα.... στρόφιον Micae, ἀναίσχυντε σύ Affini tribuit Degani
}%
\\
χαλᾷ ταχέως τὸ στρόφιον ὦναίσχυντε σύ.%
\commentaire{637-40}{κἄπειτ᾽ ἀποδύσετ᾽ ἐννέα παίδων μητέρα; ... καὶ νὴ Δία τίτθος γ᾽ ὥσπερ ἡμεῖς οὐκ ἔχει.}{%Les tripatouillomanes ont encore sévi à ce vers !.. ὦναίσχυντοι, \italique{quousque tandem...} ô Willems, combien la tâche que tu t'étais confié est loin d'être terminée !\\
%Le vers coulait pourtant d'une source claire : Clisthène est en train de déshabiller Mnésiloque ; ce dernier feint donc l'indignation (χαλᾷ est un subjonctif dit \italique{délibératif}) devant l'empressement (ταχέως) que cet homme -- car Clisthène reste un homme -- met à enlever le \italique{soutien gorge} d'une femme -- car Mnésiloque continue de jouer une femme (vertueuse même, ici, mais un peu tard). Qu'y a-t-il de plus limpide ? l'apostrophe ὦναίσχυντε σύ conclue logiquement la tirade de mère effarouchée qu'affecte le peu crédible Mnésiloque (et qui fait d'ailleurs écho au précédent ἀναίσχυντός τις εἶ lancé par le même Mnésiloque au vers 611\marqueur{1}).\\
%Qu'ont eu besoin les éditeurs, tous, de suivre Brunck\marqueur{2} ? et d'inventer un changement d'interlocuteurs là où il n'y avait pas lieu d'être -- alors que le \italique{Ravennas} est on ne peut plus fiable sur ce point -- et, \italique{mutatis mutandis}, de surajouter une autre retouche devenue nécessaire : χάλα à la place de χαλᾷ.
Le manuscrit de Ravenne porte en marge du vers 638 μν\auDessus{η}{σ}, indiquant ainsi que Mnésiloque prend la parole ; le manuscrit ne présente alors plus aucun signe de changement de locuteur jusqu'au vers 641\marqueur{1}. Or Mnésiloque %s'il est sûr que Mnésiloque prononce le vers 637, ἀποδύσετ᾽ ἐννέα παίδων μητέρα, il 
ne peut aucunement prononcer le vers 639, commentaire sur sa poitrine étonnamment plate (ὡς καὶ στιβαρά τις φαίνεται καὶ καρτερά).
Il y a donc un \italique{paragraphos} à rétablir.
Je dis \italique{un} car les éditeurs suivent quasi-unanimement Brunck en `rétablissent' deux : au début du vers 638 (Clistène) et au début du vers 639 (la femme -- Mika -- reprend la parole).
Une autre solution me semble plus légère : laisser à Mnésiloque les deux vers 637-638 et à Mika les vers 639-640.
Cette retouche comporte surtout l'avantage d'être bien meilleure quant à l'action et au jeu de scène.\\
Presque tous les éditeurs depuis Brunck, disais-je, retirent le vers 638 à Mnésiloque pour l'accorder à Clisthène. Double erreur ! \ \ \ \italique{Primo}, parce qu'ils sont alors contraints de surajouter une seconde émendation : transformer le χάλᾳ du manuscrit en un χάλα (impératif lancé par Clisthène) alors que le iota adscrit est nettement écrit. \ \ \ \italique{Secundo}, parce que Mika, au vers 636, avait déjà ordonné à Clisthène de déshabiller Mnésiloque (ἀπόδυσον αὐτόν) ; il n'y a aucune de raison de douter qu'il ne soit en train de lui obéir. Un peu d'action, que diable ! une pièce d'Aristophane ne doit pas avoir la même pesanteur que l'administration française. (Selon nos éditeurs, Mika ordonnerait à Clisthène de déshabiller Mnésiloque qui transmettrait alors l'ordre à l'intéressé... qui finirait par ne rien faire ?..)\\
Le vers 638, χαλᾷ ταχέως τὸ στρόφιον ὦναίσχυντε σύ, était pourtant excellent dans la bouche de Mnésiloque : Clisthène est en train de le déshabiller, il feint l'indignation (χαλᾷ est un subjonctif dit \italique{délibératif}) devant l'empressement (ταχέως) que cet homme -- car Clisthène reste un homme -- met à enlever le \italique{soutien gorge} d'une femme. « Il défait mon strophion avec promptitude. Impudent ! » Mnésiloque continue de jouer une femme -- vertueuse même, mais un peu tard. Qu'y a-t-il de plus limpide ? l'apostrophe ὦναίσχυντε σύ conclue logiquement la tirade de mère effarouchée qu'affecte le peu crédible Mnésiloque (et qui fait d'ailleurs écho au précédent ἀναίσχυντός τις εἶ (v. 611) lancé par le même Mnésiloque au même Clisthène\marqueur{2}).
}%
\debutNote{637-40}%
%\note{1}{D'aucuns avaient même prétexté de ce double ἀναίσχυντος pour attribuer ἀναίσχυντός τις εἶ à Clisthène, histoire que leur texte présente une homogénéité certaine : foutaises partout, justesse nulle part.}%
%\note{2}{Quelques uns, croyant faire preuve d'originalité, attribuèrent la réplique à Mika plutôt qu'à Clisthène. Mais tous tripatouillent !.. tous !!! La pire espèce étant celle qui tente de camoufler son forfait en le faisant disparaître de l'apparat critique.}%
\note{1}{Il est d'usage dans notre manuscrit que le \italique{signe} soit un παράγραφος semblable à un tiret cadratin ` --- ' lorsque le changement de locuteur se fait en début de vers. Si le changement intervient en milieu de vers, il y a plus souvent un \latin{punctum duplex} ` : ' ou bien une abréviation indiquant le nom du personnage qui prend la parole. Par abus de langage, nous utiliserons le terme de \italique{paragraphos} pour désigner indifféremment l'un de ces signes.}%
\note{2}{Quelques éditeurs avaient même prétexté de ce double ἀναίσχυντος pour attribuer l'ἀναίσχυντός τις εἶ du vers 611 aussi à Clisthène, histoire que leurs multiples tripatouillages présentent une apparence d'homogénéité...}%
\\\passeLigne}%
\end{groupe}%
%% v. 639-640
\begin{groupe}%
\personnageGrec{Μίκα}%
\trimIamb{ὡς καὶ στιβαρά τις φαίνεται καὶ καρτερά·\\%\apparat{639-40}{(Μνη.) 639 et (Γυνή) 640 R : \lecon{(Μίκα)} Bibi}\\
καὶ νὴ Δία τιτθούς γ᾽ ὥσπερ ἡμεῖς οὐκ ἔχει.\apparat{640}{τίτθος R : \lecon{τιτθούς} ed. Iuntina}\\\passeLigne}%
\end{groupe}%
%% v. 641
\begin{groupe}%
\personnageGrec{Μνησίλοχος}%
\trimIamb{στερίφη γάρ εἰμι κοὐκ ἐκύησα πώποτε.%\testimonia{641}{(στερίφη) Hesychius σ 1769 στερίφη· στεῖρα, μὴ τεκοῦσα, μηδὲ τίκτουσα $=$ Photius σ 532 ; cf. Moeris σ 46 (στερίφη $=$ στερρά)}
\apparat{641}{στεριφῇ R : \lecon{στερίφη} (sed στεριφὴ) \Mon}%
\commentaire{641}{στερίφη γάρ εἰμι.}{Jeu de mots d'Aristophane : comment Mnésiloque explique-t-il la platitude de sa poitrine ? tout simplement parce qu'il a les chairs \italique{fermes} (sens premier de στέριφος, le Scythe l'emploie avec cette acception au vers 1185). Mais στέριφος peut aussi avoir le sens dérivé de \italique{stérile} ; sens que le spectateur ne privilégiera qu'en entendant la seconde partie de la réplique de Mnésiloque (κοὐκ ἐκύησα). C'est la raison pour laquelle Aristophane recourt à un mot rare, στέριφος. (Le scholiaste de noter : ἀντὶ τοῦ στείρα, commentaire qui se retrouve dans de nombreux lexiques, cf. l'apparat des témoins\marqueur{1}.)%\\
%Y a-t-il ici aussi détournement d'un tour proverbial ? toujours est-il qu'on retrouve un tour identique chez Hipponax (fr. 122 ed. Degani) :%
%\citComm{ἀμφιδέξιος γάρ εἰμι κοὐκ ἁμαρτάνω κόπτων.}%
}%
\debutNote{641}%
\note{1}{Et pour l'apparat le plus fourni, se reporter à celui de l'édition Theodoridis de Photius.}%
\\\passeLigne}%
\end{groupe}%
%% v. 642
\begin{groupe}%
\personnageGrec{Κλεισθένης}%
\trimIamb{\ligne{642}νῦν· τότε δὲ μήτηρ ἦσθα παίδων ἐννέα.\apparat{642-3}{(Μνησ.) (641-3) %nulla nota mutatae personae 
R : %(Γυνή) (642) et (Κλεισθ.) (643) Ellebodius : 
\lecon{(Κλεισθ.) (642-3)} Bibi \ligneApp{642} δὴ R : \lecon{δὲ} \Mon}%
\commentaire{642-3}{νῦν· τότε δὲ μήτηρ ἦσθα παίδων ἐννέα. / ἀνίστασ᾽ ὀρθός. ποῖ τὸ πέος ὠθεῖς κάτω;}{Il n'y a aucun \italique{paragraphos} aux vers 642 et 643.
Pourtant, ce n'est assurément pas Mnésiloque (qui a la parole au vers précédent) qui peut prononcer ces vers. On ajoute souvent deux changement d'interlocuteurs : Mika pour le vers 642 et Clisthène pour le vers 643. Par simplicité, je préfère n'en ajouter qu'un et attribue donc le distique entier à Clisthène.}%
\commentaire{643}{ἀνίστασ᾽ ὀρθός.}{«~Redresse-toi !~» L'injonction adressée à Mnésiloque est bien-sûr à double entente, et Clisthène se fera doublement obéir. Mais interrogeons-nous : pourquoi Mnésiloque s'était-il assis ou recroquevillé ? réponse simple : pour cacher sa virilité. % (nous y reviendrons à la fin de cette note). 
Ils se redressent. Mnésiloque tente une dernière fois de cacher ce qui doit être caché et Clisthène : ποῖ τὸ πέος ὠθεῖς κάτω; «~Où repousses-tu ce \italique{péos} ?!~» Le \italique{péos} passe alors au verso de Mnésiloque, d'où peut le contempler Mika, admiratrice. Avant que ne s'ensuive la scène burlesque où Clisthène paraît fort intéressé aussi.\\
Mnésiloque s'était habillé dans cet ordre : κροκωτός, στρόφιον et ἔγκυκλον. Le στρόφιον et l'ἔγκυκλον viennent d'être retirés, mais pas le κροκωτός (au vers 941, Mnésiloque le porte toujours) ; comment se fait-il que son \italique{olisbos} soit apparent malgré la présence de cette robe safranée ? parce que la crocote est composée d'un tissu très fin qui la rend transparente\marqueur{1}}%
\debutNote{643}%
\note{1}{Cf. le cratère à figures rouges (ca. -380) représentant une scène de comédie : un vieillard entouré de deux gardes et vêtu de la crocote... qui laisse tout transparaître ; il est reproduit en annexe. Cratère «~St. Agata Antigone~», St. Agata dei Goti I, collection Rainone-Mustilli. Inventorié par Arthur D. Trendall , \italique{Phlyax Vases}, «~BICS~» suppl. LXIX (1967).}%
\\
ἀνίστασ᾽ ὀρθός. ποῖ τὸ πέος ὠθεῖς κάτω;%
\commentaire{643}{Note sur τὸ ἔγκυκλον $=$ \italique{une jupe}.}{Clisthène vient de retirer le στρόφιον de Mnésiloque, ce qui entraîna quelques remarques sur sa poitrine. Or, lors de la scène d'habillement (vers 253-261), Mnésiloque avait successivement enfilé le κροκωτός, le στρόφιον \italique{puis, à la fin,} l'ἔγκυκλον. Il résulte donc que l'ἔγκυκλον \italique{est un vêtement qui ne peut pas recouvrir la poitrine}. (Sinon Clisthène aurait dû commencer par l'encycle avant le στρόφιον.) 
Ce ne peut donc être en aucun cas un \italien{mantello ampio, pesante}\marqueur{1}.
%Et puisque, jusqu'ici, seules d'ineptes sottises ont été écrites à propos de l'ἔγκυκλον, il nous a paru profitable 
Allongeons quelque peu cette note afin que soit irréfragable l'identification de l'ἔγκυκλον avec une sorte de jupe courte.\\
%Ce vêtement qui se porte plutôt à l'intérieur aux dires d'Eustathe (voir ci-dessous).
%Ce vêtement que l'étymologie, ἔγ-κυκλον, désigne comme ce qui \italique{en-toure}, ce qui \italique{en-ceint}.
Que pourrait-il être d'autre ce vêtement qui sitôt retiré laisse tout émoustillé le mari de Praxagora ; voyez plutôt (\italique{Assemblée des femmes}, vers 535-538) :%
\citComm{Βλέπυρος -- εἶτ' οὐ τὸ σαυτῆς ἱμάτιον ἐχρῆν σ' ἔχειν;\\
ἀλλ' ἔμ' ἀποδύσασ' ἐπιβαλοῦσα \textbf{τοὔγκυκλον}\\
ᾤχου καταλιποῦσ' ὡσπερεὶ προκείμενον,\\
μόνον οὐ στεφανώσασ' οὐδ' ἐπιθεῖσα λήκυθον.\\
Alors que n'as-tu pris ton propre manteau [pour sortir] ? Mais tu me dévêts et jettes sur moi \italique{ton encycle}, puis tu t'en vas, me laissant là comme un cadavre exposé, hors que tu as omis de me couronner et de placer près de moi un lécythe.}%
Le champ lexical est celui d'une exposition funèbre\marqueur{2} comme la traduction de Willems le laisse entendre, mais le λήκυθον du vers 538 doit être pris \italique{sensu obsceno} selon la matrice métaphorique : λήκυθος $=$ ἀλάβαστος\marqueur{3} $=$ \latin{phallus}.
Blépyros reproche à sa femme de ne pas s'être acquitté du \italique{nome de mâtines} (ὄρθριος νόμος\marqueur{4}).
Ce vêtement permet à Aristophane un jeu de mot avec l'\italique{enceinte} de la ville fortifiée de Pylos\marqueur{5} -- oui ! --, la Πύλος, cette ville qui déclenche intarissablement la verve éroticomique de notre poète, jonglant avec l'autre matrice métaphorique (féminine cette fois) : πύλος $=$ πύλη $=$ θύρα $=$ \latin{pudenda muliebria} (raison pour laquelle l'Athénien ne veut pas la céder). 
%Excusez-moi des longueurs ; je me vois contraint d'expliciter quelque peu Aristophane à ses propres commentateurs devant toutes les conneries qu'ils débitent -- littéralement : ils pensent que l'\italique{encycle} désigne le \italique{con}. 
%Cela dit, revenons à \italique{Lysistrata} : 
Il y avait certainement à l'époque de la représentation de \italique{Lysistrata} -- sans doute d'une année antérieure à celle des \italique{Thesmophories} -- des tractations pour que la paix fût conclue en échange de la remise de Pylos ou, par synecdoque, de son \italique{encycle}. 
Ce n'est pas autrement qu'il faut comprendre le passage au début de la même pièce (vers 111-114) : Lysistrata demande : %
\citComm{ἐθέλοιτ᾽ ἂν οὖν, εἰ μηχανὴν εὕροιμ᾽ ἐγώ,\\
μετ᾽ ἐμοῦ καταλῦσαι τὸν πόλεμον;\\
Consentirez-vous donc, si j'en trouve le moyen, à mettre fin avec moi à la guerre ?}%
et Aristophane de répondre, par la voix d'une femme :%
\citComm{\decalage{μετ᾽ ἐμοῦ καταλῦσαι τὸν πόλεμον;}νὴ τὼ θεώ\\
ἔγωγ᾽ ἂν οὖν κἂν εἴ με χρείη \textbf{τοὔγκυκλον}\\
τουτὶ καταθεῖσαν ἐκπιεῖν αὐθημερόν.\\
Oui, par les deux déesses. J'en suis, quand il me faudrait déposer l'\italique{encycle} que voici en gage pour boire dès aujourd'hui [à la paix : les traités, σπονδαί, sont en grec, littéralement, des \italique{libations}].}%
Il faut donc bien-sûr aussi attribuer à cette même femme -- et comme ne le font pas les éditeurs d'Aristophane -- le vers 104 :%
\citComm{ὁ δ᾽ ἐμός γε τελέους ἑπτὰ μῆνας ἐν Πύλῳ.}%
(Elle se dit prête à faire tomber l'\italique{encycle} dès le retour de son mari en garnison depuis sept longs mois à Pylos.)
Aristophane, par deux fois dans la même pièce, exhorte ses contemporains à renoncer à Pylos, c'est dire l'importance qu'il accordait à l'éventualité de la paix ! On connaît la suite...
Ce vêtement, donc, qui enclot le sadinet, puisqu'il faut parler net... Alors ? que peut-il bien être ce vêtement ? S'il vous plaît, ne me répondez pas de chœur avec tous les fonctionnaires de la recherche : «~un lourd et pesant manteau~».\\
Mais délaissons un peu nos primes amours, les comédies d'Aristophane, pour aborder un terrain plus aride et cependant qui n'est certainement pas à négliger, je veux parler de la lexicologie.
Ce sens de \italique{jupe} pour l'ἔγκυκλον est confirmé, ce me semble, par le commentaire le plus précis : celui du scholiaste de \italique{Lysistrata} (vers 113), qui glose ἔγκυκλον ainsi :%
\citComm{τὸ περιβόλαιον τὸ ζωνάριον.}%
J'interprète ζωνάριον comme un adjectif se rapportant à περιβόλαιον\marqueur{6} : \italique{le vêtement / la couverture}, celle qui enceint, \italique{au niveau de la taille} ; une jupe donc. Citons aussi ce passage du copieux commentaire à l'\italique{Iliade} (976.13) d'Eustathe de Thessalonique (il commente \italique{Iliade} XIV, v. 178) :%
\citComm{ἰστέον δὲ ὅτι, καθάπερ ὁ ἑανὸς οὕτω καὶ τὸ ἔκκυκλον ἢ ἔγκυκλον, οὗ καὶ ὁ κωμικὸς μέμνηται, γυναικεῖον ἦν ποτε φόρημα, φησὶ γοῦν Παυσανίας ὅτι ἔγκυκλον περιπόρφυρον\marqueur{7} ἱμάτιον καὶ χιτὼν γυναικεῖος, ὃν ἔνδοθεν ἐνδύονται γυναῖκες, εἶτα τὸ ἔνδυμα. λέγεται δὲ φησίν [sic, je lis : σφῇσιν] ὑπ' ἐνίων ἐπιγονατίς, διότι μέχρι γονάτων διήκει.\\
Il faut savoir que, de même que l'ἑανὸς [= riche robe de femme dans le passage d'Homère qu'Eustathe commente], il y a aussi l'ἔκκυκλον (ou ἔγκυκλον), dont s'est aussi souvenu le Comique [= Aristophane]. C'était un vêtement (φόρημα) féminin porté jadis. Ce qui est sûr, c'est que Pausanias indique que l'ἔγκυκλον était un vêtement (ἱμάτιον) aux franges pourpres (περιπόρφυρον) et une tunique (χιτὼν) féminine, que les femmes revêtaient à l'intérieur, et ensuite un habit (ἔνδυμα). Il est qualifié par elles-mêmes, d'après certains, d'ἐπιγονατίς, parce qu'il atteignait les genoux.}%
Ce vêtement érotique d'intérieur peut-il être autre chose qu'une \italique{jupe sexy} ?%
}%
\debutNote{643}%
\note{1}{%Et certainement pas un \italien{mantello ampio, pesante,} 
Comme le conclue Franca Perusino dans son article «~\italien{L'ἔγκυκλον, un mantello femminile nelle commedie di Aristofane}~», \italique{Quaderni Urbinati di Cultura Classica}, 72.3 (2002), pp. 129-133. %Je ne cite cet article que par pur scrupule bibliographique : il présente tous les travers philologiques rencontrés chez les commentateurs d'Aristophane. 
Ses conclusions erronées ont malheureusement été reproduites dans la plupart des travaux portant sur le costume de l'Ancienne comédie... espérons que cette note rectificatrice mettra enfin un terme à la contagion ! %Ne soyons pas trop injuste et signalons la belle trouvaille de cet article : une probable identification d'un ἔγκυκλον sur un vase à figures rouges de la première moitié du \nieme{V} siècle. Ce vase, reproduit en annexe, montre une femme nue, un χιτών à la main, et, à côté d'elle, des bottes ainsi qu'un tissu à franges (περιπόρφυρον, cf. le commentaire d'Eustathe) posé sur un tabouret... notre ἔγκυκλον ? C'est bien ce qu'il m'en semble ! Chacun prêchant pour sa propre paroisse, j'y vois effectivement une \italique{jupe} et non un \italien{mantello ampio e pesante}. Je laisse le lecteur juge de l'iconographie ; pour ce qui est du sens d'ἔγκυκλον chez Aristophane, le doute n'est pas possible. Ce vase est le n° 204020 du \italique{Corpus Vasorum Antiquorum}, Beazley l'a recensé (378.122) dans son monumental \italique{Attic Red-Figure Vase-Painters}\expo{2} (1963).
}%
\note{2}{Cf. par exemple \italique{Lysistrata}, vv. 599 sqq.}%
\note{3}{Cf. la saillie du chaud-lapin Kinésias soumis à la torture par Myrrhine (\italique{Lysistrata}, v. 947) :%
\citComm{\textbf{Μυρρίνη} -- λαβὲ τόνδε τὸν ἀλάβαστον. --- \textbf{Κινησίας} -- ἀλλ᾽ ἕτερον ἔχω.}%
}%
\note{4}{\italique{Assemblée des femmes}, vers 741. Par jeu de mots avec l'ὄρθιος νόμος, le \italique{mode orthien}, qu'Hésychius (ο 1188) n'a pas manqué de nous définir :%
\citComm{ὄρθιον νόμον· οὕτως ὠνομάσθαι φασὶν αὐτὸν ἀπὸ τοῦ ῥυθμοῦ, καθάπερ καὶ τὸν τροχαῖον.}%
(Cf. aussi sur le mode orthien \italique{Cavaliers} 1280 ; \italique{Acharniens} 16.)\\
Les allusions d'Aristophane à cette tradition matinale à laquelle sacrifient ses personnages sont légions ; ne citons que \italique{Lysistrata} v. 966 puisque le passage est encore sous nos yeux.}%
\note{5}{\italique{Lysistrata}, vers 1162-1163. Pour conclure la paix, l'ambassadeur laconien demande à récupérer l'\italique{encyle} :%
\citComm{Λάκων -- ἁμές γε λῶμες, αἴ τις ἁμὶν τὤγκυκλον\\
λῇ τοῦτ᾽ ἀποδόμεν. --- Λυσιστράτη -- ποῖον ὦ τᾶν; --- Λάκων -- τὰν Πύλον,\\
ἇσπερ πάλαι δεόμεθα καὶ βλιμάττομες.}%
}%
\note{6}{Chantraine, page 74 de sa \italique{Formation des noms en grec ancien}, précise : «~la finale -άριον qui a fourni essentiellement des diminutifs s'observe parfois dans des adjectifs, cf. σωμα παιδαριον dans des inscription de Delphes~».}%
%\note{7}{La faussaire précitée avait mutilé ce texte, et ce afin de le tordre à ses conclusions -- selon la pente bien connue de ces gens-là. On constate qu'il n'était nul besoin de le charcuter : une minuscule émendation (ajouter un sigma à φησίν et lire σφῇσιν) suffit à le rendre intelligible.}%
\note{7}{On lit en épigraphie : ἔνκυκλον περιποίκιλο, ἔγκυκλον ποικίλον. À l'inverse, on croise aussi un ἔνκυκλον λευκ[όν].}%
%\note{9}{Cf. le cratère à figures rouges (ca. -380) représentant une scène de comédie : un vieillard entouré de deux gardes et vêtu de la crocote... qui laisse tout transparaître ; il est reproduit en annexe. Cratère «~St. Agata Antigone~», St. Agata dei Goti I, collection Rainone-Mustilli. Inventorié par Arthur D. Trendall , \italique{Phlyax Vases}, «~BICS~» suppl. LXIX (1967).}%
\\\passeLigne}%
\end{groupe}%
%% v. 644
\begin{groupe}%
\personnageGrec{Μίκα}%
\trimIamb{τοδὶ διέκυψε καὶ μάλ᾽ εὔχρων ὦ τάλαν.\apparat{644}{τὸ δὶ δὴ ἔκυψε R : \lecon{τοδὶ διέκυψε} Dobree, Dindorf}%
\commentaire{644}{τάλαν.}{Ce mot, le vocatif de τάλας, tant masculin que féminin, n'est mis par Aristophane que dans la bouche de femmes\marqueur{1}.}%
\debutNote{644}%
\note{1}{\italique{Lysistrata}, vers 102, 910 et 914 ; \italique{Grenouilles}, vers 559 ; \italique{Assemblée des femmes}, vers 124, 658 et 1005 ; \italique{Ploutos}, vers 706 et 1055.}\\\passeLigne}%
\end{groupe}%
%% v. 645a
\begin{groupe}%
\personnageGrec{Κλεισθένης}%
\trimIamb{καὶ ποῦ 'στιν;\\\passeLigne}%
\end{groupe}%
%% v. 645b
\begin{groupe}%
\personnageGrec{Μίκα}%
\trimIamb{\ligne{645}\decalage{καὶ ποῦ 'στιν;}αὖθις εἰς τὸ πρόσθεν οἴχεται.\\\passeLigne}%
\end{groupe}%
%% v. 646a
\begin{groupe}%
\personnageGrec{Κλεισθένης}%
\trimIamb{οὐκ ἐνγεταυθί.\apparat{646}{\lecon{ἐνγεταυθί} R : ἐγγεταυθί Thiersch}\\\passeLigne}%
\end{groupe}%
%% v. 646b
\begin{groupe}%
\personnageGrec{Μίκα}%
\trimIamb{\ligne{646}\decalage{οὐκ ἐνγεταυθί.}ἀλλὰ δεῦρ᾽ ἥκει πάλιν.\apparat{646}{\lecon{ἀλλὰ} R : μἀλλὰ Bentley}\\\passeLigne}%
\end{groupe}%
%% v. 647-648
\begin{groupe}%
\personnageGrec{Κλεισθένης}%
\trimIamb{ἰσθμόν τιν᾽ ἔχεις ὤνθρωπ᾽· ἄνω τε καὶ κάτω\testimonia{647-8}{Suidas ι 640}\apparat{647}{ῶνθρωπε R : ἄνθρωπε \Mon : ἄνθρωπ᾽ Suid. ι 640 : \lecon{ὤνθρωπ᾽} Ellebodius}%
\commentaire{647-8}{ἰσθμόν τιν᾽ ἔχεις ὦνθρωπ᾽· ἄνω τε καὶ κάτω / τὸ πέος διέλκεις πυκνότερον Κορινθίων.}{Clisthène fait référence au \grec{δίολκος} : un chemin dallé de 8 kilomètres qui permettait de faire traverser aux bateaux l'isthme de Corinthe grâce à une sorte de rail rudimentaire et d'éviter ainsi le long et périlleux contournement du Péloponnèse\marqueur{1}. Pendant la scène burlesque de déshabillage qui vient d'avoir lieu, Mnésiloque, pour cacher sa virilité – et certainement aussi émoustillé d'être effeuillé de la sorte –, a actionné à plusieurs reprises le mécanisme permettant de lever et baisser son \italique{phallus}. Admiratif devant de telles performances, Clisthène le flatte d'être encore mieux outillé que les machines qui permettaient le transbordement des bateaux.\\
Voilà pour l'image `τὸ πέος διέλκεις πυκνότερον Κορινθίων' ! il reste à expliquer la première partie de la réplique de Clisthène : ἰσθμόν τιν᾽ ἔχεις, «~tu possèdes un [drôle de] genre d'\italique{isthme}~».
ἰσθμός désigne proprement un \italique{passage étroit} et, par analogie, le \italique{cou}. Le spectateur athénien saisissait instamment l'image : par \italique{isthme} il fallait entendre \italique{pénis} grâce à la matrice métaphorique suivante : ἰσθμός = αὐχήν\marqueur{2} = τράχηλος\marqueur{3} = πέος.}%
\debutNote{647-8}%
\note{1}{Voir à ce sujet l'article de M. J. T. Lewis, « Railways in the Greek and Roman world », in \italique{Early Railways : a Selection of Papers from the First International Early Railways Conference}, 2001, p. 8-19. Thucydide (III, 15 ; VIII, 7) évoque le δίολκος dans sa \italique{Guerre du Péloponnèse}.}%
\note{2}{αὐχήν a le sens de πέος par exemple au vers 681 de \italique{Lysistrata} : le chœur des vieillards a maille à partir avec l'engeance féminine ; après avoir rappelé par une métaphore hippique combien la femme est portée à la luxure, il en vient à la conclusion que le meilleur moyen pour dominer ces fougueuses juments est de les chevaucher -- la fameuse \italique{phallocratie} ? Il formule cette idée à l'aide d'une image carcérale : il faut assujettir leur \italique{bois percé} (= \italique{carcan}) au \italique{cou} (αὐχήν) :
\citComm{ἀλλὰ τούτων χρῆν ἁπασῶν ἐς τετρημένον ξύλον\\
ἐγκαθαρμόσαι λαβόντας \textbf{τουτονὶ τὸν αὐχένα}.\\
Eh bien, il importe de s'assurer d'elles toutes, et d'assujettir dans un carcan \italique{ce cou-là}.}
On peut avoir un autre exemple d'αὐχήν avec ce sens-là au vers 184 du \italique{Cyclope} d'Euripide.}%
\note{3}{τράχηλος est phallossémique au vers 20 des \italique{Grenouilles}. Aristophane commence la pièce avec un rapide survol des blagues lourdes et vulgaires de ses concurrents aux concours comiques ; ce n'est pas sa muse qui s'abaisserait à de si grossières manœuvres ! non, \italique{on n'en soufflera mot}, οὐκ ἐρεῖ ! οὐκ ἐρεῖ !.. et puis au vers 19 :
\citComm{ὢ τρισκακοδαίμων ἄρ᾽ \textbf{ὁ τράχηλος οὑτοσί},\\
ὅτι θλίβεται μέν, τὸ δὲ γέλοιον οὐκ ἐρεῖ.\\
Ô trois fois malencontreuse \italique{nuque}, qui est moulue, et n'a pas le mot pour rire !}
Il faut bien-sûr lire -- je dis «~bien-sûr~», mais personne ne l'avait proposé jusqu'ici... --, il faut donc lire :
\citComm{ἄρ᾽ ὁ τράχηλος οὑτοσί,\\
ὅτι θλίβεται μέν, τὸ δὲ γέλοιον \textbf{οὐ κε ῥεῖ;}\\
Est-ce que le τράχηλος que voici, puisqu'il est pressé, pour la farce ne devrait pas couler ?}}\\
τὸ πέος διέλκεις πυκνότερον Κορινθίων.\\\passeLigne}%
\end{groupe}%
%% v. 649-650a
\begin{groupe}%
\personnageGrec{Μίκα}%
\trimIamb{ὤ μιαρὸς οὗτος· ταῦτ᾽ ἄρ᾽ ὑπὲρ Εὐριπίδου\\
ἡμῖν ἐλοιδορεῖτο.\\\passeLigne}%
\end{groupe}%
%% v. 650b-651
\begin{groupe}%
\personnageGrec{Μνησίλοχος}%
\trimIamb{\ligne{650}\decalage{ἡμῖν ἐλοιδορεῖτο.}κακοδαίμων ἐγώ,\\
εἰς οἷ᾽ ἐμαυτὸν εἰσεκύλισα πράγματα.\apparat{651}{εἰσεκύλησα R : \lecon{εἰσεκύλισα} Scaliger}%
\commentaire{651}{εἰσεκύλισα.}{Homère emploie le verbe κυλίνδειν en parlant de malheurs que le ciel \italique{fait rouler} sur quelqu'un : πῆμα θεὸς Δαναοῖσι κυλίνδει (\italique{Il}. 17, 688)\marqueur{1}. La même image se trouve chez Aristophane : νὴ τὸν Διόνυσον ἄπορά γ᾽ ἡμῖν πράγματα / δαίμων τις ἐσκεκύκληκεν ἐς τὴν οἰκίαν (\italique{Guêp.} 1474 s.), « ah! les inextricables tracas qu'un dieu a fait dévaler sur notre maison. » Le même verbe est employé d'une manière un peu différente, \italique{Thesm}. 767, mais toujours à propos de malheurs : ὁ μὲν γὰρ αἴτιος / κἄμ᾽ ἐσκυλίσας εἰς τοιαυτὶ πράγματα / οὐ φαίνεταί πω, « le coupable, celui qui m'a fait rouler dans de tels malheurs, n'apparaît toujours pas. » Cf. \italique{Thesm}. 650 s. : κακοδαίμων ἐγώ, / εἰς οἷ᾽ ἐμαυτὸν εἰσεκύλισα πράγματα. Il est clair que, dans ce cas, ἐμβάλλειν τινὰ εἰς συμφοράς n'est qu'un synonyme expressif\marqueur{2} de ἐμβάλλειν (p. ex. ἐμβάλλειν τινὰ εἰς συμφοράς, Antiphon 3, 4, 10 ; εἰς ἀτυχίας, Eschin. 3, 79 ; εἰς κακόν, Ar. \italique{Nuées} 1460). (Note de Taillardat, §558.)}%
\debutNote{651}%
\note{1}{Le moyen κυλίνδεσθαι se dit de malheur lui-même «~roulant~» sur quelqu'un : νῶιν δὴ τόδε πῆμα κυλίνδεσθαι, \italique{Il.}, 11, 347 ; cf. \italique{Od}. 2, 163 ; 8, 81.}%
\note{2}{Connu d'autres auteurs ; Phérécr. 146 [K = 156 K.-A.] : εἰκῆ μ᾿ ἐπῆρας ὄντα τηλικουτονὶ / πολλοῖς ἐμαυτὸν ἐγκυλῖσαι πράγμασιν ; Cf. Xén. \italique{Mém.} I, 2, 22 : τοὺς εἰς ἔρωτας ἐγκυλισθέντας ; Ἀπολλόδωρος ὁ Καρύστιος 5, 7 [ $=$ Athénée 280e] : εἰκῆ τέ πως / ἡμᾶς κυλίνδους᾽ ὅντιν᾽ ἂν τύχῃ τρόπον (sc. τύχη ἄγροικος).}\\\passeLigne}%
\end{groupe}%
%% v. 652a
\begin{groupe}%
\personnageGrec{Μίκα}%
\trimIamb{ἄγε δὴ τί δρῶμεν;\\\passeLigne}%
\end{groupe}%
%% v. 652b-654
\begin{groupe}%
\personnageGrec{Κλεισθένης}%
\trimIamb{\ligne{652}\decalage{ἄγε δὴ τί δρῶμεν;}τουτονὶ φυλάττετε\\
καλῶς, ὅπως μὴ διαφυγὼν οἰχήσεται·\apparat{653}{οἴχεται R : \lecon{οἰχήσεται} Faber p.788 \ligneApp{654} πρυτάνεσσιν R : \lecon{-εσιν} \Mon}\\
ἐγὼ δὲ ταῦτα τοῖς πρυτάνεσιν ἀγγελῶ.%\apparat{654}{πρυτάνεσσιν R : \lecon{πρυτάνεσιν} \Mon}
\\\passeLigne}%
\end{groupe}%
%% v. 655-662
\begin{groupe}%
\personnageGrec{Κορυφαῖος}%
\typeVers{\decalageN{4}$4$ an\cat}\tetramAnap{ἡμᾶς τοίνυν μετὰ τοῦτ᾽ ἤδη τὰς λαμπάδας ἁψαμένας χρὴ%
\commentaire{655-86}{Le chœur compris entre les vv. 655-686.}{La \nieme{19} note d'Alphonse Willems est d'une telle richesse d'explication que je ne peux que la reproduire intégralement ci-dessous. Pour ce qui touche à la colométrie, je reprends celle de Lætitia\marqueur{1}. (Le passage des vers 667-686 a de fortes correspondances métriques avec les vers 707-725).\\
Ce petit chœur n'a pas été examiné avec assez d'attention. On n'a pas remarqué qu'il s'accompagne d'un jeu de scène unique dans l'œuvre d'Aristophane. Trois morceaux le composent : \\
\ \ \ 1° Les vv. 655 à 658 appartiennent au coryphée. Mnésiloque n'est peut-être pas le seul homme qui ait pénétré parmi les femmes. Il importe de s'en assurer, et, après avoir déposé les manteaux et allumé les torches, de « tourner et virer tout à l'entour du Pnyx, et d'inspecter avec soin les tentes et les passages ».\\
\ \ \ 2° Les quatre suivants (659-662) sont des trochées, et le changement de mètre indique un changement d'interlocuteur. C'est, en effet, la réponse du \italique{parastate} ou chef du demi-chœur [le \italique{Ravennas} porte sur sa marge, η\auDessus{μ}{χ}ι γ\auDessus{υ}{ν}, écrite d'une seconde main, mais dont l'indication se retrouve sur le manuscrit de Münich]. En disant qu'il y avait lieu d'inspecter le Pnyx, les tentes et les passages, la coryphée indiquait clairement qu'il fallait quitter l'orchestre, c'est-à-dire l'enceinte sacrée adjacente au temple\marqueur{2}. La parastate annonce donc que le chœur va « s'élancer dehors (ἐξορμᾶν) d'un pied léger\marqueur{3} ». Puis elle engage celle qu'elle appelle τὴν πρώτην à commencer tout de suite la reconnaissance. Comment faut-il interpréter τὴν πρώτην ? D'y voir, comme on le fait généralement, une locution adverbiale (\latin{quam primum} ou \latin{prima quaque via}) me semble inadmissible. Avec M. Mazon, je crois que ce terme désigne une personne déterminée\marqueur{4} ; suivant moi, ce serait τὴν πρώτην ἱέρειαν, appelée aussi τὴν προστάτιν, savoir la première, il y en avait beaucoup\marqueur{5}, des prêtresses temporaires, ἀρήτειραι, de Déméter, celle qui, tout à l'heure, avait prononcé les prières d'usage (vv. 295 et suiv.), en un mot, la coryphée. 
Callimaque, parlant d'une de ces prêtresses, ἱερέη Δήμητρος, se sert d'une désignation analogue : πολλῶν προστασίη νέων γυναικῶν\marqueur{6}. \\
\ \ \ 3° Le morceau qui suit (vv. 663-684) est un \italique{hyporchême}. Une partie du chœur reste en place et chante, tandis que l'autre, en quête de quelque complice de Mnésiloque\marqueur{7}, évolue silencieusement, σιωπῇ (v. 660), va et vient, entre et sort, cherchant partout, dans l'orchestre et dehors, καὶ τὰ τῇδε καὶ τὰ δεῦρο (v. 666), jusqu'à ce qu'il déclare n'avoir rien trouvé (687-688).\\
%\eff{Après cette fine analyse de Willems, s'agit-il de roues (trokhos) cf. Daremberg Salvatio 1027a ? τὴν πρώτην τρέχειν χρῆν ὡς τάχιστ᾽ ἤδη κύκλῳ}
}%
\debutNote{655-86}%
\note{1}{\italique{The songs of Aristophanes}, pp. 418-422.}%
\note{2}{On peut conclure de là que le Thesmophorion, qu'on savait situé sur une hauteur, l'était sur la colline du Pnyx et voisin du lieu d'assemblée.}%
\note{3}{Le simple, ὀρμᾶν, qui se lit plus loin, v. 953, et \italique{Gren}., 478, signifierait « prendre sa course ».}%
\note{4}{\italique{Essai sur la composition des comédies d'Aristophane}, Paris (1904), p. 131, n. 4.}%
\note{5}{Il y en avait deux par dème (Isée, \italique{Succession de Kiron}, 19-20), donc plus de deux cents.}%
\note{6}{\italique{Anthol. Palat}., VII, 728.}%
\note{7}{εἴ τις ἐν τόποις ἑδραῖος. Qu'avons-nous affaire d'ajouter, avec MM. Blaydes et van Leeuwen, τούτοισι τοῖς devant τόποις ? La même locution adverbiale se lit dans un fragment de Nausicratès : [τρίγκλαι] ἃς κλύδων Αἰξωνικὸς ἐν τόποις παιδεύεται (\italique{Fragm. Com. Gr}., t. IV, p. 575 [= Athénée 325e = Nausicratès, fr. 1.6-11 Kassel-Austin ; le texte n'est pas établi avec certitude : on trouve aussi : ἐντόπους à la place de ἐν τόποις. Même si l'exemple n'est pas probant, Willems n'en a pas moins raison ; Lætitia remarque que l'adjonction de Blaydes accroît la disparité entre les deux stances des vers 663-666 -- là où l'on chercherait à rétablir une symétrie.]) et la locution inverse se lit dans Sophocle : ἥτις μ᾽ ἀποικίσειεν ἐκ τόπων (\italique{Trach}., 955), où jamais personne n'a trouvé à redire (ἐκ τούτων τῶν τόπων.). Je ne vois pas que ces tournures soient plus bizarres que ἐν ποσί et ἐκ ποδῶν, ou bien que ἐν χώρᾳ et ἐκ χώρας, πρὸ ou παρὰ χειρῶν, etc.}%
\\
ξυζωσαμένας εὖ κἀνδρείως τῶν θ᾽ ἱματίων ἀποδύσας%
\commentaire{656}{τῶν θ᾽ ἱματίων ἀποδύσας.}{Willems, fidèle à son poste de \italique{gardien de la tradition}, justifie contre les retoucheurs de texte l'expression ἀποδῦναι $+$ gén. (C'est sa \nieme{20} note aux \italique{Thesmophories}.)}\\
\Automatique{ζητεῖν, εἴ που κἄλλος τις ἀνὴρ ἐπελήλυθε, καὶ περιθρέξαι}\apparat{657}{εἰσελήλυθεν R : %ἀνελήλυθε Fritzsche : 
\lecon{ἐπελήλυθε} Handley \textbf{658} ἀθρῆσαι R : \lecon{διαθρῆσαι} Daubuz}%
\commentaire{657}{εἰσελήλυθεν.}{La leçon du manuscrit de Ravenne est \italique{contra metrum}. Tannegui Le Fèvre (Faber dans l'apparat), le père de Mme Dacier, avait proposé la simple correction ἐσελήλυθε. Austin\marqueur{1} la réfuta : ἐς ne précède jamais une voyelle chez Aristophane. Restent les conjectures de Fritzsche et de Handley. Fritzsche proposa de lire ἀνελήλυθε, l'expliquant par une haplographie ἀνὴρ (ἀν)ἐλήλυθε suivie d'une complétion postérieure εἰσ-. % (on reste circonspect devant de telles explications) ; 
Handley proposa ἐπελήλυθε (la confusion entre ΙC et Π en onciales étant une erreur courante) -- hypothèse beaucoup plus simple et, partant, plus satisfaisante.}%
\debutNote{657}%
\note{1}{Voir ses notes : « Textual problems in Ar. \italique{Thesm}. » (1987) ; je reprends la trame de celle qui nous intéresse présentement.}\\
\Automatique{τὴν Πύκνα πᾶσαν καὶ τὰς σκηνὰς καὶ τὰς διόδους διαθρῆσαι.}\\\passeLigne}%
\end{groupe}%
\begin{groupe}%
\personnageGrec{Κορυφ. β}%
\typeVers{$4$ troch.\cat}\tetramTroch{εἶα δὴ πρώτιστα μὲν χρὴ κοῦφον ἐξορμᾶν πόδα\testimonia{659}{(εἶα δὴ) Hesychius ε 8 ἔα δή· ἄγε δή $=$ Photius ε  1 $=$ \Syn ε 1 $=$ Suidas ε 3}\apparat{659}{(Ἡμιχ. γυν.) \Rpc \ligneApp{660} πανταχῇ post χρὴ R : \lecon{ante μόνον} Küster}\\
καὶ διασκοπεῖν σιωπῇ πανταχῇ· μόνον δὲ χρὴ\\%\apparat{660}{πανταχῇ post χρὴ R : \lecon{ante μόνον} Küster}\\
μὴ βραδύνειν, ὡς ὁ καιρός ἐστι μὴ μέλλειν ἔτι,\\
\Automatique{ἀλλὰ τὴν πρώτην τρέχειν χρῆν ὡς τάχιστ᾽ ἤδη κύκλῳ.}\apparat{662}{χρὴ R : \lecon{χρῆν} Bentley : χρῆ μ᾽ Austin}%
%\commentaire{662}{τὴν πρώτην}{La plupart des traductions et des commentaires donnent une valeur adverbiale à \grec{τὴν πρώτην}. L'interprétation de Willems\marqueur{1} me semble bien préférable ...}%
%\debutNote{662}%
%\note{1}{Cf. sa note 19 reproduite en annexe.}
\\\passeLigne}%
\end{groupe}%
%% v. 663-666
\begin{groupe}%
\personnageGrec{Χορός}%
\typeVers{$5$ troch.}\chantChoeur{\ligne{663}εἶά νυν ἴχνευε καὶ μά-}%
\end{groupe}%
%% v. 663-666
\begin{groupe}%
\decoupeVers{4a}\typeVers{}\chantChoeur{\ligne{663}\alinea τευε ταχὺ πάντ᾽, εἴ τις ἐν τόποις ἑδραῖος}%
\end{groupe}%
%% v. 663-666
\begin{groupe}%
\decoupeVers{b}\typeVers{lecyth.}\chantChoeur{ἄλλος αὖ λέληθεν ὤν.}%
\end{groupe}%
%% v. 663-666
\begin{groupe}%
\typeVers{$2$ troch.}\chantChoeur{πανταχῇ δὲ ῥῖψον ὄμμα,\apparat{665}{διάρριψον R : \lecon{δὲ ῥῖψον} Hermann \ligneApp{666} post καὶ τὰ τῇδε add. καὶ τὰ κεῖσε Kaibel}}%
\end{groupe}%
%% v. 663-666
\begin{groupe}%
\decoupeVers{a}\typeVers{$2$ troch.}\chantChoeur{καὶ τὰ τῇδε καὶ τὰ δεῦρο%\apparat{666}{Post καὶ τὰ τῇδε add. καὶ τὰ κεῖσε Kaibel}
}%
\end{groupe}%
%% v. 663-666
\begin{groupe}%
\decoupeVers{b}\typeVers{lecyth.}\chantChoeur{\ligne{666}πάντ᾽ ἀνασκόπει καλῶς.\\\passeLigne}%
\end{groupe}%
%% v. 667-686
\begin{groupe}%
\typeVers{$2$ anap.}\chantChoeur{\ligne{667}ἢν γάρ ληφθῇ δράσας ἀνόσια,\apparat{667}{μὴ λάθῃ R : \lecon{ληφθῇ} Reisig : με λάθῃ Bergk}}%
\end{groupe}%
%% v. 663-666
\begin{groupe}%
\typeVers{$2$ anap.}\chantChoeur{δώσει τε δίκην καὶ πρὸς τούτῳ}%
\end{groupe}%
%% v. 663-666
\begin{groupe}%
\typeVers{$2$ anap\cat}\chantChoeur{τοῖς ἄλλοις ἀνδράσιν ἔσται\apparat{669}{ἅπασιν \inRas{ἅπασιν}{} ἔσται \Rpc : ἅπασιν ἔσται \Mon : \lecon{ἀνδράσιν ἔσται} Beer ap. Hermann (1845) : ἔσται ἅπασιν Bergk : alii alia}}%
\end{groupe}%
%% v. 663-666
\begin{groupe}%
\typeVers{anap.}\chantChoeur{παράδειγμ᾽ ὕβρεως}%
\end{groupe}%
%% v. 663-666
\begin{groupe}%
\typeVers{$2$ anap.}\chantChoeur{ἀδίκων τ᾽ ἔργων ἀθέων τε τρόπων·}%
\end{groupe}%
%% v. 663-666
\begin{groupe}%
\typeVers{$2$ anap.}\chantChoeur{φήσει δ᾽ εἶναί τε θεοὺς φανερῶς,}%
\end{groupe}%
%% v. 663-666
\begin{groupe}%
\typeVers{anap.}\chantChoeur{δείξει τ᾽ ἤδη}%
\end{groupe}%
%% v. 663-666
\begin{groupe}%
\typeVers{$3$ troch\cat}\chantChoeur{πᾶσιν ἀνθρώποις σεβίζειν δαίμονας}%
\end{groupe}%
%% v. 663-666
\begin{groupe}%
\decoupeVers{6}\typeVers{$2$ dochm. ?}\chantChoeur{†δικαίως τ᾽ ἐφέποντας † ὅσια καὶ νόμιμα\apparat{675}{\lecon{ἐφέποντας} R : ἐφέπειν Hermann}}%
\end{groupe}%
%% v. 663-666
\begin{groupe}%
\typeVers{$2$ dochm.}\chantChoeur{\ligne{677}μηδομένους ποιεῖν ὅ τι καλῶς ἔχει.}%
\end{groupe}%
%% v. 663-666
\begin{groupe}%
\typeVers{$3$ iamb\cat}\chantChoeur{κἂν μὴ ποιῶσι ταῦτα τοιάδ᾽ ἔσται·}%
\end{groupe}%
\endnumbering
\end{Leftside}
\begin{Rightside}
\beginnumbering
%% v. 635+ (part 2)
\begin{groupe}
\trimIambFr{\ligne{635}Voilà l'homme dont tu parles.\\\passeLigne}%
\end{groupe}%
%% v. 635b
\begin{groupe}%
\personnageFr{Clisthène}%
\trimIambFr{\espace{3.7cm}Que dois-je faire ?\\\passeLigne}%
\end{groupe}%
%% v. 636
\begin{groupe}%
\personnageFr{Mika}%
\trimIambFr{Déshabille-le. Ce qu'il dit n'a pas de sens.\\\passeLigne}%
\end{groupe}%
%% v. 637-8
\begin{groupe}%
\personnageFr{Mnésiloque}%
\trimIambFr{Eh quoi! déshabiller une mère de neuf enfants !\\
Il défait mon strophion avec promptitude. Impudent !\\\passeLigne}%
\end{groupe}%
%% v. 639-640
\begin{groupe}%
\personnageFr{Mika}%
\trimIambFr{Qu'elle a l'air robuste et vigoureuse !\\
Et, par Zeus, elle n'a pas de tétons comme nous.\\\passeLigne}%
\end{groupe}%
%% v. 641
\begin{groupe}%
\personnageFr{Mnésiloque}%
\trimIambFr{C'est que je suis stérile. Je n'ai jamais conçu.\testimonia{641}{(στερίφη) Hesychius σ 1769 στερίφη· στεῖρα, μὴ τεκοῦσα, μηδὲ τίκτουσα $=$ Photius σ 532 ; cf. Moeris σ 46 (στερίφη $=$ στερρά)}\\\passeLigne}%
\end{groupe}%
%% v. 642
\begin{groupe}%
\personnageFr{Clisthène}%
\trimIambFr{À présent. Mais tout à l'heure tu étais mère de neuf enfants.
Tiens-toi debout. Où pousses-tu ton cas là-\fin dessous ?\\\passeLigne}%
\end{groupe}%
%% v. 644
\begin{groupe}%
\personnageFr{Mika}%
\trimIambFr{(Derrière Mnésiloque). Le voilà qui émerge, et beau \fin de couleur, le coquin.\\\passeLigne}%
\end{groupe}%
%% v. 645a
\begin{groupe}%
\personnageFr{Clisthène}%
\trimIambFr{(Échangeant de place avec Mika). Où est-il ?\\\passeLigne}%
\end{groupe}%
%% v. 645b
\begin{groupe}%
\personnageFr{Mika}%
\trimIambFr{\ligne{645}\espace{2.2cm}Il est reparti en avant.\\\passeLigne}%
\end{groupe}%
%% v. 646a
\begin{groupe}%
\personnageFr{Clisthène}%
\trimIambFr{(Revenant devant.) Non, pas de ce côté.\\\passeLigne}%
\end{groupe}%
%% v. 646b
\begin{groupe}%
\personnageFr{Mika}%
\trimIambFr{\ligne{646}\espace{2.2cm}Eh! le voilà revenu par ici.\\\passeLigne}%
\end{groupe}%
%% v. 647-648
\begin{groupe}%
\personnageFr{Clisthène}%
\trimIambFr{Une façon d'isthme que tu as là, homme : tu fais passer ton cas de-ci de-là,
plus dru que les Corinthiens.\\\passeLigne}%
\end{groupe}%
%% v. 649-650a
\begin{groupe}%
\personnageFr{Mika}%
\trimIambFr{Ah! la canaille! Voilà donc pourquoi au profit d'Euripide\\
il nous invectivait.\\\passeLigne}%
\end{groupe}%
%% v. 650b-651
\begin{groupe}%
\personnageFr{Mnésiloque}%
\trimIambFr{\ligne{650}\espace{2.5cm}Malheureux que je suis !\\
dans quelle méchante affaire je me suis empêtré !\\\passeLigne}%
\end{groupe}%
%% v. 652a
\begin{groupe}%
\personnageFr{Mika}%
\trimIambFr{Voyons, que faisons-nous ?\\\passeLigne}%
\end{groupe}%
%% v. 652b-654
\begin{groupe}%
\personnageFr{Clisthène}%
\trimIambFr{\ligne{652}\espace{2.6cm}Veillez bien sur lui,\\
qu'il ne se dérobe point par la fuite.\\
Moi, je vais rapporter la chose aux prytanes.\\\passeLigne}%
\end{groupe}%
%% v. 655-662
\begin{groupe}%
\personnageFr{Chœur}%
\tetramAnapFr{À nous maintenant, après avoir allumé nos torches,
serré bien et résolument nos tuniques et mis bas nos manteaux,
de chercher si d'aventure il n'est pas entré quelqu'autre homme. Tournons et virons
tout à l'entour du Pnyx, inspectons avec soin les tentes et \fin les passages.\\\passeLigne}%
\end{groupe}%
\begin{groupe}%
\tetramTrochFr{Or sus! tout d'abord élançons-nous au dehors d'un pied léger
et épions partout en silence. Surtout ne tardons pas : le cas ne souffre plus de délai.
Il faut que notre coryphée coure au plus tôt faire la ronde.
(Ici commence un chœur dansé.)\\\passeLigne}%
\end{groupe}%
%% v. 663-666
\begin{groupe}%
\chantChoeurFr{Or sus ! suis chaque piste, quête vivement,}%
\end{groupe}%
%% v. 663-666
\begin{groupe}%
\chantChoeurFr{au cas où quelqu'autre en ces lieux}%
\end{groupe}%
%% v. 663-666
\begin{groupe}%
\chantChoeurFr{se serait tapi à notre insu.}%
\end{groupe}%
%% v. 663-666
\begin{groupe}%
\chantChoeurFr{Jette les yeux partout.}%
\end{groupe}%
%% v. 663-666
\begin{groupe}%
\chantChoeurFr{Que rien ici, rien là-bas,}%
\end{groupe}%
%% v. 663-666
\begin{groupe}%
\chantChoeurFr{n'échappe à ton \oe il vigilant.\\\passeLigne}%
\end{groupe}%
%% v. 667-686
\begin{groupe}%
\chantChoeurFr{Car s'il est pris après ce sacrilège,}%
\end{groupe}%
%% v. 663-666
\begin{groupe}%
\chantChoeurFr{il paiera son crime, et de surcroît}%
\end{groupe}%
%% v. 663-666
\begin{groupe}%
\chantChoeurFr{prémunira tous les autres,}%
\end{groupe}%
%% v. 663-666
\begin{groupe}%
\chantChoeurFr{par son exemple, contre la violence,}%
\end{groupe}%
%% v. 663-666
\begin{groupe}%
\chantChoeurFr{les actes injurieux}%
\end{groupe}%
%% v. 663-666
\begin{groupe}%
\chantChoeurFr{et les façons impies.}%
\end{groupe}%
%% v. 663-666
\begin{groupe}%
\chantChoeurFr{Il confessera hautement qu'il est des dieux,}%
\end{groupe}%
%% v. 663-666
\begin{groupe}%
\chantChoeurFr{il apprendra}%
\end{groupe}%
%% v. 663-666
\begin{groupe}%
\chantChoeurFr{à tous les hommes à les honorer,}%
\end{groupe}%
%% v. 663-666
\begin{groupe}%
\chantChoeurFr{à s'efforcer, par une conduite pieuse et un constant }%
\end{groupe}%
%% v. 663-666
\begin{groupe}%
\chantChoeurFr{souci des lois, de ne rien faire que de bien.}%
\end{groupe}%
\endnumbering
\end{Rightside}
\end{pages}
\Pages
\clearpage\ifodd\value{page}\hbox{}\newpage\fi
\enTeteG{ΕΠΕΙΣΟΔΙΟΝ Γ' --- Codex Ravennas 429, f. 170v (Classense 172v) --- vv. 679–720}
\begin{adjustwidth}{-1.5cm}{-3.0cm}
\begin{adjustbox}{max totalheight=21.0cm, max width=16.5cm, keepaspectratio}
\begin{minipage}{16.5cm}
{\fontsize{7.2pt}{8.5pt}\selectfont
\noindent\textbf{Folio 17 --- Codex Ravennas 429, \texttt{f. 170v} (\texttt{Classense 172v}) : ΕΠΕΙΣΟΔΙΟΝ Γ', vv. 679–720}\hfill \textit{Témoins :} $\text{R}^1, \Rpc, \Rdeux$ (\Mon)\par\vspace{2pt}
\hrule height 0.4pt\vspace{2.5pt}
}
{\fontsize{8.1pt}{9.5pt}\selectfont
\ligneRav{679}{χορ/}{\hspace{1.1em}αὐτῶν ὅταν ληφθῇ τις †ὅσια δρῶν†,}{\textbf{679} \lecon{αὐτῶν} R : τούτων Richards \espApp \lecon{ὅσια} R : ἀνόσια Zanetti : ὅσια μὴ Burges : ἀνόσιόν τι Hermann}
\ligneRav{680}{}{\hspace{1.1em}μανίαις φλέγων λύσσῃ παράκο-}{}
\ligneRav{}{}{\hspace{1.1em} πος, εἴ τι δρῴη πᾶσιν ἐμφανὴς ὁρᾶν}{}
\ligneRav{683}{}{\hspace{1.1em}\difference{ἔσται} γυναιξὶ καὶ βροτοῖς,}{\textbf{683} ἐστὶν R : \lecon{ἔσται} Bothe \espApp γυναιξὶν : \lecon{-ξὶ} Brunck}
\ligneRav{}{}{\hspace{1.1em}ὅτι τά παράνομα τά τ᾽ ἀνόσια παρὼν}{}
\ligneRav{686}{}{\hspace{1.1em}θεὸς \difference{παραχρῆμ᾽ ἀποτίνεται}·}{\textbf{686} ἀποτίνεται· | παραχρῆμά τε τίνεται R : \lecon{παραχρῆμ᾽ ἀποτίνεται} Hermann}
\ligneRav{687}{χορ/}{ἀλλ᾽ ἔοιχ᾽ ἡμῖν ἅπαντά πως διεσκέφθαι καλῶς·}{\textbf{687} (Κορ.) add. von Velsen\espApp Κορ. add. Velsen \espApp \lecon{πως} R : που Blaydes}
\ligneRav{}{}{οὐχ ὁρῶμεν γοῦν ἔτ᾽ ἄλλον οὐδέν᾽ ἐγκαθήμενον}{}
\ligneRav{689}{\auDessus{γυ}{ν}}{\difference{ἆ ἆ}}{\textbf{689} ἀά supra lin. R : \lecon{ἆ ἆ} Thiersch \espApp ποῖ R : \lecon{ποῖ ποῖ} Biset}
\ligneRav{689}{}{\difference{ποῖ ποῖ} σὺ φεύγεις; οὗτος οὗτος οὐ μενεῖς;}{\textbf{689} ἀά supra lin. R : \lecon{ἆ ἆ} Thiersch \espApp ποῖ R : \lecon{ποῖ ποῖ} Biset}
\ligneRav{690}{}{τάλαιν᾽ ἐγὼ τάλαινα, καὶ τὸ παιδίον}{}
\ligneRav{}{}{ἐξαρπάσας μοι φροῦδος ἀπὸ τοῦ τιτθίου·}{}
\ligneRav{}{μ\auDessus{ν}{Η}}{κέκραχθι· τοῦτο δ᾽ οὐδέποτε σὺ ψωμιεῖς,}{}
\ligneRav{693}{}{τάδε δ᾽ \difference{ἀφῆτ᾽} (ἀφε\auDessus{ῖ}{ῆ}τ᾽ \Rdeux) ἤδη τὸ παιδίον;}{\textbf{693} ἀφε\auDessus{ῖ}{ῆ}τ᾽ \Rdeux : ἀφεῖτ᾽ \Mon}
\ligneRav{}{}{πληγὲν μαχαίρᾳ τῇδε φοινίας φλέβας}{}
\ligneRav{695}{—}{καθαιματώσει βωμόν. : ὢ τάλαιν᾽ ἐγώ·}{}
\ligneRav{}{}{γυναῖκες, οὐκ ἀρήξετ᾽; οὐ πολλὴν βοὴν}{}
\ligneRav{697}{}{στήσεσθε καὶ τροπαῖον, ἀλλὰ τοῦ μόνου}{\textbf{697} καὶ om. R : καὶ suppl. Σ ad \italique{Pl.} v.453 et Greg. Cor. \italique{de dialectis} p.22}
\ligneRav{}{—}{τέκνου με περιόψεσθ᾽ ἀποστερουμένην; : ἔα ἔα·}{}
\ligneRav{700}{}{\hspace{1.1em}ὦ πότνιαι Μοῖραι τί \difference{τόδε} δέρκομαι}{\textbf{700} δὲ R : \lecon{τόδε} Dobree : δὴ Reisig}
\ligneRav{}{}{\hspace{1.1em}νεοχμὸν αὖ τέρας;}{}
\ligneRav{702}{}{ὡς \difference{ἅπαντ᾽ ἄρ᾽ ἐστὶ} τόλμης †\difference{πλέα}† κἀναισχυντίας·}{\textbf{702} ἅπαν γὰρ ἐστὶν R : ἅπαντα Suid., cod. Laur. : \lecon{ἅπαντ᾽ ἄρ᾽ ἐστὶ} Bentley \espApp ἔργα R : \lecon{πλέα} Suid. : δοῦλα cod. Laur. : μεστὰ Porson}
\ligneRav{}{}{οἷον αὖ δέδρακεν ἔργον, οἷον αὖ φίλαι τόδε·}{}
\ligneRav{704}{μ\auDessus{ν}{Η}}{\difference{οἷον}; ὑμῶν \difference{ἐξαρήξω} τὴν ἄγαν αὐθαδίαν·}{\textbf{704} οἷον R : \lecon{οἷον;} Ellebodius \espApp ἐξάρξω R : ἐξαράξει Ellebodius : ἐξαράξω Bentley : \lecon{ἐξαρήξω} Bibi}
\ligneRav{705}{χορ/}{ταῦτα δῆτ᾽ οὐ δεινὰ πράγματ᾽ ἐστὶ καὶ περαιτέρω;}{}
\ligneRav{706}{\auDessus{γυ}{ν}}{δεινὰ δῆθ᾽, \difference{ὁτιή} γ᾽ ἔχει μου 'ξαρπάσας τὸ παιδίον·}{\textbf{706} ὅτι R : \lecon{ὁτιή} Faber : ὅστις Porson}
\ligneRav{}{χορ/}{\hspace{1.1em}τί ἂν οὖν εἴποι πρὸς ταῦτά τις, ὅτε}{}
\ligneRav{}{}{\hspace{1.1em}τοιαῦτα ποιῶν ὅδ᾽ ἀναισχυντεῖ;}{}
\ligneRav{}{μ\auDessus{ν}{Η}}{\hspace{1.1em}κοὔπω μέντοι γε πέπαυμαι·}{}
\ligneRav{710}{χορ/}{\hspace{1.1em}ἀλλ᾽ οὖν \difference{ἥκειν γ᾽} ὅθεν ἥκεις·}{\textbf{710} ἥκεις τ᾿ R : ἥκεις γ᾽ Suid. η 170 : ἥκων γ᾽ Willems : \lecon{ἥκειν γ᾽} Bibi ap. Phot. η 88}
\ligneRav{711}{}{\hspace{1.1em}φαύλως \difference{δ᾿} ἀποδρὰς \difference{οὐκ ἔξει}·}{\textbf{711} τ᾿ R, Suid : \lecon{δ᾿} Bentley \espApp οὐ λέξεις R, Suid. : \lecon{οὐκ ἔξει.} Bibi ap. Σ}
\ligneRav{}{}{\hspace{1.1em}οἷον δράσας διέδυς ἔργον·}{}
\ligneRav{}{}{\hspace{1.1em}λήψει δὲ κακόν·}{}
\ligneRav{}{μ\auDessus{ν}{Η}}{τοῦτο μέντοι μὴ γένοιτο μηδαμῶς, ἀπεύχομαι·}{}
\ligneRav{715}{χορ/}{\hspace{1.1em}τίς οὖν σοι, τίς ἂν σύμμαχος ἐκ θεῶν}{\textbf{715} \lecon{οὖν σοι} R : ἂν σοι Hamaker}
\ligneRav{}{}{\hspace{1.1em}ἀθανάτων ἔλθοι ξὺν ἀδίκοις ἔργοις;}{}
\ligneRav{717}{μ\auDessus{ν}{Η}}{μάτην λαλεῖτε· τὴνδ᾽ ἐγὼ οὐκ ἀφήσω·}{\textbf{717} τὴν δ᾽ R : \lecon{τήνδ᾽} Biset}
\ligneRav{}{χορ/}{\hspace{1.1em}ἀλλ᾽ οὐ μὰ τὼ θεὼ τάχ᾽ οὐ}{}
\ligneRav{719}{}{\hspace{1.1em}χαίρων ἴσως \difference{ἐνυβριεῖς}}{\textbf{719} ἐνυβρίσεις R : \lecon{ἐνυβριεῖς} anonymus Romanus}
\ligneRav{720}{}{\hspace{1.1em}λόγους τε λέξεις ἀνοσίους}{}
}
\vspace{3pt}\noindent\rule{4.5cm}{0.4pt}\par\vspace{1.5pt}
{\fontsize{7.0pt}{8.3pt}\selectfont
\noindent\textbf{Loci paralleli (vv. 679–720) ---} \textbf{696-7} (γυναῖκες -- τροπαῖον) \GregCor p.20-22 cf. schol.\expo{Ald} \italique{Plutus} 453c \espTest \textbf{696-7} (γυναῖκες -- τροπαῖον) \GregCor pp. 20-22 ; cf. scholium\expo{Ald} \italique{Plutus} 453c \espTest \textbf{702} (ἅπαντα pro ἅπαν γὰρ ἐστὶν -- κἀναισχυντίας) Suidas α 2902 (sed πλέα pro ἔργα) ; \italique{Corpus paroemiographorum Graecorum}, Appendices, I, 35 (sed ἔργα omisso) ; \italique{Laurentiamus} LVIII 24 p. 40 (ed. Cohn) (sed δοῦλα pro ἔργα) \espTest \textbf{702} (ἅπαντ' -- κἀναισχυντίας) Suidas α 2902 (sed ἅπαντα pro ἅπαν γὰρ ἐστὶν et πλέα pro ἔργα) ; \italique{Corpus paroemiographorum Graecorum}, Appendices, I, 35 (sed ἅπαντα pro ἅπαν γὰρ ἐστὶν et ἔργα omisso) ; \italique{Laurentiamus} LVIII 24 p. 40 (ed. Cohn) (sed ἅπαντα pro ἅπαν γὰρ ἐστὶν et δοῦλα pro ἔργα) \espTest \textbf{710-3} Suidas η 170 (sed λήψῃ pro λήψει) \espTest \textbf{710} Photius η 88 ἥκειν· ἀντὶ τοῦ ἥκων\par
}
\vspace{2.5pt}\noindent\rule{4.5cm}{0.4pt}\par\vspace{1.5pt}
{\fontsize{6.9pt}{8.2pt}\selectfont
\noindent\textbf{Suite (en commentaire) de l'apparat critique \& notes paléographiques (\texttt{f. 170v}) ---} \textbf{Page particulièrement riche en corruptions d'onciale et de minuscule élucidées par Louis Hauseux (\textbf{Bibi}) :} --- \textbf{v. 679 :} \lecon{αὐτῶν} R (τούτων Richards) \espApp \lecon{ὅσια} R (maintenu par L. Hauseux avec interrogation ; ἀνόσια Zanetti ; ὅσια μὴ Burges ; ἀνόσιόν τι Hermann). --- \textbf{v. 697 :} καὶ om. R : \lecon{καὶ} suppl. $\Sigma$ ad \italique{Plut.} 453 et Grégoire de Corinthe, \italique{De dialectis}, p. 22. --- \textbf{v. 704 :} οἷον R : \lecon{οἷον;} Ellebodius \espApp ἐξάρξω R (ἐξαράξει Ellebodius ; ἐξαράξω Bentley) : \lecon{ἐξαρήξω} \lecon{Bibi} ap. Hésychius $\varepsilon\ 3599$. --- \textbf{v. 710 :} ἥκεις τ' R (ἥκεις γ' Souda $\eta\ 170$ ; ἥκων γ' Willems) : \lecon{ἥκειν γ'} \lecon{Bibi} ap. Photius $\eta\ 88$ (permutation minuscule $\nu \to \varsigma$). --- \textbf{v. 711 :} τ' R, Souda : \lecon{δ'} Bentley \espApp οὐ λέξεις R, Souda : \lecon{οὐκ ἔξει.} \lecon{Bibi} ap. $\Sigma^{\text{R}}$ (ὅθεν οὐ φεύξῃ, confusion onciale Κ/Λ).\par
}
\end{minipage}
\end{adjustbox}
\end{adjustwidth}
\newpage
\enTeteD{Codex Ravennas 429 (Biblioteca Classense) --- Fac-similé du f. 170v (vv. 679–720)}
\begin{adjustwidth}{-1.5cm}{-3.0cm}
\begin{center}
\includegraphics[width=16.5cm,height=20.5cm,keepaspectratio]{/workspaces/Thesmophories/dossier_ravennas_thesmophories/folios_individuels_ff162v-177r/17_f170v_classense172v_vv0679-0720.jpg}\par
\vspace{2pt}
{\scriptsize\textbf{Codex Ravennas 429, f. 170v (Classense 172v)} --- ΕΠΕΙΣΟΔΙΟΝ Γ', \textbf{vv. 679–720} ($\text{R}^1, \Rpc, \Rdeux$ (\Mon)).}
\end{center}
\end{adjustwidth}
\newpage
\clearpage\ifodd\value{page}\hbox{}\newpage\fi
\enTeteG{ΘΕΣΜΟΦΟΡΙΑΖΟΥΣΑΙ --- Texte grec établi (f. 170v, vv. 679–720)}
\enTeteD{LES THESMOPHORIES --- Traduction française (f. 170v, vv. 679–720)}
\begin{pages}
\begin{Leftside}
\beginnumbering
%% v. 663-666
\begin{groupe}%
\typeVers{$2$ iamb...}\chantChoeur{αὐτῶν ὅταν ληφθῇ τις †ὅσια δρῶν†,\apparat{679}{\lecon{αὐτῶν} R : τούτων Richards \espApp \lecon{ὅσια} R : ἀνόσια Zanetti : ὅσια μὴ Burges : ἀνόσιόν τι Hermann}}%
\end{groupe}%
%% v. 680
\begin{groupe}%
\typeVers{$5$ iamb.}\chantChoeur{μανίαις φλέγων λύσσῃ παράκο-
\commentaire{680}{μανίαις φλέγων λύσσῃ παράκοπος.}{Παράκοπος, \italique{délirant}, est une métaphore usée\marqueur{1} et d'origine incertaine. Elle paraît sans rapport avec la frappe de la monnaie\marqueur{2} ; mais elle en a peut-être un avec la notion d'« ébranler le cerveau »\marqueur{3}, cf. Euripide \italique{Hippolyte}, v. 238, etc... παρακόπτειν φρένας. (Note de Taillardat §476.)}%
\debutNote{680}%
\note{1}{Eschyle \italique{Prom.} 581 ; Eur. \italique{Bacch.} 33, 1000, etc.}%
\note{2}{Malgré παρακόπτειν (\italique{Ach.} 517), « battre de la fausse monnaie ».}%
\note{3}{Cf. §474 du même livre : \italique{Les images d'Aristophane} de Jean Taillardat.}}%
\end{groupe}%
%% v. 681-2
\begin{groupe}%
\decoupeVers{2}\typeVers{}\chantChoeur{\alinea πος, εἴ τι δρῴη πᾶσιν ἐμφανὴς ὁρᾶν}%
\end{groupe}%
%% v. 683
\begin{groupe}%
\typeVers{$2$ iamb.}\chantChoeur{\ligne{683}ἔσται γυναιξὶ καὶ βροτοῖς,\apparat{683}{ἐστὶν R : \lecon{ἔσται} Bothe % (\italique{erit} Divus) 
\espApp γυναιξὶν : \lecon{-ξὶ} Brunck}%
\commentaire{683}{γυναιξὶ καὶ βροτοῖς.}{L'on s'attendait à la formule consacrée : θεοῖσι καὶ βροτοῖς\marqueur{1}.}%
\debutNote{683}%
\note{1}{Homère, \italique{Iliade} XXIV, 553 ; Pindare, \italique{Pythiques} 3.30 ; Sophocle, \italique{Électre} v. 199.}}%
\end{groupe}%
%% v. 684-5
\begin{groupe}%
\decoupeVers{5}\typeVers{$2$ iamb.}\chantChoeur{ὅτι τά παράνομα τά τ᾽ ἀνόσια παρὼν}%
\end{groupe}%
%% v. 686
\begin{groupe}%
\typeVers{dochm...}\chantChoeur{\ligne{686}θεὸς παραχρῆμ᾽ ἀποτίνεται.\apparat{686}{ἀποτίνεται· | παραχρῆμά τε τίνεται R : \lecon{παραχρῆμ᾽ ἀποτίνεται} Hermann \ligneApp{687} (Κορ.) add. von Velsen}\\\passeLigne}%
\end{groupe}%
%% v. 687-688
\begin{groupe}%
\personnageGrec{Κορυφαῖος}\typeVers{tetr. tr\cat}\tetramTroch{\ligne{687}ἀλλ᾽ ἔοιχ᾽ ἡμῖν ἅπαντά πως διεσκέφθαι καλῶς.\\%\apparat{687}{Κορ. add. Velsen \espApp \lecon{πως} R : που Blaydes}\\
οὐχ ὁρῶμεν γοῦν ἔτ᾽ ἄλλον οὐδέν᾽ ἐγκαθήμενον\\}%
\end{groupe}%
%% v. 689-691
\begin{groupe}%
%
\personnageGrec{Μίκα}%
\typeVers{extr. metr.}\trimIamb{\ligne{689}ἆ ἆ}%
\end{groupe}%
%% v. 689-691
\begin{groupe}%
%\personnageGrec{Μίκα}%
\typeVers{trim. iamb.}\trimIamb{\ligne{689}ποῖ ποῖ σὺ φεύγεις; οὗτος οὗτος οὐ μενεῖς;\apparat{689}{ἀά supra lin. R : % ἔα Farraeus : ἆ ab init. v. 689 Bergler : 
\lecon{ἆ ἆ} Thiersch \espApp ποῖ R : \lecon{ποῖ ποῖ} Biset}%
\commentaire{689-761}{Parodie du \italique{Télèphe} d'Euripide.}{Comme dans ses \italique{Acharniens} représentés une quinzaine d'années plus tôt -- il y a de nombreuses analogies entre ces deux pièces --, Aristophane détourne le \italique{Télèphe} d'Euripide. On ne sait pas grand'chose de cette dernière pièce, la majorité des fragments conservés étant eux-mêmes tirés d'Aristophane et de ses anciens commentateurs\marqueur{1}.
Télèphe était un héros populaire puisque les trois grands tragiques lui ont consacré pas moins de sept pièces ; il fut blessé par un coup de lance d'Achille.
Comme sa blessure ne guérissait pas, il finit par consulter l'oracle d'Apollon, réponse de ce dernier : ὁ τρώσας ἰάσεται, « Qui t'a blessé te guérira ». Il se rend donc, déguisé en mendiant, dans le camp de l'armée achéenne préparant son expédition sur Troie. Allait-il d'abord trouver Clytemnestre qui lui confiait son petit Oreste pour appuyer son plaidoyer ? toujours est-il que Télèphe se retrouva la proie exposée à la colère des chefs grecs, il se réfugiait alors en suppliant sur l'autel en prenant en otage Oreste. Le sort de l'enfant était plus heureux que dans sa thesmophorienne parodie : grâce à l'entremise d'Agamemnon, Achille se laissait convaincre de soigner Télèphe et tout était bien qui finissait bien. Outre cette scène d'enlèvement du petit Oreste, de nombreux autres passages des \italique{Thesmophories} sont sans doute aussi inspirés du \italique{Télèphe} (v. 39-70 : où l'altercation entre le héros et un gardien attirait Clytemnestre à l'extérieur et permettait sa rencontre avec Télèphe ; v. 466-470 : lorsqu'au début de son plaidoyer, le héros -- sous une fausse identité -- usait de précautions oratoires ; v. 598-600, 659-667 avec l'annonce de la présence d'un espion dans le camp grec.)}%
\debutNote{689-761}%
\note{1}{Pour une présentation détaillée du \italique{Télèphe}, le lecteur pourra consulter l'édition de François Jouan \& Herman van Looy aux Belles Lettres, tome VIII, \nieme{3} partie.}\\
τάλαιν᾽ ἐγὼ τάλαινα, καὶ τὸ παιδίον\\
ἐξαρπάσας μοι φροῦδος ἀπὸ τοῦ τιτθίου.\\\passeLigne}%
\end{groupe}%
%% v. 692-695a
\begin{groupe}%
\personnageGrec{Μνησίλοχος}\trimIamb{κέκραχθι· τοῦτο δ᾽ οὐδέποτε σὺ ψωμιεῖς,\\
ἢν μή μ᾽ ἀφῆτ᾽· ἀλλ᾽ ἐνθάδ᾽ ἐπὶ τῶν μηρίων%
\apparat{693}{ἀφε\auDessus{ῖ}{ῆ}τ᾽ \Rdeux : ἀφεῖτ᾽ \Mon}%
\commentaire{693-5}{ἐπὶ τῶν μηρίων ... βωμόν.}{Vers tragiques, très sûrement empruntés à Euripide. (Note de Willems.)}\\
πληγὲν μαχαίρᾳ τῇδε φοινίας φλέβας\\
καθαιματώσει βωμόν.\\\passeLigne}%
\end{groupe}%
%% v. 695b-698
\begin{groupe}%
\personnageGrec{Μίκα}%
\trimIamb{\ligne{695}\decalage{καθαιματώσει βωμόν.}ὢ τάλαιν᾽ ἐγώ.\\
γυναῖκες, οὐκ ἀρήξετ᾽; οὐ πολλὴν βοὴν\testimonia{696-7}{(γυναῖκες -- τροπαῖον) \GregCor p.20-22 cf. schol.\expo{Ald} \italique{Plutus} 453c}\commentaire{696-7}{οὐ πολλὴν βοὴν / στήσεσθε καὶ τροπαῖον.}{Dans le sens de μήστωρ φόβοιο d'Homère. « Un vacarme à faire sauver ». (Musset, \italique{Confess}., p. 111 ; note de Willems.)}\\
στήσεσθε καὶ τροπαῖον, ἀλλὰ τοῦ μόνου\apparat{697}{καὶ om. R : καὶ suppl. Σ ad \italique{Pl.} v.453 et Greg. Cor. \italique{de dialectis} p.22}\\
τέκνου με περιόψεσθ᾽ ἀποστερουμένην;\\\passeLigne}%
\end{groupe}%
%% v. 699-701
\begin{groupe}%
%
\ledpb%
\personnageGrec{Χορός}\typeVers{προαναφώνημα}\chantChoeur{\decalage{aa}ἔα ἔα.%\commentaire{699-725}{Chant amébée.}{\eff{Retrouver dans Willems.}}
}%
\end{groupe}%
%% v. 699-701
\begin{groupe}%
\typeVers{2 dochm.}\chantChoeur{ὦ πότνιαι Μοῖραι τί τόδε δέρκομαι\apparat{700}{δὲ R : \lecon{τόδε} Dobree : δὴ Reisig}}%
\end{groupe}%
%% v. 699-701
\begin{groupe}%
\typeVers{dochm.}\chantChoeur{νεοχμὸν αὖ τέρας;\\\passeLigne}%
\end{groupe}%
%% v. 702-703
\begin{groupe}%
\typeVers{$4$ troch\cat}%
%\ledpb%
\tetramTroch{\Automatique{ὡς ἅπαντ᾽ ἄρ᾽ ἐστὶ τόλμης †πλέα† κἀναισχυντίας.}\testimonia{702}{(ἅπαντα pro ἅπαν γὰρ ἐστὶν -- κἀναισχυντίας) Suidas α 2902 (sed πλέα pro ἔργα) ; \italique{Corpus paroemiographorum Graecorum}, Appendices, I, 35 (sed ἔργα omisso) ; \italique{Laurentiamus} LVIII 24 p. 40 (ed. Cohn) (sed  δοῦλα pro ἔργα)}\apparat{702}{ἅπαν γὰρ ἐστὶν R : ἅπαντα Suid., cod. Laur. : \lecon{ἅπαντ᾽ ἄρ᾽ ἐστὶ} Bentley \espApp ἔργα R : \lecon{πλέα} Suid. : δοῦλα cod. Laur. : μεστὰ Porson}% \ligneApp{704} οἷον R : \lecon{οἷον;} Ellebodius \espApp ἐξάρξω R : %ἐξαράξει Ellebodius : 
%ἐξαράξω Bentley : \lecon{ἐξαρήξω} Bibi}% : \lecon{διπλὰ} Bibi}%
\commentaire{702}{ὡς ἅπαν γὰρ ἐστὶν τόλμης ἔργα κἀναισχυντίας.}{Le scholiaste met en avant le tour proverbial de la formule en commentant d'un sobre «~παροιμία~». Heureusement, d'autres glosateurs furent davantage prolixes. % ; à nous donc de mener une enquête plus approfondie !
%Pour ce faire, commençons par recenser quelques autres attestations de l'expression (ou de ses proches cousines). Et tout d'abord, remarquons que les termes τόλμη et ἀναισχυντία sont souvent mis en collocation : outre le présent passage, Platon\marqueur{1} et Démosthène\marqueur{2} associent les deux mots avec un sens analogue. Isocrate nous apporte quelque précision, l'on comprend en le lisant\marqueur{3} que l'audace (τόλμα) pouvait être interprétée comme un signe (σημεῖον) d'impudence (ἀναισχυντία). Les parœmiographes n'ont pas manqué de la relever, avec de légères variations. Les voici :
On retrouve la formule chez Suidas (α 2902) : 
\citComm{ἅπαντα τόλμης πλέα κἀναισχυντίας: ἐπὶ τῶν ἰταμῶν καὶ ἀναιδῶν.}
Et dans le \italique{Corpus paroemiographorum Graecorum}\marqueur{1} :
\citComm{ἅπαντα τόλμης καὶ ἀναισχυντίας: ὅτι ἡ ἀναισχυντία νικᾷ τὰ πολλά.}
Enfin dans le codex \italique{Laurentiamus} LVIII 24\marqueur{2}, l'on tombe sur la formulation suivante :%
\citComm{ἅπαντα τόλμης δοῦλα καὶ ἀναισχυντίας.}%
Ces trois leçons sont évidemment tirées du vers 702 des \italique{Thesmophories} et des anciens commentaires afférents. 
%Ces trois versions pourraient-elles découler d'une même source ? (savoir : l'\allemand{Urtext} des \italique{Thesmophories}.)
%Il me semble que oui... \latin{sed quod gratis asseritur gratis negatur}. Aussi nous faut-il présenter un embryon de `démonstration' ; ce qui devrait infailliblement nous conduire à formuler d'aventureuses conjectures. Partons donc à l'aventure ! \\%(Je prie le lecteur rigoriste de ne pas trop m'en tenir rigueur.)\\
Le début du vers dans le \italique{Ravennas}, ὡς ἅπαν γὰρ ... ἔργα κτλ., est fautif. La correction la plus simple à apporter est de lire : ὡς ἅπαντ᾽ ἄρ᾽ κτλ. %\marqueur{3} 
Il n'y aurait alors qu'une simple confusion d'un Τ lu en Γ : ΑΠΑΝΤΑΡ $ \to $ ΑΠΑΝΓΑΡ.
Cette correction mineure fait concorder le texte d'Aristophane avec les parémiographes. % sur la présence d'ἅπαντα. 
Les variantes ne différent plus que sur un seul mot : ἔργα (dans le \italique{Ravennas}) ; πλέα (\italique{Souda}) ; δοῦλα (Laurent.) ; $\emptyset$ (\italique{Corpus paroemiographorum}).

Rejetons d'emblée la leçon du \italique{Ravennas}, qui sent son interpolation ; ἔργα n'est qu'une glose d'ἅπαντα qui aura pénétré le texte.
La même erreur se produit au vers 717 des \italique{Oiseaux} : les manuscrits comportent `πρὸς ἅπαντα τρέπεσθε' ; une glose du scholiaste substitue ἔργα à ἅπαντα et Suidas (α 1638) finit par citer le vers ainsi : `πρὸς ἔργα τράπεσθε'.

La leçon de la \italique{Souda}, πλέα, est satisfaisante quant au sens (tout \italique{plein} d'effronterie et d'audace) ; elle n'est cependant pas métriquement correcte, raison pour laquelle Porson lui substitua μεστά. 
À défaut d'explication satisfaisante, j'ai préféré gardé la leçon fautive de Suidas, non sans la frapper de l'obèle du désespoir (†).

Finissons cette note aporétique en remarquant que les termes τόλμη et ἀναισχυντία sont souvent mis en collocation : Platon\marqueur{3} et Démosthène\marqueur{4} associent les deux mots avec un sens analogue. Isocrate nous apporte quelque précision, l'on comprend en le lisant\marqueur{5} que l'audace (τόλμα) pouvait être interprétée comme un signe (σημεῖον) d'impudence (ἀναισχυντία). Les Grecs avaient fait d'\italique{Implacabilité} un dieu : θεὸς ἡ Ἀναίδεια\marqueur{6}.
%À présent, si l'on fait l'hypothèse -- avec des `si', l'on scierait le texte entier !.. -- si donc l'on fait l'hypothèse que le texte original d'Aristophane contenait à cette place : διπλά, alors les différentes leçons pourraient s'expliquer.
%Très tôt, un manuscrit aura porté un δ surmonté de deux points pour abréger le δι\marqueur{7} de διπλά. L'omission des deux points donne δπλα ; la chute du δ pouvait alors entraîner la lecture πλέα tandis qu'un π cursif ressemblant au digramme ου faisait écrire à un scribe δοῦλα. 
%Si maintenant, sur une autre branche, un autre archétype voyait sa leçon ΔΙΠΛΑ réduite en ΔΙΠΑ (chute classique lorsque l'on se trouve en présence de deux lettres consécutives identiques ou ressemblantes), le Δ lu en Ε, le Ι en Ρ et le Π en Γ (toutes erreurs fréquentes prises individuellement, peut-être induites en l'occurrence par l'ἔργον du vers suivant), cette tradition aurait bien pu donner notre leçon ἔργα. 
%\citComm{ὡς ἅπαντ᾽ ἄρ᾽ ἐστὶ τόλμης διπλὰ κἀναισχυντίας.\\
%Ô combien toutes [les actions] sont doublées d'audace et d'impudence.}%
%L'on s'aperçoit à la traduction qu'on ressent l'envie de gloser ἅπαντα par ἔργα. Mettons-nous alors à la place d'un scribe ayant devant lui la leçon déjà corrompue de διπλὰ\marqueur{8}. Devant donc cette difficulté surmontée de la glose ἔργα, il aura choisi la \latin{lectio facilior} ; c'est-à-dire recopier ἔργα. C'est ainsi que je m'explique sa présence dans le \italique{Ravennas}.\\
%La prudence m'avait dans un premier temps incité à circonscrire cette conjecture à l'apparat critique. Las ! l'esprit d'imagination se laisse si facilement séduire qu'à la relecture διπλά rejoignit l'\allemand{Urtext}.
%\eff{Il ne me semble pas impossible que ces trois versions du proverbe -- avec ἔργα (notre texte, le \italique{Ravennas}), πλέα (\italique{Souda}) et δοῦλα (cod. Laur.) -- découlent d'une même source. Qu'il me soit permis de céder à la tentation des aventureuses conjectures, en espérant qu'une fois ne soit pas coutume. 
%Supposez, pour la démonstration, que le texte original d'Aristophane contenait διπλά. Très tôt, un manuscrit aurait porté un δ surmonté de deux points pour abréger ce δι (cf. le vers 122 et la correction proposée : δινεύματα). L'omission des deux points donne δπλα ; la chute du δ pouvait alors entraîner la lecture πλέα tandis qu'un π cursif ressemblant à ου faisait écrire à un scribe δοῦλα. Si maintenant, sur une autre branche, un autre archétype voyait sa leçon ΔΙΠΛΑ réduite en ΔΙΠΑ (chute classique lorsque l'on se trouve en présence de deux lettres consécutives identiques ou ressemblantes), le Δ lu en Ε, le Ι en Ρ et le Π en Γ (toutes erreurs fréquentes prises individuellement, peut-être induites en l'occurrence par l'ἔργον du vers suivant), cette tradition aurait bien pu donner notre leçon ἔργα. 
%\citComm{ὡς ἅπαντ᾽ ἄρ᾽ ἐστὶ τόλμης διπλὰ κἀναισχυντίας.\\
%Ô combien vont de pair audace et impudence.}
%Restons néanmoins prudent : si séduisante soit-elle pour l'esprit d'imagination, cette conjecture n'en restera pas moins circonscrite à l'apparat critique.\\
%Concluons cette note en digressant légèrement sur un détail amusant et plein de sens : les Grecs avaient fait d'\italique{Implacabilité} un dieu : θεὸς ἡ Ἀναίδεια\marqueur{9}.%
%P.-S. : faisons aussi remarquer au lecteur que l'émendation de Bentley -- ἅπαντ᾽ ἄρ᾽ ἐστὶ à la place de ἅπαν γὰρ ἐστὶ ; un simple Τ lu en Γ -- ne se justifie que si le scribe recopiait un manuscrit déjà élidé.}
%\eff{C'est le verbe apantaw ?
%\italique{Ion} (940) :
%\citComm{Πρεσβύτης -- Τίν'; Ὡς ἀπαντᾷ δάκρυά μοι τοῖς σοῖς λόγοις.}
%Ménandre 706(483) K.-A.
%Loci communes (Capita theologica id est scite dicta atque electa, ex diversis, tum Christianorum, tum profanorum libris (PG 91, 721–1018).)
%κεφάλαια θεολογικὰ ἤτοι ἐκλογαὶ
%ΛΟΓΟΣ Β´.
%Περί φρονήσεως καί βουλῆς. 
%Μενάνδρου. Ἅπαντα δοῦλα τοῦ φρονεῖν καθίσταται.
%Βουλευόμενοι παραδείγματα ποιοῦνται παρεληλυθότα τῶν μελλόντων. Τό γάρ ἀφανές, ἐκ τοῦ φανεροῦ ταχίστην ἔχει τήν διάγνωσιν. 
%"Loci communes" du Pseudo-Maxime le Confesseur. Éditions :
%Florilège sacro-profane du Pseudo-Maxime (2001), Sargologos, Étienne
%IHM, Sibylle : "Ps.-Maximus Confessor. Erste kritische Edition
%einer Redaktion des sacro-profanen Florilegiums Loci Communes" (2001)}
}%
\debutNote{702}%
\note{1}{\italique{Corpus paroemiographorum Graecorum}, vol. I, Appendices, I, 35, p. 384 ($=$ Vat. K II 36).}%
\note{2}{Il s'agit d'un manuscrit conservé à la Biblioteca Medicea Laurenziana (BML) de Florence, dépôt Plutarque, sous la cote 58. 24 %(J'en profite pour signaler au lecteur un outil sur l'internet fort pratique pour tout ce qui touche aux manuscrits : Πίνακες. Il y trouvera, par exemple, l'historique de notre manuscrit ainsi que sa numérisation.) Je me reporte à l'édition qu'en fit 
et édité par Leopold Cohn en 1887 : \italique{Zu den Paroemoigraphen} (notre proverbe y figure à la page 40).}%
%\note{3}{Cette émendation, proposée par Bentley et unanimement adoptée se justifierait d'autant mieux si le scribe du \italique{Ravennas} recopiait un manuscrit déjà élidé -- du moins partiellement, le modèle restant majoritairement à \latin{scriptio plena}.}%
\note{3}{Platon, \italique{Apologie de Socrate}, 38d. Voir aussi son \italique{Théétète} (196d) où le tour proverbial est encore plus sensible : ὅμως δέ — πάντα γὰρ τολμητέον — τί εἰ ἐπιχειρήσαιμεν ἀναισχυντεῖν; «~De l'audace ! encore de l'audace ! toujours de l'audace !~»}%
\note{4}{Παραγραφὴ πρὸς Ναυσίμακον καὶ Ξενοπείθην, 5 (Νόμος).}%
\note{5}{Περὶ ἀντιδόσεως, 190 : ἔτι δὲ τὴν τόλμαν μὴ τὴν ἀναισχυντίας σημεῖον γιγνομένην, ἀλλὰ τὴν μετὰ σωφροσύνης. Voir aussi Démosthène, \italique{In Aristogitonem 1}, 68.8 : ταῦτα τολμᾷ λέγειν δι' ἀναισχυντίαν.}%
%\note{7}{Cf. l'article de Victor Coulon, «~Notes critiques et exégétiques sur Aristote, \italique{Poétique} , 1450a12, 1450a29, sur Thucydide I, 1, 2 et sur cinq passages d'Aristophane~», in \italique{Revue des Études Grecques} (1937), t. 50 fasc. 238, p. 458. (À propos du vers 122 et de la correction de Bentley δινεύματα à la place de διανεύματα R.)}%
%\note{8}{Et donc incompréhensible ; cf. la version courte ἅπαντα τόλμης καὶ ἀναισχυντίας citée plus haut : le scribe du \italique{Corpus} a tout simplement décidé de ne pas recopier le mot qu'il n'arrivait pas à lire.}%
\note{6}{\italique{Implacabilité} et non \italique{Impudence}. Cf. le \italique{Dictionnaire des antiquités grecques et romaines} (Daremberg \& Saglio), t. I, p. 398b (\latin{Areopagus}). Les Athéniens avaient dressé des autels à Ὕβρις et Ἀναίδεια dans l'Aréopage, contre lesquels s'était insurgé ce rabat-joie de Cicéron -- non sans faire un contre-sens comme le souligne l'auteur de l'article. Dans sa querelle avec Agamemnon, l'attitude d'Achille -- qui par bien des aspects est celle d'un adolescent attardé -- nous révèle en fait son côté divin : il est \italique{implacable}, tout autant que le destin qu'il connaît et a accepté. Pour revenir au proverbe `θεὸς ἡ Ἀναίδεια', certainement inspiré par ces autels de l'Aréopage, voir Zénobius (IV, 36). Érasme l'a aussi recensé dans ses \italique{Adages} (n° 2063).}\\
οἷον αὖ δέδρακεν ἔργον, οἷον αὖ φίλαι τόδε.\\\passeLigne}%
\end{groupe}%
%% v. 704
\begin{groupe}%
\personnageGrec{Μνησίλοχος}%
\tetramTroch{οἷον; ὑμῶν ἐξαρήξω τὴν ἄγαν αὐθαδίαν.\apparat{704}{οἷον R : \lecon{οἷον;} Ellebodius \espApp ἐξάρξω R : ἐξαράξει Ellebodius : ἐξαράξω Bentley : \lecon{ἐξαρήξω} Bibi}%
\commentaire{704}{οἷον ὑμῶν ἐξάρξω τὴν ἄγαν αὐθαδίαν.}{Il manque une syllabe à ἐξάρξω, la leçon du manuscrit de Ravenne. Les différents éditeurs y suppléent avec le verbe ἐξαράττειν (ἐξαράξω au futur ou, en s'éloignant davantage de la leçon -- quitte à corriger le texte... --, ἐξαράξει, pour faire de οἷον le sujet). La correction est légère et pourtant elle ne manque pas de laisser dubitatif : si l'image de \italique{briser} (ἐξαράξω) l'\italique{audace} (αὐθαδίαν) de \italique{quelqu'un} (ὑμῶν) ne paraît aucunement étrangère à une imagination française, il n'en allait peut-être pas de même pour les Grecs. L'incontournable Taillardat\marqueur{1} l'avait déjà noté : considérer l'arrogance «~comme un objet matériel qu'on peut briser [...] est une image unique~». En revanche, en grec comme en français, l'image de \italique{chasser}, \italique{pousser} hors (ἐξ-), \italique{repousser} un sentiment ou un trait de caractère est courante. Ainsi, l'on chassera l'audace par l'audace\marqueur{2} :%
\citComm{Ἀντιφάνης δέ· οἴνῳ .... τὸν οἶνον ἐξελαύνειν, [...] τριωβόλῳ δὲ πόρνην, αὐθαδίαν αὐθαδίᾳ.}%
L'on se débarrassera de son chagrin\marqueur{3} :%
\citComm{ἐκκρούει τὴν λύπην}%
comme les vieillards cagneux secouent leurs ennuis chez Aristophane\marqueur{4} :%
\citComm{ἀποσείονται δὲ λύπας.}%
Euripide met dans la bouche du fourbe Polydectès qui veut s'accaparer Danaé que\marqueur{5} :%
\citComm{πατέρα τε παισὶν ἡδέως συνεκφέρειν\\
φίλους ἔρωτας \textbf{ἐκβαλόντ' αὐθαδίαν}.\\
Un père doit supporter avec indulgence les affaires amoureuses de son fils et \italique{repousser l'entêtement}.}
L'occasion fait le larron : profitons-en pour placer cet aphorisme rapporté par Athénée\marqueur{6} illustrant le sens de la mesure grec, même quand il s'agit de boire...%
\citComm{Πάσας δ᾽ ἐκ κραδίας ἀνίας ἀνδρῶν ἀλαπάζει\\
πινόμενος κατὰ μέτρον· ὑπὲρ μέτρον δὲ χερείων.\\
\francais{Le vin est souverain pour chasser les soucis du cœur de qui n'en boit pas trop, mais l'excès nuit.}}%
(Un franc-buveur espiègle répondrait que la mesure du boire est de boire sans mesure\marqueur{7} !..)
Le verbe ἀρήγω (futur : ἀρήξω) signifie \italique{secourir} et, par suite, \italique{écarter}, \italique{repousser}. Il conviendrait tant pour le sens que paléographiquement s'il était préfixé de ἐξ-.
Comme pour nous conforter, Hésychius (ε 3599) nous atteste justement l'existence du composé :%
\citComm{ἐξαρῆξαι· ἐκφορῆσαι.}%
Lisons donc :%
\citComm{ὑμῶν ἐξαρήξω τὴν ἄγαν αὐθαδίαν\\
Je vous extirperai votre trop-plein d'audace.}
P.-S. : l'idée de recourir aux matrices métaphoriques pour le bon établissement d'un texte revient à l'immense Jean Taillardat. Je ne peux que renvoyer à son coruscant article « Images et matrices métaphoriques »\marqueur{8}. Il le termine par ce vœu :%
\citComm{En guise de conclusion, je formulerai un souhait, celui qu'on puisse, un jour, consulter un dictionnaire des matrices métaphoriques pour les langues indo-européennes vivantes (langues populaires et argots compris) et mortes. On l'imaginerait volontiers comme une sorte de dictionnaire analogique, mais réduit aux images, les notions à exprimer (par exemple « colère », « inconsistance », « jeunesse », etc.) étant classées en ordre alphabétique ; quel merveilleux outil serait ainsi offert à qui s'intéresse à la recherche étymologique et combien de tâtonnements, d'erreurs même, seraient alors évités ! Mais quelle équipe aura assez d'abnégation et de courage pour se lancer jamais dans cette énorme entreprise ?}%
Puisse ce grand et beau dictionnaire un jour voir le jour !%
}%
\debutNote{704}%
\note{1}{§628 de ses \italique{Images d'Aristophane}, ouvrage dont on ne soulignera jamais assez l'importance pour toute personne désireuse d'avoir un avant-goût du génie de la langue aristophanesque.}%
\note{2}{Athénée, \italique{Deipnosophistes}, II, 44a.}%
\note{3}{Aristote, \italique{Éthique à Nicomaque, 1154a27}.}%
\note{4}{\italique{Grenouilles}, v. 346.}%
\note{5}{Fragment du \italique{Dictys}. Stobée, IV, 26, 16 $=$ fr. 7 \italique{Dictys} Jouan \& van Looy $=$ fr. 339 Kannicht $=$ fr. 433 Mette.}
\note{6}{Athénée, II, 37b.}%
\note{7}{Aristophane, \italique{Guêpes}, vers 79-80 :%
\citComm{εἶναι φιλοπότην αὐτόν. [ ... ]\\
αὕτη γε χρηστῶν ἐστιν ἀνδρῶν ἡ νόσος.}.}%
\note{8}{\italique{Bulletin de l'Association Guillaume Budé} (décembre 1977), Lettre d'humanité n° 36, pp. 344-354.}%
%\note{9}{Dictionnaire dont le grand'œuvre de Taillardat, ses \italique{Images d'Aristophane}, en est la première pierre... et pourrait en être la pierre d'angle !}%
\\\passeLigne}%
\end{groupe}%
%% v. 705
\begin{groupe}%
\personnageGrec{Χορός}%
\tetramTroch{\Automatique{ταῦτα δῆτ᾽ οὐ δεινὰ πράγματ᾽ ἐστὶ καὶ περαιτέρω;}\\\passeLigne}%
\end{groupe}%
%% v. 706
\begin{groupe}%
\personnageGrec{Μίκα}%
\tetramTroch{\small{δεινὰ δῆθ᾽, ὁτιή γ᾽ ἔχει μου 'ξαρπάσας τὸ παιδίον.}\apparat{706}{ὅτι R : \lecon{ὁτιή} Faber : ὅστις Porson}\\\passeLigne}%
\end{groupe}%
%% v. 707-708
\begin{groupe}%
\personnageGrec{Χορός}%
\typeVers{dim. anap.}\dimAnap{τί ἂν οὖν εἴποι πρὸς ταῦτά τις, ὅτε\\
τοιαῦτα ποιῶν ὅδ᾽ ἀναισχυντεῖ;%
\commentaire{707-25}{}{Ce passage, sans correspondre parfaitement du point de vue métrique à celui des vers 667-686, présente néanmoins de troublantes ressemblances avec ce dernier. Se reporter à Lætitia\marqueur{1} pour les parallélismes qu'on peut faire avec le texte tel qu'il nous a été conservé.}%
\debutNote{707-25}%
\note{1}{\italique{The songs of Aristophanes}, pp. 421 sqq. Je reprends pour une bonne part sa colométrie.}\\\passeLigne}%
\end{groupe}%
%% v. 709
\begin{groupe}%
\personnageGrec{Μνησίλοχος}%
\typeVers{paroem.}\dimAnap{κοὔπω μέντοι γε πέπαυμαι.\\\passeLigne}%
\end{groupe}%
%% v. 710
\begin{groupe}%
\personnageGrec{Χορός}%
\typeVers{paroem.}\chantChoeur{ἀλλ᾽ οὖν ἥκειν γ᾽ ὅθεν ἥκεις.\testimonia{710-3}{Suidas η 170 (sed λήψῃ pro λήψει)}\testimonia{710}{Photius η 88 ἥκειν· ἀντὶ τοῦ ἥκων}\apparat{710}{ἥκεις τ᾿ R : ἥκεις γ᾽ Suid. η 170 : ἥκων γ᾽ Willems : \lecon{ἥκειν γ᾽} Bibi ap. Phot. η 88 \ligneApp{711} τ᾿ R, Suid : \lecon{δ᾿} Bentley \espApp οὐ λέξεις R, Suid. : \lecon{οὐκ ἔξει.} Bibi ap. Σ}%
\commentaire{710-2}{ἀλλ᾽ οὖν ἥκεις τ᾿ ὅθεν ἥκεις.}{Je reproduis ci-dessous la \nieme{21} note d'Alphonse Willems. La correction qu'il y propose fut sauvée du néant de l'oubli et adoptée par la plupart des éditeurs modernes grâce à la diligente entremise de Coulon :

«~Ce passage a été refait de cinq ou six façons, plus bizarres les unes que les autres. 
La correction me paraît élémentaire. Les manuscrits concordent, sauf sur le premier ἥκεις (ἥκεις γ' dans Suidas, ἥκεις τ' dans le \italique{Ravennas}, ἥκισθ' dans le manuscrit de Munich [qui n'est qu'une copie du manuscrit de Ravenne, comportant quelques émendations parfois bien senties]). Lisez ἥκων γ', le sens devient parfait :
\citComm{ἀλλ᾽ οὖν ἥκων γ᾽ ὅθεν ἥκεις}
Il est à peine besoin d'ajouter que ἥκων est pris dans le sens de ἀνήκων, \italique{redux}, ce qui est des plus fréquents\marqueur{1}. Et quant à ἀλλ᾽ οὖν..γε, chacun sait qu'il signifie \italique{toujours}, \italique{tout au moins}, \italique{en tous cas}\marqueur{2}.\\
%\eff{Traduisez : 
%\citComm{\francais{Toujours tu n'iras pas, revenu au lieu dont tu es parti, après une fuite facile, raconter ton crime et ton évasion...}}}
% [...]\marqueur{3}\\
C'est le mot d'Épicharme, fr. 126\marqueur{3} : κἀπῆνθεν ὅθεν ἦνθεν πάλιν ; ou de Ménandre (t. IV, p. 211 [$=$ fr. 373 K.-A.]) : ἀπῆλθεν ὅθεν ἦλθεν ταχύ. Cf. aussi Soph., \italique{Œd. Col.}, 1227 : βῆναι κεῖσ' ὁπόθεν περ ἥκει ; Eur., \italique{Her. F.}, 1247 : θανών, ὅθενπερ ἦλθον, εἶμι γῆς ὕπο, etc.~»

Le chœur enjoint Mnésiloque de cesser cette escalade de l'impudence : qu'il quitte au moins l'autel et les rejoigne\marqueur{4}.\\
À la lumière de la démonstration de Willems, je me permets d'ajouter la glose suivante du \latin{Lexicon} de Photius (η 88) :%
\citComm{ἥκειν· ἀντὶ τοῦ ἥκων.}%
Supposer que la glose de Photius se rapporte à notre vers offre l'avantage de rester plus proche du texte des manuscrits : il y aurait eu simple permutation d'un ν par un ς.
L'on peut donc lire après Willems :%
\citComm{ἀλλ᾽ οὖν ἥκειν γ᾽ ὅθεν ἥκεις.\\
\francais{Au moins, de revenir d'où tu es parti.}}%
}%
\debutNote{710-2}%
\note{1}{\italique{Paix}, 265, 275; \italique{Lysistr}., 595, 731, 924 ; \italique{Gren}., 503 ; \italique{Ploutos}, 521, 1179, etc.}%
\note{2}{\italique{Nuées}, 985, 1002 ; \italique{Guêpes}, 1129, 1190 ; Esch., \italique{Prom}., 1058 ; Soph., \italique{El}., 233 ; \italique{Ant}., 84, etc. M'éloignant de Willems quant au sens à donner à ce vers, je ne retranscris pas ici la traduction qu'il en propose. Le lecteur la trouvera en annexe où toutes les notes d'Alphonse Willems sur les \italique{Thesmophories} sont intégralement reproduites.}%
\note{3}{$[=$ Plutarque, \italique{Moralia}, Παραμυθητικὸς πρὸς Ἀπολλώνιον, 110a $=$ fr. 213(245) Epich. Kassel-Austin, \italique{Poet. Com. Gr.}, I, \italique{Comoedia Dorica}.$]$}%
\note{4}{C'est l'explication de la scholie aussi conservée par Suidas (η 170) : ἥκεις φησίν εἰς τὰς χεῖρας ἡμῶν, ὅθεν οὐ φεύξῃ.}%
}%
\end{groupe}%
%% v. 711-713
\begin{groupe}%
\typeVers{paroem.}\chantChoeur{φαύλως δ᾿ ἀποδρὰς οὐκ ἔξει.%
\commentaire{711-2}{φαύλως τ᾽ ἀποδρὰς οὐ λέξεις / οἷον δράσας διέδυς ἔργον.}{%. ARISTOPHANES Comic. Acharnenses {0019.001} Line 215
%εἴ τις οἶδ’ ὅποι τέτραπται γῆς ὁ τὰς σπονδὰς φέρων.
%Ἐκπέφευγ’, οἴχεται φροῦδος. Οἴ-    (208-209)
%μοι τάλας τῶν ἐτῶν τῶν ἐμῶν·   (210)
%οὐκ ἂν ἐπ’ ἐμῆς γε νεότητος, ὅτ’ ἐ-   (211-212)
%γὼ φέρων ἀνθράκων φορτίον    (213-214)
%ἠκολούθουν Φαΰλλῳ τρέχων, ὧδε φαύλως ἂν ὁ    (215)
%σπονδοφόρος οὗτος ὑπ’ ἐμοῦ τότε διωκόμενος    (216-217)
%ἐξέφυγεν οὐδ’ ἂν ἐλαφρῶς ἂν ἀπεπλίξατο.    (218)
%
%νῦν δ᾽ ἐπειδὴ στερρὸν ἤδη τοὐμὸν ἀντικνήμιον,
%220καὶ παλαιῷ Λακρατείδῃ τὸ σκέλος βαρύνεται,
%οἴχεται. διωκτέος δέ: μὴ γὰρ ἐγχάνῃ ποτὲ
%μηδέ περ γέροντας ὄντας ἐκφυγὼν Ἀχαρνέας.
%
%ὅστις ὦ Ζεῦ πάτερ καὶ θεοὶ τοῖσιν ἐχθροῖσιν ἐσπείσατο,
%οἷσι παρ᾽ ἐμοῦ πόλεμος ἐχθοδοπὸς αὔξεται τῶν ἐμῶν χωρίων:
%230κοὐκ ἀνήσω πρὶν ἂν σχοῖνος αὐτοῖσιν ἀντεμπαγῶ
%ὀξὺς ὀδυνηρὸς ... ἐπίκωπος, ἵνα
%μήποτε πατῶσιν ἔτι τὰς ἐμὰς ἀμπέλους.
%
%ἀλλὰ δεῖ ζητεῖν τὸν ἄνδρα καὶ βλέπειν βαλληνάδε
%235καὶ διώκειν γῆν πρὸ γῆς, ἕως ἂν εὑρεθῇ ποτέ:
%ὡς ἐγὼ βάλλων ἐκεῖνον οὐκ ἂν ἐμπλῄμην λίθοις.
%
%
%
%
%
%Βδ.        ἀπολεῖς με. ποίας χάρακας; ἀλλ’ ὡς ἢ κάπρον
%ἐδιώκαθές ποτ’ ἢ λαγών, ἢ λαμπάδα
%ἔδραμες, ἀνευρὼν ὅ τι νεανικώτατον.
%Φι.        ἐγᾦδα τοίνυν τό γε νεανικώτατον· @1   (1205)
%ὅτε τὸν δρομέα Φάυλλον ὢν βούπαις ἔτι
%εἷλον διώκων λοιδορίας ψήφοιν δυοῖν.
%Βδ.        παῦ’· ἀλλὰ δευρὶ κατακλινεὶς προσμάνθανε
%ξυμποτικὸς εἶναι καὶ ξυνουσιαστικός.
%Φι.        πῶς οὖν κατακλινῶ; φράζ’ ἁνύσας.    (1210)
%Βδ.                                        εὐσχημόνως.   (1210)
%
%
%
%faulws logizw : vers 656
% Φι.                                        παῦσαι καὶ μὴ πατέριζε.   (652)
%εἰ μὴ γάρ, ὅπως δουλεύω ’γώ, τουτὶ ταχέως με διδάξεις,
%οὐκ ἔστιν ὅπως οὐχὶ τεθνήσει, κἂν χρῇ σπλάγχνων μ’
%ἀπέχεσθαι.    (654)
%Βδ.        ἀκρόασαί νυν, ὦ παπίδιον, χαλάσας ὀλίγον τὸ μέτωπον. @1   (655)
%καὶ πρῶτον μὲν λόγισαι φαύλως, μὴ ψήφοις ἀλλ’ ἀπὸ
%χειρός,    (656)
%τὸν φόρον ἡμῖν ἀπὸ τῶν πόλεων συλλήβδην τὸν προσιόντα,
%κἄξω τούτου τὰ τέλη χωρὶς καὶ τὰς πολλὰς ἑκατοστάς,
%πρυτανεῖα, μέταλλ’, ἀγοράς, λιμένας, μισθώσεις, δημιό-
%πρατα·    (659)
%
%
%REFAIRE NOTE !!!!! Il n'y a pas lieu de corriger le texte ! (même si les précédents éditeurs disaient effectivement n'importe quoi. Il s'agit de Phaüllos ! un célèbre coureur cf Acharniens 215 , Guêpes 1206 (voir Dargemberg et Saglio, Realencyclopädie). Traduire : "Tu ne diras pas avoir décampé à la manière de Phaüllos" ?)
% Voir Nuées 167-168
Où se trouve le complément réclamé par λέγειν ? Je ne suis pas sûr qu'il soit de bonne grécité de construire la phrase ainsi %que tous les éditeurs 
: λέξεις οἷον δράσας, «~tu raconteras tout ce que tu as fait~». (Sauf erreur de ma part, le grec aurait recouru à un verbe conjugué et non au participe de δράω.)
%Une retouche du texte -- très légère, accordez-moi au moins cela avant de crier au pourfendeur des éditeurs-tripatouilleurs pris lui-même en flagrant délit de bidouillage --, 
Une légère retouche du texte permettrait une bien meilleure lecture, tant pour la syntaxe grecque qu'au niveau du sens.
Lisez : 
\citComm{φαύλως γ᾽ ἀποδρὰς οὐκ ἔξει.\\
\francais{Ce n'est pas avec une fuite facile que tu t'en iras !}}
(Confusion entre un Κ et un Λ.)
Le vers 712, οἷον δράσας διέδυς ἔργον, est en fait apposé au vers 713 et en explique la sentence : λήψει δὲ κακόν, « pour tout ce que tu as fait, tu recevras châtiment ».
%Mais es raisons exposées jusqu'ici sont trop maigres pour justifier une émendation du texte\marqueur{1}. 
%En voici une autre. 
Pour justifier cette correction, je m'appuierai sur un débris de scholie qui se trouve à l'intérieur même du texte du \italique{Ravennas}, et que j'interprète ainsi :
\citComm{ἥκεις, φησίν, εἰς τὰς χεῖρας ἡμῶν, ὅθεν οὐ φεύξῃ.\\
Tu [re]viens, dit-elle, dans nos bras, et d'où tu ne fuiras pas.}
Comprise ainsi, cette scholie est clairement une paraphrase du vers 710 pour la première partie (ἥκεις εἰς τὰς χεῖρας ἡμῶν $=$ ἥκειν γ᾽ ὅθεν ἥκεις). C'est la seconde partie de la scholie qui nous intéresse : ὅθεν οὐ φεύξῃ, «~d'où tu ne fuiras pas~». Peut-on espérer meilleure glose que `οὐ φεύξῃ' pour le `οὐκ ἔξει' que nous rétablissons ?%\\
% et du vers 711 pour la seconde (ὅθεν οὐ φεύξῃ $=$ οὐκ ἔξει), du moins tels que je les lis.
%L'autre scholie, encore plus abîmée, est extérieure ; elle se trouve dans les marges du manuscrit :
%\citComm{ἐκ τῶν ἀνδρῶν φησίν [bis repetita : je lis σφῇσιν] ὅθεν οὐκ ἔξεστιν}
%L'on comprend que ce scholiaste interprétait différemment le vers 710 -- d'accord en cela avec Willems : ὅθεν ἥκεις, "d'où tu viens", sous-entendu ἐκ τῶν ἀνδρῶν, "de chez les hommes", jusqu'à nous (σφῇσιν). La fin de la scholie, ὅθεν οὐκ ἔξεστιν, n'offre aucun sens acceptable. Il faut bien-sûr lire : 
%\citComm{ὅθεν οὐκ ἔξεσθα\\
%Et d'où tu ne sortiras pas...}
%Je plie les scholies pour les accorder à mes vues sur le texte, me direz-vous. Peut-être, mais convenez que les corrections sont minimes par rapport au gain en intelligibilité.%
}%
\debutNote{711-2}%
%\note{1}{D'autant que le manuscrit de Ravenne concorde parfaitement à ce vers avec la \italique{Souda} qui le cite.}%
%\note{2}{Il y a nombre de σφῇσιν à rétablir dans les commentaires là où des scribes peu attentifs auront laissé le trop commun φησίν. (Nous en avons déjà croisé un sur notre route chez Eustathe de Thessalonique, cf. la note au vers 643.) L'erreur s'explique d'autant plus aisément lorsque -- comme c'est le cas ici -- le mot précédent se termine par un σ.}%
%\commentaire{711}{φαύλως,}{\eff{C'est bien en mauvaise part qu'il faut prendre ici φαύλως : \italique{injurieusement}, \italique{sans façon}, \italique{avec mépris}, \italique{grossièrement}.
%Aristophane l'emploie une fois ailleurs avec un verbe qui se rapporte à la course (\italique{Acharniens}, vers 215) :
%\citComm{οὐκ ἂν ἐπ᾽ ἐμῆς γε νεότητος, ὅτ᾽ ἐ-\\
%γὼ φέρων ἀνθράκων φορτίον\\
%ἠκολούθουν Φαΰλλῳ τρέχων, ὧδε \textbf{φαύλως} ἂν ὁ\\
%σπονδοφόρος οὗτος ὑπ᾽ ἐμοῦ τότε διωκόμενος\\
%ἐξέφυγεν οὐδ᾽ ἂν ἐλαφρῶς ἂν ἀπεπλίξατο.\\
%Ce n'est pas dans mon jeune âge, quand avec une charge de charbons je luttais de vitesse avec Phaÿllos, que ce portepaix aurait si φαύλως échappé à ma poursuite, qu'il se serait si prestement dérobé.}}
%Cf. les Guêpes et Phaÿllos : pourquoi est-il accusé d'outrage ? est-ce à cause du φαύλως ? ou bien parce que Phaÿllos était proverbialement réputé parler gauloisement ? (cf. le proverbe "passer les bornes", recensé par Suidas et Érasme) ou allusion à un Phaÿllos contemporains ?
%}%
}%
\end{groupe}%
%% v. 710
\begin{groupe}%
\typeVers{2 anap.}\chantChoeur{οἷον δράσας διέδυς ἔργον·}%
\end{groupe}%
%% v. 710
\begin{groupe}%
\typeVers{anap.}\chantChoeur{λήψει δὲ κακόν.\\\passeLigne}%
\end{groupe}%
%% v. 714
\begin{groupe}%
\personnageGrec{Μνησίλοχος}%
\typeVers{$4$ troch\cat}\tetramTroch{τοῦτο μέντοι μὴ γένοιτο μηδαμῶς, ἀπεύχομαι.\\\passeLigne}%
\end{groupe}%
%% v. 715-716
\begin{groupe}%
\personnageGrec{Χορός}%
\typeVers{2 dochm.}\diDoch{τίς οὖν σοι, τίς ἂν σύμμαχος ἐκ θεῶν\apparat{715}{\lecon{οὖν σοι} R : ἂν σοι Hamaker}\\
ἀθανάτων ἔλθοι ξὺν ἀδίκοις ἔργοις;\\\passeLigne}%
\end{groupe}%
%% v. 717
\begin{groupe}%
\personnageGrec{Μνησίλοχος}%
\typeVers{trim. ia\cat}\trimIambCat{μάτην λαλεῖτε· τὴνδ᾽ ἐγὼ οὐκ ἀφήσω.\apparat{717}{τὴν δ᾽ R : \lecon{τήνδ᾽} Biset}\\\passeLigne}%
\end{groupe}%
%% v. 718-725
\begin{groupe}%
%
\personnageGrec{Χορός}%
\typeVers{2 iamb.}\chantChoeur{ἀλλ᾽ οὐ μὰ τὼ θεὼ τάχ᾽ οὐ\\
χαίρων ἴσως ἐνυβριεῖς\apparat{719}{ἐνυβρίσεις R : \lecon{ἐνυβριεῖς} anonymus Romanus}\\
λόγους τε λέξεις ἀνοσίους}
\end{groupe}%
\endnumbering
\end{Leftside}
\begin{Rightside}
\beginnumbering
%% v. 663-666
\begin{groupe}%
\chantChoeurFr{S'ils se comportent autrement, voici ce qui se \fin passera :}%
\end{groupe}%
%% v. 663-666
\begin{groupe}%
\chantChoeurFr{celui d'eux qui sera pris à faire un sacrilège,}%
\end{groupe}%
%% v. 663-666
\begin{groupe}%
\chantChoeurFr{enflammé de fureur, délirant de rage,}%
\end{groupe}%
%% v. 663-666
\begin{groupe}%
\chantChoeurFr{témoignera aux femmes et à tous les mortels}%
\end{groupe}%
%% v. 663-666
\begin{groupe}%
\chantChoeurFr{que la divinité châtie les atteintes aux lois }%
\end{groupe}%
%% v. 663-666
\begin{groupe}%
\chantChoeurFr{et à la piété et ne se fait pas attendre.\\\passeLigne}%
\end{groupe}%
%% v. 687-688
\begin{groupe}%
\personnageFr{Coryphée}%
\tetramTrochFr{\ligne{687}Mais il nous semble avoir bien cherché partout et nous ne voyons pas d'autre homme caché parmi nous.
(Mnésiloque enlève l'enfant et se réfugie sur l'autel.)\\\passeLigne}%
\end{groupe}%
%% v. 689-691
\begin{groupe}%
\personnageFr{Mika}%
\trimIambFr{Ah ! Ah !}%
\end{groupe}%
%% v. 689-691
\begin{groupe}%
%\personnageFr{Mika}%
\trimIambFr{Où ? Où te sauves-tu ? holà ! holà ! Arrête. \\
Malheur ! Ah ! malheur ! le voilà parti avec mon enfant\\
arraché de ma mamelle.\\\passeLigne}%
\end{groupe}%
%% v. 692-695a
\begin{groupe}%
\personnageFr{Mnésiloque}%
\trimIambFr{Hurle. Jamais plus tu ne l'appateras,
si vous ne me laissez pas libre. Mais ici même « sur ces oblations,\\
frappé de ce couteau, du pur sang de ses veines\\
il rougira l'autel ».\\\passeLigne}%
\end{groupe}%
%% v. 695b-698
\begin{groupe}%
\personnageFr{Mika}%
\trimIambFr{\ligne{695}\espace{3cm}Malheur à moi !\\
Femmes, au secours ! Poussez une puissante clameur,\\
une clameur à faire lâcher pied. Souffrirez-vous\\
qu'on me ravisse mon unique enfant?\\\passeLigne}%
\end{groupe}%
%% v. 699-701
\begin{groupe}%
\personnageFr{Chœur}\chantChoeurFr{Ah! Ah !}%
\end{groupe}%
%% v. 699-701
\begin{groupe}%
\chantChoeurFr{Augustes M\oe res, quel prodige}%
\end{groupe}%
%% v. 699-701
\begin{groupe}%
\chantChoeurFr{nouveau vient frapper mes yeux !\\\passeLigne}%
\end{groupe}%
%% v. 702-703
\begin{groupe}%
\tetramTrochFr{Tout n'est donc qu'audace doublée d'effronterie !\\
Quel forfait encore, amies, quel forfait il a commis !\\\passeLigne}%
\end{groupe}%
%% v. 704
\begin{groupe}%
\personnageFr{Mnésiloque}%
\tetramTrochFr{Quel forfait ? je vous extirperai votre outrecuidance !\\\passeLigne}%
\end{groupe}%
%% v. 705
\begin{groupe}%
\personnageFr{Chœur}%
\tetramTrochFr{N'est-ce pas une abomination, et pis encore ?\\\passeLigne}%
\end{groupe}%
%% v. 706
\begin{groupe}%
\personnageFr{Mika}%
\trimIambFr{Une abomination, oui, qu'on m'ait arraché mon \fin enfant.\\\passeLigne}%
\end{groupe}%
%% v. 707-708
\begin{groupe}%
\personnageFr{Chœur}%
\dimAnapFr{Que dire après cela, quand\\
l'auteur d'une telle action a perdu toute honte ?\\\passeLigne}%
\end{groupe}%
%% v. 709
\begin{groupe}%
\personnageFr{Mnésiloque}%
\dimAnapFr{Et je n'ai pas fini.\\\passeLigne}%
\end{groupe}%
%% v. 710
\begin{groupe}%
\personnageFr{Chœur}%
\chantChoeurFr{Au moins, reviens d'où tu es parti ; ce n'est pas}%
\end{groupe}%
%% v. 711-713
\begin{groupe}%
\chantChoeurFr{ avec une fuite facile que tu t'en sortiras.}%
\end{groupe}%
%% v. 710
\begin{groupe}%
\chantChoeurFr{Avec tout ce que tu as fait en t'immisçant ici,}%
\end{groupe}%
%% v. 710
\begin{groupe}%
\chantChoeurFr{malheur t'en prendra !\\\passeLigne}%
\end{groupe}%
%% v. 714
\begin{groupe}%
\personnageFr{Mnésiloque}%
\tetramTrochFr{Puisse-t-il n'en être rien, j'en conjure le ciel.\\\passeLigne}%
\end{groupe}%
%% v. 715-716
\begin{groupe}%
\personnageFr{Chœur}%
\diDochFr{Et qui donc, qui des dieux immortels s'associera\\
à des actes injurieux ?\\\passeLigne}%
\end{groupe}%
%% v. 717
\begin{groupe}%
\personnageFr{Mnésiloque}%
\trimIambCatFr{Inutiles bavardages ! Je ne la lâcherai pas.\\\passeLigne}%
\end{groupe}%
%% v. 718-725
\begin{groupe}%
\personnageFr{Chœur}%
\chantChoeurFr{Ah! par les deux déesses, peut-être \\
n'auras-tu pas lieu de rire de tes outrages \\
et des discours insensés que tu profères.}
\end{groupe}%
\endnumbering
\end{Rightside}
\end{pages}
\Pages
\clearpage\ifodd\value{page}\hbox{}\newpage\fi
\enTeteG{ΕΠΕΙΣΟΔΙΟΝ Γ' --- Codex Ravennas 429, f. 171r (Classense 173r) --- vv. 721–764}
\begin{adjustwidth}{-1.5cm}{-3.0cm}
\begin{adjustbox}{max totalheight=21.0cm, max width=16.5cm, keepaspectratio}
\begin{minipage}{16.5cm}
{\fontsize{7.2pt}{8.5pt}\selectfont
\noindent\textbf{Folio 18 --- Codex Ravennas 429, \texttt{f. 171r} (\texttt{Classense 173r}) : ΕΠΕΙΣΟΔΙΟΝ Γ', vv. 721–764}\hfill \textit{Témoins :} $\text{R}^1, \Rpc, \Rdeux$, \PsiDeux\par\vspace{2pt}
\hrule height 0.4pt\vspace{2.5pt}
}
{\fontsize{8.1pt}{9.5pt}\selectfont
\ligneRav{721}{χορ/}{\hspace{1.1em}ξὺν ἀθέοις ἔργοις}{\textbf{721} ἀθέοις ἔργοις γὰρ R : ἀθ. γὰρ ἔργ. Burges : ἐπ' ἀθ. ἔρ. γὰρ Enger : \lecon{ξὺν ἀθ. ἔργ. γὰρ} Bibi ap. Taillardat cf. v.716 \espApp γὰρ R : \lecon{καὶ γὰρ} Hermann}
\ligneRav{}{}{καὶ γὰρ ἀνταμειψόμεσθά σ᾽ ὥσπερ εἰκὸς ἀντὶ τῶνδε·}{}
\ligneRav{724}{}{\hspace{1.1em}\difference{τάχα δὲ} μεταβαλοῦσ᾽ ἐπὶ κακὸν ἑτερότρο-}{\textbf{724} τάχα δέ σε R : \lecon{τάχα δὲ} Fritzsche}
\ligneRav{725}{}{\hspace{1.1em} πον ἐπέχει \difference{τύχη}·}{\textbf{725} τις τύχη R : \lecon{τύχη} Bergk}
\ligneRav{726}{χορ/}{ἀλλὰ τάσδε μὲν λαβεῖν χρῆν σ᾽ ἐκφέρειν τε τῶν ξύλων,}{\textbf{726} \lecon{\expo{χ}ρῆν σ᾽} \Rdeux : χρὴ μ᾽ \Mon}
\ligneRav{}{}{καὶ καταίθειν τὸν πανοῦργον πυρπολεῖν θ᾽ ὅσον τάχος·}{}
\ligneRav{}{\auDessus{γυ}{ν}}{ἴωμεν ἐπὶ τὰς κληματίδας ὦ Μανία·}{}
\ligneRav{729}{}{κἀγώ σ᾽ ἀποδείξω θυμάλωπα τήμερον·}{\textbf{729} \lecon{κἀγώ} R : ἦ ᾽γώ Sommerstein}
\ligneRav{730}{μ\auDessus{ν}{Η}}{ὕφαπτε καὶ κάταιθε· \difference{σὺ δὲ τὸ} Κρητικὸν}{\textbf{730} σὺ τόδε τὸ R : \lecon{σὺ δὲ τὸ} Suid. κ 2408}
\ligneRav{}{}{ἀπόδυθι ταχέως· τοῦ θανάτου δ᾽ ὦ παιδίον}{}
\ligneRav{}{}{μόνην γυναικῶν αἰτιῶ τὴν μητέρα·}{}
\ligneRav{733}{}{τουτὶ τί ἔστιν; ἀσκὸς \difference{ἐγένεθ᾽} ἡ κόρη}{\textbf{733} ἐγένηθ᾽ R : \lecon{ἐγένεθ᾽} Gryaneaus}
\ligneRav{}{}{οἴνου πλέως καὶ ταῦτα Περσικὰς ἔχων·}{}
\ligneRav{735}{}{ὦ θερμόταται γυναῖκες, ὦ ποτίσταται}{\textbf{735} \lecon{θερμόταται} R, Suid. θ 251}
\ligneRav{}{}{κἀκ παντὸς ὑμεῖς μηχανώμεναι πιεῖν,}{}
\ligneRav{}{}{ὦ μέγα καπήλοις ἀγαθὸν ἡμῖν δ᾽ αὖ κακόν,}{}
\ligneRav{}{}{κακὸν δὲ καὶ τοῖς σκευαρίοις καὶ τῇ κρόκῃ·}{}
\ligneRav{}{\auDessus{γυ}{ν}}{παράβαλλε πολλὰς κληματίδας ὦ Μανία·}{}
\ligneRav{740}{μ\auDessus{ν}{Η}}{παράβαλλε δῆτα· σὺ δ᾽ \difference{ἀπόκριναί} μοι τοδί,}{\textbf{740} ἀπόκρινέ R : \lecon{ἀπόκριναί} Μυ2}
\ligneRav{741}{—}{τουτὶ τεκεῖν φῄς; : \difference{καὶ δέκα μῆνας} αὔτ᾽ ἐγὼ}{\textbf{741} δέκα μῆνας R : \lecon{καὶ δέκα μῆνας} Biset : δέκα μῆνας γε Scaliger}
\ligneRav{}{—}{ἤνεγκον. : ἤνεγκας σύ; : νὴ τὴν Ἄρτεμιν·}{}
\ligneRav{743}{μ\auDessus{ν}{Η}}{τρικότυλον; ἢ πῶς. : εἰπέ μοι· τί μ᾽ \difference{ἠργάσω};}{\textbf{743} \_τρικο\douteux{τ}\douteux{υ}\douteux{λ}\douteux{ο}ν·η\douteux{π}\douteux{ω}\douteux{c} \  ειπ\douteux{ε}\douteux{μ}οι\douteux{τ}[ \PsiDeux : τρικότυλον ἢ πῶς; εἰπέ μοι --- τί μ᾽ εἰργάσω; R : \lecon{ἠργάσω} Hall-Geldart}
\ligneRav{}{}{ἀπέδυσας ὦναίσχυντέ μου τὸ παιδίον}{}
\ligneRav{745}{—}{\difference{\difference{τυννοῦτο}ν} ὄν. : τυννοῦτο; μικρὸν νὴ Δία·}{\textbf{745} \lecon{τυννοῦτον} et \lecon{τυννοῦτο} \PsiDeux, Brunck : τυνοῦτον et τυνοῦτο R \espApp Γϋ. ante μικρὸν R : \lecon{nulla nota mutat. pers. } \PsiDeux}
\ligneRav{746}{}{πόσ᾽ ἔτη δὲ \difference{γέγονε}; τρεῖς χοᾶς ἢ τέτταρας;}{\textbf{746} \lecon{γέγονε} \PsiDeux, Brunck : γέγονεν R}
\ligneRav{747}{\auDessus{γυ}{ν}}{σχεδὸν τοσοῦτον \difference{χὤσον} ἐκ Διονυσίων·}{\textbf{747} καὶ ὅσον \PsiDeux, R : \lecon{χὤσον} Bentley\espApp καὶ ὅσον \PsiDeux, R : ὅσον Suid. δ 1168 : \lecon{χὤσον} Bentley}
\ligneRav{}{—}{ἀλλ᾽ ἀπόδος αὐτό. : μὰ τὸν Ἀπόλλω τουτονί·}{}
\ligneRav{749}{\auDessus{γυ}{ν}}{ἐμπρήσομεν τοίνυν σε. : πάνυ γ᾽ ἐμπίμπρατε·}{\textbf{749} ἐμπιμπρᾶται R : \lecon{ἐμπιμπράτε} (sic) \Mon}
\ligneRav{750}{}{αὕτη δ᾽ ἀποσφαγήσεται μάλ᾽ αὐτίκα·}{}
\ligneRav{}{\auDessus{γυ}{ν}}{μὴ δῆθ᾽, ἱκετεύω σ᾽· ἀλλ᾽ ἔμ᾽ ὅ τι χρῄζεις ποίει}{}
\ligneRav{}{—}{ὑπέρ γε τούτου. : φιλότεκνός τις εἶ φύσει·}{}
\ligneRav{}{}{ἀλλ᾽ οὐδὲν ἧττον ἥδ᾽ ἀποσφαγήσεται}{}
\ligneRav{754}{\auDessus{γυ}{ν}}{δὸς τὸ σφαγεῖον (σφαγ\inRas{ί}{εῖ}ον \Rpc), ὦ γύναι,}{\textbf{754} \lecon{δός τὸ} \PsiDeux, Lobeck ad Soph. \italique{Ai.} v.1066 : δός μοι τὸ R \espApp σφαγίον \Rac : \lecon{σφαγεῖον} \Rpc, \Mon}
\ligneRav{755}{}{ἵν᾽ οὖν τό γ᾽ αἷμα τοῦ τέκνου τοὐμοῦ λάβω·}{}
\ligneRav{}{μ\auDessus{ν}{Η}}{ὕπεχ᾽ αὐτό, χαριοῦμαι γὰρ ἕν γε τοῦτό σοι·}{}
\ligneRav{}{\auDessus{γυ}{ν}}{κακῶς ἀπόλοι᾽, ὡς φθονερὸς εἶ καὶ δυσμενής·}{}
\ligneRav{}{μ\auDessus{ν}{Η}}{τουτὶ τὸ δέρμα τῆς ἱερείας γίγνεται·}{}
\ligneRav{758}{\auDessus{γυ}{ν}}{τί τῆς \difference{ἱερείας} γίγνεται; : τουτί. λαβέ·}{\textbf{758-9} \lecon{ἱερείας} \PsiDeux, R : ἱερέας Austin \& Olson}
\ligneRav{760}{ἄλλη γυνή}{ταλαντάτη Μίκα τίς ἐξεκόρησέ σε;}{\textbf{760} ἐξεκόρησεν R : \lecon{ἐξεκόρησε} \Mon}
\ligneRav{761}{}{τίς τὴν ἀγαπητὴν παῖδά \difference{σοὐξηράσατο};}{\textbf{761} [\douteux{ }\douteux{ }\douteux{ }\douteux{ }\douteux{ }\douteux{ }\douteux{ }\douteux{ }]απη\douteux{τ}η\douteux{ν}\douteux{π}\douteux{ }ι\douteux{ }α\douteux{ }\douteux{ }υξ[ \PsiDeux \espApp σοὐξῃρήσατο R : \lecon{σοὐξηράσατο} Fritzsche}
\ligneRav{}{\auDessus{γυ}{ν}}{ὁ πανοῦργος οὗτος. ἀλλ᾽ ἐπειδήπερ πάρει,}{}
\ligneRav{}{}{φύλαξον αὐτόν, ἵνα λαβοῦσα Κλεισθένη}{}
\ligneRav{}{}{τοῖσιν πρυτάνεσιν ἃ πεποίηχ᾽ οὗτος φράσω·}{}
}
\vspace{3pt}\noindent\rule{4.5cm}{0.4pt}\par\vspace{1.5pt}
{\fontsize{7.0pt}{8.3pt}\selectfont
\noindent\textbf{Loci paralleli (vv. 721–764) ---} \textbf{729} (θυμάλωπα) Hesychius θ 862 θυμάλωψ ; cf. Photius θ 256 Θυμάλωψ, Hesychius θ 861 \espTest \textbf{730-1} (σὺ δὲ -- ταχέως) Suidas κ 2408 \espTest \textbf{731} (Κρητικὸν) Hesychius κ 4087 Κρητικόν· ἱματίδιον λεπτὸν καὶ βραχύ \espTest \textbf{734} (Περσικὰς) \Eustathe 1522.12 \espTest \textbf{735-7} Suidas θ 251 (ὦ et αὖ omissis, καλόν et δὲ pro ἀγαθόν et δ᾽) \espTest \textbf{735} \Eustathe 1441.23 ; \EtymSym α 241 (ed. Lasserre-Livadaras) \espTest (ποτίσταται) \EtymMagnum 31.13 \espTest \textbf{742} (ἤνεγκον) \Antiatt 98.11 ; \Phryn 73.1 \espTest \textbf{745} (τυννοῦτο) Photius τ 552 τυννοῦτο $\sim$ Hesychius τ 1645 \espTest \textbf{747} Suidas δ 1168 (καὶ omisso) \espTest \textbf{749} (ἐμπίμπρατε) Hesychius ε 2458 ἐμπίπραται $\sim$ Suidas ε 1018 $\sim$ Zonaras 709 \espTest \textbf{754} Photius σ 850 σφαγεῖον· [...] περισπᾶται ἡ δευτέρα συλλαβή\par
}
\vspace{2.5pt}\noindent\rule{4.5cm}{0.4pt}\par\vspace{1.5pt}
{\fontsize{6.9pt}{8.2pt}\selectfont
\noindent\textbf{Suite (en commentaire) de l'apparat critique \& notes paléographiques (\texttt{f. 171r}) ---} \textbf{Contexte :} Suite et dénouement de la parodie du \italique{Télèphe} : les femmes apportent du bois pour brûler Mnésiloque ; en démaillotant l'otage, Mnésiloque découvre que le « bébé » est une outre pleine de vin chaussée de petites bottines perses (\italique{persikai}), qu'il « égorge » et boit goulûment au-dessus du bassin sacrificiel (\italique{sphageion}, vv. 721–764). --- \textbf{Témoin papyrologique :} \textbf{\PsiDeux{} (\italique{P.Oxy.} LVI 3839, fr. 1)} couvre les \textbf{vv. 742–766} et apporte quatre confirmations décisives face à R (vv. 743, 745, 746, 754). --- \textbf{Confrontation directe entre R et le papyrus \PsiDeux{} (vv. 743–754) :} --- \textbf{v. 721 :} ἀθέοις ἔργοις γὰρ R (ἀθ. γὰρ ἔργ. Burges ; ἐπ' ἀθ. ἔρ. γὰρ Enger) : \lecon{ξὺν ἀθέοις ἔργοις γὰρ} \lecon{Bibi} ap. Taillardat (cf. v. 716) \espApp γὰρ R : \lecon{καὶ γὰρ} Hermann.\par
}
\end{minipage}
\end{adjustbox}
\end{adjustwidth}
\newpage
\enTeteD{Codex Ravennas 429 (Biblioteca Classense) --- Fac-similé du f. 171r (vv. 721–764)}
\begin{adjustwidth}{-1.5cm}{-3.0cm}
\begin{center}
\includegraphics[width=16.5cm,height=20.5cm,keepaspectratio]{/workspaces/Thesmophories/dossier_ravennas_thesmophories/folios_individuels_ff162v-177r/18_f171r_classense173r_vv0721-0764.jpg}\par
\vspace{2pt}
{\scriptsize\textbf{Codex Ravennas 429, f. 171r (Classense 173r)} --- ΕΠΕΙΣΟΔΙΟΝ Γ', \textbf{vv. 721–764} ($\text{R}^1, \Rpc, \Rdeux$, \PsiDeux).}
\end{center}
\end{adjustwidth}
\newpage
\clearpage\ifodd\value{page}\hbox{}\newpage\fi
\enTeteG{ΘΕΣΜΟΦΟΡΙΑΖΟΥΣΑΙ --- Texte grec établi (f. 171r, vv. 721–764)}
\enTeteD{LES THESMOPHORIES --- Traduction française (f. 171r, vv. 721–764)}
\begin{pages}
\begin{Leftside}
\beginnumbering
%% v. 718-725
\begin{groupe}%
\typeVers{dochm.}\chantChoeur{ξὺν ἀθέοις ἔργοις\apparat{721}{ἀθέοις ἔργοις γὰρ R : ἀθ. γὰρ ἔργ. Burges : ἐπ' ἀθ. ἔρ. γὰρ Enger : \lecon{ξὺν ἀθ. ἔργ. γὰρ} Bibi ap. Taillardat cf. v.716 \espApp γὰρ R : \lecon{καὶ γὰρ} Hermann \ligneApp{724} τάχα δέ σε R : \lecon{τάχα δὲ} Fritzsche \ligneApp{725} τις τύχη R : \lecon{τύχη} Bergk}% grat. metr. \ligneApp{726} \lecon{\expo{χ}ρῆν σ᾽} \Rdeux : χρὴ μ᾽ \Mon}%
\commentaire{721-5}{ἀθέοις ἔργοις γὰρ ἀνταμει- / ψόμεσθά σ᾽ ὥσπερ εἰκὸς ἀντὶ τῶνδε / τάχα δὲ \{σε\} μεταβαλοῦσ᾽ ἐπὶ κακὸν ἑτερότροπον / ἐπέχει \{τις\} τύχη.}{Le grand Jean Taillardat avait donné un exposé sur son interprétation des vers 721-725 du manuscrit de Ravenne. Il nous offre une explication limpide de ce passage fort débattu ; je transcris pieusement ci-dessous le résumé de son intervention, tout en reprenant sa traduction\marqueur{1} : \\
« Car toi, avec tes impiétés, nous te châtierons de ta conduite, comme de juste. Et puis la Fortune, changeant brusquement d'humeur, aura vite fait de poursuivre un méchant homme qui a pris le chemin du crime ».\\
Cette traduction appelle les justifications suivantes : \\
— Ἀθέοις ἔργοις : selon Jean Taillardat, il est impossible d'interpréter ce datif comme le complément de moyen ἀνταμειψόμεσθα et de traduire « \latin{\italique{Impiis factis} enim remunerabimus te sicut par est pro eis} » (Andrea Divo) ou « \latin{\italique{Nefariis} enim \italique{factis} vicissim, ut par est, haec tua ulciscemur} » (Brunck) ou « Nous y répondrons, comme de raison, \italique{par des actes sacrilèges} » (Willems)\marqueur{2}. En effet, on ne pouvait pas considérer comme sacrilège « \italique{la coutume} (qui) était d'apporter du feu près des suppliants réfugiés sur l'autel » afin de les contraindre à le lâcher (Σ Eur. \italique{Androm}. 256 : ἔθος ἦν πῦρ προσφέρειν τοῖς εἰς βωμὸν καταφεύγουσιν ; cf. Eur. \italique{Herc. fur}. 243, 285). Ἀθέοις ἔργοις est donc : \ \ \ 1) \ soit un datif-instrumental d'accompagnement lié à σ(ε), cf. Kühner-Gerth, I, p. 435; ἀθέοις ἔργοις a, dans ce cas, la même fonction syntaxique que ξὺν ἀδίκοις ἔργοις\marqueur{3} (v. 716) lié à σοι (v. 715) ; \ \ \ 2) \ soit un datif-instrumental de cause lié à ἀνταμειψόμεσθα. Jean Taillardat penche pour la première solution.\\
— Ἀντὶ τῶνδε (v. 723) $=$ ἀντὶ τῶνδε τῶν ἀθέων ἔργων « nous te payerons en retour de tes impiétés».\\
— Μεταβαλοῦσ(α), v. 724, est employé absolument, ce qui n'est pas rare, cf. L.S.J. \italique{s. v.} μεταβάλλω, A. III. 2.\\
On reconnaîtra d'autre part la locution ἐπέχειν ἐπί τινα (Hérodote 9.59) à peu près synonyme de διώκειν τινά (Id. \italique{ibid}.), car τάχα ... ἐπὶ κακὸν ἑτερότροπον ἐπέχει (présent de prophétie) τύχη équivaut à τάχα ... ἐπὶ κακὸν ἄνδρ᾽ ὃς τὴν ἑτέραν ὁδὸν ἐτράπετ' ἐπέχει τύχη.\\
Pour κακός $=$ κακός τις ἀνήρ, cf. Théognis 189, 190, 577, 972, 1112 (de même κακὴ $=$ κακή τις γυνή, id. 185). Le composé ἑτερότροπος « qui a pris la mauvaise route » (avec ἕτερος euphémique pour κακός, cf. L.S.J. III. 2) est une allusion au mythe bien connu des deux chemins, celui du bien et celui du mal (Hésiode \italique{Op}. 287-292 ; Prodicos in Xénophon \italique{Mém}. 2. 1. 21-27). \\
Quant à σε de τάχα δέ σε μεταβαλοῦσα, κτλ. (cod. R), Jean Taillardat le considère comme une glose qui expliquait κακόν ἑτερότροπον et qui s'est glissée dans le texte même. \\
Quoi qu'on en ait dit, le sens de ce passage est très clair : le chœur fait entendre au Parent d'Euripide qu'à défaut d'une punition donnée par les femmes, il recevra un châtiment infligé par le Ciel.}
\debutNote{721-5}
\note{1}{À retrouver dans la \italique{Revue des Études Grecques}, tome 101, fascicule 482-484 (juillet-décembre 1988). Actes de l'association, p. XVIII-XIX.}
\note{2}{Ou « \anglais{Since your crimes are godless, it is fitting that our punishment of you be \italique{godless} as well} » (Austin \& Olson.)}%
\note{3}{Fort de cette remarque %on-ne-peut-plus pertinente
, je [Bibi] n'ai pas hésité à introduire ce ξύν dans le texte  au vers 721, qui avait un problème de mètre : [ξὺν] ἀθέοις ἔργοις. C'est aussi \italique{gratia metri} qu'Enger proposait en son temps l'ajout d' [ἐπ'] ἀθέοις ἔργοις, conjecture acceptée par Lætitia (\italique{The Songs of Aristophanes}, p. 427). N.D.É.}}%
\end{groupe}%
%% v. 718-725
\begin{groupe}%
\typeVers{tetram. troch.}\tetramTroch{\small{καὶ γὰρ ἀνταμειψόμεσθά σ᾽ ὥσπερ εἰκὸς ἀντὶ τῶνδε.}}%
\end{groupe}%
%% v. 718-725
\begin{groupe}%
\typeVers{2 dochm.}\chantChoeur{\ligne{724}τάχα δὲ μεταβαλοῦσ᾽ ἐπὶ κακὸν ἑτερότρο-%
\commentaire{724-5}{τάχα δὲ \{σε\} μεταβαλοῦσ᾽ ἐπὶ κακὸν ἑτερότροπον ἐπέχει \{τις\} τύχη.}{Il suffit que le chœur dise `τάχα δὲ' ... et le rythme même du vers s'accélère : les trois dochmiaques des vers 724-725 contiennent force résolutions.}}%
\end{groupe}%
%% v. 718-725
\begin{groupe}%
\typeVers{dochm.}\chantChoeur{\alinea πον ἐπέχει τύχη.\\\passeLigne}%
\end{groupe}%
%% v. 726-727
\begin{groupe}%
\personnageGrec{Κορυφαῖος}%
\typeVers{$4$ troch\cat}\tetramTroch{\ligne{726}ἀλλὰ τάσδε μὲν λαβεῖν χρῆν σ᾽ ἐκφέρειν τε τῶν ξύλων,%
\commentaire{726-7}{ἀλλὰ τάσδε μὲν λαβεῖν χρῆν σ᾽ ἐκφέρειν τε τῶν ξύλων, / καὶ καταίθειν τὸν πανοῦργον πυρπολεῖν θ᾽ ὅσον τάχος.}{C'était, en effet, l'habitude de faire périr par le feu ceux qui avaient cherché refuge sur un autel. Cf. Euripide, \italique{Androm}., v. 256 ; \italique{Herc. Fur}., v. 243 et 285 et les interprètes. (Note de Willems.)\\
Cette \italique{exhortation} -- κατακελευσμός selon la dénomination consacrée de Zieliński\marqueur{1} -- de la Coryphée pour appeler à la reprise de l'action est selon l'usage un distique de tétramètres trochaïques. (Et, comme au vers 531-532, débute par `ἀλλά'.)}%
\debutNote{726-7}%
\note{1}{Voir \italique{The Verse of Greek Comedy} (1912), p. 316, de John W. White. Le lecteur pourra utilement consulter ce livre de référence, entre autre s'il s'intéresse à la structure des comédies d'Aristophane.}%
\commentaire{726}{Τάσδε}{renvoie à des compagnes de Mica et non pas à des torches, selon l'hypothèse invraisemblable de Fritzsche -- qui a malheureusement fait florès. Logique pratique élémentaire%\marqueur{1} 
: avant d'aller recueillir du bois, il faut d'abord se \italique{débarrasser} de ce qui encombre -- en l'occurrence : les torches -- et non s'en embarrasser. Ce n'est qu'une fois le bois entassé que Mica pourra effectivement se saisir d'une torche et bouter le feu au bûcher.}%
%\debutNote{726}%
%\note{1}{Il m'est utile d'avoir été professeur pour ce commentaire des \italique{Thesmophories} : on apprend à tout préciser, même des évidences telles qu'on aurait d'abord cru insulter l'intelligence de l'auditoire en s'abaissant à les souligner. Un éditeur d'Aristophane, tel le jeune élève dont l'esprit est en cours de formation, a parfois besoin d'être pris par la main.}%
%\commentaire{726}{ἐκφέρειν τε τῶν ξύλων,}{\eff{parler de la porte du fond de scène d'où serait tiré ce bois ?}}
\\
καὶ καταίθειν τὸν πανοῦργον πυρπολεῖν θ᾽ ὅσον τάχος.\\\passeLigne}%
\end{groupe}%
%% v. 728-729
\begin{groupe}%
\personnageGrec{Μίκα}%
\typeVers{trim. iamb.}\trimIamb{ἴωμεν ἐπὶ τὰς κληματίδας ὦ Μανία.\commentaire{728}{ἴωμεν ἐπὶ τὰς κληματίδας.}{Blaydes, jugeant l'article de trop, propose : ἴωμεν οὖν ἐπὶ κληματίδας. Nous répondrons par le texte suivant de Plutarque, où il s'agit, comme ici, d'un sanctuaire où des femmes se trouvent réunies pour la célébration de leurs mystères :%
\citComm{$[$ ἡ Δαμοκρίτα $]$ ξύλα ταῖς θύραις πολλὰ προσνήσασα ταῦτα δ᾽ ἦν \textbf{εἰς τὴν τῆς ἑορτῆς θυσίαν ὑπ᾽ ἐκείνων παρεσκευασμένα}, πῦρ ἐνῆκε.\\
Damocrita ayant entassé contre les portes une quantité de bois, \italique{que les femmes avaient préparé pour le sacrifice de la fête}, elle y mit le feu.}%
(Plutarque, \italique{Narrat. amour}., 5, p. 775e ; Note de Willems.)}\\
κἀγώ σ᾽ ἀποδείξω θυμάλωπα τήμερον.\testimonia{729}{(θυμάλωπα) Hesychius θ 862 θυμάλωψ ; cf. Photius θ 256 Θυμάλωψ, Hesychius θ 861}\apparat{729}{\lecon{κἀγώ} R : ἦ ᾽γώ Sommerstein}\\\passeLigne}%
\end{groupe}%
%% v. 730-732
\begin{groupe}%
%
\personnageGrec{Μνησίλοχος}%
\trimIamb{ὕφαπτε καὶ κάταιθε· σὺ δὲ τὸ Κρητικὸν\testimonia{730-1}{(σὺ δὲ -- ταχέως) Suidas κ 2408}\testimonia{731}{(Κρητικὸν) Hesychius κ 4087 Κρητικόν· ἱματίδιον λεπτὸν καὶ βραχύ}\apparat{730}{σὺ τόδε τὸ R : \lecon{σὺ δὲ τὸ} Suid. κ 2408}%
\commentaire{730}{τὸ Κρητικὸν.}{Hésychius (κ 4087) : %
\citComm{Κρητικόν· ἱματίδιον λεπτὸν καὶ βραχύ.\\
\francais{Crétique : petit habit fin et court.}}%
}\\%
ἀπόδυθι ταχέως· τοῦ θανάτου δ᾽ ὦ παιδίον\\
μόνην γυναικῶν αἰτιῶ τὴν μητέρα.}%
\end{groupe}%
%% v. 733-735
\begin{groupe}%
\trimIamb{τουτὶ τί ἔστιν; ἀσκὸς ἐγένεθ᾽ ἡ κόρη\apparat{733}{ἐγένηθ᾽ R : \lecon{ἐγένεθ᾽} Gryaneaus}%
\commentaire{733}{ἀσκὸς ἐγένεθ᾽ ἡ κόρη.}{Il était probablement interdit d'apporter du vin lors des Thesmophories, fêtes consacrées à Déméter et sa fille. (Cf. par exemple l'adage latin \latin{Cereri sacrificant}\marqueur{1}, \italique{sacrifier à Cérès}, c'est-à-dire faire abstinence de vin.) Cette interdiction théorique devait souffrir quelques écarts pratiques, si l'on en croit le compte-rendu de Mnésiloque (vers 626-633).}%
\debutNote{733}%
\note{1}{Recensé par Érasme, \italique{Adages} n° 3071. Voir aussi les adage n° 1895, 1896.}\\
οἴνου πλέως καὶ ταῦτα Περσικὰς ἔχων.\testimonia{734}{(Περσικὰς) \Eustathe 1522.12}%
\commentaire{734}{Περσικὰς ἔχων.}{On comprend d'après d'autres passages d'Aristophane que les \italique{persiques} désignent des chaussures légères et recouvrant le pied\marqueur{1}, peut-être des \italique{chaussons}\marqueur{2} ; en tout cas, ces \italique{persiques} sont portées par les femmes\marqueur{3}, ou comme ici les enfants.\\
Les «~persiques~» dont parle Mnésiloque sont visibles sur le fameux cratère de Würzbourg\marqueur{4}. Il faut savoir que les outres à l'époque étaient faites de la peau d'un animal évidé\marqueur{5}, bouc ou chèvre la plupart du temps. Ces «~persiques~» sont en réalité les pattes arrières de l'animal outré, fermées par un lacet.}%
\debutNote{734}%
\note{1}{La puce des socratiques expériences des \italique{Nuées}, v. 151, se retrouve chaussée de `persiques' après que Socrate lui ait plongé les pattes dans la cire – voir aussi Hésychius (π 2005) : \grec{περσικά· εὐτελῆ ὑποδήματα}.% \eff{N'y a-t-il pas là un proverbe ?}
}%
\note{2}{Les Athéniennes pouvaient les garder au lit, en prenant garde toutefois de maintenir les jambes en l'air, cf. \italique{Lysistrata} v. 229.}%
\note{3}{\italique{Assemblée des femmes} v. 319 et Hésychius (μ 929) : \grec{μεσοπερσικαί· ὑποδήματα γυναικεῖα}.}%
\note{4}{Reproduit en annexe. Voir la note aux vers 754-757.}%
\note{5}{Pour plus d'informations, aller consulter l'indispensable Daremberg \& Saglio, \italique{Dictionnaire des antiquités grecques et romaines} à l'entrée \latin{Uter}.% \eff{+ Realencyclopädie}
}%
\\
ὦ θερμόταται γυναῖκες, ὦ ποτίσταται\testimonia{735-7}{Suidas θ 251 (ὦ et αὖ omissis, καλόν et δὲ pro ἀγαθόν et δ᾽)}\testimonia{735}{\Eustathe 1441.23 ; \EtymSym α 241 (ed. Lasserre-Livadaras) \espTest (ποτίσταται) \EtymMagnum 31.13}%\apparat{735}{\lecon{θερμόταται} R, Suid. θ 251}% : θερμοπόται Bibi}%
\commentaire{735}{ὦ θερμόταται γυναῖκες.}{θερμός évoque l'\italique{impétuosité} féminine : \anglais{hot-headed, hasty}, l'\italique{audace} des femmes\marqueur{1}.
Le mot peut aussi désigner la \italique{fièvre} ; il n'y a cependant pas d'équivalent grec à notre \italique{fièvre de Bercy} : à ma connaissance, le grec n'a jamais rapproché la fièvre de l'ivrognerie.
Quand donc le scholiaste à notre vers écrit `ὦ μέθυσοι', il glose `ὦ ποτίσταται' et non `ὦ θερμόταται'.
Suidas (θ 251) s'y est trompé qui recopie la glose en la faisant précéder du lemme θερμόταται :%
\citComm{Θερμόταται: μέθυσοι. ὦ θερμόταται γυναῖκες, μέγα καπήλοις καλόν, ἡμῖν δὲ κακόν [...]}%
Le cas n'est pas isolé : Suidas avait déjà commis un précédent du même genre : au vers 178 (\italique{confer} le commentaire \italique{ad locum}), lisant κοινῇ à place de καινῇ, il n'hésita pas à utiliser la glose afférente, μεγάλῃ καὶ θαυμαστῇ, comme nouvelles acceptions de κοινῇ (en y ajoutant le véritable sens de κοινῇ : ἀνθρωπίνῃ).
Ce sont là des méthodes qui montrent que la \italique{Souda} est un lexique à manipuler avec précaution.
%Faut-il voir avec le \LSJ , Taillardat (§313, 351) et les différents traducteurs une évocation de l'\italique{impétuosité} féminine : \anglais{hot-headed, hasty}, de leur \italique{audace} ? Comme suivent une énumération des vertus ivrognesques du beau sexe, je serais porté à l'avis du scholiaste – qui glose \grec{ὦ μέθυσοι} –, repris par Suidas à l'entrée (θ 251) :
%\citComm{Θερμόταται: μέθυσοι. ὦ θερμόταται γυναῖκες, μέγα καπήλοις καλόν, ἡμῖν δὲ κακόν [...]}
%Faut-il comprendre qu'il s'agit d'un type très particulier de bouffées de chaleur : la \italique{fièvre de Bercy}~? «~Ô les enfiévrées de Bercy~!~»\\
%Pourtant, cette acception de θέρμος n'est attestée nulle part ailleurs et n'a pas non plus été relevée par les lexicographes...\\
%Alors, un scribe aurait-il par erreur attribué la glose au mauvais lemme ; ὦ μέθυσοι se rapportant à ὦ ποτίσταται et non ὦ θερμόταται ?\\
%Je vois encore une autre hypothèse : si le texte d'Aristophane contenait à l'origine ὦ θερμοπόται\marqueur{1}.
%\eff{θερμοπότης désigne littéralement un \italique{buveur de boisson chaude}\marqueur{θερμοποτίς étant l'un des termes pour la coupe où l'on délayait \italique{à chaud} le vin.}. Avec l'usage de mélanger le vin épais à de l'eau \italique{chaude}, ce mot a pu désigner plus particulièrement les francs buveurs (voir Athénée). serait à prendre au sens de "ardent à boire" (forgé sur le même modèle que le θερμό-βουλος d'Euripide, fr. 858, qu'Aristophane parodie au vers 119 des \italique{Acharniens}). 
%"vin béni pour les débitants"}%
}%
\debutNote{735}%
\note{1}{Cf. le \LSJ , Taillardat §§313, 351.%\eff{le citer ? + Érasme}
}%
%\note{1}{À prendre au sens de "ardent à boire" ? (forgé sur le même modèle que le θερμό-βουλος d'Euripide, fr. 858, qu'Aristophane parodie au vers 119 des \italique{Acharniens}.) Le θερμοπώλιον, ou \latin{thermopolium}, désignait une \italique{taverne} chez les latin ; le θερμοπότης, \italique{qui boit chaud}, devait être particulièrement porté sur certain type de boisson, comme il appert d'un passage d'Athénée (VIII, 352b) :τοὺς δὲ Ῥοδίους ὁ αὐτὸς Στρατόνικος σπαταλῶνας καὶ θερμοπότας θεωρῶν.}%
}%
\end{groupe}%
%% v. 736-738
\begin{groupe}%
\trimIamb{κἀκ παντὸς ὑμεῖς μηχανώμεναι πιεῖν,\\
ὦ μέγα καπήλοις ἀγαθὸν ἡμῖν δ᾽ αὖ κακόν,\\
κακὸν δὲ καὶ τοῖς σκευαρίοις καὶ τῇ κρόκῃ.\\
\commentaire{738}{κακὸν δὲ καὶ τοῖς σκευαρίοις καὶ τῇ κρόκῃ.}{%\eff{Erasme 1646 et l'athénien qui tend la main jusqu'à la mort.}
Mnésiloque reproche à ces soûlaudes d'Athéniennes de contribuer à la ruine des ménages par leur inconduite biberonnesque.
Montrons-le !
\grec{τό σκευάριον}, diminutif de σκεῦος, désigne littéralement un \italique{petit récipient} (\italique{menue vaisselle}) ou \italique{petit meuble}.
Employé au pluriel, il peut désigner plus particulièrement le meuble où l'on entasse son argent. \italique{Confer} le passage du \italique{Ploutos} ci-dessous (vers 806-811) qui nous renseigne sur l'usage du mobilier grec :%
\citComm{Ἡ μὲν σιπύη μεστή ᾽στι λευκῶν ἀλφίτων,\\
οἱ δ᾽ ἀμφορῆς οἴνου μέλανος ἀνθοσμίου.\\
Ἅπαντα δ᾽ ἡμῖν ἀργυρίου καὶ χρυσίου\\
\textbf{τὰ σκευάρια} πλήρη ᾽στίν, ὥστε θαυμάσαι.\\
Τὸ φρέαρ δ᾽ ἐλαίου μεστόν· αἱ δὲ λήκυθοι\\
μύρου γέμουσι, τὸ δ᾽ ὑπερῷον ἰσχάδων.\\
\francais{Notre huche au pain regorge d'une blanche farine, /
nos amphores d'un vin noir qui embaume. /
Tous nos \italique{meubles}} (τὰ σκευάρια) \francais{sont pleins /
d'argent et d'or, que c'est merveille. /
La citerne regorge d'huile, les lécythes /
sont emplis d'essences, le grenier de figues sèche}}
(C'est bien de \italique{pièces} d'or et d'argent dont il est question puisque les plats et autres assiettes d'airain ou d'argent sont évoqués aux vers suivants.)
Il est tout naturel que ces meubles (coffrets ?) servant de porte-monnaie aient pris par métonymie le sens de \italique{biens}, \italique{économies}.
C'est d'ailleurs si naturel que c'est effectivement le cas : aux vers 838-839, le Juste, contant ses misères passées, emploie τὰ σκευάρια en ce sens métonymique :
\citComm{Δίκαιος — κομιδῇ μὲν οὖν· / αὐχμὸς γὰρ ὢν \textbf{τῶν σκευαρίων} μ᾽ ἀπώλεσεν.\\
\francais{Le Juste — Précisément. Une sécheresse s'étant mise dans \italique{mon mobilier} ($=$ mes économies) fut ce qui me perdit.}}%
Cette acception de τὰ σκευάρια est encore plus manifeste dans un passage au début des \italique{Grenouilles}. Dionysos y cherche un moyen d'expédier en Enfer ses \italique{bagages} ; son esclave Xanthias regimbant, il hèle un mort en route pour l'Hadès (vers 170-172) :
\citComm{Καὶ γάρ τιν᾽ ἐκφέρουσι τουτονὶ νεκρόν.\\
Οὗτος, σὲ λέγω μέντοι, σὲ τὸν τεθνηκότα.\\
Ἄνθρωπε, βούλει \textbf{σκευάρι᾽} εἰς Ἅιδου φέρειν;\\
\francais{Car voilà justement un mort qu'on emporte.\\
Holà ! c'est à toi que je parle, toi le mort.\\
Homme, veux-tu porter à l'Hadès quelques \italique{petits colis} ?}}%
Il y a un évident détournement dans la plus pure veine aristophanesque de l'universel adage : on n'emporte (φέρειν) pas son \italique{argent} (σκευάρι᾽) dans l'autre monde (εἰς Ἅιδου) ; les σκευάρια de l'adage sont à prendre au sens figuré, \italique{argent} mais se retrouvent tout matériellement incarnés en \italique{bagages} dans la version illustrée d'Aristophane\marqueur{1}. Voilà pour τοῖς σκευαρίοις, passons à τῇ κρόκῃ.\\
\grec{ἡ κρόκη} est au propre \italique{le fil de trame}. 
Il possède, comme nous l'allons voir, le même sens imagé de \italique{richesses} que τὰ σκευάρια.
Le grec, comme le français, a prodigalement tiré maintes images et maints proverbes de l'activité filandière et tisserande. Observons l'un deux : σπαθᾶν τὴν κρόκην.
Au vers 55 des \italique{Nuées}, Strepsiade énumère les oppositions entre lui, le rustre économe, et sa femme, la raffinée dilapidatrice. Il pousse la comparaison jusque dans l'alcôve où sa femme souffre de besoins tyranniques ; il lui reproche de \italique{tirer trop sur la corde} -- le grec dit \italique{frapper sur le fil} :%
\citComm{ὦ γύναι, λίαν σπαθᾷς.}%
σπαθᾶν [τὴν κρόκην] désigne l'action de battre chaque fil de trame (κρόκη) au moyen de la σπάθη, ou \italique{battant}, pour rendre le tissu compact\marqueur{2} ; ce faisant, l'on utilise davantage de laine, d'où le fait que l'expression signifie \italique{dissiper l'argent}, \italique{gaspiller}.
Or l'on trouve les expressions σπαθᾶν τὰ πατρῷα\marqueur{3}, τὰ χρήματα\marqueur{4}, τάλαντα\marqueur{5}. 
Le sème commun à ces différents mots est celui de \italique{richesse} (matérielle ou pécuniaire) ; puisqu'ils ont été substitués à τὴν κρόκην, c'est donc la preuve -- taillardesque\marqueur{6} -- qu'ils partagent avec κρόκη ce sème.
%Pour s'assurer que notre propos soit bien entendu, 
Par exemple, en français l'on rencontre les expressions \italique{ourdir un complot}, \italique{une intrigue}, \italique{une conspiration}, \italique{une vengeance}. Pris dans leur contexte, le sème commun aux différents compléments d'ourdir est : \italique{machination secrète}. Or l'expression tisserandière propre est «~ourdir une trame~», qui consiste à entrecroiser les fils de trame avec les fils ourdis, c'est-à-dire disposés sur l'ourdissoir. Les termes de complot, intrigue, \&c. ne sont substituables à la trame que parce qu'ils ont en commun avec elle le sème de \italique{machination secrète} (sens figuré de trame).\\
Finissons avec un amusant proverbe sur ce fléau qu'est l'alcool, surtout pour le sexe filandier\marqueur{7} :%
\citComm{ἡ γραῦς τὰς ἀγαθίδας ἐκπίνουσα\\
La vieille boit\marqueur{8} ses pelotes.}%
ἡ ἀγαθίς est une \italique{pelote de fil}. On constate que ce mot a aussi pris le sens de \italique{bien}, \italique{richesse} (on parlerait d'un \italique{bas de laine} en français) ; par permutation métaphorique avec κρόκη ? peut-être aussi par paronomase : ἀγαθῶν ἀγαθίδες est une expression désignant une grande quantité de biens.
Notre tradition populaire n'est d'ailleurs pas en reste : voici ce que l'on chantait à l'époque où l'on chantait encore\marqueur{9} :%
\citComm{\francais{Parlons un peu des filandières.\\
Celles qui filaient l'étoupe\\
Ont mangé toute la soupe.\\
Celles qui filaient le lin\\
Ont lampé tout le vin.}}%
Et encore, à propos de filles qui se sont enivrées :%
\citComm{\francais{Elles ne l'ont pas fait en filant du lin,\\
Mais elles l'ont bien fait en buvant du vin.}}%
%\eff{Problème (Ach. 451 et Guêpes) : c'est aussi les vêtements/costumes de théâtre...
%Guêpes 794 : le coq digère même l'argent :
%ἀλεκτρυόνος μ᾽ ἔφασκε κοιλίαν ἔχειν. “ταχὺ γοῦν καθέψεις τἀργύριον,” ἦ δ᾽ ὃς λέγων.
%À moins que le coq porte une krokê naturellement : Hésychius :
%κρόκος" τὸ φυτόν. καὶ ἀλεκτρυόνες οἱ τὸν τράχηλον τοιοῦτον ἔχοντες. κρόκη
%Item : οὐδ᾽ἐτρύπησε κρόκην : οὐδ᾽εἰργάσατo
%Lysistrata 897-898 :
%ΚΙ. Ὀλίγον μέλει σοι τῆς κρόκης φορουμένης
%ὑπὸ τῶν ἀλεκτρυόνων;
%EMPORTER LE FIL : gaspiller... sachant que les coqs portent un krokos(ê) (Hésychius !) proverbe : tirer le fil du vêtement ?}%
}%
\debutNote{738}%
\note{1}{Même procédé dans les \italique{Nuées} : quand Strepsiade désabusé songe au \italique{Tourbillon} (Δῖνος) de Socrate, il ne voit plus en cette nouvelle divinité qu'une vile \italique{poterie} (δῖνος). Si l'on veut savoir ce que le véritable Socrate entendait par ce \italique{Tourbillon}, il suffit d'interroger un autre grand esprit socratique : Pierre Boutang : «~La lignée [de race et d'esprit], c'est merveilleux !.. mais un peu trop linéaire pour mon goût. Donc il me faut quelque chose de plus circulaire : des lignes qui se répètent, qui reviennent sur elles-mêmes, qui se retournent ; des tourbillons ! [ ... ] Un grand tourbillon créateur !~». Quoi qu'on en ait dit, le Socrate croqué par Aristophane est loin d'être si caricatural et malgré ses défauts qu'Aristophane grossit à la loupe, il est encore à mille coudées au-dessus Cléon, par exemple.% de notre pâle réplique moderne de Cléon, le cocaïno-mégalomane à la `pensée complexe' -- pour prendre un exemple qui nous \italique{touche}, temporairement -- heureusement !
}%
\note{2}{Avec le commentaire du scholiaste repris dans la \italique{Souda} (ε 3176) : καὶ γὰρ ἐπὶ τῶν ὑφαινομένων ἱματίων λέγομεν σπαθᾶν τὸ ἄγαν κρούειν τὴν κρόκην. Taillardat (§440) développe l'exégèse de l'image. Voir aussi le lexique de Photius (ε 2014) : Ἐσπαθημένος· πεπυκνωμένος· ἀπὸ τῶν κατακρουουσῶν σπαθῶν τὴν κρόκην.}%
\note{3}{Diphil. 43, 27.}%
\note{4}{Plut., \italique{Per.} 14. Il s'agit ici des \italique{deniers publics} : τῶν δὲ περὶ τὸν Θουκυδίδην ῥητόρων καταβοώντων τοῦ Περικλέους ὡς σπαθῶντος τὰ χρήματα καὶ τὰς προσόδους ἀπολλύντος.}%
\note{5}{Lucien, \italique{Cat.} 20.}%
\note{6}{Lire et relire l'excellentissime article du regretté Jean, « Images et matrices métaphoriques » (1977), in \italique{Bulletin de l'Association Guillaume Budé} 4, pp. 344-354.}%
\note{7}{\italique{Corpus paroemiographorum Graecorum}, Appendices, III, 3.}%
\note{8}{Sur l'image de «~manger~» son bien, voir Taillardat §439.}%
\note{9}{Doit-on dire de cette époque qu'«~elle est révolue, qu'elle a fait son temps ; qu'on n'la reverra plus... et qu'c'est bien attristant~» ? Toujours est-il que j'emprunte ces chansons à une tradition qui a depuis longtemps défunté, mais qui fut néanmoins -- et heureusement ! -- bien embaumée par notre grand Patrice Coirault, aux cotes 110-05 et 110-06 de son \italique{Répertoire des chansons françaises de tradition orale}.}%
}%
\end{groupe}%
%% v. 739
\begin{groupe}%
%
\personnageGrec{Μίκα}%
\trimIamb{\ligne{738}παράβαλλε πολλὰς κληματίδας ὦ Μανία.\\}%
\end{groupe}%
%% v. 740-741a
\begin{groupe}%
\personnageGrec{Μνησίλοχος}%
\trimIamb{παράβαλλε δῆτα· σὺ δ᾽ ἀπόκριναί μοι τοδί,\apparat{740}{ἀπόκρινέ R : \lecon{ἀπόκριναί} Μυ2}\\
τουτὶ τεκεῖν φῄς;\\\passeLigne}%
\end{groupe}%
%% v. 741b-742a
\begin{groupe}%
\personnageGrec{Μίκα}%
\trimIamb{\ligne{741}\decalage{τουτὶ τεκεῖν φῄς;}καὶ δέκα μῆνας αὔτ᾽ ἐγὼ\apparat{741}{δέκα μῆνας R : \lecon{καὶ δέκα μῆνας} Biset : δέκα μῆνας γε Scaliger}%
\commentaire{741}{δέκα μῆνας.}{Les mois lunaires grecs étant de 29 ou 30 jours, la femme portait bien environ dix mois l'enfant. Aulu-Gelle dans ses \italique{Nuits attiques} traite pendant tout un chapitre (III, 16) de la question du temps de gestation. N'en déplaise au scholiaste, la durée de 10 mois est bien attestée, notamment par Ménandre ; Aulu-Gelle cite aussi Varron expliquant que les noms des Parques \italique{Nona} et \italique{Decima} viennent des nombres de mois fixés pour l'accouchement (\latin{item Nona et Decima a partus tempestivi tempore}).}\\
ἤνεγκον.\testimonia{742}{(ἤνεγκον) \Antiatt 98.11 ; \Phryn 73.1}\\\passeLigne}%
\end{groupe}%
%% v. 742b
\begin{groupe}%
\personnageGrec{Μνησίλοχος}%
\trimIamb{\ligne{742}\decalage{ἤνεγκον.}ἤνεγκας σύ;%
\commentaire{742-66}{}{L'on possède pour ce passage un autre témoin que le \italique{Ravennas} : le premier fragment du \nieme{3839} papyrus d'Oxyrhynque\marqueur{1} (\PsiDeux \  dans l'apparat critique). Ce morceau de papyrus est bien antérieur au manuscrit de Ravenne : il est daté d'environ 200 ap. J.-C. Il présente objectivement de meilleures leçons en bien des points. Aussi, aux endroits où ces deux témoins différaient et que le bon choix à faire n'était pas évident, j'ai privilégié la version du papyrus. Nous reproduisons en annexe ce papyrus ainsi que la lecture qu'en fit Walter E. H. Coockle pour la collection \italique{The Oxyrhynchus Papyri}.%
}%
\debutNote{742-66}%
\note{1}{\italique{The Oxyrhynchus Papyri}, vol. LVI (1989), n°3839, fr. 1.}%
\\\passeLigne}%
\end{groupe}%
%% v. 742c
\begin{groupe}%
\personnageGrec{Μίκα}%
\trimIamb{\ligne{742}\decalage{ἤνεγκον.ἤνεγκας σύ;}νὴ τὴν Ἄρτεμιν.\\\passeLigne}%
\end{groupe}%
%% v. 743a
\begin{groupe}%
\personnageGrec{Μνησίλοχος}%
\trimIamb{τρικότυλον; ἢ πῶς.\apparat{743}{\_τρικο\douteux{τ}\douteux{υ}\douteux{λ}\douteux{ο}ν·η\douteux{π}\douteux{ω}\douteux{c} \  ειπ\douteux{ε}\douteux{μ}οι\douteux{τ}[ \PsiDeux : τρικότυλον ἢ πῶς; εἰπέ μοι --- τί μ᾽ εἰργάσω; R : \lecon{ἠργάσω} Hall-Geldart}
\commentaire{743}{τρικότυλον ἢ πῶς;}{C'est-à-dire du vin supérieur, à une obole les trois cotyles. Hésychius (τ 1388) : %
\citComm{\grec{τρικότυλος οἶνος· οὗ τρεῖς κοτύλαι ὀβολοῦ πωλοῦνται},}%
ce qui est cher, en effet, car cela met le métrète (144 cotyles, environ 40 litres) à $7fr68$\marqueur{1}, le double à peu près du prix ordinaire. Démosthène (1045.5) désigne de la même manière un vin plus cher encore, \grec{ὀ οἶνος δωδεκάδραχμος}, le vin à douze drachmes, c'est-à-dire à $11fr52$ le métrète\marqueur{2}. À rapprocher de τετρώβολον, désignant dans la \italique{Paix}, 254, le miel attique. [...] — Ἤ πῶς; est mis pour ἤ πῶς ὤνιος; Cf. \italique{Acharn}., 758, \italique{Caval}., 480. (Note de Willems.)\\
Mnésiloque s'était déjà adressé à Mica en lui reprochant d'avoir acheté son enfant au vers 565.\\
Comme annoncé, je suis le papyrus d'Oxyrhynque contre le \italique{Ravennas} ; on peut y déchiffrer :% 
\citComm{\_τρικο\douteux{τ}\douteux{υ}\douteux{λ}\douteux{ο}ν·η\douteux{π}\douteux{ω}\douteux{c} \ \ ειπ\douteux{ε}\douteux{μ}οι\douteux{τ}[}%
Le point médian est, dans ce papyrus, un signe de ponctuation interrogatif. Quant à l'espace blanc de deux lettres entre \douteux{π}\douteux{ω}\douteux{c} et ειπ\douteux{ε}\douteux{μ}οι, il indique que le changement d'interlocuteur a lieu à cet endroit, et non après εἰπέ μοι comme le placent le \italique{Ravennas} et de nombreux éditeurs.
La lecture du papyrus donne donc pour le vers 743 :%
\citComm{τρικότυλον; ἢ πῶς. --- εἰπέ μοι·τί μ᾽ ἠργάσω;}%
}%
\debutNote{743}%
\note{1}{$7fr68$ de feu notre franc germinal~; c'est-à-dire qu'on pouvait échanger à tout moment contre $7,68 \times 4,5 = 34,56$ grammes d'argent pur~; soit effectivement 48 oboles (de $0,72$ grammes dans le système éginète)... mais cela, c'était avant l'invention des monnaies de singe : chez nous, bien-sûr, et déjà en $-406$ à Athènes (\italique{confer} les \italique{Grenouilles}, vers 718-733).}%
\note{2}{72 oboles soit $51,84$ grammes d'argent.}\\\passeLigne}%
\end{groupe}%
%% v. 743b-745a
\begin{groupe}%
\personnageGrec{Μίκα}%
\trimIamb{\ligne{743}\decalage{τρικότυλον; ἢ πῶς.}εἰπέ μοι· τί μ᾽ ἠργάσω;\\
ἀπέδυσας ὦναίσχυντέ μου τὸ παιδίον\\
τυννοῦτον ὄν.\testimonia{745}{(τυννοῦτο) Photius τ 552 τυννοῦτο $\sim$ Hesychius τ 1645}\\\passeLigne}%
\end{groupe}%
%% v. 745b-746
\begin{groupe}%
\personnageGrec{Μνησίλοχος}%
\trimIamb{\ligne{745}\decalage{τυννοῦτον ὄν.}τυννοῦτο; μικρὸν νὴ Δία.\apparat{745}{\lecon{τυννοῦτον} et \lecon{τυννοῦτο} \PsiDeux, Brunck : τυνοῦτον et τυνοῦτο R \espApp Γϋ. ante μικρὸν R : \lecon{nulla nota mutat. %ae
pers.%onae
} \PsiDeux \ligneApp{746} \lecon{γέγονε} \PsiDeux, Brunck : γέγονεν R \ligneApp{747} καὶ ὅσον \PsiDeux, R %: ὅσον Suid. δ 1168 
: \lecon{χὤσον} Bentley}\\
πόσ᾽ ἔτη δὲ γέγονε; τρεῖς χοᾶς ἢ τέτταρας;%
\commentaire{746}{πόσ᾽ ἔτη δὲ γέγονε; τρεῖς Χοᾶς ἢ τέτταρας;}{S'il faut en croire le scholiaste, l'on pouvait compter l'âge à partir du nombre de Grandes Dionysies écoulées (\grec{τοῦτο ἔθος ἦν Ἀττικοῖς, λέγειν τὰ ἔτη κατὰ τὸν ὑπερπίπτοντα ἀριθμὸν ἀπὸ τῶν Διονυσίων}) ; usage qui correspond chez nous à celui de compter les \italique{printemps}. Mnésiloque, filant la métaphore de l'outre et l'enfant, préfère se référer malicieusement aux \italique{Choés}, journée la plus arrosée des Anthestéries\marqueur{1} ; d'autant qu'un \grec{χοῦς} est aussi une unité de volume\marqueur{2}. Cela permet insidieusement à Mnésiloque de préciser que l'outre n'est pas si petite : elle est d'une capacité d'une dizaine de litres\marqueur{3}.}%
\debutNote{746}%
\note{1}{L'on trouvera d'abondants et utiles renseignements sur cette fête à l'article \latin{Dionysia} du \italique{Dictionnaire des antiquités grecques et romaines} de Daremberg \& Saglio. Cette fête des Anthestéries correspond à peu près à nos Saint-Vincent, Percée du Vin Jaune ou autres fêtes du beaujolais nouveau.}%
\note{2}{Le \grec{χοῦς} – ou \italique{conge} romain — correspond à 12 cotyles, soit environ 3 litres ¼.}%
\note{3}{D'une capacité \italique{théorique} de trois ou quatre conges : il est évident que Mica ne s'est pas encombrée du fardeau que serait cette outre pleine (Mnésiloque aurait d'ailleurs eu le plus grand mal à la lui chaparder). Une telle capacité peut étonner mais il faut bien voir que ces outres étaient faites de la peau complète d'un animal (souvent un bouc ou une chèvre) : ainsi, Strepsiade (\italique{Nuées}, v. 1238), trouvant que son créancier ferait une belle outre, l'estime de six conges. Il y avait des outres d'une capacité impressionnante pour jouer au jeu de l'\grec{ἀσκωλιασμός}, consistant à sauter -- \italique{clocher}, comme l'étymologie l'indique -- et tenir en équilibre sur une outre graissée. Voir l'article \latin{Uter} du \italique{Dictionnaire des antiquités grecques et romaines} de Daremberg \& Saglio. Il y a de toute façon une exagération comique : l'outre représentée sur le cratère de Würzbourg n'a certainement pas cette capacité et Mnésiloque est dans l'outrance verbale.}\\\passeLigne}%
\end{groupe}%
%% v. 747-748a
\begin{groupe}%
\personnageGrec{Μίκα}%
\trimIamb{σχεδὸν τοσοῦτον χὤσον ἐκ Διονυσίων.\testimonia{747}{Suidas δ 1168 (καὶ omisso)}%\apparat{747}{καὶ ὅσον \PsiDeux, R : ὅσον Suid. δ 1168 : \lecon{χὤσον} Bentley}%
\commentaire{747}{σχεδὸν τοσοῦτον χὤσον ἐκ Διονυσίων.}{L'on traduit généralement ce vers par quelque chose d'équivalent peu ou prou à : « C'est cela, plus le temps écoulé depuis les Dionysies ». Outre qu'inventer un anniversaire à l'outre est une billevesée indigente relevant davantage de l'humour des philologues que de celui d'Aristophane, il présente une impossibilité : essayez donc de conserver pendant trois ou quatre \italique{printemps} une outre de vin dans les conditions de l'Athènes de l'an $-2300$ avant le Réfrigérateur !
Et puisque Mica ne peut pas parler de l'âge de son outre, elle parle donc de sa \italique{capacité}. En l'occurrence, elle essaye de se défendre du chef d'accusation de biberonne porté par Mnésiloque~: sa chère petite outre, «~peut-être est-elle de cette taille [c'est-à-dire d'une capacité d'une dizaine de litres] mais elle est telle [c'est-à-dire aussi peu remplie] depuis les Dionysies~». Autrement dit, Mica au jour des Thesmophories prétend ne pas avoir touché au contenu de l'outre depuis les dernières Dionysies, ce qui représente une abstinence de sept mois ! (Bluff !)}\\
ἀλλ᾽ ἀπόδος αὐτό.\\\passeLigne}%
\end{groupe}%
%% v. 748b
\begin{groupe}%
\personnageGrec{Μνησίλοχος}%
\trimIamb{\ligne{748}\decalage{ἀλλ᾽ ἀπόδος αὐτό.}μὰ τὸν Ἀπόλλω τουτονί.
\commentaire{748}{μὰ τὸν Ἀπόλλω τουτονί.}{Contre les commentateurs qui se contentent d'une banale formule de ce genre pour inventer sur scène la présence d'une statue d'Apollon omniprésente,
% (alors que – faut-il le rappeler ? – Mnésiloque se trouve dans l'enceinte du Thesmophorion, consacré à Déméter et Perséphone et est déjà juché sur l'autel qu'il aurait donc bien du mal à faire passer pour un Ἀγυιεύς ; qu'à cela ne tienne, nos professeurs inventèrent une statue d'Apollon omniprésente pour le décor fixe des pièces de théâtre, statue qui se surajoute au décor particulier à chaque pièce), 
Willems avait déjà répondu par une note au vers 83 des \italique{Nuées}. Comme ce genre d'idées saugrenues se retrouve encore dans les dernières éditions les plus autorisées, nous avons cru faire œuvre utile en reproduisant ladite note en annexe. %Il est à espérer que sa lecture dissipe définitivement les nuées de Philologus. % ce genre de nuées. %(Il n'est d'ailleurs pas impossible qu'Aristophane fasse jurer Mnésiloque sur Apollon simplement pour citer un passage correspondant du \italique{Télèphe}.)
}\\\passeLigne}%
\end{groupe}%
%% v. 749a
\begin{groupe}%
\personnageGrec{Μίκα}%
\trimIamb{ἐμπρήσομεν τοίνυν σε.\\\passeLigne}%
\end{groupe}%
%% v. 749b-750
\begin{groupe}%
\personnageGrec{Μνησίλοχος}%
\trimIamb{\ligne{749}\decalage{ἐμπρήσομεν τοίνυν σε.}πάνυ γ᾽ ἐμπίμπρατε·\testimonia{749}{(ἐμπίμπρατε) Hesychius ε 2458 ἐμπίπραται $\sim$ Suidas ε 1018 $\sim$ Zonaras 709}\apparat{749}{ἐμπιμπρᾶται R : \lecon{ἐμπιμπράτε} (sic) \Mon}\\
αὕτη δ᾽ ἀποσφαγήσεται μάλ᾽ αὐτίκα.\\\passeLigne}%
\end{groupe}%
%% v. 751-752a
\begin{groupe}%
\personnageGrec{Μίκα}%
\trimIamb{μὴ δῆθ᾽, ἱκετεύω σ᾽· ἀλλ᾽ ἔμ᾽ ὅ τι χρῄζεις ποίει\\
ὑπέρ γε τούτου.\\\passeLigne}%
\end{groupe}%
%% v. 752b-753
\begin{groupe}%
\personnageGrec{Μνησίλοχος}%
\trimIamb{\ligne{752}\decalage{ὑπέρ γε τούτου.}φιλότεκνός τις εἶ φύσει.\\
ἀλλ᾽ οὐδὲν ἧττον ἥδ᾽ ἀποσφαγήσεται\\\passeLigne}%
\end{groupe}%
%% v. 754-755
\begin{groupe}%
\personnageGrec{Μίκα}%
\trimIamb{οἴμοι τέκνον. δός τὸ σφαγεῖον Μανία,\testimonia{754}{Photius σ 850 σφαγεῖον· [...] περισπᾶται ἡ δευτέρα συλλαβή}\apparat{754}{\lecon{δός τὸ} \PsiDeux, Lobeck ad Soph. \italique{Ai.} v.1066 : δός μοι τὸ R \espApp σφαγίον \Rac : \lecon{σφαγεῖον} \Rpc, \Mon}
\commentaire{754-7}{Scène d'égorgement de l'outre nourrissonne.}{Cette scène est illustrée sur le cratère de Würzbourg\marqueur{1} (reproduit en annexe). Amusante mise en abîme : le `σφαγεῖον' que tient Mica sur cette représentation n'est autre qu'un cratère de forme identique au cratère de Würbourg lui-même -- c'est-à-dire un vase propre à recevoir du vin pour l'y mélanger et non une coupe servant à recevoir le sang des victimes (σφαγεῖον).
% Les détournements de ce genre sont fort nombreux dans les \italique{Thesmophories}.
}%
\debutNote{754-7}%
\note{1}{Cratère apulien d'environ 370 av. J.-C., Martin von Wagner Museum der Universität Würzburg, Antiken-Sammlung, H5697. L'on pourra se reporter pour plus de détails à l'article d'Eric Csapo, \italique{A Note on the Würzburg Bell-Crater H5697} («~Telephus Travestitus~») in \italique{Phoenix}, Vol. 40, No. 4 (Winter, 1986), p. 379-392.}\\
\commentaire{754}{δός μοι τὸ σφαγεῖον.}{Le papyrus d'Oxyrhynque n°3839 ne comporte pas le μοι. Nous avons dit précédemment suivre le papyrus contre le \italique{Ravennas} et nous nous y sommes tenus. Mais il y a d'excellentes raisons à cela, outre bien-sûr le privilège de l'ancienneté.
D'abord, il n'est rien de plus courant que l'interpolation de petits mots tels que μοι (cf. le vers 1047, où un autre papyrus, daté du \nieme{II} siècle ap. J.-C., ne présente pas non plus le μοι du manuscrit de Ravenne).
Comment ce μοι est-il apparu ? il est vraisemblable que ce soit pour de pseudo-raisons métriques.
Faisons l'hypothèse que d'anciens manuscrits comportaient un \nieme{754} vers ainsi troussé :%
\citComm{οἴμοι τέκνον δός τὸ σφαγίον Μανία.}%
Noter la graphie σφαγίον (et non σφαγεῖον) : c'est celle de la première version du \italique{Ravennas} (avant correction en σφαγεῖον, faite d'une autre main s'il faut en croire Coockle, \latin{scholiorum manus, aequalis fere primae manui} selon von Velsen\marqueur{1}). σφαγίον était donc scandé : $ \cup \cup - $ ; il manquait par conséquent une syllabe longue à placer entre τέκνον et σφαγίον pour que le vers soit d'équerre. L'ajout d'un μοι était la solution la plus simple et à la portée de n'importe quel scribe.\\
Ce sont là des considérations métriques qui devaient passer bien au-dessus de la tête des scribes ayant retranscrit les \italique{Thesmophories} ? % me direz-vous peut-être. 
Détrompez-vous : la correction même du \italique{Ravennas} -- alors qu'à l'époque les deux graphies se prononçaient de la même façon -- montre bien l'intérêt pour la scansion ; sur le papyrus, Coockle releva une «~\anglais{unexplained ink like an oddly shaped perispomen accent in a second hand above} γε~» (le γε de σφαγειον) ; l'explication me paraît simple : l'annotateur confirmait que la deuxième syllabe de σφαγεῖον est longue. C'est d'ailleurs l'objet d'une glose de Photius (σ 850) :%
\citComm{σφαγεῖον· [...] περισπᾶται ἡ δευτέρα συλλαβή.}%
}%
\debutNote{754}%
\note{1}{\italique{Aristophanis Thesmophoriazusae} (1883). Adolphe von Velsen a établi définitivement dans cette édition que le second manuscrit par lequel nous sont connues les \italique{Thesmophories}, le \italique{Monacensis gr. 492} (G dans son apparat critique, \Mon dans le nôtre), n'était qu'une copie du \nieme{XV} siècle du \italique{Ravennas}, apportant quelques émendations bien vues. Son édition compare exhaustivement le \italique{Ravennas} et le \italique{Monacensis} ; ce qui est d'un grand intérêt pour déterminer si une correction du \italique{Ravennas} est ancienne ou plus récente : il suffit d'aller comparer avec la leçon du manuscrit de Münich. Son apparat critique est le plus précis et permet d'avoir la meilleure idée possible du manuscrit de Ravenne -- avec celui de l'édition Prato ; en partie une reprise du premier mais revu et amélioré en de nombreux points. Pour qui voudrait pousser le scrupule plus loin, et puisque le manuscrit de Ravenne n'a toujours pas été numérisé -- ô scandale ! le manuscrit sans qui nous ne connaîtrions que neuf des onze pièces survivantes d'Aristophane !.. --, Van Leeuwen édita en fac-similé le \italique{Ravennas} : \italique{Comoediae Undecim cum Scholiis : Codex Ravennas 137, 4, A} [$=$ \italique{Rav.} 429 (nouvelle cote)], \italique{phototypice editus} in \italique{Codices Graeci et Latini photographice depicti}, t. IX (1904).}%
ἵν᾽ οὖν τό γ᾽ αἷμα τοῦ τέκνου τοὐμοῦ λάβω.\\%
\commentaire{755-7}{τό γ᾽ αἷμα τοῦ τέκνου τοὐμοῦ λάβω --- ὕπεχ᾽ αὐτό, χαριοῦμαι γὰρ ἕν γε τοῦτό σοι --- κακῶς ἀπόλοι᾽, ὡς φθονερὸς εἶ καὶ δυσμενής.}{%En prononçant le vers 756, les commentateurs avaient imaginé jusqu'ici que : \ \ soit Mnésiloque déversait le contenu de l'outre sur l'autel\marqueur{1} ; \ \ soit il buvait d'un trait le vin\marqueur{2}. 
La mise en scène de ces vers n'a pas bien été comprise\marqueur{1}.
%Mais ces visions de la scène rentrent toutes deux en contradiction avec le texte d'Aristophane : 
Mica a demandé à \italique{récupérer} le sang de son enfant (τό γ᾽ αἷμα τοῦ τέκνου τοὐμοῦ λάβω) et puisque Mnésiloque lui concède cette unique faveur (χαριοῦμαι γὰρ ἕν γε τοῦτό σοι), il faut l'en croire.
%Or la seconde hypothèse est incompatible avec la réplique de Mnésiloque, qui lui concède cette unique faveur, χαριοῦμαι γὰρ ἕν γε τοῦτό σοι -- et l'ironie a bon dos quand il s'agit d'ignorer tout bonnement le texte que nous livre Aristophane !.. En outre, si Mnésiloque, déversait effectivement l'outre sur l'autel, comment pourrait-il dire : ὕπεχ᾽ αὐτό, "mets-le [le σφαγεῖον] sous"... sous quoi ? sous l'\italique{autel} ?! les commentateurs qui proposent cela savent-ils bien à quoi ressemble l'autel en question ? et quand bien même : il serait alors impossible à Mica d'en récupérer plus qu'une infime partie et Mnésiloque de ne pas tenir sa promesse.\\
Seulement, Mnésiloque prend au mot Mica. Le \italique{mot}, en l'occurrence, c'est λάβω : \italique{prendre}, \italique{récupérer}, dans la bouche de Mica. Mnésiloque entend \italique{recevoir}. C'est entendu : elle n'a qu'à \italique{présenter} (ὕπεχ᾽ αὐτό) le σφαγεῖον. Il lui jette alors à la figure le contenu de l'outre. C'est précisément cet instant, où Mica tend le σφαγεῖον, qui est saisi sur le cratère de Würzbourg. On voit qu'il suffit à Mnésiloque de balancer le bras pour asséner l'outre à fragmentation liquide à la tête de sa mère.
C'est ce comportement dépassant, outrepassant la mesure -- μηδὲν ἄγαν, Mnésiloque, μηδὲν ἄγαν. Ne te l'a-t-on pas assez répété ? --, qui explique l'emportement de Mica : κακῶς ἀπόλοι᾽, ὡς φθονερὸς εἶ καὶ δυσμενής. Devant Mica maculée du «~sang~» de sa fille, ajoutez la confusion induite par les nombreux quiproquos\marqueur{2}, et vous comprendrez aisément les quelques précisions que demande Critylla déboussolée à son arrivée au vers 759 : `τί τῆς ἱερείας γίγνεται;' puis `ταλαντάτη Μίκα τίς ἐξεκόρησέ σε;'.
(L'exégèse de ce passage est d'ailleurs la preuve \italique{in texto} que Mica présente des signes \italique{visibles} d'agression -- agression qui n'est donc pas uniquement psychologique. Ces signes, comme nous l'allons voir à la note s'y rapportant, sont les taches de vin qui imprègnent son vêtement.)}%
\debutNote{755-7}%
\note{1}{Les commentateurs avaient imaginé jusqu'ici que : \ \ soit Mnésiloque déversait le contenu de l'outre sur l'autel ; \ \ soit il buvait d'un trait le vin. Cette dernière option est notamment l'opinion du scholiaste : τοῦτο ἅμα εἰπὼν ἐκχέει τὸν οἶνον ὁ κηδεστὴς ἐν τῷ βωμῷ. Je passe sur les autres explications alambiquées du même acabit. %pour metteur en scène payé à l'heure de spectacle produit. (Mnésiloque en boit un peu, en verse par terre, mais pris de scrupules, finit par en tendre les dernières gouttes à Mica. Je ne caricature qu'à peine...) Soyons sérieux : une représentation d'Aristophane n'est pas le Festival de Cannes.
Austin \& Olson, conformément à leur pente, ne se prononcent pas et avancent les deux hypothèses. Il faudrait cesser de croire que la pusillanimité est une preuve d'esprit critique. «~La recherche de la vérité n'empêche pas de prendre parti~», comme disait l'autre imbécile.%
}%
%\note{2}{Austin \& Olson, conformément à leur pente, ne se prononcent pas et avancent les deux hypothèses. Il faudrait cesser de croire que la pusillanimité est une preuve d'esprit critique. «~La recherche de la vérité n'empêche pas de prendre parti~», comme disait l'autre imbécile.}%
\note{2}{Le grec dirait qu'on surajoute \italique{du feu au feu}, πῦρ ἐπὶ πῦρ.}%
\passeLigne}%
\end{groupe}%
%% v. 756
\begin{groupe}%
\personnageGrec{Μνησίλοχος}%
\trimIamb{ὕπεχ᾽ αὐτό, χαριοῦμαι γὰρ ἕν γε τοῦτό σοι.\\
\\\passeLigne}%
\end{groupe}%
%% v. 757
\begin{groupe}%
\personnageGrec{Μίκα}%
\trimIamb{κακῶς ἀπόλοι᾽, ὡς φθονερὸς εἶ καὶ δυσμενής.\\\passeLigne}%
\end{groupe}%
%% v. 758
\begin{groupe}%
\personnageGrec{Μνησίλοχος}%
\trimIamb{\ligne{757}τουτὶ τὸ δέρμα τῆς ἱερείας γίγνεται.\apparat{758-9}{\lecon{ἱερείας} \PsiDeux, R : ἱερέας Austin \& Olson}%
\commentaire{758}{τουτὶ τὸ δέρμα τῆς ἱερείας γίγνεται.}{τὸ δέρμα désigne l'outre vide – qui rappelons-le était faite d'une peau d'animal évidé. L'usage était que le prêtre reçoive certaines parties de la victime (\grec{τά ἱερώσυνα}).}\\\passeLigne}%
\end{groupe}%
%% v. 759a
\begin{groupe}%
%
\personnageGrec{Κρίτυλλα}%
\trimIamb{\ligne{758}τί τῆς ἱερείας γίγνεται;\apparat{759}{$\emptyset$ R : (ἄλλη γυνή) \Rpc (ad v.760) : \lecon{(Κρίτ.)}. Sommerstein \espApp τουτί R : \lecon{τουτί.} Bergk}%
\commentaire{759}{}{Je suis Sommerstein qui donne la parole pour ce vers à Critylla, quoiqu'aucun \italique{paragraphos} indiquant un changement d'interlocuteur ne soit présent sur le manuscrit de Ravenne (le papyrus est inexploitable sur ce point).
Faut-il prendre au premier degré la déclaration de Mnésiloque et en inférer que c'est Critylla qui jouait le rôle de prêtresse ?%
}\\\passeLigne}%
\end{groupe}%
%% v. 759b
\begin{groupe}%
\personnageGrec{Μνησίλοχος}%
\trimIamb{\ligne{759}\decalage{τί τῆς ἱερείας γίγνεται;}τουτί. λαβέ.%\apparat{759}{post τουτί interpunxit Bergk}
\\\passeLigne}%
\end{groupe}%
%% v. 760-761
\begin{groupe}%
\personnageGrec{Κρίτυλλα}%
\trimIamb{ταλαντάτη Μίκα τίς ἐξεκόρησέ σε;\apparat{760}{ἐξεκόρησεν R : \lecon{ἐξεκόρησε} \Mon \ligneApp{761}[\douteux{ }\douteux{ }\douteux{ }\douteux{ }\douteux{ }\douteux{ }\douteux{ }\douteux{ }]απη\douteux{τ}η\douteux{ν}\douteux{π}\douteux{ }ι\douteux{ }α\douteux{ }\douteux{ }υξ[ \PsiDeux \espApp σοὐξῃρήσατο R : \lecon{σοὐξηράσατο} Fritzsche}\\
τίς τὴν ἀγαπητὴν παῖδά σοὐξηράσατο;
\commentaire{760-1}{τίς ἐξεκόρησέ σε; / τίς τὴν ἀγαπητὴν παῖδά σου ᾿ξῃρήσατο;}{Faute d'avoir compris que Mnésiloque venait d'asperger Mica du «~sang~» de sa propre fille, ce passage n'a pas non plus été entendu.
Même le grand Willems -- une fois n'est pas coutume -- fut pris la main dans le sac à ordures : trompé par un jeu de mots en latin\marqueur{1} qu'il a cru pouvoir rapprocher de ce distique, il proposa comme interprétation une paillardise qui n'a pas lieu d'être, et qui d'ailleurs n'y est pas.\\
Figurons-nous la scène : l'outre a été égorgée et déversée sur Mica ; arrive Critylla à qui Mnésiloque jette l'outre vide. Critylla avec l'«~enfant~» éviscéré ; Mica maculée de sang face à elle.%
\citComm{Critylla --- τίς ἐξεκόρησέ σε;}%
Le verbe ἐκκορέω, \italique{balayer}, est fréquemment utilisé comme synonyme de \italique{maltraiter}, \italique{dévaster}. Cf., par exemple\marqueur{2}, Hésychius (ε 1453 ; κ 3067) :%
\citComm{ἐκκοροῦσι· φθείρουσιν, ἐκκαλλύνουσιν.\\
κορέων, κορέσκων· ἐξυβρίζων.}%
Photius (ε 412) :%
\citComm{ἐκκορηθῆναι· ἐκκαλλυνθῆναι· τὸ γὰρ κορεῖν καλλύνειν καὶ σαίρειν.}%
Ainsi dans la \italique{Paix} (vers 58-59), Trygée s'écrie\marqueur{3} :%
\citComm{Ὦ Ζεῦ, τί ποτε βουλεύει ποεῖν;\\
Κατάθου τὸ κόρημα· μὴ ᾽κκόρει τὴν Ἑλλάδα\\
\francais{Ô Zeus, que prétends-tu faire ?\\
Dépose ton balai : ne balaie pas l'Hellade de la surface de la Terre.}}%
Il n'y a à notre \nieme{760} aucun sens obscène, ni là ni ailleurs\marqueur{4}. La question de Mica revient à ceci : «~Qui t'a fourragée\marqueur{5} ?~», «~Qui t'a ruiné [ton habit] ?~»\\
Mais ἐξεκόρησε n'est pas mis là fortuitement ; il fait entendre le mot κόρη, la «~fille~» de Mica, avec la reprise au vers suivant : τὴν ἀγαπητὴν παῖδά σου.
Le sens premier de la question de Mica est donc de savoir qui l'a ainsi salie ; le sens allusif qui l'a \italique{privée de sa fille} (ἐξ-κόρη).\\
Pour être sûr que le spectateur saisisse bien l'amphibologie, Aristophane revient à la charge au vers 761 : τίς τὴν ἀγαπητὴν παῖδά σου ᾿ξηράσατο\marqueur{6};
Le pronom σύ a volontairement été rejeté avant le verbe pour laisser toute ambiguïté à la crase : σοὐξηράσατο $=$ σοι ἐξηράσατο\marqueur{7} d'une part ; σοὐξῃρήσατο $=$ σοῦ ἐξῃρήσατο d'autre part.\\
σοι ἐξερᾶν est le pendant de ἐκκορεῖν σε du vers précédent : Critylla demande qui a osé \italique{vider}, \italique{déverser}\marqueur{8} sur Mica ce que sa fille chérie avait dans le ventre.\\
σοῦ ἐξαιρεῖν : en raison de la proximité phonétique entre le `η' et le `α', le spectateur devait aisément entendre -- pour peu que l'acteur de Critylla s'y prêtât : τίς τὴν ἀγαπητὴν παῖδά σου ἐξῃρήσατο; «~Qui t'a ravi ta fille chérie ?~» ; écho au ἐκ-κόρη subrepticement ouï au vers précédent.\\
Puisque le sens premier de la question de Critylla %-- celui de ἐκκορεῖν σε et celui en conformité avec l'action -- 
est bien σοι ἐξερᾶν, je reprends la conjecture de Friztsche (lire σοὐξηράσατο à la place du σοὐξῃρήσατο du \italique{Ravennas}) d'autant que ἐξῃρήσατο serait un aoriste inhabituel\marqueur{9}.%
}%
\debutNote{760-1}%
\note{1}{\latin{Quis te devirginavit} ?}%
\note{2}{Pour d'autres exemples, qui ne manquent pas, se reporter aux \italique{testimonia} des éditions des lexiques sous-cités}%
\note{3}{Voir Taillardat (§629) pour un exposé plus fourni sur cette image.}%
\note{4}{Suétone (τρὶς κεκορημένε· οἱονεὶ πολλάκις ἐκσεσαρωμένος) a donc raison contre Taillardat (Περὶ βλασφημιῶν, β 260 $=$ Eustathe p.1542.47 $=$ fr. Anacréon 366 \italique{Poetae Melici Graeci} Page) et beaucoup d'autres : l'imprécation d'Anacréon `τρὶς κεκορημένε' lancée au Thrace Smerdiès signifie chastement «~triple ordure~».}%
\note{5}{Je m'amuse, mais comme j'appréhende déjà que cette badinerie soit prétexte à la meute libidineuse des commentamenteurs pour déchaîner leur rut philololubrique, je m'empresse de préciser : à aucun moment Critylla n'évoque l'idée que le sang qui couvre Mica puisse être hymenéen. Et d'abord parce que ce serait inepte : Mica est mère. 
Et que ceux qui tiennent tant à cette théorie exhibent seulement un passage -- ou même seulement un indice autre que des polyçonneries -- dans lequel le mot παῖδα ait pu désigner l'hymen. En attendant, il serait bon prophylactiquement d'expurger -- ἐξερᾶν, mais j'anticipe... -- du Bailly, à l'entrée ἐκκορέω, l'acception \italique{déflorer} afférente à ce passage.
%Mais si vous tenez tant à cette théorie, ô Philologus, avant d'ouvrir cette bouche qui vomit en surabondance ce que vous devriez celer de votre cœur, mentionnez-moi seulement un passage -- ou même, exhibez-moi seulement un indice, à défaut de vos polyçonneries -- dans lequel le mot παῖδα ait pu désigner l'hymen. En attendant, il serait bon prophylactiquement d'expurger -- ἐξερᾶν, mais j'anticipe... -- du Bailly, à l'entrée ἐκκορέω, l'acception \italique{déflorer} afférente à ce passage.
}%
\note{6}{Je fais mienne la conjecture de Fritzsche : 'ξηράσατο (de ἐξερᾶν) à la place du 'ξῃρήσατο (de ἐξαιρεῖν) présent dans le manuscrit de Ravenne. (Je m'en explique en fin de note.)}%
\note{7}{Cf. les vers 158 et 624.}%
\note{8}{Peut-être même faut-il prendre ἐξερᾶν au sens de \italique{régurgiter}, \italique{recracher}, et comprendre que Mnésiloque, ajoutant l'insulte à l'outrage, après avoir empli sa bouche de vin, l'aura \italique{recraché} à la tête de Mica.}%
%\note{7}{\^ \^}%
\note{9}{Impossible selon le rigoriste O. Lautensach, p. 128 de sa grammaire \italique{Die Aoriste bei den Attischen Tragikern und Komikern} (1911) : \allemand{Aoriste auf} -ησα\allemand{. Der im Etymologicum Magnum p.777, 53 als "ungriechisch" verworfene und von Lobeck zu Phrynicus p. 716ff. aus der späteren Gräcität nachgewiesene Aor. }ᾕρησα \allemand{ist unmöglich bei Ar. \italique{Thesm.} 761 in} σοὐξῃρήσατο\allemand{ zu dulden, wofür Meineke} σου διεχρήσατο\allemand{, Fritzsche }σου ᾽ξηράσατο\allemand{, Dindorf} σου ᾽ξαφείλετο\allemand{ schreibt, während Lobeck, v. Velsen und Blaydes den Vers für unecht erklären.}}\\\passeLigne}%
\end{groupe}%
%% v. 762-764
\begin{groupe}%
\personnageGrec{Μίκα}%
\trimIamb{ὁ πανοῦργος οὗτος. ἀλλ᾽ ἐπειδήπερ πάρει,\\
φύλαξον αὐτόν, ἵνα λαβοῦσα Κλεισθένη\\
τοῖσιν πρυτάνεσιν ἃ πεποίηχ᾽ οὗτος φράσω.\\\passeLigne}%
\end{groupe}%
\endnumbering
\end{Leftside}
\begin{Rightside}
\beginnumbering
%% v. 718-725
\begin{groupe}%
\chantChoeurFr{Car toi, avec tes impiétés, nous te châtierons }%
\end{groupe}%
%% v. 718-725
\begin{groupe}%
\chantChoeurFr{de ta conduite, comme de juste.}%
\end{groupe}%
%% v. 718-725
\begin{groupe}%
\tetramTrochFr{Et puis la Fortune, changeant brusquement d'humeur aura vite fait de poursuivre un méchant homme }%
\end{groupe}%
%% v. 718-725
\begin{groupe}%
\chantChoeurFr{qui a pris le chemin du crime.\\\passeLigne}%
\end{groupe}%
%% v. 726-727
\begin{groupe}%
\personnageFr{Coryphée}%
\tetramTrochFr{\ligne{726}\small{Mais prends ces femmes avec toi, et va chercher du bois.\\
Il faut le brûler vif, ce fourbe, l'anéantir au plus tôt par le \fin feu.}\\\passeLigne}%
\end{groupe}%
%% v. 728-729
\begin{groupe}%
\personnageFr{Mika}%
\trimIambFr{Allons chercher les fagots, Mania.\\
Je ferai de toi un tison aujourd'hui..\\\passeLigne}%
\end{groupe}%
%% v. 730-732
\begin{groupe}%
\personnageFr{Mnésiloque}%
\trimIambFr{Allumez, brûlez. Toi, ma petite, ôte vite ta crétique\\
et n'accuse de ta mort que ta mère.\\
Qu'est ceci ? La fillette n'est plus qu'une outre}%
\end{groupe}%
%% v. 733-735
\begin{groupe}%
\trimIambFr{pleine de vin, ... et avec cela chaussée de persiques.\\
Ô femmes archi-enfiévrées et archi-biberonnes,\\
ingénieuses à trouver partout matière à boire !}%
\end{groupe}%
%% v. 736-738
\begin{groupe}%
\trimIambFr{Ô providence des débitants, et notre fléau à nous !\\
Fléau aussi des ustensiles de ménage, et fléau de la trame !
(Les femmes reviennent avec des fagots.)\\\passeLigne}%
\end{groupe}%
%% v. 739
\begin{groupe}%
\personnageFr{Mika}%
\trimIambFr{Entasse force fagots autour de lui, Mania.\\\passeLigne}%
\end{groupe}%
%% v. 740-741a
\begin{groupe}%
\personnageFr{Mnésiloque}%
\trimIambFr{Entasse, c'est bon. Mais toi, réponds à ceci.\\
Tu te dis mère de cet enfant ?\\\passeLigne}%
\end{groupe}%
%% v. 741b-742a
\begin{groupe}%
\personnageFr{Mika}%
\trimIambFr{\ligne{741}\espace{2.8cm}Dix mois durant\\
je l'ai porté.\\\passeLigne}%
\end{groupe}%
%% v. 742b
\begin{groupe}%
\personnageFr{Mnésiloque}%
\trimIambFr{\ligne{742}\espace{1.5cm}Tu l'as porté, toi ?\\\passeLigne}%
\end{groupe}%
%% v. 742c
\begin{groupe}%
\personnageFr{Mika}%
\trimIambFr{\ligne{742}\espace{3.5cm}Oui, par Artémis.\\\passeLigne}%
\end{groupe}%
%% v. 743a
\begin{groupe}%
\personnageFr{Mnésiloque}%
\trimIambFr{(Il porte l'outre à ses lèvres.) Du trois cotyles ? Ou de \fin quel prix ?\\\passeLigne}%
\end{groupe}%
%% v. 743b-745a
\begin{groupe}%
\personnageFr{Mika}%
\trimIambFr{\ligne{743}\espace{4.2cm}Dis-moi : quel tour vas-tu me faire ?\\
Impudent, tu as dévêtu mon enfant,\\
qui est si menu.\\\passeLigne}%
\end{groupe}%
%% v. 745b-746
\begin{groupe}%
\personnageFr{Mnésiloque}%
\trimIambFr{\ligne{745}\espace{2.2cm}Si menu ? (Il le regarde.) Il est petit, par Zeus.
Et quel âge a-t-il ? trois ou quatre choés ?\\\passeLigne}%
\end{groupe}%
%% v. 747-748a
\begin{groupe}%
\personnageFr{Mika}%
\trimIambFr{À peu près, mais elle est ainsi depuis les dionysiaques.\\
Allons ! rends-le-moi.\\\passeLigne}%
\end{groupe}%
%% v. 748b
\begin{groupe}%
\personnageFr{Mnésiloque}%
\trimIambFr{\ligne{748}\espace{2.8cm}Non, par le glorieux Apollon.%
%\commentaire{748}{μὰ τὸν Ἀπόλλω τουτονί.}{Cf. la note sur Nuées, v. 83, t. I, p.380.\eff{À reproduire en annexe !}}.
\\\passeLigne}%
\end{groupe}%
%% v. 749a
\begin{groupe}%
\personnageFr{Mika}%
\trimIambFr{Alors nous allons te brûler.\\\passeLigne}%
\end{groupe}%
%% v. 749b-750
\begin{groupe}%
\personnageFr{Mnésiloque}%
\trimIambFr{\ligne{749}\espace{3cm}Parfaitement, brûlez.\\
Mais j'égorge celle-ci à l'instant.\\\passeLigne}%
\end{groupe}%
%% v. 751-752a
\begin{groupe}%
\personnageFr{Mika}%
\trimIambFr{Non, non, je t'en conjure. Fais à moi tout ce que tu voudras, 
plutôt qu'à elle.\\\passeLigne}%
\end{groupe}%
%% v. 752b-753
\begin{groupe}%
\personnageFr{Mnésiloque}%
\trimIambFr{\ligne{752}\espace{2.5cm}Tu as pour tes enfants une tendresse de naturel.
Mais elle n'en sera pas moins égorgée.\\\passeLigne}%
\end{groupe}%
%% v. 754-755
\begin{groupe}%
\personnageFr{Mika}%
\trimIambFr{Ah! mon enfant ! Donne-moi une patère, Mania,\\
que je recueille au moins le sang de ma fille..\\\passeLigne}%
\end{groupe}%
%% v. 756
\begin{groupe}%
\personnageFr{Mnésiloque}%
\trimIambFr{Tiens-la là-dessous. Je veux bien t'accorder cette seule grâce.
(Il asperge Mika du vin de l'outre.)\\\passeLigne}%
\end{groupe}%
%% v. 757
\begin{groupe}%
\personnageFr{Mika}%
\trimIambFr{Puisses-tu périr de male mort! Que tu es jaloux et \fin mauvais !\\\passeLigne}%
\end{groupe}%
%% v. 758
\begin{groupe}%
\personnageFr{Mnésiloque}%
\trimIambFr{Cette peau revient à la prêtresse.\\\passeLigne}%
\end{groupe}%
%% v. 759a
\begin{groupe}%
\personnageFr{Kritylla}%
\trimIambFr{\ligne{758}Qu'est-ce qui revient à la prêtresse ?\\\passeLigne}%
\end{groupe}%
%% v. 759b
\begin{groupe}%
\personnageFr{Mnésiloque}%
\trimIambFr{\ligne{759}\espace{3cm}Tiens, prends. (Il lui jette l'outre vide.)\\\passeLigne}%
\end{groupe}%
%% v. 760-761
\begin{groupe}%
\personnageFr{Kritylla}%
\trimIambFr{Infortunée Mica, qui t'a souillée ?\\
Qui t'a soustrait ta fille chérie ?\\\passeLigne}%
\end{groupe}%
%% v. 762-764
\begin{groupe}%
\personnageFr{Mika}%
\trimIambFr{C'est ce fourbe. Mais puisque te voilà,\\
garde-le ; que j'aille, en compagnie de Clisthène,\\
raconter aux prytanes ce qu'il a fait.\\\passeLigne}%
\end{groupe}%
\endnumbering
\end{Rightside}
\end{pages}
\Pages
\clearpage\ifodd\value{page}\hbox{}\newpage\fi
\enTeteG{ΕΠΕΙΣΟΔΙΟΝ Γ' / ΠΑΡΑΒΑΣΙΣ --- Codex Ravennas 429, f. 171v (Classense 173v) --- vv. 765–799}
\begin{adjustwidth}{-1.5cm}{-3.0cm}
\begin{adjustbox}{max totalheight=21.0cm, max width=16.5cm, keepaspectratio}
\begin{minipage}{16.5cm}
{\fontsize{7.2pt}{8.5pt}\selectfont
\noindent\textbf{Folio 19 --- Codex Ravennas 429, \texttt{f. 171v} (\texttt{Classense 173v}) : ΕΠΕΙΣΟΔΙΟΝ Γ' / ΠΑΡΑΒΑΣΙΣ, vv. 765–799}\hfill \textit{Témoins :} $\text{R}^1, \Rpc, \Rdeux$ (\Mon)\par\vspace{2pt}
\hrule height 0.4pt\vspace{2.5pt}
}
{\fontsize{8.1pt}{9.5pt}\selectfont
\ligneRav{765}{μ\auDessus{ν}{Η}}{ἄγε δὴ τίς ἔσται μηχανὴ σωτηρίας;}{}
\ligneRav{}{}{τίς πεῖρα, τίς ἐπίνοι᾽; ὁ μὲν γὰρ αἴτιος}{}
\ligneRav{767}{}{κἄμ᾽ \difference{εἰσκυλίσας εἰς} τοιαυτὶ πράγματα}{\textbf{767} ἐσκυλίσας ἐς R : \lecon{εἰσκυλίσας εἰς} edd. pler.}
\ligneRav{768}{}{οὐ φαίνετ᾽, οὔπω. φέρε \difference{τίν᾽ οὖν ἂν} ἄγγελον}{\textbf{768} τίν᾽ οὖν R : τίν᾽ οὖν νῦν Biset : \lecon{τίν᾽ οὖν ἂν} Porson (\italique{Aristophanica} p.218)}
\ligneRav{}{}{πέμψαιμ᾽ ἐπ᾽ αὐτόν; οἶδ᾽ ἐγὼ καὶ δὴ πόρον}{}
\ligneRav{770}{}{ἐκ τοῦ Παλαμήδους· ὡς ἐκεῖνος, τὰς πλάτας}{}
\ligneRav{771}{}{ῥίψω γράφων. ἀλλ᾽ οὐ πάρεισιν αἱ πλάται·}{\textbf{771} \lecon{οὐ πάρεισιν} R : οὐ γάρ εἰσιν Fritzsche}
\ligneRav{772}{}{πόθεν οὖν γένοιτ᾽ ἄν μοι πλάται. πόθεν; πόθεν;}{\textbf{772} \lecon{γένοιτ᾽} R : γένοιντ᾽ Grynaeus \espApp πόθεν; R : \lecon {πόθεν; πόθεν;} Scaliger}
\ligneRav{773}{}{τί δ᾽ ἂν \difference{εἰ ταδὶ τἀγάλματ᾽} ἀντὶ τῶν πλατῶν}{\textbf{773} εἶτα διὰ τἀγάλματ᾽ R : ταδὶ τἀγάλματ᾽ Suid. : \lecon{εἰ ταδὶ τἀγάλματ᾽} Ellebodius}
\ligneRav{}{}{γράφων διαρρίπτοιμι; βέλτιον πολύ·}{}
\ligneRav{775}{}{ξύλον γέ τοι καὶ ταῦτα κἀκεῖν᾽ ἦν ξύλον·}{}
\ligneRav{}{}{\hspace{1.1em}ὦ χεῖρες ἐμαὶ}{}
\ligneRav{777}{}{\hspace{1.1em}ἐγχειρεῖν \difference{χρῆν} ἔργῳ πορίμῳ·}{\textbf{777} χρή R, Suid. π 2064 : \lecon{χρῆν} Bentley}
\ligneRav{}{}{\hspace{1.1em}ἄγε δὴ πινάκων ξεστῶν δέλτοι,}{}
\ligneRav{}{}{\hspace{1.1em}δέξασθε σμίλης ὁλκοὺς}{}
\ligneRav{780}{}{\hspace{1.1em}κήρυκας ἐμῶν μόχθων· οἴμοι}{}
\ligneRav{781}{}{\hspace{1.1em}τουτὶ \difference{τορῶ} μοχθηρόν·}{\textbf{781} τὸ ῥῶ R~: \lecon{τορῶ} Bibi ap. Hesych. τ 1162}
\ligneRav{782}{}{\hspace{1.1em}`\difference{χώρει}', `χώρει'. ποίαν αὔλακα;}{\textbf{782} \lecon{χώρει} bis R : χωρεῖ bis Brunck}
\ligneRav{783}{}{\hspace{1.1em}βάσκετ᾽ ἐπείγετε πάσας \difference{καθ᾽ ὁδοὺς}}{\textbf{783} καθόδους R : \lecon{καθ᾽ ὁδοὺς} Biset}
\ligneRav{784}{}{\hspace{1.1em}κείνᾳ ταύτᾳ· ταχέως χρή·}{\textbf{784} ταῦτα R : ταύτᾳ Grynaeus (\italique{hac} Divus)}
\ligneRav{785}{χορ/}{ἡμεῖς τοίνυν ἡμᾶς αὐτὰς εὖ λέξωμεν παραβᾶσαι,}{}
\ligneRav{}{}{καίτοι πᾶς τις τὸ γυναικεῖον φῦλον κακὰ πόλλ᾽ ἀγορεύει,}{}
\ligneRav{}{}{ὡς πᾶν ἐσμὲν κακὸν ἀνθρώποις κἀξ ἡμῶν ἐστιν ἅπαντα,}{}
\ligneRav{788}{}{ἔριδες νείκη \difference{στάσις ἀργαλέα} λύπη πόλεμος. φέρε δή νυν,}{\textbf{788} στάσεισ· ἀργαλέα· R : \lecon{στάσις ἀργαλέα} Scaliger \espApp εἰ καὶ κακόν R : \lecon{εἰ κακόν} Zanetti}
\ligneRav{789}{}{\difference{εἰ κακόν} ἐσμεν, τί γαμεῖθ᾽ ἡμᾶς, εἴπερ ἀληθῶς κακόν ἐσμεν,}{\textbf{789} εἰ καὶ κακόν R : \lecon{εἰ κακόν} Zanetti}
\ligneRav{790}{}{κἀπαγορεύετε μήτ᾽ ἐξελθεῖν μήτ᾽ ἐκκύψασαν ἁλῶναι,}{\textbf{790} ἐγκύψασαν R : \lecon{ἐκκύψασαν} Reiske}
\ligneRav{}{}{ἀλλ᾽ οὑτωσὶ πολλῇ σπουδῇ τὸ κακὸν βούλεσθε φυλάττειν;}{}
\ligneRav{792}{}{κἂν ἐξέλθῃ τὸ γύναιόν ποι, \difference{κᾆθ᾽ εὕρητ᾽} αὐτὸ θύρασιν,}{\textbf{792} κᾆθ᾽ εὕροιτ' R : \lecon{κᾆθ᾽ εὕρητ᾽} Porson}
\ligneRav{793}{}{μανίας \difference{μαίνεσθ᾽}, οὓς \difference{χρῆν} σπένδειν καὶ χαίρειν, εἴπερ ἀληθῶς}{\textbf{793} μαίνεθ᾽ R : \lecon{μαίνεσθ᾽} \Mon \espApp χρή R : \lecon{χρῆν} Brunck}
\ligneRav{794}{}{ἔνδοθεν \difference{ηὕρετε} φροῦδον τὸ κακὸν καὶ μὴ \difference{κατελαμβάνετ᾽} ἔνδον·}{\textbf{794} εὕρετε R : \lecon{ηὕρετε} Meineke \espApp καταλαμβάνετ᾽ R : \lecon{κατελαμβάνετ᾽} Brunck}
\ligneRav{795}{}{κἂν \difference{καταδάρθωμεν} ἐν ἀλλοτρίων παίζουσαι καὶ κοπιῶσαι,}{\textbf{795} καταδαρθῶμεν R : \lecon{καταδάρθωμεν} Dindorf}
\ligneRav{}{}{πᾶς τις τὸ κακὸν τοῦτο ζητεῖ περὶ τὰς κλίνας περινοστῶν·}{}
\ligneRav{797}{}{κἂν ἐκ θυρίδος παρακύπτωμεν, \difference{τὸ κακὸν ζητεῖτε θεᾶσθαι}·}{\textbf{797} ζητεῖ τὸ κακὸν τεθεᾶσθαι R : \lecon{τὸ κακὸν ζητεῖτε θεᾶσθαι} J. Kaye}
\ligneRav{798}{}{κἂν \difference{αἰσχυνθεῖσ᾽} (αἰσχυνθ\inRas{η}{εὶ}ς \Rpc) ἀνακύψῃ,}{\textbf{798} αἰσχυνθεὶς (ex αἰσχυνθῆς) \Rpc : \lecon{αἰσχυνθεῖσ᾽} Zanetti}
\ligneRav{799}{}{αὖθις \difference{τὸ κακὸν παρακύψαν ἰδεῖν}. οὕτως ἡμεῖς ἐπιδήλως}{\textbf{799} παρακύψασαν ἰδεῖν τὸ κακὸν R : \lecon{τὸ κακὸν παρακύψαν ἰδεῖν} Dindorf}
}
\vspace{3pt}\noindent\rule{4.5cm}{0.4pt}\par\vspace{1.5pt}
{\fontsize{7.0pt}{8.3pt}\selectfont
\noindent\textbf{Loci paralleli (vv. 765–799) ---} \textbf{769-75} Suidas π 45 (οἶδα δ' pro οἶδ᾽, πότερον pro πόρον et γ' ἔκειντο pro γένοιντ᾽ ἄν) \espTest \textbf{770-1} Suidas π 1702 \espTest \textbf{777} Suidas π 2064 \espTest \textbf{778-80} (πινάκων -- μόχθων) Pollux X 58 (ἄγε δὴ et δέλτοι omissis) \espTest \textbf{778-9} Suidas σ 742\par
}
\vspace{2.5pt}\noindent\rule{4.5cm}{0.4pt}\par\vspace{1.5pt}
{\fontsize{6.9pt}{8.2pt}\selectfont
\noindent\textbf{Suite (en commentaire) de l'apparat critique \& notes paléographiques (\texttt{f. 171v}) ---} \textbf{Témoin papyrologique :} \textbf{\PsiDeux{} (\italique{P.Oxy.} LVI 3839, fr. 1)} s'achève au \textbf{v. 766}. --- \textbf{Conjecture éclatante de \textbf{Bibi} au v. 781 (τὸ ῥῶ R $\to$ τορῶ \textbf{Bibi}) :} Là où tous les éditeurs suivaient la coupure de mots de R (τὸ ῥῶ, « ce rhô »), Louis Hauseux restitue sans modifier une seule lettre de l'onciale ΤΟΡΩ le verbe \lecon{τορῶ} (« je grave, je creuse au burin », d'après la glose d'Hésychius $\tau\ 1162$ τορεῖ· τορνεύει, γλύφει, τιτρώσκει), répondant directement à la question ποίαν αὔλακα; (« Quel sillon ? — Je grave : \italique{Va, va !} »). --- \textbf{Correction \italique{in rasura} au v. 798 :} \Rac{} avait écrit αἰσχυνθῆς, corrigé sur grattage par \Rpc{} en \lecon{αἰσχυνθ\inRas{η}{εὶ}ς} (pour le participe féminin élidé αἰσχυνθεῖσ' restitué par Zanetti). --- \textbf{v. 772 :} \lecon{γένοιτ'} R (γένοιντ' Grynaeus) \espApp πόθεν; R (haplographie finale) : \lecon{πόθεν; <πόθεν;>} Scaliger. --- \textbf{v. 781 :} τὸ ῥῶ R : \lecon{τορῶ} \lecon{Bibi} ap. Hésychius $\tau\ 1162$. --- \textbf{v. 797 :} ζητεῖ τὸ κακὸν τεθεᾶσθαι R (inversion et médivision) : \lecon{τὸ κακὸν ζητεῖτε θεᾶσθαι} J. Kaye.\par
}
\end{minipage}
\end{adjustbox}
\end{adjustwidth}
\newpage
\enTeteD{Codex Ravennas 429 (Biblioteca Classense) --- Fac-similé du f. 171v (vv. 765–799)}
\begin{adjustwidth}{-1.5cm}{-3.0cm}
\begin{center}
\includegraphics[width=16.5cm,height=20.5cm,keepaspectratio]{/workspaces/Thesmophories/dossier_ravennas_thesmophories/folios_individuels_ff162v-177r/19_f171v_classense173v_vv0765-0799.jpg}\par
\vspace{2pt}
{\scriptsize\textbf{Codex Ravennas 429, f. 171v (Classense 173v)} --- ΕΠΕΙΣΟΔΙΟΝ Γ' / ΠΑΡΑΒΑΣΙΣ, \textbf{vv. 765–799} ($\text{R}^1, \Rpc, \Rdeux$ (\Mon)).}
\end{center}
\end{adjustwidth}
\newpage
\clearpage\ifodd\value{page}\hbox{}\newpage\fi
\enTeteG{ΘΕΣΜΟΦΟΡΙΑΖΟΥΣΑΙ --- Texte grec établi (f. 171v, vv. 765–799)}
\enTeteD{LES THESMOPHORIES --- Traduction française (f. 171v, vv. 765–799)}
\begin{pages}
\begin{Leftside}
\beginnumbering
%% v. 765-775
\begin{groupe}%
%
\personnageGrec{Μνησίλοχος}%
\trimIamb{ἄγε δὴ τίς ἔσται μηχανὴ σωτηρίας;\\
τίς πεῖρα, τίς ἐπίνοι᾽; ὁ μὲν γὰρ αἴτιος\\
κἄμ᾽ εἰσκυλίσας εἰς τοιαυτὶ πράγματα\apparat{767}{ἐσκυλίσας ἐς R : \lecon{εἰσκυλίσας εἰς} edd. pler. \ligneApp{768} τίν᾽ οὖν R : τίν᾽ οὖν νῦν Biset : \lecon{τίν᾽ οὖν ἂν} Porson (\italique{Aristophanica} p.218)}%
\commentaire{767}{ἐσκυλίσας ἐς τοιαυτὶ πράγματα.}{Pour la métaphore de «~rouler dans les malheurs~», voir le vers 650 ainsi que ce qu'en dit Taillardat (§558).}\\
οὐ φαίνετ᾽, οὔπω. φέρε τίν᾽ οὖν ἂν ἄγγελον\\%\apparat{768}{%\lecon{φαίνετ᾽, οὔπω} R : φαίνεταί πω Dobree \espApp τίν᾽ οὖν R : τίν᾽ οὖν νῦν Biset : \lecon{τίν᾽ οὖν ἂν} Porson (\italique{Aristophanica} p.218)}\\% : alii alia}\\
πέμψαιμ᾽ ἐπ᾽ αὐτόν; οἶδ᾽ ἐγὼ καὶ δὴ πόρον\testimonia{769-75}{Suidas π 45 (οἶδα δ' pro οἶδ᾽, πότερον pro πόρον et γ' ἔκειντο pro γένοιντ᾽ ἄν)}\\
ἐκ τοῦ Παλαμήδους· ὡς ἐκεῖνος, τὰς πλάτας\testimonia{770-1}{Suidas π 1702}%
\commentaire{770}{ἐκ τοῦ Παλαμήδους·}{Euripide représenta son \italique{Palamède} en 415 av. J.-C. dans une tétralogie comportant aussi l'\italique{Alexandros}, les \italique{Troyennes} et le drame satyrique \italique{Sysiphe} (il fut classé second). Eschyle et Sophocle représentèrent aussi un \italique{Palamède}. Ce héros aux-mille-tours ingénieux de l'épopée troyenne (et pourtant absent de l'\italique{Iliade} ; sans doute pour ne pas faire d'ombre à Ulysse) avait procuré à l'armée grecque d'« immenses bienfaits »\marqueur{1} : il apaisa une famine à Aulis et des troubles causés à propos du partage des provisions, en enseignant aux Grecs les « lettres phéniciennes », leur apprit le jeu de dés pour se distraire, les mesures et les jetons de vote, ce qui lui valut un grand renom. Le scholiaste de l'\italique{Oreste} (v. 432) continue :%
\citComm{ἐπὶ τούτῳ δὲ φθονήσαντες οἱ περὶ Ἀγαμέμνονα καὶ Ὀδυσσέα καὶ Διομήδην τοιόνδε τι σκευωροῦσι κατ᾽ αὐτοῦ. λαβόντες γὰρ Φρύγα αἰχμάλωτον χρυσίον κομίζοντα Σαρπηδόνι ἠνάγκασαν γράψαι Φρυγίοις γράμμασι περὶ προδοσίας ὡς παρὰ Πριάμου πρὸς Παλαμήδην. καὶ τοῦτον μὲν φονεύουσι, θεράποντα δὲ Παλαμήδους πείθουσι χρήμασιν ἅμα τοῖς Τρωϊκοῖς χρήμασι καὶ τὸ γραφὲν πινάκιον ὑπὸ τὴν κλίνην θέσθαι Παλαμήδους. αὐτοὶ δὲ παρελθόντες προδοσίαν κατήγγελλον τοῦ ἥρωος καὶ φωραθῆναι τὴν σκηνὴν ἐκέλευον. εὑρεθέντος δὲ τοῦ πινακίου καὶ τῶν χρημάτων ὑπὸ τὴν κλίνην λίθοις φονεύεται Παλαμήδης\\
\francais{Devenus jaloux de ce fait, Agamemnon, Ulysse et Diomède imaginèrent ce genre de procédé : ayant fait prisonnier un Phrygien qui apportait de l'or à Sarpédon, ils le forcèrent à écrire en caractères phrygiens une lettre parlant de trahison, comme si elle avait été adressée par Priam à Palamède, puis ils le tuèrent. Ils corrompirent alors un serviteur de Palamède, pour qu'il place en même temps que l'or phrygien la tablette écrite sous la tente de son maître. Eux-mêmes, ayant comparu [dans l'assemblée des chefs], accusèrent le héros de trahison et demandèrent qu'on fouille la tente. Après la découverte de la tablette et de l'or sous la tente, Palamède fut lapidé.}}%
Anxieux de faire connaître la mort injuste de son fils à Nauplios, Œax, le frère de Palamède, imaginait le stratagème ici repris par Mnésiloque\marqueur{2} :
\citComm{ὁ γὰρ Ἐὐριπίδης ἐν τῷ Παλαμήδει ἐποίησε τὸν Οἴακα τὸν ἀδελφὸν Παλαμήδους ἐπιγράψαι εἰς τὰς ναῦς τὸν θάνατον αὐτοῦ, ἵνα φερόμεναι αὗται ἔλθωσιν εἰς τὸν Ναύπλιον τὸν πατέρα αὐτοῦ καὶ ἀπαγγείλωσι τὸν θάνατον αὐτοῦ.\\
\francais{Euripide, dans le \italique{Palamède}, a représenté Oiax, frère de Palamède, gravant sur des navires [le récit de] sa mort, afin qu'ils soient portés jusque chez Nauplios et lui apprennent sa mort.}}}%
\debutNote{770}
\note{1}{On retrouve dans une scholie à son \italique{Oreste}, au vers 432, les grandes lignes de la version euripidéenne du drame. J'en reprends la traduction de Jouan \& van Looy, à retrouver dans leur édition des fragments d'Euripide aux Belles Lettres, tome VIII, \nieme{2} partie.}
\note{2}{Scholie au vers 771 des \italique{Thesmophories}.}
\commentaire{770}{ὡς ἐκεῖνος,}{Œax, frère de Palamède.}\\
ῥίψω γράφων. ἀλλ᾽ οὐ πάρεισιν αἱ πλάται.\\%\apparat{771}{\lecon{οὐ πάρεισιν} R : οὐ γάρ εἰσιν Fritzsche}\\
πόθεν οὖν γένοιτ᾽ ἄν μοι πλάται. πόθεν; πόθεν;\apparat{772}{\lecon{γένοιτ᾽} R : %γ᾽ ἔκειντό Suid. π 45 : 
γένοιντ᾽ Grynaeus %(\italique{fierent} Divus) 
\espApp πόθεν; R : %πόθεν πλάται; Biset : 
\lecon {πόθεν; πόθεν;} Scaliger %: πλάται πόθεν; Sandbach : alii alia
\ligneApp{773} εἶτα διὰ τἀγάλματ᾽ R : ταδὶ τἀγάλματ᾽ Suid. %σ 742 
: \lecon{εἰ ταδὶ τἀγάλματ᾽} Ellebodius}%
\commentaire{772}{πόθεν οὖν γένοιντ᾽ ἄν μοι πλάται. πόθεν;}{Le vers est incomplet. Scaliger proposa la solution la plus instinctive : redoubler le πόθεν en fin de vers : πόθεν ; <πόθεν;> Cette solution est vraisemblablement la bonne. \ \ \ 1° \ Une simple haplographie l'explique. \ \ \ 2° \ Elle est une nouvelle occurrence des railleries artiphanesques quant à la manie du redoublement chez les tragiques. Le double πόθεν se trouve d'ailleurs dans l'\italique{Ion} (vers 1453 sq.) :%
\citComm{ἰὼ γύναι, πόθεν πόθεν ἔλαβες ἐμὸν\\
βρέφος ἐς ἀγκάλας;}%
Même si ce qui faisait vraiment battre le cœur d'Euripide, c'était le redoublement du ποῦ (\italique{Électre}, v. 487 sq.) :%
\citComm{ποῦ ποῦ νεᾶνις πότνι᾽ ἐμὴ δέσποινά τε,\\
Ἀγαμέμνονος παῖς, ὅν ποτ᾽ ἐξέθρεψ᾽ ἐγώ;}%
%Et encore 
%Devant ces épouvantables redoublement de ποῦ, le sang d'Aristophane n'a fait qu'un tour. La caricature ne s'est pas fait attendre :%
%Pour se borner au \italique{Palamède}\marqueur{1}, on trouve le redoublement de l'adverbe πάλαι :%
%\citComm{πάλαι πάλαι δή σ᾿ ἐξερωτῆσαι θέλων, σχολή μ᾿ ἀπεῖργε.}%
Dans les \italique{Oiseaux} (vers 921), c'est le redoublement de πάλαι qui est moqué : l'ennuyeux poète qui débite%
\citComm{πάλαι, πάλαι δὴ τήνδ᾽ ἐγὼ κλῄζω πόλιν}%
ne pénétrera pas plus dans Νεφελοκοκκυγίᾳ qu'il n'aura de chance de demeurer dans la Cité de Platon.%
}%
%\debutNote{772}%
%\note{1}{Fr. 2 Jouan--van Loy $=$ 579 Kannicht $=$ 776 Mette.}
\\
τί δ᾽ ἂν εἰ ταδὶ τἀγάλματ᾽ ἀντὶ τῶν πλατῶν\\
γράφων διαρρίπτοιμι; βέλτιον πολύ.%
\commentaire{773-4}{τί δ᾽ ἂν εἶτα διὰ τἀγάλματ᾽ ἀντὶ τῶν πλατῶν / γράφων διαρρίπτοιμι;}{Willems (\nieme{22} note, reproduite en annexe) a expliqué en quel sens il fallait prendre \grec{τὰ ἀγάλματα} : \grec{πᾶν ἐφ᾽ ᾧ τις ἀγάλλεται} (scholie à notre vers), \italique{tout ce qu'on offre et consacre aux dieux}. En l'occurrence, ce sont les \grec{πόπανα} que Mnésiloque a apportés (au vers 285) et non pas des \italique{ex-voto} comme on en voit encore apparaître dans les modernes traductions\marqueur{1}.\\
Y a-t-il ici une allusion au proverbe grec ἄλλο γένος κώπης, «~un autre genre de rame~»\marqueur{2} ?}%
\debutNote{773-4}%
\note{1}{Après avoir encombré la scène d'une statue d'Apollon permanente d'après le vers 748, les commentateurs rajoutent ici des tablettes votives pour le décor.}%
\note{2}{Rapporté par Plutarque, \italique{Proverbes alexandrins}, 2, 52. Ou aussi dans \italique{Zu den Paroemiographen} de Cohn, p.42. Érasme l'a incorporé à ses \italique{Adages} (n°1621).}\\
ξύλον γέ τοι καὶ ταῦτα κἀκεῖν᾽ ἦν ξύλον.\\\passeLigne}%
\end{groupe}%
%% v. 776-784
\begin{groupe}%
%
%\personnageGrec{Μνησίλοχος}%
\typeVers{anap.}\anap{ὦ χεῖρες ἐμαὶ}%
\end{groupe}%
%% v. 776-784
\begin{groupe}%
\typeVers{dim. anap.}\anap{ἐγχειρεῖν χρῆν ἔργῳ πορίμῳ.\testimonia{777}{Suidas π 2064}\apparat{777}{χρή R, Suid. π 2064 : \lecon{χρῆν} Bentley}\\
ἄγε δὴ πινάκων ξεστῶν δέλτοι,\testimonia{778-80}{(πινάκων -- μόχθων) Pollux X 58 (ἄγε δὴ et δέλτοι omissis)}\testimonia{778-9}{Suidas σ 742}}%
\end{groupe}%
%% v. 776-784
\begin{groupe}%
\typeVers{(catalect.)}\anap{δέξασθε σμίλης ὁλκοὺς\\
κήρυκας ἐμῶν μόχθων· οἴμοι}%
\end{groupe}%
%% v. 776-784
\begin{groupe}%
\typeVers{(catalect.)}\anap{τουτὶ τορῶ μοχθηρόν·\apparat{781}{τὸ ῥῶ R~: \lecon{τορῶ} Bibi ap. Hesych. τ 1162}% \ligneApp{782} \lecon{χώρει} bis R : χωρεῖ bis Brunck}
\commentaire{781}{τουτὶ τὸ ῥῶ μοχθηρόν·}{« Voilà un `Rhô' défectueux ! »... telle est la leçon de notre unique manuscrit et telles aussi les traductions embarrassées. Car la question fuse : pourquoi Mnésiloque nous parle-t-il d'un `ρ' ? Le scholiaste toujours prompt à lancer une ânerie propose d'y voir la troisième lettre du destinataire, Εὐρῑπίδη... ce qui ne répond pas à la question : pourquoi la troisième lettre ? % pourquoi le Rhô ? il ne me semble pas plus difficile que les autres lettres à tracer. 
Passons sur les autres explications du même acabit (si Aristophane avait un peu plus de glosateurs, tous les mots du dictionnaire comportant un rhô y seraient effectivement passés : Εὐριπίδη, χώρει, χαίρειν, \&c.)
Lisez simplement : %
\citComm{τουτὶ τορῶ μοχθηρόν.}%
L'aoriste τορεῖν est homérique et fait partie d'un groupe dialectal, également homérique, d'aoristes radicaux thématiques où s'est dégagé un vocalisme en `o' : ἔμολον de βλώσκω, ἔθορον de θρώσκω et ἔτορον de τείρω. Ces homérismes se sont introduits chez les Tragiques, et peuvent aussi se rencontrer dans certains passages d'Aristophane. On l'avait déjà noté pour ἔμολον : on peut maintenant y ajouter ἔτορον\marqueur{1}. Hésychius (τ 1162) -- qui nous fut d'un grand secours comme à son habitude – l'avait déjà relevé :%
\citComm{τορεῖν· τορῆσαι, τρῆσαι, τεμεῖν ;}%
le verbe signifie \italique{percer}, \italique{perforer}, \italique{couper}.
Reconstituons la scène : Mnésiloque cherche un expédient pour envoyer son message. Heureusement, il connaît son euripide par cœur : il décide de procéder à la façon d'Œax dans le \italique{Palamède}. Manque de chance : il manque les rames (ἀλ ̓ οὐ πάρεισιν αἱ πλάται) pour graver dessus son message !.. qu'à cela ne tienne : les ustensiles de théâtres pour figurer les πόπανα feront bien l'affaire (ils sont de bois, après tout !) ; voilà donc nos morceaux de bois promus en tablettes (δέλτοι). Mnésiloque peut maintenant se mettre à graver la “cire” avec son stylet – un stylet ? quel stylet ? Nouveau détournement : il appelle `stylet' le coutelas avec lequel il vient d'égorger l'outre de vin. Seulement, à force de prendre des vessies pour des lanternes, Mnésiloque se fait rattraper par la dure sylvestre réalité : graver un message au couteau sur du bois se révèle bien plus difficile qu'avec un stylet sur une tablette de cire ; cela explique son cri de découragement : τουτὶ τορῶ μοχθηρόν, « ceci que je perce [il ne précise pas quoi, et pour cause : c'est à la fois un πόπανον, une πλάτη et un δέλτος]... est de mauvaise qualité ! »}%
\debutNote{781}%
\note{1}{Je dois cette remarque à Jean-Victor Vernhes qui, entre autres qualités, est un très fin connaisseur de la langue homérique. Que le lecteur n'hésite pas à consulter le Dictionnaire étymologique de la langue grecque de Pierre Chantraine à l'entrée τορεῖν pour plus de détails.}\\
`χώρει', `χώρει'. ποίαν αὔλακα;\\%\apparat{782}{\lecon{χώρει} bis R : χωρεῖ bis Brunck}\\%
\commentaire{782}{χώρει, χώρει.}{%Le doublement du verbe χωρεῖν se trouve aussi au vers 1076 de \italique{Médée}, grand moment d'euripidéenne psychologie où l'«~héroïne~», prise d'un sursaut d'instinct maternel, s'interroge : est-il bien sain de massacrer ses propres enfants par vengeance ? Ne vaudrait-il pas mieux les laisser s'échapper : χωρεῖτε χωρεῖτ᾽, «~fuyez, fuyez !~». Ah ! ces fameux dilemmes euripidéens !.. \anglais{to kill ? or not to kill ?..}
%Mais prenons garde à ne pas fonctionner par système ; ici, l'explication du double χωρεῖν est tout autre : 
Mnésiloque est en train de graver son message : χώρει, «~accours !~», enjoint-il à Euripide ; χώ...ρει χώ...ρει, «~Ac...cours ... Ac...cours~», se dicte-t-il à lui-même. Il l'écrit sur deux rames différentes, comme le pluriel αἱ πλάται le laissait présager et comme c'était aussi le cas dans la tragédie originale ; le scholiaste de notre \italique{Ravennas} précise bien : %
\citComm{ὁ γὰρ Οἴαξ ἐγχαράττει \textbf{πολλαῖς πλάταις} τὰ περὶ τὸν Παλαμήδην καὶ ἀφίησιν εἰς θάλασσαν, ὥστε μιᾷ γέ τινι τὸν Ναύπλιον περιπεσεῖν.\\
Œax grave effectivement de \italique{nombreuses rames} avec les circonstances touchant à $[$la mort de$]$ Palamède et les lance vers la mer de sorte que Nauplios tombe sur l'une d'elles.}%
À propos des ἀγάλματα, un scholiaste avait aussi écrit :%
\citComm{πίνακες γὰρ ἦσαν ἐν τῷ ἱερῷ, καὶ ἀπὸ τούτων λαβὼν γράφει, καὶ ῥίπτει αὐτὰς λέγων· `ἀπέλθετε, σημάνατε Ἐὐριπίδῄ.\\
C'était des tablettes dans le temple. Les prenant, il écrivit puis les lança, en disant : «~Allez ! Informez Euripide !~»}%
}%
\commentaire{782-4}{ποίαν αὔλακα / βάσκετ᾽ ἐπείγετε πάσας καθόδους / κείνᾳ ταύτᾳ· ταχέως χρή.}{Les différents traducteurs voient dans la formule `ποίαν αὔλακα' une exclamation d'auto-contentement de Mnésiloque : «~Quelle belle lettre ! Quel sillon !~» Αὖλαξ est en effet attesté dans l'\italique{Anthologie palatine} (6, 68) au sens de sillon creusé par le stylet, c'est-à-dire le \italique{trait d'écriture}. Soit. Mais cette interprétation confère à ποίαν un sens qu'il n'a pas. Cet adjectif interrogatif sert à demander des précisions : \italique{de quelle nature ?} \italique{de quel genre ?} (Aristophane l'emploie fréquemment dans ses dialogues pour débrouiller -- ou souligner ? -- un quiproquo.) Il est utilisé en reprise d'un terme, reprise directe : Ἀγάθων ὁ κλεινὸς --- Ποῖος οὗτος Ἁγάθων; (vers 30) ; Τὸν δεῖνα γιγνώσκεις --- Τὸν δεῖνα; ποῖον; (vers 621) ;  Πρωτέως τάδ᾽ ἐστὶ μέλαθρα --- Ποίου Πρωτέως; (vers 874) ou implicite : Πρὸς ποῖον ἔργον; (vers 586) ; Χρείᾳ δὲ ποίᾳ τόνδ᾽ ἐπεισφέρεις λόγον; (vers 1163). Mnésiloque mentionne pour la première fois l'αὖλαξ ; mais reprise il y a bien\marqueur{1} : le mot devait sans aucun doute apparaître dans la tragédie d'Euripide. 
%Réfléchissons ! 
Où en est Mnésiloque de l'action ? Il vient de graver en deux exemplaires l'appel au secours, χώρει, qu'il compte faire parvenir à Euripide... comment ? %telle est la question ! 
Assurément, cette αὖλαξ désigne l'expédient qu'Œnax avait trouvé pour emporter ses messages.
On peut raisonnablement penser qu'il les déposait dans un \italique{cours d'eau} (qui envoyait ensuite les πλάται où l'une d'entre elles finissait par arriver à bon port, comme l'indique le scholiaste. Ça ne paraît pas très vraisemblable, mais la vraisemblance et une pièce d'Euripide font deux.)
C'est bien en ce sens qu'il faut prendre αὖλαξ, le \italique{sillon}, et, par extension, la \italique{rigole} ou le \italique{ruisseau}\marqueur{2}.\\
Mnésiloque ne s'embarrasse pas vraiment du lièvre qu'il vient de soulever : la question n'est là que pour souligner les invraisemblances de la mise en scène, comme Aristophane aime bien les relever. Il décide d'aller au plus rapide : il jette avec force les deux πλάται. (Une de chaque côté, vers chaque εἴσοδος, comme le laisse soupçonner πάσας καθ᾽ ὁδοὺς.) Cette \italique{nouveauté} (κείνᾳ ταύτᾳ) par rapport au modèle original d'Euripide a l'avantage d'accélérer l'action ; c'est ainsi qu'il faut comprendre le vers 784 : κείνᾳ ταύτᾳ· ταχέως χρή.%
}%
\debutNote{782-4}%
\note{1}{Et si l'on en doutait, l'on pourrait toujours corriger le texte -- ce dont on se gardera bien ! Lire ποῖ ποτ' ἂν αὔλακα revient au même quant à l'action et dispense de supposer que l'αὖλαξ était le terme employé par Euripide. (ποτε \italique{gratia metri} mais dont la disparition relèverait d'une haplographie bien explicable. Au reste, la collocation `ποῖ ποτε' est très fréquente.)}%
\note{2}{Voir le \LSJ \ et le \espagnol{Diccionario Griego-Español} IV à l'entrée αὖλαξ.}\\
βάσκετ᾽ ἐπείγετε πάσας καθ᾽ ὁδοὺς\apparat{783}{καθόδους R : \lecon{καθ᾽ ὁδοὺς} Biset}}\\%
\end{groupe}%
%% v. 776-784
\begin{groupe}%
%
\typeVers{(catalect.)}\anap{κείνᾳ ταύτᾳ· ταχέως χρή.\apparat{784}{ταῦτα R : ταύτᾳ Grynaeus (\italique{hac} Divus)}%
\\\passeLigne}%
\end{groupe}%
%% v. 785-843
\begin{groupe}%
%
\personnageGrec{Χορός}%
\typeVers{\decalageN{a}$4$ an\cat}\tetramAnap{\ligne{785}ἡμεῖς τοίνυν ἡμᾶς αὐτὰς εὖ λέξωμεν παραβᾶσαι,
\commentaire{785-845}{Parabase.}{La parabase des \italique{Thesmophories}, à l'instar de \italique{Lysistrata} et des \italique{Grenouilles}, ne présente ni ᾠδή/ἀντῳδή ni ἀντεπίρρημα\marqueur{1}.}%
\debutNote{785-845}%
\note{1}{Mon but n'est pas de proposer une description exhaustive de la structure des comédies d'Aristophane, ni même des \italique{Thesmophories} en particulier. Ce travail a déjà été très bien fait par White, \italique{The Verse of Greek Comedy} (1912). On lira aussi avec un grand intérêt les travaux de Jean Irigoin regroupés dans son anthologie posthume \italique{Le poète grec au travail} (2009) et notamment son article «~Prologue et \italique{parodos} dans la comédie d'Aristophane (\italique{Acharniens}, \italique{Paix}, \italique{Thesmophories})~» :  «~[Aristophane], ce poète comique à l'imagination débridée, s'astreignait, dans certaines parties de ses pièces, à respecter des règles concernant le choix des vers, leur disposition, et même leur nombre~». Le lecteur intéressé se reportera à ces travaux -- et à nombre d'autres -- auxquels je ne saurais me substituer.}\\
καίτοι πᾶς τις τὸ γυναικεῖον φῦλον κακὰ πόλλ᾽ ἀγορεύει,\\
ὡς πᾶν ἐσμὲν κακὸν ἀνθρώποις κἀξ ἡμῶν ἐστιν ἅπαντα,\\
\Automatique{ἔριδες νείκη στάσις ἀργαλέα λύπη πόλεμος. φέρε δή νυν,}\apparat{788}{στάσεισ· ἀργαλέα· R : \lecon{στάσις ἀργαλέα} Scaliger \espApp εἰ καὶ κακόν R : \lecon{εἰ κακόν} Zanetti \ligneApp{790} ἐγκύψασαν R : \lecon{ἐκκύψασαν} Reiske}%
\commentaire{788}{ἔριδες νείκη στάσις ἀργαλέα}{Le vocabulaire est Homérique( \italique{Iliade}, XVII, v. 384-5) :%
\citComm{τοῖς δὲ πανημερίοις ἔριδος μέγα νεῖκος ὀρώρει / ἀργαλέης·\\
\francais{Durant tout le jour se prolongea cette lutte sanglante et terrible.}}%
}\\
\Automatique{εἰ κακόν ἐσμεν, τί γαμεῖθ᾽ ἡμᾶς, εἴπερ ἀληθῶς κακόν ἐσμεν,}%\apparat{789}{εἰ καὶ κακόν R : \lecon{εἰ κακόν} Zanetti}
\commentaire{789}{εἰ καὶ κακόν ἐσμεν, τί γαμεῖθ᾽ ἡμᾶς, εἴπερ ἀληθῶς κακόν ἐσμεν}{La perfidie d'Aristophane est à son comble quand il donne la parole aux femmes dans la parabase où elles assurent leur propre défense. « Si nous sommes un fléau, disent-elles aux hommes, pourquoi nous épousez-vous ? » Mais le spectateur, qui connaît, depuis l'école, son Hésiode par cœur et qui a déjà vu bien des comédies, répond intérieurement, comme mu par un ressort : « Parce qu'on n'échappe pas si facilement au piège tendu par Zeus, parce que la femme est un fléau nécessaire pour la perpétuation de l'espèce\marqueur{1} ». Le poète comique ne fait ici, sous des apparences trompeuses, qu'entretenir le vieux réflexe de misogynie. Autre argument employé par les femmes pour montrer leur valeur : celui de la jalousie qu'elles excitent\marqueur{2}. Tout commentaire est ici inutile... (Note d'Alain Blanchard\marqueur{3}.)}
\debutNote{789}
\note{1}{Hésiode, les \italique{Travaux et les jours}, v. 83-84 ; Susarion, fr. 1, 3-4 Kassel-Austin. Voir aussi ce que dit Renzo Tosi (§284 et §285 de son monumental \italique{Dictionnaire des sentences latines et grecques}, 2010) : Κακὸν ἀναγκαῖον γυνή ou \latin{Femina dulce malum} : la femme est un καλὸν κακόν (\italique{Théogonie}, v. 587). Quant au vers précédent, il y une allusion à peine voilée aux pandoriennes calamités déclenchées par la mythique -- et belle, surtout, pour nos malheurs ! -- Hélène. (Voir à ce propos Tosi, §350 : \latin{Nam fuit ante Helenam cunnus taeterrima belli causa}.) Il est donc tout naturel que Mnésiloqie réponde à cette parabase par une reprise de l'\italique{Hélène}.}
\note{2}{Aux vers 792-794.}
\note{3}{Tirée de son remarquable article «~L'art de la mauvaise foi dans les \italique{Thesmophories} d'Aristophane~», in \italique{Dieux, héros et médecins grecs. Hommage à Fernand Robert.} (2001), p. 97-106. Tout amateur d'Aristophane ne peut pas ne pas aller prestement lire ces pages si riches (j'émets seulement quelques réserves sur le tout dernier paragraphe).
Blanchard y montre que « la mauvaise foi d'Aristophane s'est exercée, en ce qui concerne Euripide, sur un thème moins traditionnel que celui de la misogynie et qu'elle permet de mieux comprendre le sujet même de sa comédie des \italique{Thesmophories}. » De même que la défaite athénienne à Délion avait profité à la popularité de Socrate, qui s'était conduit héroïquement dans la débâcle comme on le sait par Platon, et que c'est ce bonheur d'une réputation rehaussée dans le malheur athénien qui avait provoqué Aristophane à l'écriture vengeresse des \italique{Nuées} ; de même, les \italique{Thesmophories} furent écrites peu après la désastreuse expédition de Sicile : alors qu'Athènes vivait ses heures les plus sombres, Euripide en profita pour redorer son blason : on fit appel à lui pour écrire un chant funèbre en l'honneur des Athéniens morts à Syracuse (Plutarque nous en a conservé deux vers). 
Anecdote plus significative encore : certains Athéniens, réduits en esclavage, ne durent leur salut qu'à Euripide, en récitant de ses prologues, si prisés des Syracusains. Que la poésie d'Euripide pût sauver qui que ce soit ?.. alors qu'elle est justement l'une des causes de la dégénérescence athénienne !? voilà une nouvelle fumisterie qu'Aristophane mettra à mal tout au long des \italique{Thesmophories} : les différents détournements des tragédies d'Euripide se heurteront toutes au mur de l'incompréhension des geôliers de Mnésiloque. Une fois bien cerné le contexte historique, l'on comprend mieux Aristophane : ce « jeu de massacre » n'est motivé que par la passion -- brûlante chez Aristophane -- de la justice : il cogne, cogne, et cogne encore contre tous ces charognards des misères athéniennes. \anglais{Now is a time to storm !} Pour conclure : si Euripide est misogyne, c'est là son moindre défaut.}\\
\Automatique{κἀπαγορεύετε μήτ᾽ ἐξελθεῖν μήτ᾽ ἐκκύψασαν ἁλῶναι,}\\%\apparat{790}{ἐγκύψασαν R : \lecon{ἐκκύψασαν} Reiske}\\
\Automatique{ἀλλ᾽ οὑτωσὶ πολλῇ σπουδῇ τὸ κακὸν βούλεσθε φυλάττειν;}\\
\Automatique{κἂν ἐξέλθῃ τὸ γύναιόν ποι, κᾆθ᾽ εὕρητ᾽ αὐτὸ θύρασιν,}\apparat{792}{κᾆθ᾽ εὕροιτ' R : \lecon{κᾆθ᾽ εὕρητ᾽} Porson \ligneApp{793} μαίνεθ᾽ R : \lecon{μαίνεσθ᾽} \Mon \espApp χρή R : \lecon{χρῆν} Brunck \ligneApp{794} εὕρετε R : \lecon{ηὕρετε} Meineke \espApp καταλαμβάνετ᾽ R : \lecon{κατελαμβάνετ᾽} Brunck}\\
\Automatique{μανίας μαίνεσθ᾽, οὓς χρῆν σπένδειν καὶ χαίρειν, εἴπερ ἀληθῶς}\\%\apparat{793}{μαίνεθ᾽ R : \lecon{μαίνεσθ᾽} \Mon \espApp χρή R : \lecon{χρῆν} Brunck \ligneApp{794} εὕρετε R : \lecon{ηὕρετε} Meineke \espApp καταλαμβάνετ᾽ R : \lecon{κατελαμβάνετ᾽} Brunck}\\
\Automatique{ἔνδοθεν ηὕρετε φροῦδον τὸ κακὸν καὶ μὴ κατελαμβάνετ᾽ ἔνδον.}\\
\Automatique{κἂν καταδάρθωμεν ἐν ἀλλοτρίων παίζουσαι καὶ κοπιῶσαι,}\\%\apparat{795}{καταδαρθῶμεν R : \lecon{καταδάρθωμεν} Dindorf}\\
\Automatique{πᾶς τις τὸ κακὸν τοῦτο ζητεῖ περὶ τὰς κλίνας περινοστῶν.}}
\end{groupe}%
%% v. 785-843 (part 1)
\begin{groupe}\tetramAnap{\Automatique{κἂν ἐκ θυρίδος παρακύπτωμεν, τὸ κακὸν ζητεῖτε θεᾶσθαι·}\apparat{797}{ζητεῖ τὸ κακὸν τεθεᾶσθαι R : \lecon{τὸ κακὸν ζητεῖτε θεᾶσθαι} J. Kaye 
\ligneApp{799} παρακύψασαν ἰδεῖν τὸ κακὸν R : \lecon{τὸ κακὸν παρακύψαν ἰδεῖν} Dindorf}
\commentaire{797}{τὸ κακὸν ζητεῖτε θεᾶσθαι}{Suivant la belle correction de Kaye : τὸ κακὸν ζητεῖτε θεᾶσθαι, v. 797, au lieu de : ζητεῖ τὸ κακὸν τε θεᾶσθαι. (Note de Willems.)}\\
\Automatique{κἂν αἰσχυνθεῖσ᾽ ἀναχωρήσῃ, πολὺ μᾶλλον πᾶς ἐπιθυμεῖ}\\
\Automatique{αὖθις τὸ κακὸν παρακύψαν ἰδεῖν. οὕτως ἡμεῖς ἐπιδήλως}\\\passeLigne}%
\end{groupe}%
\endnumbering
\end{Leftside}
\begin{Rightside}
\beginnumbering
%% v. 765-775
\begin{groupe}%
\personnageFr{Mnésiloque}%
\trimIambFr{\ligne{765}Voyons, quel sera l'instrument de mon salut ?\\
Quelle tentative? Quelle combinaison ? Car l'auteur vrai,
celui qui m'impliqua dans cette méchante affaire,
n'apparaît pas encore. Allons, quel messager
lui dépêcherai-je ? (Il réfléchit)... Il me vient à l'esprit un expédient
tiré du \italique{Palamède}. Comme cet autre, je lancerai mes avirons
chargés d'écritures... Mais je ne les ai pas sous la main, ces avirons...
Comment donc faire pour me procurer des avirons ? Comment ?...
Mais quoi, si à défaut d'avirons, j'écrivais sur ces offrandes
et les jetais de tous côtés ? Cela vaut bien mieux.
Elles sont de bois après tout; ceux-là aussi étaient de bois.\\\passeLigne}%
\end{groupe}%
%% v. 776-784
\begin{groupe}%
\anapFr{À l'œuvre, ô mes mains,
}%
\end{groupe}%
%% v. 776-784
\begin{groupe}%
\anapFr{il s'agit d'un travail expédient.\\
Allons, feuillets de mes tablettes polies,}%
\end{groupe}%
%% v. 776-784
\begin{groupe}%
\anapFr{\ligne{779}recevez les entailles du ciseau,\\
messagères de mes misères. Aïe,}%
\end{groupe}%
%% v. 776-784
\begin{groupe}%
\anapFr{c'est pénible à percer !..
«~Ac...cours~», «~Ac...cours !~» et quel sera notre cours d'eau ?
(Il se tourne et lance les tablettes des deux côtés.)
\ligne{782}Partez, volez par toutes les voies}%
\end{groupe}%
%% v. 776-784
\begin{groupe}%
%
\anapFr{\ligne{784}avec cette nouvelle façon de faire : il faut \fin aller vite.\\\passeLigne}%
\end{groupe}%
%% v. 785-843
\begin{groupe}%
\personnageFr{Chœur}%
\tetramAnapFr{Maintenant nous allons faire nous-mêmes notre éloge dans la parabase.
Chacun pourtant médit à l'envi de l'espèce féminine :
nous sommes un pur fléau pour les hommes ; de nous viennent tous les maux,
disputes, querelles, affreuses séditions, chagrins, guerres. Çà, voyons,
si nous sommes un fléau, pourquoi nous épousez-vous, ce fléau que nous sommes ?
Pourquoi nous interdire de sortir, de nous laisser prendre à pousser la tête dehors ?
Pourquoi vous piquer de le garder avec un soin si jaloux, ce fléau ?
\ligne{792}Arrive-t-il à votre petite femme d'être en course, et que vous la trouviez absente,
vous entrez dans des fureurs, quand il vous faudrait faire libation et vous réjouir
de ce qu'il a disparu de chez vous, le fléau, et que vous ne le trouviez pas au logis.
Nous endormons-nous chez les autres, fatiguées de danser,
chacun se met après le fléau et rode autour \fin des lits.}
\end{groupe}%
%% v. 785-843 (part 1)
\begin{groupe}\tetramAnapFr{\ligne{796}Nous penchons-nous par une fenêtre, vous cherchez à le dévisager, le fléau,
et si par pudeur elle se retire, chacun n'en est que plus impatient
\ligne{797}de la voir se pencher de nouveau, ce fléau. Tant il est vrai que nous valons bien mieux\\\passeLigne}%
\end{groupe}%
\endnumbering
\end{Rightside}
\end{pages}
\Pages
\clearpage\ifodd\value{page}\hbox{}\newpage\fi
\enTeteG{ΠΑΡΑΒΑΣΙΣ --- Codex Ravennas 429, f. 172r (Classense 174r) --- vv. 800–837}
\begin{adjustwidth}{-1.5cm}{-3.0cm}
\begin{adjustbox}{max totalheight=21.0cm, max width=16.5cm, keepaspectratio}
\begin{minipage}{16.5cm}
{\fontsize{7.2pt}{8.5pt}\selectfont
\noindent\textbf{Folio 20 --- Codex Ravennas 429, \texttt{f. 172r} (\texttt{Classense 174r}) : ΠΑΡΑΒΑΣΙΣ, vv. 800–837}\hfill \textit{Témoins :} $\text{R}^1, \Rpc, \Rdeux$, \PsiUn\par\vspace{2pt}
\hrule height 0.4pt\vspace{2.5pt}
}
{\fontsize{8.1pt}{9.5pt}\selectfont
\ligneRav{800}{χορ/}{ὑμῶν ἐσμεν πολὺ βελτίους, βάσανός τε πάρεστιν ἰδέσθαι·}{}
\ligneRav{}{}{βάσανον δῶμεν πότεροι χείρους. ἡμεῖς μὲν γάρ φαμεν ὑμᾶς,}{}
\ligneRav{802}{}{\difference{ὑμεῖς δ᾽ ἡμᾶς}. \difference{σκεψώμεθα} δὴ κἀντιτιθῶμεν πρὸς ἕκαστον,}{\textbf{802} ὑμεῖς δ᾽ ὑμᾶς R : \lecon{ὑμεῖς δ᾽ ἡμᾶς} \Mon \espApp σκεψώμεσθα R : \lecon{σκεψώμεθα} \Mon}
\ligneRav{803}{}{παραβάλλουσαι τῆς τε γυναικὸς καὶ τἀνδρὸς τοὔνομ᾽ \difference{ἑκάστου}·}{\textbf{803} ἕκαστον R : \lecon{ἑκάστου} Grynaeus}
\ligneRav{804}{}{Ναυσιμάχης \difference{μέν γ᾽} ἥττων \difference{ἐστὶν} Χαρμῖνος· \difference{δῆλα δὲ} τἄργα·}{\textbf{804} μέν R : \lecon{μέν γ᾽} Dobree \espApp ἐστὶ R : \lecon{ἐστὶν} Küster \espApp δηλαδῆ R : \lecon{δῆλα δὲ} Grynaeus}
\ligneRav{805}{}{καὶ μὲν δὴ καὶ Κλεοφῶν \difference{χείρων} πάντως δήπου Σαλαβακχοῦς·}{\textbf{805} χεῖρον R : \lecon{χείρων} Zanetti}
\ligneRav{}{}{πρὸς Ἀριστομάχην δὲ χρόνου πολλοῦ, πρὸς ἐκείνην τὴν Μαραθῶνι,}{}
\ligneRav{807}{}{καὶ \difference{Στρατονίκην} ὑμῶν οὐδεὶς οὐδ᾽ ἐγχειρεῖ πολεμίζειν·}{\textbf{807} \lecon{Στρατονίκην} \PsiUn : στρατωνίκην R}
\ligneRav{808}{}{ἀλλ᾽ Εὐβούλης τῶν πέρυσίν \difference{τις βουλευτής} ἐστιν ἀμείνων}{\textbf{808} τίς βουλευτής R : \lecon{τις βουλευτής} Küster}
\ligneRav{809}{}{παραδοὺς ἑτέρῳ τὴν βουλείαν; οὐδ᾽ αὐτὸς τοῦτό γε \difference{φήσει}·}{\textbf{809} \lecon{αὐτὸς} R : Ἄνυτος Maas \espApp φήσεις R : \lecon{φήσει} Küster}
\ligneRav{810}{}{οὕτως ἡμεῖς πολὺ βελτίους τῶν ἀνδρῶν \difference{εὐχόμεθ᾽} εἶναι·}{\textbf{810} εὐχόμεσθ᾽ R : \lecon{εὐχόμεθ᾽} \Mon}
\ligneRav{811}{}{οὐδ᾽ ἂν κλέψασα γυνὴ ζεύγει κατὰ πεντήκοντα τάλαντα}{\textbf{811} ζεύγει post γυνὴ R : post οὐδ᾽ ἂν van Herwerden}
\ligneRav{}{}{εἰς πόλιν ἔλθοι τῶν δημοσίων· ἀλλ᾽ ἢν τὰ μέγισθ᾽ ὑφέληται}{}
\ligneRav{813}{}{φορμὸν πυρῶν τἀνδρὸς κλέψασ᾽, αὐθημερὸν αὔτ᾽ ἀπέδωκεν}{\textbf{813} \lecon{αὔτ᾽ ἀπέδωκεν} R, M\expo{pc} : ἀνταπέδωκεν M\expo{ac}, Bentley}
\ligneRav{}{—}{\hspace{1.1em}ἀλλ᾽ ἡμεῖς ἂν πολλοὺς τούτων}{}
\ligneRav{815}{}{\hspace{1.1em}\difference{ἀποδείξαιμεν} ταῦτα \difference{ποιοῦντας}·}{\textbf{815} ἀποδείξομεν R : \lecon{ἀποδείξαιμεν} Bentley \espApp ποθοῦντας R : \lecon{ποιοῦντας} Zanetti (\italique{facientes} Divus)}
\ligneRav{}{}{\hspace{1.1em}καὶ πρὸς τούτοις γάστριδας ἡμῶν}{}
\ligneRav{}{}{\hspace{1.1em}ὄντας μᾶλλον καὶ λωποδύτας}{}
\ligneRav{}{}{\hspace{1.1em}καὶ βωμολόχους κἀνδραποδιστάς·}{}
\ligneRav{819}{}{\hspace{1.1em}καὶ μὲν δήπου καὶ τὰ πατρῷά γε}{\textbf{819} δήπου R : δή που Denniston}
\ligneRav{820}{}{\hspace{1.1em}χείρους ἡμῶν εἰσιν σῴζειν·}{}
\ligneRav{}{}{\hspace{1.1em}ἡμῖν μὲν γὰρ σῶν ἔτι καὶ νῦν}{}
\ligneRav{822}{}{\hspace{1.1em}\difference{ταὐτὸ κανοῦν καὶ τὸ σκιάδειον}·}{\textbf{822-3} τἀντίον ὁ κανὼν οἱ καλαθίσκοι / τὸ σκιάδειον· R : \lecon{ταὐτὸ κανοῦν καὶ τὸ σκιάδειον·} Bibi}
\ligneRav{824}{}{\hspace{1.1em}τοῖς δ᾽ ἡμετέροις \difference{ἀνδράσι} τούτοις}{\textbf{824} ἀνδράσιν R : \lecon{ἀνδράσι} \Mon}
\ligneRav{825}{}{\hspace{1.1em}\difference{ἀπόλωλεν} μὲν πολλοῖς ὁ κανὼν}{\textbf{825} ἀπόλωλε R : \lecon{ἀπόλωλεν} Daubuz, Bentley}
\ligneRav{}{}{\hspace{1.1em}ἐκ τῶν οἴκων αὐτῇ λόγχῃ,}{}
\ligneRav{}{}{\hspace{1.1em}πολλοῖς δ᾽ ἑτέροις ἀπὸ τῶν ὤμων}{}
\ligneRav{}{}{\hspace{1.1em}\{ ἐν ταῖς στρατιαῖς \}}{}
\ligneRav{}{}{\hspace{1.1em}ἔρριπται τὸ σκιάδειον·}{}
\ligneRav{830}{χορ/}{πόλλ᾽ ἂν αἱ γυναῖκες ἡμεῖς ἐν δίκῃ μεμψαίμεθ᾽ ἂν}{}
\ligneRav{}{}{τοῖσιν ἀνδράσιν δικαίως, ἓν δ᾽ ὑπερφυέστατον·}{}
\ligneRav{832}{}{\difference{χρῆν} γάρ, ἡμῶν εἰ τέκοι τις ἄνδρα χρηστὸν τῇ πόλει,}{\textbf{832} χρὴ R : \lecon{χρῆν} Brunck}
\ligneRav{}{}{ταξίαρχον ἢ στρατηγόν, λαμβάνειν τιμήν τινα,}{}
\ligneRav{834}{}{προεδρίαν τ᾽ αὐτῇ δίδοσθαι \difference{Στηνίοισι} καὶ Σκίροις}{\textbf{834} Τηνίοισι R : \lecon{Στηνίοισι} Σ, Suid. σ 1089}
\ligneRav{835}{}{ἔν τε ταῖς ἄλλαις ἑορταῖς αἷσιν ἡμεῖς ἤγομεν·}{}
\ligneRav{836}{}{\difference{εἰ δὲ} δειλὸν καὶ πονηρὸν ἄνδρα τις τέκοι γυνή,}{\textbf{836} ἤν δὲ R : \lecon{εἰ δὲ} Brunck}
\ligneRav{}{}{ἢ τριήραρχον πονηρὸν ἢ κυβερνήτην κακόν,}{}
}
\vspace{3pt}\noindent\rule{4.5cm}{0.4pt}\par\vspace{1.5pt}
{\fontsize{7.0pt}{8.3pt}\selectfont
\noindent\textbf{Loci paralleli (vv. 800–837) ---} \textbf{800, 810} (βελτίους $=$ βελτίονες) Moeris β 8 \espTest \textbf{805} (Σαλαβακχοῦς) Hésychius (σ 88) Σαλαβακχώ· πόρνης ὄνομα. ἀπὸ δὲ ταύτης καὶ τὰς καταφερεῖς εἰς τὰ ἀφροδίσια οὕτως ἔλεγον Ἀττικοί \espTest \textbf{808-9} (τις βουλευτής -- βουλείαν) Suidas β 431 (ἐστι pro ἐστιν et ἀμείνων omisso) \espTest \textbf{809} (παραδοὺς ἑτέρῳ τὴν βουλείαν) \Antiatt 84.30 \espTest (βουλείαν) \Harp B 19 ; β Photius β 233 βουλία ; Hesychius β 921 βουλεῖα \espTest \textbf{816} (γάστριδας) Photius γ 40 γάστρις $=$ Suidas γ 72 ; cf. \AelDion γ 3 \espTest \textbf{822} %(τἀντίον ὁ κανὼν) Pollux VII 35 κανὼν ἱστοῦ τὸ καλούμενον ἀντίον \espTest \textbf{822} (τἀντίον ὁ κανὼν) Pollux VII 35 κανὼν ἱστοῦ τὸ καλούμενον ἀντίον \espTest \textbf{834} (Στηνίοισι) Suidas σ 1089 Στήνια καὶ Σκίρα ($=$ Scholium) ; Hesychius σ 1825 στήνια\par
}
\vspace{2.5pt}\noindent\rule{4.5cm}{0.4pt}\par\vspace{1.5pt}
{\fontsize{6.9pt}{8.2pt}\selectfont
\noindent\textbf{Suite (en commentaire) de l'apparat critique \& notes paléographiques (\texttt{f. 172r}) ---} \textbf{Témoin papyrologique :} \textbf{\PsiUn{} (\italique{PSI} XI 1194)} couvre les \textbf{vv. 804–809} et confirme Στρατονίκην au v. 807. --- \textbf{v. 815 (ποθοῦντας R pour πο(ι)οῦντας) :} L. Hauseux consacre un important excursus paléographique à cette faute d'onciale (ΠΟΟΥΝΤΑϹ $\to$ ΠΟΘΟΥΝΤΑϹ, confusion des lettres rondes Ο/Θ sur la forme attique sans iota ποοῦντας). --- \textbf{vv. 822–823 (restitution de \textbf{Bibi}) :} τἀντίον ὁ κανὼν οἱ καλαθίσκοι / τὸ σκιάδειον· R $\to$ \lecon{ταὐτὸ κανοῦν καὶ τὸ σκιάδειον·} \lecon{Bibi} (οἱ καλαθίσκοι est une glose marginale intruse expliquant κανοῦν, corrompu en κανὼν). --- \textbf{v. 834 :} Τηνίοισι R (chute du $\Sigma$ initial par haplographie après δίδοσθαι) $\to$ \lecon{Στηνίοισι} $\Sigma^{\text{R}}$, Souda $\sigma\ 1089$ (fête des \italique{Stènia}). --- \textbf{vv. 822–823 :} τἀντίον ὁ κανὼν οἱ καλαθίσκοι / τὸ σκιάδειον· R : \lecon{ταὐτὸ κανοῦν καὶ τὸ σκιάδειον·} \lecon{Bibi}. --- \textbf{v. 834 :} Τηνίοισι R : \lecon{Στηνίοισι} $\Sigma$, Souda $\sigma\ 1089$.\par
}
\end{minipage}
\end{adjustbox}
\end{adjustwidth}
\newpage
\enTeteD{Codex Ravennas 429 (Biblioteca Classense) --- Fac-similé du f. 172r (vv. 800–837)}
\begin{adjustwidth}{-1.5cm}{-3.0cm}
\begin{center}
\includegraphics[width=16.5cm,height=20.5cm,keepaspectratio]{/workspaces/Thesmophories/dossier_ravennas_thesmophories/folios_individuels_ff162v-177r/20_f172r_classense174r_vv0800-0837.jpg}\par
\vspace{2pt}
{\scriptsize\textbf{Codex Ravennas 429, f. 172r (Classense 174r)} --- ΠΑΡΑΒΑΣΙΣ, \textbf{vv. 800–837} ($\text{R}^1, \Rpc, \Rdeux$, \PsiUn).}
\end{center}
\end{adjustwidth}
\newpage
\clearpage\ifodd\value{page}\hbox{}\newpage\fi
\enTeteG{ΘΕΣΜΟΦΟΡΙΑΖΟΥΣΑΙ --- Texte grec établi (f. 172r, vv. 800–837)}
\enTeteD{LES THESMOPHORIES --- Traduction française (f. 172r, vv. 800–837)}
\begin{pages}
\begin{Leftside}
\beginnumbering
%% v. 800+ (part 2)
\begin{groupe}
\tetramAnap{\ligne{800}\Automatique{ὑμῶν ἐσμεν πολὺ βελτίους, βάσανός τε πάρεστιν ἰδέσθαι.}\testimonia{800, 810}{(βελτίους $=$ βελτίονες) Moeris β 8}\\
\Automatique{βάσανον δῶμεν πότεροι χείρους. ἡμεῖς μὲν γάρ φαμεν ὑμᾶς,}%
\commentaire{801-2}{ἡμεῖς μὲν γάρ φαμεν ὑμᾶς, ὑμεῖς δ᾽ ἡμᾶς}{Molière, \italique{Impromptu de Versailles}.%
\citComm{\francais{Il gage que c'est moi, et moi je gage que c'est lui.}}%
}}
\end{groupe}%
%% v. 785-843
\begin{groupe}%
%
\tetramAnap{\Automatique{ὑμεῖς δ᾽ ἡμᾶς. σκεψώμεθα δὴ κἀντιτιθῶμεν πρὸς ἕκαστον,}\apparat{802}{ὑμεῖς δ᾽ ὑμᾶς R : \lecon{ὑμεῖς δ᾽ ἡμᾶς} \Mon \espApp σκεψώμεσθα R : \lecon{σκεψώμεθα} \Mon \ligneApp{803} ἕκαστον R : \lecon{ἑκάστου} Grynaeus \ligneApp{804} μέν R : \lecon{μέν γ᾽} Dobree \espApp ἐστὶ R : \lecon{ἐστὶν} Küster \espApp δηλαδῆ R : \lecon{δῆλα δὲ} Grynaeus \ligneApp{805} χεῖρον R : \lecon{χείρων} Zanetti}\\
\Automatique{παραβάλλουσαι τῆς τε γυναικὸς καὶ τἀνδρὸς τοὔνομ᾽ ἑκάστου.}\\
\Automatique{Ναυσιμάχης μέν γ᾽ ἥττων ἐστὶν Χαρμῖνος· δῆλα δὲ τἄργα.}%
\commentaire{804}{Ναυσιμάχης μέν ἥττων ἐστὶν Χαρμῖνος·}{Charminos, stratège athénien, qui venait de se faire battre sur mer près de Symé\marqueur{1}, est rapproché de Nausimaché (\italique{bataille navale}). En quoi Charminos est-il inférieur (ἥττων) à Nausimaché ? Certes par le nom -- comme Aristophane n'avait pas manqué de prévenir : après les déboires qu'il fit subir à la marine athénienne, il serait bien indigne d'en porter un de ce genre !.. mais en rester à cet unique rapprochement serait bien léger. % \marqueur{2}. 
La précision de la scholie -- Ναυσιμάχη πόρνη καὶ Σαλαβακχώ -- n'est pour une fois pas une fantaisie de Byzantin ; Nausimaché ne surpasse pas Charminos dans le \italique{combat naval} que par le nom : de par sa profession, elle devait aussi y exceller \italique{en pratique} (τἄργα) -- métaphoriquement, s'entend\marqueur{2} --, ce qui lui valut à la fois son renom et son surnom.}%
\debutNote{804}%
\note{1}{Cf. Thucydide, VIII, 41.}%
%\note{2}{À l'instar des précédents commentateurs qui, du haut de leur tour d'érudition, ignorent l'essentiel : Aristophane est drôle !}%
\note{2}{La langue grecque, et celle d'Aristophane en particulier, fourmille de métaphores nautiques pour rendre de façon imagée les corps-à-corps charnels ; ces images remplissent les sections 179 à 184 des \italique{Images d'Aristophane} de Taillardat et les pages 161-166 de la \italique{Maculate Muse} d'Henderson. Pour faire écho à Nausimaché, ne citons que le verbe ναυμαχεῖν usité par Aristophane aux vers 430 des \italique{Grenouilles} et 675 de \italique{Lysistrata}. Quant à τἄργα, on a déjà souligné l'acception scabreuse que peut avoir ce mot (cf. le vers 204 : ἔργα νυκτερείσια).}\\
\Automatique{καὶ μὲν δὴ καὶ Κλεοφῶν χείρων πάντως δήπου Σαλαβακχοῦς.}\testimonia{805}{(Σαλαβακχοῦς) Hésychius (σ 88) Σαλαβακχώ· πόρνης ὄνομα. ἀπὸ δὲ ταύτης καὶ τὰς καταφερεῖς εἰς τὰ ἀφροδίσια οὕτως ἔλεγον Ἀττικοί}%
\commentaire{805}{Κλεοφῶν χείρων πάντως δήπου Σαλαβακχοῦς.}{Salabachô est un nom générique de prostituée\marqueur{1}. Aristophane le mentionne aussi au vers 765 des \italique{Cavaliers} -- c'est le Paphlagonien qui parle :
\citComm{Κλέων --- τῇ μὲν δεσποίνῃ Ἀθηναίᾳ τῇ τῆς πόλεως μεδεούσῃ\\
εὔχομαι, εἰ μὲν περὶ τὸν δῆμον τὸν Ἀθηναίων γεγένημαι\\
βέλτιστος ἀνὴρ μετὰ Λυσικλέα καὶ Κύνναν καὶ Σαλαβακχώ,\\
ὥσπερ νυνὶ μηδὲν δράσας δειπνεῖν ἐν πρυτανείῳ·\\
Cléon --- \francais{À notre maîtresse Athéna, qui régit cette cité,\\
je fais cette prière ; si pour le peuple des Athéniens je me suis montré\\
le plus méritant des hommes après Lysiclès, Cynna et Salabacchô,\\
puissé-je comme maintenant, sans avoir rien fait, dîner dans le Prytanée.}}%
L'exégèse de ce passage nous permettra de comprendre du même coup celui qui nous intéresse.
Willems avait justement remarqué que les trois noms étaient mis ici par jeu de consonance avec les trois représentants de l'Athènes impérialiste : Périclès, Cimon et Thémistocle\marqueur{2}. 
Pour le sens, voici ce qu'il en est : la modestie -- toujours fausse -- de Cléon, qui se place après ces trois-là, n'est que pure affectation : il fait comprendre qu'il cumule en fait les mérites des trois.
Λυσι-κλῆς, il l'est depuis sa victoire de Sphactérie : il acquit la \italique{célébrité} en \italique{déliant} la situation empêtrée du siège de Sphactérie.
Κύννα, il prétend l'être, lui qui ne cesse de proclamer être le \italique{chien} de garde de la cité\marqueur{3}.
Σαλα-βακχώ se laisse aisément comprendre ; on rapprochera la première partie du nom d'une glose de Photius (σ 43) :%
\citComm{σάλα· φροντίς. οὕτως Αἰσχύλος.}%
On reprochait au braillard Cléon d'être un excité ; il a beau jeu de prétendre l'être en vertu de sa \italique{préoccupation}, du \italique{soin} qu'il porte à la cité. La seconde partie du nom, βακχώ, étant une allusion à l'agitation de tous les fidèles de Bacchus, et notamment les bacchantes. Suétone, dans son Περὶ βλασφημιῶν, \italique{Des termes injurieux}, cite βακχός au chapitre εἰς ἀλαζόνας, «~envers les bavards, les prétentieux~»\marqueur{4}.
%C'est pour la même raison qu'Aristophane compare Salabacchô et Cléophon ; l'avantage est bien-sûr donné à Salabacchô.
\textbf{Κλεοφῶν} était un démagogue dont la trop longue `carrière' politique fut brutalement interrompue en 404 av. J.-C. par une condamnation à mort. Dans les \italique{Grenouilles}, Aristophane charge ce va-t-en-guerre\marqueur{5} dont la langue par trop pendue et la lippe par trop pendante n'émettaient qu'un babil barbarisant criant ses barbaresques origines\marqueur{6}. (Pour un dossier complet sur Cléophon, le lecteur pourra consulter le commentaire de Fritzsche \italique{ad loc.}) %La glossolalie de Cléon trouve un écho en Σαλαβακχώ -- mélioratif pour Σαλαβακχώ puisque Cléon n'hésite pas à se comparer à elle -- c'est ce en quoi Σαλαβακχώ emporte l'avantage sur Cléophon.
}%
\debutNote{805}%
\note{1}{Hésychius (σ 88) : Σαλαβακχώ· πόρνης ὄνομα. ἀπὸ δὲ ταύτης καὶ τὰς καταφερεῖς εἰς τὰ ἀφροδίσια οὕτως ἔλεγον Ἀττικοί.}%
\note{2}{Platon prit les trois mêmes comme exemples d'hommes d'état honorés alors qu'ils sont en partie responsables de la crise athénienne (\italique{Gorgias}, 519a) :%
\citComm{ἄνευ γὰρ σωφροσύνης καὶ δικαιοσύνης λιμένων καὶ νεωρίων καὶ τειχῶν καὶ φόρων καὶ τοιούτων φλυαριῶν ἐμπεπλήκασι τὴν πόλιν.\\
sans se préoccuper de la sagesse ni de la justice, ils ont gorgé la ville de ports, d'arsenaux, de murs, de tributs \& autres niaiseries.}%
%Le détournement d'Aristophane paraît aussi irrévérencieux envers ces trois-là.
}%
\note{3}{Même jeu de mots avec le nom de Κύννα au vers 1032 des \italique{Guêpes} : ἀπ᾽ ὀφθαλμῶν Κύννης rappelle l'expression homérique des \italique{yeux d'un chien} (\italique{Iliade}, I, 225).}%
\note{4}{Περὶ βλασφημιῶν, Περὶ παιδιῶν, édité en 1967 (aux Belles Lettres) par Jean Taillardat -- taillardesque édition, c'est tout dire ! βακχός est la \nieme{127} βλασφημία.}%
\note{5}{Vers 1504 et 1532-3.}%
\note{6}{Vers 678-682. L'on se reportera pour l'exégèse au toujours brillant Alphonse Willems, t. III, p. 115-117.}\\
\Automatique{πρὸς Ἀριστομάχην δὲ χρόνου πολλοῦ, πρὸς ἐκείνην τὴν Μαραθῶνι,}%
\commentaire{806-8}{Ἀριστομάχη, Στρατονίκη et Εὐβούλη}{ sont des noms attestés au moment de la représentation des \italique{Thesmophories}\marqueur{1}. Ces noms sont parlants, bien-sûr, mais l'allusion ; la raison pour laquelle Aristophane les cite, m'échappe. Les deux premières étaient-elles de puissantes princesses contre lesquelles les Athéniens n'auraient pas osé se frotter\marqueur{2} ?}
\debutNote{806-8}
\note{1}{Voir \italique{A Lexicon of Greek Personal Names} de P. M. Fraser et E. Matthews.}%
\note{2}{Στρατονίκη est porté par plusieurs membres de la famille royale macédonienne, dont l'une fut donnée en mariage à l'un des rois de Thrace Σεύθης, contemporain des \italique{Thesmophories} (voir Thucydide, II, 101 et IV, 101). Cicéron (\italique{Tusculanes}, V, 59) nous apprend que Denys lorsqu'il devint Tyran de Syracuse en 406 comptait une certaine \latin{Aristomache} parmi ses épouses.}}
\end{groupe}%
%% v. 785-843
\begin{groupe}%
%
\tetramAnap{\Automatique{καὶ Στρατονίκην ὑμῶν οὐδεὶς οὐδ᾽ ἐγχειρεῖ πολεμίζειν.}\apparat{807}{\lecon{Στρατονίκην} \PsiUn : στρατωνίκην R}\\
\Automatique{ἀλλ᾽ Εὐβούλης τῶν πέρυσίν τις βουλευτής ἐστιν ἀμείνων}\testimonia{808-9}{(τις βουλευτής -- βουλείαν) Suidas β 431 (ἐστι pro ἐστιν et ἀμείνων omisso)}%\apparat{808}{τίς βουλευτής R : \lecon{τις βουλευτής} Küster}
\commentaire{808-9}{ἀλλ᾽ Εὐβούλης τῶν πέρυσίν τις βουλευτής ἐστιν ἀμείνων / παραδοὺς ἑτέρῳ τὴν βουλείαν; οὐδ᾽ αὐτὸς τοῦτό γε φήσεις.}{Ces deux vers ne sont qu'allusifs et ont été diversement interprétés. Ils sont d'autant plus importants que c'est peut-être le seul passage de la comédie qui puisse nous aider véritablement à dater les \italique{Thesmophories}. (Je renvoie à l'Introduction sur la question de la date de notre pièce.)
À quel événement font-ils allusion ? Deux écoles s'affrontent :\\
\ \ \ 1°) \ celle de Wilamowitz-Willems\marqueur{1} qui situe la représentation des \italique{Thesmophories} aux Dionysies de 411 : le Conseil des Cinq-Cents bouleutes fut rendu tout-à-fait caduc par l'instauration d'un conseil de dix vieillards, les πρόβουλοι\marqueur{2}, magistrats à compétence presque illimitée qui prirent en main la direction des affaires. Willems continue : «~C'était délivrer au Conseil des Cinq-Cents un certificat d'incapacité  $[\ldots]$. Aussi, quand il sortit de charge en skirophorion (juin) 412, n'existait-il plus que de nom. %Et si, ce que nous ignorons, l'usage prévalait déjà, comme au temps de Démosthène, de décerner à la Boulè démissionnaire une couronne d'or pour prix de ses services, il est bien sûr que cette fois nul ne se leva à l'Ecclésie pour proposer pareil décret 
[...] C'est à cette transmission annuelle de pouvoirs (παράδοσις, c'est le terme consacré), par laquelle chaque bouleute sortant cédait sa place à son successeur, que le chœur fait allusion~» ;\\
\ \ \ 2°) \ celle de Rogers -- et à laquelle je me rattache -- qui considère que la pièce est d'un an plus tardive : en mai 411, les Quatre Cents conjurés favorables à l'oligarchie se rendirent au βουλευτήριον et versèrent leur salaire aux bouleutes alors en séance. Les bouleutes se virent intimer l'ordre de se retirer, ce qu'ils firent sans même que les conjurés eussent à recourir à la violence\marqueur{3}. Ainsi périt une première fois la démocratie athénienne. C'est ce manque de courage (ἀμείνων) qu'Aristophane brocarde : pas un n'aura défendu la charge qu'ils avaient reçue. Que ces piteux représentants de la démocratie fussent ceux de l'\italique{an dernier} (τῶν πέρυσίν) place donc notre pièce aux Dionysies de l'an 410.}%
%{Des différentes explications apportées à ce passage, nous ne retiendrons que celle de Willems\marqueur{1} -- la seule plausible\marqueur{2} : le Conseil des Cinq-Cents bouleutes fut rendu tout-à-fait caduc par l'instauration d'un conseil de dix vieillards, les πρόβουλοι\marqueur{3}, magistrats à compétence presque illimitée qui prirent en main la direction des affaires. Willems continue : «~C'était délivrer au Conseil des Cinq-Cents un certificat d'incapacité, le droit de formuler un avis préalable étant précisément l'essentiel de sa prérogative. Aussi, quand il sortit de charge en skirophorion (juin) 412, n'existait-il plus que de nom. Et si, ce que nous ignorons, l'usage prévalait déjà, comme au temps de Démosthène, de décerner à la Boulè démissionnaire une couronne d'or pour prix de ses services, il est bien sûr que cette fois nul ne se leva à l'Ecclésie pour proposer pareil décret [...] C'est à cette transmission annuelle de pouvoirs (παράδοσις, c'est le terme consacré), par laquelle chaque bouleute sortant cédait sa place à son successeur, que le chœur fait allusion~»}%
\debutNote{808-9}%
\note{1}{Voir la \nieme{23} note de Willems, reproduite entièrement à la fin.}%
%\note{2}{D'aucuns, trompés par le scholiaste, ont prétendu qu'il y avait dans ces vers une pointe d'Aristophane contre un bouleute précis, non nommé... Mais ne voit-on pas qu'avec la construction τις ... ἀμείνων, τις ne peut avoir qu'une valeur indéfinie ? « Citez m'en un seulement de meilleur », et que si Aristophane pensait à un bouleute en particulier, il aurait dû user d'une construction analogue aux vers précédents avec ἥττων, χείρων ou tout autre synonyme.}%
\note{2}{Voir Thucydide VIII, 1 et la \italique{Constitution d'Athènes}, XXIX ; Aristophane met en scène un Probule dans \italique{Lysistrata} (v. 387-613), jouée quelques mois auparavant -- là est la question : combien ? deux mois comme le pense Willems ?.. ou quatorze comme le pense Rogers ?}%
\note{3}{Événement conté par Thucydide aux chapitre 69-70 de son livre VIII.}\\
\Automatique{παραδοὺς ἑτέρῳ τὴν βουλείαν; οὐδ᾽ αὐτὸς τοῦτό γε φήσει.}\testimonia{809}{(παραδοὺς ἑτέρῳ τὴν βουλείαν) \Antiatt 84.30 \espTest (βουλείαν) \Harp B 19 ; β Photius β 233 βουλία ; Hesychius β 921 βουλεῖα}\apparat{809}{\lecon{αὐτὸς} R : Ἄνυτος Maas \espApp φήσεις R : \lecon{φήσει} Küster \ligneApp{810} εὐχόμεσθ᾽ R : \lecon{εὐχόμεθ᾽} \Mon}\\
\Automatique{οὕτως ἡμεῖς πολὺ βελτίους τῶν ἀνδρῶν εὐχόμεθ᾽ εἶναι.}\\
\Automatique{οὐδ᾽ ἂν κλέψασα γυνὴ ζεύγει κατὰ πεντήκοντα τάλαντα}%\apparat{811}{ζεύγει post γυνὴ R : post οὐδ᾽ ἂν van Herwerden}
\commentaire{811-2}{οὐδ᾽ ἂν κλέψασα γυνὴ ζεύγει κατὰ πεντήκοντα τάλαντα / εἰς πόλιν ἔλθοι τῶν δημοσίων·}{Aristophane use de la même formule κατὰ πεντήκοντα τάλαντα, «~cinquante talents d'un coup~», dans les \italique{Guêpes} : Bdélycléon tance son paternel berné par les démagogues qui, non contents de dissiper l'argent de la cité, détournent le tribut des autres cités (vers 666-671) : %
\citComm{Φιλοκλέων --- καὶ ποῖ τρέπεται δὴ ᾽πειτα τὰ χρήματα τἄλλα;\\
Βδελυκλέων --- εἰς τούτους τοὺς “οὐχὶ προδώσω τὸν Ἀθηναίων κολοσυρτόν,\\
ἀλλὰ μαχοῦμαι περὶ τοῦ πλήθους ἀεί”. σὺ γάρ, ὦ πάτερ, αὐτοὺς\\
ἄρχειν αἱρεῖ σαυτοῦ τούτοις τοῖς ῥηματίοις περιπεφθείς.\\
κᾆθ᾽ οὗτοι μὲν δωροδοκοῦσιν \textbf{κατὰ πεντήκοντα τάλαντα}\\
ἀπὸ τῶν πόλεων ἐπαπειλοῦντες τοιαυτὶ κἀναφοβοῦντες·\\
“δώσετε τὸν φόρον, ἢ βροντήσας τὴν πόλιν ὑμῶν ἀνατρέψω.”\\
Philocléon --- \francais{Et où passe alors le reste de l'argent ?}\\
Bdélycléon --- \francais{À ceux qui nous disent: « Non, je ne trahirai point la masse athénienne, toujours je combattrai pour le populaire ! » Car c'est toi, mon père, qui te les donnes pour maîtres, enjôlé par leurs boniments. Eux cependant extorquent aux cités des cinquante talents à la fois, en les menaçant, en les terrifiant de cette sorte : « Vous paierez le tribut, ou je tonne contre votre ville et la renverse. »}}%
Hérodote (VI, 132-136) rapporte une histoire semblable : celle de Miltiade. Le vainqueur de Marathon, promit aux Athéniens quantité d'or s'ils lui remettaient soixante-dix vaisseaux avec troupe et argent. Les Athéniens lui fournirent le tout et Miltiade de cingler vers Paros pour en faire le siège. Il envoya un héraut réclamer 100 talents avec le même genre de menaces que sus-citées. L'expédition fut un échec et l'on intenta un procès à Miltiade pour avoir trompé la cité (τῆς Ἀθηναίων ἀπάτης εἵνεκεν). Or Miltiade se mourait d'une mauvaise blessure. 
Compte tenu des services rendus à la nation, et que c'était certainement peine perdue de transporter ce semi-cadavre jusqu'au \italique{barathre}\marqueur{1}, l'agonisant ne fut condamné qu'à restituer \italique{cinquante talents}.
%Avez-vous, vous aussi, été saisi par 
L'histoire de Miltiade et les deux passages d'Aristophane coïncident jusque dans le chiffre du détournement !..
Miltiade, le héros de Marathon -- τὸν Μαραθῶνι, c'est son épithète, quasi-homérique, qui lui est inlassablement accolée --, profita du crédit qu'il avait auprès de Dêmos, le cajolant des souvenirs de Marathon\marqueur{2}, pour mener une expédition dont le but véritable était d'enflouer ses propres caisses.\\
En regard des déboires de Miltiade, la construction de l'argumentation du chœur féminin se comprend parfaitement : «~Devant Aristomaque, notre héroïne de Marathon (τὴν Μαραθῶνι), aucun d'entre vous n'entreprendrait de se battre~», et, parant d'avance à l'objection que les hommes ont aussi leur marathonien héros en la personne de Miltiade, l'on précise aussitôt : «~Une femme ayant volé cinquante talents d'un (bon) coup n'aurait pas l'outrecuidance de venir parader sur un char lors de son retour à la cité~».\\
\italique{Confer} sur ces deux vers 811-812 la note 24 d'Alphonse Willems (reproduite en annexe) : à la fête des Panathénées, la file brillante des chars était ce qui attirait le plus les regards dans la procession (Schoemann, \italique{Antiquités grecques}, t. II, p. 549).
}%
\debutNote{811-2}%
\note{1}{Platon (\italique{Gorgias}, 516e) nous apprend que Miltiade devait être initialement précipité vers le \italique{barathre} (εἰς τὸ βάραθρον ἐμβαλεῖν), illustrant l'adage latin qu'il n'y a pas loin du Capitole à la roche tarpéienne. (Pour en apprendre davantage sur ce gouffre, que le lecteur se reporte au \italique{Dictionnaire des antiquités grecques et romaines}, à l'entrée \latin{Barathron}.)
Si Miltiade s'est vu gracier -- la Parque se consola néanmoins très rapidement : la grâce ne fut qu'un sursis de quelques jours --, la raison principale est peut-être qu'il bénéficiait de nombreuses relations, critère qui distingue souvent les condamnés à mort exécutés de ceux qui voient leur peine commuée, selon le mot de Gide.}%
\note{2}{«~Flattez-les, payez-les d'agréables mensonges : /	Quelque indignation dont leur cœur soit rempli, / Ils goberont l'appât ; vous serez leur ami.~»}%
\\%
\Automatique{εἰς πόλιν ἔλθοι τῶν δημοσίων· ἀλλ᾽ ἢν τὰ μέγισθ᾽ ὑφέληται}\\
%
\Automatique{φορμὸν πυρῶν τἀνδρὸς κλέψασ᾽, αὐθημερὸν αὔτ᾽ ἀπέδωκεν}\apparat{813}{\lecon{αὔτ᾽ ἀπέδωκεν} R, M\expo{pc} : ἀνταπέδωκεν M\expo{ac}, Bentley}\\\passeLigne}%
\end{groupe}%
%% v. 814-829
\begin{groupe}%
%
\personnageGrec{\strophe{πνῖγος}}\typeVers{\decalage{πνῖγος \ }dim. anap.}\anap{ἀλλ᾽ ἡμεῖς ἂν πολλοὺς τούτων%
\commentaire{814}{τούτων.}{Les spectateurs.}\\
ἀποδείξαιμεν ταῦτα ποιοῦντας.\apparat{815}{ἀποδείξομεν R : \lecon{ἀποδείξαιμεν} Bentley \espApp ποθοῦντας R : \lecon{ποιοῦντας} Zanetti (\italique{facientes} Divus)}%\espApp ταῦτα R : ταὐτά` Bibi 
\commentaire{815}{ποθοῦντας.}{L'erreur manifeste du scribe du \italique{Ravennas}, copiant ποθοῦντας à la place de πο(ι)οῦντας, m'est l'occasion -- ὁ καιρός ! -- d'une courte digression portant sur les différentes graphies ποι- ou πο- lorsque la lettre qui suit est un `ο' ou un `ω'\marqueur{1}.
Je m'en tiens aux leçons de l'épigraphie attique ainsi qu'à l'usage scolaire : savoir, écrire systématiquement ποι-. Toutefois, le \italique{Ravennas} adopte majoritairement la graphie πο- lorsque la syllabe est brève. Cette convention, reprise par la plupart des éditions actuelles des \italique{Thesmophories}, est celle de nombreux manuscrits remontant à la période hellénistique. La confusion du scribe -- qui aurait lu rapidement une suite de `Ο' en y ajoutant un `Θ' -- tend à faire penser que cette règle était aussi de mise pour le manuscrit qu'il recopiait.}%
\debutNote{815}%
\note{1}{Débat qui a fait déjà couler beaucoup d'encre et que je ne compte ni trancher, ni même exposer en détail. Le lecteur intéressé pourra se reporter aux pages 226-230 que lui consacre Willem J. W. Koster dans son livre \italique{Autour d'un manuscrit d'Aristophane écrit par Demetrius Triclinius} (1957).}\\
καὶ πρὸς τούτοις γάστριδας ἡμῶν\testimonia{816}{(γάστριδας) Photius γ 40 γάστρις $=$ Suidas γ 72 ; cf. \AelDion γ 3}\\
ὄντας μᾶλλον καὶ λωποδύτας\\
καὶ βωμολόχους κἀνδραποδιστάς.\\
καὶ μὲν δήπου καὶ τὰ πατρῷά γε\apparat{819}{δήπου R : δή που Denniston}%
\commentaire{819-23}{}{%\eff{Euripide Medée fr. 4 ?γυνὴ γὰρ ἐξελθοῦσα πατρῴων δόμων / οὐ τῶν τεκόντων ἐστίν, ἀλλὰ τοῦ λέχους. ἀνδρὸς Ἕλληνος  Γυνὴ μὲν“ἐξελϑοῦσαπατρῴωνδόμων οὐ τῶν “τεκόντων ἐστίν, ἀλλὰ τοῦ λέχους. τὸ δ᾽ ἄρσενἕστηκ'ἐν δόμοιςἀεὶ γέγος ϑεῶνπατρῴωνκαὶ τάφωντιμᾶ ὅρον}
Le manuscrit présente le texte suivant :%
\citComm{καὶ μὲν δήπου καὶ τὰ πατρῷά γε\\
χείρους ἡμῶν εἰσιν σῴζειν·\\
ἡμῖν μὲν γὰρ σῶν ἔτι καὶ νῦν\\
τἀντίον ὁ κανὼν οἱ καλαθίσκοι\\
τὸ σκιάδειον·}%
Les trois premiers vers ne posent pas de difficulté. Et encore : à condition de garder à l'esprit les divers sens de τὰ πατρῷα, \italique{l'héritage} : désigne les \italique{biens transmis par les ancêtres} ; \italique{héritage} tout matériel que seul une femme sera à même de conserver\marqueur{1}, mais héritage aussi immatériel, \italique{les traditions}\marqueur{2}, que les femmes perpétuent tout-de-même moins mal que les hommes (χείρους ἡμῶν εἰσιν σῴζειν).
Suit une liste de ces biens que les femmes sauvent mais que leur mari, eux, perdent. C'est là que le bât blesse.
D'abord, cette liste est étonnamment hétéroclite ; si les trois premiers éléments qu'elle comporte sont en rapport avec le filage et le tissage : τὸ ἀντίον, \italique{l'ensouple} (cyclindre du métier à tisser) ; ὁ κανὼν, \italique{barre}, \italique{règle} (avec le sens spécialisé de \italique{tige de quenouille}), οἱ καλαθίσκοι, \italique{petits paniers} (où les femmes déposaient la laine filée), l'on se demande par quel ressort \italique{l'ombrelle}, τὸ σκιάδειον, vient clôturer cette liste.
Plus important : sur ces quatre éléments, seuls deux, κανών et σκιάδειον, subsistent dans la comparaison avec ce qu'en font les hommes. Pourquoi Aristophane nous parle-t-il alors de l'ἀντίον et des καλαθίσκοι, inutiles -- donc nuisibles -- si ce n'est pas pour filer la métaphore quelques vers plus loin ? (Je ne parle que de ce qui achoppe au niveau du sens. Quant à la forme, l'on avait déjà relevé l'inhabituelle succession de quatre brèves au vers 822.)\\
%Comme souvent en pareille situation, le mieux est encore de se mettre à la place d'Aristophane. 
Aristophane veut nous peindre la sotte dissipation des hommes ; Interrogeons les parémiographes\marqueur{3} pour savoir comment on évoque une prodigalité de cigale ou une gestion de fourmi dans sa langue imagée et proverbiale :%
\citComm{Αὐτῷ κανῷ: αὕτη τάττεται κατὰ τῶν ἄρδην ὁτιοῦν λυμαινομένων. Φησὶ γὰρ ὁ Δίδυμος ὅτι τὰ δεῖπνα ἐπὶ κανῶν ἐκομίζετο. Καθάπερ ἐν Γρυψὶ λέγει καὶ Πλάτων·\\
Ἀνήρπακας ἁπαξάπαντ᾿ \textbf{αὐτῷ κανῷ}.\\
Et la corbeille elle-même : ce [proverbe] se range contre ceux prompts à se ruiner de quelque manière. Didyme dit que les repas étaient emportés dans des paniers. Ainsi Platon dans ses \italique{Griffons} :\\
Tu as tout emporté, tout, \italique{et le panier avec} !..}%
Le grec ajoute «~et le panier~» pour désigner tous les biens jusqu'au dernier\marqueur{4}, comme en français on ajoute «~et son père~» pour désigner «~tout le monde~». Dans les \italique{Grenouilles} (vers 560), ce n'est pas autrement qu'il faut entendre le désarroi de l'hôtesse face à Hercule, cet ogre qui a liquidé toutes ses provisions :%
\citComm{Μὰ Δί᾽ οὐδὲ τὸν τυρόν γε τὸν χλωρόν, τάλαν,\\
ὃν οὗτος \textbf{αὐτοῖς τοῖς ταλάροις} κατήσθιεν.\\
Par Zeus, et ce fromage frais, pauvre amie, qu'il dévorait \italique{avec les éclisses elles-mêmes}.}%
C'est peut-être encore ce même proverbe qu'on retrouve au vers 170 des \italique{Guêpes} dans la bouche de Philocléon :%
\citComm{μὰ τὸν Δί᾽ οὐ δῆτ᾽, ἀλλ᾽ ἀποδόσθαι βούλομαι\\
τὸν ὄνον ἄγων \textbf{αὐτοῖσι τοῖς κανθηλίοις}·\\
νουμηνία γάρ ἐστιν.\\
Par Zeus, certes non ! Je veux aller vendre l'âne \italique{avec ses paniers} ; car c'est nouvelle lune.}%
L'«~αὐτῇ λόγχῃ~» du vers 826 -- noter la présence du pronom --, mis comme permutant métaphorique d'«~αὐτῷ κανόνι~», rappelle le détournement du proverbe avec jeu de mots sur κανών / κανοῦν.\\
Notre conclusion est que le texte original des \italique{Thesmophories} -- l'\allemand{Urtext} -- comportait certainement au vers 822 le terme κανοῦν, \italique{corbeille}, et non κανών. %Cette hypothèse permet de faire coup double et d'expliquer la singulière présence d'
La singulière présence des καλαθίσκοι, \italique{les corbeilles à laine}, s'explique alors par l'insertion fort importune dans le texte d'une glose de κανοῦν. %de qui n'est que le fruit d'une glose de κανοῦν fort importunément ajoutée au texte.
\\
L'image ne porte pas sur l'activité filandière ; la présence simultanée de corbeilles et d'une ombrelle ne laisse pas de doute possible : Aristophane nous décrit une procession religieuse\marqueur{5}. 
Au sein de ces cortèges traditionnels, on voyait défiler les \italique{cané-}phores, qui avaient la charge des corbeilles à offrandes, et qu'elles portaient sur la tête\marqueur{6}, quand d'autres, les \italique{skiro-}phores, les accompagnaient\marqueur{7} en tenant de deux mains fermes et assurées, la lourde \italique{ombrelle} consacrée, τὸ σκίρον ou τὸ σκιάδειον. 
Et pour Aristophane, sa foi, ces parades d'Athéniennes défilant cahin-caha n'étaient pas tout-à-fait indignes des traditions reçues des ancêtres (τὰ πατρῷα).
À la place des vers 822-823, je propose donc de lire :%
\citComm{καὶ μὲν δήπου καὶ τὰ πατρῷά γε\\
χείρους ἡμῶν εἰσιν σῴζειν·\\
ἡμῖν μὲν γὰρ σῶν ἔτι καὶ νῦν\\
ταὐτὸ κανοῦν\marqueur{8} [καὶ] τὸ σκιάδειον·\\
Et assurément aussi, ils sont moins aptes que nous\\
à garder les traditions\\
qui sont, par notre fait, sauves encore aujourd'hui\\
-- même le panier [et] l'ombrelle.}%
%ταὐτό
%Hésychius ἐπίρρηται: ἐπικάτακται
%τανταλίζεται: σαλεύεται (Αr fr. 963 Κ.-Α.
}%
\debutNote{819-23}%
\note{1}{L'homme qui avait aliéné son patrimoine, ὁ τὰ πατρῷα κατεδηδοκώς, était frappé d'ἀτιμία. Eschine, \italique{In Timarchum}, 30 ; Diogène Laërce, \italique{Solon}, I, 55. Cette loi qui n'était certainement plus observée au temps d'Eschine subsistait pour la forme, comme un vestige de l'antique règle ; il y eut toujours une δίκη κατεδηδοκέναι τὰ πατρῷα (Bekker, \italique{Anecdota}, p. 199 et 310 ; note de Fustel de Coulanges).}%
\note{2}{Certains commentateurs n'ont pas manqué de noter que les Athéniennes n'auraient pas eu légalement de droit de propriété. Cette remarque aurait pu -- aurait dû -- leur mettre la puce à l'oreille : c'est bien en tant que championnes des \italique{traditions} que se posent les femmes du chœur dans ce passage.}%
\note{3}{Leutsch \& Schneidewim, \italique{Corpus Paroemiographorum Graecorum}, vol. I, Zenobius, II, 31. Recensé par Érasme (\italique{Adages}, n° 1088). Kassel \& Austin l'ont intégré comme \nieme{15} fragment de Platon le Comique. Suidas (α 4465 ; α 4466 ; α 4540) a relevé des expressions analogues : αὐταῖς ἀμάλαις ; αὐταῖσι διαβολαῖς ; αὐτοῖς ὅπλοις ; αὐτοῖς ταλάροις.}%
\note{4}{Biens dévorés, en l'occurrence. Mais l'image du panier comme signe d'abondance se retrouve dans un autre proverbe, transmis par Hésychius (ε 7077) : εὔπλουτον κανοῦν.}%
\note{5}{La fête des Σκίρα est évoquée quelques vers plus loin. Cf. les articles \latin{Umbella} et \latin{Skirophoria} du \italique{Dictionnaire des antiquités grecques et romaines} de Daremberg \& Saglio.}%
\note{6}{Ces canéphores ont certainement servi de modèle aux fameuses Caryatides, portant, elles, un chapiteau sur la tête.}%
\note{7}{Voir, par exemple, le vers 1550 des \italique{Oiseaux}. Les ombrelles tenaient effectivement un rôle important dans la liturgie.}%
\note{8}{Au vers 164, il y a confusion de finale -ων / -ουν du même genre entre les leçons du \italique{Ravennas} (διεκίνων) et de la \italique{Souda} ε 989 (διεκίνουν).}%
\\
χείρους ἡμῶν εἰσιν σῴζειν·}
\end{groupe}%
%% v. 814-829
\begin{groupe}%
%
\anap{ἡμῖν μὲν γὰρ σῶν ἔτι καὶ νῦν\\
ταὐτὸ κανοῦν καὶ τὸ σκιάδειον·\testimonia{822}{%(τἀντίον ὁ κανὼν) 
Pollux VII 35 κανὼν ἱστοῦ τὸ καλούμενον ἀντίον}\\
% \commentaire{821-3}{τἀντίον ὁ κανὼν οἱ καλαθίσκοι / τὸ σκιάδειον·}{Sur le sens exact de ces mots, voir l'excellente thèse de M. Fr. Sosset, \italique{Le Tissage dans la Grèce antique}, Bruxelles (1896), pp. 11 à 25. (Extrait de la \italique{Revue de l'Université de Bruxelles}. Note de Willems.)}%
\apparat{822-3}{τἀντίον ὁ κανὼν οἱ καλαθίσκοι / τὸ σκιάδειον· R : \lecon{ταὐτὸ κανοῦν καὶ τὸ σκιάδειον·} Bibi \ligneApp{824} ἀνδράσιν R : \lecon{ἀνδράσι} \Mon \ligneApp{825} ἀπόλωλε R : \lecon{ἀπόλωλεν} Daubuz, Bentley}%
\ligne{824}τοῖς δ᾽ ἡμετέροις ἀνδράσι τούτοις\\
ἀπόλωλεν μὲν πολλοῖς ὁ κανὼν\\
ἐκ τῶν οἴκων αὐτῇ λόγχῃ,%
\commentaire{826}{ἐκ τῶν οἴκων.}{Précisons le sens «~ἐκ τῶν οἴκων~». Ce ne sera pas du luxe. Car il ne peut être question de rompre des lances \italique{à l'extérieur des maisons} plutôt qu'à \italique{à l'intérieur} ? % : les traductions actuelles nous font comprendre que le reproche des femmes porte sur ce que les maris rompent leur lance \italique{à l'extérieur des maisons}. Préféreraient-elles qu'ils le fassent \italique{à l'intérieur} ? Les commentateurs vont encore plus loin : ils nous 
%Entre les explications expliquent que le fameux héritage paternel, τὰ πατρῷα, dissipé par les hommes n'est autre que la panoplie militaire -- ô scandale justifiant de couvrir d'infamie tous ces fils prodigues indignes !! 
Devrons-nous aussi commenter les commentaires qui voient en cette panoplie militaire le fameux héritage paternel, τὰ πατρῷα, que les fils prodigues indignes dissiperaient -- ô scandale !
Au garde-à-vous, je félicite ces brillants universitaires pour leurs éclaircissements : sans eux, je crois que les douceurs martiales de la poésie d'Aristophane m'auraient échappé. Notre poète roulant des variations sur le thème : «~Tu seras un soldat, mon fils !~», c'est à vous tirer des larmes (de guerre, bien-sûr !)\\
%Bref !.. Sortons d'Oxford cinq minutes et remettons les pieds sur Terre. 
Οἶκος, ce mot rare et qui renvoie à des réalités n'existant pas ailleurs qu'en Grèce classique a tout-de-même donné le terme οἰκο-νομία, qui me semble être la bonne gestion d'une \italique{maison}, perçue comme un \italique{fonds} ou un \italique{capital}. «~ἐκ τῶν οἴκων~» désigne tout ce que l'on \italique{tire de} ce fonds, ou de ce capital. La bonne équation est évidemment celle-ci : τὰ πατρῷα $ = $ ἐκ τῶν οἴκων.
La même tournure avec le même sens se trouve chez Xénophon dans son ... \italique{Économique} (XX, 21) :%
\citComm{τὰ οὖν συντρίβοντα τοὺς οἴκους πολῦ μᾶλλον ταῦτά ἐστιν ἢ αἱ λίαν ἀνεπιστημοσύναι. τὸ γὰρ τὰς μὲν δαπάνας χωρεῖν ἐντελεῖς \textbf{ἐκ τῶν οἴκων}, τὰ δὲ ἔργα μὴ τελεῖσθαι λυσιτελούντως πρὸς τὴν δαπάνην\\
Voilà qui ruine une maison beaucoup plus sûrement qu'une excessive ignorance : d'un côté, dépenser complètement \italique{ce qu'on tire des maisons}, de l'autre des travaux qui ne sont pas profitables à la dépense}%
Traduisez donc :%
\citComm{τοῖς δ᾽ ἡμετέροις ἀνδράσι τούτοις\\
ἀπόλωλεν μὲν πολλοῖς ὁ κανὼν\\
ἐκ τῶν οἴκων αὐτῇ λόγχῃ\marqueur{1}.\\
Du fait de nos ci-devant maris (elles les désignent),\\
la lance a anéanti pour beaucoup\\
ce qu'on tirait de nos maisons -- lance comprise\marqueur{1}.}%
}%
\debutNote{826}%
\note{1}{Écho au proverbe : αὐτῷ κανῷ... sauf que le \italique{canon} est remplacé par la \italique{lance}. Autre sexe, autre proverbe :%
\citComm{ἦν γάρ τις αἶνος, ὡς γυναιξὶ μὲν τέχναι\\
μέλουσι, \textbf{λόγχῃ δ ̓ ἄνδρες} εὐστοχώτεροι.\\
εἰ γὰρ δόλοισιν ἦν τὸ νικητήριον,\\
ἡμεῖς ἂν ἀνδρῶν εἴχομεν τυραννίδα.\\
Il existe en effet un proverbe qui dit qu'à la femme les artifices reviennent, tandis que \italique{les hommes sont plus habiles à atteindre le but par la lance}. Si le prix de la victoire allait à la ruse, c'est nous qui régnerions sur les hommes.}%
(Euripide, fr. 6 \italique{Danaé} Jouan-van Looy $=$ fr. 321 Kannicht $=$ fr. 413 Mette.)
}\\
πολλοῖς δ᾽ ἑτέροις ἀπὸ τῶν ὤμων\\
\{ ἐν ταῖς στρατιαῖς \}%
}%
\end{groupe}%
%% v. 814-829
\begin{groupe}%
\typeVers{(catalect.)}\anap{%
ἔρριπται τὸ σκιάδειον.%
\commentaire{827-9}{πολλοῖς δ᾽ ἑτέροις ἀπὸ τῶν ὤμων / ἐν ταῖς στρατιαῖς / ἔρριπται τὸ σκιάδειον.}{Un bouclier est une \italique{ombrelle}. Cette métaphore unique est très claire : un bouclier est rond (cf. κύκλος, Taillardat § 642) et bombé comme une ombrelle. Mais on peut se demander si l'argot des soldats athéniens ne connaissait aussi pas cette métaphore unique\marqueur{1} : le bouclier protège des coups comme l'ombrelle du soleil, cf. Σ : τὴν ἀσπίδα λέγει· σκιάζει\marqueur{2} γὰρ τὸν ἄνδρα. Cette hypothèse est malheureusement invérifiable.
(Note de Taillardat, §643. L'utilisation deux vers plus haut de κανών comme synonyme de \italique{lance}, λόγχη, pourrait relever de ce même argot militaire dont Taillardat fait la très astucieuse hypothèse.)\\
La charge soldatesque des hommes ? elle leur passe par-dessus la tête et les épaules lorsqu'ils en délestent leur bouclier (ἀπὸ τῶν ὤμων ἔρριπται τὸ σκιάδειον), bien encombrant il est vrai lorsqu'il faut prendre la poudre d'escampette en plein milieu des combats... 
Mica, la femme de Cléonyme -- Cléonyme qui semble avoir fait école avec sa célèbre ῥίψασπις (voir le commentaire au vers 605) -- devait probablement à ce moment être pointée, au moins des yeux.}%
\debutNote{827-9}%
\note{1}{À propos de θύλακος (§268), on [l'«~on~», ou l'«~om~» en question n'est autre que le grand Taillardat !] a déjà relevé une métaphore ironique qui a dû appartenir à l'argot des soldats athéniens. Cette tendance à ironiser sur ses armes a toujours existé : on sait que les soldats de César appelaient, par plaisanterie, \latin{cippi} et \latin{lilia} certaines défenses qu'ils construisaient devant Alésia (\italique{Bell. Gall.} 7, 73, 4 et 8). C'est de la même façon que dans l'argot militaire français, un fusil peut s'appeler une \italique{arbalète}, une \italique{seringue}, une \italique{canne à pêche}, etc. (voir A. Dauzat, \italique{L'argot de la guerre}\expo{2} (1919).)}%
\note{2}{Σκιαζειν, \italique{couvrir} : Hésiod. \italique{Théog.} 716 ; Hdt. 6, 117, etc.}%
\\\passeLigne}%
\end{groupe}%
%% v. 830-845 (part 1)
\begin{groupe}\typeVers{$4$ troch\cat}\personnageGrec{Κορυφαῖος}\tetramTroch{πόλλ᾽ ἂν αἱ γυναῖκες ἡμεῖς ἐν δίκῃ μεμψαίμεθ᾽ ἂν\\
τοῖσιν ἀνδράσιν δικαίως, ἓν δ᾽ ὑπερφυέστατον.\\
\Automatique{χρῆν γάρ, ἡμῶν εἰ τέκοι τις ἄνδρα χρηστὸν τῇ πόλει,}\apparat{832}{χρὴ R : \lecon{χρῆν} Brunck}\\
ταξίαρχον ἢ στρατηγόν, λαμβάνειν τιμήν τινα,\\
\Automatique{προεδρίαν τ᾽ αὐτῇ δίδοσθαι Στηνίοισι καὶ Σκίροις}\testimonia{834}{(Στηνίοισι) Suidas σ 1089 Στήνια καὶ Σκίρα ($=$ Scholium) ; Hesychius σ 1825 στήνια}\apparat{834}{Τηνίοισι R : \lecon{Στηνίοισι} Σ, Suid. σ 1089}
\commentaire{834}{Στηνίοισι καὶ Σκίροις.}{Fêtes athéniennes réservées aux femmes. C'est notamment aux Sciries que fut élaboré le projet politique de Praxagora (\italique{Assemblée des femmes}, vers 18).}\\
ἔν τε ταῖς ἄλλαις ἑορταῖς αἷσιν ἡμεῖς ἤγομεν·\\
εἰ δὲ δειλὸν καὶ πονηρὸν ἄνδρα τις τέκοι γυνή,\apparat{836}{ἤν δὲ R : \lecon{εἰ δὲ} Brunck}\\
ἢ τριήραρχον πονηρὸν ἢ κυβερνήτην κακόν,
\commentaire{833-7}{ταξίαρχον ἢ στρατηγόν ... ἢ τριήραρχον πονηρὸν ἢ κυβερνήτην κακόν}{Les exemples d'Aristophane sont tirés de l'actualité la plus brûlante. 
Les mauvais triérarques et commandants de vaisseaux sont ceux responsables de la défaite navale d'Érétrie à l'été 411, défaite qui plongea les Athéniens dans une consternation, nous rapporte Thucydide (VIII, 95 et 96), plus grande encore que celle consécutive au désastre de Sicile. Défaite cinglante qui ne rattrape pas aux yeux d'Aristophane la récente et lointaine victoire de Κυνὸς Σῆμα (VIII, 104-108) dans l'Hellespont. (Cette victoire devait être bien amère à Aristophane, puisque son contemporain Eupolis y mourut. Nous y reviendrons.)}\\\passeLigne}%
\end{groupe}%
\endnumbering
\end{Leftside}
\begin{Rightside}
\beginnumbering
%% v. 800+ (part 2)
\begin{groupe}
\tetramAnapFr{\ligne{800}que vous. Si l'on veut une preuve, il n'y a qu'à regarder.
Venons-en à la preuve et sachons qui sont les pires. Car nous disons que c'est vous,}
\end{groupe}%
%% v. 785-843
\begin{groupe}%
\tetramAnapFr{et vous dites que c'est nous. Vérifions donc et mettons-les en présence isolément
en rapprochant chaque fois le nom de la femme de celui de l'homme.
Charminos est inférieur à Nausimaché : les faits parlent d'eux-mêmes,
et à coup sûr aussi Cléophon est pire en somme que Salabaccho.
Avec Aristomaché, l'héroïne de \fin Marathon,}
\end{groupe}%
%% v. 785-843
\begin{groupe}%
\tetramAnapFr{\ligne{806}et avec Stratonice, il y a beau jour que nul d'entre vous ne tenterait de se mesurer.
\ligne{807}Et Euboulé, parmi les bouleutes de l'an dernier, en est-il un qui,
après avoir résigné à un autre son office, l'emporte sur elle ? Lui-même s'en défendra.
Ainsi nous nous prévalons d'une grande supériorité sur les hommes.
Joignez qu'une femme, après avoir diverti par cinquante talents les deniers publics,
ne se rendra pas à l'Acropole sur un char ; mais eût-elle commis le plus grave des larcins,
%
un boisseau de froment volé à son mari, elle le restituera le jour même.\\\passeLigne}%
\end{groupe}%
%% v. 814-829
\begin{groupe}%
\anapFr{\ligne{814}Mais nous pourrions en désigner plusieurs \fin parmi ceux-là\\
qui en font autant,\\
et qui, par là-dessus, sont gloutons\\
plus que nous, tire-laine,\\
charlatans et marchands d'hommes.\\
Et assurément aussi ils savent moins que nous\\
conserver les traditions}
\end{groupe}%
%% v. 814-829
\begin{groupe}%
\anapFr{qui sont par notre fait sauves encore aujourd'hui\\
– panier et ombrelle compris.\\
Au contraire, du fait de nos nombreux maris,\\
la lance a ruiné\\
notre capital – lance comprise.\\
Et maint autre\\
$\{$dans ses campagnes$\}$}%
\end{groupe}%
%% v. 814-829
\begin{groupe}%
\anapFr{%
\ligne{829}a jeté à bas de ses épaules... l'ombrelle.\\\passeLigne}%
\end{groupe}%
%% v. 830-845 (part 1)
\begin{groupe}\personnageFr{Coryphée}
\tetramTrochFr{Nous autres femmes aurions droit en toute justice
d'imputer aux hommes bien des torts, un surtout passant toute mesure.
Il serait de votre devoir, quand l'une de nous a donné à la patrie un homme utile,
\ligne{833}taxiarque ou stratège, d'avoir pour elle quelque considération,
de lui accorder la proédrie dans les Sténies, les Scirophories
et les autres fêtes que nous célébrons.
Celle, au contraire, qui aurait donné le jour à un pleutre, à un propre à rien,
méchant triérarque ou lâche commandant de navire,\\\passeLigne}%
\end{groupe}%
\endnumbering
\end{Rightside}
\end{pages}
\Pages
\clearpage\ifodd\value{page}\hbox{}\newpage\fi
\enTeteG{ΠΑΡΑΒΑΣΙΣ / ΕΠΕΙΣΟΔΙΟΝ Δ' --- Codex Ravennas 429, f. 172v (Classense 174v) --- vv. 838–875}
\begin{adjustwidth}{-1.5cm}{-3.0cm}
\begin{adjustbox}{max totalheight=21.0cm, max width=16.5cm, keepaspectratio}
\begin{minipage}{16.5cm}
{\fontsize{7.2pt}{8.5pt}\selectfont
\noindent\textbf{Folio 21 --- Codex Ravennas 429, \texttt{f. 172v} (\texttt{Classense 174v}) : ΠΑΡΑΒΑΣΙΣ / ΕΠΕΙΣΟΔΙΟΝ Δ', vv. 838–875}\hfill \textit{Témoins :} $\text{R}^1, \Rpc, \Rdeux$, \PsiUn\par\vspace{2pt}
\hrule height 0.4pt\vspace{2.5pt}
}
{\fontsize{8.1pt}{9.5pt}\selectfont
\ligneRav{838}{χορ/}{ὑστέρ\auDessus{ο}{α}ν αὐτὴν καθῆσθαι σκάφιον ἀποκεκαρμένην}{\textbf{838} \lecon{ὑστέρ\auDessus{ο}{α}ν} \Rpc : ὑστάτην Σ}
\ligneRav{839}{}{τῆς τὸν ἀνδρεῖον τεκούσης· τῷ γὰρ εἰκὸς \difference{ωὖπολι}}{\textbf{839} ὦ πόλεις R : ὦ πόλις Gelenius : \lecon{ωὖπολι} Bibi}
\ligneRav{840}{}{τὴν Ὑπερβόλου καθῆσθαι μητέρ᾽ ἠμφιεσμένην}{}
\ligneRav{}{}{λευκὰ καὶ κόμας καθεῖσαν πλησίον τῆς Λαμάχου·}{}
\ligneRav{842}{}{καὶ δανείζειν \difference{χρήμαθ᾽}, ἥν χρῆν, εἰ \difference{δανείσειέν} τινι}{\textbf{842} χρήματα τ' R : \lecon{χρήμαθ᾽} Scaliger \espApp \lecon{ἥν χρῆν} R : ᾗ χρῆν Scaliger \espApp χρῆν ἦν R : \lecon{χρῆν εἰ} Brunck \espApp δανείσειε R : \lecon{δανείσειέν} Biset}
\ligneRav{}{}{καὶ τόκον πράττοιτο, διδόναι μηδέν᾽ ἀνθρώπων τόκον,}{}
\ligneRav{844}{}{ἀφαιρεῖσθαι τ\auDessus{ὰ}{ο}δὶ τὰ χρήματα,}{\textbf{844} ἀφερεῖσθαι R : \lecon{ἀφαιρεῖσθαι} \Mon \espApp χρήματατ᾽ R : \lecon{χρήματ᾽} Grynaeus \espApp τ\auDessus{ὰ}{ο}δὶ \Rpc : \lecon{τοδί} \Mon}
\ligneRav{845}{}{`ἀξία γοῦν εἶ τόκου τεκοῦσα τοιοῦτον τόκον.'}{}
\ligneRav{}{μ\auDessus{ν}{Η}}{ἰλλὸς γεγένημαι προσδοκῶν· ὁ δ᾽ οὐδέπω·}{}
\ligneRav{}{}{τί δῆτ᾽ ἂν εἴη τοὐμποδών; οὐκ ἔσθ᾽ ὅπως}{}
\ligneRav{848}{}{οὐ τὸν \difference{Παλαμήδη} ψυχρὸν ὄντ᾽ αἰσχύνεται·}{\textbf{848} Παλαμήδην R : \lecon{Παλαμήδη} Elmsley (\italique{Quart. Rev.} 7 (1812) p.454)}
\ligneRav{}{}{τῷ δῆτ᾽ ἂν αὐτὸν προσαγαγοίμην δράματι;}{}
\ligneRav{850}{}{ἐγᾦδα· τὴν καινὴν Ἑλένην μιμήσομαι·}{}
\ligneRav{851}{}{\difference{πάντως δ᾽} ὑπάρχει μοι γυναικεία στολή·}{\textbf{851} πάντως R : \lecon{πάντως δ᾽} Bentley, cf. v.1012}
\ligneRav{852}{\auDessus{γυ}{ν}}{τί αὖ σὺ κυρκανᾷς; \difference{τί κοικύλλεις} ἔχων;}{\textbf{852} ἢ τί κοικύλλεις R, Suid. : \lecon{τί κοικύλλεις} Faber}
\ligneRav{853}{}{πικρὰν \difference{ἑλένην} ὄψει τάχ᾽, εἰ μὴ κοσμίως}{\textbf{853} Ἑλένη R : \lecon{ἑλένην} Bibi ap. Hesych. ε 1995}
\ligneRav{}{}{ἕξεις, ἕως ἂν τῶν πρυτάνεών τις φανῇ·}{}
\ligneRav{855}{μ\auDessus{ν}{Η}}{Νείλου μὲν αἵδε καλλιπάρθενοι ῥοαί,}{}
\ligneRav{}{}{ὃς ἀντὶ δίας ψακάδος Αἰγύπτου πέδον}{}
\ligneRav{}{}{λευκῆς νοτίζει μελανοσυρμαῖον λεών·}{}
\ligneRav{}{\auDessus{γυ}{ν}}{πανοῦργος εἶ νὴ τὴν Ἑκάτην τὴν φωσφόρον·}{}
\ligneRav{}{μ\auDessus{ν}{Η}}{ἐμοὶ δὲ γῆ μὲν πατρὶς οὐκ ἀνώνυμος}{}
\ligneRav{860}{—}{ὅστις γε τολμᾷς σῆμα \difference{σοί} (σ\auDessus{ί}{ύ} \Rdeux) κλέπτειν γυνή·}{\textbf{860} σύ (ex σί) \Rdeux : \lecon{σοί} Zanetti (\italique{tibi} Divus) \espApp τ' R : \lecon{γ᾽} Gelenius}
\ligneRav{861}{}{πατὴρ ἐκεῖνός \difference{ἐστι}; Φρυνώνδας μὲν οὖν·}{\textbf{861} ἐστιν R : \lecon{ἐστι} \Mon}
\ligneRav{862}{μ\auDessus{ν}{Η}}{Ἑλένη δ᾽ ἐκλήθην. : αὖθις αὖ \difference{γίγνει} γυνή,}{\textbf{862} γίγνη R : \lecon{γίγνει} Brunck}
\ligneRav{863}{}{πρὶν τῆς ἑτέρας δοῦναι γυναικίσεως δίκην;}{\textbf{863} \lecon{δοῦναι post ἑτέρας} R : post δίκην Suid. γ 496}
\ligneRav{864}{μ\auDessus{ν}{Η}}{ψυχαὶ δὲ πολλαὶ δι᾽ ἔμ᾽ ἐπὶ \difference{Σκαμανδρίοις}}{\textbf{864} Σκαμανδρίαις R : \lecon{Σκαμανδρίοις} Fritzsche ex Eur.\italique{Hel} 52 et 609}
\ligneRav{865}{—}{ῥοαῖσιν ἔθανον. : ὤφελες δὲ καὶ σύ γε·}{}
\ligneRav{}{μ\auDessus{ν}{Η}}{κἀγὼ μὲν ἐνθάδ᾽ εἴμ᾽· ὁ δ᾽ ἄθλιος πόσις}{}
\ligneRav{867}{}{οὑμὸς \difference{Μενέλεως} οὐδέπω προσέρχεται·}{\textbf{867} Μενέλαος R : \lecon{Μενέλεως} Brunck, cf. v.901}
\ligneRav{868}{—}{τί οὖν ἔτι ζῶ; : τῶν κοράκων πονηρίᾳ·}{\textbf{868} τῶν κορ. πον. continuat (Μνη.) R : \lecon{(Γυνή)} Σ, Ellebodius}
\ligneRav{}{μ\auDessus{ν}{Η}}{ἀλλ᾽ ὥσπερ αἰκάλλει τι καρδίαν ἐμήν·}{}
\ligneRav{870}{}{μὴ ψεῦσον ὦ Ζεῦ τῆς ἐπιούσης ἐλπίδος·}{\textbf{870} \lecon{ἐλπίδος} R : ἡμέρας ap. Greg. Cor. p.14-15, Σ ad Eur. \italique{Hec}.225 \espApp post v. 870 < οἶσθ᾽ οὖν ὃ δρᾶσον; $ - \cup - | \times - \cup - $ > \lecon{om.} R : restituit ap. Σ Bibi (cf. Eur. \italique{Hel.}325)}
\ligneRav{871}{\auDessus{εὐ}{ρ}}{τίς τῶνδ᾽ ἐρυμνῶν δωμάτων ἔχει κράτος,}{\textbf{871} ὡς Περσεύς in mg. R : \lecon{ὡς Μενέλαος} Zanetti}
\ligneRav{872}{}{ὅστις \difference{ξένους} δέξαιτο ποντίῳ σάλῳ}{\textbf{872} ξένος R : \lecon{ξένους} Zanetti}
\ligneRav{873}{}{κάμνοντας ἐν χειμῶνι καὶ ναυαγίαις;}{\textbf{873} \lecon{κάμνοντας} R : καμόντας Lenting}
\ligneRav{}{μ\auDessus{ν}{Η}}{Πρωτέως τάδ᾽ ἐστὶ μέλαθρα. : ποίου Πρωτέως;}{}
\ligneRav{875}{}{\difference{ὤ τρισκακοδαίμων}, ψεύδεται νὴ τὼ θεώ,}{\textbf{875} ὦ τρισκακόδαιμον R : \lecon{ὤ τρισκακοδαίμων} Wilamowitz ad \italique{Lys.}449}
}
\vspace{3pt}\noindent\rule{4.5cm}{0.4pt}\par\vspace{1.5pt}
{\fontsize{7.0pt}{8.3pt}\selectfont
\noindent\textbf{Loci paralleli (vv. 838–875) ---} \textbf{838} (σκάφιον ἀποκεκαρμένην) Photius σ 291 σκάφιον ἀποκεκαρμένην ; cf. Hesychius σ 879 σκάφιον ; Photius σ 292 ($\simeq$ Scholium) ; \Eustathe 907.40, 1292.61 ; \AelDion σ 23 \espTest \textbf{846} Suidas ι 323 $\sim$ Zonaras 1103 ἰλλός \espTest (ἰλλός) Pollux II 52 ; \Eustathe 768.9 ; \AelDion ι 6 ; cf. Moeris ι 21 (ἰλλόν $=$ στραβόν) ; Hesychius ι 578 ἰλλός ; Hesychius η 257 (ἠίλλοντο $=$ ἐστράβιζον) \espTest \textbf{852} Suidas κ 2534 \espTest (κοικύλλεις) Hesychius κ 3232 κοικύλλειν $\sim$ Photius κ 860 \espTest \textbf{861} (Φρυνώνδας) Suidas φ 700 ; \Syn φ 211 Φρυνώνδας $\sim$ Hesychius φ 938 \espTest \textbf{863} Suidas γ 496 γυναίκισις et Zonaras 459 πρὶν τῆς ἑτέρας γυναικίσεως δίκην δοῦναι \espTest \textbf{868} Suidas τ 648 \espTest \textbf{869} (αἰκάλλει τι καρδίαν ἐμήν) Suidas αι 161 ; Zonaras 91 (τι omisso) \espTest \textbf{870} \GregCor p.14-15 ($=$ scholium\expo{Mosch.} Euripides, \italique{Hecub.} 225) (sed ἡμέρας pro ἐλπίδος) ; Suidas ψ 62 \espTest (μὴ ψεῦσον) \Antiatt 107.31 \espTest \textbf{870} \GregCor pp. 14-15 ($=$ scholium\expo{Mosch.} Euripides, \italique{Hecub.} 225) (sed ἡμέρας pro ἐλπίδος) ; Suidas ψ 62 \espTest (μὴ ψεῦσον) \Antiatt 107.31\par
}
\vspace{2.5pt}\noindent\rule{4.5cm}{0.4pt}\par\vspace{1.5pt}
{\fontsize{6.9pt}{8.2pt}\selectfont
\noindent\textbf{Suite (en commentaire) de l'apparat critique \& notes paléographiques (\texttt{f. 172v}) ---} \textbf{Corrections supralinéaires de $\text{R}^{\text{pc}} / \text{R}^2$ :} Au \textbf{v. 838}, \lecon{ὑστέρ\auDessus{ο}{α}ν} (\Rpc{} trace $\alpha$ au-dessus de $\text{ο}$ ; $\Sigma^{\text{R}}$ lisait ὑστάτην) ; au \textbf{v. 844}, \lecon{τ\auDessus{ὰ}{ο}δὶ} (\Rpc{}, τοδί \Mon{}) ; au \textbf{v. 860}, \lecon{σύ} corrigé par \Rdeux{} à partir de σί (pour σοί Zanetti). --- \textbf{Deux conjectures capitales de \textbf{Bibi} (vv. 839 et 853) :} --- \textbf{v. 838 :} \lecon{ὑστέρ\auDessus{ο}{α}ν} \Rpc{} : ὑστάτην $\Sigma$. --- \textbf{v. 839 :} ὦ πόλεις R (ὦ πόλις Gelenius) : \lecon{ωὖπολι} \lecon{Bibi}. --- \textbf{v. 853 :} Ἑλένη R : \lecon{ἑλένην} \lecon{Bibi} ap. Hésychius $\varepsilon\ 1995$. --- \textbf{v. 868 :} τῶν κορ. πον. continuat (Μνη.) R : \lecon{(Γυνή)} $\Sigma$, Ellebodius.\par
}
\end{minipage}
\end{adjustbox}
\end{adjustwidth}
\newpage
\enTeteD{Codex Ravennas 429 (Biblioteca Classense) --- Fac-similé du f. 172v (vv. 838–875)}
\begin{adjustwidth}{-1.5cm}{-3.0cm}
\begin{center}
\includegraphics[width=16.5cm,height=20.5cm,keepaspectratio]{/workspaces/Thesmophories/dossier_ravennas_thesmophories/folios_individuels_ff162v-177r/21_f172v_classense174v_vv0838-0875.jpg}\par
\vspace{2pt}
{\scriptsize\textbf{Codex Ravennas 429, f. 172v (Classense 174v)} --- ΠΑΡΑΒΑΣΙΣ / ΕΠΕΙΣΟΔΙΟΝ Δ', \textbf{vv. 838–875} ($\text{R}^1, \Rpc, \Rdeux$, \PsiUn).}
\end{center}
\end{adjustwidth}
\newpage
\clearpage\ifodd\value{page}\hbox{}\newpage\fi
\enTeteG{ΘΕΣΜΟΦΟΡΙΑΖΟΥΣΑΙ --- Texte grec établi (f. 172v, vv. 838–875)}
\enTeteD{LES THESMOPHORIES --- Traduction française (f. 172v, vv. 838–875)}
\begin{pages}
\begin{Leftside}
\beginnumbering
%% v. 838+ (part 2)
\begin{groupe}
\tetramTroch{\ligne{838}\Automatique{ὑστέραν αὐτὴν καθῆσθαι σκάφιον ἀποκεκαρμένην}\testimonia{838}{(σκάφιον ἀποκεκαρμένην) Photius σ 291 σκάφιον ἀποκεκαρμένην ; cf. Hesychius σ 879 σκάφιον ; Photius σ 292 ($\simeq$ Scholium) ; \Eustathe 907.40, 1292.61 ; \AelDion σ 23}\apparat{838}{\lecon{ὑστέρ\auDessus{ο}{α}ν} \Rpc : ὑστάτην Σ \ligneApp{839} ὦ πόλεις R : ὦ πόλις Gelenius : \lecon{ωὖπολι} Bibi}%
\commentaire{838}{σκάφιον ἀποκεκαρμένην.}{C'est l'humiliation que le tyran de Cumes, Aristodème, avait fait subir aux filles de son pays : περιτρόχαλα κείρεσθαι, tandis qu'il forçait les garçons à porter des cheveux longs. Plutarque, \italique{De mul. virtut}, 26, p. 361f. (Note de Willems.)\\
La scholie précise à propos de σκάφιον : %
\citComm{εἶδος κουρᾶς δουλικῆς.\\
\francais{Genre de coupe rase qui convient à l'esclave.}}%
}\\
τῆς τὸν ἀνδρεῖον τεκούσης· τῷ γὰρ εἰκὸς ωὖπολι%\apparat{839}{ὦ πόλεις R : ὦ πόλις Gelenius : \lecon{ωὖπολι} Bibi}%
\commentaire{839}{ὦ πόλεις.}{Dans son livre \italique{Aristophanes und Eupolis} (2007), Natalia Kyriakidi a montré qu'Eupolis, le poète comique rival d'Aristophane, est selon toute vraisemblance mort en 411 lors de la victoire navale de Κυνὸς Σῆμα (cf. comm. vv. 833-7). Pour arriver à cette conclusion, elle récuse très justement la plupart des témoignages, qui exhalent l'anecdote recyclée, pour ne conserver que le témoignage de Suidas (ε 3657) -- ou de la \italique{Souda}, si l'on préfère :%
\citComm{Εὔπολις, Σωσίπολις, Ἀθηναῖος, κωμικός. ὃς ἑπτακαίδεκα ἔτη γεγονὼς ἤρξατο ἐπιδείκνυσθαι: καὶ ἐδίδαξε μὲν δράματα ιζ, ἐνίκησε δὲ ζ. καὶ ἀπέθανε ναυαγήσας κατὰ τὸν Ἑλλήσποντον ἐν τῷ πρὸς Λακεδαιμονίους πολέμῳ: καὶ ἐκ τούτου ἐκωλύθη στρατεύεσθαι ποιητήν. %τὰ δὲ δράματα αὐτοῦ Αἶγες, Ἀστράτευτος ἢ Ἀνδρόγυνοι, καὶ ἄλλα.
\\
Fils de Sosipolis, Athénien, un poète comique. Il a commencé à monter des pièces de théâtre à l'âge de dix-sept ans. Il a produit 17 pièces et en a gagné 7. Il est mort dans le naufrage de son navire dans l'Hellespont pendant la guerre contre les Lacédémoniens ; à la suite de cela, il fut interdit à un poète de servir en campagne. %Ses pièces sont les \italique{Chèvres}, l'\italique{Exempté} ou les \italique{Androgynes}, et d'autres.
}%
Natalia continue (p. 11) : \allemand{Die zwei großen Schlachten, die am Hellespont während des Peloponnesischen Krieges stattfanden, sind der Sieg der Athener bei} Κυνὸς Σῆμα \allemand{411 und die endgültige Niederlage bei }Αἰγὸς Ποταμοί \allemand{405. Der Tod des Eupolis während einer Seeschlacht könnte auch die Entstehung der Geschichte vom Ertränken des Eupolis durch Alkibiades erklären} [$=$ Eupolis, selon certaines légendes, aurait été noyé en mer par Alcibiade au retour de Sicile], \allemand{weil beide Geschichten als Vorfälle im Meer bei einer Kriegshandlung wenige Jahre vor dem Ende des Peloponnesischen Krieges überliefert sind. Storey vermutet, dass Ailians aitiologische Erklärung des Ortsnamens} Κυνὸς Θρῆνος [$=$ le lieu commun du chien qui meurt de chagrin sur la tombe de son maître] \allemand{auf }Κυνὸς Σῆμα \allemand{zurückzuführen ist. Da kein Werk des Eupolis nach 411 zu datieren ist und für die Zeit danach keine Indizien für eine weitere Tätigkeit des Dichters überliefert sind, erscheint die Information der Suda als das einzige zuverlässige Zeugnis über den Tod des Eupolis. So enden sein Leben und seine dichterische Karriere wahrscheinlich im Sommer 411 während der Seeschlacht bei} Κυνὸς Σῆμα.\\
Le texte du manuscrit de \italique{Ravenne} comporte au vers 839 une adresse à la Cité :%
\citComm{τῷ γὰρ εἰκὸς ὦ πόλεις (sic.)}%
% "Est-ce à raison, ô ville, etc." dont l'emphase conviendrait parfaitement à un pastiche de tragédie, mais qui ne paraît point à sa place ici. 
La vraisemblable quasi-coïncidence des dates entre la mort d'Eupolis et la représentation des \italique{Thesmophories}\marqueur{1} m'incite à penser qu'il y a là plutôt une évocation d'Eupolis, auteur intarissable sur la mère d'Hyperbolos, et à laquelle il a même consacré une pièce entière : son \italique{Maricas} (voir le commentaire du vers 840).
S'il faut attribuer l'idée d'accorder les sièges des Athéniennes avec la valeur de leur progéniture à feu Eupolis (et peut-être aussi celle de les tondre le cas décevant échéant), il n'aura pas manqué de s'indigner lui aussi que la mère d'Hyperbolos spécule avec le fruit de la rapine de son fils. 
%Cette hypothèse se peut-elle prouver ? 
On ne possède malheureusement que peu de choses d'Eupolis. % et du \italique{Maricas} en particulier. 
On sait que le \italique{Maricas} dut être représenté en $-421$, que le chœur était formé de deux demi-chœurs : celui des riches et celui des pauvres citoyens% -- qui s'opposent en cela aux \italique{Cavaliers} d'Aristophane ? --
, chœurs qui débattaient avec le démagogue Hyperbolos.
La comédie est une violente charge contre ce dernier.
Elle se terminait par une danse licencieuse de la mère du démagogue.
Nombre de commentateurs, à l'instar des scholiastes, ont rapproché les \italique{Cavaliers} du \italique{Maricas}... mais pas les \italique{Thesmophories} -- plus pauvres en scholies anciennes.
Pourtant, certaines idées d'Eupolis semblent en être reprises. Nous l'allons voir tout-à-l'heure.
Mettons en garde tout-de-suite : les hypothèses -- comme les explications -- seront fort spéculatives.\\
Un commentaire du \italique{Maricas}, daté du \nieme{II} siècle ap. J.-C., nous est parvenu de façon très morcelaire. On y lit cette citation (fr. 192.96-99 Kassel-Austin) :%
\citComm{τοῦτ᾽ ἐκδανείζει καὶ κυκᾷς [τὸ]\douteux{ν} ναυτικὸν \douteux{ἄ}[ργυρον \ \ \ \ \ (\textbf{ἄργυρον} Bibi)\\
ἐπὶ τῷ πέμπτῳ μέρε[ι \ \ \ \ \ \ \ ]τοὺς τόκ[ους οἱ\\
ναυτικοί, ἡμεῖς δ᾽ ἄρ᾽ οἴκ[ \ \ \ \ \ \ \ \ \ \ \ ]εν· ὁ τῶν\\
πλουσίων λέγει χορός.}%
%(J'interprète sans véritable boussole et navigue à vue ; des erreurs se glissent certainement : lecteurs, soyez attentifs !)
Le chœur des riches (ὁ τῶν πλουσίων χορός) reproche à Maricas ($=$ Hyperbolos) de prêter (ἐκδανείζει) en concoctant (κυκᾷς) des prêts (τὸν ναυτικὸν ἄργυρον) au cinquième (ἐπὶ τῷ πέμπτῳ μέρει\marqueur{2}), c'est-à-dire avec un taux d'intérêt de $20 \%$.
Il est questions ensuite des τόκους, `enfants' ($=$ intérêts) de ces ναυτικοί (cf. vers 845).
Continuant le jeu de mot nautique (les prêts sont littéralement de l'argent \italique{marin}), est-ce à ce moment qu'Eupolis, par la voix du chœur, proposait de faire tondre \italique{nautiquement} cet usurier ou sa mère et de lui faire porter une \italique{barque} (σκάφιον vers 838) ?\\
Revenons aux \italique{Thesmophories} : en guise d'oraison funèbre, Aristophane rend hommage à la pertinence des propositions d'Eupolis :%
%Ce \italique{Maricas} remonte à la période où Lamachos était encore perçu par les comiques comme un bravache. Eupolis aura pu mettre dans le même panier la mère de Lamachos avec celle d'Hyperbolos. Rétrospectivement, au vu des derniers événements, voilà qui semble bien peu pertinent.
%S'il faut attribuer l'idée d'accorder les sièges des Athéniennes (et peut-être aussi celle de les tondre) avec la valeur de leur progéniture à feu Eupolis, 
%Une preuve supplémentaire -- et donc, nous espérons, définitive -- de la quasi-coïncidence de la mort d'Eupolis et de la représentation des \italique{Thesmophories} est l'adresse d'Aristop  :% dont l'oraison funèbre est réduite à la portion congrue
\citComm{τῷ γὰρ εἰκὸς\marqueur{3} ωὖπολι\\
Comment tolérer\marqueur{3}, ô Eupolis, que \&c.}%
%Aristote multiples exemples : τῷ γὰρ τὸ πρᾶγμα εἶναι ἢ μὴ εἶναι
%(L'expression `τῷ γὰρ εἰκὸς' se trouve telle quelle aussi au vers 703 des \italique{Acharniens}.)
Dernier coup de chapeau d'Aristophane à son défunt rival, mort des suites de la folie des Athéniens et de l'incompétence des amiraux.%
}%
\debutNote{839}%
\note{1}{Voir l'introduction et la question de la date. Avec Rogers, je penche pour l'hypothèse que les \italique{Thesmophories} furent représentées aux Dionysies de 410. Si nous avons vu juste avec Natalia, Eupolis serait donc mort l'été précédent la première représentation de la pièce.}%
\note{2}{Sens bien attesté d'ἐπί. Cf. Dém. 1250, 20 : ἐπ᾽ ὀκτὼ ὀβολοῖς δανείζειν τοῦ μηνὸς τὴν μνᾶν, «~prêter de l'argent moyennant un intérêt de huit oboles la mine par mois~», c'est-à-dire à $16 \%$.}%
\note{3}{Le tour elliptique apparaît aussi au vers 703 des \italique{Acharniens} où il est glosé `πῶς δίκαιόν ἐστιν' par une scholie. Le ton est effectivement indigné : «~comment peut-il être juste/raisonnable que...~», c'est-à-dire «~quelle indignité que...~». Voir aussi Sophocle (\italique{Ajax}, vers 1366-1367), non indigné mais ironique :%
\citComm{Ἀγαμέμνων -- ἡ πάνθ᾽ ὅμοια πᾶς ἀνὴρ αὑτῷ πονεῖ.\\
Ὀδυσσεύς -- τῷ γάρ με μᾶλλον εἰκὸς ἢ ᾽μαυτῷ πονεῖν;\\
Agamemnon -- En toute chose, chacun ne travaille que pour lui-même.\\
Ulysse -- Comment peut-il être plus raisonnable que de travailler pour moi-même ?}%
}}
\end{groupe}%
%% v. 830-845
\begin{groupe}%
%
\tetramTroch{τὴν Ὑπερβόλου καθῆσθαι μητέρ᾽ ἠμφιεσμένην%
\commentaire{840-1}{Ὑπερβόλου ... Λαμάχου,}{Hyperbolos, le démagogue athénien -- marchand de lampes de son état\marqueur{1} -- devint démagogue en chef après la mort de Cléon devant Amphipolis en 422 av. J.-C.\marqueur{2} ; sa mère et lui furent lors une cible privilégiée d'Aristophane et des poètes comiques\marqueur{3} (suite sans doute à quelque lapsus à relent freudien du démagogue, lapsus qui fut repris à l'encan par les Comiques). Cet ὀξίνην Ὑπερβόλον\marqueur{4} réunissait tout pour être détesté d'Aristophane : de caractère exécrable, il était un partisan, comme son mentor Cléon, de la guerre à outrance\marqueur{5} ; c'était un intenteur de procès\marqueur{6}, un fourbe\marqueur{7}, qui s'est enrichi par ses coquineries\marqueur{8}. Il eut l'insigne honneur d'être le dernier des Athéniens ostracisés en 417\marqueur{9}. Il fut assassiné à Samos, son lieu d'exil\marqueur{10}, et son corps jeté en mer\marqueur{11}... cette \italique{inhudation} lui valut une dernière raillerie de la part d'Aristophane\marqueur{12}. 
%\eff{Qu'est-ce qui valut à la mère d'Hyperbolos d'être raillée ainsi par les poètes comiques ? Il me semble qu'on peut inférer un début de réponse du fameux passage des \italique{Nuées} (vers 551-556) où Aristophane fait preuve d'une belle assurance de posséder sa Muse, pourtant encore bien neuve : %
%\citComm{οὗτοι δ’, ὡς ἅπαξ παρέδωκεν λαβὴν Ὑπέρβολος,\\
%τοῦτον δείλαιον κολετρῶσ’ ἀεὶ καὶ τὴν μητέρα.\\
%Εὔπολις μὲν τὸν Μαρικᾶν πρώτιστον παρείλκυσεν\\
%ἐκστρέψας τοὺς ἡμετέρους Ἱππέας κακὸς κακῶς,\\
%προσθεὶς αὐτῷ γραῦν μεθύσην τοῦ κόρδακος οὕνεχ’, ἣν\\
%Φρύνιχος πάλαι πεπόηχ’, ἣν τὸ κῆτος ἤσθιεν.\\
%εἶθ᾽ Ἕρμιππος αὖθις ἐποίησεν εἰς Ὑπέρβολον,\\
%ἄλλοι τ᾽ ἤδη πάντες ἐρείδουσιν εἰς Ὑπέρβολον,\\
%τὰς εἰκοὺς τῶν ἐγχέλεων τὰς ἐμὰς μιμούμενοι.\\
%ὅστις οὖν τούτοισι γελᾷ, τοῖς ἐμοῖς μὴ χαιρέτω:\\
%ἢν δ᾽ ἐμοὶ καὶ τοῖσιν ἐμοῖς εὐφραίνησθ᾽ εὑρήμασιν,\\
%εἰς τὰς ὥρας τὰς ἑτέρας εὖ φρονεῖν δοκήσετε.\\
%Mes rivaux, depuis qu'une fois Hyperbolos leur a donné prise, ne cessent de piétiner le malheureux et sa mère. Eupolis, tout le premier, traîna sur la scène son \italique{Maricas}, où, piètre qu'il est, il retourna piètrement nos \italique{Cavaliers}, en y ajoutant en vue de la cordace une vieille drôlesse, celle-là même qu'imagina jadis Phrvnichos et que le monstre marin voulait manger. Puis, à son tour, Hermippos prit à partie Hyperbolos, et tous les autres de se buter contre Hyperbolos, en m'empruntant ma comparaison des anguilles. Que celui donc qui rit à leurs pièces ne se plaise pas aux miennes ! Mais si vous vous laissez égayer par moi et mes inventions, dans les temps à venir vous passerez pour sensés}%
%La clef est dans "παρέδωκεν λαβὴν Ὑπέρβολος"}\\
\textbf{Lamachos} était un stratège athénien\marqueur{13} et l'un des représentants athéniens de la Paix de Nicias\marqueur{14}. Avec Nicias et Alcibiabe, il est investi des pleins pouvoirs lors de l'expédition de Sicile\marqueur{15} où il trouve la mort\marqueur{16}. Lamachos présente une originalité unique parmi les personnages historiques d'Aristophane puisque le poète opère un volte-face à son endroit : du va-t-en-guerre au panache ridicule dépeint dans les \italique{Acharniens} et la \italique{Paix}\marqueur{17} (accusé même de s'enrichir grâce à la guerre\marqueur{18}), Aristophane lui rend hommage ici comme dans les \italique{Grenouilles}\marqueur{19}. La mort héroïque de Lamachos\marqueur{20} est sans doute pour beaucoup dans ce changement radical d'opinion.}
\debutNote{840-1}%
\note{1}{\italique{Cavaliers}, v. 739 (non nommé) ; \italique{Nuées}, v. 1064 ; \italique{Paix}, v. 690.}%
\note{2}{Voir la \italique{Paix}, v. 679-692, Paix qui fout le camp, révoltée que le peuple se soit donné un si indigne patron.}%
\note{3}{Cf. les \italique{Nuées}, v. 551-562 : Aristophane daube sur ses plagiaires de contemporains, notamment Eupolis réchauffant leurs \italique{Cavaliers} pour son \italique{Maricas} et y rajoutant un autre plat déjà concocté par Phrynicos en son temps : une mère d'Hyperbolos, ivrognesse et dansant le licencieux cordax. Sur les attaques des comiques contre la mère d'Hyperbolos, présentée comme une débauchée notoire, voir aussi les fragments 7-12 des \italique{Vendeuses de pain} (αἱ Ἀρτοπώλιδες) d'Hermippos (dans l'édition Kassel-Austin) ainsi que les fragments de Platon le Comique dans sa pièce éponyme \italique{Hyperbolos}.}%
\note{4}{\italique{Cavaliers}, v. 1303. Thucydide, qui n'estimait guère plus ce μοχθὴρον ἄνθρωπον, ne parle qu'une seule fois d'Hyperbolos : pour nous annoncer sa mort et en brosser un portrait expéditif (VIII, 73.3) : %
\citComm{καὶ Ὑπέρβολόν τέ τινα τῶν Ἀθηναίων, μοχθηρὸν ἄνθρωπον, ὠστρακισμένον οὐ διὰ δυνάμεως καὶ ἀξιώματος φόβον, ἀλλὰ διὰ πονηρίαν καὶ αἰσχύνην τῆς πόλεως, ἀποκτείνουσι μετὰ Χαρμίνου\\
S'en prenant à Hyperbolos, un Athénien méprisable qu'on avait ostracisé non par peur de son influence et de son prestige, mais parce que c'était un malhonnête homme qui déshonorait la cité, ils [les conjurés favorables à l'oligarchie] l'assassinèrent avec l'aide de Charminos.}%
}%
\note{5}{\italique{Cavaliers}, v. 1303-1315 ; il cherchait à renforcer la flotte pour que l'expédition de Sicile fût poussée jusqu'à Carthage.}%
\note{6}{\italique{Acharniens}, v. 846-847.}%
\note{7}{\italique{Guêpes}, v. 1007.}%
\note{8}{\italique{Nuées}, v. 1064-1066. Tel fils, telle mère ? L'usure brocardée par Aristophane dont fait preuve la mère d'Hyperbolos n'était-elle pas la continuation des activités louches du fils, depuis que ce dernier se trouvait en exil à Samos ?}%
\note{9}{Plut., \italique{Nicias} 11 ; \italique{Alcib.} 13.}%
\note{10}{%\italique{Après} la représentation des \italique{Thesmophories}. Willems avait déjà remarqué que l'attitude de la mère d'Hyperbolos, prenant part en costume d'apparat aux fêtes célébrées par les femmes, était incompatible avec le deuil d'une mère éplorée. Ce détail confirme, ce nous semble, la date de 411 pour la première représentation des \italique{Thesmophories} et infirme celle de 410 soutenue notamment par Rogers.
Thucydide, VIII, 73.3.}%
\note{11}{Théopompe 115F69, in \italique{Die Fragmente der griechischen Historiker}, Teil 2, \italique{Zeitgeschichte - B}.}%
\note{12}{\italique{Grenouilles}, v. 570.}%
\note{13}{Thucydide IV, 75. 1 : il commande une dizaine de navires dans le Pont.}%
\note{14}{Thucydide V 19.2 et 24.1}%
\note{15}{Thucydide (VI, 8) ; il s'opposa à Alcibiade sur la stratégie (V 49) : Lamachos voulait aller droit au but et attaquer Syracuse sitôt le débarquement fait ; il finit par se ranger à l'opinion d'Alcibiade (VI, 50).}%
\note{16}{Thucydide (VI, 101) : Lamachos courant à la rescousse de l'aile droite de l'armée athénienne se retrouva isolé avec quelques-uns de ses camarades ; il fut tué par les Syracusains qui récupérèrent son corps avant de le restituer (VI, 103) aux Athéniens. Nicias fut dès lors le seul commandant en Sicile.}%
\note{17}{Lamachos arborait sur son casque trois aigrettes qui lui donnaient des faux airs de Gorgone (\italique{Acharniens}, v. 574 ; v. 964-965). Aristophane rapproche cet accoutrement d'épouvante à d'autre monstre-croquemitaine au nom voisin (Μορμώ, \italique{Acharniens} v. 582 ; \italique{Paix} v. 474). Enfin, cette crinière en bataille -- de circonstance pour un militaire ! -- permettait un jeu analogue en grec entre λόφοι et λόχοι, trop tentant pour qu'Aristophane n'y cèdât pas (\italique{Acharniens}, v. 567, 574, 1074).}%
\note{18}{Il est traité de μισθαρχίδης au v. 597 des \italique{Acharniens} ; attaque renouvelée au vers 619.}%
\note{19}{\italique{Grenouilles}, v. 1039.}%
\note{20}{Voir -- outre le passage de Thucydide (VI, 101) -- Plutarque, \italique{Nicias} 18 (et aussi le \italique{Lachès} (19c) de Platon où il est question du courage de Lamachos.)}\\
\Automatique{λευκὰ καὶ κόμας καθεῖσαν πλησίον τῆς Λαμάχου·}\\
\Automatique{καὶ δανείζειν χρήμαθ᾽, ἥν χρῆν, εἰ δανείσειέν τινι}\apparat{842}{χρήματα τ' R : %χρήματ' \Mon : 
\lecon{χρήμαθ᾽} Scaliger \espApp \lecon{ἥν χρῆν} R : ᾗ χρῆν Scaliger %(\italique{cui} Divus) 
\espApp χρῆν ἦν R : \lecon{χρῆν εἰ} Brunck \espApp δανείσειε R : \lecon{δανείσειέν} Biset \ligneApp{844} ἀφερεῖσθαι R : \lecon{ἀφαιρεῖσθαι} \Mon \espApp χρήματατ᾽ R : \lecon{χρήματ᾽} Grynaeus \espApp τ\auDessus{ὰ}{ο}δὶ \Rpc : \lecon{τοδί} \Mon}%
%\commentaire{842}{ᾗ χρῆν}{\eff{Imprécation ? (Austin)}}
\\
\Automatique{καὶ τόκον πράττοιτο, διδόναι μηδέν᾽ ἀνθρώπων τόκον,}\\
ἀλλ᾽ ἀφαιρεῖσθαι βίᾳ τὰ χρήματ᾽ εἰπόντας τοδί,\\
`ἀξία γοῦν εἶ τόκου τεκοῦσα τοιοῦτον τόκον.'
\commentaire{845}{ἀξία γοῦν εἶ τόκου τεκοῦσα τοιοῦτον τόκον.}{Jeu de mots qui roule sur le double sens de τόκος, \latin{usura} et \latin{proles}. C'est ainsi que l'orateur Eschine disait de Démosthène qui, souffleté par Midias, s'était désisté moyennant 30 mines : %
\citComm{ὁ γὰρ ἄνθρωπος οὐ κεφαλήν, ἀλλὰ πρόσοδον κέκτηται.\\
Ce n'est pas une tête (κεφαλήν), mais un capital que possède cet homme.}%
(Cf. \italique{Contre Ctésiphon}, 212 ; note de Willems.)}\\\passeLigne}%
\end{groupe}%
%% v. 846-851
\begin{groupe}%
%
\personnageGrec{Μνησίλοχος}%
\typeVers{trim. iamb.}\trimIamb{ἰλλὸς γεγένημαι προσδοκῶν· ὁ δ᾽ οὐδέπω.\testimonia{846}{Suidas ι 323 $\sim$ Zonaras 1103 ἰλλός \espTest (ἰλλός) Pollux II 52 ; \Eustathe 768.9 ; \AelDion ι 6 ; cf. Moeris ι 21 (ἰλλόν $=$ στραβόν) ; Hesychius ι 578 ἰλλός ; Hesychius η 257 (ἠίλλοντο $=$ ἐστράβιζον)}\\
τί δῆτ᾽ ἂν εἴη τοὐμποδών; οὐκ ἔσθ᾽ ὅπως\\
οὐ τὸν Παλαμήδη ψυχρὸν ὄντ᾽ αἰσχύνεται.\apparat{848}{Παλαμήδην R : \lecon{Παλαμήδη} Elmsley (\italique{Quart. Rev.} 7 (1812) p.454)}\\
τῷ δῆτ᾽ ἂν αὐτὸν προσαγαγοίμην δράματι;\\
ἐγᾦδα· τὴν καινὴν Ἑλένην μιμήσομαι.
\commentaire{850}{τὴν καινὴν Ἑλένην.}{« La nouvelle \italique{Hélène} ». \italique{Nouvelle} parce qu'Euripide l'avait présentée il y a peu de temps, sous l'archontat de Cléocrite (printemps 412). Mais l'adjectif καινή est aussi à prendre au sens d'\italique{innovante}, \italique{imprévue}, \italique{étrange}. Voici ce que dit Victor-Henry Debidour de cette « \italique{Hélène} nouveau style »\marqueur{1} :\\
Dans la tradition légendaire, Hélène de Sparte, femme de Ménélas, était le parangon de la beauté. Mais aussi le plus scandaleux exemple d'infidélité féminine, cause, par son inconduite, des malheurs sans nombre de la guerre de Troie. Le personnage s'était stéréotypé en « femme fatale », d'autant plus haïssable qu'elle était plus prestigieuse ; et \italique{Euripide lui-même ne la traite pas autrement} -- sauf justement dans la tragédie très peu tragique qui porte son nom. Certes, elle y reste, bien-sûr, la belle Hélène, et si charmeuse que le fils du sage Protée l'aime comme un fou, et comme un sot. Mais la voici transformée en épouse sans tache, victime innocente de la longue perfidie des deux déesses, dont l'une l'a prostituée et l'autre séquestrée : elles l'ont engagée dans une rocambolesque histoire où elle mérite toute notre sympathie. Pour se faire désigner par Pâris comme « la plus belle », Cypris avait promis au prince-berger troyen de lui donner Hélène ; Héra, furieuse, contre sa rivale et son juge, fit en sorte que Pâris n'emmenât à Troie qu'un fantôme d'Hélène, une « nuée » toute semblable à la vraie et suffisamment consistante pour leurrer le ravisseur pendant les dix ans qu'il la garda, et le légitime époux après ce qu'il crut être sa reconquête. La vraie Hélène, elle, miraculeusement transportée en Égypte, y était tenue cachée durant toutes ces années, confiée à la vertueuse protection de Protée. Ainsi, c'est pour un « crime sans criminelle » que des milliers de preux se sont entre-tués et que Troie a été anéantie.\\
Telle est l'\italique{étonnante} variante que recueille Euripide sur la plus célèbre entreprise guerrière de la légende grecque. Car il ne l'a pas inventée : Hérodote en fait état bien avant lui\marqueur{2}. De ce point de départ, Euripide a tiré un canevas sur lequel il a brodé une aventure romanesque et humoristique qui souvent ferait croire qu'avant d'être parodiée elle-même par Aristophane, l'\italique{Hélène} donnait dans la parodie.}
\debutNote{850}
\note{1}{Dans l'introduction à sa traduction de l'\italique{Hélène}. Oui ! car Debidour le grand traducteur d'Aristophane, qui a su rendre un peu de la verve de notre poète dans la langue de Rabelais, Debidour ne s'est pas limité à la Comédie Ancienne ; il est aussi l'auteur d'une traduction des pièces conservées des trois grands Tragiques : \italique{Les Tragiques -- Eschyle, Sophocle, Euripide, théâtre complet}, La Pochothèque (1999). Traduction à laquelle s'ajoute l'excellent dossier et les index fort utiles de Paul Demont et Anne Lebeau. Victor-Henry Debidour vérifie ainsi le paradoxe de Socrate dont Platon n'a pas cru utile de nous donner la démonstration à la fin de son \italique{Banquet}.}
\note{2}{Hérodote II, 112-120. Par la suite, du reste, le thème de la réhabilitation d'Hélène fut maintes fois repris en exercice d'école : la paradoxe devient à son tour lieu commun.}\\
πάντως δ᾽ ὑπάρχει μοι γυναικεία στολή.\apparat{851}{πάντως R : \lecon{πάντως δ᾽} Bentley, cf. v.1012}\\\passeLigne}%
\end{groupe}%
%% v. 852-854
\begin{groupe}%
\personnageGrec{Κρίτυλλα}%
\trimIamb{τί αὖ σὺ κυρκανᾷς; τί κοικύλλεις ἔχων;\testimonia{852}{Suidas κ 2534 \espTest (κοικύλλεις) Hesychius κ 3232 κοικύλλειν $\sim$ Photius κ 860}\apparat{852}{ἢ τί κοικύλλεις R, Suid. : \lecon{τί κοικύλλεις} Faber \ligneApp{853} Ἑλένη R : %Ἑλένην Bonini : 
\lecon{ἑλένην} Bibi ap. Hesych. ε 1995}%
\commentaire{852}{τί κοικύλλεις ἔχων;}{Κοικύλλειν est un verbe à redoublement expressif\marqueur{1} qui tire son origine de τὰ κύλα, \italique{les poches sous les yeux}, d'où le sens d'« écarquiller les yeux »\marqueur{2}. Voici pour le sens en lequel l'emploie Critylla.
Le verbe est particulièrement bien adapté à la situation : notre anti-héros vient de lancer `à la mer' ses `rames' dans l'espoir de rameuter Euripide ; il est là maintenant, le regard hagard, devant cette mer imaginaire, à attendre... pareil à ce personnage comique, surnommé Κοικυλίων, type de l'idiot qui comptait les vagues de la mer\marqueur{3}, le comble de l'inutilité pour un grec : κύματα μετρεῖν était même passé en proverbe (détourné par Aristophane au vers 1257 de la \italique{Paix}).}%
\debutNote{852}%
\note{1}{Sur le même modèle que μοιμύλλειν, δενδίλλειν, δανδαίνειν, παπταίνειν, \&c.}%
\note{2}{Cf. Hésychius κ 3232 : κοικύλλειν· περιβλέπειν· κύλα γὰρ τὰ ἐπάνω τῶν ὀφθαλμῶν [$=$ la partie supérieure des paupières, mais voir aussi la glose d'Hésychius sous κύλα, où le terme désigne la partie \italique{inférieure} des paupières]. οἱ δὲ τὸ περιστρέφεσθαι καὶ διατρίβειν περιβλεπόμενον. ἄλλοι φθονεῖν.}%
\note{3}{Ælien, \italique{Varia Historia}, XIII, 15 ($=$ \latin{Adespota comica}, fr. 72 K.-A.) : 
\citComm{φασὶ παχύτατον γενέσθαι τὴν διάνοιαν οἱ τῆς κωμῳδίας ποιηταὶ Πολύδωρόν τινα, καὶ ἄλλον Κοικυλίωνα ὄνομα, ὅσπερ τὰ κύματα ἠρίθμει ὑπὸ τῆς ἄγαν ἀνοίας.\\
\francais{Les poètes comiques disent qu'un certain Polydore (aussi [connu sous] un autre nom : Cœcylion) affolit tant qu'il comptait les vagues sous le coup de son trop plein de folie.}}%
Ce passage, déniché par Fritzsche (à l'érudition et la bibliothèque toujours impressionnantes !), est à rapprocher du chapitre du Περὶ βλασφημιῶν de Suétone (édité, pour ce qu'il a pu en reconstituer, en 1967 aux Belles Lettres par le grand Taillardat) consacré à quelques idiots célèbres. Περὶ βλασφημιῶν, β 189 $=$ Eustathe \italique{Comm. Od.} p. 1669.53 : ὥσπερ δὲ τουτούς ἡ ἱστορία λόγου ἠξίωσεν ἐν οἷς καὶ τὸν τὰ κύματα μετροῦντα Πολύ[δ]ωρον.}\\
πικρὰν ἑλένην ὄψει τάχ᾽, εἰ μὴ κοσμίως
\commentaire{853}{πικρὰν ἑλένη.}{Le détournement du \italique{Palamède} ayant échoué (après celui du \italique{Télèphe}), Mnésiloque songe maintenant à recourir à la récente \italique{Hélène} d'Euripide. Seulement, de toutes les \italique{Thesmophories}, aucun des geôliers de Mnésiloque ne rentre jamais dans le jeu de ses parodies d'Euripide. %Il est étonnant que pas un des nombreux commentateurs d'Aristophane ne s'en soit avisé~: 
Quand Kritylla prononce `πικρὰν ἑλένην ὄψει τάχ᾽', elle ne songe pas un seul instant à la mythique Hélène, mise nouvellement en scène par Euripide avec réécriture du scénario. C'est Hésychius encore une fois – décidément le plus utile des lexicographes pour les \italique{Thesmophories}~! – qui nous le confirme (ε 1995)~:%
\citComm{ἑλένη· λαμπάς, δετή\\
\francais{hélène~: flambeau, torche.}}%
Kritylla menace Mnésiloque d'un bon coup de flambeau, flambeau qu'elle avait lors de la procession et l'arrivée au temple. (Coup de flambeau qu'elle assénera au vers 858 en évoquant Hécate \textbf{φωσ}φόρος.)}\\
ἕξεις, ἕως ἂν τῶν πρυτάνεών τις φανῇ.\\\passeLigne}%
\end{groupe}%
%% v. 855-857
\begin{groupe}%
\personnageGrec{Μνησίλοχος}%
\trimIamb{Νείλου μὲν αἵδε καλλιπάρθενοι ῥοαί,
\commentaire{855-8}{Début du prologue de l'\italique{Hélène}.}{Ainsi commence la véritable pièce d'Euripide :
\citComm{Ἑλένη --- Νείλου μὲν αἵδε καλλιπάρθενοι ῥοαί,\\
ὃς ἀντὶ δίας ψακάδος Αἰγύπτου πέδον\\
λευκῆς \difference{τακείσης χιόνος ὑγραίνει γύας.}\\
\francais{Du Nil voici les rives aux belles Nymphes,\\
Nil qui à la place de la pluie dispensée par Zeus\\
arrose la plaine d'Égypte une fois la blanche neige fondue et irrigue les terres.}}%
Comme tous les semi-habiles – euphémisme pascalien : comprenez `semi-débiles' – Euripide aime à bavasser ; il ne peut s'empêcher dès le troisième vers de la pièce d'exposer son opinion sur la raison des crues du Nil – Ah ça ! se croit-il obligé d'en avoir une quel que soit le sujet ? et croit-il par-dessus le marché qu'elle nous intéresse ? Selon donc le très savant Euripide, ces crues sont dues à la fonte des neiges. Théorie fantaisiste et traitée avec dédain par Hérodote (II, 19-25), condescendant vis-à-vis de cette thèse soutenue par Anaxagore\marqueur{1}, le maître d'Euripide – raison pour laquelle ce dernier se croit tenu de la défendre ? 
En tout cas il récidive, imperturbable, quelques années après l'\italique{Hélène} et les \italique{Thesmophories}, l'an 408/407, dans \italique{Archélaos}\marqueur{2}, tragédie aujourd'hui perdue. (Remarquons que dans le prologue d'\italique{Hélène} comme dans le fragment d'\italique{Archélaos}, les ῥοαί ne doivent pas se traduire par \italique{flots} mais bien par \italique{rives} ou \italique{berges}\marqueur{3} et que l'expression « καλλιπάρθενοι ῥοαί » suggère peut-être un jeu de mots avec les Ῥοαί, les nymphes des grenadiers\marqueur{4}.)

Attaquons-nous au détournement du prologue d'\italique{Hélène} par Aristophane (\italique{Thesm.} vv. 855-857). Pour traîner dans la boue – ou pis – la muse d'Euripide, notre poète ne recule devant rien. La syrméa dont il est question au vers 857 était une espèce de rave très abondante en Égypte. On en tirait un suc portant le même nom utilisé comme purgatif\marqueur{5}. Suivant Hérodote\marqueur{6}, les Égyptiens se purgeaient menstruellement, trois jours de suite. Ce curieux hygiénisme est assez plaisant pour qu'Aristophane n'ait pas manqué de le railler.

Aristophane avait certainement mis en scène une confrontation entre le vendeur de syrméa et le vendeur de pâtisseries dans ses \italique{Danaïdes} (comédie aujourd'hui perdue) – les gâteaux de miel étaient en effet réputés astringents\marqueur{7}.

Dans la \italique{Paix}, Trygée suggère au marchand d'armes menacé par la ruine de couper ses lances pour les proposer à la vente comme échalas. Autre idée : recycler ses casques d'une bien étrange manière (vers 1253-1254) :
\citComm{Τρυγαῖος – Πώλει βαδίζων αὐτὰ τοῖς Αἰγυπτίοις·\\
ἔστιν γὰρ ἐπιτήδεια συρμαίαν μετρεῖν.\\
Trygée – Va les vendre en Égypte.\\
Ils sont bons à mesurer la syrméa.}%
(C'est-à-dire à mesurer les \italique{effets} de la syrméa ; Trygée ne voit donc en ces casques que de vulgaires pots de chambre. En fait ils ne sont même bons à rien ; car Trygée fait entendre le proverbe κύματα μετρεῖν, compter les vagues, comble de l'inutilité\marqueur{8}.)

Même genre de détournement dans le présent passage des \italique{Thesmophories} : les Égyptiens se servent du Nil pour se \italique{nettoyer} (ici νοτίζειν = βαλανεύειν, σπογγίζειν) des désagréments de leur régime.

On peut conclure de ce passage qu'en grec le mot πέδον, le sol ; littéralement, ce sur quoi repose le pied, a fourni la même image qu'en français le \italique{fondement} et sert à désigner ἡ πυγή\marqueur{9}. 
Revenons à notre passage des \italique{Thesmophories} et proposons une traduction :%
\citComm{Μνησίλοχος – Νείλου μὲν αἵδε καλλιπάρθενοι ῥοαί,\\
ὃς ἀντὶ δίας ψακάδος Αἰγύπτου πέδον\\
λευκῆς νοτίζει μελανοσυρμαῖον λεών.\\
Mnésiloque – (Rôle d'Hélène.) « Du Nil ce sont ici les rives aux belles Nymphes,\\
Nil qui à la place de l'eau du ciel, en humecte le fondement\\
d'Égypte la blanche » et du peuple à la noire syrméa.}%
On comprend la réaction indignée de la gardienne de Mnésiloque (\italique{Thesm}. v. 858) :
\citComm{Κρίτυλλα – πανοῦργος εἶ νὴ τὴν Ἑκάτην τὴν φωσφόρον.\\
Critylla – Tu es un effronté, par Hécate la porte-torche.\\
(Elle lui porte un coup de sa torche.)}%
Cf. le vers 853 et le commentaire \italique{ad loc.} :
Critylla menaçait Mnésiloque d'un bon coup de flambeau. Et c'est au présent vers 858 qu'elle l'assène en évoquant Hécate \textbf{φωσ}φόρος.
%Ainsi commence la véritable pièce d'Euripide :%
%\citComm{Ἑλένη --- Νείλου μὲν αἵδε καλλιπάρθενοι ῥοαί,\\
%ὃς ἀντὶ δίας ψακάδος Αἰγύπτου πέδον\\
%λευκῆς \difference{τακείσης χιόνος ὑγραίνει γύας.}\\
%\francais{Du Nil voici les rives aux belles Nymphes,\\
%Nil qui à la place de la pluie dispensée par Zeus\\
%arrose la plaine d'Égypte une fois la blanche neige fondue et irrigue les terres.}}%
%Comme tous les semi-habiles -- euphémisme pascalien : comprenez `semi-débiles' --, Euripide aime à bavasser ; il ne peut s'empêcher dès le troisième vers de la pièce d'exposer son opinion sur la raison des crues du Nil -- Ah ça ! se croit-il obligé d'en avoir une quel que soit le sujet ? et croit-il par-dessus le marché qu'elle nous intéresse ? Selon donc le très savant Euripide, ces crues sont dues à la fonte des neiges. 
%Théorie fantaisiste et traitée avec dédain par Hérodote (II, 19-25), condescendant vis-à-vis de cette thèse soutenue par Anaxagore\marqueur{1}, le maître d'Euripide -- raison pour laquelle ce dernier se croit tenu de la défendre ? En tout cas il récidive quelques années après l'\italique{Hélène} et les \italique{Thesmophories}, l'an 408/407, dans \italique{Archélaos}\marqueur{2}:%
%\citComm{Δαναὸς ὁ πεντήκοντα θυγατέρων πατὴρ\\
%Νείλου λιπὼν καλλήϊον Γαιᾶ ὕδωρ, \ \ \ \ κάλλιστον ἐκ γαίας (vel γαίης) codd. : \lecon{καλλήϊον Γαιᾶ} Bibi\\
%ὃς ἐκ μελαμβρότοιο πληροῦται ῥοὰς\\
%Αἰθιοπίδος γῆς, ἡνίκ' ἂν τακῇ χιὼν\\
%τέθριπφ' ἱέντος ἡλίου κατ' αἰθέρα \ \ \ \ \ \ τέθριππεύοντος : \lecon{τέθριπφ' ἱέντος} Papathomopoulos\\
%Danaos, le père de cinquante filles $[ =$ les Danaïdes ; un oracle avait prédit que Danaos serait tué par un de ses gendres, d'où son exil d'Égypte pour ne pas que ses filles épousent leurs cousins et leur arrivée en \italique{Suppliantes} à Argos$]$,
%du Nil abandonnant l'eau fertile du Géon,
%Nil qui emplit ses berges depuis la terre
%au visage sombre d'Éthiopie chaque fois que vient à fondre la neige
%lorsque le Soleil élance ses quadriges au fil de l'éther.}%
%(Il y a quelques problèmes textuels : au second vers, les manuscrits donnent †κάλλιστον ἐκ γαίας† (ou ἐκ γαίης) ; Euripide use encore d'un néologisme formé sur κάλλι- comme avec καλλιπάρθενος ou καλλιδόναξ. Il faut rétablir καλλήϊον, \italique{à la belle récolte}, \italique{fructueuse}, d'après la glose d'Hésychius (κ 468) elle-même corrompue :%
%\citComm{†κάλλης· καρπίμου}%
%(Dans un cas comme dans l'autre, l'origine de la corruption vient d'un iota I mélu en sigma lunaire C.)\\
%La fin du vers se corrige sans difficulté à l'aide d'une autre glose d'Hésychius (γ 66) :%
%\citComm{Γαιών· ποταμός, παρ᾽ Ἕλλησι Νεῖλος}%
%Au cinquième vers, le †τέθριππεύοντος† fautif doit se lire τέθριπφ' ἱέντος conformément à l'image grecque de \italique{lancer} une cavalcade\marqueur{3} et l'usage d'Euripide d'employer τέθριππον au pluriel.\\
%On constate au passage qu'encore une fois il n'y avait nul besoin de tournebouler le texte transmis pour le remettre d'aplomb.)\\
%Remarquons qu'en ces deux passages, les ῥοαί ne doivent pas se traduire par \italique{flots} mais bien par \italique{rives} ou \italique{berges}\marqueur{4}. (L'expression «~καλλιπάρθενοι ῥοαί~» du prologue d'\italique{Hélène} suggère peut-être un jeu de mots avec les Ῥοαί, les \italique{nymphes des grenadiers}\marqueur{5}.)\\
%Attaquons-nous enfin au détournement du prologue par Aristophane !
%Pour traîner dans la boue -- ou pis -- la muse d'Euripide, notre poète ne recule devant rien. 
%La syrméa dont il est question au vers 857 était une espèce de rave très abondante en Égypte. 
%On en tirait un suc portant le même nom utilisé comme purgatif\marqueur{6}. Suivant Hérodote\marqueur{7}, les Égyptiens se purgeaient menstruellement, trois jours de suite. 
%Ce curieux hygiénisme est assez plaisant pour qu'Aristophane n'ait pas manqué de le railler.\\
%Au vers 570, nous avons vu que les gâteaux de miel grecs avaient, au contraire, des vertus astringentes, ce qui expliquait qu'Aristophane ait certainement mis en scène une confrontation entre le vendeur de syrméa et le vendeur de pâtisseries dans ses \italique{Danaïdes}, comédie aujourd'hui perdue.\\
%De même, dans la \italique{Paix}, Trygée suggère au marchand d'armes menacé par la ruine de couper ses lances pour les proposer à la vente comme échalas. Autre idée : recycler ses casques d'une bien étrange manière (vers 1253-1254) :%
%\citComm{Πώλει βαδίζων αὐτὰ τοῖς Αἰγυπτίοις·\\
%ἔστιν γὰρ ἐπιτήδεια συρμαίαν μετρεῖν.\\
%Va les vendre en Égypte.\\
%Ils sont bons à mesurer la syrméa.}%
%(C'est-à-dire à mesurer \italique{les effets} de la syrméa ; Trygée ne voit donc en ces casques que de vulgaires pots de chambre. Je dirais même plus : ils ne sont bons à rien, car Trygée fait entendre le proverbe κύματα μετρεῖν, \italique{compter les vagues}, comble de l'inutilité.)\\
%Même genre de détournement dans le présent passage des \italique{Thesmophories} : les Égyptiens se servent du Nil pour se \italique{nettoyer} (ici νοτίζειν $=$ βαλανεύειν, σπογγίζειν) des désagréments de leur régime. 
%On peut conclure de ce passage qu'en grec le mot πέδον, le \italique{sol} ; littéralement, ce sur quoi repose le pied, a fourni la même image qu'en français le \italique{fondement} et sert à désigner ἡ πυγή\marqueur{8}.\\
%\eff{
%Ce sens permet de comprendre le passage suivant des \italique{Oiseaux} (vers 504-507) où il est encore question d'Égypte : CITER COIRAULT%
%\citComm{Πισθέταιρος --- Αἰγύπτου δ’ αὖ καὶ Φοινίκης πάσης κόκκυξ βασιλεὺς ἦν·\\
%χὠπόθ’ ὁ κόκκυξ εἴποι «κόκκυ,» τότ’ <ἂν> οἱ Φοίνικες ἅπαντες\\
%τοὺς πυροὺς ἂν καὶ τὰς κριθὰς ἐν τοῖς πεδίοις ἐθέριζον.\\
%Εὐελπίδης --- Τοῦτ’ ἄρ’ ἐκεῖν’ ἦν τοὔπος ἀληθῶς· «Κόκκυ, ψωλοί, πεδίονδε.»\\
%Pisthétaire --- En Égypte, d'autre part, et dans toute la Phénicie, le coucou était roi ; et quand le coucou criait : coucou ! tous les Phéniciens allaient moissonner le froment et l'orge dans les champs\\
%Évelpide --- Le voilà donc le dicton dans toute sa vérité : coucou, déprépucés, aux champs}%
%Voici, selon Willems, comment il faut entendre ce passage. Le coucou apparaissait en Grèce au printemps. Pline nous apprend que les laboureurs qui n'avaient pas fini de tailler leurs vignes à l’équinoxe de cette saison étaient en butte à des propos salés, \latin{petulantiæ sales}, qu’on leur jetait en imitant le chant du coucou\marqueur{\italique{Hist. nat.}, XVII, 66.}. 
%Le dicton consacré par l'usage, nous l’avons ici en original : κόκκυ, ψωλοὶ, πεδίονδε. Il avait cours, ce qui n'a rien d'étonnant, en Italie aussi bien qu'en Grèce, et Horace y fait également allusion\marqueur{\italique{Satir}., I, 7, 29}. I'o)ó; était le terme cru pour désigner un débauché ; les Latins disaient verpus, et nous renverrons pour le sens du mot, tant au propre qu'au figuré, à la définition de Forcellini citée ci-dessus. En français on traduirait assez exactement : « Coucou, déprépucés, aux champs! (*) » On conçoit que pareil lardon, jeté en passant auvigneron négligent, soit appelé par Pline une exprobratio fœda. C'est, n'en doutons pas, pour amener cette équivoque cynique, familière à la généralité de ses auditeurs, qu'Aristophane a fait du coucou le roi des Égyptiens et des Phéniciens, lesquels, comme on le sait par Hérodote, étaient des circoncis (!o)oi, au sens de repitogo.). Il ne sera pas inutile d'observer, ce que certains paraissent avoir oublié, qu'Évelpide est un pur campagnard, habitant à Halimus (vv. 496 et 585), et que les exemples qu'il cite sont empruntés à la vie rustique. Nous savons, depuis Mannhardt, combien les mêmes croyances et pratiques se retrouvent chez les populations rurales de tout legroupe indo-européen. C'est ainsi que dans certaines parties de l'Allemagne les paysans ont encore coutume de se rouler par terre à la première apparition, non du milan, comme dans Aristophane (v. 502), mais du coucou (º). La présente note complète ce que dit Mannhardt à propos de l'usage mentionné par Pline (º). Entre autres rapprochements, le savant mytho graphe cite l'épithète de Horbuck lancée aux passants par les ruraux du Schleswig. Cemot depatois (en allemand Hurenbock) est bien l'équivalent du po)ó; des Grecs.
%Mais avec le sophiste des sophistes Euripide, les mots perdent tout sens ; la cause du discrédit total de la parole publique et de tous les malheurs qui s'ensuivent pour Athènes, c'est dans ses tragédies qu'il fut aller les chercher. (Aristophane rajouterait aussi dans le même panier de crabes tous les socratiques ergoteurs.) Quelle sera l'épithète assez vigoureuse pour qualifier le scribouilleur qui, pour décrire les flots (ῥοαί) du Nil, cet afflux fangeux charriant des tonnes de limon noir, recourt à l'adjectif ... καλλιπάρθενοι, \italique{à la beauté virginale} ? Il a osé ! c'est trop fort ! je vais te le réécrire, moi, ton prologue, mon salaud. Ça va être du propre ! parce que j'en ai une autre d'explication sous la main. Ta masse noirâtre, des \italique{ondes virginales} ? il nous prend vraiment pour des perdrix ! je vais te dire ce qu'est c'est : des  ῥοαί ? que non pas ; des διάρροιαι, oui !\\
%Le lecteur sent venir l'explication alternative d'Aristophane (qui n'a jamais eu peur d'appeler un chat un chat et une pièce d'Euripide de la merde). (L'image du flux d'un fleuve comparable à une colique se trouve dans un fragment du \italique{Gérytadès}\marqueur{Athénée, XII, 551a $=$ fr. 156(149.150) Kassel-Austin.} : τούτους γὰρ, ἢν πολλῷ ξυνέλθῃ, ξυλλαβὼν ὁ τῆς διαρροίας ποταμὸς οἰχήσεται, "le fleuve de fange, s'il est beaucoup gonflé, va les emporter au passage [les poètes souffreteux]".)\\
%Le verbe νοτίζειν, \italique{mouiller}, est attesté avec le sens de οὐρειν\marqueur{Voir le \LSJ }. Ici, νοτίζειν $=$ χέζειν. 
%Aristophane reprend-il une hypothèse plus populaire que celle d'Hérodote pour expliquer la crue du Nil ? ou use-t-il de mots crus pour une invention de son cru tout-à-fait incongrue ? En tout cas, le texte est clair : comme dans la chanson, les Égyptiens, ce \italique{peuple à la noire syrméa}, μελανοσυρμαῖον λεών, auront tant consommé de syrméa « que trois moulins en ont tourné / Petits ruisseaux, grandes rivières / Pendant trois jours ont débordé. »}%
}%
\debutNote{855-8}%
%\note{1}{59A91 D.-K. $=$ fr. 300 Eschyle (ed. Radt). Eschyle avait en son temps aussi évoqué cette fonte des neiges dans les \italique{Suppliantes}, v. 559 sq.% \eff{ainsi qu'au fragment}
%}%
%\note{2}{fr. 1 Ἀρχέλαος Jouan-van Looy $=$ fr. 228 Kannicht. Commentaire de Tibérius le rhéteur : οἱ γὰρ τρεῖς περὶ τοῦ ποταμοῦ στίχοι περιττοί.}%
%\note{3}{Dans le jeu enfantin de la χελιχελώνη, la «~tortue~» bondissait sur la troupe aux cris de «~λευκᾶν ἀφ' ἵππων εἰς θάλασσαν ἅλατο~», «~du haut de ses blanches cavales il a bondi dans la mer~». Voir l'édition des Περὶ βλασφημιῶν et Περὶ παιδιῶν par Taillardat, χελιχελώνη $=$ π 19.}%
%\note{4}{Hésychius (ρ 415) :%
%\citComm{ῥοία· κυλίστρα τῶν ἵππων παρὰ τῷ ποταμῷ καὶ ψάμμῳ.}%
%Rien que pour l'\italique{Hélène}, les ῥοαί sont à prendre au sens de \italique{rives} au moins aux vers 53, 124 (avec jeu de mots sur Εὐρώτας), 250, 492.}%
%\note{5}{Eustathe, \italique{Commentaire à l'}Odyssée, p. 1572, 35 : ῥοιαὶ ὁμωνύμως τῷ δένδρῳ καὶ αἱ περὶ αὐτὰς μυθικαὶ νύμφαι.}%
%\note{6}{\italique{Souda} (σ 1667) : 
%\citComm{Συρμαία: τὸ καθαρτικὸν δι' ἐμέτων καὶ κοιλίας. καὶ Συρμαί̈ζειν τὸ κενοῦσθαι, ἐκ ζειῶν καὶ ὕδατος. οὕτως Ἡρόδοτος. λάχανον σελίνῳ ἐοικός: ἢ βρῶμα διὰ στέατος καὶ μέλιτος.}%
%Et voici ce qu'en dit Hésychius (σ 2782) :
%\citComm{σύρμαια· [...] ἔστι δὲ βρωμάτιον διὰ στέατος καὶ μέλιτος. λέγεται δὲ καὶ συρμαϊσμός, καὶ ἔστι πρὸς κάθαρσιν. καὶ πόμα δι᾽ ὕδατος καὶ ἁλῶν.}}%
%\note{7}{Hérodote II, 77 ; et aussi II, 88 ; II, 125.}%
%\note{8}{Noter que si πέδον, la \italique{plaine}, renvoie au verso ; πεδίον, la \italique{petite plaine}, elle, fait allusion au recto. Cf. \italique{Lysistrata}, v. 88 ; \italique{Oiseaux}, v. 507 avec le commentaire de Willems. Voir aussi l'emploi de λειμών au vers 171 du \italique{Cyclope} d'Euripide.}
\note{1}{59A91 D.-K. = fr. 300 Eschyle (ed. Radt). Eschyle avait en son temps aussi évoqué cette fonte des neiges dans les \italique{Suppliantes}, v. 559 sq.}%
\note{2}{Fr. 1 \italique{Ἀρχέλαος} Jouan-van Looy = fr. 228 ed. Kannicht.}%
\note{3}{Hésychius (ρ 415) : ῥοία· κυλίστρα τῶν ἵππων παρὰ τῷ ποταμῷ καὶ ψάμμῳ. Rien que pour l'\italique{Hélène}, les ῥοαί sont à prendre au sens de rives au moins aux vers 53, 124 (avec jeu de mots sur Εὐρώτας), 250, 492.}%
\note{4}{Eustathe, \italique{Commentaire à l'Odyssée}, p. 1572, 35 : ῥοιαὶ ὁμωνύμως τῷ δένδρῳ καὶ αἱ περὶ αὐτὰς μυθικαὶ νύμφαι.}%
\note{5}{Souda (σ 1667) : Συρμαία : τὸ καθαρτικὸν δι' ἐμέτων καὶ κοιλίας. καὶ Συρμαίζειν τὸ κενοῦσθαι, ἐκ ζειῶν καὶ ὕδατος. οὕτως Ἡρόδοτος. λάχανον σελίνῳ ἐοικός : ἢ βρῶμα διὰ στέατος καὶ μέλιτος.\\
Et voici ce qu'en dit Hésychius (σ 2782) : σύρμαια· […] ἔστι δὲ βρωμάτιον διὰ στέατος καὶ μέλιτος. λέγεται δὲ καὶ συρμαϊσμός, καὶ ἔστι πρὸς κάθαρσιν. καὶ πόμα δι' ὕδατος καὶ ἁλῶν.}%
\note{6}{Hérodote II, 77 ; et aussi II, 88 ; II, 125.}%
\note{7}{C'est ce qui ressort d'un passage des \italique{Nuées} (vers 506-508) : Strepsiade hésite à entrer dans l'antre de Socrate, qu'il compare aimablement à celui de Trophonius : εἰς τὼ χεῖρέ νυν / δός μοι μελιτοῦτταν πρότερον, ὡς δέδοικ' ἐγὼ / εἴσω καταβαίνων ὥσπερ εἰς Τροφωνίου. « Mets-moi d'abord en main un gâteau au miel ; car j'ai peur de descendre, comme si je pénétrais dans l'antre de Trophonios. » Il faut bien sûr prononcer d'un trait `εἰς Τροφωνίου' pour faire entendre στρόφος, la colique (la peur se manifestant chez les Grecs par des troubles du bas-ventre ; en français : \italique{avoir les pétoches}). En passant, Euripide suggère le même jeu de mots scatologique : au vers 300 de son \italique{Ion} : Créuse dit à Ion qu'elle est venue avec son mari, mais qu'il a dû s'arrêter en chemin (ἐκστρέφει Τροφωνίου ; Euripide, ainsi qu'il est coutumier, joue à différents niveaux de compréhension)… Le tour était proverbial et Érasme l'a recensé (\italique{Adages} n° 677). Sur ces gâteaux de miel qui préservaient des dérangements intestinaux, voir aussi \italique{Thesm}. 570. On peut maintenant comprendre le bon mot de table qu'Athénée (\italique{Deipnosophistes}, I, 4d) rapporte : ἀμίδα δότω τις ἢ πλακοῦντά τις δότω, « Qu'on m'apporte un pot de chambre ou bien un gâteau ! » Voilà aussi pourquoi Aristophane fait intervenir dans ses \italique{Danaïdes} deux vendeurs opposés : celui de syrméa et celui d'ἐγκρίδες (Pollux, \italique{Onomasticon}, VII 199 = fr. 276(265) Kassel-Austin) : ἐν δ' Ἀριστοφάνης Δαναΐσι συρμαιοπῶλαι καὶ ἐγκριδοπῶλαι. Athénée (XIV, 645e) définit les encrides comme des gâteaux frits à l'huile et emmiellés : ΕΓΚΡΙΔΕΣ. πεμμάτιον ἑψόμενον ἐν ἐλαίῳ καὶ μετὰ τοῦτο μελιτούμενον.}%
\note{8}{Voir Érasme, \italique{Adages} (n° 345, \latin{Undas numeras}) ou Suétone, Περὶ βλασφημιῶν (éd. Taillardat), β 188.}%
\note{9}{Noter que si πέδον, \italique{la plaine}, renvoie au verso ; πεδίον, \italique{la petite plaine}, elle, fait allusion au recto. Cf. \italique{Lysistrata}, v. 88 ; \italique{Oiseaux}, v. 507 avec le commentaire de Willems. Voir aussi l'emploi de λειμών au vers 171 du \italique{Cyclope} d'Euripide.}%
\\
ὃς ἀντὶ δίας ψακάδος Αἰγύπτου πέδον\\
λευκῆς νοτίζει μελανοσυρμαῖον λεών.\\\passeLigne}%
\end{groupe}%
%% v. 858
\begin{groupe}%
\personnageGrec{Κρίτυλλα}%
\trimIamb{πανοῦργος εἶ νὴ τὴν Ἑκάτην τὴν φωσφόρον.%
\commentaire{858}{τὴν Ἑκάτην τὴν φωσφόρον.}{Ce n'est pas par hasard que Critylla adjoint à Hécate l'épithète de φωσφόρος, \italique{porte-torche} : elle est elle-même équipée d'une telle torche et il est à parier qu'en jurant de la sorte, Critylla met sa menace du vers 853 à exécution en assénant à Mnésiloque le coup de flambeau promis. %\marqueur{1}.
Cf. \italique{Lysistrata} vers 453 ; fr. 592.5 Kassel-Austin [$=$ Athénée XV, 701b] : ἐκφέρετε πεύκας κατ᾽ Ἀγάθωνα φωσφόρους.}%
%\debutNote{858}%
%\note{1}{Cf. \italique{Lysistrata} vers 453 ; fr. 592.5 Kassel-Austin [$=$ Athénée XV, 701b] : ἐκφέρετε πεύκας κατ᾽ Ἀγάθωνα φωσφόρους.}
\\\passeLigne}%
\end{groupe}%
%% v. 859-860a
\begin{groupe}%
\personnageGrec{Μνησίλοχος}%
\trimIamb{ἐμοὶ δὲ γῆ μὲν πατρὶς οὐκ ἀνώνυμος%
\commentaire{859-60}{}{\italique{Hélène}, v. 16-17a :%
\citComm{\difference{ἡμῖν} δὲ γῆ μὲν πατρὶς οὐκ ἀνώνυμος\\
Σπάρτη, πατὴρ δὲ Τυνδάρεως.}%
}\\
Σπάρτη, πατὴρ δὲ Τυνδάρεως.\\\passeLigne}%
\end{groupe}%
%% v. 860b-861
\begin{groupe}%
\personnageGrec{Κρίτυλλα}%
\trimIamb{\ligne{860}\decalage{Σπάρτη, πατὴρ δὲ Τυνδάρεως.}σοί γ᾽ ὦλεθρε\apparat{860}{σύ (ex σί) \Rdeux : \lecon{σοί} Zanetti (\italique{tibi} Divus) \espApp τ' R : \lecon{γ᾽} Gelenius \ligneApp{861} ἐστιν R : \lecon{ἐστι} \Mon}\\
πατὴρ ἐκεῖνός ἐστι; Φρυνώνδας μὲν οὖν.\testimonia{861}{(Φρυνώνδας) Suidas φ 700 ; \Syn φ 211 Φρυνώνδας $\sim$ Hesychius φ 938}%
\commentaire{861}{Φρυνώνδας.}{Phrynondas et Eurybatos, deux coquins dont le nom était devenu appellatif, comme en français Cartouche et Mandrin. (Note de Willems ; qui a en tête le \italique{Protagoras} (327d) de Platon lorsque ces deux truands de la tradition populaire sont rapprochés.)\\
\italique{Souda}, φ 770 :%
\citComm{Φρυνώνδας· τῶν ἐπὶ πονηρίᾳ διαβεβοημένων. ὃς ξένος ὢν κατὰ τὰ Πελοποννησιακὰ διέτριβεν Ἀθήνησιν. Ἀριστοφάνης Ἀμφιάρεῳ· ὦ μιαρὲ καὶ Φρυνώνδα καὶ πονηρὲ σύ. ἐκ τούτου τοὺς πονηροὺς Φρυνώνδας καλοῦσι. κακοήθης ὁ Φρυνώνδας οὗτος καὶ πανοῦργος, μισθοῦ τινας μηχανὰς πραγμάτων ἐπὶ κακοῖς ῥᾳδίως συντιθείς, ὥς φησιν Εὔδημος. καὶ αὖθις· ἀλλ' οἶμαι οὔτε Φρυνώνδας οὔτ' Εὐρύβατος οὐδ' ἄλλος πώποτε τῶν πάλαι πονηρῶν, τοιοῦτος μάγος καὶ γόης.\\
\francais{Phrynondas : de ceux connus pour leur méchanceté. C'était un étranger qui vivait à Athènes au moment de la guerre du Péloponnèse. Aristophane dans son \italique{Amphiaraos} : « ô misérable et Phrynondas et vil que tu es ». On appelle les méchants de son nom : Phrynondas.
Ce Phrynondas [était] mauvais et méchant ; quelqu'un qui, pour de l'argent, élaborait facilement quelque intrigue pour de mauvaises fins, comme le dit Eudemos.
Et ailleurs : « mais je pense que ni Phrynondas ni Eurybatos ni aucun autre des vilains d'autrefois ne se sont jamais comportés comme un filou et un tricheur ».}%
}%
}\\\passeLigne}%
\end{groupe}%
%% v. 862a
\begin{groupe}%
%
\personnageGrec{Μνησίλοχος}%
\trimIamb{Ἑλένη δ᾽ ἐκλήθην.%
\commentaire{862}{}{\italique{Hélène}, v. 22 :%
\citComm{Ἑλένη δ᾽ ἐκλήθην. ἃ δὲ πεπόνθαμεν κακὰ}%
}\\\passeLigne}%
\end{groupe}%
%% v. 862b-863
\begin{groupe}%
\personnageGrec{Κρίτυλλα}%
\trimIamb{\ligne{862}\decalage{Ἑλένη δ᾽ ἐκλήθην.}αὖθις αὖ γίγνει γυνή,\apparat{862}{γίγνη R : \lecon{γίγνει} Brunck \ligneApp{863} \lecon{δοῦναι post ἑτέρας} R : post δίκην Suid. γ 496}\\
πρὶν τῆς ἑτέρας δοῦναι γυναικίσεως δίκην;\testimonia{863}{Suidas γ 496 γυναίκισις et Zonaras 459 πρὶν τῆς ἑτέρας γυναικίσεως δίκην δοῦναι}%\apparat{863}{\lecon{δοῦναι post ἑτέρας} R : post δίκην Suid. γ 496}
\\\passeLigne}%
\end{groupe}%
%% v. 864-865a
\begin{groupe}%
\personnageGrec{Μνησίλοχος}%
\trimIamb{ψυχαὶ δὲ πολλαὶ δι᾽ ἔμ᾽ ἐπὶ Σκαμανδρίοις\apparat{864}{Σκαμανδρίαις R : \lecon{Σκαμανδρίοις} Fritzsche ex Eur.\italique{Hel} 52 et 609}%
\commentaire{864-5}{}{\italique{Hélène}, v. 52-53a :%
\citComm{ψυχαὶ δὲ πολλαὶ δι᾽ ἔμ᾽ ἐπὶ Σκαμανδρίοις\\
ῥοαῖσιν ἔθανον· \difference{ἡ δὲ πάντα τλᾶσ᾽ ἐγὼ}}%
}\\
ῥοαῖσιν ἔθανον.\\\passeLigne}%
\end{groupe}%
%% v. 865b
\begin{groupe}%
\personnageGrec{Κρίτυλλα}%
\trimIamb{\ligne{865}\decalage{ῥοαῖσιν ἔθανον.}ὤφελες δὲ καὶ σύ γε.\\\passeLigne}%
\end{groupe}%
%% v. 866-868a
\begin{groupe}%
\personnageGrec{Μνησίλοχος}%
\trimIamb{κἀγὼ μὲν ἐνθάδ᾽ εἴμ᾽· ὁ δ᾽ ἄθλιος πόσις\\
οὑμὸς Μενέλεως οὐδέπω προσέρχεται.\apparat{867}{Μενέλαος R : \lecon{Μενέλεως} Brunck, cf. v.901}
\commentaire{866-8}{}{\italique{Hélène}, v. 49 :%
\citComm{κἀγὼ μὲν ἐνθάδ᾽ εἴμ᾽, ὁ δ᾽ ἄθλιος πόσις}%
\italique{Hélène}, v. 56a : %
\citComm{τί δῆτ᾽ ἔτι ζῶ; }%
}\\
τί οὖν ἔτι ζῶ;\\\passeLigne}%
\end{groupe}%
%% v. 868b
\begin{groupe}%
\personnageGrec{Κρίτυλλα}%
\trimIamb{\ligne{868}\decalage{τί οὖν ἔτι ζῶ;}τῶν κοράκων πονηρίᾳ.\testimonia{868}{Suidas τ 648}\apparat{868}{τῶν κορ. πον. continuat (Μνη.) R : \lecon{(Γυνή)} Σ, Ellebodius}\\\passeLigne}%
\end{groupe}%
%% v. 869-870
\begin{groupe}%
\personnageGrec{Μνησίλοχος}%
\trimIamb{ἀλλ᾽ ὥσπερ αἰκάλλει τι καρδίαν ἐμήν.\testimonia{869}{(αἰκάλλει τι καρδίαν ἐμήν) Suidas αι 161 ; Zonaras 91 (τι omisso)}\\
μὴ ψεῦσον ὦ Ζεῦ τῆς ἐπιούσης ἐλπίδος.\testimonia{870}{\GregCor p.14-15 ($=$ scholium\expo{Mosch.} Euripides, \italique{Hecub.} 225) (sed ἡμέρας pro ἐλπίδος) ; Suidas ψ 62 \espTest (μὴ ψεῦσον) \Antiatt 107.31}%\apparat{post v. 870}{< οἶσθ᾽ οὖν ὃ δρᾶσον; $ - \cup - | \times - \cup - $ > \lecon{om.} R : restituit ap. Σ Bibi (cf. Eur. \italique{Hel.}325)}%\apparat{870}{\lecon{ἐλπίδος} R : ἡμέρας ap. Greg. Cor. p.14-15, Σ ad Eur. \italique{Hec}.225 \espApp post v. 870 < οἶσθ᾽ οὖν ὃ δρᾶσον; $ - \cup - | \times - \cup - $ > \lecon{om.} R : restituit ap. Σ Bibi (cf. Eur. \italique{Hel.}325)}%
%οἶσθ᾽ οὖν ὃ δρᾶσον; $ - \cup - | \times - \cup - $ %
\commentaire{870}{μὴ ψεῦσον ὦ Ζεῦ τῆς ἐπιούσης ἐλπίδος.}{Voici la scholie au vers 870 :% -- malheureusement lacunaire :%-- laisse à penser que le texte lui-même du manuscrit de Ravenne admet aussi quelque lacune. Voici la scholie :%
\citComm{σύνηθες τὸ σχῆμα. Μένανδρος·\\
\alinea οἶσθ᾽ ὃ ποίησον.\\
Σοφοκλῆς Πηλεῖ·\\
\alinea μὴ ψεῦσον ὦ Ζεῦ· μή μ᾽ ἕλης ἄνευ δορός.\\
\alinea οἶσθ᾽ οὖν ὃ δρᾶσον.\\
Εὐριπίδης Πολυΐδῳ.\\
\francais{La construction est usuelle. Ménandre : 
« Tu sais quoi faire. » 
Sophocle dans son \italique{Pelée} : 
« Ne me trompe pas, ô Zeus ; tu ne me prendras pas sans lutte. /
Tu sais donc ce qu'il te reste à faire. »
(Euripide dans son \italique{Polyidos}.)}}%
Cette scholie ne se peut comprendre qu'en la mettant en regard d'un passage du Περὶ διαλέκτων de Grégoire de Corinthe (\nieme{XII} siècle, p. 14 ed. Schaefer) ou de la scholie de l'\italique{Hécube} d'Euripide (vers 225 -- nos trois commentaires remontent assurément à la même source ; cf. aussi Suidas ψ 62). À propos de la tournure `οἶσθ᾽ οὖν ὃ δρᾶσον;', voici ce que le scholiaste de l'\italique{Hécube} écrit :%
\citComm{δρᾶσον: τὰ καλὰ τῶν ἀντιγράφων δρᾶσον ἔχει. οἱ Ἀττικοὶ χρῶνται τοῖς προστακτικοῖς ἀντὶ ὑποτακτικῶν ἀπαγορευτικῶν, καὶ αὐτοῖς τοῖς προστακτικοῖς ἀντὶ ὁριστικῶν.\\
τοῦ μὲν οὖν πρώτου τὸ Ἀριστοφάνους ἐν Θεσμοφοριαζούσαις (v. 870) παράδειγμα `μὴ ψεῦσον, ὦ Ζεῦ, τῆς ἐπιούσης ἡμέρας (sic).'\\
τοῦ δευτέρου δὲ Μένανδρος `οἶσθ' ὅ τι ποίησον' ἀντὶ 
%<note type="footnote">21.
%ποιήσεις Matth. pro ποιήσῃς. 27. οἱ Ἀττικοὶ —] Scholion ex 26.
%δρᾶσον, superscripto εις, in Greg. Cor. p. 14. compositum. textu
%habet M. In Gr. est τὰ ἀκριβέστερα 30. Θεσμοφ.] 870. ubi ἐπιούσης
%τῶν ἀντιγράφων οἶσθ᾿ οὖν ὃ ἐλπίδος. δρᾶσον ἔχουσιν.</note>
τοῦ `ὅ τι ποιήσεις'.\\
τὸν αὐτὸν τρόπον καὶ Εὐριπίδης `οἶσθ᾿ οὖν ὃ δρᾶσον' ἀντὶ τοῦ `δράσεις' ἐξήνεγκεν.\\
δρᾶσον: les meilleures copies ont `δρᾶσον'.
Les Attiques recourent aux impératifs à place de subjonctifs de prohibition. Et de nouveau aux (je lis : καὶ αὖ τοῖς) impératifs à la place d'indicatifs.\\
L'exemple d'Aristophane sur le premier point : `μὴ ψεῦσον κτλ.'\\
Ménandre sur le second point : `οἶσθ' ὅ τι ποίησον' à la place de `ὅ τι ποιήσεις'.\\
Euripide a le même tour : `οἶσθ᾿ οὖν ὃ δρᾶσον' à la place de `δράσεις'.}%
Noter que ce tour d'Euripide se trouve justement dans l'\italique{Hélène} (vers 325) :%
\citComm{Χορός --- οἶσθ᾽ οὖν ὃ δρᾶσον; μνήματος λιποῦσ᾽ ἕδραν.\\
\francais{Hé bien, sais-tu ce que tu dois faire ? Quitte ton poste près du tombeau.}}%
%Faut-il en induire qu'Aristophane le détournait à ce moment des \italique{Thesmophories} ? 
%Cela ne me paraît pas impossible : ces mots seraient même tout-à-fait de saison dans la bouche de Mnésiloque : Euripide et lui ne se sont pas revus depuis la scène de travestissement chez Agathon ; comment Euripide saurait-il le plan échafaudé par Mnésiloque (détourner de ses pièces pour s'évader) ? il me paraît vraisemblable que Mnésiloque voyant apparaître Euripide lui demande s'il connaît bien son rôle -- rôle de Ménélas qu'il connaît forcément, l'ayant lui-même écrit. Façon habile de suggérer à Euripide quoi faire sans éveiller les soupçons de Critylla. Un vers de l'\allemand{Urtext} des \italique{Thesmophories} se sera-t-il perdu dans les scholies ? la conjecture est malheureusement invérifiable et restera circonscrite au commentaire :%
%\citComm{Μνησίλοχος -- ἀλλ᾽ ὥσπερ αἰκάλλει τι καρδίαν ἐμήν.\\
%μὴ ψεῦσον ὦ Ζεῦ τῆς ἐπιούσης ἐλπίδος.\\
%οἶσθ᾽ οὖν ὃ δρᾶσον; $ - \cup - | \times - \cup - $}%
%Le point commun de ces citations est qu'elles comportent toutes deux « οἶσθ᾽ὃ ποίησον / δρᾶσον» ; si le scholiaste les a rapprochées du présent passage, il paraît probable que l'expression se fût aussi trouvée dans le texte original. 
%D'ailleurs, ces mots seraient tout-à-fait de saison dans la bouche de Mnésiloque : Euripide et lui ne se sont pas revus depuis la scène de travestissement chez Agathon ; comment Euripide saurait-il le plan échafaudé par Mnésiloque (détourner de ses pièces pour s'évader) ? il nous paraît vraisemblable que Mnésiloque  voyant apparaître Euripide lui demande s'il connaît bien son rôle -- rôle de Ménélas qu'il connaît forcément, l'ayant lui-même écrit. Façon habile de suggérer à Euripide quoi faire sans éveiller les soupçons de Critylla.\\
%Peut-être Aristophane détournait-il légèrement le vers 325 de l'\italique{Hélène} :
%\citComm{Χορός --- οἶσθ᾽ οὖν ὃ δρᾶσον; μνήματος λιποῦσ᾽ ἕδραν.\\
%\francais{Hé bien, sais-tu ce que tu dois faire ? Quitte ton poste près du tombeau.}}%
}%\\
\\\passeLigne}%
\end{groupe}%
%% v. 871-873
\begin{groupe}%
%
\personnageGrec{Εὐριπίδης}%
\trimIamb{\ligne{871}τίς τῶνδ᾽ ἐρυμνῶν δωμάτων ἔχει κράτος,%\apparat{871}{ὡς Περσεύς in mg. R : \lecon{ὡς Μενέλαος} Zanetti \ligneApp{872} ξένος R : \lecon{ξένους} Zanetti \ligneApp{873} \lecon{κάμνοντας} R : καμόντας Lenting}%
\commentaire{871}{}{\italique{Hélène}, v. 68 :%
\citComm{Τεῦκρος --- τίς τῶνδ᾽ ἐρυμνῶν δωμάτων ἔχει κράτος;}%
\italique{Hélène}, v. 449 :%
\citComm{Μενελέως --- ναυαγὸς \difference{ἥκω} ξένος, \difference{ἀσύλητον γένος}.}}\\
ὅστις ξένους δέξαιτο ποντίῳ σάλῳ\apparat{872}{ξένος R : \lecon{ξένους} Zanetti}\\
κάμνοντας ἐν χειμῶνι καὶ ναυαγίαις;\apparat{873}{\lecon{κάμνοντας} R : καμόντας Lenting}\\\passeLigne}%
\end{groupe}%
%% v. 874a
\begin{groupe}%
\personnageGrec{Μνησίλοχος}%
\trimIamb{Πρωτέως τάδ᾽ ἐστὶ μέλαθρα.%
\commentaire{874-8}{}{\italique{Hélène}, v. 460 : %
\citComm{Πρωτεὺς τάδ᾽ οἰκεῖ δώματ᾽, Αἴγυπτος δὲ γῆ.}}\\\passeLigne}%
\end{groupe}%
%% v. 874b-876 (part 1)
\begin{groupe}\personnageGrec{Κρίτυλλα}
\trimIamb{\ligne{874}\decalage{Πρωτέως τάδ᾽ ἐστὶ μέλαθρα.}ποίου Πρωτέως;
\commentaire{874-6}{ποίου Πρωτέως; ... ἐπεὶ τέθνηκε Πρωτέας ἔτη δέκα.}{Elle confond Protée avec Protéas\marqueur{1}, stratège athénien dont il est question dans Thucydide\marqueur{2}. (Note de Willems.)
}
\debutNote{874-6}
\note{1}{Identiques au génitif.}
\note{2}{Thucydide I, 45.2 ; II, 23.2.}\\
ὤ τρισκακοδαίμων, ψεύδεται νὴ τὼ θεώ,\apparat{875}{ὦ τρισκακόδαιμον R : \lecon{ὤ τρισκακοδαίμων} Wilamowitz ad \italique{Lys.}449}\\\passeLigne}%
\end{groupe}%
\endnumbering
\end{Leftside}
\begin{Rightside}
\beginnumbering
%% v. 838+ (part 2)
\begin{groupe}
\tetramTrochFr{\ligne{838}s'assiérait, les cheveux taillés au \fin bol, derrière}
\end{groupe}%
%% v. 830-845
\begin{groupe}%
\tetramTrochFr{\ligne{839}la femme qui enfanta le brave. Comment tolérer, ô Eupolis,
que la mère d'Hyperbolos siège, vêtue
de blanc, les cheveux flottants, à côté de la mère de Lamachos,
et qu'elle prête de l'argent à usure... le faut-il supporter ? Si elle prête de l'argent
et exige un profit, loin de lui payer un intérêt,
il faudrait lui ôter de force ses richesses, en lui disant :
« Tu vaux bien que ton argent produise, toi qui produisis un tel fruit. »\\\passeLigne}%
\end{groupe}%
%% v. 846-851
\begin{groupe}%
\personnageFr{Mnésiloque}%
\trimIambFr{\ligne{846}Je suis devenu louche à lorgner de ce côté. Et lui ?... \fin pas encore.\\
Qu'est-ce donc qui l'empêche ? Nul doute,\\
le \italique{Palamède} est froid et il en a honte.\\
Par quel drame alors l'attirer ...\\
J'y suis : je vais contrefaire sa récente \italique{Hélène}\\
Aussi bien, je suis accoutré d'un costume de femme.\\\passeLigne}%
\end{groupe}%
%% v. 852-854
\begin{groupe}%
\personnageFr{Kritylla}%
\trimIambFr{\ligne{851}Que manigances-tu encore ? Qu'as-tu à écarquiller les yeux ?
On t'en fera voir une Hélène, amère celle-là, si tu ne te tiens décemment
jusqu'à la venue de l'un des \fin prytanes.\\\passeLigne}%
\end{groupe}%
%% v. 855-857
\begin{groupe}%
\personnageFr{Mnésiloque}%
\trimIambFr{(Rôle d'Helène.) « Du Nil ce sont ici les rives aux belles \fin Nymphes,\\
Nil qui à la place de l'eau du ciel, humecte le fonde-\fin ment d'Égypte\\
la blanche » et du peuple à la noire syrméa.\\\passeLigne}%
\end{groupe}%
%% v. 858
\begin{groupe}%
\personnageFr{Kritylla}%
\trimIambFr{Tu es un effronté, par Hécate la porte-torche. (Elle lui porte un coup de sa torche.)\\\passeLigne}%
\end{groupe}%
%% v. 859-860a
\begin{groupe}%
\personnageFr{Mnésiloque}%
\trimIambFr{« J'ai pour patrie un lieu non sans renom,\\
c'est Sparte, et mon père est Tyndare. »\\\passeLigne}%
\end{groupe}%
%% v. 860b-861
\begin{groupe}%
\personnageFr{Kritylla}%
\trimIambFr{\ligne{860}\espace{4.5cm}Lui ton père, être de malheur ?\\
Phrynondas, veux-tu dire.\\\passeLigne}%
\end{groupe}%
%% v. 862a
\begin{groupe}%
\personnageFr{Mnésiloque}%
\trimIambFr{« Je fus nommée Hélène. »\\\passeLigne}%
\end{groupe}%
%% v. 862b-863
\begin{groupe}%
\personnageFr{Kritylla}%
\trimIambFr{\ligne{862}\espace{2.8cm}Voilà qu'encore tu redeviens femme,\\
avant d'avoir expié ton premier déguisement.\\\passeLigne}%
\end{groupe}%
%% v. 864-865a
\begin{groupe}%
\personnageFr{Mnésiloque}%
\trimIambFr{« Et à cause de moi, bien des guerriers sont morts\\
aux rives du Scamandre. »\\\passeLigne}%
\end{groupe}%
%% v. 865b
\begin{groupe}%
\personnageFr{Kritylla}%
\trimIambFr{\ligne{865}\espace{2.4cm}Que n'as-tu fait comme eux ?\\\passeLigne}%
\end{groupe}%
%% v. 866-868a
\begin{groupe}%
\personnageFr{Mnésiloque}%
\trimIambFr{« Moi, me voici. Mais lui l'infortuné », mon époux\\
Ménélas n'apparaît toujours point.\\
« Ah ! pourquoi vis-je encore? »\\\passeLigne}%
\end{groupe}%
%% v. 868b
\begin{groupe}%
\personnageFr{Kritylla}%
\trimIambFr{\ligne{868}\espace{2.3cm}C'est la faute aux corbeaux.\\\passeLigne}%
\end{groupe}%
%% v. 869-870
\begin{groupe}%
\personnageFr{Mnésiloque}%
\trimIambFr{(Mnésiloque aperçoit Euripide.)
Mais quelque chose \fin vient flatter mon cœur.\\
Ô Zeus, ne démens pas ma naissante espérance !%\\
%(À Euripide :) Sais-tu donc ce qu'il te reste à faire ? [ ... ] 
\\\passeLigne}%
\end{groupe}%
%% v. 871-873
\begin{groupe}%
\personnageFr{Euripide}%
\trimIambFr{(Rôle de Ménélas.) « Quel est le chef puissant de ces fortes demeures ? » 
Voudra-t-il accueillir des hôtes dérangés par la houle marine
au cours d'une tempête et \fin d'un naufrage ?\\\passeLigne}%
\end{groupe}%
%% v. 874a
\begin{groupe}%
\personnageFr{Mnésiloque}%
\trimIambFr{C'est à Protée qu'est ce palais.\\\passeLigne}%
\end{groupe}%
%% v. 874b-876 (part 1)
\begin{groupe}\personnageFr{Kritylla}
\trimIambFr{\ligne{874}\espace{4.2cm}Qui ça, Protée ?\\
Oh le triple misérable! Il ment par les deux déesses.\\\passeLigne}%
\end{groupe}%
\endnumbering
\end{Rightside}
\end{pages}
\Pages
\clearpage\ifodd\value{page}\hbox{}\newpage\fi
\enTeteG{ΕΠΕΙΣΟΔΙΟΝ Δ' (Ἑλένη) --- Codex Ravennas 429, f. 173r (Classense 175r) --- vv. 876–919}
\begin{adjustwidth}{-1.5cm}{-3.0cm}
\begin{adjustbox}{max totalheight=21.0cm, max width=16.5cm, keepaspectratio}
\begin{minipage}{16.5cm}
{\fontsize{7.2pt}{8.5pt}\selectfont
\noindent\textbf{Folio 22 --- Codex Ravennas 429, \texttt{f. 173r} (\texttt{Classense 175r}) : ΕΠΕΙΣΟΔΙΟΝ Δ' (Ἑλένη), vv. 876–919}\hfill \textit{Témoins :} $\text{R}^1, \Rpc, \Rdeux$ (\Mon)\par\vspace{2pt}
\hrule height 0.4pt\vspace{2.5pt}
}
{\fontsize{8.1pt}{9.5pt}\selectfont
\ligneRav{876}{\auDessus{γυ}{ν}}{ἐπεὶ τέθνηκε Πρωτέας ἔτη δέκα·}{}
\ligneRav{}{\auDessus{εὐ}{ρ}}{ποίαν δὲ χώραν εἰσεκέλσαμεν σκάφει;}{}
\ligneRav{}{μ\auDessus{ν}{Η}}{Αἴγυπτον. : ὦ δύστηνος οἷ πεπλώκαμεν·}{}
\ligneRav{879}{μ\auDessus{ν}{Η}}{πείθει τι \difference{τῷ κακῷ} κακῶς ἀπολουμένῳ}{\textbf{879} τῷ R : τούτῳ τῷ Grynaeus : \lecon{τῷ κακῷ} Scaliger}
\ligneRav{880}{}{ληροῦντι λῆρον; Θεσμοφόριον τουτογί·}{\textbf{880} \lecon{Θεσμοφόριον} R, cf.278}
\ligneRav{881}{\auDessus{εὐ}{ρ}}{αὐτὸς δὲ Πρωτεὺς ἔνδον ἔστ᾽ \difference{ἢ 'ξώπιος};}{\textbf{881} ἠξώπιος R, id. v.884 : \lecon{ἢ 'ξώπιος;} Biset}
\ligneRav{}{\auDessus{γυ}{ν}}{οὐκ ἔσθ᾽ ὅπως οὐ ναυτιᾷς ἔτ᾽ ὦ ξένε,}{}
\ligneRav{883}{}{\difference{ὅστις γ᾽} ἀκούσας ὅτι τέθνηκε Πρωτέας}{\textbf{883} ὅστις R : \lecon{ὅστις γ᾽} Scaliger}
\ligneRav{}{}{ἔπειτ᾽ ἐρωτᾷς `ἔνδον ἔστ᾽ ἢ 'ξώπιος;'}{}
\ligneRav{885}{\auDessus{εὐ}{ρ}}{αἰαῖ \difference{τέθνηκε}. ποῦ δ᾽ ἐτυμβεύθη τάφῳ;}{\textbf{885} τέθνηκεν R : \lecon{τέθνηκε} Bentley}
\ligneRav{}{μ\auDessus{ν}{Η}}{τόδ᾽ ἐστὶν αὐτοῦ σῆμ᾽, ἐφ᾽ ᾧ καθήμεθα·}{}
\ligneRav{887}{\auDessus{γυ}{ν}}{κακῶς τ᾽ ἄρ᾽ ἐξόλοιο κἀξολεῖ γέ τοι,}{\textbf{887} \lecon{κακῶς τ᾽ ἄρ᾽} R : κακῶς γ᾽ ἄρ᾽ Denniston (\italique{CR} XLIV(1930), p.215)}
\ligneRav{}{}{ὅστις γε τολμᾷς σῆμα τὸν βωμὸν καλεῖν·}{}
\ligneRav{889}{\auDessus{εὐ}{ρ}}{\difference{τί δὴ σὺ} θάσσεις τάσδε τυμβήρεις ἕδρας}{\textbf{889} τί δὲ σὺ R : τί δαὶ σὺ Faber : \lecon{τί δὴ σὺ} Scaliger}
\ligneRav{890}{—}{φάρει καλυπτὸς ὦ ξένη; : βιάζομαι}{}
\ligneRav{891}{}{γάμοισι Πρωτέως παιδὶ \difference{συμμεῖξαι} λέχος·}{\textbf{891} συμμῖξαι R : \lecon{συμμεῖξαι} van Leeuwen}
\ligneRav{892}{\auDessus{γυ}{ν}}{τί ὦ κακόδαιμον ἐξαπατᾷς \difference{αὖ τὸν} ξένον;}{\textbf{892} αὐτὸν (ex αὖτὸν) \Rpc : \lecon{αὖ τὸν} Casaubon}
\ligneRav{}{}{οὗτος πανουργῶν δεῦρ᾽ ἀνῆλθεν ὦ ξένε}{}
\ligneRav{}{}{ὡς τὰς γυναῖκας ἐπὶ κλοπῇ τοῦ χρυσίου·}{}
\ligneRav{895}{μ\auDessus{ν}{Η}}{\difference{κόνυζα} τοὐμὸν σῶμα βάλλουσα ψόγῳ·}{\textbf{895} βαϋζα R : βάϋζε Zanetti : \lecon{κόνυζα} Bibi}
\ligneRav{}{\auDessus{εὐ}{ρ}}{ξένη τίς ἡ γραῦς ἡ κακορροθοῦσά σε;}{}
\ligneRav{}{μ\auDessus{ν}{Η}}{αὕτη Θεονόη Πρωτέως. : μὰ τὼ θεὼ}{}
\ligneRav{}{}{εἰ μὴ Κρίτυλλά γ᾽ Ἀντιθέου Γαργηττόθεν·}{}
\ligneRav{}{—}{σὺ δ᾽ εἶ πανοῦργος. : ὁπόσα τοι βούλει λέγε·}{}
\ligneRav{900}{}{οὐ γὰρ γαμοῦμαι σῷ κασιγνήτῳ ποτὲ}{}
\ligneRav{901}{}{προδοῦσα \difference{Μενέλεων} τὸν ἐμὸν ἐν Τροίᾳ πόσιν·}{\textbf{901} Μενέλαον R : \lecon{Μενέλεων} Scaliger}
\ligneRav{}{\auDessus{εὐ}{ρ}}{γύναι τί εἶπας; στρέψον ἀνταυγεῖς κόρας·}{}
\ligneRav{}{μ\auDessus{ν}{Η}}{αἰσχύνομαί σε τὰς γνάθους ὑβρισμένη·}{}
\ligneRav{}{\auDessus{εὐ}{ρ}}{τουτὶ τί ἔστιν; ἀφασία τίς τοί μ᾽ ἔχει·}{}
\ligneRav{905}{}{ὦ θεοὶ τίν᾽ ὄψιν εἰσορῶ; τίς εἶ γύναι;}{}
\ligneRav{906}{μ\auDessus{ν}{Η}}{σὺ δ᾽ εἶ τίς; \difference{αὑτὸς} γὰρ σὲ κἄμ᾽ ἔχει λόγος·}{\textbf{906} αὐτὸς R : \lecon{αὑτὸς} Boissonade}
\ligneRav{}{\auDessus{εὐ}{ρ}}{Ἑλληνὶς εἶ τις ἢ 'πιχωρία γυνή;}{}
\ligneRav{}{μ\auDessus{ν}{Η}}{Ἑλληνίς. ἀλλὰ καὶ τὸ σὸν θέλω μαθεῖν·}{}
\ligneRav{909}{\auDessus{εὐ}{ρ}}{\difference{Ἑλένῃ σ᾽} ὁμοίαν δὴ \difference{μάλιστ᾽ εἶδον} γύναι·}{\textbf{909} Ἐλένης R : \lecon{Ἑλένῃ σ᾽} Eur. \italique{Hel.}563 \espApp μάλισθ᾽ ἴδον R : \lecon{μάλιστ᾽ εἶδον} \italique{Hel}.563}
\ligneRav{910}{μ\auDessus{ν}{Η}}{ἐγὼ δὲ \difference{Μενελέῳ σ᾽} ὅσα γ᾽ ἐκ τῶν \difference{λαιφίων}·}{\textbf{910} Μενέλαον R : Μενέλεων Suid. : \lecon{Μενελέῳ σ᾽} Thiersch \espApp ἀφύων R : ἰφύων Σ, Suid. : ἀμφίων Grégoire : \lecon{λαιφίων} Bibi\espApp Μενέλαον R : Μενέλεων Suid.}
\ligneRav{}{\auDessus{εὐ}{ρ}}{ἔγνως ἄρ᾽ ὀρθῶς ἄνδρα δυστυχέστατον·}{}
\ligneRav{}{μ\auDessus{ν}{Η}}{ὦ χρόνιος ἐλθὼν σῆς δάμαρτος ἐσχάρας}{}
\ligneRav{914}{}{\hspace{1.1em}λαβέ με λαβέ με πόσι, \difference{περίβαλε} δὲ χέρας·}{\textbf{914} περίβαλλε R : \lecon{περίβαλε} Biset}
\ligneRav{915}{}{\hspace{1.1em}φέρε σὲ κύσω. ἄπαγέ μ᾽ ἄπαγ᾽ ἄπαγέ με}{}
\ligneRav{}{—}{λαβὼν ταχὺ πάνυ. : κλαύσετ᾽ ἄρα νὴ τὼ θεὼ}{}
\ligneRav{}{}{ὅστις σ᾽ ἀπάξει τυπτόμενος τῇ λαμπάδι·}{}
\ligneRav{}{\auDessus{εὐ}{ρ}}{σὺ τὴν ἐμὴν γυναῖκα κωλύεις ἐμέ,}{}
\ligneRav{}{}{τὴν Τυνδάρειον παῖδ᾽, ἐπὶ Σπάρτην ἄγειν;}{}
}
\vspace{3pt}\noindent\rule{4.5cm}{0.4pt}\par\vspace{1.5pt}
{\fontsize{7.0pt}{8.3pt}\selectfont
\noindent\textbf{Loci paralleli (vv. 876–919) ---} \textbf{878} Suidas οι 1 ὦ δύστηνος, οἷ πεπλώκαμεν; \espTest \textbf{882} (ναυττιᾶν $=$ ναυσιᾶν) Moeris ν 9 \espTest \textbf{902} (στρέψον ἀνταυγεῖς κόρας) Suidas α 2617 \espTest (ἀνταυγεῖς κόρας) Photius α 2059 ἀνταυγεῖς κόρας \espTest \textbf{904} \Antiatt 83.9 ; cf. Hesychius α 8572 ἀφασία \espTest \textbf{910} Suidas ι 778 (sed Μενέλεων pro Μενέλαον, ὄσον pro ὅσα γ᾽ et ἰφύων pro ἀφύων) \espTest \textbf{910} Suidas ι 778 (sed Μενέλεων pro Μενέλαον et ἰφύων pro ἀφύων) \espTest \textbf{912} (ἐσχάρας) Photius ε 2044 ἐσχάρας $\subset$ Hesychius ε 6447 ἐσχάραι\par
}
\vspace{2.5pt}\noindent\rule{4.5cm}{0.4pt}\par\vspace{1.5pt}
{\fontsize{6.9pt}{8.2pt}\selectfont
\noindent\textbf{Suite (en commentaire) de l'apparat critique \& notes paléographiques (\texttt{f. 173r}) ---} \textbf{Contexte :} Suite et sommet de la \textbf{parodie d'\italique{Hélène}} : stichomythie de reconnaissance entre Euripide-Ménélas et Mnésiloque-Hélène sous le regard incrédule et furieux de la gardienne (\textbf{Kritylla}, vv. 876–919). --- \textbf{Deux émendations paléographiques majeures de \textbf{Bibi} (vv. 895 et 910) :} --- \textbf{v. 895 :} βαϋζα R (βάϋζε Zanetti) : \lecon{κόνυζα} \lecon{Bibi} (confusion minuscule $\kappa \to \beta$). --- \textbf{v. 910 :} Μενέλαον R (Μενέλεων Souda) : \lecon{Μενελέῳ σ'} Thiersch \espApp \textbf{ἀφύων R} (ἰφύων $\Sigma$, Souda ; ἀμφίων Grégoire) : \lecon{λαιφίων} \lecon{Bibi} ap. Hésychius $\lambda\ 171$ (chute de $\Lambda$ devant $\text{Α}$ en onciale).\par
}
\end{minipage}
\end{adjustbox}
\end{adjustwidth}
\newpage
\enTeteD{Codex Ravennas 429 (Biblioteca Classense) --- Fac-similé du f. 173r (vv. 876–919)}
\begin{adjustwidth}{-1.5cm}{-3.0cm}
\begin{center}
\includegraphics[width=16.5cm,height=20.5cm,keepaspectratio]{/workspaces/Thesmophories/dossier_ravennas_thesmophories/folios_individuels_ff162v-177r/22_f173r_classense175r_vv0876-0919.jpg}\par
\vspace{2pt}
{\scriptsize\textbf{Codex Ravennas 429, f. 173r (Classense 175r)} --- ΕΠΕΙΣΟΔΙΟΝ Δ' (Ἑλένη), \textbf{vv. 876–919} ($\text{R}^1, \Rpc, \Rdeux$ (\Mon)).}
\end{center}
\end{adjustwidth}
\newpage
\clearpage\ifodd\value{page}\hbox{}\newpage\fi
\enTeteG{ΘΕΣΜΟΦΟΡΙΑΖΟΥΣΑΙ --- Texte grec établi (f. 173r, vv. 876–919)}
\enTeteD{LES THESMOPHORIES --- Traduction française (f. 173r, vv. 876–919)}
\begin{pages}
\begin{Leftside}
\beginnumbering
%% v. 876+ (part 2)
\begin{groupe}
\trimIamb{\ligne{876}ἐπεὶ τέθνηκε Πρωτέας ἔτη δέκα.\\\passeLigne}%
\end{groupe}%
%% v. 877
\begin{groupe}%
\personnageGrec{Εὐριπίδης}%
\trimIamb{ποίαν δὲ χώραν εἰσεκέλσαμεν σκάφει;%\commentaire{877}{εἰσεκέλσαμεν}{\eff{lire un epekelsamen ; cf. Hésychius ? erreur inverse de l'autre ? attention à la scansion !}}
\\\passeLigne}%
\end{groupe}%
%% v. 878a
\begin{groupe}%
\personnageGrec{Μνησίλοχος}%
\trimIamb{Αἴγυπτον.\\\passeLigne}%
\end{groupe}%
%% v. 878b
\begin{groupe}%
\personnageGrec{Εὐριπίδης}%
\trimIamb{\ligne{878}\decalage{Αἴγυπτον.}ὦ δύστηνος οἷ πεπλώκαμεν.\testimonia{878}{Suidas οι 1 ὦ δύστηνος, οἷ πεπλώκαμεν;}%
\commentaire{878}{}{\italique{Hélène}, v. 461 :%
\citComm{Αἴγυπτος; ὦ δύστηνος, οἷ \difference{πέπλευκ᾽ ἄρα.}}}\\\passeLigne}%
\end{groupe}%
%% v. 879-880
\begin{groupe}%
\personnageGrec{Μνησίλοχος}%
\trimIamb{πείθει τι τῷ κακῷ κακῶς ἀπολουμένῳ\apparat{879}{τῷ R : τούτῳ τῷ Grynaeus %(1544) 
: \lecon{τῷ κακῷ} Scaliger \ligneApp{880} \lecon{Θεσμοφόριον} R, cf.278% : Θεσμοφορεῖον Meineke
}\\
ληροῦντι λῆρον; Θεσμοφόριον τουτογί.\\\passeLigne}%
\end{groupe}%
%% v. 881
\begin{groupe}%
\personnageGrec{Εὐριπίδης}%
\trimIamb{αὐτὸς δὲ Πρωτεὺς ἔνδον ἔστ᾽ ἢ 'ξώπιος;\apparat{881}{ἠξώπιος R, id. v.884 : \lecon{ἢ 'ξώπιος;} Biset}%
\commentaire{881}{}{\italique{Hélène}, v. 465 :%
\citComm{Μενελέως --- \difference{ποῦ δῆτ᾽ ἂν εἴη; πότερον ἐκτὸς ἢ 'ν δόμοις;}}}\\\passeLigne}%
\end{groupe}%
%% v. 882-884
\begin{groupe}%
\personnageGrec{Κρίτυλλα}%
\trimIamb{οὐκ ἔσθ᾽ ὅπως οὐ ναυτιᾷς ἔτ᾽ ὦ ξένε,\testimonia{882}{(ναυττιᾶν $=$ ναυσιᾶν) Moeris ν 9}\\
ὅστις γ᾽ ἀκούσας ὅτι τέθνηκε Πρωτέας\apparat{883}{ὅστις R : \lecon{ὅστις γ᾽} Scaliger}\\
ἔπειτ᾽ ἐρωτᾷς `ἔνδον ἔστ᾽ ἢ 'ξώπιος;'\\\passeLigne}%
\end{groupe}%
%% v. 885
\begin{groupe}%
\personnageGrec{Εὐριπίδης}%
\trimIamb{αἰαῖ τέθνηκε. ποῦ δ᾽ ἐτυμβεύθη τάφῳ;\apparat{885}{τέθνηκεν R : \lecon{τέθνηκε} Bentley}\\\passeLigne}%
\end{groupe}%
%% v. 886
\begin{groupe}%
\personnageGrec{Μνησίλοχος}%
\trimIamb{τόδ᾽ ἐστὶν αὐτοῦ σῆμ᾽, ἐφ᾽ ᾧ καθήμεθα.%
\commentaire{886}{}{\italique{Hélène}, v. 466 :%
\citComm{Γραῦς --- τόδ᾽ ἐστὶν αὐτοῦ \difference{μνῆμα, παῖς δ᾽ ἄρχει χθονός}.}}\\\passeLigne}%
\end{groupe}%
%% v. 887-888
\begin{groupe}%
\personnageGrec{Κρίτυλλα}%
\trimIamb{κακῶς τ᾽ ἄρ᾽ ἐξόλοιο κἀξολεῖ γέ τοι,\apparat{887}{\lecon{κακῶς τ᾽ ἄρ᾽} R : κακῶς γ᾽ ἄρ᾽ Denniston (\italique{CR} XLIV(1930), p.215)}\\
ὅστις γε τολμᾷς σῆμα τὸν βωμὸν καλεῖν.\\\passeLigne}%
\end{groupe}%
%% v. 889-890a
\begin{groupe}%
\ledpb%
\personnageGrec{Εὐριπίδης}%
\trimIamb{τί δὴ σὺ θάσσεις τάσδε τυμβήρεις ἕδρας\apparat{889}{τί δὲ σὺ R : τί δαὶ σὺ Faber : \lecon{τί δὴ σὺ} Scaliger}\\
φάρει καλυπτὸς ὦ ξένη;\\\passeLigne}%
\end{groupe}%
%% v. 890b-891
\begin{groupe}%
\personnageGrec{Μνησίλοχος}%
\trimIamb{\ligne{890}\decalage{φάρει καλυπτὸς ὦ ξένη;}βιάζομαι\\
γάμοισι Πρωτέως παιδὶ συμμεῖξαι λέχος.\apparat{891}{συμμῖξαι R : \lecon{συμμεῖξαι} van Leeuwen}\\\passeLigne}%
\end{groupe}%
%% v. 892-894
\begin{groupe}%
\personnageGrec{Κρίτυλλα}%
\trimIamb{τί ὦ κακόδαιμον ἐξαπατᾷς αὖ τὸν ξένον;\apparat{892}{αὐτὸν (ex αὖτὸν) \Rpc : \lecon{αὖ τὸν} Casaubon}\\
οὗτος πανουργῶν δεῦρ᾽ ἀνῆλθεν ὦ ξένε\\
ὡς τὰς γυναῖκας ἐπὶ κλοπῇ τοῦ χρυσίου.%
\commentaire{894}{ὡς τὰς γυναῖκας ἐπὶ κλοπῇ τοῦ χρυσίου,}{«~Subtiliser les bijoux des femmes~» ; voilà un chef d'accusation nouveau contre Mnésiloque d'autant plus étonnant que rien jusqu'ici ne l'a laissé présager. 
Le fait que Critylla porte un flambeau indique pourtant qu'elle est présente depuis l'arrivée de la procession. Ment-elle pour éviter d'avoir à raconter la trop longue affaire ? il pourrait y avoir une autre explication. %D'abord, l'on dispose de lois du \nieme{V} siècle av. J.-C. réglementant l'habillement lors des cultes à Déméter\marqueur{1} :
%\citComm{[Εἰ γυ]νὰ Fέσετοι ζτεραῖον λ\chap{ο}πος, | [ἰερὸ]ν \chap{ἐ}ναι ταῖ Δάματρι ταῖ Θεσμοφόροι· | [εἰ δὲ] μὲ ὑνιερόσει δυμενὲς ἔασα ἐπὲ Fέργο | [κακ\chap{ο}ς] ζ᾽ ἐξόλοιτυ, κὰ ὄζις τότε δαμιοFοργ\chap{ε} | [ἀφάε]σται δαρχμὰς τριάκοντα· εἰ δὲ μὲ ἀφάετοι, | [μεδ᾽ ἰν] τὰν ἀσέβειαν ἔχε ὄδε κῦρος, δέκο Fέτεα \chap{ἐ}να[ι | ἰναγὲς] τ\chap{ο}δε.\\
%Lorsqu'une femme met un vêtement de couleur\marqueur{2}, celui-ci doit être dédié à Déméter Thesmophore. Mais si elle ne le consacre pas, en étant hostile à la cérémonie, elle périra, et celui qui sera alors le fonctionnaire civil aura trente drachmes. Mais si ce dernier n'a pas l'autorité pour empêcher l'indignité, il sera souillé par elle pendant dix ans.}
Il semble que seules les femmes tenant un rôle actif lors de ces cérémonies pouvaient porter des bijoux (prêtresses, canéphores, \&c.) ; c'est ce qui ressort notamment d'un passage de \italique{Lysistrata} :
\citComm{στρωμάτων δὲ ποικίλων καὶ /
χλανιδίων καὶ ξυστίδων καὶ /
\textbf{χρυσίων}, ὅσ᾽ ἐστί μοι, /
οὐ φθόνος ἔνεστί μοι πᾶσι παρέχειν φέρειν /
τοῖς παισίν, \textbf{ὁπόταν} τε θυγάτηρ τινὶ \textbf{κανηφορῇ}. \\
Et de surtouts brodés, petites chlamides, xystides, \italique{bijoux}, tout ce que je possède, il ne me coûte pas d'offrir à chacun de l'emporter, pour ses enfants et pour sa fille, \italique{si elle devenait canéphore}.}
En fait, le terme χρυσίον, \italique{bijou}, \italique{ornement} semble avoir parfois un sens imagé chez Aristophane et désigner les \italique{appâts} féminins. Ainsi, lorsque Pisthétaire s'exclame à propos de la belle \& court vêtue joueuse de flûte Procné (\italique{Oiseaux}, v. 669) :%
\citComm{Πισθέταιρος --- ὅσον δ᾽ ἔχει τὸν χρυσόν, ὥσπερ παρθένος,}%
il est bien évident qu'il fait l'éloge de sa beauté, son galbe, ou que sais-je encore\marqueur{1} et non pas de sa verroterie. Ce blason sera d'ailleurs continué quelques vers plus loin par le chœur (vers 676-684). 
Autre exemple : lorsque Dicéopolis fait venir la jeune canéphore en lui prodiguant ce conseil (vers 257-258) :%
\citComm{πρόβαινε, κἀν τὤχλῳ φυλάττεσθαι σφόδρα\\
μή τις λαθών σου περιτράγῃ τὰ χρυσία.\\
\francais{Avance, et surtout prends garde que dans la foule\\
on ne te grignote tes bijoux.}}%
Certes, comme on l'a vu, les canéphores portaient des bijoux ; mais on devine aisément avec Aristophane qu'en voyant cette jeune beauté dont Dicéopolis vient de dire : « Heureux celui qui t'épousera et te procréera des belettes », la foule va ressentir quelque pulsion ne relevant pas uniquement de la cleptomanie.
Le verbe περιτρώγειν paraît d'ailleurs bien impropre pour une image de vol à la tire\marqueur{2}.
Quant à κλοπῇ, Mnésiloque nous a déjà fait remarquer au vers 205 que le verbe κλέπτειν pouvait par euphémisme désigner une activité qui n'a de commun avec le vol que de se passer au plus noir de la nuit.\\
Terminons en indiquant que cette accusation de tentative de viol lors d'une réunion exclusivement féminine se retrouve ailleurs ; ces assemblées féminines à huis-clos et que la rumeur publique voyait comme des séances de débauche sexuelle avaient tout pour attiser la curiosité de quelque luxurieux. C'est à une telle tentative d'un certain Critias qu'un fragment d'Hipponax\marqueur{3} fait allusion :%
\citComm{οὔ μοι δικαίως μοιχὸς ἁλῶναι δοκέει\\
Κριτίης ὁ Χῖος ἐν $ \cup - \times - $ δούμῳ\\
Ce n'est pas selon moi par justice\marqueur{4} qu'en tant que débauché fût saisi\\
Critias de Khios dans $[$ … $]$ une assemblée féminine.}}%
\debutNote{894}%
%\note{1}{R. Koerner, \italique{Inschriftliche Gesetztexte der Früher Grieschischen Polis} (1993), 35 $=$ \italique{Lois sacrées des cités grecques. Supplément} 32 $=$ \italique{Supplementum Epigraphicum Graecum} XI, 1112. Je reprends le texte établi par Koerner.}%
%\note{2}{Et pourtant Mnésiloque rejoint les Thesmophories avec sa crocote couleur safran !.. au milieu d'une féminine assemblée sobrement vêtue ? (de crocotes et de mitres tout-de-même, cf. le vers 941, (ἐ)ν κροκωτοῖς καὶ μίτραις ; mais la crocote était-elle nécessairement de la couleur du κρόκος dont elle tire son nom ?) 
%Quoi qu'il en soit, cette mise en scène me paraît assez cocasse pour être vraisemblable : Mnésiloque vêtu d'une robe chatoyante dénotant tout-à-fait avec la sobriété du reste de l'assemblée.}%
%\note{3}{Ce un passage de \italique{Lysistrata} (v. 1189-1193) : c'est le chœur des femmes qui parle :%
%\citComm{στρωμάτων δὲ ποικίλων καὶ /
%χλανιδίων καὶ ξυστίδων καὶ /
%\textbf{χρυσίων}, ὅσ᾽ ἐστί μοι, /
%οὐ φθόνος ἔνεστί μοι πᾶσι παρέχειν φέρειν /
%τοῖς παισίν, \textbf{ὁπόταν} τε θυγάτηρ τινὶ \textbf{κανηφορῇ}. /
%Et de surtouts brodés, petites chlamides, xystides, \italique{bijoux}, tout ce que je possède, il ne me coûte pas d'offrir à chacun de l'emporter, pour ses enfants et pour sa fille, \italique{si elle devenait canéphore}.}}%
%\note{4}{\italique{Oiseaux}, v. 669.}%
\note{1}{Hésychius (χ 792) : χρυσίον· τὸ τῶν παιδίων αἰδοῖον. Henderson (\italique{Maculate Muse}, p. 131) y flaire un jeu de mot entre τὸν χρυσόν et τὸν κυσόν.}%
%\note{6}{Vers 676-684.}%
\note{2}{Si impropre que l'on en serait tenté de comprendre λαβών à la place de λαθών.}%
\note{3}{Fragment 30 de l'édition au fort utile commentaire d'Olivier Masson ($=$ fragment 42 de l'édition Degani). Je reprends la correction proposée par Masson : δούμῳ à la place de δούλῳ, qui s'appuie sur une correction identique chez Hésychius (δ 2256) :%
\citComm{δοῦμος· ἡ οἰκία, ἢ τὴν ἐπὶ τὸ αὐτὸ συνέλευσιν τῶν γυναικῶν,\\
doumos : la maison. \italique{Item} : une réunion en commun des femmes.}}%
\note{4}{Je lis : οὔ μοι δίκᾳ ὡς μοιχὸς ἁλῶναι δοκέει κτλ.}%
%\note{1}{M. Dillon, \italique{Girls and Women in Classical Greek Religion}, 2004, p. 261-262}
\\\passeLigne}%
\end{groupe}%
%% v. 895
\begin{groupe}%
\personnageGrec{Μνησίλοχος}%
\trimIamb{κόνυζα τοὐμὸν σῶμα βάλλουσα ψόγῳ.\apparat{895}{βαϋζα R : βάϋζε Zanetti : \lecon{κόνυζα} Bibi}%
\commentaire{895}{βαϋζα τοὐμὸν σῶμα βάλλουσα ψόγῳ.}{Taillardat (§569) avait observé que le tour βάλλειν τινὰ ψόγῳ, « accabler quelqu'un de reproches » est d'un ton tragique fait sur le modèle de βάλλειν τινὰ λίθοις\marqueur{1}. (Et réciproquement, les mots sont bien souvent considérés comme des projectiles, cf. Taillardat §502.) 
De même, la périphrase τοὐμὸν σῶμα pour ἐγώ est d'un style élevé et particulièrement prisée d'Euripide\marqueur{2}.

La correction de la leçon du \italique{Ravennas} universellement admise, βαϋζα en βάϋζε, est inintelligible. (Les dictionnaires se sont d'ailleurs bien gardés de citer notre vers 895 à l'entrée βαΰζω.)
Faisons l'hypothèse que la fin du mot transmis soit correcte. Le suffixe -υζᾰ a été productif de quelques noms de végétaux en grec\marqueur{3} ; la \italique{conyze}, κόνυζα convient parfaitement : les minuscules `κ' et `β' sont très semblables sur de nombreux manuscrits\marqueur{4} -- dont notre \italique{Ravennas}. 
Le sens s'éclaire ainsi : Mnésiloque se compare à une plante -- une fleur dira-t-on en français -- en proie aux jets de reproches de Critylla. C'est là le point : ψόγῳ à la fin du vers est mis πάρα προσδοκίαν pour λίθῳ. Or le proverbe grec « jeter une pierre à une herbe » (Apostolius, XIV, 72) :%
\citComm{πόᾳ λίθον ἐκβάλλει· ἐπὶ τῶν τἀναντία τοῖς ἐπιβούλοις ποιούντων·\\
Il jette une pierre à une herbe\marqueur{5}.}%
%Élien : Aelian, De Natura Animalium I 45 : γυπῶν πτερὰ εἰ θυμιάσειέ τις, ὡς ἀκούω, καὶ ἐκ φωλεῶν καὶ ἐξ εἰλυῶν τοὺς ὄφεις προάξει ῥᾷστα. τὸ ζῷον ὁ δρυοκολάπτης ἐξ οὗ δρᾷ καὶ κέκληται. ἔχει μὲν γὰρ ῥάμφος ἐπίκυρτον, κολάπτει δὲ ἄρα τούτῳ τὰς δρῦς, καὶ ἐνταυθοῖ ὡς ἐς καλιὰν τοὺς νεοττοὺς ἐντίθησιν, οὐ δεηθεὶς καρφῶν καὶ τῆς ἐξ αὐτῶν πλοκῆς καὶ οἰκοδομίας οὐδὲ ἕν. οὐκοῦν εἴ τις λίθον ἐνθεὶς ἐπιφράξειε τῷ ὀρνέῳ τῷ προειρημένῳ τὴν ἔσδυσιν, ὃ δὲ συμβαλὼν τὴν ἐπιβουλὴν [p. 25] κομίζει πόαν ἐχθρὰν τῷ λίθῳ: ὃ δὲ οἷα βαρούμενος καὶ μὴ φέρων ἐξάλλεται, καὶ ἀνέῳγεν αὖθις τῷ προειρημένῳ ἡ φίλη ὑποδρομή.
%J’ai appris qu’en faisant des fumigations de plumes de vautour on parvient sans difficulté à faire sortir les serpents de leurs repaires ou de leurs tanières. L’animal qu’on nomme pivert doit son nom à ses façons 44 : il a un bec recourbé à son extrémité, dont il se sert pour becqueter les chênes, et c’est à l’intérieur de l’arbre qu’il loge ses petits comme si c’était un nid ; et il n’a nul besoin de se procurer des brindilles, de les tresser ensemble et de construire un nid. Or, si l’on obstrue l’ouverture que s’est ménagée notre oiseau en y mettant une pierre, il prend acte de la manœuvre et va chercher une herbe qui a un pouvoir antipathique sur la pierre, et il la place dessus. La pierre, accablée par ce poids et incapable de le supporter, s’éjecte, et voilà de nouveau ouvert à notre oiseau l’accès à son précieux abri.
%
% Las ! on ne détruit que ce que l'on remplace...%\\
%(3287) ἄλυζα· ἄλυπον Hésychius
%(3144) κνύζα· [κνυζηθμός S, ἀπὸ τοῦ κνυζᾶσθαι· ἐπὶ τῶν κυνῶν. καὶ] ἡ κόνυζα κατὰ συγκοπήν
%Σκόνυζαν: καὶ χωρὶς τοῦ σ κόνυζαν λέγουσι· Φερεκράτης· Ἀρκεῖ μία σκόνυζα καὶ θύμω δύο. Petite fleur ?
%Μουσουργοί· ψάλτριαι. αἱ δὲ μουσουργοὶ βάρβαροι ἦσαν γυναῖκες· ὄνομα αὐταῖς ἐπιχώριον βάρζα. καὶ τούτων αἱ μὲν αὐλοῦσιν, αἱ δὲ ψάλλουσι πενταχόρδῳ, αἱ δὲ ἑπταχόρδῳ ψαλτηρίῳ· αἱ αὐταὶ δὲ καὶ προσᾴδουσι τῷ ψαλμῷ. ὁ ψαλμὸς δὲ οὐ πρὸς τὴν ᾠδὴν ἔχει πεποιημένον τὸ μέλος καθάπερ παρ’ Ἕλλησιν, ἀλλ’ ἔστιν οἷος ἐνδόσιμος  εἶναι τῇ ᾠδῇ. καὶ λήγοντι μὲν οὖν τοῦ πότου βασιλεῖ καὶ ἡ ᾠδὴ ξυνεπαύετο· ὁπότε δὲ ἐγχέαι κελεύσειεν, αὖθις ἤρχοντο ᾄδειν καὶ ψάλλειν. ὁ δὲ Ἀρβαζάκιος μουσουργοῖς συνέζη τοσαύταις, ὅσας οὔτ’ ἐκεῖνος ἀριθμεῖν εἶχεν οὔτε οἱ προσήκοντες λογισταί. τῶν γὰρ στρατιωτῶν ᾔδεσαν τὸ πλῆθος· τῶν ἑταιρῶν δὲ καὶ τὸν ἀριθμὸν τὸν  ἐκ τῶν χειρῶν διέφυγεν. ὥσπερ οὖν Ὀρόντην τὸν Πέρσην φασὶν εἰπεῖν, ὅτι τῶν δακτύλων ὁ μικρότατος καὶ μυρία σημαίνει καὶ ἕνα ἀριθμόν, οὕτω κἀκεῖνοι τὰς ἑταίρας κατὰ μονάδα καὶ μυριάδα ἠρίθμουν. Suidas
%\eff{ Hésychius 352 βᾶϋ· εἶδος ἄνθους
%ζᾶ: μέγα.ἰσχυρόν.πολύ. ὅθενκαὶ ζάπλουτοςὁ μεγαλόπλουτος Cf. Photius zêta 12}
%Anthologie Palatine : τέχνας εἵνεκα σεῖο καὶ ἁ λίθος οἶδε βρυάζειν,
%Πραξίτελες: λῦσον, καὶ πάλι κωμάσομαι.
%νῦν δ᾽ ἡμῖν οὐ γῆρας ἔτ᾽ ἀδρανές, ἀλλ᾽ ὁ πεδητὰς
%Σειληνοῖς κώμων βάσκανός ἐστι λίθος.
%βρυάζει:πάλιν τρυφερῶςδιάκειται,θάλλει,εὐφραίνεται Synagoge cf Photius
%	βρύα: ἃ γίνεται μὲν ἐπὶ πετρῶν, λέγεται δὲ καὶ σφακχὸςἢ
%	Τσκαφίς.ἐπετίθετοδὲ ἁγνισμοῦχάριν
%	αἰσχρὰ καθʼ ἡμετέρης ἔβλυσε παρθενίης => BLUZA
%	βαυκαλίζω jouer à la prude
%	ὄβρυζος
%	βρῦ(ν) ζαλάω  βρῦν εἰπεῖν Ar.Nu.1382
%	βαυβίζω
%	Archil. 46
%	χορυζᾷ Τξαίνει.μωραίνει vere en zato avec haplographie ? ZATOTOEMON
%Le βαϋζα du \italique{Ravennas} a été universellement corrigé en βάϋζε, correction qui soulève plusieurs objections. \ \ \  1° \ \ \ Au niveau du sens : après les graves accusations que Critylla vient de porter à l'encontre de Mnésiloque-Hélène, on s'attend à une défense vigoureuse : iel vient d'être diffamé.e ! le simple reproche de forme : βάϋζε τοὐμὸν σῶμα, "tu jappes un peu fort après moi" nous paraît une défense un peu légère. \ \ \ 2° \ \ \ Si la confusion αι/ε, est courante (en raison d'une prononciation identique), ce vers serait l'unique exemple de la pièce d'une substitution finale de ε par α.\\
%Faisons l'hypothèse que la fin du mot transmis soit correcte. Le suffixe -υζᾰ a été productif de quelques noms de végétaux en grec\marqueur{ἄλυζα $=$ ἄλυπον ; κνύζα	 $=$ κόνυζα $=$ σκόνυζα, \italique{conyze} ou \italique{herbe aux puces} ou \italique{encensière} (\latin{erigeron viscosum L.}) ; ἡ μώλυζα, \italique{la tête d'ail} ; ἡ ὄρυζα, \italique{le riz}.} ; ἄλυζα concorderait avec la leçon du manuscrit. C'est chez Hésychius qu'on trouve le mot (α 3287) :%
%\citComm{ἄλυζα· ἄλυπον\\
%Alyza : Turbith blanc (\latin{Globularia alypum})}%
%Le sens s'éclaire ainsi : Mnésiloque se compare à une plante -- une fleur dirai-t-on en français -- en proie aux jets de reproches de reproches de Critylla. C'est là le point : ψόγῳ à la fin du vers est mis πάρα προσδοκίαν pour λίθῳ. Or il nous est justement parvenu un proverbe grec :%
%\citComm{πόᾳ λίθον ἐκβάλλει· ἐπὶ τῶν τἀναντία τοῖς ἐπιβούλοις ποιούντων}%
%\eff{Hésychius : ἄλυζα: ἄλυπον   et  κνύζα: beaucoup de noms de fleurs en -uza
%οὔτ᾽ἐκ χερὸς μεθέντα καρτερὸνλίθον on ne peut pas plus rattraper le mot lâché que la pierre (Eur.)
%Jeter une pierre à une fleur Πόᾳ λίϑον ἐκβάλλει: ἐπὶ τῶνταναντίατοῖςἐστιβούλοιςποιούντων. Ὃ γὰρ δρυοκολάπτης τὸ ζῷον,
%ἔχειμὲν ῥάμφοςἐπίκυρτον,κολάσιτεεδὲ ἄρα τούτῳτὰς
%δρῦς,καὶ ἐνταυϑοῖκοιλάναςτὸν τόπον»,ὡς δἰς καλλιὰ}
%
%Élien : Aelian, De Natura Animalium I 45 : γυπῶν πτερὰ εἰ θυμιάσειέ τις, ὡς ἀκούω, καὶ ἐκ φωλεῶν καὶ ἐξ εἰλυῶν τοὺς ὄφεις προάξει ῥᾷστα. τὸ ζῷον ὁ δρυοκολάπτης ἐξ οὗ δρᾷ καὶ κέκληται. ἔχει μὲν γὰρ ῥάμφος ἐπίκυρτον, κολάπτει δὲ ἄρα τούτῳ τὰς δρῦς, καὶ ἐνταυθοῖ ὡς ἐς καλιὰν τοὺς νεοττοὺς ἐντίθησιν, οὐ δεηθεὶς καρφῶν καὶ τῆς ἐξ αὐτῶν πλοκῆς καὶ οἰκοδομίας οὐδὲ ἕν. οὐκοῦν εἴ τις λίθον ἐνθεὶς ἐπιφράξειε τῷ ὀρνέῳ τῷ προειρημένῳ τὴν ἔσδυσιν, ὃ δὲ συμβαλὼν τὴν ἐπιβουλὴν [p. 25] κομίζει πόαν ἐχθρὰν τῷ λίθῳ: ὃ δὲ οἷα βαρούμενος καὶ μὴ φέρων ἐξάλλεται, καὶ ἀνέῳγεν αὖθις τῷ προειρημένῳ ἡ φίλη ὑποδρομή.
%J’ai appris qu’en faisant des fumigations de plumes de vautour on parvient sans difficulté à faire sortir les serpents de leurs repaires ou de leurs tanières. L’animal qu’on nomme pivert doit son nom à ses façons 44 : il a un bec recourbé à son extrémité, dont il se sert pour becqueter les chênes, et c’est à l’intérieur de l’arbre qu’il loge ses petits comme si c’était un nid ; et il n’a nul besoin de se procurer des brindilles, de les tresser ensemble et de construire un nid. Or, si l’on obstrue l’ouverture que s’est ménagée notre oiseau en y mettant une pierre, il prend acte de la manœuvre et va chercher une herbe qui a un pouvoir antipathique sur la pierre, et il la place dessus. La pierre, accablée par ce poids et incapable de le supporter, s’éjecte, et voilà de nouveau ouvert à notre oiseau l’accès à son précieux abri.
%\\
}%
\debutNote{895}%
\note{1}{Βάλλειν κακοῖς (Sophocle, \italique{Ajax} 1244), αἰτίᾳ (\italique{id}. \italique{Trach.} 940), φθόνῳ (Eur. \italique{Élect.} 902). Cf. ἀτιμίῃσιν ἰάλλειν (Homère, \italique{Odyssée} XIII, 142), λόγοις ἰάπτειν (Sophocle, \italique{Ajax} 501).}%
%\note{2}{ de ses \italique{Images d'Aristophane}.}%
\note{2}{Voir le \LSJ \ à l'entrée σῶμα II.}%
\note{3}{ἄλυζα $=$ ἄλυπον ; κνύζα	 $=$ κόνυζα $=$ σκόνυζα, \italique{conyze} ou \italique{herbe aux puces} ou \italique{encensière} (\latin{erigeron viscosum L.}) ; ἡ μώλυζα, \italique{la tête d'ail} ; ἡ ὄρυζα, \italique{le riz}.}%
\note{4}{Ex. : dans les \italique{Euménides} d'Eschyle, lire νεβρόν au lieu de νεκρόν v. 246.}%
\note{5}{L'explication qui suit de Michel Apostolius devrait plutôt nous conduire à traduire ainsi : « Il retire une pierre à l'aide d'une herbe. » Mais la tentative d'Apostolius de trouver une exégèse savante au proverbe est vouée à l'échec -- comme pour tout ce qui touche à la tradition orale. Le rapprochement que fait Apostolius avec un passage du précieux Élien (\italique{De Natura Animalium}, I, 45), pour astucieux qu'il soit, n'est pas convaincant.}\\\passeLigne}%
\end{groupe}%
%% v. 896
\begin{groupe}%
\personnageGrec{Εὐριπίδης}%
\trimIamb{ξένη τίς ἡ γραῦς ἡ κακορροθοῦσά σε;
\commentaire{896}{κακορροθοῦσά.}{Le vocabulaire est euripidéen -- du véritable Euripide s'entend. %\marqueur{1}. 
Le verbe κακορροθεῖν, substitut expressif de κακολογεῖν, est un terme qu'on retrouve exclusivement dans ses tragédies -- sauf reprise aristophanesque (\italique{Acharniens}, v. 576).}%
%\debutNote{896}%
%\note{1}{Aristophane possède un talent sans pareil pour nous croquer ses contemporains -- au moins ceux qui en valaient la peine à ses yeux. Pour s'en assurer, on peut relire le \italique{Banquet}, l'\italique{Apologie} et les autres pièces de Platon où Socrate apparaît pour constater avec amusement que les traits grossiers qu'Aristophane dessine, Platon les a esquissés discrètement. Mais lus avec attention, ses dialogues nous révèlent le même Socrate que celui des \italique{Nuées} : un va-nu-pieds à l'hygiène si douteuse qu'il lui faut des circonstances exceptionnelles pour prendre un bain ; un caractère irascible et violent ; menant une vie essentiellement recluse ; faisant montre d'un mépris affiché envers les «~faiseurs de comédies~» ; etc., etc. Oui : Aristophane était dans le vrai dans ses \italique{Nuées} ! et Socrate n'a été grand que dans son dernier combat pour affronter la rage plébéienne, dernier sursaut d'humeur d'Athènes la décadente qui ignorait alors qu'elle avait déjà sombré définitivement, coulée par sa démocratie aux relents impérialistes -- démocratie en rut qui avait au moins le mérite de donner quelque signe de vie contrairement à celles qu'on connaît. Si Socrate n'avait pas été héroïquement retranché du jeune Platon, il ne fait à mon sens aucun doute que le plus grand des Grecs aurait fini par s'éloigner du maître de sa jeunesse. Au lieu de devenir Platon, serait-il devenu Shakespeare ? y aurions-nous gagné ? Il ne se serait en tout cas certainement pas abîmé avec l'âge arrivant dans une métaphysique aporétique, et si détachée des réalités qu'elle révolta Aristote -- comment comprendre la boulimie de ce dernier pour les sciences naturelles, si ce n'est qu'il allait à rebours de son maître ? Ce très curieux phénomène psychologique, ce mysticisme de la soixantaine qui s'élève au-dessus de nos réalités bassement terrestres -- chez des hommes d'action (!) et d'une trempe exceptionnelle (!) -- a touché de la même manière notre Platon français, Pierre Boutang. À ma connaissance, la plus belle description qui en a été faite -- outre l'expérience personnelle qui ne se partage pas -- se trouve dans le chef-d'œuvre de Jean Boudou, \italique{Lo libre de Catòia}, avec le personnage du grand-père.}%
\\\passeLigne}%
\end{groupe}%
%% v. 897a
\begin{groupe}%
\personnageGrec{Μνησίλοχος}%
\trimIamb{αὕτη Θεονόη Πρωτέως.\\\passeLigne}%
\end{groupe}%
%% v. 897b-899a
\begin{groupe}%
\personnageGrec{Κρίτυλλα}%
\trimIamb{\ligne{897}\decalage{αὕτη Θεονόη Πρωτέως.}μὰ τὼ θεὼ\\
εἰ μὴ Κρίτυλλά γ᾽ Ἀντιθέου Γαργηττόθεν·
\commentaire{898}{εἰ μὴ Κρίτυλλά γ᾽ Ἀντιθέου Γαργηττόθεν·}{Je signale à l'attention d'un éventuel exégète parcourant ces pages que la réplique de Critylla trouve un écho dans un fragment de la comédie perdue d'Aristophane des Δαιταλεῖς, \italique{Convives}, transmise par Pollux (X, 158 $=$ fr. 226(217) Kassel-Austin) :%
\citComm{εἰ μὴ δικῶν τε γυργάθους ψηφισμάτων τε θωμοὺς\\
< φέροντες >}%
Le rapport est flagrant : même formule commençant par «~εἰ μὴ~» ; Κρίτ-υλλά, nom formé du suffixe -υλλα\marqueur{1} et du radical κρίτ- qui fait entendre le \italique{jugement} a le même champs sémantique que les δικῶν ; Γαργηττόθεν et γυργάθους sont très proches phonétiquement ainsi que -τιθέου (de Ἀντιθέου) et τε θωμοὺς.\\
Athénée (I, 4d) évoque ce même fragment d'Aristophane dans son premier livre ; ce qui nous permet d'en préciser le contexte. Dans un chapitre assez décousu\marqueur{2} consacré aux propos de table, on comprend qu'il est certainement question de sycophantes -- pique-assiettes et autres croque-lardons --, dont la vie est douce, bonne et grasse :%
\citComm{τὸν βίον, φησίν, εὐσταθεῖς· οὐκ ἐγχειρογάστορες, γυργάθους ψηφισμάτων φέροντες, Ἀριστοφάνης φησίν}%
%Avec ces faibles indices, on peut raisonnablement penser que -τιθέου et τε θωμοὺς font entendre un "τις θεός" ou "τις θεῶν" dans la locution originale. C'est bien cela : Critylla jure ses grands dieux et invoque la justice... On en arrive à la phase la plus délicate : à quoi renvoie Γαργηττόθεν ou γυργάθους ? (noter le suffixe -θεν de Γαργηττόθεν, sûrement mis à dessein.) 
%Or l'oreille grecque devait percevoir comme particulièrement proches les sons `γ' et `κ'\marqueur{\eff{Retrouver ; c'est dans Érasme avec des histoires de carcinos et de gargalos karkanos ?}}. On rétablit donc, d'après le sens et l'homophonie, la présence d'un κατώρθωσεν dans le proverbe détourné. \\
Comme l'oreille grecque percevait comme particulièrement proches les sons `γ' et `κ', faut-il entendre la locution grecque :%
\citComm{οὐκ ... εἰ μὴ τὴν δίκην τις θεῶν κατώρθωσεν,\\
À moins qu'un dieu ne rétablisse la justice,}%
conservée par Isocrate\marqueur{3} et auquel Aristophane fait vraisemblablement allusion en \italique{Cavaliers} 1360.
%\citComm{οὐκ ἄν ποτέ σοι γενέσθαι τηλικαῦτα τὸ μέγεθος, \textbf{εἰ μή τις θεῶν} αὐτὰ \textbf{συγκατώρθωσεν}\\
%\italique{Orationes}, V (Φίλιππος), 152.}%
%\citComm{οὐ \textbf{δίκαιός} ἐστιν μετέχειν \textbf{εἴ τι Τιμόθεος} πράττων \textbf{μὴ κατώρθωσεν}\\
%\italique{Orationes}, XV (Περὶ ἀντιδόσεως), 105.}%
%Chance : Isocrate semble avoir fait allusion à la même locution :%
%\citComm{οὐκ εἴ τις μιᾷ τύχῃ τηλικοῦτόν τι κατώρθωσεν\\
%\italique{Orationes}, XV (Περὶ ἀντιδόσεως), 128.}%
%Qu'en peut-on conclure ? pas grand-chose pour l'heure... %, si ce n'est qu'il devait exister un tour grec, s'en remettant à la justice divine ressemblant peu ou prou à : %
%Il est vraisemblable qu'Aristophane y fasse aussi allusion au vers 1360 des \italique{Cavaliers}, quand un synégore passé maître dans le chantage menace :%
%\citComm{Οὐκ ἔστιν ὑμῖν τοῖς δικασταῖς ἄλφιτα,\\
%εἰ μὴ καταγνώσεσθε ταύτην τὴν δίκην}%
%\citComm{οὐκ ... εἰ μὴ τὴν δίκην τις θεῶν κατώρθωσεν\\
%À moins qu'un dieu ne rétablisse la justice.}%
%Si nous avons vu juste, alors il est vraisemblable qu'Aristophane y fasse aussi allusion au vers 1360 des \italique{Cavaliers}, quand un synégore passé maître dans le chantage menace :%
%\citComm{Οὐκ ἔστιν ὑμῖν τοῖς δικασταῖς ἄλφιτα,\\
%εἰ μὴ καταγνώσεσθε ταύτην τὴν δίκην}%
%Je vous passe le relais, ô lecteur(s), et m'en remets à vous pour pousser plus avant la recherche. 
%Cavaliers 1360 εἰ μὴ καταγνώσεσθε ταύτην τὴν δίκην / dikê = trois(oboles) ?
%NONLa métaphore semble halieutique : ἐγχειρογάστορες, \italique{ceux qui vivent du travail de leurs mains}, est certainement à prendre au sens spécialisé de \italique{pêcheurs}\marqueur{Hésychius}. Le 
%εἰ μή τις θεῶν αὐτὰ συγκατώρθωσεν isocrate
%εἰ μὴ δικῶν τε γυργαθοὺς ψηφισμάτων τε θωμοὺς φέροντες
%et des Cavaliers 979
% Καίτοι πρεσβυτέρων τινῶν
%οἵων ἀργαλεωτάτων
%ἐν τῷ δείγματι τῶν δικῶν
%ἤκουσ’ ἀντιλεγόντων, 
%Athénée  :
%On devine que les ἐγχειρογάστορες au sens de pêcheurs, le gurgathos étant un diktuon
%τὸν βίον, φησίν, εὐσταθεῖς, οὐκ ἐγχειρογάστορες. γυργάθους ψηφισμάτων φέροντες, Ἀριστοφάνης φησίν
%	Ach 684 : οὐχ ὁρῶντες οὐδὲν εἰ μὴ τῆς δίκης τὴν ἠλύγην.
%	Cav 1360  «Οὐκ ἔστιν ὑμῖν τοῖς δικασταῖς ἄλφιτα,
%	εἰ μὴ καταγνώσεσθε ταύτην τὴν δίκην»
%Κρίτυλλά
%Aristophane emploie ailleurs le prénom de Critylla : au vers 324 de \italique{Lysistrata} :
%\citComm{Χορὸς Γυναικῶν --- πέτου πέτου Νικοδίκη,\\
%πρὶν ἐμπεπρῆσθαι Καλύκην\\
%τε καὶ Κρίτυλλαν περιφυσήτω\\
%\francais{$[$Demi-$]$chœur des femmes --- Vole, vole, Nicodicé,\\
%avant que, léchées de toutes parts par les flammes, Calycé\\
%et Critylla périssent.}}
%Le demi-ch\oe ur des hommes est en train de bouter le feu aux barricades féminines ; le demi-ch\oe ur des femmes sonne alors le tocsin et prend ses dispositions « avant que ne soient brûlées Calycé et Critylla ».\\
%Calycé est mise ici car son nom évoque le bois\marqueur{1}, bois qui ne va pas tarder à se consumer si les femmes ne réagissent promptement. Il en va de même pour Critylla : son nom est formé du suffixe -υλλα\marqueur{2} et du radical κρίτ- qui fait entendre le pistachier térébinthe (\latin{Pistacia Terebinthus}) : Hésychius 4154 :
%\citComm{κρίτανος\marqueur{3}· τέρμινθος\\
%\francais{Critanos : térébinthe.}}
%Le pistachier térébinthe -- bois dur et dense -- en font entre autres choses un bois de chauffage. C'est pour cette raison que le choeur des femmes l'évoque dans \italique{Lysistrata} et pour la même raison qu'Aristophane la baptise ainsi ici. Car Critylla est la manieuse de torche\marqueur{4} ; et elle a déjà montré à Mnésiloque de quel bois elle se chauffait avec sa torche qu'elle a reconverti en gourdin.
}%
\debutNote{898}%
%\note{1}{\eff{peuplier noir ? Daremberg ?}}
\note{1}{Le suffixe en -υλλα/ιλλα est courant pour les noms de femme : Κύδιλλα, Ἥρυλλα, ou Ξένυλλα qui est apparue plus haut \&c. (voir \italique{La formation des noms en grec ancien} de Chantraine, p. 252.)}%
\note{2}{Le premier livre, duquel le passage est tiré, ne nous est connu que par l'\italique{épitomé} des \italique{Deipnosophistes}, ce qui complexifie encore la délimitation et l'attribution des citations... Voir l'article de Benoît Louyest, «~L'épitomé du \italique{Banquet des sophistes} d'Athénée~», in \italique{Rursus} 8 (2012).}%
\note{3}{Isocrate, \italique{Orationes}, V (Φίλιππος), 152 ; XV (Περὶ ἀντιδόσεως), 105.}%
%\note{3}{Noter l'emplacement de l'accent : que Κρίτυλλα soit proparoxyton permet de rapprocher son nom de κρίτανος \eff{la ville de Critalla ?} tandis que d'autres noms comme Κιρτίας ou Κριτόλα sont à rapprocher du verbe κρίνειν.\eff{Et Cri'ton ?}}
%\note{4}{Voir les vers 853 et 858.}
\commentaire{898}{Γαργηττόθεν·}{La scholie :%
\citComm{ἀπὸ δήμου τῆς Αἰγηίδος φυλῆς, ἔνθα τὸ Εὐρυσθέως ἀπόκειται σῶμα ἐν τῷ Ἀρατηρίῳ. ἐπίρρημα δὲ τὸ Γαργηττόθεν.\\
D'un dème de la tribu Égéide, c'est là que repose le corps d'Eurysthée dans l'Arêtérion. D'où le nom de `Gargêttothen'.}%
}\\
σὺ δ᾽ εἶ πανοῦργος.\\\passeLigne}%
\end{groupe}%
%% v. 899b-901
\begin{groupe}%
\personnageGrec{Μνησίλοχος}%
\trimIamb{\ligne{899}\decalage{σὺ δ᾽ εἶ πανοῦργος.}ὁπόσα τοι βούλει λέγε.\\
οὐ γὰρ γαμοῦμαι σῷ κασιγνήτῳ ποτὲ\\
προδοῦσα Μενέλεων τὸν ἐμὸν ἐν Τροίᾳ πόσιν.\apparat{901}{Μενέλαον R : \lecon{Μενέλεων} Scaliger}%
%\commentaire{901}{προδοῦσα Μενέλαον τὸν ἐμὸν ἐν Τροίᾳ πόσιν.}{\italien{Staccato e allegro} : comptez donc le nombre de résolutions de ce vers. Μενέλεων à la place de Μενέλαον -- avec synalèphe, synérèse, synizèze, ou tout ce que vous voudrez -- est la solution la plus simple pour rétablir le mètre.}
\\\passeLigne}%
\end{groupe}%
%% v. 902
\begin{groupe}%
\personnageGrec{Εὐριπίδης}%
\trimIamb{γύναι τί εἶπας; στρέψον ἀνταυγεῖς κόρας.\testimonia{902}{(στρέψον ἀνταυγεῖς κόρας) Suidas α 2617 \espTest (ἀνταυγεῖς κόρας) Photius α 2059 ἀνταυγεῖς κόρας}\\\passeLigne}%
\end{groupe}%
%% v. 903
\begin{groupe}%
\personnageGrec{Μνησίλοχος}%
\trimIamb{αἰσχύνομαί σε τὰς γνάθους ὑβρισμένη.
\commentaire{903}{αἰσχύνομαί σε τὰς γνάθους ὑβρισμένη.}{Rappelons que lors de la séance de travestissement, le pauvre et anciennement barbu Mnésiloque s'était fait flamber et raser la barbe. Sur un fameux cratère\marqueur{1} représentant la scène d'égorgement de la `petite' de Mika, on le voit en porter de très visibles stigmates. Et le scholiaste de commenter ce vers :%
\citComm{ἐπειδὴ ξυρηθεὶς ἦν,\\
puisqu'il a été rasé.}%
}%
\debutNote{903}%
\note{1}{Le « cratère de Würzburg », Martin-von-Wagner-Museum der Universitàt Wurzburg, H 5697. Reproduit en annexe ; cf. la note aux vers 754-757.}\\\passeLigne}%
\end{groupe}%
%% v. 904-905
\begin{groupe}%
%
\personnageGrec{Εὐριπίδης}%
\trimIamb{τουτὶ τί ἔστιν; ἀφασία τίς τοί μ᾽ ἔχει.\testimonia{904}{\Antiatt 83.9 ; cf. Hesychius α 8572 ἀφασία}%
\commentaire{904-5}{}{\italique{Hélène}, v. 549 : %
« \difference{ἔκπληξιν ἡμῖν} ἀφασίαν \difference{τε προστίθης} » ; v. 557 : « τίς εἶ; τίν᾽ ὄψιν \difference{σήν}, γύναι, \difference{προσδέρκομαι}; » ; v. 72 : « ὦ θεοί, τίν᾽ \difference{εἶδον} ὄψιν; \difference{ἐχθίστην ὁρῶ} ».}\\
ὦ θεοὶ τίν᾽ ὄψιν εἰσορῶ; τίς εἶ γύναι;\\\passeLigne}%
\end{groupe}%
%% v. 906
\begin{groupe}%
\personnageGrec{Μνησίλοχος}%
\trimIamb{σὺ δ᾽ εἶ τίς; αὑτὸς γὰρ σὲ κἄμ᾽ ἔχει λόγος.\apparat{906}{αὐτὸς R : \lecon{αὑτὸς} Boissonade}%
\commentaire{906-12}{}{\italique{Hélène}, v. 558 sq. :%
\citComm{Ἑλένη --- σὺ δ᾽ εἶ τίς; αὑτὸς γὰρ σὲ κἄμ᾽ ἔχει λόγος.\\
\difference{Μενελέως --- οὐπώποτ᾽ εἶδον προσφερέστερον δέμας.\\
Ἑλένη --- ὦ θεοί· θεὸς γὰρ καὶ τὸ γιγνώσκειν φίλους.}\\
< Μενελέως --- Ἑλληνὶς εἶ τις ἢ ἐπιχωρία γυνή; > (Ce vers a été restitué à la tragédie d'après Aristophane.)\\
Ἑλένη --- Ἑλληνίς· ἀλλὰ καὶ τὸ σὸν θέλω μαθεῖν.\\
Μενελέως --- Ἑλένῃ σ᾽ ὁμοίαν δὴ μάλιστ᾽ εἶδον, γύναι.\\
Ἑλένη --- ἐγὼ δὲ Μενέλεῴ γε σέ· \difference{οὐδ᾽ ἔχω τί φῶ}.\\
Μενελέως --- ἔγνως γὰρ ὀρθῶς ἄνδρα δυστυχέστατον.\\
Ἑλένη --- ὦ χρόνιος ἐλθὼν σῆς δάμαρτος \difference{ἐς χέρας}.}}\\\passeLigne}%
\end{groupe}%
%% v. 907
\begin{groupe}%
\personnageGrec{Εὐριπίδης}%
\trimIamb{Ἑλληνὶς εἶ τις ἢ 'πιχωρία γυνή;\\\passeLigne}%
\end{groupe}%
%% v. 908
\begin{groupe}%
\personnageGrec{Μνησίλοχος}%
\trimIamb{Ἑλληνίς. ἀλλὰ καὶ τὸ σὸν θέλω μαθεῖν.\\\passeLigne}%
\end{groupe}%
%% v. 909
\begin{groupe}%
\personnageGrec{Εὐριπίδης}%
\trimIamb{Ἑλένῃ σ᾽ ὁμοίαν δὴ μάλιστ᾽ εἶδον γύναι.\apparat{909}{Ἐλένης R : \lecon{Ἑλένῃ σ᾽} Eur. \italique{Hel.}563 \espApp μάλισθ᾽ ἴδον R : \lecon{μάλιστ᾽ εἶδον} \italique{Hel}.563 \ligneApp{910} Μενέλαον R : Μενέλεων Suid. : \lecon{Μενελέῳ σ᾽} Thiersch \espApp ἀφύων R : ἰφύων Σ, Suid. : ἀμφίων Grégoire : \lecon{λαιφίων} Bibi}\\\passeLigne}%
\end{groupe}%
%% v. 910
\begin{groupe}%
\personnageGrec{Μνησίλοχος}%
\trimIamb{ἐγὼ δὲ Μενελέῳ σ᾽ ὅσα γ᾽ ἐκ τῶν λαιφίων.\testimonia{910}{Suidas ι 778 (sed Μενέλεων pro Μενέλαον, ὄσον pro ὅσα γ᾽ et ἰφύων pro ἀφύων)}%\apparat{910}{Μενέλαον R : Μενέλεων Suid. %ι 778 : \lecon{Μενελέῳ σ᾽} Thiersch %(cf. Eur. \italique{Hel}.564) \espApp \lecon{ὅσα γ᾽} R : ὄσον Suid. \espApp ἀφύων R : ἰφύων Σ, Suid. : ἀμφίων Grégoire : φυκίων Austin \& Olson : \lecon{λαιφίων} Bibi}%
\commentaire{910}{ἐγὼ δὲ Μενέλαον ὅσα γ᾽ ἐκ τῶν ἰφύων.}{Scène de reconnaissance. Lorsqu'elle avait reconnu son vieux mari face à elle, la véritable Hélène d'Euripide s'était exclamé (v. 564) :%
\citComm{Ἐγὼ δὲ Μενέλεῴ γε σέ· οὐδ' ἔχω τί φῶ.\\
Nul ne fut plus pareil à Ménélas. Que dire ?\\
(traduction d'Henri Grégoire.)}
Chez Aristophane, à quoi Hélène reconnaît-elle son naufragé de mari ? « Οὐδ' ἔχω τί φῶ » devient dans le manuscrit de Ravenne : ὅσα γ᾽ ἐκ τῶν ἀφύων, « d'après ton menu fretin ». 
On lit dans le lexique de la \italique{Souda} (ι 778) : ὄσον ἐκ τῶν ἰφύων, « à tes lavandes\marqueur{1} ». 
(C'est cette dernière version que commente le scholiaste, recopié en marge du manuscrit de Ravenne.)
Qu'il s'agisse de \italique{menu fretin} ou de \italique{lavande}, de telles leçons sont ineptes\marqueur{2} !
Colin Austin a bien résumé la situation\marqueur{3} :
« L'alpha d'ἀφύη étant bref, cette leçon est métriquement impossible. Quant à la ``lavande'' de la \italique{Souda}, le scholiaste y voit une allusion aux herbes de la mère d'Euripide, ce qui semble bien tiré par les cheveux\marqueur{4} ! Comme le note H. Grégoire\marqueur{5} « les herbes sont ici plus ridicules encore que les anchois », moins ridicules cependant que les \italique{oignons marins} de Gag***\marqueur{6} (τιφύων, « \italique{your squills}... \anglais{here a metaphor for male genitals} ») : cette saloperie-là\marqueur{7} se réfute d'elle-même. En fait la solution du problème se trouve au v. 935. Là Euripide est accusé d'être un ἱστιορράφος, \italique{un ravaudeur de voiles}. Fritzsche (p. 363)\marqueur{8} avait bien vu qu'Aristophane se moque ici de l'accoutrement de Ménélas dans l'\italique{Hélène} d'Euripide: tout au long de cette pièce il est fait mention des haillons, des lambeaux de voile dont le roi de Sparte est pitoyablement affublé, comme un gueux, après son naufrage en Égypte (v. 416 τὰς ἐμὰς δυσχλαινίας, 422 ναὸς ἐκβόλοις ἃ ἀμπίσχομαι, 544 sq. ἄγριος δέ τις / μορφὴν ὅδ᾽ ἐστίν, 554 στολήν γ᾽ ἄμορφον, 790 ὥσπερ πτωχὸς, 1204 ὡς ἐσθῆτι δυσμόρφῳ πρέπει, 1282 τῆς ἀχλαινίας, 1297 ἐσθῆτά τ᾽ ἐξάλλαξον). En particulier au ν. 423 il s'écrie: « Mes vêtements d'antan, mes habits somptueux, mes anciennes splendeurs, la vague m'a tout pris », πέπλους δὲ τοὺς πρὶν λαμπρά τ᾽ ἀμφιβλήματα / χλιδάς τε πόντος ἥρπασ᾽, et au ν. 1079 « les haillons qui m'affublent d'eux-mêmes parleront, en témoins du naufrage », καὶ μὴν τάδ᾽ ἀμφίβληστρα σώματος ῥάκη / ξυμμαρτυρήσει ναυτικῶν ἐρειπίων. Dans les \italique{Grenouilles} (ν. 842) Eschyle a bien raison d'appeler Euripide \italique{un faiseur de mendiants}, \italique{un ravaudeur de guenille}, καὶ πτωχοποιὲ καὶ ῥακιοσυρραπτάδη. Pour en revenir au v. 910 des « Thesmo », on est bien tenté de croire que c'est à ses \italique{hardes} que le pauvre naufragé est reconnu par sa femme. »\\ 
Grégoire\marqueur{9} proposa de lire : `ἐκ τῶν ἀμφίων' ; τὸ ἄμφιον étant une abréviation de ἀμφίεσμα, le \italique{vêtement}. Cette proposition expliquerait la présence des ἀφύων, par simple chute d'une lettre suivie d'iotacisme. En revanche, elle passe à l'as la leçon `ἰφύων' conservée par le scholiaste du manuscrit de Ravenne et par Suidas (ι 778), témoins tout aussi dignes de confiance.
(Re-)prends ton habit Grégoire ! Mais creusons son idée, qui est la bonne ; il nous suffit de trouver un synonyme de vêtement (avec, de préférence, une connotation péjorative) qui puisse rendre compte des deux leçons\marqueur{10}. 
Ce bon vieil Hésychius ne nous ayant jamais trahi, c'est encore à lui que nous avons posé la question. Cet affable vieillard n'a pas tardé à répondre à notre requête (λ 171) :%
\citComm{λαίφια· ῥάκη.\\
Laiphia : hardes.}%
Hésychius nous confirme l'existence d'un diminutif de λαῖφος, \italique{vêtement en lambeaux}, \italique{tissu}, \italique{voile}, ainsi que son sens.\\
Tentons à présent une explication paléographique aux ἀφύων et autres ἰφύων : si l'on fait la simple hypothèse de la disparition de la première lettre du mot λαιφίων\marqueur{11}, le lambda, le scribe se retrouvait alors devant l'inconnu αιφιων, qu'il lisait indifféremment – iotacisme oblige – αιφυων. En butte à ce mot qui n'existe pas, il pouvait alors raisonnablement écrire, selon l'inspiration du moment, ἀφύων ou ἰφύων.}%
\debutNote{910}%
\note{1}{%Dans un article passionnant sur quelques plantes apparaissant chez Aristophane, mêlant avec un rare bonheur de savantes études de textes et des connaissances botaniques très pointues, la regrettée 
Suzanne Amigues identifiait en fait l'ἴφυον avec la nielle des blés (\latin{Agrostemma githago L.}). Voir son article « Quelques légumes de disette chez Aristophane et Plutarque » paru en 1988 dans le \italique{Journal des Savants}, pp. 157-171.}%
\note{2}{Il se trouve néanmoins de nombreux défenseurs de la leçon « ἐκ τῶν ἰφύων », leçon qui fait beaucoup rire de savants hellénistes (?!) – décidément très bon public !%Nous avons expliqué, aux vers 455-456, les λάχανα ἄγρια, le σκάνδιξ et le verbe διασκανδικίζειν ; la présence de ces ἴφυα, elle, ne s'explique que par une mauvaise compréhension des quolibets des Comiques visant la mère d'Euripide.
}%
\note{3}{«~Observations critiques sur les \italique{Thesmophories} d'Aristophane~», in \italique{Dodone} 19.2 (1990), p. 27.}%
\note{4}{C'est pourtant cette leçon « bien tirée par les cheveux » qu'Austin fit imprimer dans son édition des \italique{Thesmophories} ; il faut croire que, ici comme en d'autres occasions, son acolyte Olson fut de mauvais conseil.}%
\note{5}{Tome V de l'édition d'Euripide aux Belles Lettres (1950), p. 65 n. 1.}%
\note{6}{« Une telle saloperie ?.. ou plutôt un Gag, non ? » C'est ce mot d'ordre que je reprends à Austin, modifiant au passage le nom – \latin{nomen omen} ! – du farceur qui va recevoir sa fessée, conformément à ce que réclame la justice !}%
\note{7}{Pour réussir à faire perdre son sang-froid au pourtant toujours modéré Colin Austin, il fallait effectivement que ce fut une belle salauderie… Nous aurions volontiers laissé la misérable hypo-foutaise « qui se réfute d'elle-même » de Gag*** dans un oubli bien mérité, malheureusement, il nous a été infligé de la voir exhumée et prise au sérieux dans plusieurs articles récents… %je crois donc de santé publique de l'autopsier une bonne fois pour toute et d'enterrer – définitivement, espérons le ! – cette charogne ! Ne nous attardons pas dans cette opération et allons droit au but : la saloperie de Gag*** présente à la fois le quadruple exploit : 
Enterrons une bonne fois pour toute -- espérons-le ! -- cette charogne : la saloperie de Gag*** présente à la fois le quadruple exploit \ \ \ 1°) \ de ne pas rendre compte de la leçon `ἀφύων' du témoin principal de la pièce : le manuscrit de Ravenne ; \ \ \ 2°) \ d'élucubrer sur un soi-disant sens obscène de τίφυον qui n'est attesté nulle part ; \ \ \ 3°) \ d'être intrinsèquement absurde : si Ménélas est réputé avoir bien peu d'ardeur (raison du départ de l'insatisfaite madame ?) ; comment se fait-il qu'il soit en pleine \italique{érection} et que ce soit à cela qu'Hélène le reconnaisse ? \ \ \ 4°) \ \italique{last but not least}, d'être anti-comique à souhait : un cocu impuissant n'est drôle \italique{que} parce qu'il est impuissant : représentez-le en pleine érection et vous gâterez à coup sûr la farce ! \ \ \ \
Reconnaissons au moins à Gag*** le mérite (pédagogique pour l'édification du lecteur) d'avoir réussi à concentrer en un seul passage les quatre vices cardinaux du \italique{misthophore} de la recherche, savant homme qui se paye – littéralement ! – de mots. Qu'y a-t-il d'étonnant à ce qu'un individu salarié pour palabrer palabre ? Rien, à priori, mais on leur saurait tout-de-même gré d'occuper leur jactance à autre chose que de salir Aristophane.}%
\note{8}{Voir aussi G. Lange, \italique{Quaest. in Ar. Thesm., Diss. Gott}. (1891), p. 8.}%
\note{9}{Au tome V de son édition d'Euripide aux Belles Lettres, déjà mentionnée. Mais son hypothèse est de douze ans plus ancienne : il l'avait présentée en anglais dans l'article « Modern Greek in Aristophanes » paru dans le premier numéro de \italique{The Link, A Review of Mediaeval and Modern Greek} (1938), pp. 16-20. Une anecdote qui montre comme le monde grec est petit : Henri Grégoire n'est autre que celui qui reprit la suite du grand Alphonse Willems à la chaire de philologie grecque de l'Université libre de Bruxelles en 1909, lorsque ce dernier dut prendre sa retraite.}%
\note{10}{Il y avait une autre possibilité de recherche : trouver un mot dont les multiples sens auraient pu être glosés à la fois par ἴφυον et ἀφύη ; les gloses explicatives ayant fini par prendre la place de la leçon originale (le cas s'est déjà vu). Creusant dans cette voie, Austin \& Olson avaient suggéré une émendation très ingénieuse : φυκίων. Émendation ingénieuse sur le plan paléographique, beaucoup moins sur le plan comique…}%
\note{11}{Erreur courante dans notre manuscrit de Ravenne. Citons pour exemple le vers 834 où le manuscrit comporte « δίδοσθαι \textbf{Τ}ηνίοισι » à la place de « δίδοσθαι \textbf{Στ}ηνίοισι ». Cette erreur est d'autant plus fréquente lorsque il y a deux lettres identiques consécutives, ainsi au vers 555 du manuscrit de Ravenne : « ὡς \textbf{τ}λεγγίδας λαβοῦσαι » en place de « ὡς \textbf{στ}λεγγίδας λαβοῦσαι ». 
Une telle faute s'explique aisément si l'on imagine le manuscrit copié en onciales et où les mots n'étaient pas même espacés. Dans le mot « λαιφίων », les deux premières lettres, le `λ' et le `α' sont de formes semblables en onciales. La `chute' du `λ' me paraît donc une hypothèse loin d'être invraisemblable.}\\\passeLigne}%
\end{groupe}%
%% v. 911
\begin{groupe}%
%
\personnageGrec{Εὐριπίδης}%
\trimIamb{ἔγνως ἄρ᾽ ὀρθῶς ἄνδρα δυστυχέστατον.\\\passeLigne}%
\end{groupe}%
%% v. 912-916a
\begin{groupe}%
\personnageGrec{Μνησίλοχος}%
\trimIamb{ὦ χρόνιος ἐλθὼν σῆς δάμαρτος ἐσχάρας\testimonia{912}{(ἐσχάρας) Photius ε 2044 ἐσχάρας $\subset$ Hesychius ε 6447 ἐσχάραι}%
\commentaire{912}{ἐλθὼν σῆς δάμαρτος ἐσχάρας / λαβέ.}{Le vers correspondant de l'\italique{Hélène} (le vers 566) finit par : ἐς χέρας ; Hélène reproche à Ménélas d'avoir trop tardé à rejoindre ses \italique{mains} (ses \italique{bras}, dira-t-on en français). Occasion par trop belle pour Aristophane de renouer avec la tradition populaire qui ne voit en Hélène qu'une pandorienne dévergondée, à l'origine de tant de maux : il fait donc prononcer à Mnésiloque `ἐσχάρας'. Le mot ἐσχάρα désigne proprement le \italique{foyer} ou le \italique{réchaud}\marqueur{1}. Ici comme ailleurs\marqueur{2}, le terme peut désigner le sexe de la femme. Jean Taillardat\marqueur{3} rapproche cette métaphore d'un passage de la \italique{Paix}\marqueur{4} où le mot ὀπτάνιον, \italique{cuisine} est pris comme synonyme de κύσθος.}
\debutNote{912}
\note{1}{\italique{Acharniens}, v. 888 ; \italique{Guêpes}, v. 938 ; \&c.}
\note{2}{\italique{Cavaliers} v. 1286 : κυκῶν τὰς ἐσχάρας avec le commentaire du scholiaste : τὰ χείλη τῶν γυναικείων αἰδοίων.}
\note{3}{\italique{Les images d'Aristophane}, §117.}
\note{4}{Vers 891.}}
\end{groupe}%
%% v. 913-915
\begin{groupe}%
\typeVers{2. dochm.}\diDoch{\ligne{914}λαβέ με λαβέ με πόσι, περίβαλε δὲ χέρας.\apparat{914}{περίβαλλε R : \lecon{περίβαλε} Biset}
\commentaire{913-5}{λαβέ με λαβέ με πόσι, / περίβαλε δὲ χέρας. / φέρε σὲ κύσω. ἄπαγέ μ᾽ ἄπαγ᾽ ἄπαγέ με.}{
%Ἄπαγε ξένον ἐν χειμῶνι· παροιμία ἐπὶ τῶν ὀχληρῶν.
Les quatre dochmiaques sont tous intégralement résolus (et Lætitia\marqueur{1} note la «~\anglais{correption }κύσ\lettreBreve{ω}. ἄπαγε \anglais{harsh and unusual in dochmiacs}~») ; signe d'une accélération -- voire une précipitation en l'occurrence -- du rythme.\\
Aristophane nous présente Hélène pour ce qu'elle est : une nymphomane. (Loin donc de la pseudo-réhabilitation d'Euripide, plus risible que crédible.)
Comment Hélène nous fait-elle entendre les besoins tyranniques qui la tiraillent ? par un procédé vieux comme le monde et qu'Aristophane emploie aussi au début des \italique{Cavaliers} (vers 26) : répéter un bout de phrase, en l'accélérant... et faire entendre ainsi un autre motif, bien différent de l'initial. 
La luronne détachait-elle le `πάγε με' de la fin de son exhortation afin de suggérer le verbe πήγνυμι, \italique{ficher} ?\\
La source d'inspiration : \italique{Hélène}, v. 625 sq. :%u
\citComm{Ἑλένη --- ὦ φίλτατ᾽ ἀνδρῶν Μενέλεως, ὁ μὲν χρόνος\\
παλαιός, ἡ δὲ τέρψις ἀρτίως πάρα.\\
ἔλαβον ἀσμένα πόσιν ἐμόν, φίλαι,\\
περί τ᾽ ἐπέτασα χέρα\\
φίλιον ἐν μακρᾷ φλογὶ φαεσφόρῳ.\\
Μενελέως --- κἀγὼ σέ· πολλοὺς δ᾽ ἐν μέσῳ λόγους ἔχων\\
οὐκ οἶδ᾽ ὁποίου πρῶτον ἄρξωμαι τὰ νῦν.\\
Ἑλένη --- γέγηθα, κρατὶ δ᾽ ὀρθίους ἐθείρας\\
ἀνεπτέρωκα καὶ δάκρυ σταλάσσω,\\
περὶ δὲ γυῖα χέρας ἔβαλον, ἡδονάν,\\
ὦ πόσις, ὡς λάβω.}%
}%
\debutNote{913-5}%
\note{1}{\italique{The Songs of Aristophanes}, p. 429.}\\
φέρε σὲ κύσω. ἄπαγέ μ᾽ ἄπαγ᾽ ἄπαγέ με%
%\commentaire{914}{ἄπαγέ μ᾽ ἄπαγ᾽ ἄπαγέ με.}{Il n'est pas impossible qu'Aristophane recoure au même procédé }%
}
\end{groupe}%
%% v. 912-916a
\begin{groupe}%
\typeVers{trim. iamb.}\trimIamb{λαβὼν ταχὺ πάνυ.\\\passeLigne}%
\end{groupe}%
%% v. 916b-917
\begin{groupe}%
\personnageGrec{Κρίτυλλα}%
\trimIamb{\ligne{916}\decalage{λαβὼν ταχὺ πάνυ.}κλαύσετ᾽ ἄρα νὴ τὼ θεὼ\\
ὅστις σ᾽ ἀπάξει τυπτόμενος τῇ λαμπάδι.%\commentaire{917}{τῇ λαμπάδι.}{ \eff{Scholie} : didascalie d'Aristophane ?}
\\\passeLigne}%
\end{groupe}%
%% v. 918-919
\begin{groupe}%
\personnageGrec{Εὐριπίδης}%
\trimIamb{σὺ τὴν ἐμὴν γυναῖκα κωλύεις ἐμέ,\\
τὴν Τυνδάρειον παῖδ᾽, ἐπὶ Σπάρτην ἄγειν;\\\passeLigne}%
\end{groupe}%
\endnumbering
\end{Leftside}
\begin{Rightside}
\beginnumbering
%% v. 876+ (part 2)
\begin{groupe}
\trimIambFr{\ligne{876}Il est mort, Protéas, depuis dix ans.\\\passeLigne}%
\end{groupe}%
%% v. 877
\begin{groupe}%
\personnageFr{Euripide}%
\trimIambFr{Et quel est ce pays où ma nef a pris terre ?\\\passeLigne}%
\end{groupe}%
%% v. 878a
\begin{groupe}%
\personnageFr{Mnésiloque}%
\trimIambFr{L'Égypte.\\\passeLigne}%
\end{groupe}%
%% v. 878b
\begin{groupe}%
\personnageFr{Euripide}%
\trimIambFr{\ligne{878}\espace{1.7cm}« Vers quel lieu, malheureux, nous avons \fin navigué ! »\\\passeLigne}%
\end{groupe}%
%% v. 879-880
\begin{groupe}%
\personnageFr{Kritylla}%
\trimIambFr{T'en fies-tu à ce coquin — la peste soit de lui —
qui te raconte des bourdes ? Ceci, c'est le Thesmophorion.\\\passeLigne}%
\end{groupe}%
%% v. 881
\begin{groupe}%
\personnageFr{Euripide}%
\trimIambFr{Et lui, Protée, est-il céans, ou s'il est hors ?\\\passeLigne}%
\end{groupe}%
%% v. 882-884
\begin{groupe}%
\personnageFr{Kritylla}%
\trimIambFr{Il faut que tu aies encore mal au c\oe ur, étranger.\\
On te dit qu'il est mort, Protéas,\\
et tu demandes s'il est céans ou dehors.\\\passeLigne}%
\end{groupe}%
%% v. 885
\begin{groupe}%
\personnageFr{Euripide}%
\trimIambFr{Hélas! il a vécu. Où donc est le sépulcre où il fut \fin inhumé ?\\\passeLigne}%
\end{groupe}%
%% v. 886
\begin{groupe}%
\personnageFr{Mnésiloque}%
\trimIambFr{C'est ici son tombeau, tu nous y vois assise.\\\passeLigne}%
\end{groupe}%
%% v. 887-888
\begin{groupe}%
\lednopb%
\personnageFr{Kritylla}%
\trimIambFr{Puisses-tu périr misérablement ! Et tu périras,\\
toi qui oses appeler l'autel une tombe.\\\passeLigne}%
\end{groupe}%
%% v. 889-890a
\begin{groupe}%
\personnageFr{Euripide}%
\trimIambFr{D'où vient qu'assise sur ce siège sépulcral,\\
tu dérobes tes traits sous un voile, étrangère ?\\\passeLigne}%
\end{groupe}%
%% v. 890b-891
\begin{groupe}%
\personnageFr{Mnésiloque}%
\trimIambFr{\espace{3.8cm}On prétend sans mon v\oe u\\
que le fils de Protée me reçoive en son lit.\\\passeLigne}%
\end{groupe}%
%% v. 892-894
\begin{groupe}%
\personnageFr{Kritylla}%
\trimIambFr{Pourquoi, misérable, tromper encore cet étranger ?\\
Il est entré en fraude, ô étranger,
auprès des femmes afin de toucher subrepticement à leur bijou.\\\passeLigne}%
\end{groupe}%
%% v. 895
\begin{groupe}%
\personnageFr{Mnésiloque}%
\trimIambFr{Je suis la conyze qu'on accable d'un blâme.\\\passeLigne}%
\end{groupe}%
%% v. 896
\begin{groupe}%
\personnageFr{Euripide}%
\trimIambFr{Étrangère, dis-moi, qui donc est cette vieille, qui t'incrimine \fin ainsi ?\\\passeLigne}%
\end{groupe}%
%% v. 897a
\begin{groupe}%
\personnageFr{Mnésiloque}%
\trimIambFr{Elle ? C'est Théonoé, la fille de Protée.\\\passeLigne}%
\end{groupe}%
%% v. 897b-899a
\begin{groupe}%
\personnageFr{Kritylla}%
\trimIambFr{\ligne{897}\espace{3.6cm}Non, par les déesses,
je ne suis autre que Critylla, fille d'Antithéos, du dème de Gargettos.\\
Toi, tu es un fripon.\\\passeLigne}%
\end{groupe}%
%% v. 899b-901
\begin{groupe}%
\personnageFr{Mnésiloque}%
\trimIambFr{\ligne{899}\espace{3cm}Dis ce que tu voudras.\\
En épousant ton frère, je ne trahirai point
Ménélas, mon époux, qui reste à Troie.\\\passeLigne}%
\end{groupe}%
%% v. 902
\begin{groupe}%
\personnageFr{Euripide}%
\trimIambFr{Qu'as-tu dit ? Femme, tourne vers moi ta brillante \fin prunelle.\\\passeLigne}%
\end{groupe}%
%% v. 903
\begin{groupe}%
\personnageFr{Mnésiloque}%
\trimIambFr{J'ai honte, après l'outrage imprimé sur mes joues.%\commentaire{903}{}{On n'oublie pas qu'il a perdu sa barbe.}.
\\\passeLigne}%
\end{groupe}%
%% v. 904-905
\begin{groupe}%
\personnageFr{Euripide}%
\trimIambFr{Qu'est-ce là ? Une stupeur me prend.\\
Ô dieux, que vois-je ? Femme, qui donc es-tu?\\\passeLigne}%
\end{groupe}%
%% v. 906
\begin{groupe}%
\personnageFr{Mnésiloque}%
\trimIambFr{Qui donc es-tu toi-même ? « C'est une question qui \fin m'appartient aussi ».%
\commentaire{906}{« C'est une question qui m'appartient aussi ».}{J'emprunte ce vers à A. de Musset, \italique{À quoi rêvent...}, acte II, scène 3. (Note de Willems.)}\\\passeLigne}%
\end{groupe}%
%% v. 907
\begin{groupe}%
\personnageFr{Euripide}%
\trimIambFr{« Es-tu Hellène, ou bien es-tu de ce pays ? »\\\passeLigne}%
\end{groupe}%
%% v. 908
\begin{groupe}%
\personnageFr{Mnésiloque}%
\trimIambFr{« Hellène. Et à mon tour je veux savoir ta race. »\\\passeLigne}%
\end{groupe}%
%% v. 909
\begin{groupe}%
\personnageFr{Euripide}%
\trimIambFr{« Ô femme, je te vois tout pareille à Hélène. »\\\passeLigne}%
\end{groupe}%
%% v. 910
\begin{groupe}%
\personnageFr{Mnésiloque}%
\trimIambFr{« Toi-même à Ménélas, si j'en juge par ces hardes.%\commentaire{}{}{Le texte est visiblement corrompu. Porson, et bien d'autres après lui, lisent σέ τῷ 'κ ἰφύων, « le Ménélas du marché aux légumes », plaisanterie qui, même dans une parodie, paraît par trop intempestive surtout venant de Mnésiloque. Aussi Droysen la met-il dans la bouche de Critylla.}. 
»\\\passeLigne}%
\end{groupe}%
%% v. 911
\begin{groupe}%
\personnageFr{Euripide}%
\trimIambFr{« Oui, tu l'as reconnu, ce grand infortuné. »\\\passeLigne}%
\end{groupe}%
%% v. 912-916a
\begin{groupe}%
\personnageFr{Mnésiloque}%
\trimIambFr{« Ô toi, le tard venu dans les bas de ta femme »,}
\end{groupe}%
%% v. 912-916a
\begin{groupe}%
\diDochFr{prends, prends-moi, mon époux. Presse-moi dans tes bras.
Voyons, que je te baise. Emmène-moi, emmène, emmène-moi, et au plus vite.}
\end{groupe}%
%% v. 912-916a
\begin{groupe}%
\trimIambFr{(Il descend de l'autel, les femmes l'entourent.)\\\passeLigne}%
\end{groupe}%
%% v. 916b-917
\begin{groupe}%
\personnageFr{Kritylla}%
\trimIambFr{\ligne{916}\espace{2.8cm}Ah! il pleurera, par les deux déesses,\\
celui qui t'emmènera. Je le frapperai de ma torche.\\\passeLigne}%
\end{groupe}%
%% v. 918-919
\begin{groupe}%
\personnageFr{Euripide}%
\trimIambFr{Quoi ? mon épouse à moi, fille de Tyndare,\\
tu voudrais m'empêcher de la conduire à Sparte ?\\\passeLigne}%
\end{groupe}%
\endnumbering
\end{Rightside}
\end{pages}
\Pages
\clearpage\ifodd\value{page}\hbox{}\newpage\fi
\enTeteG{ΕΠΕΙΣΟΔΙΟΝ Δ' / ΧΟΡΙΚΟΝ --- Codex Ravennas 429, f. 173v (Classense 175v) --- vv. 920–958}
\begin{adjustwidth}{-1.5cm}{-3.0cm}
\begin{adjustbox}{max totalheight=21.0cm, max width=16.5cm, keepaspectratio}
\begin{minipage}{16.5cm}
{\fontsize{7.2pt}{8.5pt}\selectfont
\noindent\textbf{Folio 23 --- Codex Ravennas 429, \texttt{f. 173v} (\texttt{Classense 175v}) : ΕΠΕΙΣΟΔΙΟΝ Δ' / ΧΟΡΙΚΟΝ, vv. 920–958}\hfill \textit{Témoins :} $\text{R}^1, \Rpc, \Rdeux$, \PsiDeux\par\vspace{2pt}
\hrule height 0.4pt\vspace{2.5pt}
}
{\fontsize{8.1pt}{9.5pt}\selectfont
\ligneRav{920}{\auDessus{γυ}{ν}}{οἴμ᾽ ὡς πανοῦργος καὐτὸς εἶναί μοι δοκεῖς}{}
\ligneRav{}{}{καὶ τοῦδέ τις ξύμβουλος. οὐκ ἐτὸς πάλαι}{}
\ligneRav{}{}{ᾐγυπτιάζετ᾽. ἀλλ᾽ ὅδε μὲν δώσει δίκην·}{}
\ligneRav{}{}{προσέρχεται γὰρ ὁ πρύτανις χὠ τοξότης·}{}
\ligneRav{}{\auDessus{εὐ}{ρ}}{τουτὶ πονηρόν· ἀλλ᾽ ὑπαποκινητέον·}{}
\ligneRav{925}{μ\auDessus{ν}{Η}}{ἐγὼ δ᾽ ὁ κακοδαίμων τί δρῶ; : μέν᾽ ἥσυχος·}{}
\ligneRav{}{}{οὐ γὰρ προδώσω σ᾽ οὐδέποτ᾽, ἤνπερ ἐμπνέω,}{}
\ligneRav{927}{}{ἢν μὴ \difference{προλίπωσ᾽} αἱ μυρίαι με μηχαναί·}{\textbf{927} προλείπωσ᾽ R : \lecon{προλίπωσ᾽} Scaliger}
\ligneRav{928}{μ\auDessus{ν}{Η}}{αὕτη μὲν ἡ μήρινθος οὐδὲν ἔσπασεν·}{\textbf{928} (Γυ.) R : \lecon{(Μνη.)} Scaliger}
\ligneRav{929}{πρυτ/}{ὅδ᾽ \difference{ἔσθ᾽} ὁ πανοῦργος ὃν \difference{ἔλεγ᾽} ἡμῖν Κλεισθένης;}{\textbf{929} ἔστ᾽ R : \lecon{ἔσθ᾽} \Mon \espApp ἔλεγεν R : \lecon{ἔλεγ᾽} \Mon}
\ligneRav{930}{}{οὗτος τί κύπτεις; δῆσον αὐτὸν εἰσάγων}{}
\ligneRav{}{}{ὦ τοξότ᾽ ἐν τῇ σανίδι, κἄπειτ᾽ ἐνθαδὶ}{}
\ligneRav{}{}{στήσας φύλαττε καὶ προσιέναι μηδένα}{}
\ligneRav{}{}{ἔα πρὸς αὐτόν, ἀλλὰ τὴν μάστιγ᾽ ἔχων}{}
\ligneRav{934}{—}{παῖ᾽ ἢν προσίῃ τις. : νὴ Δί᾽ ὡς \difference{νῦν δή γ᾽} ἀνὴρ}{\textbf{934} νῦν δῆτ' R : \lecon{νῦν δή γ᾽} Dobree : νυνδή γ' Cobet}
\ligneRav{935}{}{ὀλίγου μ᾽ ἀφείλετ᾽ αὐτὸν ἱστιορράφος·}{}
\ligneRav{}{μ\auDessus{ν}{Η}}{ὦ πρύτανι πρὸς τῆς δεξιᾶς, ἥνπερ φιλεῖς}{}
\ligneRav{}{}{κοίλην προτείνειν ἀργύριον ἤν τις διδῷ,}{}
\ligneRav{}{}{χάρισαι βραχύ τί μοι καίπερ ἀποθανουμένῳ·}{}
\ligneRav{939}{πρυτ/}{τί σοι \difference{χαρίσωμαι}; : γυμνὸν ἀποδύσαντά με,}{\textbf{939} χαρίσομαι R : \lecon{χαρίσωμαι} Casaubon}
\ligneRav{940}{}{κέλευε πρὸς τῇ σανίδι δεῖν τὸν τοξότην,}{}
\ligneRav{941}{}{ἵνα \difference{μὴ 'ν} κροκωτοῖς καὶ μίτραις γέρων ἀνὴρ}{\textbf{941} μὴν R : \lecon{μὴ 'ν} Küster}
\ligneRav{942}{}{μὴ \inRas{γε}{} τοῖς κροκωτοῖς καὶ τῇ μίτρᾳ}{\textbf{942} \inRas{γε}{} τοῖς \Rdeux}
\ligneRav{943}{πρυτ/}{ἔχοντα ταῦτ᾽ \difference{ἔδοξε} τῇ βουλῇ σε δεῖν,}{\textbf{943} ἔδοξεν R : \lecon{ἔδοξε} Biset}
\ligneRav{944}{}{ἵνα τοῖς παροῦσι δῆλος ᾖς πανοῦργος ὤν·}{\textbf{944} \lecon{παροῦσι} R : παριοῦσι Brunck}
\ligneRav{945}{μ\auDessus{ν}{Η}}{ἰαππαπαιάξ· ὦ κροκώθ᾽ οἷ᾽ εἴργασαι·}{\textbf{945} \lecon{ἰαππαπαιάξ} R, Suid. ε 2807 : ἰατταταιάξ Biset}
\ligneRav{946}{}{κοὐκ \difference{ἔστ᾽} ἔτ᾽ ἐλπὶς οὐδεμία σωτηρίας·}{\textbf{946} ἔστιν R : \lecon{ἔστ᾽} Bentley}
\ligneRav{947}{χορ/}{ἄγε νυν ἡμεῖς \difference{παίσωμεν} ἅπερ νόμος ἐνθάδε \difference{ταῖσι} γυναιξίν,}{\textbf{947} πέσωμεν R : \lecon{παίσωμεν} Zanetti, cf. Σ \espApp ταῖσ R : \lecon{ταῖσι} Daubuz}
\ligneRav{948}{}{παραχώρει ὅταν ὄργια σεμνὰ θεαῖν ἱεραῖς ὥραις ἀνέχωμεν, ἅπερ καὶ}{\textbf{948} \lecon{θεαῖν} R : θεοῖν Cobet}
\ligneRav{}{}{\hspace{1.1em}Παύσων σέβεται καὶ νηστεύει,}{}
\ligneRav{950}{}{\hspace{1.1em}πολλάκις αὐταῖν ἐκ τῶν ὡρῶν}{\textbf{950} \lecon{αὐταῖν} R : αὐτοῖν Cobet}
\ligneRav{}{}{\hspace{1.1em}ἐς τὰς ὥρας ξυνεπευχόμενος}{}
\ligneRav{952}{}{\hspace{1.1em}τοιαῦτα \difference{μέλειν} θάμ᾽ ἑαυτῷ·}{\textbf{952} μέλλειν R : \lecon{μέλειν} Zanetti (\italique{curae esse} Divus) : μολεῖν Casaubon}
\ligneRav{}{χορ/}{\hspace{1.1em}ὅρμα χώρει·}{}
\ligneRav{}{}{\hspace{1.1em}κοῦφα ποσὶν ἄγ᾽ ἐς κύκλον,}{}
\ligneRav{955}{}{\hspace{1.1em}χειρὶ σύναπτε χεῖρα,}{\textbf{955} \lecon{χειρὶ} R : χερὶ Dobree}
\ligneRav{956}{}{\hspace{1.1em}†ῥυθμὸν χορείας ὕπαγε πᾶσα†·}{\textbf{956} \lecon{ὕπαγε πᾶσα} R : πᾶσ' ὕπαγε Austin}
\ligneRav{956}{}{\hspace{1.1em}βαῖνε καρπαλίμοιν ποδοῖν·}{\textbf{956} \lecon{ὕπαγε πᾶσα} R : πᾶσ' ὕπαγε Austin}
\ligneRav{}{}{\hspace{1.1em}ἐπισκοπεῖν δὲ πανταχῇ}{}
\ligneRav{}{}{\hspace{1.1em}κυκλοῦσαν ὄμμα χρὴ χοροῦ κατάστασιν·}{}
}
\vspace{3pt}\noindent\rule{4.5cm}{0.4pt}\par\vspace{1.5pt}
{\fontsize{7.0pt}{8.3pt}\selectfont
\noindent\textbf{Loci paralleli (vv. 920–958) ---} \textbf{922} (ᾐγυπτιάζετ᾽) \SynB α 494 αἰγυπτιάζειν $=$ Photius α 518 $\subset$ Suidas αι 75 \espTest \textbf{926} (ἤνπερ ἐμπνέω) Suidas η 379 (ἤν pro ἤνπερ) \espTest \textbf{928} Suidas α 4475, μ 978 ; \Eustathe 760.44 ; scholium\expo{RV} \italique{Vespae} 175b \espTest \textbf{931} (σανίδι) Photius σ 68 σανίς, σ 65 σανίδα ; Hesychius σ 169 σανίς ; \AelDion σ 7 \espTest \textbf{935} Suidas ι 691 ὀλίγου ταύτην ἱστιορράφος \espTest (ἱστιορράφος) Pollux VII 161 \espTest \textbf{940} (πρὸς τῇ σανίδι δεῖν) Hesychius δ 926 δῆσαι πρὸς σανίδι· Ἀττικῶς· βασανίσαι ἐν τῷ δεσμωτηρίῳ \espTest \textbf{942} Suidas ε 3215 \espTest (ἑστιῶν) Hesychius ε 6411 ἑστιῶν· ἀριστοποιῶν $=$ Photius ε 2031 $=$ \Syn ε 875 $=$ Suidas ε 3227 ; cf. \Phryn 66.7 \espTest \textbf{945} (ἰαππαπαιάξ) Suidas ε 2807 \espTest \textbf{947} (ὅταν -- ἀνέχωμεν) Suidas α 2364 (θεαῖν codd. GFS\expo{a.c.} vel θεῶν codd. ATS\expo{p.c} vel θεεῖν cod. I vel θεαῖνων cod. M pro θεοῖν)\par
}
\vspace{2.5pt}\noindent\rule{4.5cm}{0.4pt}\par\vspace{1.5pt}
{\fontsize{6.9pt}{8.2pt}\selectfont
\noindent\textbf{Suite (en commentaire) de l'apparat critique \& notes paléographiques (\texttt{f. 173v}) ---} \textbf{Témoin papyrologique :} \textbf{\PsiDeux{} (\italique{P.Oxy.} LVI 3839, fr. 2)} couvre les \textbf{vv. 941–953}. --- \textbf{Grattage \italique{in rasura} au v. 942 :} \Rac{} avait inséré la particule γε devant τοῖς ; \textbf{\Rdeux{} l'a grattée (\inRas{γε}{} τοῖς)} pour rétablir le trimètre iambique. --- \textbf{Bêtacisme vocalique au v. 947 :} πέσωμεν R pour \lecon{παίσωμεν} (Zanetti, confirmé par $\Sigma^{\text{R}}$). --- \textbf{v. 927 :} προλείπωσ' R (iotacisme) : \lecon{προλίπωσ'} Scaliger. --- \textbf{v. 947 :} πέσωμεν R : \lecon{παίσωμεν} Zanetti, cf. $\Sigma$ \espApp ταῖσ R : \lecon{ταῖσι} Daubuz.\par
}
\end{minipage}
\end{adjustbox}
\end{adjustwidth}
\newpage
\enTeteD{Codex Ravennas 429 (Biblioteca Classense) --- Fac-similé du f. 173v (vv. 920–958)}
\begin{adjustwidth}{-1.5cm}{-3.0cm}
\begin{center}
\includegraphics[width=16.5cm,height=20.5cm,keepaspectratio]{/workspaces/Thesmophories/dossier_ravennas_thesmophories/folios_individuels_ff162v-177r/23_f173v_classense175v_vv0920-0958.jpg}\par
\vspace{2pt}
{\scriptsize\textbf{Codex Ravennas 429, f. 173v (Classense 175v)} --- ΕΠΕΙΣΟΔΙΟΝ Δ' / ΧΟΡΙΚΟΝ, \textbf{vv. 920–958} ($\text{R}^1, \Rpc, \Rdeux$, \PsiDeux).}
\end{center}
\end{adjustwidth}
\newpage
\clearpage\ifodd\value{page}\hbox{}\newpage\fi
\enTeteG{ΘΕΣΜΟΦΟΡΙΑΖΟΥΣΑΙ --- Texte grec établi (f. 173v, vv. 920–958)}
\enTeteD{LES THESMOPHORIES --- Traduction française (f. 173v, vv. 920–958)}
\begin{pages}
\begin{Leftside}
\beginnumbering
%% v. 920-923
\begin{groupe}%
%
\personnageGrec{Κρίτυλλα}%
\trimIamb{οἴμ᾽ ὡς πανοῦργος καὐτὸς εἶναί μοι δοκεῖς\\
καὶ τοῦδέ τις ξύμβουλος. οὐκ ἐτὸς πάλαι\\
ᾐγυπτιάζετ᾽. ἀλλ᾽ ὅδε μὲν δώσει δίκην.\testimonia{922}{(ᾐγυπτιάζετ᾽) \SynB α 494 αἰγυπτιάζειν $=$ Photius α 518 $\subset$ Suidas αι 75}%
\commentaire{922}{ᾐγυπτιάζετ᾽.}{Αἰγυπτιάζειν, faire l'Égyptien, comme Σιφνιάζειν, Χιάζειν, \italique{faire le Chiote}, Ἑλληνίζειν, \italique{faire l'Hellène}, Λεσβιάζειν, Βακίζειν, \italique{faire le Bakis} (\italique{Paix}, 1072). Les Égyptiens avaient la réputation d'être très rusés. On a conservé un vers d'Eschyle\marqueur{1} : %
\citComm{δεινοί πλέκειν τοι μηχανὰς Αἰγύπτιοι\\
Les Égyptiens habiles à ourdir des trames.}%
Cf. Properce (III, 9, 33) : \latin{Noxia Alexandria, dolis aptissima lellus.} (Note de Willems ; cf. aussi Taillardat §409.)}%
\debutNote{922}%
\note{1}{Tiré d'une scholie à Théocrite (15, 48b). C'est le fragment 373 de l'édition de Radt.}\\
προσέρχεται γὰρ ὁ πρύτανις χὠ τοξότης.\\\passeLigne}%
\end{groupe}%
%% v. 924
\begin{groupe}%
\personnageGrec{Εὐριπίδης}%
\trimIamb{τουτὶ πονηρόν· ἀλλ᾽ ὑπαποκινητέον.\\\passeLigne}%
\end{groupe}%
%% v. 925a
\begin{groupe}%
\personnageGrec{Μνησίλοχος}%
\trimIamb{ἐγὼ δ᾽ ὁ κακοδαίμων τί δρῶ;\\\passeLigne}%
\end{groupe}%
%% v. 925b-927
\begin{groupe}%
\personnageGrec{Εὐριπίδης}%
\trimIamb{\ligne{925}\decalage{ἐγὼ δ᾽ ὁ κακοδαίμων τί δρῶ;}μέν᾽ ἥσυχος.\\
οὐ γὰρ προδώσω σ᾽ οὐδέποτ᾽, ἤνπερ ἐμπνέω,\testimonia{926}{(ἤνπερ ἐμπνέω) Suidas η 379 (ἤν pro ἤνπερ)}\\
ἢν μὴ προλίπωσ᾽ αἱ μυρίαι με μηχαναί.\apparat{927}{προλείπωσ᾽ R : \lecon{προλίπωσ᾽} Scaliger}\\\passeLigne}%
\end{groupe}%
%% v. 928
\begin{groupe}%
\personnageGrec{Μνησίλοχος}%
\trimIamb{%
αὕτη μὲν ἡ μήρινθος οὐδὲν ἔσπασεν.\testimonia{928}{Suidas α 4475, μ 978 ; \Eustathe 760.44 ; scholium\expo{RV} \italique{Vespae} 175b}\apparat{928}{(Γυ.) R : \lecon{(Μνη.)} Scaliger}
\commentaire{928}{αὕτη μὲν ἡ μήρινθος οὐδὲν ἔσπασεν.}{Le tour est proverbial et l'expression nous est bien connue grâce à une scholie au vers 175 des \italique{Guêpes} où Aristophane joue avec l'halieutique proverbe :%
\citComm{παρὰ τὴν παροιμίαν « αὕτη μὲν ἡ μήρινθος οὐδὲν ἔσπασεν ». ἐλέγετο δὲ ἐπὶ τῶν πείρᾳ καθιέντων ἄγκιστρον μετὰ σχοίνου, καὶ μὴ ἐπιτυγχανόντων ἰχθύος τινός.\\
\francais{En regard du proverbe : « voici une ligne qui n'a rien tiré hors de l'eau ». On le dit à propos de ceux qui jettent un appât à la pointe acérée (un hameçon avec un roseau), mais qui ne rencontrent pas un poisson.}}%
Ce proverbe est dit d'une ruse qui échoue\marqueur{1}. En français : un échec sur toute la ligne.}
\debutNote{928}
\note{1}{Pour un exposé plus développé sur ces deux passages avec à la clef l'explication du jargon halieutique, le lecteur pourra aller consulter le §395 des \italique{Images d'Aristophane} de J. Taillardat. Érasme a recensé l'adage (n° 845).}\\\passeLigne}%
\end{groupe}%
%% v. 929-934a
\begin{groupe}%
%
\personnageGrec{Πρύτανις}%
\trimIamb{ὅδ᾽ ἔσθ᾽ ὁ πανοῦργος ὃν ἔλεγ᾽ ἡμῖν Κλεισθένης;\apparat{929}{ἔστ᾽ R : \lecon{ἔσθ᾽} \Mon \espApp ἔλεγεν R : \lecon{ἔλεγ᾽} \Mon}\\
οὗτος τί κύπτεις; δῆσον αὐτὸν εἰσάγων%\commentaire{}{}{\eff{ἔκυψεν ὥσπερ Τήλεφος} cf Taillardar 349}
\\
ὦ τοξότ᾽ ἐν τῇ σανίδι, κἄπειτ᾽ ἐνθαδὶ\testimonia{931}{(σανίδι) Photius σ 68 σανίς, σ 65 σανίδα ; Hesychius σ 169 σανίς ; \AelDion σ 7}\\
στήσας φύλαττε καὶ προσιέναι μηδένα\\
ἔα πρὸς αὐτόν, ἀλλὰ τὴν μάστιγ᾽ ἔχων\\
παῖ᾽ ἢν προσίῃ τις.\\\passeLigne}%
\end{groupe}%
%% v. 934b-935
\begin{groupe}%
\personnageGrec{Κρίτυλλα}%
\trimIamb{\ligne{934}\decalage{παῖ᾽ ἢν προσίῃ τις.}νὴ Δί᾽ ὡς νῦν δή γ᾽ ἀνὴρ\apparat{934}{νῦν δῆτ' R : \lecon{νῦν δή γ᾽} Dobree : νυνδή γ' Cobet}\\
ὀλίγου μ᾽ ἀφείλετ᾽ αὐτὸν ἱστιορράφος.\testimonia{935}{Suidas ι 691 ὀλίγου ταύτην ἱστιορράφος \espTest (ἱστιορράφος) Pollux VII 161}%
\commentaire{935}{ἱστιορράφος.}{Allusion aux hardes dont était couvert le prétendu Ménélas. (Note de Willems ; ἱστιορράφος est mis παρὰ προσδοκίαν pour μηχανορράφος\marqueur{1}.)}%
\debutNote{935}%
\note{1}{Cf. Euripide, \italique{Andromaque}, vv. 447, 1116 ; Sophocle, \italique{Œdipe Roi}, v. 387.}\\\passeLigne}%
\end{groupe}%
%% v. 936-938
\begin{groupe}%
\personnageGrec{Μνησίλοχος}%
\trimIamb{ὦ πρύτανι πρὸς τῆς δεξιᾶς, ἥνπερ φιλεῖς\\
κοίλην προτείνειν ἀργύριον ἤν τις διδῷ,%
\commentaire{936-7}{ὦ πρύτανι πρὸς τῆς δεξιᾶς, ἥνπερ φιλεῖς / κοίλην προτείνειν ἀργύριον ἤν τις διδῷ.}{%\eff{Σ + Hésychius : c'était l'usage d'attacher nu au sanis ?}
Précis à l'usage du justiciable confronté à un prytane.
Les deux répliques que Mnésiloque adresse au prytane nous donnent un bon aperçu de l'estime en laquelle Aristophane tenait ces magistrats -- car nous allons voir que Mnésiloque tente de le soudoyer, comme cela semble l'usage à Athènes.
On trouve le même reproche de concussion dans la \italique{Paix} (vers 905-908) :%
\citComm{Ἀλλ᾽, ὦ πρυτάνεις, δέχεσθε τὴν Θεωρίαν.\\
Θέασ᾽ ὡς προθύμως ὁ πρύτανις παρεδέξατο.\\
Ἀλλ᾽ οὐκ ἄν, εἴ τι προῖκα προσαγαγεῖν σ᾽ ἔδει,\\
ἀλλ᾽ ηὗρον ἄν σ᾽ ὑπέχοντα τὴν ἐκεχειρίαν.\\
Allons, prytanes, recevez Théoria. Voyez l'accueil empressé que lui fait ce prytane. Quelle différence s'il t'avait fallu introduire une affaire gratis. Je t'aurais vu te retrancher sur les vacances.}%
(ὑπέχοντα τὴν ἐκεχειρίαν, «~alléguant la trêve~», est la version bégayée de ce qu'il fallait vraiment entendre : ὑπέχοντα τὴν χεῖρα, «~tendant la main~».)\\
En effet, pour soumettre une requête à la Boulé l'on devait préalablement obtenir l'aval des prytanes ; le cas échéant, les bouleutes discutaient et rédigeaient une proposition de loi (προβούλευμα) soumise ensuite à l'Ecclésia. L'auteur de la \italique{République des Athéniens} (III, 1) constate qu'il fallait parfois patienter jusqu'à une année avant que les prytanes ne donnent leur accord.
De même, pour intenter une action en justice, un citoyen devait régler des \italique{frais de justice}, τά πρυτανεῖα\marqueur{1}, dont le nom indique clairement à qui cette somme devait être remise.\\ 
Par ses fonctions, le prytane \italique{s'engageait} à servir d'intermédiaire : ce genre de serment passait très certainement par une \italique{poignée de main}\marqueur{2}.
C'est le sens de «~ἥνπερ φιλεῖς / κοίλην προτείνειν~», «~[cette main] que tu as l'habitude de tendre ouverte~».\\
Mais toutes ces transactions étaient autant d'occasions pour les prytanes de détourner à leur profit une partie du bien public.
C'est entendu : le prytane est concussionnaire mais il serait suicidaire pour Mnésiloque de le lui dire !%
\citComm{Flattez-les, payez-les d'agréables mensonges,\\
Quelque indignation dont leur cœur soit rempli,\\
Ils goberont l'appât, vous serez leur ami.}%
Au contraire : Mnésiloque tente de mettre le prytane dans sa poche -- il y parvient d'ailleurs au vu de la bienveillante réponse qu'il reçoit : τί σοι χαρίσωμαι;\\
Pour cela, jeu de scène\marqueur{3} : tout en déclamant «~πρὸς τῆς δεξιᾶς κτλ.~», Mnésiloque, qui n'est pas encore lié, tend sa main vers le prytane et s'offre par ce geste en suppliant\marqueur{4}. 
(Et d'ajouter -- en aparté, évidemment ! mais ce que le prytane a bien compris et qui explique ses bonnes dispositions à l'égard de Mnésiloque -- : «~ἀργύριον ἤν τις διδῷ~», «~[cette main que tu aimes à tendre] lorsque l'on t'offre de l'argent~».)\\
Le prytane accordera bien cette petite faveur à Mnésiloque (χάρισαι βραχύ τί μοι), à lui qui est sur le point de mourir ; et Mnésiloque d'insister sur ce dernier point (καίπερ ἀποθανουμένῳ).
Aristophane joue avec une locution proverbiale, selon sa pente maintenant bien établie par notre commentaire : les Athéniens, même à l'article de la mort, étaient réputés continuer de quémander la main tendue\marqueur{5} :%
\citComm{Ἀττικὸς ὑπέχει τὴν χεῖρα ἀποθνήσκων· ἐπὶ τῶν φιλαργύρων. φιλοκερδεῖς γὰρ οἱ Ἀθηναῖοι.}%
Le \italique{verre du condamné} que demande Mnésiloque (καίπερ ἀποθανουμένῳ) n'est autre que cette faveur mendiée la main tendue vers le prytane.}%
\debutNote{936-7}%
\note{1}{Voir à ce sujet le Daremberg \& Saglio à l'entrée \latin{Prytaneia}. Ces informations aident à interpréter la glose d'Hésychius (π 1511) : %
\citComm{πεπρυτανευκώς· ἐπιμεμηνιευκώς. καὶ προανηλωκώς,}%
rangée parmi les fragments comiques adespotes (408) dans l'édition Kassel-Austin.}%
\note{2}{On pourra lire l'article de Francesco Mari : «~Les sens de la poignée de main en Grèce ancienne du VIIIe au Ve siècle avant J.-C.~», in \italique{Ktèma}, 43 (2018), pp. 105-131.}%
\note{3}{Il faut rendre aux jeux de scènes toute l'importance qu'ils avaient à l'époque ; empêtrés dans leurs costumes et d'imposants masques de théâtre, il est à parier que les acteurs compensaient cette gêne par des jeux de mains des plus expressifs.}%
\note{4}{(La \italique{droite} en question est bien celle de Mnésiloque et non du prytane.) Pour cette signification de la main tendue, cf. \LSJ \ à l'entrée προτείνω II.}%
\note{5}{Recensé par Diogénianus (III, 12), \italique{Corpus paroemiographorum Graecorum}, I, p. 216. C'est l'adage n° 1646 d'Érasme. À noter qu'Aristophane y fait aussi une possible allusion aux vers 170 et suivants des \italique{Grenouilles}, lorsque Dionysos hèle un mort pour qu'il porte ses bagages, ce dernier lui répond : Δύο δραχμὰς μισθὸν τελεῖς; Voir plus haut le commentaire du vers 738.}\\
χάρισαι βραχύ τί μοι καίπερ ἀποθανουμένῳ.\\\passeLigne}%
\end{groupe}%
%% v. 939a
\begin{groupe}%
\personnageGrec{Πρύτανις}%
\trimIamb{τί σοι χαρίσωμαι;\apparat{939}{χαρίσομαι R : \lecon{χαρίσωμαι} Casaubon}\\\passeLigne}%
\end{groupe}%
%% v. 939b-942
\begin{groupe}%
\personnageGrec{Μνησίλοχος}%
\trimIamb{\ligne{939}\decalage{τί σοι χαρίσωμαι;}γυμνὸν ἀποδύσαντά με,\\
κέλευε πρὸς τῇ σανίδι δεῖν τὸν τοξότην,\testimonia{940}{(πρὸς τῇ σανίδι δεῖν) Hesychius δ 926 δῆσαι πρὸς σανίδι· Ἀττικῶς· βασανίσαι ἐν τῷ δεσμωτηρίῳ}\\
ἵνα μὴ 'ν κροκωτοῖς καὶ μίτραις γέρων ἀνὴρ\apparat{941}{μὴν R : \lecon{μὴ 'ν} Küster}\\
γέλωτα παρέχω τοῖς κόραξιν ἑστιῶν.\testimonia{942}{Suidas ε 3215 \espTest (ἑστιῶν) Hesychius ε 6411 ἑστιῶν· ἀριστοποιῶν $=$ Photius ε 2031 $=$ \Syn ε 875 $=$ Suidas ε 3227 ; cf. \Phryn 66.7}\apparat{942}{\inRas{γε}{} τοῖς \Rdeux}% \espApp ἑστιῶν R : \lecon{ἡστίων} Bibi}%
\commentaire{939-42}{γυμνὸν ἀποδύσαντά με / κέλευε πρὸς τῇ σανίδι δεῖν τὸν τοξότην, / ἵνα μὴ 'ν κροκωτοῖς καὶ μίτραις γέρων ἀνὴρ / γέλωτα παρέχω τοῖς κόραξιν ἑστιῶν.}{%
Ces quatre vers sont peut-être parmi les plus mécompris de toute la comédie.
L'erreur vient de ce que `πρός τῇ σανίδι δεῖν' signifie littéralement « lier \italique{contre} la planche » et non pas « lier à la planche » ; cette dernière expression se construisant en grec avec la préposition ἐν et non πρός\marqueur{1}.
Que dit en apparence Mnésiloque ? il fait profession honorable de pudeur ; ce pudibond de la onzième heure tient à cacher ce qui doit être caché. 
Voilà pourquoi il demande à être retourné \italique{contre} (πρός) \italique{la planche} du supplice (τῇ σανίδι)\marqueur{2}.
À qui tient-il à cacher son anatomie ? aux femmes évidemment\marqueur{3} !.. Au vers 942, les `κροκωτοῖς καὶ μίτραις' sont les \italique{crocotes} et \italique{mitres} des Athéniennes réunies, rappelons-le, en assemblée dans le but de perdre Euripide. Dans ce contexte, ἐν + datif a le sens qu'il a toujours\marqueur{4} : \italique{au milieu de}, \italique{parmi}, \italique{devant} : «~ἐν κροκωτοῖς καὶ μίτραις~» $=$ «~au milieu des (en)crocotées et des (en)mitrées~».
Voici pour le sens clair.

Mais n'oublions pas que Mnésiloque pratique avec un art consommé le double langage : toute parole qu'il adresse au prytane est à double sens : un clair, où il est dit que le sens véritable est caché\marqueur{5}. 
Que demande réellement Mnésiloque au prytane ?
L'expression `δῆσαι πρὸς σανίδι' a heureusement été relevée par Hésychius (δ 926) -- et c'est le seul des nombreux lexique consultés à avoir relevé ce tour :%
\citComm{δῆσαι πρὸς σανίδι· Ἀττικῶς· βασανίσαι ἐν τῷ δεσμωτηρίῳ.\\
Lier face à la planche : Atticisme. Interroger dans la prison.}%
Mnésiloque veut un procès équitable (il sent que la justice d'exception risque de lui être par trop expéditive) : il est préférable pour lui de s'en remettre aux masculines institutions. Il demande donc à être conduit en prison pour y être interrogé.

Réponse négative du prytane : c'est sur place, \italique{devant les personnes présentes} et non pas au δεσμωτήριον, qu'il devra expier son forfait :%
\citComm{Πρύτανις --- ἔχοντα ταῦτ᾽ ἔδοξε τῇ βουλῇ σε δεῖν,\\
ἵνα τοῖς παροῦσι δῆλος ᾖς πανοῦργος ὤν.\\}%
La σανίς de l'expression rapportée par Hésychius désigne sans doute la \italique{planche} sur laquelle était inscrite l'ordre du jour des tribunaux.
Cf. en \italique{Guêpes} (v. 848) et Photius (σ 68) : %
\citComm{σανίς· καὶ ἐν ᾗ τοὺς κακούργους ἔδουν καὶ ἐν ᾗ τὰς δίκας ἐνεχάραττον· τὸ µὲν ἐν Θεσµοφοριαζούσαις, τὸ δὲ ἐν Σφηξὶ (848) δεδήλωται.}%
Aristophane, ainsi qu'il aime à le faire assez souvent, détourne un mot ou une expression imagée qu'il interprète littéralement.

Devant le refus du prytane, Mnésiloque se retourne contre sa crocote -- au cœur des péripéties de la pièce -- qu'il accuse d'être à l'origine de la situation : il la gronde comme on gronderait un enfant\marqueur{6} :%
\citComm{Μνησίλοχος --- ἰαππαπαιάξ· ὦ κροκώθ᾽ οἷ᾽ εἴργασαι.}%

C'est encore Hésychius (κ 4170) qui nous permet de confirmer qu'il y a probablement là un jeu de mots κρόκαλ' $\sim$ κροκώθ' :%
\citComm{κρόκαλον· τὸ πανοῦργον παιδίον.}%

%1004 γυμνῷ φυλακὴν ἐπιτάττειν· παροιμία ἐπὶ τῶν <μὴ> δεομένων προστάξεως, διὰ τὸ ἀναγκαίως καὶ χωρὶς ἐπικελεύσεως τοῦτο πράττειν. Φερεκράτης Τυραννίδι (fr.154 Κ.Ἁ.)
%Suidas : Γυμνῷ φυλακὴν ἐπιτάττεις: ἐπὶ τῶν ἀδυνάτως ἐχόντων τὰ προστεταγμένα πληροῦν, διὰ τὰ ἀναγκαῖα. ὅτι τὸ Γυμνῶ μετὰ γενικῆς: μετὰ αἰτιατικῆς δὲ, ὡς τὸ, γυμνώσαντες αὐτὸν ἐξήγαγον ἔξω. καὶ, γυμνῶν τὰ ἐσφραγισμένα.
%
%\marqueur{Pour se rendre bien compte de ce que pouvait être cette omniprésence envahissante -- et salissante ! -- des oiseaux lors des sacrifices ou, comme c'est notre cas, d'exécution à ciel ouvert, le lecteur pourra se reporter au Daremberg \& Saglio, à l'entrée \latin{Meniskos} (Μηνίσκος).}
%\eff{Hésychius δῆσαι πρὸς σανίδι:᾿Αττικῶς:βασανίσαιἐν τῷ δεσμωτηρίῳ (Αr. Τhesm. 940 ?) en fait Cratinus cf. scholiaste de notre passage  ἡλωτός: cloué }
%	pros tê sanidi : face à la planche, pour cacher sa nudité... mais voir aussi gerron Ael. Dion.
%\eff{fr. 639 K.-A. : ἐπὶ τῶι ταρίχει τὸν γέλωτα κατέδομαι.\\
%fr. 803 γέρρον ...  ̓Αριοτοφάνης ὧς ἐπὶ φυλακῆς τινος καὶ μοχλοῦ τροπικῶς παρέλαβε τὴν λέξιν}
%Ces quatre vers sont peut-être parmi les plus mécompris de toute la comédie. 
%Les commentateurs avaient effectivement bien peu de chance de les comprendre puisqu'ils commencent par un contre-sens : ils construisent κέλευε τὸν τοξότην γυμνὸν ἀποδύσαντά με, «~ordonne à l'archer de me déshabiller tout nu~»...
%Comment Mnésiloque pourrait-il faire cette requête alors que Clisthène l'a déjà mis à nu\marqueur{1} ? En réalité, le groupe `γυμνὸν ἀποδύσαντά με' est un accusatif absolu et se traduit par : «~puisque j'ai été mis à nu, ordonne à l'archer \&c.~»\\
%Mnésiloque pratique avec un art consommé le double langage : toute parole qu'il adresse au prytane est à double sens : un clair, mais où il est dit que le sens véritable est caché. 
%Que dit en apparence Mnésiloque ? il fait profession honorable de pudeur ; ce pudibond de la onzième heure tient à cacher ce qui doit être caché. 
%Voilà pourquoi il demande à être retourné \italique{contre} (πρός) \italique{la planche} du supplice (τῇ σανίδι)\marqueur{2}. %\marqueur{Le scholiaste ne disait pas autre chose NON : `ἐπίτρεψον οὖν τῷ τοξότῃ', φησὶ, `γυμνόν με δῆσαι πρὸς τῇ σανίδι'.}.
%À qui tient-il à cacher son anatomie ? aux femmes évidemment -- et non aux corbeaux\marqueur{3} !.. Au vers 942, les `κροκωτοῖς καὶ μίτραις' sont les \italique{crocotes} et \italique{mitres} des femmes réunies, rappelons-le, en assemblée dans le but de perdre Euripide.\marqueur{4}. Dans ce contexte, ἐν + datif a le sens qu'il a toujours\marqueur{6} : \italique{au milieu de}, \italique{parmi}, \italique{devant} : «~ἐν κροκωτοῖς καὶ μίτραις~» $=$ «~au milieu des (en)crocotées et des (en)mitrées~».
%Voici pour le sens clair.\\
%Mais que demande réellement Mnésiloque au prytane ?
%L'expression `δῆσαι πρὸς σανίδι' a heureusement été relevée par Hésychius (δ 926) -- qui se montre encore une fois beaucoup plus pertinent que notre scholiaste :%
%\citComm{δῆσαι πρὸς σανίδι· Ἀττικῶς· βασανίσαι ἐν τῷ δεσμωτηρίῳ.\\
%Lier à la planche : Atticisme. Interroger dans la prison.}%
%Mnésiloque veut un procès équitable (il sent que la justice d'exception risque de lui être par trop expéditive) : il est préférable pour lui de s'en remettre aux masculines institutions.
%La σανίς de l'expression désigne sans doute la \italique{planche} sur laquelle était écrite inscrite l'ordre du jour des tribunaux.
%Cf. en \italique{Guêpes} (v. 848) et Photius (σ 68) : %
%\citComm{σανίς· καὶ ἐν ᾗ τοὺς κακούργους ἔδουν καὶ ἐν ᾗ τὰς δίκας ἐνεχάραττον· τὸ µὲν ἐν Θεσµοφοριαζούσαις, τὸ δὲ ἐν Σφηξὶ (848) δεδήλωται.}%
%L'on reste consterné devant l'interprétation des professionnels du commentaire pour ce passage : Mnésiloque se proposerait de graisser la patte du prytane – qui semble avoir quelque inclination naturelle pour la chose – en avançant la raison officielle qu'il ne veut pas prêter à rire aux corbeaux (sic !?!)... \\
%Soufflé, bouche bée, cloué au bec par une telle exégèse !
%Il y a effectivement de quoi rire, mais ce n'est pas les corbeaux qui le feront.
%Et d'abord, leur fantaisie ne tient même pas du point de vue du grec. Nous ne sommes dans un pastiche de style tragique que je sache ; le pluriel ἐν κροκωτοῖς καὶ μίτραις n'aurait donc rien à faire ici s'il désignait \italique{la}\marqueur{1} crocote ou \italique{la} mitre de Mnésiloque.
%Et si ces crocotes et mitres ne sont pas celles de Mnésiloque ; ce sont donc celles des femmes réunies, rappelons-le, en assemblée dans le but de perdre Euripide. Dans ce contexte, ἐν + datif a le sens qu'il a toujours : \italique{au milieu de}, \italique{parmi}, \italique{devant} : "ἐν κροκωτοῖς καὶ μίτραις" $=$ "au milieu des (en)crocotées et des (en)mitrées".\\
%Ce n'est donc évidemment pas devant les corbeaux mais bien devant les \italique{femmes} que Mnésiloque dit craindre d'être ridicule.\\
%Il n'y a plus qu'à rétablir un augment oublié à la fin du vers 942\marqueur{2} : ἡστίων (imparfait) en place de l'incompréhensible ἑστιῶν et voilà notre passage parfaitement intelligible : %
%\citComm{γυμνὸν ἀποδύσαντά με\\
%κέλευε πρὸς τῇ σανίδι δεῖν τὸν τοξότην,\\
%ἵνα μὴ 'ν κροκωτοῖς καὶ μίτραις γέρων ἀνὴρ\\
%γέλωτα παρέχω τοῖς κόραξιν ἡστίων.\\
%Ordonne à l'archer que je sois déshabillé avant que d'être attaché à la planche afin qu'au milieu des porteuses de crocotes et de mitres je ne prête pas à rire avec les corbeaux qui dîneraient.}%
%Figure porte absence et présence, dit Pascal ; l'on soupçonne un double sens : un clair, que nous venons d'établir, mais où il est dit que le sens est caché. 
%Puisqu'il ne semble pas qu'il y ait d'allusion proverbiale, littéraire ou à la chronique du temps, reste le jeu de mots (cas le plus favorable puisque nous possédons tout le matériel pour le saisir).\\
%Ne perdons pas de vue l'objectif de Mnésiloque : s'enfuir ; or par analogie κόραξ désigne tout instrument de forme recourbée comme un bec de corbeau et peut être synonyme de κατακλείς : \italique{cliquet}, \italique{crochet} ; cette fois, nous tenons le bon bout ! Mnésiloque demande bien-sûr, en loucédé, à être libéré de ses entraves.\\
%Poursuivons. Le verbe παρέχω, \italique{fournir}, ne diffère que bien peu du verbe παρα-χόω, \italique{verser sur le côté}. 
%En faisant quelques efforts de désarticulation, le protagoniste pouvait certainement faire entendre à la fois παρέχω, \italique{que je prête} [à rire] et παραχοῖ\marqueur{3}, \italique{qu'il verse de côté} [mes crochets].\\
%Le double sens se précise ; il demeure cependant le dernier mot, ἑστιῶν. 
%Notre correction en ἡστίων nous facilite la tâche : il suffit de faire la liaison avec τοῖς κόραξιν, la prononciation devait être très proche de celle de τοῖς κόραξι νήστιων, \italique{les crochets [au niveau] de mes viscères}. 
%(Pour cette acception de νῆστις, Photius la donne à l'entrée νῆστιν en citant Aristophane lui-même : νῆστιν· τὸ ἔντερον· `ἢ σπλῆνα ἢ γλῶτταν'. Ἀριστοφάνης\marqueur{4}.)\\
%Sans forcer le texte, ce nous semble, voici donc une lecture alternative :%
%\citComm{παρέχω τοῖς κόραξιν ἡστίων\\
%παραχοῖ τοῖς κόραξι νήστιων\\
%Qu'il [l'archer] me débarrasse  des crochets [au niveau] de mes reins.}%
%Admirez, s'il vous plaît, ce petit bijou d'ambiguïté qui ne peut être comparé, chez nous, qu'aux perles du génial Raymond Devos.
}%
\debutNote{939-42}%
\note{1}{Cf. les vers 934 (δῆσον αὐτὸν [...] ἐν τῇ σανίδι) et 1165 (οὑν [$=$ ὁ ἐν] τῇ σανίδι).}%
\note{2}{En évitant par la même occasion que les corbeaux ne l'énucléent et ne se repaissent de ses yeux... spectacle qu'on image sans peine peu ragoûtant et fort désagréable pour le supplicié.}%
\note{3}{Et non aux corbeaux, comme on peut le lire dans les éditions les plus récentes des \italique{Thesmophories}...}%
\note{4}{\italique{Confer} les vers 90-91.}%
\note{5}{Quelques vers auparavant, il avait demandé une faveur au prytane, tout en lui faisant bien comprendre que cette faveur mériterait salaire. Sur la concussion des prytanes, voir aussi la \italique{Paix} (vers 905-908).}%
\note{6}{Voir par exemple, \italique{Philoctète}, v. 940 : οἷ᾽ ἔργ᾽ ὁ παῖς μ᾽ ἔδρασεν οὑξ Ἀχιλλέως.}%
%\note{1}{Cf. vers 637 sqq. et les notes afférentes ; Mnésiloque n'a plus sur lui que sa crocote tout-à-fait transparente.}%
%\note{2}{L'expression grecque correcte pour dire \italique{lier à la planche} est ἐν σανίδι et non πρός σανίδι, cf. vers 934 et 1165.}%
%\note{3}{Je prie le lecteur d'excuser la présence de remarques aussi évidentes mais il faut bien rectifier la clique des commentateurs capable d'écrire à l'unisson et sans sourciller ce genre d'âneries.}%
%\note{4}{Et non les habits de Mnésiloque comme le prétend l'internationale du commentaire !.. d'abord, on les lui a retirés ; ensuite cette interprétation fantaisiste ne tient même pas du point de vue du grec : nous ne sommes dans un pastiche de style tragique que je sache ; le pluriel ἐν κροκωτοῖς καὶ μίτραις n'aurait donc rien à faire ici s'il désignait \italique{la}\marqueur{5} crocote ou \italique{la} mitre de Mnésiloque.}%
%\note{5}{Ou \italique{le} si l'on voulait rester conforme au genre du mot grec ; mais l'oreille française interdit toute féminisation de «~crocote~».}%
%\note{6}{Cf. les vers 90-91 et la note qui s'y rapporte.}%
%\note{2}{Même oubli que, par exemples, aux vers 488, 521 ou 794.}%
%\note{3}{D'autant que les règles classiques de contraction auraient dû fournir la forme παραχῷ.}%
%\note{4}{Photius, ν 201. Cf. Athénée III, 96c $=$ fr. 520 Kassel-Austin.}
\\\passeLigne}%
\end{groupe}%
%% v. 943-944
\begin{groupe}%
\personnageGrec{Πρύτανις}%
\trimIamb{ἔχοντα ταῦτ᾽ ἔδοξε τῇ βουλῇ σε δεῖν,\apparat{943}{ἔδοξεν R : \lecon{ἔδοξε} Biset}%
\commentaire{943-56.}{Le papyrus d'Oxyrhynque n° 3839\marqueur{1},}{dont il a été déjà question aux vers 742-766 contient contient également un fragment de ce passage. Il est si gâté qu'il ne nous apporte aucune information (si ce n'est une colométrie que Lætitia suivit \anglais{without strong conviction}\marqueur{2}). Le lecteur le trouvera néanmoins reproduit en annexe ; gage à lui d'en tirer quelque chose.}%
\debutNote{943-56}%
\note{1}{\italique{The Oxyrhynchus Papyri}, vol. LVI (1989), n°3839, fr. 2. Analysé par W. E. H. Coockle avec le concours de Colin Austin.}%
\note{2}{L. P. E. Parker, \italique{The Songs of Aristophanes} (1997), p. 430.}\\
ἵνα τοῖς παροῦσι δῆλος ᾖς πανοῦργος ὤν.\apparat{944}{\lecon{παροῦσι} R : παριοῦσι Brunck}\\\passeLigne}%
\end{groupe}%
%% v. 945-946
\begin{groupe}%
\personnageGrec{Μνησίλοχος}%
\trimIamb{ἰαππαπαιάξ· ὦ κροκώθ᾽ οἷ᾽ εἴργασαι·\testimonia{945}{(ἰαππαπαιάξ) Suidas ε 2807}\apparat{945}{\lecon{ἰαππαπαιάξ} R, Suid. ε 2807 : ἰατταταιάξ Biset \ligneApp{946} ἔστιν R : \lecon{ἔστ᾽} Bentley}%
\commentaire{945}{ὦ κροκώθ᾽ οἷ᾽ εἴργασαι·}{Y a-t-il un jeu de mots entre κρόκαλον et κροκωτός ? 
Cf. Hésychius (κ 4170 et 4270) :%
\citComm{κρόκαλον· τὸ πανοῦργον παιδίον.}%
}\\
κοὐκ ἔστ᾽ ἔτ᾽ ἐλπὶς οὐδεμία σωτηρίας.%\apparat{946}{ἔστιν R : \lecon{ἔστ᾽} Bentley}
\\\passeLigne}%
\end{groupe}%
%% v. 947-952
\begin{groupe}%
%
\personnageGrec{Κορυφαῖος}%
\typeVers{\decalageN{4}$4$ an\cat}\tetramAnap{\Automatique{ἄγε νυν ἡμεῖς παίσωμεν ἅπερ νόμος ἐνθάδε ταῖσι γυναιξίν,}\apparat{947}{πέσωμεν R : \lecon{παίσωμεν} Zanetti, cf. Σ \espApp ταῖσ R : \lecon{ταῖσι} Daubuz \ligneApp{948} \lecon{θεαῖν} R : θεοῖν Cobet}%
\commentaire{947-1000}{}{Pour une analyse plus fine du chant mêlé de danses de tout ce passage, je renvoie à l'article d'Ester Cerbo\marqueur{1}.
Quant à sa structure, celle d'Enger, assez communément admise -- et notamment par Lætitia\marqueur{2} --, me semble juste  :%
\citComm{--- anapestes (v. 947-952) ;\\
--- proode (v. 953-958) ;\\
--- première paire strophique (v. 959-962 $=$ v. 963-965) ;\\
--- première mésode (v. 966-968) ;\\
--- seconde paire strophique (v. 969-976 $=$ v. 977-984) ;\\
--- seconde mésode (v. 985-989) ;\\
--- tierce paire strophique (v. 990-994 $=$ v. 995-1000).}%
}%
\debutNote{947-1000}%
\note{1}{«~Uno spettacolo di danza alla festa delle donne: a proposito di Ar. \italique{Thesm.} 947-1000~» (2021).}%
\note{2}{Dans son indispensable \italique{The Songs of Aristophanes}, pp. 428-436.}\\
\Automatique{παραχώρει ὅταν ὄργια σεμνὰ θεαῖν ἱεραῖς ὥραις ἀνέχωμεν, ἅπερ καὶ}\testimonia{947}{(ὅταν -- ἀνέχωμεν) Suidas α 2364 (θεαῖν codd. GFS\expo{a.c.} vel θεῶν codd. ATS\expo{p.c} vel θεεῖν cod. I vel θεαῖνων cod. M pro θεοῖν)}}%
\end{groupe}%
\begin{groupe}%
\typeVers{dim. anap.}\dimAnap{Παύσων σέβεται καὶ νηστεύει,%
\commentaire{949-52}{Παύσων σέβεται καὶ νηστεύει κτλ.}{Pauson était connu pour son indigence. Les parémiographes nous ont rapporté à son sujet deux proverbes\marqueur{1} :%
\citComm{Παύσωνος πτωχότερος· οὗτος ζωγράφος ἦν καὶ ἐπὶ πενίᾳ διατεθρύλλητο.\\
Παύσων καὶ Ἶρος: ὀνόματα κύρια πενήτων. ὁ Ἀσκληπιὸς Παύσωνα καὶ Ἶρον κἂν ἄλλον τινὰ τῶν ἀπόρων ἰάσαιτο. κτλ.\\
(Iros, le mendiant odieux d'Homère, \italique{Odyssée} XVIII.)}%
L'on trouvera tous les témoignages sur Pauson d'Athènes dans le \italique{Recueil Milliet} d'Adolphe Reinach\marqueur{2}. Pour ce qui est d'Aristophane, la scholie à notre passage indique :%
\citComm{Παύσων ὁ ζωγράφος εἰς πενίαν διαβάλλεται, καὶ διὰ τοῦτο εἶπε `νηστεύει'.}%
Pauson l'importun est aussi mentionné au vers 853 des \italique{Acharniens} :%
\citComm{οὐδ᾽ αὖθις αὖ σε σκώψεται\\
Παύσων ὁ παμπόνηρος\\
Ni les brocards de Pauson, le gueux achevé, tu ne les essuieras [sur le marché].\\
Σ : οὗτος ὁ Παύσων ζωγράφος πένης σχωπτολόγος.}%
Enfin au vers 602 du \italique{Ploutos} :%
\citComm{Χρέμυλος --- Παρὰ τῆς Ἑκάτης ἔξεστιν τοῦτο πυθέσθαι,\\
εἴτε τὸ πλουτεῖν εἴτε τὸ πεινῆν βέλτιον. Φησὶ γὰρ αὕτη\\
τοὺς μὲν ἔχοντας καὶ πλουτοῦντας δεῖπνον κατὰ μῆν᾽ ἀποπέμπειν,\\
τοὺς δὲ πένητας τῶν ἀνθρώπων ἁρπάζειν πρὶν καταθεῖναι.\\
Ἀλλὰ φθείρου καὶ μὴ γρύξῃς\\
ἔτι μηδ᾽ ὁτιοῦν.\\
Οὐ γὰρ πείσεις, οὐδ᾽ ἢν πείσῃς.\\
Πενία --- Ὦ πόλις Ἄργους, κλύεθ᾽ οἷα λέγει.\\
Χρέμυλος --- \textbf{Παύσωνα} κάλει τὸν ξύσσιτον.\\
Chrémyle --- Il n'est que de s'informer chez Hécate, lequel vaut mieux d'être riche ou d'être pauvre. Car elle prétend que les riches et les opulents lui font hommage chaque mois d'un dîner, et que les pauvres s'en emparent avant qu'on l'ait mis là. Mais la peste soit de toi ; ne souffle plus le mot. Car tu ne me persuaderas pas, quand tu me persuaderais.\\
La Pauvreté --- « Cité des Argiens, écoutez ce qu'il dit ! »\\
Chrémyle --- Invoque \textbf{Pauson}, ton commensal.\\
Σ : ὁ Παύσων δὲ ἐπὶ πενίᾳ κωμῳδεῖται ζωγράφος ὦν. --- μετακαλοῦ σύντροφον τὸν Παύσωνα.}%
On constate que seuls les commentateurs indiquent que Pauson fût peintre (ζωγράφος) : il n'en n'est jamais question chez Aristophane.
Il paraît douteux que ce Pauson-là, le mendiant fâcheux à la dévotion fort intéressée\marqueur{4}, fût le peintre connu qu'évoque Aristote\marqueur{3} ; la confusion venant peut-être de ce que la façon de peindre les figures du ζωγράφος ainsi que la seule anecdote nous ayant été rapportée à son sujet témoignent d'un caractère porté à la misanthropie\marqueur{5}.}%
\debutNote{949-52}%
\note{1}{Respectivement Michel Apostolius (XIII, 2) et Suidas (π 823).}%
\note{2}{\italique{Recueil Milliet : Textes grecs et latins relatifs à l'histoire de la peinture ancienne}, tome I, publiés, traduits et commentés par Adolphe Reinach (1921) ; §§ 131-132 ; 173-179.}%
\note{3}{Aristote, \italique{Politique}, 1340a36 ; \italique{Poétique}, 1448a6.}%
\note{4}{Si sa dévotion prend les formes habituelles du souhait de renouvellement des prières (cf. le vers 287 ; Euripide, \italique{Électre}, vers 805), on se doute bien de ce qui attirait réellement Pauson-crève-la-faim. Voir ce que nous avons déjà dit aux vers 558-9 de la κρεανομία -- ou orgie institutionnalisée -- qui seule explique la réputation de piété des Athéniens.}%
\note{5}{Pauson le peintre peignait des figures pires que ce qu'elles étaient en réalité et Aristote le déconseille pour l'éducation des jeunes gens. On lui avait commandé le tableau d'un cheval se roulant dans la poussière ; l'artiste dessine un cheval courant et soulevant un nuage de poussière autour de lui. Il y travaillait quand survient l'amateur. Celui-ci récrimine : ce n'est pas le sujet commandé. Pauson fait retourner par son garçon le tableau sens dessus dessous ; on voit alors un cheval renversé, qui se roule dans la poussière.}\\
πολλάκις αὐταῖν ἐκ τῶν ὡρῶν\apparat{950}{\lecon{αὐταῖν} R : αὐτοῖν Cobet}\\
ἐς τὰς ὥρας ξυνεπευχόμενος}
\end{groupe}%
\begin{groupe}%
\typeVers{paroem.}\dimAnap{τοιαῦτα μέλειν θάμ᾽ ἑαυτῷ.\apparat{952}{μέλλειν R : \lecon{μέλειν} Zanetti (\italique{curae esse} Divus) : μολεῖν Casaubon}\\\passeLigne}%
\end{groupe}%
%% v. 953-958
\begin{groupe}%
%
\personnageGrec{Χορός}%
\typeVers{anap.}\chantChoeur{ὅρμα χώρει·}
\end{groupe}%
%% v. 953-958
\begin{groupe}%
\typeVers{lecyth.}\chantChoeur{κοῦφα ποσὶν ἄγ᾽ ἐς κύκλον,}
\end{groupe}%
%% v. 953-958
\begin{groupe}%
\typeVers{aristoph.}\chantChoeur{χειρὶ σύναπτε χεῖρα,\apparat{955}{\lecon{χειρὶ} R : χερὶ Dobree}%
\commentaire{955}{χειρὶ σύναπτε χεῖρα.}{La danse que commande le chœur est probablement ce que Maurice Emmanuel\marqueur{1} appelle une \italique{danse des mains jointes}. Il s'agit donc bien d'une danse individuelle et non en couple : chaque danseur joignant la main à son autre main.}%
\debutNote{955}%
\note{1}{\italique{La danse grecque antique} (1896), §319.}}%
\end{groupe}%
%% v. 953-958
\begin{groupe}%
\chantChoeur{†ῥυθμὸν χορείας ὕπαγε πᾶσα†·\apparat{956}{\lecon{ὕπαγε πᾶσα} R : πᾶσ' ὕπαγε Austin}}
\end{groupe}%
%% v. 953-958
\begin{groupe}%
\typeVers{glyc.}\chantChoeur{\ligne{956}βαῖνε καρπαλίμοιν ποδοῖν.}
\end{groupe}%
%% v. 953-958
\begin{groupe}%
\typeVers{dim. iamb.}\chantChoeur{ἐπισκοπεῖν δὲ πανταχῇ}
\end{groupe}%
%% v. 953-958
\begin{groupe}%
\typeVers{trim. iamb.}\chantChoeur{κυκλοῦσαν ὄμμα χρὴ χοροῦ κατάστασιν.\\\passeLigne}%
\end{groupe}%
\endnumbering
\end{Leftside}
\begin{Rightside}
\beginnumbering
%% v. 920-923
\begin{groupe}%
\personnageFr{Kritylla}%
\trimIambFr{Ah! toi aussi, que tu m'as l'air d'un fripon
et d'être son conseiller. Rien d'étonnant que depuis un temps
vous jouiez les Égyptiens. Mais celui-là en portera la peine.
Car voici venir le prytane et l'archer.\\\passeLigne}%
\end{groupe}%
%% v. 924
\begin{groupe}%
\personnageFr{Euripide}%
\trimIambFr{Voilà qui ne vaut rien. Il est temps d'esquiver.\\\passeLigne}%
\end{groupe}%
%% v. 925a
\begin{groupe}%
\personnageFr{Mnésiloque}%
\trimIambFr{Et moi, malheureux, que vais-je faire ?\\\passeLigne}%
\end{groupe}%
%% v. 925b-927
\begin{groupe}%
\personnageFr{Euripide}%
\trimIambFr{\ligne{925}\espace{4.2cm}Sois en paix.\\
Jamais je ne te trahirai, tant que j'aurai le souffle,\\
à moins que mes mille stratagèmes ne viennent à me \fin faillir.\\\passeLigne}%
\end{groupe}%
%% v. 928
\begin{groupe}%
\personnageFr{Mnésiloque}%
\trimIambFr{%
Voilà une amorce qui n'a pas pris.\\\passeLigne}%
\end{groupe}%
%% v. 929-934a
\begin{groupe}%
\personnageFr{Prytane}%
\trimIambFr{Est-ce là le fripon dont nous parlait Clisthène ?\\
Holà, pour quoi courbes-tu le front? Archer, rentre avec lui et attache-le\\
à l'ais. Puis expose-le ici,
et garde-le bien. Ne laisse approcher personne,\\
mais frappe de ton fouet\\
\ligne{934}quiconque s'avancera.\\\passeLigne}%
\end{groupe}%
%% v. 934b-935
\begin{groupe}%
\personnageFr{Kritylla}%
\trimIambFr{\ligne{934}\espace{2.9cm}Bien, par Zeus, car tout à l'heure\\
un rapetasseur de voiles a failli me l'enlever.\\\passeLigne}%
\end{groupe}%
%% v. 936-938
\begin{groupe}%
\personnageFr{Mnésiloque}%
\trimIambFr{Ô prytane, au nom de cette droite dont tu aimes\\
à tendre le creux (À part :) quand on t'offre de l'argent...\\ 
Accorde-moi une petite faveur, à moi qui suis sur le point de mourir.\\\passeLigne}%
\end{groupe}%
%% v. 939a
\begin{groupe}%
\personnageFr{Prytane}%
\trimIambFr{\ligne{939}Quelle faveur ?\\\passeLigne}%
\end{groupe}%
%% v. 939b-942
\begin{groupe}%
\personnageFr{Mnésiloque}%
\trimIambFr{\ligne{939}\espace{2.8cm}ordonne à l'archer\\
de m'attacher nu face à l'ais.\\
Je crains d'être la risée des crocotes et des mitres\\
avec les corbeaux que je vais régaler.\\\passeLigne}%
\end{groupe}%
%% v. 943-944
\begin{groupe}%
\personnageFr{Prytane}%
\trimIambFr{La boulé a décidé qu'on t'attacherait dans ce costume,\\
pour témoigner aux personnes présentes que tu es un \fin fripon.\\\passeLigne}%
\end{groupe}%
%% v. 945-946
\begin{groupe}%
\personnageFr{Mnésiloque}%
\trimIambFr{Aïe, aïe! aïe, aïe! Ô crocote, que de malheurs tu as causés !
Il n'est plus pour moi le moindre espoir de \fin salut.\\\passeLigne}%
\end{groupe}%
%% v. 947-952
\begin{groupe}%
\personnageFr{Coryphée}%
\tetramAnapFr{Allons, livrons-nous à nos danses, comme c'est ici la coutume des femmes,
quand aux dates saintes nous célébrons les augustes mystères des deux déesses.}%
\end{groupe}%
\begin{groupe}%
\dimAnapFr{Pauson aussi les sanctifie. Il jeûne\\
et souvent prie avec nous que d'année\\
en année il puisse longtemps}
\end{groupe}%
\begin{groupe}%
\dimAnapFr{encore leur témoigner même dévotion.\\\passeLigne}%
\end{groupe}%
%% v. 953-958
\begin{groupe}%
\personnageFr{Chœur}%
\chantChoeurFr{Pars, élance-toi}
\end{groupe}%
%% v. 953-958
\begin{groupe}%
\chantChoeurFr{d'un pied léger, forme le cercle,}
\end{groupe}%
%% v. 953-958
\begin{groupe}%
\chantChoeurFr{les mains enlacées aux mains,}
\end{groupe}%
%% v. 953-958
\begin{groupe}%
\chantChoeurFr{marque le rythme de la danse.}
\end{groupe}%
%% v. 953-958
\begin{groupe}%
\chantChoeurFr{Avance de tes pieds alertes.}
\end{groupe}%
%% v. 953-958
\begin{groupe}%
\chantChoeurFr{Mais il faut que le chœur soit aux aguets}
\end{groupe}%
%% v. 953-958
\begin{groupe}%
\chantChoeurFr{et tourne la vue de toutes parts.\\\passeLigne}%
\end{groupe}%
\endnumbering
\end{Rightside}
\end{pages}
\Pages
\clearpage\ifodd\value{page}\hbox{}\newpage\fi
\enTeteG{ΧΟΡΙΚΟΝ --- Codex Ravennas 429, f. 174r (Classense 176r) --- vv. 959–1000}
\begin{adjustwidth}{-1.5cm}{-3.0cm}
\begin{adjustbox}{max totalheight=21.0cm, max width=16.5cm, keepaspectratio}
\begin{minipage}{16.5cm}
{\fontsize{7.2pt}{8.5pt}\selectfont
\noindent\textbf{Folio 24 --- Codex Ravennas 429, \texttt{f. 174r} (\texttt{Classense 176r}) : ΧΟΡΙΚΟΝ, vv. 959–1000}\hfill \textit{Témoins :} $\text{R}^1, \Rpc, \Rdeux$ (\Mon)\par\vspace{2pt}
\hrule height 0.4pt\vspace{2.5pt}
}
{\fontsize{8.1pt}{9.5pt}\selectfont
\ligneRav{959}{χορ/}{\hspace{1.1em}ἅμα δὲ καὶ γένος Ὀλυμπίων θεῶν}{}
\ligneRav{960}{}{\hspace{1.1em}μέλπε καὶ γέραιρε φωνῇ πᾶσα χορομανεῖ τρόπῳ·}{}
\ligneRav{}{}{\hspace{1.1em}εἰ δέ τις προσδοκᾷ κακῶς ἐρεῖν}{}
\ligneRav{965}{}{\hspace{1.1em}ἐν ἱερῷ γυναῖκά μ᾽ οὖσαν ἄνδρας, οὐκ ὀρθῶς φρονεῖ·}{}
\ligneRav{966}{}{\hspace{1.1em}ἀλλὰ χρή}{\textbf{966} \lecon{χρὴ} R : χρῆν Bothe : χρή μ᾽ Meineke}
\ligneRav{967}{}{\hspace{1.1em}ὥσπερ ἔργον αὖ τι καινὸν}{\textbf{967} \lecon{ὥσπερ} R : ὡς πρὸς Austin}
\ligneRav{968}{}{\hspace{1.1em}{\Taille{11.5}πρῶτον εὐκύκλου χορείας \difference{εὐφυᾶ} στῆσαι βάσιν.}}{\textbf{968} εὐφυῆ R : \lecon{εὐφυᾶ} Brunck}
\ligneRav{969}{}{\hspace{1.1em}πρόβαινε \difference{ποσὶ} τὸν \difference{εὐλύραν}}{\textbf{969} ποσὶν R : \lecon{ποσὶ} Bentley \espApp ἐλύραν R : \lecon{εὐλύραν} Zanetti}
\ligneRav{970}{}{\hspace{1.1em}μέλπουσα καὶ τὴν τοξοφόρον}{}
\ligneRav{}{}{\hspace{1.1em}Ἄρτεμιν ἄνασσαν ἁγνήν·}{}
\ligneRav{}{}{\hspace{1.1em}χαῖρ᾽ ὦ ἑκάεργε,}{}
\ligneRav{}{}{\hspace{1.1em}ὄπαζε δὲ νίκην·}{}
\ligneRav{}{}{\hspace{1.1em}Ἥραν τε τὴν τελείαν}{}
\ligneRav{975}{}{\hspace{1.1em}μέλψωμεν ὥσπερ εἰκός,}{}
\ligneRav{}{}{\hspace{1.1em}ἣ πᾶσι τοῖς χοροῖσιν ἐμπαίζει τε καὶ}{}
\ligneRav{}{}{\hspace{1.1em}κλῇδας γάμου φυλάττει·}{}
\ligneRav{}{}{\hspace{1.1em}Ἑρμῆν τε νόμιον ἄντομαι}{}
\ligneRav{}{}{\hspace{1.1em}καὶ Πᾶνα καὶ Νύμφας φίλας}{}
\ligneRav{}{}{\hspace{1.1em}ἐπιγελάσαι προθύμως}{}
\ligneRav{980}{}{\hspace{1.1em}ταῖς \difference{ἡμετέραισι}}{\textbf{980} ἡμετέραις R : \lecon{ἡμετέραισι} Hermann}
\ligneRav{}{}{\hspace{1.1em}χαρέντα χορείαις·}{}
\ligneRav{982}{}{\hspace{1.1em}ἔξαιρε δὴ προθύμως}{\textbf{982} χαίρειν R : \lecon{χάριν} Ellebodius}
\ligneRav{982}{}{\hspace{1.1em}διπλῆν \difference{χάριν} χορείας·}{\textbf{982} χαίρειν R : \lecon{χάριν} Ellebodius}
\ligneRav{}{}{\hspace{1.1em}παίσωμεν ὦ γυναῖκες οἷάπερ νόμος,}{}
\ligneRav{984}{}{\hspace{1.1em}\difference{νηστεύομεν} δὲ πάντως·}{\textbf{984} νηστεύωμεν R : \lecon{νηστεύομεν} Bentley}
\ligneRav{985}{}{\hspace{1.1em}ἀλλ᾽ \difference{εἶα πάλλ᾽} ἀνάστρεφ᾽ εὐρύθμῳ ποδί,}{\textbf{985} ειαπαλλ᾽ R : \lecon{εἶα πάλλ᾽} Blaydes}
\ligneRav{}{}{\hspace{1.1em}τόρευε πᾶσαν ᾠδήν·}{}
\ligneRav{987}{}{\hspace{1.1em}ἡγοῦ δέ γ᾽ ὧδ᾽ αὐτὸς σὺ}{\textbf{987} γ᾽ ὧδ᾽ suspectum}
\ligneRav{}{}{\hspace{1.1em}κισσοφόρε Βακχεῖε}{}
\ligneRav{989}{}{\hspace{1.1em}δέσποτ᾽· ἐγὼ δὲ κώμοις}{\textbf{989} φιλοχόροισιν R : \lecon{φιλοχόροισι} Scaliger}
\ligneRav{989}{}{\hspace{1.1em}σὲ \difference{φιλοχόροισι} μέλψω·}{\textbf{989} φιλοχόροισιν R : \lecon{φιλοχόροισι} Scaliger}
\ligneRav{990}{}{\hspace{1.1em}†εὔιον ὦ Διόνυσε}{\textbf{990} \lecon{Διόνυσε} R : Διὸς σὺ Enger : Διός τε Fritzsche}
\ligneRav{}{}{\hspace{1.1em}Βρόμιε καὶ Σεμέλας παῖ,}{}
\ligneRav{}{}{\hspace{1.1em}χοροῖς τερπόμενος}{}
\ligneRav{993}{}{\hspace{1.1em}κατ᾽ ὄρεα νυμφᾶν ἐρατοῖς ἐν ὕμνοις,}{\textbf{993} \lecon{ἐρατοῖς ἐν ὕμνοις} R : ἐρατοῖσιν ὕμνοις Wilamowitz}
\ligneRav{994}{}{\hspace{1.1em}†εὔιον εὔιον εὐοῖ}{\textbf{994} lacunam statuit Hermann : ἥκεις γὰρ von Velsen : ἡδόμενος Austin cf. \italique{Plut}.288\espApp \lecon{εὔιον εὔιον} R : ὦ Εὔι Εὔι' Hermann \espApp lacunam statuit Hermann : ἥκεις γὰρ von Velsen : ἡδόμενος Austin cf. \italique{Plut}.288}
\ligneRav{994}{}{\hspace{1.1em}< $ \times \ \widearc{\cup \cup} \ \cup $ > ἀναχορεύων·}{\textbf{994} lacunam statuit Hermann : ἥκεις γὰρ von Velsen : ἡδόμενος Austin cf. \italique{Plut}.288\espApp \lecon{εὔιον εὔιον} R : ὦ Εὔι Εὔι' Hermann \espApp lacunam statuit Hermann : ἥκεις γὰρ von Velsen : ἡδόμενος Austin cf. \italique{Plut}.288}
\ligneRav{995}{}{\hspace{1.1em}ἀμφὶ δὲ \difference{σοὶ κτυπεῖται}}{\textbf{995} συῒ κτυπεῖται R : \lecon{σοὶ κτυπεῖται} Zanetti}
\ligneRav{996}{}{\hspace{1.1em}\difference{Κιθαιρώνιος} ἠχώ,}{\textbf{996} Κιθαρώνιος R : \lecon{Κιθαιρώνιος} Zanetti}
\ligneRav{}{}{\hspace{1.1em}μελάμφυλλά τ᾽ ὄρη}{}
\ligneRav{998}{}{\hspace{1.1em}δάσκια †καὶ νάπαι πετρώδεις† βρέμονται·}{\textbf{998} \lecon{καὶ νάπαι πετρώδεις} R : πετρώδεις τε νάπαι Enger}
\ligneRav{}{}{\hspace{1.1em}κύκλῳ δὲ περὶ σὲ κισσὸς}{}
\ligneRav{1000}{}{\hspace{1.1em}εὐπέταλος ἕλικι θάλλει·}{}
}
\vspace{3pt}\noindent\rule{4.5cm}{0.4pt}\par\vspace{1.5pt}
{\fontsize{7.0pt}{8.3pt}\selectfont
\noindent\textbf{Loci paralleli (vv. 959–1000) ---} \textbf{968} (εὐκύκλου) Hesychius ε 6929 εὔκυκλος· εὔχωρος $=$ Photius ε 2236 \espTest \textbf{973} (τελείαν) Hesychius τ 400 Τελεία \espTest \textbf{986} Suidas τ 789\par
}
\vspace{2.5pt}\noindent\rule{4.5cm}{0.4pt}\par\vspace{1.5pt}
{\fontsize{6.9pt}{8.2pt}\selectfont
\noindent\textbf{Suite (en commentaire) de l'apparat critique \& notes paléographiques (\texttt{f. 174r}) ---} \textbf{Contexte :} Grande ode chorale de danse en ronde (\italique{choros kyklikos}) célébrant Apollon, Artémis, Héra, Hermès, Pan, les Nymphes et Dionysos couronné de lierre (vv. 959–1000). --- Page entièrement lyrique où la colométrie de $\text{R}^1$ remonte directement à l'édition alexandrine d'Aristophane de Byzance (avec \italique{koronis} et \italique{paragraphoi} strophiques). --- \textbf{Fautes phonétiques pittoresques de $\text{R}^1$ :} Au \textbf{v. 969}, haplographie vocalique ἐλύραν pour \lecon{εὐλύραν} (Zanetti) ; au \textbf{v. 982}, χαίρειν pour \lecon{χάριν} (Ellebodius) ; au \textbf{v. 995},iotacisme/bêtacisme συῒ κτυπεῖται (« retentit pour un porc » !) au lieu de \lecon{σοὶ κτυπεῖται} (« retentit pour toi [Dionysos] », Zanetti).\par
}
\end{minipage}
\end{adjustbox}
\end{adjustwidth}
\newpage
\enTeteD{Codex Ravennas 429 (Biblioteca Classense) --- Fac-similé du f. 174r (vv. 959–1000)}
\begin{adjustwidth}{-1.5cm}{-3.0cm}
\begin{center}
\includegraphics[width=16.5cm,height=20.5cm,keepaspectratio]{/workspaces/Thesmophories/dossier_ravennas_thesmophories/folios_individuels_ff162v-177r/24_f174r_classense176r_vv0959-1000.jpg}\par
\vspace{2pt}
{\scriptsize\textbf{Codex Ravennas 429, f. 174r (Classense 176r)} --- ΧΟΡΙΚΟΝ, \textbf{vv. 959–1000} ($\text{R}^1, \Rpc, \Rdeux$ (\Mon)).}
\end{center}
\end{adjustwidth}
\newpage
\clearpage\ifodd\value{page}\hbox{}\newpage\fi
\enTeteG{ΘΕΣΜΟΦΟΡΙΑΖΟΥΣΑΙ --- Texte grec établi (f. 174r, vv. 959–1000)}
\enTeteD{LES THESMOPHORIES --- Traduction française (f. 174r, vv. 959–1000)}
\begin{pages}
\begin{Leftside}
\beginnumbering
%% v. 959-961 strophe a
\begin{groupe}%
\Strophe{(στροφή α)}\typeVers{cret. $+$ lec.}\chantChoeur{\ligne{959}ἅμα δὲ καὶ \ligne{960}γένος Ὀλυμπίων θεῶν}
\end{groupe}%
%% v. 959-961 strophe a
\begin{groupe}%
\typeVers{$4$ troch\cat}\chantChoeur{μέλπε καὶ γέραιρε φωνῇ πᾶσα χορομανεῖ τρόπῳ.\\\passeLigne}%
\end{groupe}%
%% v. 962-965 antistrophe a
\begin{groupe}%
\Strophe{(ἀντιστρ. α)}\typeVers{cret. $+$ lec.}\chantChoeur{εἰ δέ τις \ligne{963}προσδοκᾷ κακῶς ἐρεῖν}
\end{groupe}%
%% v. 962-965 antistrophe a
\begin{groupe}%
\typeVers{$4$ troch\cat}\chantChoeur{ἐν ἱερῷ γυναῖκά μ᾽ οὖσαν ἄνδρας,\ligne{965} οὐκ ὀρθῶς φρονεῖ.\\\passeLigne}%
\end{groupe}%
%% v. 966-968 mesodos a
\begin{groupe}%
\Strophe{(μεσῳδός α)}\typeVers{cret.}\chantChoeur{ἀλλὰ χρή\apparat{966}{\lecon{χρὴ} R : χρῆν Bothe : χρή μ᾽ Meineke}}
\end{groupe}%
%% v. 966-968 mesodos a
\begin{groupe}%
\typeVers{dim. troch.}\chantChoeur{ὥσπερ ἔργον αὖ τι καινὸν%\apparat{967}{\lecon{ὥσπερ} R : ὡς πρὸς Austin}
}
\end{groupe}%
%% v. 966-968 mesodos a
\begin{groupe}%
\typeVers{$4$ troch\cat}\chantChoeur{{\Taille{11.5}πρῶτον εὐκύκλου χορείας εὐφυᾶ στῆσαι βάσιν.}\testimonia{968}{(εὐκύκλου) Hesychius ε 6929 εὔκυκλος· εὔχωρος $=$ Photius ε 2236}\apparat{968}{εὐφυῆ R : \lecon{εὐφυᾶ} Brunck}\\\passeLigne}%
\end{groupe}%
%% v. 969-976 strophe b
\begin{groupe}%
%
\Strophe{(στροφή β)}\typeVers{dim. iamb.}\chantChoeur{\ligne{969}πρόβαινε ποσὶ τὸν εὐλύραν\apparat{969}{ποσὶν R : \lecon{ποσὶ} Bentley \espApp ἐλύραν R : \lecon{εὐλύραν} Zanetti}}
\end{groupe}%
%% v. 969-976 strophe b
\begin{groupe}%
\typeVers{dim. iamb.}\chantChoeur{μέλπουσα καὶ τὴν τοξοφόρον}
\end{groupe}%
%% v. 969-976 strophe b
\begin{groupe}%
\typeVers{$2$ iamb\cat}\chantChoeur{Ἄρτεμιν ἄνασσαν ἁγνήν.}
\end{groupe}%
%% v. 969-976 strophe b
\begin{groupe}%
\typeVers{reiz.}\chantChoeur{χαῖρ᾽ ὦ ἑκάεργε,}
\end{groupe}%
%% v. 969-976 strophe b
\begin{groupe}%
\typeVers{reiz.}\chantChoeur{ὄπαζε δὲ νίκην·}
\end{groupe}%
%% v. 969-976 strophe b
\begin{groupe}%
\typeVers{$2$ iamb\cat}\chantChoeur{Ἥραν τε τὴν τελείαν\testimonia{973}{(τελείαν) Hesychius τ 400 Τελεία}}
\end{groupe}%
%% v. 969-976 strophe b
\begin{groupe}%
\typeVers{$2$ iamb\cat}\chantChoeur{μέλψωμεν ὥσπερ εἰκός,}
\end{groupe}%
%% v. 969-976 strophe b
\begin{groupe}%
\typeVers{trim. iamb.}\chantChoeur{ἣ πᾶσι τοῖς χοροῖσιν ἐμπαίζει τε καὶ}
\end{groupe}%
%% v. 969-976 strophe b
\begin{groupe}%
\typeVers{$2$ iamb\cat}\chantChoeur{κλῇδας γάμου φυλάττει.\\\passeLigne}%
\end{groupe}%
%% v. 977-984 antistrophe b
\begin{groupe}%
\Strophe{(ἀντιστρ. β)}\typeVers{dim. iamb.}\chantChoeur{Ἑρμῆν τε νόμιον ἄντομαι}
\end{groupe}%
%% v. 977-984 antistrophe b
\begin{groupe}%
\typeVers{dim. iamb.}\chantChoeur{καὶ Πᾶνα καὶ Νύμφας φίλας}
\end{groupe}%
%% v. 977-984 antistrophe b
\begin{groupe}%
\typeVers{$2$ iamb\cat}\chantChoeur{\ligne{979}ἐπιγελάσαι προθύμως}
\end{groupe}%
%% v. 977-984 antistrophe b
\begin{groupe}%
\typeVers{reiz.}\chantChoeur{ταῖς ἡμετέραισι\apparat{980}{ἡμετέραις R : \lecon{ἡμετέραισι} Hermann}}
\end{groupe}%
%% v. 977-984 antistrophe b
\begin{groupe}%
\typeVers{reiz.}\chantChoeur{χαρέντα χορείαις.}
\end{groupe}%
%% v. 977-984 antistrophe b
\begin{groupe}%
\typeVers{$2$ iamb\cat}\chantChoeur{ἔξαιρε δὴ προθύμως}
\end{groupe}%
%% v. 977-984 antistrophe b
\begin{groupe}%
\typeVers{$2$ iamb\cat}\chantChoeur{\ligne{982}διπλῆν χάριν χορείας.\apparat{982}{χαίρειν R : \lecon{χάριν} Ellebodius}%
\commentaire{982}{διπλῆν.}{Pollux (IV, 105) compte la διπλῆ parmi les figures de danse pour la tragédie (τραγικῆς ὀρχήσεως σχήματα) :%
\citComm{Καὶ μὲν τραγικῆς ὀρχήσεως σχήματα σιμή χείρ, καλαθίσκος, χεὶρ καταπρανής, ξύλου παράληψις, \textbf{διπλῆ}, θερμαυστρίς, κυβίστησις, παραβῆναι τὰ τέτταρα}%
Hésychius (δ 1948) ne nous apprend rien davantage :%
\citComm{διπλῆ· ὀρχήσεως εἶδος, ἢ κρούματος.}%
Il nous est difficile de reconstituer, même grossièrement, la chorégraphie de ce choral puisqu'Aristophane recourt au vocabulaire technique de la danse de son époque, dont on ne sait plus grand chose.}}
\end{groupe}%
%% v. 977-984 antistrophe b
\begin{groupe}%
\typeVers{trim. iamb.}\chantChoeur{παίσωμεν ὦ γυναῖκες οἷάπερ νόμος,}
\end{groupe}%
%% v. 977-984 antistrophe b
\begin{groupe}%
\typeVers{$2$ iamb\cat}\chantChoeur{νηστεύομεν δὲ πάντως.\apparat{984}{νηστεύωμεν R : \lecon{νηστεύομεν} Bentley}\\\passeLigne}%
\end{groupe}%
%% v. 985-989 mesodos b
\begin{groupe}%
\Strophe{(μεσῳδός β)}\typeVers{trim. iamb.}\chantChoeur{\ligne{985}ἀλλ᾽ εἶα πάλλ᾽ ἀνάστρεφ᾽ εὐρύθμῳ ποδί,\apparat{985}{ειαπαλλ᾽ R : \lecon{εἶα πάλλ᾽} Blaydes}}
\end{groupe}%
%% v. 985-989 mesodos b
\begin{groupe}%
\typeVers{$2$ iamb\cat}\chantChoeur{τόρευε πᾶσαν ᾠδήν·\testimonia{986}{Suidas τ 789}}
\end{groupe}%
%% v. 985-989 mesodos b
\begin{groupe}%
\typeVers{ia. $+$ mol.}\chantChoeur{ἡγοῦ δέ γ᾽ ὧδ᾽ αὐτὸς σὺ%\apparat{987}{γ᾽ ὧδ᾽ suspectum}
}
\end{groupe}%
%% v. 985-989 mesodos b
\begin{groupe}%
\typeVers{cret. $+$ mol.}\chantChoeur{κισσοφόρε Βακχεῖε}
\end{groupe}%
%% v. 985-989 mesodos b
\begin{groupe}%
\typeVers{aristoph.}\chantChoeur{δέσποτ᾽· ἐγὼ δὲ κώμοις}
\end{groupe}%
%% v. 985-989 mesodos b
\begin{groupe}%
\typeVers{$2$ iamb\cat}\chantChoeur{\ligne{989}σὲ φιλοχόροισι μέλψω.\apparat{989}{φιλοχόροισιν R : \lecon{φιλοχόροισι} Scaliger}\\\passeLigne}%
\end{groupe}%
%% v. 990-994 strophe g
\begin{groupe}%
\Strophe{(στροφή γ)}\typeVers{aristoph.}\chantChoeur{\ligne{990}†εὔιον ὦ Διόνυσε\apparat{990}{%\lecon{εὔιον} R : Εὔιε Enger : Εὐοῖ Hermann (1845) \espApp 
\lecon{Διόνυσε} R : Διὸς σὺ Enger : Διός τε Fritzsche}%
%\commentaire{990}{}{\eff{Hésychius ε 6909 εὐΐων· Διονυσιακῶν} (Ευρ. Βαϲϲ 791...) εὐϊάζω = εὐάζω Eur. Cycl. 495.}
}
\end{groupe}%
%% v. 990-994 strophe g
\begin{groupe}%
\typeVers{pherecr.}\chantChoeur{Βρόμιε καὶ Σεμέλας παῖ,}
\end{groupe}%
%% v. 990-994 strophe g
\begin{groupe}%
\typeVers{hexasyll. aeol.}\chantChoeur{χοροῖς τερπόμενος}
\end{groupe}%
%% v. 990-994 strophe g
\begin{groupe}%
\typeVers{ia. $+$ arist.}\chantChoeur{κατ᾽ ὄρεα νυμφᾶν ἐρατοῖς ἐν ὕμνοις,\apparat{993}{\lecon{ἐρατοῖς ἐν ὕμνοις} R : ἐρατοῖσιν ὕμνοις Wilamowitz %(\italique{Gr. Versk. p.476}) 
\ligneApp{994} %\lecon{εὔιον εὔιον} R : ὦ Εὔι Εὔι' Hermann \espApp 
lacunam statuit Hermann : ἥκεις γὰρ von Velsen : ἡδόμενος Austin cf. \italique{Plut}.288}}
\end{groupe}%
%% v. 990-994 strophe g
\begin{groupe}%
\typeVers{$2$ iamb\cat}\chantChoeur{†εὔιον εὔιον εὐοῖ%\apparat{994}{\lecon{εὔιον εὔιον} R : ὦ Εὔι Εὔι' Hermann \espApp lacunam statuit Hermann : ἥκεις γὰρ von Velsen : ἡδόμενος Austin cf. \italique{Plut}.288}
}
\end{groupe}%
%% v. 990-994 strophe g
\begin{groupe}%
%\times \  \overarc{\smile \smile} \widearc{\smile \smile} \overset{\huge{\frown}}{\cup \cup} \overset{\huge{\frown}}{\smile \smile \:} \overset{\huge{\frown}}{\smile \smile} \ \cup
% \widearc{\smile \smile}
\typeVers{$2$ iamb\cat}\chantChoeur{\ligne{994}< $ \times \ \widearc{\cup \cup} \ \cup $ > ἀναχορεύων.\\\passeLigne}%
\end{groupe}%
%% v. 995-1000 antistrophe g
\begin{groupe}%
\Strophe{(ἀντιστρ. γ)}\typeVers{aristoph.}\chantChoeur{ἀμφὶ δὲ σοὶ κτυπεῖται\apparat{995}{συῒ κτυπεῖται R : \lecon{σοὶ κτυπεῖται} Zanetti %: συγκτυπεῖται Wilamowitz 
\ligneApp{996} Κιθαρώνιος R : \lecon{Κιθαιρώνιος} Zanetti}}
\end{groupe}%
%% v. 995-1000 antistrophe g
\begin{groupe}%
\typeVers{pherecr.}\chantChoeur{Κιθαιρώνιος ἠχώ,}
\end{groupe}%
%% v. 995-1000 antistrophe g
\begin{groupe}%
\typeVers{hexasyll. aeol.}\chantChoeur{μελάμφυλλά τ᾽ ὄρη}
\end{groupe}%
%% v. 995-1000 antistrophe g
\begin{groupe}%
\typeVers{(?)}\chantChoeur{δάσκια †καὶ νάπαι πετρώδεις† βρέμονται·\apparat{998}{\lecon{καὶ νάπαι πετρώδεις} R : πετρώδεις τε νάπαι Enger}}
\end{groupe}%
%% v. 995-1000 antistrophe g
\begin{groupe}%
\typeVers{$2$ iamb\cat}\chantChoeur{κύκλῳ δὲ περὶ σὲ κισσὸς}
\end{groupe}%
%% v. 995-1000 antistrophe g
\begin{groupe}%
\typeVers{$2$ iamb\cat}\chantChoeur{εὐπέταλος ἕλικι θάλλει.%
\\\passeLigne}%
\end{groupe}%
\endnumbering
\end{Leftside}
\begin{Rightside}
\beginnumbering
%% v. 959-961 strophe a
\begin{groupe}%
\chantChoeurFr{Et ensemble célèbre et honore de la voix la gent des dieux olympiens,}
\end{groupe}%
%% v. 959-961 strophe a
\begin{groupe}%
\chantChoeurFr{d'une âme transportée par la danse.\\\passeLigne}%
\end{groupe}%
%% v. 962-965 antistrophe a
\begin{groupe}%
\chantChoeurFr{Que si tel croit\ligne{963} qu'étant femme je dois médire}
\end{groupe}%
%% v. 962-965 antistrophe a
\begin{groupe}%
\chantChoeurFr{des hommes pendant nos rites,\ligne{965} il s'abuse.\\\passeLigne}%
\end{groupe}%
%% v. 966-968 mesodos a
\begin{groupe}%
\chantChoeurFr{Mais d'abord il convient,}
\end{groupe}%
%% v. 966-968 mesodos a
\begin{groupe}%
\chantChoeurFr{c'est la règle,}
\end{groupe}%
%% v. 966-968 mesodos a
\begin{groupe}%
\chantChoeurFr{de procéder aux figures élégantes de la danse cyclique. (Danse.)\\\passeLigne}%
\end{groupe}%
%% v. 969-976 strophe b
\begin{groupe}%
\chantChoeurFr{\ligne{969}Pars du pied en célébrant le dieu à la belle lyre}
\end{groupe}%
%% v. 969-976 strophe b
\begin{groupe}%
\chantChoeurFr{et Artémis armée de l'arc, la chaste souveraine.}
\end{groupe}%
%% v. 969-976 strophe b
\begin{groupe}%
\chantChoeurFr{Salut, dieu tutélaire,}
\end{groupe}%
%% v. 969-976 strophe b
\begin{groupe}%
\chantChoeurFr{et donne-nous la victoire.}
\end{groupe}%
%% v. 969-976 strophe b
\begin{groupe}%
\chantChoeurFr{Puis célébrons, comme il sied,}
\end{groupe}%
%% v. 969-976 strophe b
\begin{groupe}%
\chantChoeurFr{Héra la conjugale,}
\end{groupe}%
%% v. 969-976 strophe b
\begin{groupe}%
\chantChoeurFr{qui se récrée}
\end{groupe}%
%% v. 969-976 strophe b
\begin{groupe}%
\chantChoeurFr{avec tous les chœurs et}
\end{groupe}%
%% v. 969-976 strophe b
\begin{groupe}%
\chantChoeurFr{garde les clefs de l'hymen.\\\passeLigne}%
\end{groupe}%
%% v. 977-984 antistrophe b
\begin{groupe}%
\chantChoeurFr{Et Hermès pastoral,}
\end{groupe}%
%% v. 977-984 antistrophe b
\begin{groupe}%
\chantChoeurFr{et Pan et les Nymphes chéries,}
\end{groupe}%
%% v. 977-984 antistrophe b
\begin{groupe}%
\chantChoeurFr{je les prie de sourire avec bienveillance}
\end{groupe}%
%% v. 977-984 antistrophe b
\begin{groupe}%
\chantChoeurFr{et de se plaire}
\end{groupe}%
%% v. 977-984 antistrophe b
\begin{groupe}%
\chantChoeurFr{à nos danses.}
\end{groupe}%
%% v. 977-984 antistrophe b
\begin{groupe}%
\chantChoeurFr{Allons, commence avec ardeur}
\end{groupe}%
%% v. 977-984 antistrophe b
\begin{groupe}%
\chantChoeurFr{la diple, joie de la danse.}
\end{groupe}%
%% v. 977-984 antistrophe b
\begin{groupe}%
\chantChoeurFr{\ligne{983}Ébattons-nous, femmes, suivant l'usage.}
\end{groupe}%
%% v. 977-984 antistrophe b
\begin{groupe}%
\chantChoeurFr{Aussi bien nous sommes à jeun.\\\passeLigne}%
\end{groupe}%
%% v. 985-989 mesodos b
\begin{groupe}%
\chantChoeurFr{Or sus! marque une autre mesure de ton pied \fin cadencé.}
\end{groupe}%
%% v. 985-989 mesodos b
\begin{groupe}%
\chantChoeurFr{Chante à pleine voix.}
\end{groupe}%
%% v. 985-989 mesodos b
\begin{groupe}%
\chantChoeurFr{Et guide-nous toi-même,}
\end{groupe}%
%% v. 985-989 mesodos b
\begin{groupe}%
\chantChoeurFr{couronné de lierre, seigneur Bacchus.}
\end{groupe}%
%% v. 985-989 mesodos b
\begin{groupe}%
\chantChoeurFr{\ligne{988}Je te célébrerai}
\end{groupe}%
%% v. 985-989 mesodos b
\begin{groupe}%
\chantChoeurFr{en des chants bien rythmés.\\\passeLigne}%
\end{groupe}%
%% v. 990-994 strophe g
\begin{groupe}%
\chantChoeurFr{Évios! fils de Zeus}
\end{groupe}%
%% v. 990-994 strophe g
\begin{groupe}%
\chantChoeurFr{et de Sémélé, ô Bromios,}
\end{groupe}%
%% v. 990-994 strophe g
\begin{groupe}%
\chantChoeurFr{toi que les chœurs délectent,}
\end{groupe}%
%% v. 990-994 strophe g
\begin{groupe}%
\chantChoeurFr{et qui dans les montagnes, parmi les hymnes \fin aimables des Nymphes,}
\end{groupe}%
%% v. 990-994 strophe g
\begin{groupe}%
\chantChoeurFr{Évios ! Évios ! Évoé !}
\end{groupe}%
%% v. 990-994 strophe g
\begin{groupe}%
\chantChoeurFr{\ligne{994}te livres à la danse.\\\passeLigne}%
\end{groupe}%
%% v. 995-1000 antistrophe g
\begin{groupe}%
\chantChoeurFr{Autour de toi}
\end{groupe}%
%% v. 995-1000 antistrophe g
\begin{groupe}%
\chantChoeurFr{l'écho du Cithéron résonne,}
\end{groupe}%
%% v. 995-1000 antistrophe g
\begin{groupe}%
\chantChoeurFr{et retentissent les montagnes ombreuses}
\end{groupe}%
%% v. 995-1000 antistrophe g
\begin{groupe}%
\chantChoeurFr{et les ravins rocheux,}
\end{groupe}%
%% v. 995-1000 antistrophe g
\begin{groupe}%
\chantChoeurFr{tandis qu'à tes côtés le lierre}
\end{groupe}%
%% v. 995-1000 antistrophe g
\begin{groupe}%
\chantChoeurFr{aux belles feuilles pousse sa vrille tout à l'entour.%
\\\passeLigne}%
\end{groupe}%
\endnumbering
\end{Rightside}
\end{pages}
\Pages
\clearpage\ifodd\value{page}\hbox{}\newpage\fi
\enTeteG{ΕΠΕΙΣΟΔΙΟΝ Ε' (Ἀνδρομέδα) --- Codex Ravennas 429, f. 174v (Classense 176v) --- vv. 1001–1043}
\begin{adjustwidth}{-1.5cm}{-3.0cm}
\begin{adjustbox}{max totalheight=21.0cm, max width=16.5cm, keepaspectratio}
\begin{minipage}{16.5cm}
{\fontsize{7.2pt}{8.5pt}\selectfont
\noindent\textbf{Folio 25 --- Codex Ravennas 429, \texttt{f. 174v} (\texttt{Classense 176v}) : ΕΠΕΙΣΟΔΙΟΝ Ε' (Ἀνδρομέδα), vv. 1001–1043}\hfill \textit{Témoins :} $\text{R}^1, \Rpc, \Rdeux$, \PsiQuatre\par\vspace{2pt}
\hrule height 0.4pt\vspace{2.5pt}
}
{\fontsize{8.1pt}{9.5pt}\selectfont
\ligneRav{1001}{τοξ/}{ἐνταῦτα νῦν \difference{οἰμῶξι} πρὸς τὴν αἰτρίαν·}{\textbf{1001} οἰμώξει R : \lecon{οἰμῶξι} Brunck}
\ligneRav{1002}{μ\auDessus{ν}{Η}}{ὦ τοξόθ᾽ ἱκετεύω σε. : μή μ᾽ ἰκετεῦσι σύ·}{\textbf{1002} ἱκετεῦση R : \lecon{ἱκετεῦσι} Brunck}
\ligneRav{1003}{μ\auDessus{ν}{Η}}{χάλασον τὸν ἧλον. : ἀλλὰ ταῦτα \difference{δρᾶσ᾽} ἐγώ·}{\textbf{1003} δρᾷς R : \lecon{δρᾶσ᾽} Voss}
\ligneRav{1004}{μ\auDessus{ν}{Η}}{οἴμοι κακοδαίμων, μᾶλλον \difference{ἐπικρούεις} σύ γε·}{\textbf{1004} ἐπικρούσεις R : \lecon{ἐπικρούεις} \Mon}
\ligneRav{1005}{τοξ/}{ἔτι \difference{μᾶλλο} βοῦλις; : ἀτταταῖ \difference{ἰατταταῖ}·}{\textbf{1005} μᾶλλον R : \lecon{μᾶλλο} Bentley\espApp τατταταί R : \lecon{ἰατταταῖ} Faber}
\ligneRav{1006}{—}{κακοδαίμω\expo{ν} γέρ\inRas{ω}{ο}ν·}{\textbf{1006} κακοδαίμω\expo{ν} γέρ\inRas{ω}{ο}ν \Rpc, M\expo{pc}}
\ligneRav{1007}{}{πέρ᾽ ἐγὼ \difference{'ξινίγκι} πορμός, \difference{ἴνα} πυλάξι σοι·}{\textbf{1007} ξεινίγκι R : \lecon{'ξινίγκι} Enger \espApp ἵνα R : \lecon{ἴνα} Brunck \espApp \lecon{πυλάξι} R : πυλάξῃ Σ}
\ligneRav{}{μ\auDessus{ν}{Η}}{ταυτὶ τὰ βέλτιστ᾽ ἀπολέλαυκ᾽ Εὐριπίδου·}{}
\ligneRav{1009}{}{ἔα· θεοί, Ζεῦ σῶτερ, εἰσὶν ἐλπίδες·}{\textbf{1009} εἰσὶν R : εἴς ἔτ᾽ Naber (Faber ?) \cf 946}
\ligneRav{1010}{}{\difference{ἁνὴρ} ἔοικεν οὐ προδώσειν, ἀλλά μοι}{\textbf{1010} ἀνὴρ R : \lecon{ἁνὴρ} Bothe}
\ligneRav{1011}{}{σημεῖον \difference{ὑπεδήλωσε} Περσεὺς ἐκδραμών,}{\textbf{1011} ὑπεδήλωσεν R : \lecon{ὑπεδήλωσε} Biset}
\ligneRav{}{}{ὅτι δεῖ με γίγνεσθ᾽ Ἀνδρομέδαν· πάντως δέ μοι}{}
\ligneRav{1013}{}{τὰ δέσμ᾽ ὑπάρχει. δῆλον \difference{οὖν τοῦτ᾽} ἔσθ᾽ ὅτι}{\textbf{1013} οὖν R : οὖν νῦν Bizet : \lecon{οὖν τοῦτ᾽} Dobree}
\ligneRav{1014}{}{ἥξει με \difference{σώσων}· οὐ γὰρ ἂν παρέπτετο·}{\textbf{1014} σῶσον R : \lecon{σώσων} \Mon}
\ligneRav{1015}{\auDessus{εὐ}{ρ}}{\hspace{1.1em}φίλαι παρθένοι φίλαι,}{\textbf{1015-21} (Εὐρ.) R : \lecon{(Μνη.)} Tyrwhitt}
\ligneRav{1016}{}{\hspace{1.1em}πῶς ἂν ἀπέλθοιμι καὶ}{\textbf{1016} \lecon{ἀπέλθοιμι} R : ἐπέλθοιμι Ellebodius}
\ligneRav{1017}{}{\hspace{1.1em}τὸν Σκύθην \difference{λάθοιμι};}{\textbf{1017} λάβοιμι R : \lecon{λάθοιμι} Ellebodius}
\ligneRav{1019}{}{\hspace{1.1em}κλύεις; ὦ \difference{προσᾴδουσ'} \difference{ἀϋτὰς} ἐν ἄντροις,}{\textbf{1019} προσᾳδοῦσσαι τὰς R, Σ : \lecon{προσᾴδουσ'} Elmsley, \lecon{ἀϋτὰς} Burges (ἀϋταῖς Sommerstein) : alii alia}
\ligneRav{1020}{}{\hspace{1.1em}κατάνευσον, ἔασον ὡς}{}
\ligneRav{}{}{\hspace{1.1em}τὴν γυναῖκά μ᾽ ἐλθεῖν·}{}
\ligneRav{1022}{μ\auDessus{ν}{Η}}{\hspace{1.1em}ἄνοικτος ὅς μ᾽ ἔδησε τὸν}{\textbf{1022} (Μνη.) praefixit \Rpc, \Mon}
\ligneRav{1023}{}{\hspace{1.1em}πολυπονώτατον βροτῶν·}{\textbf{1023} \lecon{πολυπονώτατον} R, Σ : πολυστονώτατον Burges}
\ligneRav{}{}{\hspace{1.1em}μόλις δὲ γραῖαν ἀποφυγὼν}{}
\ligneRav{1025}{}{\hspace{1.1em}σαπρὰν ἀπωλόμην ὅμως·}{}
\ligneRav{1026}{}{\hspace{1.1em}ὅδε γὰρ ὁ Σκύθης φύλαξ πάλαι ἐφέστηκ᾽· ὀλοὸν}{\textbf{1026} ὁ Σκύθης R : del. Jackson}
\ligneRav{}{}{\hspace{1.1em}ἄφιλον ἐκρέμασε κόραξι δεῖπνον·}{}
\ligneRav{1030}{}{\hspace{1.1em}ὁρᾷς; οὐ χοροῖσιν οὐδ᾽ ὑφ᾽ ἡλίκων νεανίδων}{}
\ligneRav{1031}{}{\hspace{1.1em}ψῆφον \difference{ἕστηκ᾽ ἔχουσα κημόν},}{\textbf{1031} κημὸν ἕστηκ᾽ ἔχουσ᾽ R : \lecon{ἕστηκ᾽ ἔχουσα κημόν} Bibi, cf. Σ ad \italique{Vesp.} 674c}
\ligneRav{}{}{\hspace{1.1em}ἀλλ᾽ ἐν πυκνοῖς δεσμοῖσιν ἐμπεπλεγμένη}{}
\ligneRav{}{}{\hspace{1.1em}κήτει βορὰ Γλαυκέτῃ πρόκειμαι,}{}
\ligneRav{1034}{}{\hspace{1.1em}γαμηλίῳ μὲν οὐ ξὺν}{\textbf{1034-5} ξυμπαιῶνι R : \lecon{ξὺν παιῶνι} Zanetti (\italique{cum carmuine} Divus)}
\ligneRav{1035}{}{\hspace{1.1em}παιῶνι δεσμίῳ δὲ}{}
\ligneRav{}{}{\hspace{1.1em}γοᾶσθέ μ᾽ ὦ γυναῖκες ὡς}{}
\ligneRav{}{}{\hspace{1.1em}μέλεα μὲν πέπονθα μέλεος,}{}
\ligneRav{}{}{\hspace{1.1em}ὢ τάλας ἐγὼ τάλας,}{}
\ligneRav{1039}{}{\hspace{1.1em}ἀπὸ δὲ συγγόνων \difference{ἀλλ᾽} ἄνομα πάθεα,}{\textbf{1039} ἀλλὰν R : \lecon{ἀλλ᾽} Scaliger}
\ligneRav{1040}{}{\hspace{1.1em}†φῶτά \difference{λιτομένα}†,}{\textbf{1040} λιτομέναν R : ἀντομένα (vel -αν) Σ : \lecon{λιτομένα} Enger}
\ligneRav{1041}{}{\hspace{1.1em}πολυδάκρυτον Ἀίδα γόον \difference{ψεφαυγοῦς}·}{\textbf{1041} φεύγουσαν R : χέουσα Rau (\italique{Paratragödia}, p.75) : φλέγουσαν Mugrave (ad Eur. \italique{Or.}1402$=$1396) : \lecon{ψεφαυγοῦς} Bibi, cf Σ ad v.1054, Hesych. ψ 135\espApp φεύγουσαν R : χέουσα Rau (\italique{Paratragödia}, p.75) : φλέγουσαν Mugrave (ad Eur. \italique{Or.}1402$=$1396) : \lecon{ψεφαυγοῦς} Bibi, cf Σ ad v.1054 e Hesych. ψ 135}
\ligneRav{}{}{\hspace{1.1em}αἰαῖ αἰαῖ ἒ ἒ,}{}
\ligneRav{}{}{\hspace{1.1em}ὃς ἔμ᾽ ἀπεξύρησε πρῶτον,}{}
}
\vspace{3pt}\noindent\rule{4.5cm}{0.4pt}\par\vspace{1.5pt}
{\fontsize{7.0pt}{8.3pt}\selectfont
\noindent\textbf{Loci paralleli (vv. 1001–1043) ---} Tradition directe : Codex Ravennas 429 (R) et scholies anciennes ($\Sigma^{\text{R}}$).\par
}
\vspace{2.5pt}\noindent\rule{4.5cm}{0.4pt}\par\vspace{1.5pt}
{\fontsize{6.9pt}{8.2pt}\selectfont
\noindent\textbf{Suite (en commentaire) de l'apparat critique \& notes paléographiques (\texttt{f. 174v}) ---} \textbf{Témoin papyrologique :} \textbf{\PsiQuatre{} (\italique{P.Oxy.} LXXIII 4935, fr. 1)} débute au \textbf{v. 1043}. --- \textbf{Traitement du grec barbare de l'Archer scythe (vv. 1001–1007) :} $\text{R}^1$ tend parfois à normaliser les finales scythes en grec standard (οἰμώξει, ἱκετεῦση, μᾶλλον, ἵνα) ; au \textbf{v. 1006}, \Rpc{} (suivi de \Mon\expo{p.c.}) ajoute un $\nu$ en exposant et corrige \italique{in rasura} l'oméga de γέρων en omicron : \lecon{κακοδαίμω\expo{ν} γέρ\inRas{ω}{ο}ν}. --- \textbf{Attribution des vv. 1015–1021 :} R attribue le début de la monodie à (Εὐρ.), corrigé en \lecon{(Μνη.)} par Tyrwhitt, tandis qu'au \textbf{v. 1022} \Rpc{} (et \Mon{}) préfixe (Μνη.). --- \textbf{Deux interventions décisives de \textbf{Bibi} (vv. 1031 et 1041) :} --- \textbf{v. 1007 :} ξεινίγκι R : \lecon{'ξινίγκι} Enger \espApp ἵνα R : \lecon{ἴνα} Brunck \espApp \lecon{πυλάξι} R (πυλάξῃ $\Sigma$). --- \textbf{v. 1017 :} λάβοιμι R (confusion minuscule $\theta/\beta$) : \lecon{λάθοιμι} Ellebodius. --- \textbf{v. 1019 :} προσᾳδοῦσσαι τὰς R, $\Sigma$ : \lecon{προσᾴδουσ'} Elmsley, \lecon{ἀϋτὰς} Burges (ἀϋταῖς Sommerstein). --- \textbf{v. 1023 :} \lecon{πολυπονώτατον} R, $\Sigma$ (πολυστονώτατον Burges). --- \textbf{v. 1031 :} κημὸν ἕστηκ' ἔχουσ' R : \lecon{ἕστηκ' ἔχουσα κημόν} \lecon{Bibi}, cf. $\Sigma$ ad \italique{Vesp.} 674c. --- \textbf{v. 1040 :} λιτομέναν R (ἀντομένα vel -αν $\Sigma$, confusion onciale ΛΙ/ΑΝ) : \lecon{λιτομένα} Enger. --- \textbf{v. 1041 :} φεύγουσαν R (χέουσα Rau ; φλέγουσαν Musgrave) : \lecon{ψεφαυγοῦς} \lecon{Bibi}, cf. $\Sigma$ ad v. 1054, Hésychius $\psi\ 135$.\par
}
\end{minipage}
\end{adjustbox}
\end{adjustwidth}
\newpage
\enTeteD{Codex Ravennas 429 (Biblioteca Classense) --- Fac-similé du f. 174v (vv. 1001–1043)}
\begin{adjustwidth}{-1.5cm}{-3.0cm}
\begin{center}
\includegraphics[width=16.5cm,height=20.5cm,keepaspectratio]{/workspaces/Thesmophories/dossier_ravennas_thesmophories/folios_individuels_ff162v-177r/25_f174v_classense176v_vv1001-1043.jpg}\par
\vspace{2pt}
{\scriptsize\textbf{Codex Ravennas 429, f. 174v (Classense 176v)} --- ΕΠΕΙΣΟΔΙΟΝ Ε' (Ἀνδρομέδα), \textbf{vv. 1001–1043} ($\text{R}^1, \Rpc, \Rdeux$, \PsiQuatre).}
\end{center}
\end{adjustwidth}
\newpage
\clearpage\ifodd\value{page}\hbox{}\newpage\fi
\enTeteG{ΘΕΣΜΟΦΟΡΙΑΖΟΥΣΑΙ --- Texte grec établi (f. 174v, vv. 1001–1043)}
\enTeteD{LES THESMOPHORIES --- Traduction française (f. 174v, vv. 1001–1043)}
\begin{pages}
\begin{Leftside}
\beginnumbering
%% v. 1001
\begin{groupe}%
%
\personnageGrec{Τοξότης}%
\typeVers{trim. iamb.}\trimIamb{ἐνταῦτα νῦν οἰμῶξι πρὸς τὴν αἰτρίαν.\apparat{1001}{οἰμώξει R : \lecon{οἰμῶξι} Brunck}\\\passeLigne}%
\end{groupe}%
%% v. 1002a
\begin{groupe}%
\personnageGrec{Μνησίλοχος}%
\trimIamb{ὦ τοξόθ᾽ ἱκετεύω σε.%
\commentaire{1001}{τοξότης.}{Aristophane a donné à l'archer un jargon mi-scythe mi-grec, une sorte de parler nègre. On sent qu'il est impossible de rendre cela dans une traduction. (Note de Willems.)\\
À Athènes, comme partout ailleurs, les personnes chargées des missions de police étaient recrutées parmi la lie de l'humanité. En l'occurrence, il y avait en guise de police un contingent de 300 mercenaires Scythes aux ordres des prytanes\marqueur{1}. C'est cette inquiétant contraste, entre l'importance de la mission (les forces de l'«~ordre~», comme on dit...) et la bassesse de ceux qui y sont préposés, qui inspirera à Platon de si longs développements sur le bon choix et l'éducation des gardiens de sa Cité. \latin{Quis custodiet custodes?}}%
\debutNote{1001}%
\note{1}{Que les archers Scythes fussent directement aux ordres des prytanes, c'est assez bien établi par notre pièce ainsi que le vers 54 des \italique{Acharniens}, cf. aussi \italique{Protagoras} (319c).
On pourra lire l'article de Jean-Christophe Couvenhes «~L'introduction des archers scythes, esclaves publics, à Athènes~», in \italique{Transferts culturels et droits dans le monde grec et hellénistique} (2012), pp. 99-118, ou l'article plus ancien d'Agnès Plassart «~Les archers d'Athènes~», in \italique{Revue des Études Grecques}, t. 26 (1913), pp. 151-213 dont quelques détails sont discutés dans le premier article.}%
%\note{2}{}%
\commentaire{1001}{ἐνταῦτα νῦν οἰμῶξι πρὸς τὴν αἰτρίαν}{$=$ ἐνταῦθα νῦν οἴμωζε πρὸς τὴν αἰθρίαν.
(Noter le sadisme proprement argousin de l'archer.)% : il dit à Mnésiloque qu'il va pouvoir gémir au grand air (οἰμῶξι πρὸς τὴν αἰτρίαν) au moment même où il lui coupe le sifflet et serre la vis qui l'étrangle.
}\\\passeLigne}%
\end{groupe}%
%% v. 1002b
\begin{groupe}%
\personnageGrec{Τοξότης}%
\trimIamb{\ligne{1002}\decalage{ὦ τοξόθ᾽ ἱκετεύω σε.}μή μ᾽ ἰκετεῦσι σύ.\apparat{1002}{ἱκετεῦση R : \lecon{ἱκετεῦσι} Brunck}\\\passeLigne}%
\end{groupe}%
%% v. 1003a
\begin{groupe}%
\personnageGrec{Μνησίλοχος}%
\trimIamb{χάλασον τὸν ἧλον.%\commentaire{}{}{\eff{Souda} Σχάσον}
\\\passeLigne}%
\end{groupe}%
%% v. 1003b
\begin{groupe}%
\personnageGrec{Τοξότης}%
\trimIamb{\ligne{1003}\decalage{χάλασον τὸν ἧλον.}ἀλλὰ ταῦτα δρᾶσ᾽ ἐγώ.\apparat{1003}{δρᾷς R : \lecon{δρᾶσ᾽} Voss}\\\passeLigne}%
\end{groupe}%
%% v. 1004
\begin{groupe}%
\personnageGrec{Μνησίλοχος}%
\trimIamb{οἴμοι κακοδαίμων, μᾶλλον ἐπικρούεις σύ γε.\apparat{1004}{ἐπικρούσεις R : \lecon{ἐπικρούεις} \Mon}\\\passeLigne}%
\end{groupe}%
%% v. 1005a
\begin{groupe}%
\personnageGrec{Τοξότης}%
\trimIamb{ἔτι μᾶλλο βοῦλις;\apparat{1005}{μᾶλλον R : \lecon{μᾶλλο} Bentley}\\\passeLigne}%
\end{groupe}%
%% v. 1005b-1006a
\begin{groupe}%
\personnageGrec{Μνησίλοχος}%
\trimIamb{\ligne{1005}\decalage{ἔτι μᾶλλο βοῦλις;}ἀτταταῖ ἰατταταῖ·\apparat{1005}{τατταταί R : \lecon{ἰατταταῖ} Faber}\\
κακῶς ἀπόλοιο.\\\passeLigne}%
\end{groupe}%
%% v. 1006b-1007
\begin{groupe}%
\personnageGrec{Τοξότης}%
\trimIamb{\ligne{1006}\decalage{κακῶς ἀπόλοιο.}σῖγα κακοδαίμων γέρον.\apparat{1006}{κακοδαίμω\expo{ν} γέρ\inRas{ω}{ο}ν \Rpc, M\expo{pc} \ligneApp{1007} ξεινίγκι R : \lecon{'ξινίγκι} Enger \espApp ἵνα R : \lecon{ἴνα} Brunck \espApp \lecon{πυλάξι} R : πυλάξῃ Σ}\\
πέρ᾽ ἐγὼ 'ξινίγκι πορμός, ἴνα πυλάξι σοι.%\apparat{1007}{ξεινίγκι R : \lecon{'ξινίγκι} Enger \espApp ἵνα R : \lecon{ἴνα} Brunck \espApp \lecon{πυλάξι} R : πυλάξῃ Σ}%
\commentaire{1007}{πέρ᾽ ἐγὼ 'ξινίγκι πορμός, ἴνα πυλάξι σοι}{$=$ φέρ᾽ ἐγὼ 'ξενέγκω φορμόν, ἵνα φυλάξω σε.}\\\passeLigne}%
\end{groupe}%
%% v. 1008-1014
\begin{groupe}%
\personnageGrec{Μνησίλοχος}%
\trimIamb{ταυτὶ τὰ βέλτιστ᾽ ἀπολέλαυκ᾽ Εὐριπίδου.\\
ἔα· θεοί, Ζεῦ σῶτερ, εἰσὶν ἐλπίδες.%\apparat{1009}{εἰσὶν R : εἴς ἔτ᾽ Naber (Faber ?) \cf 946}
\\
ἁνὴρ ἔοικεν οὐ προδώσειν, ἀλλά μοι\apparat{1010}{ἀνὴρ R : \lecon{ἁνὴρ} Bothe %\espApp \lecon{ἔοικεν} R : μ᾽ ἔοικε von Velsen 
\ligneApp{1011} ὑπεδήλωσεν R : \lecon{ὑπεδήλωσε} Biset}\\
σημεῖον ὑπεδήλωσε Περσεὺς ἐκδραμών,\\
ὅτι δεῖ με γίγνεσθ᾽ Ἀνδρομέδαν· πάντως δέ μοι\\
τὰ δέσμ᾽ ὑπάρχει. δῆλον οὖν τοῦτ᾽ ἔσθ᾽ ὅτι\apparat{1013}{οὖν R : οὖν νῦν Bizet : \lecon{οὖν τοῦτ᾽} Dobree}\\
ἥξει με σώσων· οὐ γὰρ ἂν παρέπτετο.%\apparat{1014}{σῶσον R : \lecon{σώσων} \Mon}
\\\passeLigne}%
\end{groupe}%
%% v. 1015-1055
\begin{groupe}%
%
\personnageGrec{Μνησίλοχος}%
\typeVers{bacch. $+$ ia.}\paratragiques{%\ledpb%
\ligne{1015}φίλαι παρθένοι φίλαι,\apparat{1015-21}{(Εὐρ.) R : \lecon{(Μνη.)} Tyrwhitt}%
\commentaire{1015-135}{%Hésychius + Synagoge : αἰδέσασθαι: τὸ ἔνοχον ὄντα [καὶ] φόνῳ ἀκουσίῳ καὶ πεφυκαδευμένον ἐφ᾽ ὡρισμένῳ χρόνῳ... τοῦ τετελευτηκότος ἐξιλᾶσθαι, ὡς εἰληφότος ἤδη τιμωρίαν
%αἰδῶσσα· les murs d'une cour Hsch
%ἴτας:ὅρκος.καὶ ὁ “Αιδης
Parodie de l'\italique{Andromède}.}{La pièce perdue d'Euripide \italique{Andromède} fut peut-être représentée la même année que l'\italique{Hélène}\marqueur{1}, en 412, un avant les \italique{Thesmophories}. L'édition des fragments d'Euripide de Jouan \& van Looy\marqueur{2} propose un dossier complet auquel le lecteur pourra se reporter. Le récit le plus circonstancié du mythe d'Andromède se trouve chez Apollodore (II, 4, 3) :%
\citComm{Παραγενόμενος δὲ εἰς Αἰθιοπίαν, ἧς ἐβασίλευε Κηφεύς, εὗρε τὴν τούτου θυγατέρα Ἀνδρομέδαν παρακειμένην βορὰν θαλασσίῳ κήτει. Κασσιέπεια γὰρ ἡ Κηφέως γυνὴ Νηρηίσιν ἤρισε περὶ κάλλους, καὶ πασῶν εἶναι κρείσσων ηὔχησεν· ὅθεν αἱ Νηρηίδες ἐμήνισαν, καὶ Ποσειδῶν αὐταῖς συνοργισθεὶς πλήμμυράν τε ἐπὶ τὴν χώραν ἔπεμψε καὶ κῆτος. Ἄμμωνος δὲ χρήσαντος τὴν ἀπαλλαγὴν τῆς συμφορᾶς, ἐὰν ἡ Κασσιεπείας θυγάτηρ Ἀνδρομέδα προτεθῇ τῷ κήτει βορά, τοῦτο ἀναγκασθεὶς ὁ Κηφεὺς ὑπὸ τῶν Αἰθιόπων ἔπραξε, καὶ προσέδησε τὴν θυγατέρα πέτρα. Ταύτην θεασάμενος ὁ Περσεὺς καὶ ἐρασθεὶς ἀναιρήσειν ὑπέσχετο Κηφεῖ τὸ κῆτος, εἰ μέλλει σωθεῖσαν αὐτὴν αὐτῷ δώσειν γυναῖκα. Ἐπὶ τούτοις γενομένων ὅρκων, ὑποστὰς τὸ κῆτος ἔκτεινε καὶ τὴν Ἀνδρομέδαν ἔλυσεν. Ἐπιβουλεύοντος δὲ αὐτῷ Φινέως, ὃς ἦν ἀδελφὸς τοῦ Κηφέως ἐγγεγυημένος πρῶτος τὴν Ἀνδρομέδαν, μαθὼν τὴν ἐπιβουλήν, τὴν Γοργόνα δείξας μετὰ τῶν συνεπιβουλευόντων αὐτὸν ἐλίθωσε παραχρῆμα.\\
Parvenu en Éthiopie, dont Céphée était roi, [Persée] trouva sa fille Andromède exposée pour être dévorée par un monstre marin. Cassiopée, épouse de Céphée, avait osé se comparer aux Néréides pour la beauté, et s'était même vantée de l'emporter sur elles. Les Néréides en furent irritées. Neptune partagea leur indignation, submergea le pays, et y envoya un monstre marin. L'oracle d'Ammon ayant annoncé que ces désastres cesseraient si on exposait Andromède, fille de Cassiopée, pour être dévorée par le monstre, les Éthiopiens forcèrent Céphée à faire ce que l'oracle ordonnait et à attacher sa fille à un rocher. Persée l'ayant vu, il devint amoureux, et promit à Céphée de tuer le monstre s'il voulait la lui donner en mariage. Céphée s'y étant engagé par serment, il attendit le monstre, le tua et délivra Andromède. Phinée, frère de Céphée, à qui Andromède avait été promise avant cet événement, conspira contre lui pour le faire périr ; mais Persée l'ayant découvert, lui montra la tête de la Gorgone, et le changea en pierre, ainsi que tous ceux qui avaient pris part à son complot.}%
}%
\debutNote{1015-135}%
\note{1}{C'est ce qu'indiquent les scholies au vers 1012, πιθανῶς· συνδεδίδακται γὰρ τῇ Ἑλένη, ainsi que celle du vers 53 des \italique{Grenouilles}. Le début du prologue d'\italique{Andromède} est repris au vers 1065, «~ὦ νὺξ ἱερὰ~», -- oui : Aristophane a chamboulé l'ordre de la pièce ; scholie à ce vers 1065 : τοῦ προλόγου Ἀνδρομέδας εἰσβολή --, laisse à penser qu'\italique{Andromède} fut jouée au lever du jour, avant donc \italique{Hélène}. Matthew Whright (\italique{Euripides' Escape-Tragedies}, 2005) fait l'hypothèse contraire et avance qu'\italique{Iphigénie en Tauride} concluait la trilogie d'Euripide de 412, agrémenté du drame satyrique du \italique{Cyclope}. Voir la recension que feu Alain Blanchard en fit dans la \italique{Revue des Études Grecques}, tome 119 (2006), pp. 798-800.}%
\note{2}{Tome VIII (fragments), première partie, de l'édition d'Euripide aux Belles Lettres (1998).}%
}
\end{groupe}%
%% v. 1015-1055
\begin{groupe}%
\typeVers{chor. $+$ cret.}\paratragiques{πῶς ἂν ἀπέλθοιμι καὶ\apparat{1016}{\lecon{ἀπέλθοιμι} R : ἐπέλθοιμι Ellebodius}% \espApp ἂν <οὖν> Hermann ut numeri dochmiaci fierent}
\commentaire{1015-55}{Monodie d'Andromède-Mnésiloque.}{La structure de ce chant a fait couler beaucoup d'encre... et nous échappe encore\marqueur{1}. Mais bien loin de l'anarchie décrite par les commentateurs, je suis tenté d'y voir une illustration sonore des échos : chaque phrase musicale semble être répétée : l'on peut ainsi regrouper par paires successives les différentes structures métriques (le thème puis son écho, de mètre identique ou très proche). On regroupe ainsi les vers : 1015-1071 \& 1018-1021 ; 1022-1023 \& 1024-1025 ; 1026-1029 \& 1030-1031 ; 1032-1033 \& 1034-1035 ; 1036-1037 \& 1038-1039a ; 1043 \& 1045 ; 1050-1051 \& 1052-1053.}%
\debutNote{1015-55}%
\note{1}{On pourra aller voir ce qu'en dit Lætitia, pp. 436-445.}%
}
\end{groupe}%
%% v. 1015-1055
\begin{groupe}%
\typeVers{$+$ ityph.}\paratragiques{τὸν Σκύθην λάθοιμι;\apparat{1017}{λάβοιμι R : \lecon{λάθοιμι} Ellebodius}}
\end{groupe}%
%% v. 1015-1055
\begin{groupe}%
\typeVers{$4$ bacch.}\paratragiques{\ligne{1019}κλύεις; ὦ προσᾴδουσ' ἀϋτὰς ἐν ἄντροις,\apparat{1019}{προσᾳδοῦσσαι τὰς R, Σ : \lecon{προσᾴδουσ'} Elmsley, \lecon{ἀϋτὰς} Burges (ἀϋταῖς Sommerstein) : alii alia}%
\commentaire{1018}{ὦ προσᾳδοῦσσαι τὰς ἐν ἄντροις.}{Comme dans la pièce d'Euripide, la nymphe Écho répond aux plaintes d'Andromède. (Dans la reprise comique, c'est Euripide qui se métamorphose en Écho.)
Willems a consacré sa \nieme{25} note (reproduite en annexe) à ce vers et aux diverses retouches qu'il a subies.
Je reprends la lecture de Burges\marqueur{1}, προσᾴδουσ' ἀϋτάς, satisfaisante quant au sens et qui offre l'avantage de rester très proche du texte du manuscrit.%\eff{αἱ Λιταί}
}%
\debutNote{1018}%
\note{1}{G. Burges, \italique{Classical Journal} 14 (1816) 217. Ou bien voir W. Mitsdörffer, «~Das Mnesichoslied in Aristophanes' \italique{Thesmophoriazusen}~», \italique{Philologus} 98 (1954-55), p. 69.}}
\end{groupe}%
%% v. 1015-1055
\begin{groupe}%
\typeVers{telesill.}\paratragiques{κατάνευσον, ἔασον ὡς}
\end{groupe}%
%% v. 1015-1055
\begin{groupe}%
\typeVers{$+$ ithyph.}\paratragiques{τὴν γυναῖκά μ᾽ ἐλθεῖν.}
\end{groupe}%
%% v. 1015-1055
\begin{groupe}%
\typeVers{dim. iamb.}\paratragiques{ἄνοικτος ὅς μ᾽ ἔδησε τὸν\apparat{1022}{(Μνη.) praefixit \Rpc, \Mon}}
\end{groupe}%
%% v. 1015-1055
\begin{groupe}%
\typeVers{$+$ lecyth.}\paratragiques{πολυπονώτατον βροτῶν·\apparat{1023}{\lecon{πολυπονώτατον} R, Σ : πολυστονώτατον Burges}}
\end{groupe}%
%% v. 1015-1055
\begin{groupe}%
\typeVers{dim. iamb.}\paratragiques{\ligne{1024}μόλις δὲ γραῖαν ἀποφυγὼν}
\end{groupe}%
%% v. 1015-1055
\begin{groupe}%
\typeVers{dim. iamb.}\paratragiques{σαπρὰν ἀπωλόμην ὅμως.}
\end{groupe}%
%% v. 1015-1055
\begin{groupe}%
\typeVers{do. $+$ ia. $+$ do.}\paratragiques{\ligne{1026}ὅδε γὰρ ὁ Σκύθης φύλαξ πάλαι ἐφέστηκ᾽· ὀλοὸν\apparat{1026}{ὁ Σκύθης R : del. Jackson}%
\commentaire{1026-8}{ὅδε γὰρ ὁ Σκύθης φύλαξ / πάλαι ἐφέστηκ᾽ ὀλοὸν / ἄφιλον ἐκρέμασε κόραξι δεῖπνον.}{Ce passage est sans doute légèrement corrompu. Les tentatives pour le corriger ont toutes été infructueuses\marqueur{1}. }%
\debutNote{1026-8}%
\note{1}{Cf. l'inventaire sans concession qu'en fait Lætitia \italique{The Songs of Aristophanes}, p. 443-444 : son regard infaillible de métricienne aguerrie détecte les invraisemblances métriques des «~corrections~». Sa proposition n'en est pas moins aporétique : elle voit dans ces vers une suite de trois dochmiaques se terminant par un dimètre iambique catalectique tout en concédant que la correption πάλα\lettreBreve{ι} (\latin{vocalis ante vocalem corripitur}) est suspecte : les seuls autres exemples chez Aristophane étant au vers 915, «~κύσ\lettreBreve{ω}. ἄπαγε~», déjà évoqué plus haut, et le peu convaincant «~ἰὼ ἰώ~» au vers 1170 des \italique{Nuées}.}%
}
\end{groupe}%
%% v. 1015-1055
\begin{groupe}%
\typeVers{$2$ cr. $+$ ba.}\paratragiques{\ligne{1029}ἄφιλον ἐκρέμασε κόραξι δεῖπνον.\apparat{1027}{\lecon{ἐφέστηκ᾽ ὁλοὸν} R : ἐφεστὼς ὀλοὸν Meineke : ἐφέστηκε κὠλοὸν Austin}% \ligneApp{1029} \lecon{κρέμασε} R : κρέμασεν Bothe : κρέμασε με Mehler}
}
\end{groupe}%
%% v. 1015-1055
\begin{groupe}%
\typeVers{do. $+$ ia. $+$ do.}\paratragiques{ὁρᾷς; οὐ χοροῖσιν οὐδ᾽ ὑφ᾽ ἡλίκων νεανίδων}
\end{groupe}%
%% v. 1015-1055
\begin{groupe}%
\typeVers{$2$ cr. $+$ ba.}\paratragiques{\ligne{1031}ψῆφον ἕστηκ᾽ ἔχουσα κημόν,\apparat{1031}{κημὸν ἕστηκ᾽ ἔχουσ᾽ R : \lecon{ἕστηκ᾽ ἔχουσα κημόν} Bibi, cf. Σ ad \italique{Vesp.} 674c}%
\commentaire{1031}{ψῆφον κημὸν ἕστηκ᾽ ἔχουσ᾽.}{On mettra ce vers en regard d'un fragment d'Hermippe :%
\citComm{καὶ πρὸς κύβους ἕστηκ' ἔχων <τὸ> κήθιον.}%
%ξιφισμὸς (σχῆμα ἕστηκ᾽ ἔχουσ᾽)
τὸ κήθιον et ὁ κημός sont synonymes : ils désignent dans l'univers ludique le \italique{cornet} dans lequel on secouait %(λαγαρίζεσθαι)\eff{Non !!} 
les dés\marqueur{1}. ἡ ψῆφος, comme λίθος, est dans ce contexte un \italique{pion}.
On constate avec ces deux exemples que `ἕστηκεν ἔχων' devait être l'expression consacrée pour le joueur qui s'apprêtait à lancer les dés.
On a soupçonné une altération du vers du fait de l'étrange succession initiale de trois longues. 
L'exemple d'Hermippe suscite naturellement l'émendation suivante :%
\citComm{ὁρᾷς; οὐ χοροῖσιν οὐδ᾽ ὑφ᾽ ἡλίκων ν\synizese{εα}νίδων\\
ψῆφον ἕστηκ᾽ ἔχουσα κημόν.}%
(κημός placé après ἔχω, les mots sont ainsi agencés dans le même ordre que dans le fragment\marqueur{2}.)\\
L'avantage est qu'on obtient (avec la synizèze de ν\synizese{εα}νίδων) le schéma métrique suivant :%
\citComm{$ \cup \ - \ - \ \cup \ - \ | \ \cup \ - \ \cup \ - \ | \ \cup \ - \ - \ \cup \ -$\\
$ - \ \cup \ - \ | \ - \ \cup \ - \ | \ \cup \ - \ \smile $\\
dochmiaque $+$ iambe $+$ dochmiaque \\
crétique $+$ crétique $+$ bacchée,}%
de scansion identique à celle des deux vers précédents, annulant du même coup le surprenant abrègement de πάλα\lettreBreve{ι} en seconde position d'une résolution (qui le faisait scander $ \widearc{\cup \cup}$, cf. commentaire v. 1026-8 n.1). 
Je ne sais ce que vaut cette émendation ; je laisse le soin d'en juger à des métriciens compétents.\\
Il est en tout cas bien évident qu'Andromède-Mnésiloque, se lamentant de son sort présent et regrettant les douces années de son enfance, ne fait aucunement allusion à une quelconque \italique{urne de vote} (κημός ?). Les chœurs de danse et les jeux enfantins, voilà qui suffit à sa mélancolie. Le marigot électoral ? alors que la jeune fille ne pouvait même pas y participer activement ? Mais où donc l'ensemble des commentateurs -- à commencer par le scholiaste -- est-il allé cherché cette abracadabrantesque explication ?
Qu'Aristophane ait pu utiliser le mot κημός, nom d'une sorte d'entonnoir posé sur l'urne pour en celer le vote, comme synecdoque pour l'\italique{urne} elle-même, cela est vrai mais avec une réserve de taille ayant échappé aux commentateurs : c'est pour permettre un jeu de mots, comme nous l'allons voir. (Ce qui explique qu'il soit le seul auteur recourant à ce vocable avec le sens d'\italique{urne}.)\\
% l'inadmissible traduction du terme κημός par "urne" [pour le vote]
Cette inadmissible traduction de κημός par «~urne~» [pour le vote] nous oblige à rebondir sur un passage des \italique{Guêpes}. %, étudié lors de l'un de vos cours. 
%Le proposez-vous encore à vos élèves ? auquel cas, j'ai quelques éclaircissement à vous partager -- et si vous ne l'étudiez plus, tout ce qui est gratuit est de toute façon plus beau !
Il s'agit du début des \italique{Guêpes} avec la traditionnelle scène des deux esclaves qui présentent la situation : en l'occurrence, le mal dont souffre leur maître Philocléon (et une bonne partie des Athéniens), pris qu'il est de \italique{philéliastite}. Et lorsqu'il voit une dédicace au beau Dêmos, le fils de Pyrilampe, notre malade ne peut s'empêcher d'écrire en marge (παρέγραψε) : κημὸς καλός. Vers 97-99 :
\citComm{καὶ νὴ Δί᾽ ἢν ἴδῃ γέ που γεγραμμένον\\
υἱὸν Πυριλάμπους ἐν θύρᾳ Δῆμον καλόν,\\
ἰὼν παρέγραψε πλησίον `κημὸς καλός.'}%
Qui était ce Dêmos Πυριλάμπους ? On possède quelques autres témoignages à son sujet : Calliclès, l'horrible démagogue du \italique{Gorgias} (481d), en avait fait son éromène. 
Platon souligne malicieusement que, à l'instar de l'autre dêmos que lorgne le démagogue, Calliclès finissait toujours par céder devant son cadet aimé (αἰσθάνομαι οὖν σου ἑκάστοτε, καίπερ ὄντος δεινοῦ, ὅτι ἂν φῇ σου τὰ παιδικὰ καὶ ὅπως ἂν φῇ ἔχειν, οὐ δυναμένου ἀντιλέγειν). 
Eupolis (fr. 227(213) Kassel-Austin) nous le décrit comme \italique{bouché} ; littéralement : il est plein de cérumen (κυψέλη) : 
\citComm{καὶ τῷ Πυριλάμπους ἆρα Δήμῳ κυψέλη ἔνεστιν; }%
Enfin, Athénée (IX, 397c-d), après avoir décrit le paon comme un animal bruyant qui réveille les dormeurs (τοιόνδε ταὧν, ὃς τοὺς εὕδοντας ἐγείρει) nous apprend que le fils de Pyrilampe en faisait élevage, dont il était d'ailleurs très fier...\\
Que peut-on en conclure ? assurément, que le beau Dêmos était un invétéré babillard, obtus à tout conseil, qui en remontrait à Calliclès au rebours d'une relation maître-élève classique et qui, pas plus que les paons qu'il élevait, ne s'arrêtait jamais de parler.
Comment lire maintenant le graffiti : %
\citComm{κημὸς καλός Δῆμον καλόν}% ? 
(noter l'accusatif de Δῆμον καλόν) ?
Tout simplement en conformité avec le sens premier de κημός, savoir : une \italique{muselière}.%
\citComm{Une bonne muselière [fera taire] le beau Dêmos.}%
Que la muselière (κημός), ait pu désigner aussi une \italique{urne de vote}, voilà qui devait laisser Aristophane bien rêveur !..
La pique contre la démocratie est évidente :
\citComm{Une bonne muselière [musellera] le bon peuple [pour cinq ans, pourrait-on ajouter à propos de la nôtre, encore plus décadente que celle d'Athènes].}%
Mais Philocléon dans sa naïveté -- qui serait touchante si elle n'était pas criminelle -- ne voit tout ce qu'il y a d'avilissant dans cette mascarade démocratique ; il est satisfait que le bon peuple puisse urner à son aise :
κημὸς καλός δῆμον καλόν.\\
Autre exemple, un peu plus loin dans la même pièce : Philocléon préside au procès du général Lachès\marqueur{3}, animalisé pour l'occasion en chien \italique{Chapardeur} (Λάβης). L'accusateur public n'est autre que le chien de Dêmos, Cléon. Bdélycléon fait entrer le prévenu et son accusateur ; c'est-à-dire les deux chiens, la canaille et le roquet. Réaction de Philocléon devant ces deux patibulaires (\italique{Guêpes}, vers 852-853): %
\citComm{Φιλοκλέων -- ἐς κόρακας, ὡς ἄχθομαι\\
ὁτιὴ 'πελαθόμην τοὺς καδίσκους ἐκφέρειν.\\
Aux corbeaux ! Quel ennui ! j'ai oublié d'apporter les καδίσκοι.}%
On traduit καδίσκοι par \italique{urnes} ; et il n'est pas douteux que dans la parodie de procès, les καδίσκοι soient mis en place d'\italique{urnes}... mais on voit que Philocléon frémit tout simplement devant les deux molosses qu'on lui présente et regrette l'absence de \italique{muselières}, c'est sans doute aussi le sens propre de καδίσκος.\\
%\eff{Qu'il me soit permis pour cela un petit détour par les autres comédies d'Aristophane. Commençons par le passage des \italique{Guêpes} déjà évoqué plus haut\marqueur{Voir le commentaire des vers 811-812, à propos de l'histoire de Miltiade qui détourna à son profit cinquante talents d'un coup, κατὰ πεντήκοντα τάλαντα.} dont c'est justement une scholie qui contient le fragment d'Hermippe cité au début. Ce passage est fort mécompris (vers 669-676) : Bdélycléon parle des démagogues qui se payent sur la bête athénienne -- le gros animal aurait dit Platon :%
%\citComm{κᾆθ᾽ οὗτοι μὲν δωροδοκοῦσιν κατὰ πεντήκοντα τάλαντα\\
%ἀπὸ τῶν πόλεων ἐπαπειλοῦντες τοιαυτὶ κἀναφοβοῦντες,\\
%`δώσετε τὸν φόρον, ἢ βροντήσας τὴν πόλιν ὑμῶν ἀνατρέψω.'\\
%σὺ δὲ τῆς ἀρχῆς ἀγαπᾷς τῆς σῆς τοὺς ἀργελόφους περιτρώγων.\\
%οἱ δὲ ξύμμαχοι, ὡς ᾔσθηνται τὸν μὲν σύρφακα τὸν ἄλλον\\
%ἐκ κηθαρίου λαγαριζόμενον (voir Taillardat § 130) καὶ τραγαλίζοντα τὸ μηδέν,\\ 
%σὲ μὲν ἡγοῦνται Κόννου ψῆφον, τούτοισι δὲ δωροφοροῦσιν\\
%ὕρχας, οἶνον, δάπιδας, τυρόν, μέλι, σήσαμα, προσκεφάλαια, κτλ.\\
%Eux [les démagogues] cependant extorquent aux cités des cinquante talents à la fois,\\
%en les menaçant, en les terrifiant de cette sorte :\\
%« Vous paierez le tribut, ou je tonne contre votre ville et la renverse. »\\
%Toi, quoique tu sois le chef, tu te contentes de grignoter les rognures.\\
%Et quant aux alliés, te voyant buvotant le reste de lie\\
%de la saucière et te nourrissant de rien,\\
%ils t'estiment autant que la voix de Konnos. Tandis qu'à eux ils offrent\\
%salaisons, vin, tapis, fromage, sésame, coussins, etc.}%
%Cette traduction appelle quelques explications. \\
%%Κόννου ψῆφον : par l'intermédiaire du scholiaste\marqueur{L'érudite scholie -- cela nous change de celles des \italique{Thesmophories} ! -- est de rédactions multiples et d'avis contradictoires, altérée par endroits et commentant des textes différents (que le lecteur aille consulter une édition avec apparat critique) ... aussi est-il délicat d'en piocher le bon pour en laisser le rebut. Au le lecteur qui voudra juger de mes choix, la voici telle que le manuscrit nous l'a transmise : \citComm{\eff{scholie}}}, Callistrate nous apprend l'existence d'un proverbe : Κόννου θρῖον (ἔστι δὲ παροιμία, ἐπὶ τῶν μηδὲν ἀξίων). τὸ θρῖον, \italique{la feuille de figuier}, et par extension le vain bruit qu'elle fait en brûlant, pouvait être permuté métaphoriquement avec ψόφος\marqueur{Κοννᾶς est décrit comme un vieillard divaguant au vers 536 des \italique{Cavaliers}. Cf. θρίων ψόφους et Taillardat §509, à propos de  vers 436 : ὡς ἐγὼ πολλῶν ἀκούσας οἶδα θρίων τὸν ψόφον. Mais il y a un jeu de mots non relevé : θρίων est aussi le génitif des θριαί, \italique{galets}, écho à la grève que cultive Philocléon pour n'être jamais à court de cailloux pour le vote.}. Si notre hypothèse est juste, le proverbe Κόννου ψόφος devait aussi exister, et c'est bien ce dernier que fait entendre Bdélycléon.\\
%τό κηθάριον, un hapax vu comme le diminutif de κηθίς, est au propre une \italique{saucière} pour le vinaigre\marqueur{Cf. la scholie afférente, de rédactions multiples et contradictoires, mais dont l'avis le plus pertinent semble être : ... ἣ κηθ<άρ>ια τὰ ὀξύβαφα, ἐξ ὧν τοὺς κύβους ἠ\{μ\}φίεσαν. πλεκτὰ δὲ ἦν ταῦτα. Ἄλλως. ἐκ κηθαρίου· τὰ ἐκ πέταλα τρύβλια... C'est cette interprétation qu'a retenue Hésychius : κήθεια· κηθάρια. τὰ ὀξύβαφα, ἐν οἷς τοὺς κύβους ἔβαλλον.} ; κηθίς, κήθιον ou κηθίδιον ayant pris par la suite le sens spécialisé de \italique{cornet à dés}.\\
%λαγαρίζεσθαι, c'est \italique{jouer des coudes}, peut-être \italique{gratter} ou \italique{renverser}\marqueur{Cf. Hésychius (λ 37) : λαγαριζόμενοι· σκαλεύοντες. δηλοῖ δὲ τὸ πρὸς τὰς λαγόνας [τὰ λαγαρὰ ?] τὸν ἀγκῶνα προσάγειν, πυκνὰ διατείνοντα τὴν χεῖρα. Photius (λ 15), λαγαριζόμενοι: σκαλεύοντες, τὰς λαγόνας τύπτοντες, est déjà plus vague : \italique{battre les flancs}... d'où le sens imagé de \italique{secouer} un récipient, c'est-à-dire lui remuer les flancs ? Voir la scholie au vers 674 \italique{Guêpes} ($=$ fr. 126(121) Phérécrate Kassel-Austin) qui précise : τὸ δὲ λαγαρίζεσθαι [λαγαρύζεσθαι cod.] ἕτερόν τι [autre que l'explication précédente] ἔοικεν ἐμφαίνειν. Φερεκράτης Μυρμηκανθρώποις. --- ἀντὶ τοῦ σιδαλεύην [σκαλεύειν Meineke] : ἄρα ποθ᾽ὑμεῖς καὶ τῆς ὀροφῆς τὸν χοῦν κατὰ τῆς κεφαλῆς καταμήσονται λαγαριζόμενοι. [fragment et donc d'interprétation peu sûre ; χοῦς, la \italique{conge} (unité de mesure) semble ici désigner un \italique{amas} de terre, comme la congère en français est un amas de neige.]}. Le scholiaste nous semble être dans le vrai lorsqu'il sous-entend que λαγαριζόμενον est ici un synonyme de ῥοφοῦντα\marqueur{ἐκ κηθαρίου: Τὰ ἐκ πέταλα τρύβλια, <ἃ> Ἐὐφρόνιος κήθια. μικρὸν οὖν ῥοφοῦντα ἐκ τρυβλίου καὶ τρωγαλίζοντα. Et encore : καὶ τὸ λαγαρίζειν [λαρυγίζειν cod.], τὸ ἀπο τρώγοντα θηλάζειν.}, \italique{sirotant}, \italique{buvotant}, \italique{sifflant} ; \italique{levant le coude}, dira-t-on.\\
%Bdélycléon reproche à son père un comportement de crève-la-faim : il grignote les rognures (τοὺς ἀργελόφους περιτρώγων) et infuse le reste des ordure (τὸν μὲν σύρφακα τὸν ἄλλον) dans du vinaigre qu'il sirote ensuite (ἐκ κηθαρίου λαγαριζόμενον).\\
%Pour comprendre la réaction qu'inspire une telle attitude, il faut avoir en tête une anecdote rapportée par Théophraste (Περὶ φυτικῶν ἱστορίων, IX, 17, 2-3) :%
%\citComm{Εὔδημος γοῦν ὁ φαρμακοπώλης εὐδοκιμῶν σφόδρα κατὰ τὴν τέχνην συνθέμενος μηδὲν πείσεσθαι πρὸ ἡλίου δύναντος κατέφαγε [sc. ἐλλέβορον] μέτριον πάνυ καὶ οὐ κατέσχεν οὐδ' ἐκράτησεν. ὁ δὲ Χῖος Εὔδημος\marqueur{} πίνων ἐλλέβορον οὐκ ἐκαθαίρετο. καί ποτε ἔφη πιεῖν ἐν μιᾷ ἡμέρᾳ δύο καὶ εἴκοσι πόσεις ἐν τῇ ἀγορᾷ καθήμενος ἐπὶ τῶν σκευῶν καὶ οὐκ ἐξαναστῆναι πρὸ τοῦ δείλην γενέσθαι· τότε δ' ἐλθὼν καὶ λούσασθαι καὶ δειπνεῖν ὥσπερ εἰώθει καὶ οὐκ ἐξεμέσαι· πλὴν οὗτός γε βοήθειάν τινα παρασκευασάμενος κατέσχε· \textbf{κίσσηριν} γὰρ ἐπιπάττων ἐπ' ὄξος δριμὺ πιεῖν ἔφη μετὰ τὴν ἑβδόμην πόσιν, καὶ πάλιν ὕστερον ἐν οἴνῳ τὸν αὐτὸν τρόπον.\\
%Eudême, le vendeur de drogue, qui avait une haute réputation dans son commerce, après avoir fait le pari qu'il n'éprouverait aucun effet avant le coucher du soleil, but une dose assez modérée [d'hellébore, plante purgative, cf. La Fontaine et le proverbe "se purger avec quatre grains d'hellébore"], et elle s'avéra trop forte pour son pouvoir de résistance. Mais une fois le Chiote Eudême\marqueur{La critique a vu en ce `ὁ δὲ Χῖος Εὔδημος' la mention d'un second Eudême, homonyme au premier, le φαρμακοπώλης, et qui se serait lui-aussi lancé dans des expériences pour contrer les remontées helléboriennes... fameuses coïncidences !.. il est beaucoup plus simple de considérer que c'est bien le même Eudême, mais qu'après son premier échec, Aristote se soit plu à préciser son origine, Χῖος, qui désigne l'\italique{as} du dé, c'est-à-dire le coup perdant. Apollonius le paradoxographe ne s'y était pas trompé, lui qui a compilé cette anecdote dans ses ἱστορίαι θαυμάσιαι, et qui qualifie Eudême (le manuscrit porte fautivement Εὔνυμος) de Chios de φαρμακοπώλης. Cf. l'édition de O. Keller, \italique{Rerum Naturalium Scriptores, Apollonii} Histor. mir. \italique{L}, p. 53.} a pris un bouillon d'hellébore et n'a pas été purgé. À une occasion, il a dit avoir bu en une seule journée vingt-deux potions [d'hellébore] sur la place du marché alors qu'il était assis à son étal, et n'a pas quitté l'endroit jusqu'au soir, puis il est rentré chez lui, a pris un bain et dîné comme à son habitude et n'a pas rendu. Cependant cet homme a pu tenir parce qu'il s'était muni d'un antidote ; car il dit qu'après la septième potion, il prit une gorgée de vinaigre acidulé saupoudrée de \italique{pierre ponce}, et plus tard en prit une gorgée dans du vin de la même manière.}%
%Cet Eudême de Chios, φαρμακοπώλης de son état, est mentionné ailleurs par Aristophane (\italique{Ploutos}, vers 884). L'anecdote contée par Théophraste devait être fameuse pour qu'à son époque encore il en ait connaissance. Les vertus de la pierre ponce, κίσηρις, étaient déjà connues des Anciens\marqueur{À la suite du passage cité, Théophraste parle du fait que quelque grains de pierre ponce permettent de réguler l'ébullition d'un liquide. Cette propriété est encore enseignée et utilisée dans nos lycées au laboratoires de chimie.}. Il n'est à mon avis pas fortuit qu'Aristophane juxtapose dans un de ses fragments\marqueur{fr. 332(320) Kassel-Austin $=$ Pollux, VII, 95.} κίσηρις et στρόφος, dont l'un des multiples sens est la \italique{colique}.\\
%Avec l'image que développe Bdélycléon, voyant son père tremper étrangement les rognures les plus immangeables, semblables à de petits cailloux, dans une solution vinaigrée, puis de siroter cette "potion", Philocléon devait évoquer celle bien connue du fantasque Eudême et de ses décoctions à la pierre ponce. Il nous paraît vraisemblable qu'au lieu de la leçon des manuscrits, Κόννου ψῆφον, Aristophane avait écrit Χίων ψῆφον, cf. Paix 171 et Plutarque 232f
%
%Le scholiaste a pressenti qu'Aristophane filait la métaphore avec la cubomanie : que dit-il en substance ? que Philocléon est un \italique{petit joueur}, lui qui se contente du triobole.
%(Ajoutez à cela un qu'un autre hapax -- décidément ! --, "καὶ τραγαλίζοντα τὸ μηδέν", est peut-être bien mis là pour faire entendre "κἀστραγαλίζοντα τὸ μηδέν", "jouer aux astragales pour rien"\marqueur{κυβεύειν ne désigne pas le seul fait de jouer aux dés puisqu'on le trouve employé pour des jeux ne nécessitant pas de dés. Stephen Kidd propose de voir derrière κυβεία un terme générique pour tous les jeux d'argent quand ἀστραγαλίζειν signifierait \italique{jouer pour du beurre}, cf. son article : "How to Gamble in Greek : the Meaning of Kubeia", \italique{Journal of Hellenic Studies} 137 (2017), p. 119-134.}.)\\
%C'est en ce sens qu'il faut comprendre la fin de la tirade de Bdélycléon et la réplique de Philocléon (vers 679-680) :%
%\citComm{οὐδεὶς οὐδὲ σκορόδου κεφαλὴν τοῖς ἑψητοῖσι δίδωσιν.\\
%Φιλοκλέων --- μὰ Δί᾽ ἀλλὰ παρ᾽ Εὐχαρίδου καὐτὸς τρεῖς γ᾽ ἄγλιθας μετέπεμψα.\\
%Pas un ne te donne ne fût-ce qu'une tête d'ail pour assaisonner ta poissonnaille.\\
%Philocléon --- Mais non par Zeus : }}%
%
%
%
%σφυρά: τὰ περὶ τοὺς ἀστραγάλους, αἱ περιφέρειαι τῶν ποδῶν ἢ τῶν ὀρῶν τὰ ἐπίπεδα καὶ κατώτατα ἢ ἡ παρά τισι λεγομένη κέστρα, ἰχθὺς θαλάσσιος
%σάγλη : schlechter Wurf Charis : le récipient
%se nourrissant de rien en  creuse, ils font état de toi comme du suffrage de Connos (2); tandis qu'à eux ils envoient en présent jarrons de salaison, vin, lapis, fromage, miel, sésame, coussins, coupes, manteaux, couronnes, colliers, vases à boire, richesse et santé. Et de ceux à qui tu commandes, « après tant de labeurs sur la terre et sur l'onde », pas un no te donnera pour la poissonnaille une simple tète d'ail.%
%
%
%fr. Ar. 556(543)
%\marqueur{On se trouve bien en peine de déterminer son mètre ; il y a cet énigmatique `παρὰ τὴν ξιφοθήκην' en marge du manuscrit de Ravenne. Est-ce une simple déformation de `παρὰ τὴν ψηφοθήκην' glosant κημόν ?  Ou faut-il en conclure qu'un ξιφο- se trouvait dans le texte originel des \italique{Thesmophories} ?}, 
%Κυβευτοῦ σκεύη · ἄβαξ ἀβάκιον, κόσκινον, κύβοι, διάσειστοι κύβοι, ἀστράγαλοι, φιμοί, κημοί, τηλία, κήθια κηθίδια, ψῆφοι, πεττοί.
%.
%
%Objets du joueur de dés : plateau, petit plateau, crible, dés, dés à secouer, osselets, cornets à dés, entonnoirs couvrant l’urne, table à rebord, cornets à secouer les dés, petits cornets à secouer les dés, pions, jetons.
%(Les cornets désignent aussi des muselières.)
%Il nous
% Voir dans ce vers une allusion aux activités philéliastes avec les diversLes commentateurs, à commencer par le scholiaste
%
%
%
%κήτειον:μέγαν. [καὶ (εἰς) ὃ τὰς ψήφους διωθοῦσιν ἐν τοῖς κληρωτηρίοις].καὶ ἐμετήριονἀπὸ λαχάνου ἀγρίου,ὅπερ ἀντὶ πτερῶνκαθίεσανεἰςτὸ στόμα,ὡςΚρατῖνοςἐν “Ὥραις(΄.292 Κ.-Α.) σημαίνει,καὶ ἥδυσμα. καὶ τὸ νεαρὸν χρόμμυον. [ἔνιοι μέλοςτι] οἱ δὲτὰὈψαρτυτικὰ συνθέντεςδιὰτοῦγ φασὶγήτιο (Hésychius)
%		Voir Lamer Realencyclopedia : 
%		κηθίδιον, Poll. 10.150.—Wrongly written κητίον, Alciphr. 1.39, Ath. 11.477d
%φράζε τοίνυν: ὡς ἐγώ σοι πᾶς ἀνέρριμμαι κύβος Aristophane fr. 929(673) où pas a le sens de décisif comme dans panta lithon kinein)
%
%Grenouilles 968-970 :
%Διόνυσος --- Θηραμένης; σοφός γ᾽ ἀνὴρ καὶ δεινὸς ἐς τὰ πάντα,
%ὃς ἢν κακοῖς που περιπέσῃ καὶ πλησίον παραστῇ,
%970πέπτωκεν ἔξω τῶν κακῶν, οὐ Χῖος ἀλλὰ Κεῖος.
%
%
%proskeimai : être face ontre terre? cf. Oiseaux Χίοισιν ἥσθην πανταχοῦ προσκειμένοις. + Souda Κῷος πρὸς Χῖον
%ἀστράγαλος = σφόνδυλος (vertèbres et non kondulois)
%Hésychius :  
%ἀγάς· ἡ πτῶσις τοῦ ἀστραγάλου
%ἀστραγάλας· τὰς πονηρὰς κύνας. κακοήθεις κύνες
%βουβάλιες: οἱ βούβαλοι(Ηdt.,4,192,1).ἢ ἀστράγαλοι
%γλάματα: ἀστράγαλοι
%γραῦς:τὸ συναγόμενον καὶ ἐπιπηγνύμενον ἐπάνω, ὅταν γάλα ἕψηται, ἤ τι τοιοῦτον ὑγρόν (Αρ.ΡΙυτ. 1206).καὶ ἡ ἐν τοῖς χείλεσι τῶν ποτηρίων γραμμή. καὶ βόλος τις ἀστραγάλων
%δικέντρων: βόλος τις ἀστραγάλων οὕτως ἐκαλεῖτο
%δορχελοί:ἀστράγαλοι
%Crux : ἴδα· τὸ εἰκὴν ἐν τῷ ἀστρωγαλίζειν (keinon : le 1 = Xhios ?)
%κίλλαι:ἀστράγαλοι.ἢ ὄνοι
%κόνδοι:κεραῖαι.ἀστράγαλο
%κυνοῦραι: ἀστράγαλοι
%κῳάζειν:ἀστραγαλίζειν.ἐνεχυράζειν
%κῷοι ἀστράγαλοι
%Κῶος Χῖον: ὁ Κῶος ἀστράγαλος, ὁ ἕξ. ὁ μὲν γὰρ Χῖος ἐδύνατο ἕν ὁ δὲ Κῶος ἕξ
%Crux λίγδοι:οἱ προστετριμμένοι τῶν ἀστραγάλων
%λίσποι· οἱ ἐκτετριμμένοι τῶν ἀστραγάλων. καὶ ᾿Αθηναῖοι: ἀπὸ Θησέως, ἀπὸ τοῦ ἐν Ἅιδου προσσχεθῆναι τῇ πέτρᾳ ἀπὸ τῶν γλουτῶν
%λωγάλιοι: ἀστράγαλοι.ἢ πόρνοι
%Crux μελαιαί:ἀστράγαλοι.ἢ νωθροί
%-> μήλιαι: ἀστράγαλοι (kata- ?)
%πέλλυτρα: οἱ δεσμοὶ οἱ περὶ τὰ[ς] σφυρὰ (spheiras ?) καὶ τοὺς ἀστραγάλους τῶν δρομέων περιελισσόμενοι εἰς τὸ μὴ ἐκ<σ>τραφῆναι
%πλειστοβολεῖν: παίζειν ἀστραγάλοις
%στερεαί:τινὲς τῶν δικαστικῶν ψήφων [ἄλλοι δὲ ἀστέρας astragales ?]
%σφυρά: τὰ περὶ τοὺς ἀστραγάλους, αἱ περιφέρειαι τῶν ποδῶν ἢ τῶν ὀρῶν τὰ ἐπίπεδα καὶ κατώτατα ἢ ἡ παρά τισι λεγομένη κέστρα, ἰχθὺς θαλάσσιος
%σφυρόν:[σφοδρόν.ἢ]τὸ ἐπάνω τῶν ἀστραγάλωn
%ταλαρίς: ἐξ οὗ βάλλουσι τοὺς ἀστραγάλους
%τρόπα: εἶδος παιδιᾶς, καθ᾽ἣν στρέφουσι τοὺς ἀστραγάλους εἰς τὸ ἕτερον μέρος (Krati fr 180 Κ.Ἁ.)
%ὑπόκυκλα: τοὺς ἀστραγάλους τοὺς ὑποτιθεμένους τῷ πυθμένι τῶν τριπόδων (Σ 375 ν. 1.)
%La femme qui découpe son mari ? ἀγερσικύβηλις: Κρατῖνος ἐν Δραπέτισιν (. 66 Κ.Ἁ.) ἐπὶ  Λάμπωνος. τὸν αὐτὸν ἀγύρτην καὶ κυβηλιστὴν εἶπεν, οἱονεὶ θύτην καὶ μάντιν. Κύβηλιν γὰρ ἔλεγοντὸν πέλεκυν ... ὅθεν καὶ Λύσιππος ἐν Βάκχαις (.6 Κ.Ἁ.) τὸν αὐτὸν (ὡς) ἀγύρτην κωμωδεῖ [σικυβη ?] cf. suivant
%Sur surphetwdes mendiant : ἀγύρτης: ὀχλαγωγός ,προσαίτης, ἐπαίτης, συρφετώδης, ἰδιώτης(Εὐρ. Ρες. 503). ἔστι δὲ καὶ ὁ ἀγείρων ὄχλον. καὶ ὄνομα βόλου κυβευτικοῦ (Ευβ  .57,5 Κ.Ἁ.)
%-> À PROPOS DE AÏX : αἴγλη: λαμπηδών , αὐγή ,φῶς, λαμπρότης  ... ἔστι δὲ καὶ βόλος φαῦλος κυβευτικός (Κρατ. 405 Κ Α.)
%--> κριοῦ διακονία: ὅταν προστάσσωμεν παιδίῳ διακονῆσαι, λέγομεν· δώσω σοι εἰ τύχοι ἀστραγάλους, ἢ ἰσχάδας. καὶ ὁ κριὸς οὖν εἶπε· καταλείψω τοὺς ἀστραγάλους σοι [δώσω]
%--> Konnos : Βερενίκης [Βερονίκη: πόλις] πλόκαμος: τοῦτον κατηστερίσθαι φησὶ Κόνων. καὶ βόλος δέ τις ἀστραγάλων οὕτω καλεῖται
%᾿Αργεῖος: κυβευτικοῦ βόλου ὄνομα(Ευbul.57,3Κ.-Α.)
%--> λαγαριζόμενοι: σκαλεύοντες. δηλοῖ δὲ τὸ πρὸς τὰς λαγόνας (creux de chaque côté) τὸν ἀγκῶνα (courbure => coude) προσάγειν,πυκνὰ διατείνοντα τὴν χεῖρα(Ρher.126,4Κ.-Α.)
%Photius : λαγαριζόμενοι: σκαλεύοντες, τὰς λαγόνας (les flancs) τύπτοντες.
%(Pollux II 142 : καὶ τὸ παίειν τῷ ἀγκῶνι κυβιτίζειν λέγει)
%γυλλός:κύβος,ἢ τετράγωνος λίθος
%ἐκκεκομμένος: ὁ διὰ τοῦ κυβεύειν τὰ αὑτοῦ ἀπολέσας, ὅπερ νῦν ἐκκεκοπισμένος φαμέν. ἢ ὁ ἀποκεχωρισμένος
%“ἢ τρὶς ἕξ, ἢ τρεῖς κύβοι᾽ παροιμία ἐπὶ τῶν ἡ ἐπιτυγχανόντων ἢ ἀποτυγχανόντων. πάλαι γὰρ τρισὶν ἐχρῶντο πρὸς τὰς παιδιὰς κύβοις, καὶ οὐχ, ὡς οἱ νῦν, δύο. ἔστι δὲ ὁμωνυμία
%κήθεια: κηθάρια.τὰ ὀξύβαφα, ἐν οἷς τοὺς κύβους ἔβαλλον
%κίγκασος:κυβευτικός τις βόλος
%κυβεῦσαι: παῖξαι.ταβλίσαι, χλευάσαι
%κυβευτήν: σοφὸν κυβευτήν, ἀλλὰ μὴ στένειντύχην' (Σοφ. 947) τὸν τοῖς κύβοις χρώμενον. ἐπιστήμην γάρ τινα ἐδόκουν τὸ [δὲ] δεξιῶς χρῆσθαι. ἐπιστημονικωτέραν τὴν πεττείαν τῆς κυβείας ἐνόμιζον
%Crux : (κυ?)βῆναις: γλαῦξαις
%κύων: ὁ ἐλαυνομένου τοῦ σιδήρου τοῦ ἀργοῦ ἐξ αλλόμενος σπινθήρ. οἱ δὲ τὴν Ἐρινῦν. οἱ δὲ τὴν νόσον τὴν οὕτως λεγομένην κύνα. δηλοῖ δὲ καὶ τὸ ἀνδρεῖον μόριον. καὶ τὸ ὑλακτοῦν ζῶον. καὶ τὸν ἀναιδῆ. καὶ τὸ ἄστρον.καὶ τὸ θαλάσσιον ζῶον. καὶ ὑπὸ τῶν κυβευτῶν βόλον τινά, ὃς Χῖος καλεῖται
%-> Pour Mnésiloque/andromède : κύων ἐπὶ δεσμά: παροιμία ἐπὶ τῶν ἑαυτοὺς εἰς κολάσεις ἐπιδιδόντων. λέγεται δὲ καὶ βοῦς ἐπὶ δεσμά
%κύων: παρ᾽ ἐντέροισι παροιμία ἐπὶ τῶν μὴ δυναμένων ἀπολαύειν τῶν παρακειμένων
%-> μάγνης: κυβευτικοῦ βόλου προσηγορία καὶ ἄλλων πολλῶν λεγόμενα ὀνόματα βόλων
%Μάνης: κυβευτικοῦ βόλου ὄνομα [καὶ] βαρβαρικόν,καθάπερ Εὔβουλος ἐν Κυβευταῖς(fr 59 Κ.-Α.): παρίστησι καὶ ἄλλων κατα λεγομένων βόλων <ὀνόματα>
%μίδας: κυβευτικοῦ βόλου ὄνομα. καὶ θηρίδιόν τι διεσθίον τοὺς κυάμους. καὶ ὁ πλούσιος: ἀπὸ Μίδα, τοῦ βασιλέως
%οἴνη:ἡ ἄμπελος .καὶ κυβευτικὸς λέγεται βόλος,ὁ κενὸς καὶ ἀντικείμενος τῷ ἕξ, ὃ καὶ Χῖος. οἰνίζειν γὰρ τὸ μονάζειν κατὰ γλῶσσαν
%-> Πέρσης: κυβευτικοῦ βόλου ὄνομα
%-> qui donne raison à Taillardat sur la diff petteia/kubeia (cf. aussi Aristote) : πεσσὰ πεντέ γραμμα[τα]καὶ κύβων βολαί: Σοφοκλῆς Ναυπλίῳ Πυρκαεῖ (φρ. 429 Ἀ.). παρόσον πέντε γραμμαῖς ἔπαιζον. διαφέρει δὲ πεττεία κυβείας. ἐν ἣ μὲν γὰρ τοὺς κύβους ἀναρρίπτουσιν: ἐν δὲ τῇ πεττείᾳ αὐτὸ μόνον τὰς ψήφους μετακινοῦσι
%Mais : πεττεία: διὰ κύβων παιδιά(Ρlat. Leg.820d)
%σκειράφιον: τὸ κυβευτήριον ,ἴσως διὰ τὸ ἐν Σκ[ε]ίρῳ τὴν διατριβὴν ἔχειν
%σχαδών: κυβευτικὸς βόλος
%τρὶς ἕξ:νικητήριος βόλος(Αesch. Αg. 33). καὶ παροιμία: τρὶς ἕξ ἢ τρεῖς κύβοι. ἐλέγετο δὲ ἐπὶ τοῦ εὐστόχου, ἀπὸ τῶν κυβευτῶν (Pherecr.129Κ.- Α.)
%-> Sur mesos : τροχαῖα: μέσα ἐν κύβοις. ἢ ὁδός, ὡς Ῥίνθων (fr 23 Κ.-Α.)
%χαλκίνδα: τὸ εἰς χαλκὸν κυβεύειν
%ἀναθέσθαι: ὁμολογῆσαι ἐπὶ τῶν πεσσῶν ἐλέγετο: οἱ γὰρ παίζοντες καὶ τὰς ψήφους διώρθουν
%ἀρτοπόπος: ὁ πέσσων ἐν ἐργαστηρίῳ
%Crux ἄρτος:βόλος τις
%-> sur surfax :  παλινδορία: τὸ σκῦτος. Πλάτων Σύρφαχι (φρ. 180Κ.Ἀ.):  'σὲ μὲν, ὦ μοχθηρέ, παλινδορίᾳ παίσας αὐτοῦ καταθήσω᾽
%ET : ἐνκατασυρφάζειν: ἐνυβρίζειν
%ἐς βόθυνον: εἶδος παιδιᾶς παρὰ Κρατίνῳ (fr. 219 Κ.-Α.)
%-> έκδότω δέ τις 
%καί ψηφολόγιον ώδε καί δίφρω δύο (Ar. fr. ...)
%Diphilus 77 K.-A. : ἕλκ᾽ἐς μέσον τὸν φιμὸν ὡς ἂν ἐμβάλῆ
%Photius : Τρὶς ἓξ ἢ τρεῖς κύβοι: οἱ μὲν τρὶς ἓξ νίκην· οἱ δὲ τρεῖς κύβοι κενοί· καὶ τὸ παρ’ Εὐριπίδηι τοιοῦτον·   (10) Βέβληκ’ Ἀχιλλεὺς δύο κύβω καὶ τέτταρα· τριῶν γὰρ ὄντων τῶν ἀναρριπτουμένων βόλων, δύο μὲν κενοὺς αὐτὸν φησὶ βάλλειν· ἕνα δὲ τὸν τέταρτον· διττὸς οὖν ὁ κύβος ὅ, τε ἀναρριπτούμενος, εἴτε κενὸς, εἴτε πλήρης εἴη· καὶ ἰδίως ὁ κενός· ὅτι δὲ ὁ    (15) τρὶς ἓξ βαλὼν κατώρθου, καὶ Αἰσχύλος ἐν Ἀγα- μέμνονι παρίστησι· Τὰ δεσποτῶν γὰρ εὖ πεσόντα θήσομαι, Τρὶς ἓξ βαλούσης τῆς ἐμῆς φρυκ- τωρίας.
%Λάκωνες: κυβευτικοῦ βόλου ὄνομα. Phot
%Ploutos 816 artiasmos (pair ou impair)
%σκιραφεῖα ̓ τὰ κυβεῖα (les gains ?) Photius
%σκιρζάφιτζα ̓ τὰ κυβευτήριά φησιν ̓ἐπειδὴ διέτριβον ἐν Σκίρῳοἱ κυβεύοντες, ὡς Θεόπομπος(FrgHist 115F228)
%Photius : ἐκκόψαι· ἀντὶ τοῦ νικῆσαι. οὕτως Ἄλεξις
%
%Σ Vesp. 674 (il y a un unique psêfoforwn !) avec cette citation : καὶ πρὸς κύβους ἕστηκ' ἔχων τὸ κήθιον. (Hermippe, fr. 27 K.-A.) :
%ἐκ κηθαρίου: Σύρφακα. τὸ ἱκανὸν πλῆθος τῶν δικαστῶν τὸ συρφετῶδες, ἐξ εὐτελῶν τρεφόμενον. 
%κηθάριον γὰρ πλέγμα ἐστὶ κανισκῶδες [ἐστὶ κοσκινῶδες Dobrée], ἐπιτιθέμενον τῇ κληρωτ<ρ>ίδι τῶν ψήφων. λαγαριζόμενον δὲ τὰ λάγαρα ἐσθίοντα, ὅ ἐστιν εὔθρα<υ>στα καὶ εὐτελῆ τινα. ἣ κηθ<άρ>ια τὰ ὀξύβαφα, ἐξ ὧν τοὺς κύβους ἠ\{μ\}φίεσαν. πλεκτὰ δὲ ἦν ταῦτα. Ἄλλως. ἐκ κηθαρίου: Τὰ ἐκ πέταλα τρύβλια, <ἃ> Ἐὐφρόνιος κήθια. μικρὸν οὖν ῥοφοῦντα ἐκ τρυβλίου καὶ τρωγαλίζοντα, οἷον κυάμους καὶ ὅσα τοιαῦτα. οὐκ ἔστι δὲ, ἀλλ᾽ἀγγεῖον πλεκτὸν, εἰς ὃ τὰς ψήφους καθιᾶσι, τὸ κήθιον, οἷον ὡς ἀπὸ τοῦ ψηφηφορεῖν μόνον. καὶ Χαῖρις δὲ τὸ κήθιον τρυβλίου τι εἶδος. καὶ τὸ λαγαρίζειν [λαρυγίζειν V.], τὸ ἀπο τρώγοντα θηλάζειν. μήποτε μέντοι καὶ πρὸς τὴν τῶν κύβων παιδιὰν ὁμοίως ἐχρῶντό τινι βάλλοντες ἐξ αὐτοῦ. καὶ γὰρ Ἕρμιππος ἐν Θεοῖς 
%καὶ πρὸς κύβους ἕστηκ' ἔχων <τὸ> κήθιον.
%τὸ δὲ λαγαρίζεσθαι [λαγαρύζεσθαι V. Ald.] ἕτερόν τι ἔοικεν ἐμφαίνειν.
%Φερεκράτης Μυρμηκανθρώποις. --- ἀντὶ τοῦ σιδαλεύειν [ σκαλεύειν ;] ἄρα ποθ᾽ὑμεῖς καὶ τῆς ὀροφῆς τὸν χοῦν κατὰ τῆς κεφαλῆς καταμήσονται λαγαριζόμενοι. Υ. --- ἴσως ἂν οὖν λέγοι ἦτοι ἐκ τῶν ἀγγείων τῆς ψήφου ἀποξύοντα τὸν μισθὸν καὶ ἐξ ἄγγους τινὸς πλεχτοῦ, οἷον ὁ κάρταλος, τρωγαλιζόμενον, ὡς ἂν πένητα. λαγαριζόμενον δὲ ἀντὶ τοῦ ἀποξύοντα, ὡς ἀπὸ τῶν δικαστικῶν ἀγγείων. ταῦτα δὲ ὡς περὶ πενήτων.
%Surfax : 
%
%
%
%
%(1) Le sens de cette locution a été contesté. Il me semble pleinement confirmé par le vers 562 du Ploutos. 
%(2) 11 s'agit de ce pauvre diable de Connos, l'ancien maître de musique de Socrate, dont il a été question dans les Cavaliers (v. 534) et qui certes ne s'inquiétait guère de voter. (Peut-etre allusion à Koos aux dés ?)
%
%
%SUR KÊMOS : 
%
%
%Acharniens :
%Χορός
%ἀπολεῖς ἄρ᾽ ὁμήλικα τόνδε φιλανθρακέα;
%Δικαιόπολις
%οὐδ᾽ ἐμοῦ λέγοντος ὑμεῖς ἀρτίως ἠκούσατε.
%Χορός
%ἀλλὰ γὰρ νῦν λέγ᾽, εἴ σοι δοκεῖ, τόν τε Λακεδαιμόνιον
%† αὐτὸν ὅτι τῷ † τρόπῳ σοὐστὶ φίλος:
%340ὡς τόδε τὸ λαρκίδιον οὐ προδώσω ποτέ.
%Δικαιόπολις
%τοὺς λίθους νύν μοι χαμᾶζε πρῶτον ἐξεράσατε.
%Χορός
%οὑτοιί σοι χαμαί, καὶ σὺ κατάθου πάλιν τὸ ξίφος.
%Δικαιόπολις
%ἀλλ᾽ ὅπως μὴ ν᾽ τοῖς τρίβωσιν ἐγκάθηνταί που λίθοι.
%Χορός
%ἐκσέσεισται χαμᾶζ᾽: οὐχ ὁρᾷς σειόμενον;
%345ἀλλὰ μή μοι πρόφασιν, ἀλλὰ κατάθου τὸ βέλος.
%ὡς ὅδε γε σειστὸς ἅμα τῇ στροφῇ γίγνεται.
%
%Equites : 
%
%ἐν τῇ πυκνί, κᾆθ’ ὅταν
%μή σοι τύχῃ ὄψον ὄν, @1
%τούτων ὃς ἂν ᾖ παχύς,
%θύσας ἐπιδειπνεῖς.   (1140)
%ΔΗ.                       Σκέψασθε δέ μ’, εἰ σοφῶς
%αὐτοὺς περιέρχομαι
%τοὺς οἰομένους φρονεῖν
%κἄμ’ ἐξαπατύλλειν.
%Τηρῶ γὰρ ἑκάστοτ’ αὐ-   (1145)
%τοὺς οὐδὲ δοκῶν ὁρᾶν
%κλέπτοντας· ἔπειτ’ ἀναγ-
%κάζω πάλιν ἐξεμεῖν
%ἅττ’ ἂν κεκλόφωσί μου, (keklotosi : secouer [le cornet à dés] ?)
%κημὸν καταμηλῶν.    (1150)  (mêlos : la luette (pour se faire vomir) ? kêmmmon ici la muselière ? pour Willems, elle désigne l'enquête judiciaire + jeu de mots : cf. Σ guêpes 674 : τὸν χοῦν κατὰ τῆς κεφαλῆς καταμήσονται λαγαριζόμενοι :  en renversant l'urne ?)
%Ar. fr. 625(614) K.-A. : τὴν φάρυγα μηλῶν δύο δραχμὰς ἕξει μόνας
%Phot. = Suid. μ 940 μηλῶσαι· τὸ καθεῖναί τι εἰς βάθος. καὶ τὴν φάρυγγα μηλῶσαι,τὸ διαχρῖςαι (-ῆσαι Phot.) τῷ δακτύλῳ.᾿Αριsτοφάνης· τὴν ----- μόνας
%
%
%ΠΑ.                 Ἄπαγ’ ἐς μακαρίαν ἐκποδών.    (1150)
%ΑΛ.                                            Σύ γ’, ὦ φθόρε.    (1151)
%ΠΑ.                 Ὦ Δῆμ’, ἐγὼ μέντοι παρεσκευασμένος
%τρίπαλαι κάθημαι βουλόμενός σ’ εὐεργετεῖν.
%ΑΛ.                 Ἐγὼ δὲ δεκάπαλαι γε καὶ δωδεκάπαλαι
%
%
%Vespae : 
%
%C'est aussi avec le double sens de \italique{muselière} qu'il faut prendre κημός au vers 99 des \italique{Guêpes} : Dêmos le fils de Pyrilampe était si beau que Calliclès, le démagogue du \italique{Gorgias} (481d), en avait fait son éromène. Mais Dêmos était aussi un obtus babillard : Platon le souligne malicieusement : Calliclès, malgré qu'il en ait, finissait toujours par céder devant son cadet aimé (αἰσθάνομαι οὖν σου ἑκάστοτε, καίπερ ὄντος δεινοῦ, ὅτι ἂν φῇ σου τὰ παιδικὰ καὶ ὅπως ἂν φῇ ἔχειν, οὐ δυναμένου ἀντιλέγειν). 
%Voilà pourquoi Eupolis (fr. 227(213) Kassel-Austin) nous le décrit comme \italique{bouché} ; littéralement : il est plein de cérumen (κυψέλη) : καὶ τῷ Πυριλάμπους ἆρα Δήμῳ κυψέλη ἔνεστιν;
%Il n'est donc pas indifférent qu'Athénée (IX, 397c-d), après avoir décrit le paon comme un animal bruyant qui réveille les dormeurs (τοιόνδε ταὧν, ὃς τοὺς εὕδοντας ἐγείρει) nous apprenne que le fils de Pyrilampe en faisait élevage, dont il était très fier.
%On comprend mieux les raisons -- outre sa philéliastite -- qui poussèrent Philocléon à rajouter (παρέγραψε πλησίον) : "κημὸς καλός" à côté de "Δῆμον καλόν" (noter l'accusatif) : une bonne muselière fera utilement taire le beau Démos.
%φιληλιαστής ἐστιν ὡς οὐδεὶς ἀνήρ·
%ἐρᾷ τε τούτου τοῦ δικάζειν, καὶ στένει
%ἢν μὴ ’πὶ τοῦ πρώτου καθίζηται ξύλου.   (90)
%ὕπνου δ’ ὁρᾷ τῆς νυκτὸς οὐδὲ πασπάλην·
%ἢν δ’ οὖν καταμύσῃ κἂν ἄχνην, ὅμως ἐκεῖ
%ὁ νοῦς πέτεται τὴν νύκτα περὶ τὴν κλεψύδραν.
%ὑπὸ τοῦ δὲ τὴν ψῆφόν γ’ ἔχειν εἰωθέναι
%τοὺς τρεῖς ξυνέχων τῶν δακτύλων ἀνίσταται,   (95)
%ὥσπερ λιβανωτὸν ἐπιτιθεὶς νουμηνίᾳ.
%καὶ νὴ Δί’ ἢν ἴδῃ γέ που γεγραμμένον
%υἱὸν Πυριλάμπους ἐν θύρᾳ Δῆμον καλόν, 
%ἰὼν παρέγραψε πλησίον “κημὸς καλός”.
%τὸν ἀλεκτρυόνα δ’, ὃς ᾖδ’ ἀφ’ ἑσπέρας, ἔφη @1    (100)
%ὄψ’ ἐξεγείρειν αὐτὸν ἀναπεπεισμένον,
%παρὰ τῶν ὑπευθύνων ἔχοντα χρήματα·
%εὐθὺς δ’ ἀπὸ δορπηστοῦ κέκραγεν ἐμβάδας,
%κἄπειτ’ ἐκεῖσ’ ἐλθὼν προκαθεύδει πρῲ πάνυ,
%
%Non, par le chien, ô Nicostratos, car Pliiloxène est un ruffian. — Vaine dépense de paroles, vous ne trouverez pas. Si vous tenez à savoir, faites silence, je vais vous expliquer la maladie du maître. Il a le goût de l'Héliée, comme personne. Son envie, ce n'est que cela : être dicaste ; et il se lamente s'il ne siège pas sur le premier banc. De sommeil, il n'en goûte pas un brin de la nuit, ou s'il ferme l'œil, fût-ce un instant, son esprit voltige là-bas, nuitamment, autour de la clepsydre. Oui, habitué qu'il est à tenir le suffrage, il se lève en serrant les trois doigts, comme s'il mettait l'encens sur l'autel à la nouvelle lune. Et, par Zeus, voit-il quelque part inscrit sur une porte le nom du fils de Pyrilampe: « Bravo, Démos! » (2), il s'en va écrire à côté: « Bravo, Kêmos! ». Le coq, qui chantait à l'entrée de la nuit (4), il prétendit qu'il s'était laissé séduire
%
%
%(2) L'épithète kalos, ajoutée à un nom, se rencontre fréquemment sur les vases peints et ne signifie nullement beau, pas plus que dans l'expression kalos xàgagôc, littér. bon et brave Elle correspond exactement au bravo des Italiens, et se dit pour marquer qu'on fait grand cas d'un homme ou d'une chose. On pourrait aussi traduire : « Vive Démos! » 
%
%
%κριοῦ διακονία: ὅταν προστάσσωμεν παιδίῳ διακονῆσαι, λέγομεν· δώσω σοι εἰ τύχοι ἀστραγάλους, ἢ ἰσχάδας. καὶ ὁ κριὸς οὖν εἶπε· καταλείψω τοὺς ἀστραγάλους σοι
%Prov. : Κόννου θρῖον (mis pour ἰσχάς ), bagatelle,  θρίου ψόφος (Vespae 436 : menaces vides)
%Postérieur : Βερενίκης [Βερονίκη: πόλις] πλόκαμος: τοῦτον κατηστερίσθαι φησὶ Κόνων. καὶ βόλος δέ τις ἀστραγάλων οὕτω καλεῖται
%Φιλοκλέων
%καὶ ποῖ τρέπεται δὴ 'πειτα τὰ χρήματα τἄλλα;
%Βδελυκλέων
%ἐς τούτους τοὺς ‘οὐχὶ προδώσω τὸν Ἀθηναίων κολοσυρτόν,
%ἀλλὰ μαχοῦμαι περὶ τοῦ πλήθους ἀεί.’ σὺ γὰρ ὦ πάτερ αὐτοὺς
%ἄρχειν αἱρεῖ σαυτοῦ τούτοις τοῖς ῥηματίοις περιπεφθείς.
%κᾆθ᾽ οὗτοι μὲν δωροδοκοῦσιν κατὰ πεντήκοντα τάλαντα
%670ἀπὸ τῶν πόλεων ἐπαπειλοῦντες τοιαυτὶ κἀναφοβοῦντες,
%‘δώσετε τὸν φόρον, ἢ βροντήσας τὴν πόλιν ὑμῶν ἀνατρέψω.’
%σὺ δὲ τῆς ἀρχῆς ἀγαπᾷς τῆς σῆς τοὺς ἀργελόφους περιτρώγων.
%οἱ δὲ ξύμμαχοι ὡς ᾔσθηνται τὸν μὲν σύρφακα τὸν ἄλλον
%ἐκ κηθαρίον λαγαριζόμενον καὶ τραγαλιζοντα τὸ μηδέν,
%675σὲ μὲν ἡγοῦνται Κόννου ψῆφον, τούτοισι δὲ δωροφοροῦσιν (sur le verbe hegoumai, voir Lamer 1939)
%ὕρχας οἶνον δάπιδας τυρὸν μέλι σήσαμα προσκεφάλαια
%φιάλας χλανίδας στεφάνους ὅρμους ἐκπώματα πλουθυγιείαν:
%σοὶ δ᾽ ὧν ἄρχεις, πολλὰ μὲν ἐν γῇ πολλὰ δ᾽ ἐφ᾽ ὑγρᾷ πιτυλεύσας,
%οὐδεὶς οὐδὲ σκορόδου κεφαλὴν τοῖς ἑψητοῖσι δίδωσιν.
%
%
%
%
%Φι.        ἰώ μοί μοι.   (750)
%Βδ.        οὗτος, τί βοᾷς;
%Φι.        μή μοι τούτων μηδὲν ὑπισχνοῦ.
%κείνων ἔραμαι, κεῖθι γενοίμαν,
%ἵν’ ὁ κῆρύξ φησι· “τίς ἀψήφιστος; ἀνιστάσθω.”  (Hésych. : ἄψηφον: πολύ. μέγα.ἰσχυρόν)
%κἀπισταίην ἐπὶ τοῖς κημοῖς    (755)
%ψηφιζομένων ὁ τελευταῖος. (de la dernière chance ? cf. Taillardat p. 153 n.1)
%“σπεῦδ’, ὦ ψυχή—”. ποῦ μοι ψυχή; “πάρες, ὦ σκιερά—”. (..fon : c'est le lieu de jeu ? voir Lamer 1936... ou un élément de Jeu ? (taillardat))
%μὰ τὸν Ἡρακλέα μή νυν ἔτ’ ἐγὼ ’ν τοῖσι δικασταῖς
%κλέπτοντα Κλέωνα λάβοιμι.
%Βδ.        ἴθ’, ὦ πάτερ, πρὸς τῶν θεῶν ἐμοὶ πιθοῦ.   (759)
%
%
%
%       ἐγὼ γὰρ ἀπεδάκρυσα νῦν γνώμην ἐμὴν
%οὐδέν ποτέ γ’ ἀλλ’ ἢ τῆς φακῆς ἐμπλήμενος.
%Βδ.        οὔκουν ἀποφεύγει δῆτα;    (985)
%Φι.                              χαλεπὸν εἰδέναι.   (985)
%Βδ.        ἴθ’, ὦ πατρίδιον, ἐπὶ τὰ βελτίω τρέπου.
%τηνδὶ λαβὼν τὴν ψῆφον ἐπὶ τὸν ὕστερον
%μύσας παρᾷξον κἀπόλυσον, ὦ πάτερ.
%Φι.        οὐ δῆτα· κιθαρίζειν γὰρ οὐκ ἐπίσταμαι.
%Βδ.        φέρε νύν σε τῃδὶ τὴν ταχίστην περιάγω.    (990)
%Φι.        ὅδ’ ἔσθ’ ὁ πρότερος;
%Βδ.                            οὗτος.    (991)
%
%
%Βδ.        ἀπολεῖς με. ποίας χάρακας; ἀλλ’ ὡς ἢ κάπρον
%ἐδιώκαθές ποτ’ ἢ λαγών, ἢ λαμπάδα
%ἔδραμες, ἀνευρὼν ὅ τι νεανικώτατον.
%Φι.        ἐγᾦδα τοίνυν τό γε νεανικώτατον· @1   (1205)
%ὅτε τὸν δρομέα Φάυλλον ὢν βούπαις ἔτι
%εἷλον διώκων λοιδορίας ψήφοιν δυοῖν.
%Βδ.        παῦ’· ἀλλὰ δευρὶ κατακλινεὶς προσμάνθανε
%ξυμποτικὸς εἶναι καὶ ξυνουσιαστικός.
%Φι.        πῶς οὖν κατακλινῶ; φράζ’ ἁνύσας.    (1210)
%Βδ.                                        εὐσχημόνως.   (1210)
%Φι.        ὡδὶ κελεύεις κατακλινῆναι;
%
%
%
%ἁθρόοι γὰρ ἥξομέν σε προσκαλούμενοι.
%Φι.        ἰὴ ἰεῦ, “καλούμενοι”.   (1335)
%ἀρχαῖά γ’ ὑμῶν. ἆρά γ’ ἴσθ’
%ὡς οὐδ’ ἀκούων ἀνέχομαι
%δικῶν; ἰαιβοῖ αἰβοῖ.
%τάδε μ’ ἀρέσκει· βάλλε κημούς.
%οὐκ ἄπει; ποῦ ’στ’ ἠλιαστής; ἐκποδών.   (1340)
%ἀνάβαινε δεῦρο, χρυσομηλολόνθιον,  (il y a un jeu qui un nom proche)
%τῇ χειρὶ τουδὶ λαβομένη τοῦ σχοινίου.
%ἔχου· φυλάττου δ’, ὡς σαπρὸν τὸ σχοινίον·
%ὅμως γε μέντοι τριβόμενον οὐκ ἄχθεται.
%
%
%
%Ranea : compter avec des pierres ?
%
%       ΔΙ.                 Καὶ μὴν λογιοῦμαι ταῦτα τῶν ψήφων λαβών.
%ΕΥ.               Φθιῶτ’ Ἀχιλλεῦ, τί ποτ’ ἀνδροδάικτον ἀκούων @1
%ἰὴ κόπον οὐ πελάθεις ἐπ’ ἀρωγάν;
%Ἑρμᾶν μὲν πρόγονον τίομεν γένος οἱ περὶ λίμναν.   (1265)
%Ἰὴ κόπον οὐ πελάθεις ἐπ’ ἀρωγάν;
%ΔΙ.                   Δύο σοι κόπω, Αἰσχύλε, τούτω.
%
%LSJ lagara : τὴν πόλιν ἀντὶ λαγαρᾶς καὶ ὑποσόμφου μεστὴν ἐποίησεν ἀγλαΐας Them. Or. 18.222d
%
%Platon Lysis 206e (sur l'artiasmos) : Ταῦτα, ἦν δ' ἐγώ, χρὴ ποιεῖν. Καὶ ἅμα λαβὼν τὸν [206e] Κτήσιππον προσῇα εἰς τὴν παλαίστραν· οἱ δ' ἄλλοι ὕστεροι ἡμῶν ᾖσαν. Εἰσελθόντες δὲ κατελάβομεν αὐτόθι τεθυκότας τε τοὺς παῖδας καὶ τὰ περὶ τὰ ἱερεῖα σχεδόν τι ἤδη πεποιημένα, ἀστραγαλίζοντάς τε δὴ καὶ κεκοσμημένους ἅπαντας. Οἱ μὲν οὖν πολλοὶ ἐν τῇ αὐλῇ ἔπαιζον ἔξω, οἱ δέ τινες τοῦ ἀποδυτηρίου ἐν γωνίᾳ ἠρτίαζον ἀστραγάλοις παμπόλλοις, ἐκ φορμίσκων τινῶν προαιρούμενοι· τούτους δὲ περιέστασαν ἄλλοι θεωροῦντες. Ὧν δὴ καὶ ὁ Λύσις ἦν, καὶ εἱστήκει ἐν [207a] τοῖς παισί τε καὶ νεανίσκοις ἐστεφανωμένος
%
%Vesp. 74
%
%
%Ploutos :
%
%Χρεμύλος
%555ὡς μακαρίτην ὦ Δάματερ τὸν βίον αὐτοῦ κατέλεξας,
%εἰ φεισάμενος καὶ μοχθήσας καταλείψει μηδὲ ταφῆναι.
%Πενία
%σκώπτειν πειρᾷ καὶ κωμῳδεῖν τοῦ σπουδάζειν ἀμελήσας,
%οὐ γιγνώσκων ὅτι τοῦ Πλούτου παρέχω βελτίονας ἄνδρας
%καὶ τὴν γνώμην καὶ τὴν ἰδέαν. παρὰ τῷ μὲν γὰρ ποδαγρῶντες
%560καὶ γαστρώδεις καὶ παχύκνημοι καὶ πίονές εἰσιν ἀσελγῶς,
%παρ᾽ ἐμοὶ δ᾽ ἰσχνοὶ καὶ σφηκώδεις καὶ τοῖς ἐχθροῖς ἀνιαροί.
%Χρεμύλος
%ἀπὸ τοῦ λιμοῦ γὰρ ἴσως αὐτοῖς τὸ σφηκῶδες σὺ πορίζεις.
%Πενία
%περὶ σωφροσύνης ἤδη τοίνυν περανῶ σφῷν κἀναδιδάξω
%ὅτι κοσμιότης οἰκεῖ μετ᾽ ἐμοῦ, τοῦ Πλούτου δ᾽ ἐστὶι ὑβρίζειν.
%Χρεμύλος
%565πάνυ γοῦν κλέπτειν κόσμιόν ἐστιν καὶ τοὺς τοίχους διορύττειν.
%Βλεψίδημος
%†νὴ τὸν Δί᾽, εἰ δεῖ λαθεῖν αὐτόν, πῶς οὐχὶ κόσμιόν ἐστι;†
%Πενία
%σκέψαι τοίνυν ἐν ταῖς πόλεσιν τοὺς ῥήτορας, ὡς ὁπόταν μὲν
%ὦσι πένητες, περὶ τὸν δῆμον καὶ τὴν πόλιν εἰσὶ δίκαιοι,
%πλουτήσαντες δ᾽ ἀπὸ τῶν κοινῶν παραχρῆμ᾽ ἄδικοι γεγένηνται,
%570ἐπιβουλεύουσί τε τῷ πλήθει καὶ τῷ δήμῳ πολεμοῦσιν.
%
%
%Est-ce la TROPA qui est décrite ?? 
%
%Pollux IX 103 :
%τὸ μὲν οὖν ἀστραγάλοις παίζειν καὶ ἀστραγαλίζειν καὶ ἀστερίζειν ἔνιοι τῶν ποιητῶν εἰρήκασιν, ὅτι τοὺς ἀστραγάλους καὶ ἀστρίας εἰσὶν οἳ ὠνόμαζον, Ἀντιφάνης δὲ καὶ ἀστρίχους. 100 Τὸ δὲ σχῆμα τοῦ κατὰ τὸν ἀστράγαλον πτώματος ἀριθμοῦ δόξαν εἶχεν, καὶ τὸ μὲν μονάδα δηλοῦν καλεῖται κύων, τὸ δὲ ἀντικείμενον χιάς, καὶ Χῖος οὗτος ὁ βόλος. Δυὰς δὲ καὶ πεντὰς ἐν ἀστραγάλοις, ὥσπερ ἐν κύβοις, οὐκ ἔνεστιν. Οἱ δὲ πλείους τὸν μὲν ἑξίτην Κῷον, τὸν δὲ κύνα Χῖον καλεῖσθαι λέγουσιν. Καὶ μὴν καὶ Στησίχορος ἐκαλεῖτό τις παρὰ τοῖς ἀστραγαλίζουσιν ἀριθμός, ὃς ἐδήλου τὰ ὀκτώ· τὸν γὰρ ἐν Ἱμέρᾳ τοῦ ποιητοῦ τάφον ἐξ ὀκτὼ πάντων συντεθέντα πεποιηκέναι τὴν 'πάντ'ὀκτώ΄ φασι παροιμίαν9. 101 Ἐπεὶ δὲ τοῖς τετταράκοντα τοῖς μετὰ τοὺς τριάκοντα προστᾶσιν Ἀθήνησι συνῆρξεν Εὑριππίδης, εἰ τετταράκοντα συνήθροιζεν ἀστραγάλων βολή, τὸν ἀριθμὸν τοῦτον Εὑριππίδην ὠνόμαζον. Καὶ μὴν καὶ ἀρτιάζειν ἀστραγάλους ἐκ φορμίσκων κατερωμένους ἐν τῷ ἀποδυτηρίῳ τοὺς παῖδας ὁ Πλάτων (Lys. 206Ε) ἔφη· τὸ δ΄ἀρτιάζειν ἐν ἀστραγάλων πλήθει, κεκρυμμένων ὑπὸ ταῖν χεροῖν, μαντεῖαν εἶχε τῶν ἀρτίων ἢ καὶ περιττῶν. Τἀυτὸν δὲ τοῦτο καὶ κυάμοις ἢ καρύοις ἢ ἀμυγδάλαις, οἱ δὲ καὶ ἀργυρίῳ πράττειν ἠξίουν, εἰ πιστὸς Ἀριστοφάνης ἐν Πλούτῳ (816) λέγων· στατῆρσι δ'οἱ θεράποντες ἀρτιάζομεν. 102 Εἰ μὲν οὖν κύκλου περιγραφέντος ἀφιέντες ἀστράγαλον ἐστοχάζοντο τοῦ μεῖναι τὸν βληθέντα ἐν τῷ κύκλῳ, ταύτην εἰς ὤμιλλαν τὴν παιδιὰν ὠνόμαζον. […] 103 ἡ δὲ τρόπα καλουμένη παιδιὰ γίνεται μὲν ὡς τὸ πολὺ δι΄ἀστραγάλων, οὓς ἀφιέντες στοχάζονται βόθρου τινὸς εἰς ὑποδοχὴν τῆς τοαιύτης ῥίψεως ἐξεπίτηδες πεποιημένου· πολλάκις δὲ καὶ ἀκύλοις καὶ βαλάνοις ἀντὶ τῶν ἀστραγάλων οἱ ῥίπτοντες ἐχρῶντο.
%Alcuni poeti hanno chiamato anche astrizein il fatto di giocare agli astragali, perché chiamano gli astragali astragaloi e astriai10, Antifane dice anche astrichoi. 100 Lo schema secondo la maniera di cadere dell’astragalo ha un valore numerico e il lato dell’unità si chiama il cane (kyôn), quello opposto chias e questo colpo Chio. Non ci sono valori di due o di cinque negli astragali e nemmeno nei dadi. I più dicono che si dà alla faccia che vale 6 il nome di Koos e al cane quello di Chios. E un numero era anche detto Stesicoro dai giocatori agli astragali, era l´otto. Si dice, infatti, che la tomba del poeta a Imera era un assemblaggio di otto parti, donde il proverbio: «tutto fa otto.» 101 Poiché Eurippide dirigeva i Quaranta che governavano ad Atene dopo i Trenta, se il tiro di astragali dava in totale 40, era detto Eurippide. E Platone ha detto che i bambini giocavano a pari e dispari nel vestiario dei bagni con gli astragali che prendevano nei bussolotti. Giocare a pari e dispari sulla somma degli astragali nascosti nelle mani era la divinazione dei numeri pari e dispari. E alcuni facevano lo stesso con fave, noci, mandorle o monete, se si dà credito ad Aristofane, che dice nel Pluto: «Noialtri, gli schiavi giochiamo a pari e dispari con gli stateri». 102 Se si descrive un cerchio, lanciando un astragalo con tanta abilità per farlo restare nel cerchio, questo gioco prende nome di ômilla. [...] 103 Il gioco detto tropa si fa prevalentemente con gli astragali, che vengono lanciati nello sforzo di raggiungere una fossetta scavata in modo da poter contenere tutti gli astragali che vi sono gettati dentro. Sovente quelli che lanciano si servono di ghiande e castagne invece che di astragali. (traduzione mia)
%
%
%
%Sur les faces de l'osselet :
%Aristotele nella nota discussione anatomica dell’Historia Animalium II 1 499b 26-31: 
%Πάντα δὲ τὰ ἔχοντα ἀστραγάλους ἐν τοῖς ὄπισθεν ἔχει σκέλεσιν. Ἔχει δ’ ὀρθὸν τὸν ἀστράγαλον ἐν τῇ καμπῇ, τὸ μὲν πράνης ἔξω, τὸ δ'ὕπτιον εἴσω, καὶ τὰ μὲν κῷα ἐντὸς ἐστραμμένα πρὸς ἄλληλα, τὰ δὲ χῖα καλούμενα ἔξω, καὶ τὰς κεραίας ἄνω. Ἡ μὲν οὖν θέσις τῶν ἀστραγάλων τοῖς ἔχουσι πᾶσι τοῦτον ἔχει τὸν τρόπον. 
%Tutti gli animali che hanno gli astragali li presentano nelle zampe posteriori. Hanno l’astragalo posto verticalmente nell’articolazione, la parte superiore posizionata verso l’esterno, quella inferiore verso l’interno; i lati di Coo sono rivolti all’interno uno verso l’altro; i lati che sono detti di Chio verso l’esterno, e le parti con le corna verso l’alto. La posizione degli astragali è dunque di tal sorta in tutti gli animali che ne siano provvisti. (traduzione mia)
%
%Ancora nel De caelo II, 12, 292a 28-292b 1 una similitudine sulla legge delle probabilità e dei grandi numeri è desunta dal gioco degli astragali: 
%Ἔστι δὲ τὸ κατορθοῦν χαλεπὸν ἢ τὸ πολλὰ ἢ τὸ πολλάκις, οἷον μυρίους ἀστραγάλους Χίους βαλεῖν ἀμήχανον, ἀλλ’ ἕνα ἢ δύο ῥᾷον. Καὶ πάλιν ὅταν τοδὶ μὲν δέῃ τοῦδ’ ἕνεκα ποιῆσαι, τοῦτο δ’ ἄλλου καὶ τοῦτο ἑτέρου, ἐν μὲν ἑνὶ ἢ δυσὶ ῥᾴδιον ἐπιτυχεῖν, ὅσῳ δ’ ἂν διὰ πλειόνων, χαλεπώτερον. 
%È difficile, peraltro, avere un buon esito in molte cose o di frequente, così per esempio agli astragali non riesce di fare innumerevoli volte il colpo di Chio, mentre è facile farlo una volta o due. E ancora di nuovo, quando si deve fare una tal cosa a un tal fine, la talaltra ad un secondo fine ed un’altra ancora ad un fine ancora differente, è facile riuscirvi in un caso o due, ma ciò è tanto più difficile, quanto più si moltiplicano gli impegni da assolvere. (traduzione mia)
%Dunque, per riassumere i quattro lati dell’astragalo di forma pressoché cubica ricevono differenti denominazioni e punteggi che possiamo schematizzare (cfr. Nollé 2007: 8-9): 
%1) Χῖος, Chios, o κύων, cane, lato stretto a forma di S): 1 punto; 
%2) ὕπτιος dorso (lato largo, concavo): 3 punti; 
%3) πρανής, ventre (lato largo, convesso): 4 punti; 
%4) Κῷος, Cos, 6 punti (lato stretto, piatto): 6 punti.
%Moirai, alle Tychai (Tyche egemone, Tyche Salvatrice o ancora Tyche Felice o fonte di felicità) e alle Nikai (Vincitrice e Vittoria gioiosa). Altre divinità del destino si riconoscono in Daimon, Agathos Daimon e Nemesis,
%
%
%Τὰς ἐν τῇ φαρέτρᾳ ψηφῖδας: Φύλαρχος φησὶ τοὺς Σκύθας μέλλοντας καθεύδειν ἄγειν τὴν φαρέτραν, καὶ εἰ μὲν ἀλύπως τύχοιεν τὴν ἡμέραν ἐκείνην διαγαγόντες, καθιέναι εἰς τὴν φαρέτραν ψηφῖδα λευκήν· εἰ δὲ ὀχληρῶς, μέλαιναν. Ἐπὶ τοίνυν τῶν ἀποθνησκόντων ἐκφέρειν τὰς φαρέτρας καὶ ἀριθμεῖν τὰς ψήφους· καὶ εἰ εὑρεθείησαν πλείους αἱ λευκαί, εὐδαιμονίζειν τὸν ἀπογενόμενον. Ὅθεν παροιμιασθῆναι, τὴν ἀγαθὴν ἥμῶν ἡμέραν λεγόντων τῶν ἐκ τῆς φαρέτρας εἶναι. Καὶ Μένανδρος δὲ φησὶν ἐν Λευκαδίᾳ, τὴν ἀγαθὴν ἡμέραν λευκὴν καλεῖσθαι.
%Xifothếkè dans l'original euripidéen ?
}%
\debutNote{1031}%
\note{1}{Cf. Pollux (X, 150) sur les ustensiles de jeu : %
\citComm{Κυβευτοῦ σκεύη · ἄβαξ ἀβάκιον, κόσκινον, κύβοι, διάσειστοι κύβοι, ἀστράγαλοι, φιμοί, κημοί, τηλία, κήθια κηθίδια, ψῆφοι, πεττοί.}%
φιμός et κημός sont littéralement des \italique{muselières}, avec le sens spécialisé de \italique{cornets} (à dés ou astragales) et sont synonymes de κήθιον. On pourra par exemple se reporter à l'article de Michel Casevitz, «~Sur les termes de jeu : compléments d'histoire des mots~», \italique{Kentron} 36 (2021), p. 209-212. Pour des ouvrages beaucoup plus conséquents, le lecteur pourra lire \italique{Les jeux des Anciens}\expo{2} (1873) de Louis Becq de Fouquières, auteur à la vaste érudition et au style des plus agréables comme le \nieme{XIX} siècle en a été prodigue ; l'édition par Jean Taillardat déjà mentionnée du Περὶ βλασφημιῶν, \italique{Des termes injurieux} et du Περὶ παιδιῶν, \italique{Des jeux grecs}, de Suétone ; la somme de référence qu'est l'article de Hans Lamer «~\latin{Lusoria tabula}~» de la \italique{Paulys Realencyclopädie der classischen Altertumswissenschaft}, \allemand{Band} XIII 2 (1927), col. 1900-2029 ; plus récemment, les \latin{excerpta de ludis} chez Pollux ont été publiés par Salvatore Costanza : \italique{Giulio polluce, onomasticon: excerpta de ludis} accompagnés d'une traduction italienne. Outre l'archéologie, nos témoins principaux quant aux jeux grecs sont Pollux, Hésychius et Eustathe (qui s'appuie sur les traités perdus de Suétone) -- non exempts de contradictions entre eux d'ailleurs.}%
\note{2}{Van Leeuwen a recensé d'autres erreurs similaires où l'ordre des mots est inversé dans le \italique{Ravennas} (la comparaison avec les neuf pièces du \italique{Ravennas} qui nous sont connues par d'autres témoins est préférable : la correction est davantage sûre). Cf. l'édition en fac-similé qu'il fit de notre manuscrit : \italique{Comoediae Undecim cum Scholiis : Codex Ravennas 137, 4, A} [$=$ \italique{Rav.} 429 nouvelle cote], \italique{phototypice editus} in \italique{Codices Graeci et Latini photographice depicti}, t. IX (1904), pp. XVIII-XIX.}%
\note{3}{La scène suivante est une amusante parodie du procès intenté, ou censé intenté, au stratège athénien Lachès, qui était en effet du dême d'Æxone. Le chien du peuple, c'est Cléon, du dême de Cydathène. Lachès avait commandé la flotte envoyée en Sicile en 427. Rappelé au commencement de 425, il fut accusé d'avoir malversé. Le procès (s'il eut lieu) fut présidé suivant l'usage par un thesmothète (v. 935). Cléon, comme synégore, soutint l'accusation. Lachès fut acquitté, car il remplit encore diverses fonctions publiques. Il est intéressant de constater que si le poète n'épargne pas Cléon, il n'est guère plus indulgent pour Lachès, encore que celui-ci appartînt au parti aristocratique. (Note de Willems.)}}
\end{groupe}%
%% v. 1015-1055
\begin{groupe}%
\typeVers{$+$ trim. iamb.}\paratragiques{ἀλλ᾽ ἐν πυκνοῖς δεσμοῖσιν ἐμπεπλεγμένη}
\end{groupe}%
%% v. 1015-1055
\begin{groupe}%
\typeVers{$3$ ia. sync\cat}\paratragiques{\ligne{1033}κήτει βορὰ Γλαυκέτῃ πρόκειμαι,%
\commentaire{1033}{κήτει βορὰ Γλαυκέτῃ πρόκειμαι.}{On sait par le scholiaste des \italique{Oiseaux} (vers 348) que la pièce originale d'Euripide comportait à cet endroit :%
\citComm{ἐκθεῖναι κήτεϊ φορβάν,}%
Γλαυκέτης, l'Athénien glouton (τένθης) est cité au vers 1008 de la \italique{Paix}.
%\citComm{καὶ τὴν ἀγορὰν ἡμῖν ἀγαθῶν\\
%ἐμπλησθῆναι, μεγάλων σκορόδων,\\
%σικύων πρῴων, μήλων, ῥοιῶν.\\
%δούλοισι χλανισκιδίων μικρῶν·\\
%κἀκ Βοιωτῶν γε φέροντας ἰδεῖν\\
%χῆνας νήττας φάττας τροχίλους·\\
%καὶ Κωπᾴδων ἐλθεῖν σπυρίδας,\\
%καὶ περὶ ταύτας ἡμᾶς ἁθρόους\\
%ὀψωνοῦντας τυρβάζεσθαι\\
%Μορύχῳ Τελέᾳ \textbf{Γλαυκέτῃ}, ἄλλοις\\
%τένθαις πολλοῖς.\\
%Fais que notre marché regorge de bonnes choses : denrées de Mégare, aulx, concombres précoces, coings, grenades, petits manteaux pour esclaves ; de chez les Béotiens, qu'on y voie affluer des gens chargés d'oies, de canards, de ramiers, de pluviers ; que les anguilles du Copaïs viennent par paniers et que, massés autour d'elles pour faire nos provisions, nous nous coudoyions avec Morychos, Téléas, \textbf{Glaucetès} et bien d'autres gourmands.}%
Trygée y fait une prière pour que la paix revienne à Athènes ; la ville est sens dessus dessous du fait de la guerre : les embargos et autres razzias ont littéralement renversé le marché : au lieu de toutes les bonnes denrées, l'Athénien n'y rencontre que des rappels de ses privations. 
Autre inversion : si la paix revenait, ce serait alors aux citoyens de bousculer les trois parasites du marché aux anguilles. Aristophane joue sur une image qu'il a développée ailleurs\marqueur{1} : comme le chasseur d'anguille trouble l'eau pour capturer sa proie, le politicien véreux trouble les affaires publiques pour mieux s'enrichir. (Remarquer l'utilisation du verbe τυρβάζεσθαι, qu'Aristophane emploie aussi au vers 257 des \italique{Guêpes} : τὸν πηλὸν ὥσπερ ἀτταγᾶς τυρβάσεις βαδίζων.) 
On peut mettre notre main au feu que le groupe des trois affreux étaient de francs coquins qui se gobergeaient sur le dos d'Athènes pour mener une vie douce, bonne et grasse\marqueur{2}.
Γλαυκέτης est mis par paronomase ; \italique{Nomen, omen} : lorsqu'on s'appelle Γλαυκέτης, on est prédestiné à devenir ce poisson vorace γλαῦκος-κῆτος\marqueur{3} !
(L'image de la femme livrée en pâture (βορά) à la poissonnaille est récurrente chez les Comiques, \italique{confer} Athénée (VII, 315b) citant Platon le Comique ($=$ fr. 57(56) Kassel-Austin)) :%
\citComm{σὲ γάρ, γραῦ, συγκατῴκισεν σαπρὰν\\
ὀρφῷσι σελαχίοις τε καὶ φάγροις βοράν}%
(βορά désigne ici la menue fretaille qui sert de pâture aux poissons carnivores comme l'orphe (ὄρφος / ὀρφώς), qu'Athénée dans le même chapitre rapproche du \italique{glauque}, γλαῦκος.)\\
}%
\debutNote{1033}%
\note{1}{Cavaliers, vers 864-867. Voir aussi le vers 361. Érasme a relevé l'image dans ses \italique{Adages} : n° 2576 (\latin{Anguillas captare}).}%
\note{2}{Cf. pour Morykhos les \italique{Acharniens}, v. 887 ainsi que les \italique{Guêpes}, v. 506 \& 1142. On notera au passage qu'un Téléas fils de Télénikos apparaît à la même époque dans les \italique{Inscriptiones Graecae} (I\expo{3}, 307 ; 308 ; 331 ; 332 ; 370) à la secrétairie (ἐγραμμάτευε) du Trésor (ℎοι ταμίαι το͂ν ℎιερο͂ν χρεμάτον τε͂ς Ἀθεναίας). Voir ce qu'en disent les scholies aux vers 167-8 des \italique{Oiseaux}.}%
\note{3}{Autre jeu du même acabit sur le nom de Glaucétès : (Platon le Comique, fr. 114(106) K.-A.) :%
\citComm{ὦ θεῖε Μόρυχε· πῶς γὰρ οὐ δαίμων ἔφυς;\\
καὶ Γλαυκέτης ἡ ψῆττα, καὶ Λεωγόρας,\\
οἵ ζῆτε τερπνῶς οὐδὲν ἐνθυμούμενοι.}}}
\end{groupe}%
%% v. 1015-1055
\begin{groupe}%
\typeVers{dim. ia\cat}\paratragiques{γαμηλίῳ μὲν οὐ ξὺν\apparat{1034-5}{ξυμπαιῶνι R : \lecon{ξὺν παιῶνι} Zanetti (\italique{cum carmuine} Divus)}}
\end{groupe}%
%% v. 1015-1055
\begin{groupe}%
\typeVers{$+$ 2 ia\cat}\paratragiques{παιῶνι δεσμίῳ δὲ}
\end{groupe}%
%% v. 1015-1055
\begin{groupe}%
\typeVers{dim. iamb.}\paratragiques{γοᾶσθέ μ᾽ ὦ γυναῖκες ὡς}
\end{groupe}%
%% v. 1015-1055
\begin{groupe}%
\typeVers{$2$ troch.}\paratragiques{μέλεα μὲν πέπονθα μέλεος,}
\end{groupe}%
%% v. 1015-1055
\begin{groupe}%
\typeVers{lecyth.}\paratragiques{ὢ τάλας ἐγὼ τάλας,}
\end{groupe}%
%% v. 1015-1055
\begin{groupe}%
\typeVers{dochm. $+$ ...}\paratragiques{\ligne{1039}ἀπὸ δὲ συγγόνων ἀλλ᾽ ἄνομα πάθεα,\apparat{1039}{ἀλλὰν R : \lecon{ἀλλ᾽} Scaliger}}
\end{groupe}%
%% v. 1015-1055
\begin{groupe}%
\typeVers{}\paratragiques{†φῶτά λιτομένα†,\apparat{1040}{λιτομέναν R : ἀντομένα (vel -αν) Σ : \lecon{λιτομένα} Enger \ligneApp{1041} φεύγουσαν R : χέουσα Rau (\italique{Paratragödia}, p.75) : φλέγουσαν Mugrave (ad Eur. \italique{Or.}1402$=$1396) : \lecon{ψεφαυγοῦς} Bibi, cf Σ ad v.1054, Hesych. ψ 135}% : ἀλιτόμην ἄν Bibi}%
%\commentaire{1040}{φῶτα λιτομέναν.}{
%ἠλιτόμηνιen celui qui appelle en vain Hésychius ἠλιτομήνιεν
%264 φέγγουσαν· λαμπρύνουσαν
%1104 φῶντα:λάμποντα
%2839 πολυδάκρυτον· πολλῶν δακρύων αἴτιον
%Alceste :
%  εἴθ’ ἐπ’ ἐμοὶ μὲν εἴη,       [στρ.    (454)
%δυναίμαν δέ σε πέμψαι    (455)
%φάος ἐξ Ἀίδα τεράμνων
%καὶ Κωκυτοῖο ῥεέθρων
%Hippolyte :
%κῶμος λέλακεν, Ἄρτεμιν τιμῶν θεὰν    (55)
%ὕμνοισιν· οὐ γὰρ οἶδ’ ἀνεῳγμένας πύλας
%Ἅιδου, φάος δὲ λοίσθιον βλέπων τόδε.
%Hécube :
%  ὡς πρὶν σφαγῆναί γ’ ἐκτέτηκα καρδίαν
%θρήνοισι μητρὸς τήνδε τ’ ἐκτήκω γόοις.
%ὦ φῶς· προσειπεῖν γὰρ σὸν ὄνομ’ ἔξεστί μοι,    (435)
%μέτεστι δ’ οὐδὲν πλὴν ὅσον χρόνον ξίφους
%βαίνω μεταξὺ καὶ πυρᾶς Ἀχιλλέως.
%
%Une scholie nous apprend une variante parmi les anciens manuscrits :%
%\citComm{γρ[άφεται] καὶ `φῶτα ἀντομένα'}%
%%Et une autre :%
%%\citComm{καὶ τοῦτο ἔχει νοῦν πρὸς τὰ ἑξῆς, `φῶτ' ἀντομέναν, ὃς ἔμ᾽ ἀπεξύρησε πρῶτον }%
%Précieuse scholie, car si le nu Ν de ἀντομένα provient très certainement d'une mélecture du digramme ΛΙ ; elle nous apprend que ledit nu était précédé d'un double alpha ΑΑ dont seul un a subsisté dans la leçon du \italique{Ravennas}. (On a déjà évoqué ce genre de coquilles, très fréquent dans le manuscrit de Ravenne. Cf. par exemples les vers 834 ou 555.)\\
%Je propose donc pour l'\allemand{urtext} des \italique{Thesmophories} :%
%\citComm{ΦΩΤΑΑΛΙΤΟΜΗΝΑΝ\\
%φῶτα ἀλιτόμην ἄν}%
%Cette leçon -- qu'on espère être l'originale -- ayant dérivé en :%
%\citComm{ΦΩΤΑΛΙΤΟΜΕΝΑΝ (\italique{Ravennas})\\
%φῶτα λιτομέναν.\\
%ΦΩΤΑΑΝΤΟΜΕΝΑ(Ν) (Scholiaste)\\
%φῶτα ἀντομένα(ν)}%
%ἀλιτόμην est l'aoriste du verbe ἀλιταίνω, \italique{commettre une faute}, pris ici au sens d'\italique{expier} (\eff{Cf. Chantraine à  aleiths et Aristophane et Hésych alitwmai : atimw}). fwta (accusatif de durée) est pris dans son acception courante de \italique{jours}.
%Traduisez donc :%
%\citComm{ἀπὸ δὲ συγγόνων ἀλλ᾽ ἄνομα πάθεα\\
%φῶτα ἀλιτόμην ἄν,\\
%πολυδάκρυτον Ἀίδα γόον φρύγουσαν\marqueur{Pour la correction : φρύγουσαν à la place de la leçon du manuscrit φεύγουσαν, voir le commentaire au vers suivant.}.\\
%Par la faute de mes parents, j'endurerais des tortures impies pendant des jours, me desséchant pour l'Hadès d'un torrent de larmes.}%
%vἀλιτόμην ἄν Par la faute de sa mère/Euripide...
%	j'expierais des jours durant
%Eschyle suppliantes 416 : ὃς οὐδ’ ἐν Ἅιδου τὸν ἀλιτόντ’ ἐλευθεροῖ. (ed. Murray du TLG... délire ? ton thanonta alt. edd.)
%Hésychius : “ἐκ δ᾽ἄγαγε πρὸ φόωσδε᾽: προ εξήγαγε δὲ εἰς τὸ φῶς τὸν ἀλιτόμηνον Homéri
% ἠλιτόμηνα: ἐλαττούμενα τοῖς μησίν
% ἠλῐτό-μηνις· ὁ μάτην ἐγκαλῶν
%Crux : ἀφήγουσα:μεταστρέφουσα
%φεγγουσα
%φρύγ
%φθεγγουσα <- n'existe pas
%L'expression πολυδάκρυτος γόος est homérique (\italique{Odyssée}, XIX, 213).
%}%
}
\end{groupe}%
%% v. 1015-1055
\begin{groupe}%
\typeVers{do. $+$ cr. $+$ ba.}\paratragiques{πολυδάκρυτον Ἀίδα γόον ψεφαυγοῦς.%\apparat{1041}{φεύγουσαν R : χέουσα Rau (\italique{Paratragödia}, p.75) : φλέγουσαν Mugrave (ad Eur. \italique{Or.}1402$=$1396) : \lecon{ψεφαυγοῦς} Bibi, cf Σ ad v.1054 e Hesych. ψ 135}%
\commentaire{1041}{Ἀίδα,}{Aristophane use du même génitif dorique au vers 1333 des \italique{Grenouilles} (parodie d'Euripide).}%
\commentaire{1041}{πολυδάκρυτον Ἀίδα γόον φεύγουσαν.}{Le manuscrit porte l'incompréhensible †φεύγουσαν†. Y suppléer par χέουσα(ν), à l'instar de certains éditeurs, c'est faire fi de sa leçon. Je dirais même plus : χέουσα est le mot banal qu'on attendait, il paraît alors très difficile qu'un scribe s'y soit trompé. (σφριγῶσα, \italique{gonflée}, aurait au moins le mérite de se justifier paléographiquement mais le terme paraît impropre en grec à l'image d'yeux gonflés de larmes.)
Une autre proposition, `φλέγουσα(ν)', a aussi emporté la faveur de nombreuses éditions. Elle est reprise par Jean Taillardat à propos de l'image de la Muse-flamme (§736 n.2)... mais le grand Taillardat me semble s'abuser : si un poème, un chant, un \italique{thrêne} (qui est avant tout sonore ; la flamme étant celle de la voix selon moi) peuvent être enflammés, cela ne prouve pas pour autant que des pleurs peuvent l'être.\\
Le scholiaste nous fournit une information qui pourrait peut-être bien débrouiller l'écheveau: une glose, σκοτεινήν, \italique{obscur}, se trouve au vers 1054 pour qualifier soi-disant αἰόλαν, \italique{rapide}, \italique{mobile}, \italique{changeant}. 
Sauf qu'αἰόλος n'a jamais eu le sens de σκοτεινός et que Willems a clairement montré qu'il fallait prendre αἰόλος en son sens courant\marqueur{1}.
Suidas (αι 245) -- ou l'un des lexicographes de la \italique{Souda}, si l'on préfère -- n'a pas manqué néanmoins de relever la glose à l'entrée Αἰόλαν. Mais Suidas est coutumier de ce genre de fausses attributions (voir la note au vers 735).
Or en cherchant des mots proches de †φεύγουσαν†, je suis retombé sur ce cher Hésychius (ψ 135) :%
\citComm{ψεφαυγοῦς· σκοτεινῆς,}%
Génitif dorique (comme Ἀίδα) d'un certain *ψεφαυγής, \italique{au sombre regard}.
Euripide a une prédilection pour forger des composés en -αυγής. (\italique{Confer} le Chantraine, \italique{sub voce} αὐγή.) 
Si l'on fait l'hypothèse de la chute de `ψε', le φαυγους restant aurait bien pu être complété d'un `αν' en regard de l'autre participe †λιτομέναν† juste au-dessus, et donner ainsi notre leçon †φεύγουσαν†.
%Point tout-à-fait exacte, car ce n'est pas l'image du \italique{feu} qu'évoque Andromède, mais celle du \italique{dessèchement}.
%Lisez donc `φρύγουσα' ; un epsilon Ε se sera substitué au rhô Ρ\marqueur{Se reporter au tableau comparatif de l'évolution des crusives grecques du livre déjà cité\eff{The Greek Manuscripts ?} pour constater à quel point ces deux lettres ont pu se ressembler sur certains manuscrits}.\\
%Le verbe φρύγειν, littéralement : \italique{faire griller}, est d'un emploi assez large. Ainsi, les dérivés φρύγιον et φρύγανον désignent du \italique{bois sec}. Une épigramme funéraire de l'\italique{Anthologie Palatine} (VII, 293) fait de vents le sujet du verbe φρύγειν, qui ont \italique{desséché} Nicophềme au point de le faire périr de soif :%
%\citComm{οὐ χεῖμα Νικόφημον, οὐκ ἄστρων δύσις ἁλὸς Λιβύσσης κύμασιν κατέκλυσεν ἀλλ ἐν γαλήνῃ, φεῦ τάλας,\\
%ἀνηνέμῳ πλόῳ πεδηθείς, \textbf{ἐφρύγη δίψευς ὕπο}}%
%frugein = mourir ?
%Asséchant son torrent lacrimal (πολυδάκρυτος γόος\marqueur{L'expression est homérique (\italique{Odyssée}, XIX, 213) ; on trouve chez Eschyle (\italique{Choéphores}, vers 449) :%
%\citComm{χέουσα πολύδακρυν γόον κεκρυμμένα.}%
%Et chez Euripide (\italique{Suppliantes}, vers 773):%
%\citComm{Ἅιδου τε μολπὰς ἐκχέω δακρυρρόους.}}) qu'Andromède déverse l'\italique{assèchent l}
%σφριγάω σφριγῶσα
%Hésychius : σφριγῶσα· πεπυκνωμένη, συνεστηκυῖα, ἰσχύουσα
%cf. Taillardat §62.
}%
\debutNote{1041}%
\note{1}{Note n° 27, reproduite ci-après.}}
\end{groupe}%
%% v. 1015-1055
\begin{groupe}%
\typeVers{extr. metr.}\paratragiques{αἰαῖ αἰαῖ ἒ ἒ,}
\end{groupe}%
%% v. 1015-1055
\begin{groupe}%
\typeVers{dim .troch.}\paratragiques{ὃς ἔμ᾽ ἀπεξύρησε πρῶτον,%
\commentaire{1043-51}{Papyrus d'Oxyrhynque n° 4935.}{Un maigre fragment (n° 1) du Papyrus d'Oxyrhynque 4935\marqueur{1} (\nieme{II} siècle ap. J.-C.), reproduit en annexe, nous apporte une leçon nouvelle au vers 1047 : le `μοι' du \italique{Ravennas} n'y figure pas. De plus, il confirme -- comme l'avait déjà fait Willems en avance sur son temps -- le `βάρβαρον' du vers 1051 qui laissait de nombreux éditeurs sceptiques. (Voir sa \nieme{26} note, très instructive sur ce passage et le rôle de l'étoile Sirius.)}%
\debutNote{1043-51}%
\note{1}{\italique{The Oxyrhynchus Papyri}, LXXIII (2009), ed. A. Benaissa.}}
\end{groupe}%
\endnumbering
\end{Leftside}
\begin{Rightside}
\beginnumbering
%% v. 1001
\begin{groupe}%
\personnageFr{Archer Scythe}%
\trimIambFr{Ici maintenant, tu vas crier, en plein air.\\\passeLigne}%
\end{groupe}%
%% v. 1002a
\begin{groupe}%
\personnageFr{Mnésiloque}%
\trimIambFr{Archer, je t'en supplie...\\\passeLigne}%
\end{groupe}%
%% v. 1002b
\begin{groupe}%
\personnageFr{Archer}%
\trimIambFr{\ligne{1002}\espace{3.1cm}Ne me supplie pas.\\\passeLigne}%
\end{groupe}%
%% v. 1003a
\begin{groupe}%
\personnageFr{Mnésiloque}%
\trimIambFr{Desserre la cheville.\\\passeLigne}%
\end{groupe}%
%% v. 1003b
\begin{groupe}%
\personnageFr{Archer}%
\trimIambFr{\ligne{1003}\espace{2.9cm}C'est ce que je vais faire\\\passeLigne}%
\end{groupe}%
%% v. 1004
\begin{groupe}%
\personnageFr{Mnésiloque}%
\trimIambFr{Aïe, misère! tu la pousses davantage.\\\passeLigne}%
\end{groupe}%
%% v. 1005a
\begin{groupe}%
\personnageFr{Archer}%
\trimIambFr{Encore plus, veux-tu?\\\passeLigne}%
\end{groupe}%
%% v. 1005b-1006a
\begin{groupe}%
\personnageFr{Mnésiloque}%
\trimIambFr{\ligne{1005}\espace{2.7cm}Aie, aïe, aïe, aie, aïe, aie !\\
la peste soit de toi.\\\passeLigne}%
\end{groupe}%
%% v. 1006b-1007
\begin{groupe}%
\personnageFr{Archer}%
\trimIambFr{\ligne{1006}\espace{2.5cm}Silence, misérable vieillard.\\
Allons, je vais, moi, apporter une natte, pour faire ma garde. (Il sort.)\\\passeLigne}%
\end{groupe}%
%% v. 1008-1014
\begin{groupe}%
\personnageFr{Mnésiloque}%
\trimIambFr{Ces douceurs-là j'en ai l'obligation à Euripide.\\
(Il aperçoit dans la coulisse Euripide costumé en Persée.) Hein ! dieux ! Zeus sauveur ! il y a espoir.
L'homme ne semble pas vouloir m'abandonner. Mais il m'a\\
donné à entendre, en apparaissant en Persée,\\
qu'il me faut devenir Andromède. Aussi bien,\\
les chaînes, j'en suis muni. Il est donc clair\\
qu'il viendra me délivrer. Autrement il n'aurait pas dirigé son vol de ce côté.\\\passeLigne}%
\end{groupe}%
%% v. 1015-1055
\begin{groupe}%
\personnageFr{Mnésiloque}%
\paratragiquesFr{\ligne{1015}Vierges, vierges amies, puissé-je}
\end{groupe}%
%% v. 1015-1055
\begin{groupe}%
\paratragiquesFr{partir d'ici et}
\end{groupe}%
%% v. 1015-1055
\begin{groupe}%
\paratragiquesFr{tromper le Scythe !}
\end{groupe}%
%% v. 1015-1055
\begin{groupe}%
\paratragiquesFr{M'entends-tu, ô toi qui en des antres répètes mes chants ?}
\end{groupe}%
%% v. 1015-1055
\begin{groupe}%
\paratragiquesFr{acquiesce, souffre que}
\end{groupe}%
%% v. 1015-1055
\begin{groupe}%
\paratragiquesFr{j'aille retrouver ma femme.}
\end{groupe}%
%% v. 1015-1055
\begin{groupe}%
\paratragiquesFr{Sans merci est celui qui m'enchaîna,}
\end{groupe}%
%% v. 1015-1055
\begin{groupe}%
\paratragiquesFr{des mortels le plus éprouvé.}
\end{groupe}%
%% v. 1015-1055
\begin{groupe}%
\paratragiquesFr{Échappé à grand'peine à une vieille}
\end{groupe}%
%% v. 1015-1055
\begin{groupe}%
\paratragiquesFr{guenipe, je ne laissai pas d'être perdu.}
\end{groupe}%
%% v. 1015-1055
\begin{groupe}%
\paratragiquesFr{Car depuis un temps ce Scythe se tient en faction : il m'a suspendu, brisé, sans amis,}
\end{groupe}%
%% v. 1015-1055
\begin{groupe}%
\paratragiquesFr{en pâture aux corbeaux. Vois-tu, ce n'est pas}
\end{groupe}%
%% v. 1015-1055
\begin{groupe}%
\paratragiquesFr{pour des chœurs de danse}
\end{groupe}%
%% v. 1015-1055
\begin{groupe}%
\paratragiquesFr{ni parmi des jeunes filles de mon âge}
\end{groupe}%
%% v. 1015-1055
\begin{groupe}%
\paratragiquesFr{que je suis ici debout avec cornet et jeton.%\commentaire{1029}{κημὸν ἕστηκ᾽ ἔχουσ᾽ }{A comparer avec le bas-relief de la galerie Gustiniani et le vase d'argent du musée Corsini, représentant le jugement d'Oreste devant l'Aréopage, et où une jeune fille se tient en effet debout derrière l'urne (SAGLIO, I, p.399). Était-ce l'usage reçu à l'Aréopage?}.
}
\end{groupe}%
%% v. 1015-1055
\begin{groupe}%
\paratragiquesFr{Mais enchaînée en d'étroits liens,}
\end{groupe}%
%% v. 1015-1055
\begin{groupe}%
\paratragiquesFr{je me vois exposée à la voracité du monstre Glaucétès.%\commentaire{1031}{Γλαυκέτῃ}{Glaucétès l'omnivore, le glouton, figure ici le monstre marin qui doit dévorer Andromède.}
}
\end{groupe}%
%% v. 1015-1055
\begin{groupe}%
\paratragiquesFr{Lamentez sur moi un péan, non d'hymen, mais de chaînes,}
\end{groupe}%
%% v. 1015-1055
\begin{groupe}%
\paratragiquesFr{ô femmes, car,}
\end{groupe}%
%% v. 1015-1055
\begin{groupe}%
\paratragiquesFr{pitoyable, j'ai pitoyablement souffert.}
\end{groupe}%
%% v. 1015-1055
\begin{groupe}%
\paratragiquesFr{Malheureux que je suis, malheureux !}
\end{groupe}%
%% v. 1015-1055
\begin{groupe}%
\paratragiquesFr{Et entre d'autres tortures impies infligées par mes proches,}
\end{groupe}%
%% v. 1015-1055
\begin{groupe}%
\paratragiquesFr{implorant l'homme}
\end{groupe}%
%% v. 1015-1055
\begin{groupe}%
\paratragiquesFr{d'une lamentation pleine de larmes de l'obscur Hadès.}
\end{groupe}%
%% v. 1015-1055
\begin{groupe}%
\paratragiquesFr{ah ah! ah ah!}
\end{groupe}%
%% v. 1015-1055
\begin{groupe}%
\paratragiquesFr{l'homme qui me rasa d'abord,}
\end{groupe}%
\endnumbering
\end{Rightside}
\end{pages}
\Pages
\clearpage\ifodd\value{page}\hbox{}\newpage\fi
\enTeteG{ΕΠΕΙΣΟΔΙΟΝ Ε' (Ἠχώ) --- Codex Ravennas 429, f. 175r (Classense 177r) --- vv. 1044–1081}
\begin{adjustwidth}{-1.5cm}{-3.0cm}
\begin{adjustbox}{max totalheight=21.0cm, max width=16.5cm, keepaspectratio}
\begin{minipage}{16.5cm}
{\fontsize{7.2pt}{8.5pt}\selectfont
\noindent\textbf{Folio 26 --- Codex Ravennas 429, \texttt{f. 175r} (\texttt{Classense 177r}) : ΕΠΕΙΣΟΔΙΟΝ Ε' (Ἠχώ), vv. 1044–1081}\hfill \textit{Témoins :} $\text{R}^1, \Rpc, \Rdeux$, \PsiQuatre\par\vspace{2pt}
\hrule height 0.4pt\vspace{2.5pt}
}
{\fontsize{8.1pt}{9.5pt}\selectfont
\ligneRav{1044}{μ\auDessus{ν}{Η}}{\hspace{1.1em}ὃς ἐμὲ κροκόεντ᾽ ἐνέδυσεν·}{\textbf{1044} \lecon{κροκόεντ᾽} R : κροκόεν τόδ' Bergk \espApp ἐνέδυσεν R : ἀμφέδυσεν Hermann ap Σ}
\ligneRav{1045}{}{\hspace{1.1em}ἐπὶ δὲ τοῖσδε τόδ᾽ ἀνέπεμψεν}{}
\ligneRav{}{}{\hspace{1.1em}ἱερόν, ἔνθα γυναῖκες·}{}
\ligneRav{1047}{}{\hspace{1.1em}ἰώ μοι μοίρας}{\textbf{1047} [ι]\douteux{ω}μοι\douteux{ρ}[ας \PsiQuatre : ἰώ μοι μοίρας R : ἰώ μοίρας Hermann \espApp δαίμων R : δαῖμον Burges \espApp ἀνέτικτε R : \lecon{ἄτεγκτε} Ellebodius : ἅν ἔτικτε Casaubon\espApp δαίμων R : δαῖμον Burges \espApp ἀνέτικτε R : \lecon{ἄτεγκτε} Ellebodius : ἅν ἔτικτε Casaubon}
\ligneRav{}{}{\hspace{1.1em}ἄτεγκτε δαίμων·}{}
\ligneRav{}{}{\hspace{1.1em}ὦ κατάρατος ἐγώ·}{}
\ligneRav{}{}{\hspace{1.1em}τίς ἐμὸν οὐκ ἐπόψεται}{}
\ligneRav{}{}{\hspace{1.1em}πάθος ἀμέγαρτον ἐπὶ κακῶν παρουσίᾳ;}{}
\ligneRav{1050}{}{\hspace{1.1em}εἴθε με πυρφόρος αἰθέρος ἀστὴρ—}{}
\ligneRav{1051}{}{\hspace{1.1em}τὸν \difference{βάρβαρον} ἐξολέσειεν·}{\textbf{1051} \lecon{βάρβαρον} \PsiQuatre{}, R, Willems}
\ligneRav{1052}{}{\hspace{1.1em}οὐ γὰρ ἔτ᾽ ἀθανάταν φλόγα \difference{λεύσσειν}}{\textbf{1052} λεύσειν R : \lecon{λεύσσειν} Biset}
\ligneRav{}{}{\hspace{1.1em}ἐστὶν ἐμοὶ φίλον, ὡς ἐκρεμάσθην,}{}
\ligneRav{1054}{}{\hspace{1.1em}λαιμότμητ᾽ ἄχη δία μένων αἰόλαν}{\textbf{1054} δαιμόνων R, Σ : δαιμονῶν Brunck : δαιμόνι᾽ amicus Fritzschii : \lecon{διὰ μένων} Willems}
\ligneRav{1055}{}{\hspace{1.1em}νέκυσιν ἐπὶ πορείαν·}{}
\ligneRav{1056}{ευρ\auDessus{ι}{π} ηχω}{χαῖρ᾽ ὦ φίλη παῖ· τὸν δὲ πατέρα Κηφέα}{\textbf{1056} ευρ\auDessus{ι}{π} ηχω in mg. \Rpc, \Mon}
\ligneRav{}{}{ὅς σ᾽ ἐξέθηκεν ἀπολέσειαν οἱ θεοί·}{}
\ligneRav{1058}{μ\auDessus{ν}{Η}}{σὺ δ᾽ εἶ τίς ἥτις τοὐμὸν \difference{ᾤκτιρας} \difference{πάθος};}{\textbf{1058} ὤκτειρας R : \lecon{ᾤκτιρας} van Herwerden \espApp πάθος R : \lecon{πάθος;} Biset}
\ligneRav{}{\auDessus{εὐ}{ρ}}{Ἠχὼ, λόγων ἀντῳδὸς ἐπικοκκάστρια,}{}
\ligneRav{1060}{}{ἥπερ πέρυσιν ἐν τῷδε \difference{ταὐτῷ} χωρίῳ}{\textbf{1060} ταυτῷ R : \lecon{ταὐτῷ} Dindorf (1830)}
\ligneRav{1061}{}{Εὐριπίδῃ \difference{καὐτὴ} ξυνηγωνιζόμην·}{\textbf{1061} καὶ αὐτὴ R : \lecon{καὐτὴ} Biset}
\ligneRav{1062}{}{ἀλλ᾽ ὦ τέκνον σὲ μὲν \difference{τὸ σαυτῆς} χρὴ ποιεῖν,}{\textbf{1062} τοσαύτης R : \lecon{τὸ σαυτῆς} Dindorf (1839)}
\ligneRav{1063}{—}{\difference{κλαίειν} ἐλεινῶς. : σὲ δ᾽ ἐπικλαίειν ὕστερον·}{\textbf{1063} \lecon{κλαίειν} et \lecon{ἐπικλαίειν} R : κλάειν et ἐπικλάειν Brunck}
\ligneRav{}{\auDessus{εὐ}{ρ}}{ἐμοὶ μελήσει ταῦτά γ᾽. ἀλλ᾽ ἄρχου λόγων·}{}
\ligneRav{1065}{ὡς ἀνδρομ.}{\hspace{1.1em}ὦ νὺξ ἱερὰ}{\textbf{1065} ὁ μνησίλοχος ὡς ἀνδρομέδα in mg. \Rpc, \Mon}
\ligneRav{}{}{\hspace{1.1em}ὡς μακρὸν ἵππευμα διώκεις}{}
\ligneRav{}{}{\hspace{1.1em}ἀστεροειδέα νῶτα διφρεύουσ᾽}{}
\ligneRav{}{}{\hspace{1.1em}αἰθέρος ἱερᾶς}{}
\ligneRav{1069}{—}{\hspace{1.1em}τοῦ σεμνοτάτου δι᾽ Ὀλύμπου; : †δι᾽ Ὀλύμπου;†}{\textbf{1069} \lecon{δι᾽ Ὀλύμπου} R : †δι᾽† ὀλυμπίας Bibi ap. Hesych. (ο 649)}
\ligneRav{1070}{μ\auDessus{ν}{Η}}{\hspace{1.1em}τί ποτ᾽ Ἀνδρομέδα \difference{περίαλλα} κακῶν}{\textbf{1070} περὶ ἄλλα R : \lecon{περίαλλα} Ellebodius}
\ligneRav{1071}{—}{\hspace{1.1em}μέρος ἐξέλαχον; : \difference{μόρος ἔξελ᾽ ἄχος}·}{\textbf{1071} μέρος ἐξέλαχον; R : \lecon{μόρος ἔξελ᾽ ἄχος.} Bibi}
\ligneRav{1072}{μ\auDessus{ν}{Η}}{\hspace{1.1em}θανάτου τλήμων : †θανάτου τλήμων†·}{\textbf{1072} θανάτου τλήμων R : θαλάττης τέμνον Bibi}
\ligneRav{1073}{μ\auDessus{ν}{Η}}{\hspace{1.1em}ἀπολεῖς μ᾽ ὦ γραῦ \difference{στωμυλλομένη}·}{\textbf{1073} στωμυλλωμένη R : \lecon{στωμυλλομένη} Grynaeus}
\ligneRav{1074}{\auDessus{εὐ}{ρ}}{\hspace{1.1em}\difference{στωμυλιομενής}·}{\textbf{1074} στωμυλωμένη R : \lecon{στωμυλιομενής} Bibi}
\ligneRav{1075}{μ\auDessus{ν}{Η}}{\hspace{1.1em}νὴ Δί᾽, ὀχληρά γ᾽ εἰσήρρηκας}{}
\ligneRav{}{—}{\hspace{1.1em}λίαν. : †λίαν†·}{}
\ligneRav{}{μ\auDessus{ν}{Η}}{\hspace{1.1em}ὦγάθ᾽ ἔασόν με μονῳδῆσαι,}{}
\ligneRav{1078}{—}{\hspace{1.1em}καὶ χαριεῖ μοι παῦσαι. : ᾽παΰσαι·}{\textbf{1078} Μνη. χαριεῖ μοι παῦσαι --- Εὐρ. παῦσαι R : \lecon{Μνη. χαριεῖ μοι ᾽παΰσαι --- simul Μνη. et Εὐρ. παῦσαι} Bibi\espApp (Εὐρ.) παῦσαι R : \lecon{(Εὐρ.) ᾽παΰσαι} Bibi}
\ligneRav{1079}{μ\auDessus{ν}{Η}}{\hspace{1.1em}βάλλ᾽ ἐς κόρακας. : \difference{βάλλε σύ κόρακας}·}{\textbf{1079} βάλλ᾽ ἐς κόρακας R : \lecon{βάλλε σύ κόρακας} Bibi}
\ligneRav{1080}{μ\auDessus{ν}{Η}}{\hspace{1.1em}†\difference{τί τὸ κακόν}†; : †τί τὸ κακόν;† : ληρεῖς. : †ληρεῖς†·}{\textbf{1080} \lecon{τί τὸ κακόν} bis R : τί κακόν bis Bentley}
\ligneRav{}{μ\auDessus{ν}{Η}}{\hspace{1.1em}οἴμωζ᾽. : †οἴμωζ᾽†·}{}
}
\vspace{3pt}\noindent\rule{4.5cm}{0.4pt}\par\vspace{1.5pt}
{\fontsize{7.0pt}{8.3pt}\selectfont
\noindent\textbf{Loci paralleli (vv. 1044–1081) ---} \textbf{1049} (παρουσίᾳ) Moeris π 14 \espTest \textbf{1059} Suidas η 686 (ἐπικοκκάστρια codd. GM\expo{a.c.} vel ἐπικοκάστρια cod. I vel ἐπικαγχάστρια cod. A vel ἐπικαχάστρια codd. FV\expo{mg.} vel ἐπικοστράτ textu V) \espTest (ἀντῳδὸς ἐπικοκκάστρια) \Eustathe 1761.26 ἐκτίθεται καινοφώνους λέξεις οἷον [ ... ] καὶ τὸ ἐπικοκκάζειν ὕβρει ἤ κατά τινας ἐπιμηκάζειν, ὃθεν παρὰ Ἀριστοφάνει τὸ `ἀοιδὸς ἐπικοκκάστρια' \espTest \textbf{1066-8} (ὡς μακρὸν -- αἰθέρος) Suidas ι 531 $\sim$ Zonaras 1117 \espTest \textbf{1067} (νῶτα) Moeris ν 6 \espTest \textbf{1067-8} Suidas αι 122 $\sim$ Zonaras 67, 78 (sed ἀστρο- pro ἀστερο-)\par
}
\vspace{2.5pt}\noindent\rule{4.5cm}{0.4pt}\par\vspace{1.5pt}
{\fontsize{6.9pt}{8.2pt}\selectfont
\noindent\textbf{Suite (en commentaire) de l'apparat critique \& notes paléographiques (\texttt{f. 175r}) ---} \textbf{Témoin papyrologique :} \textbf{\PsiQuatre{} (\italique{P.Oxy.} LXXIII 4935, fr. 1)} couvre les \textbf{vv. 1043–1051}. --- \textbf{Deux indications scéniques marginales de \Rpc{} (copiées par \Mon{}) :} --- \textbf{Apport du papyrus \PsiQuatre{} (vv. 1047 et 1051) :} Au v. 1047, \PsiQuatre{} lit [ι]ω μοιρ[ας sans la dittographie μοι de R (ἰώ μοι μοίρας) ; au v. 1051, \PsiQuatre{} confirme βάρβαρον de R (défendu par Alphonse Willems et L. Hauseux). --- \textbf{Cinq restitutions et conjectures de \textbf{Bibi} dans le dialogue avec Écho (vv. 1069, 1071, 1072, 1074, 1078, 1079) :} --- \textbf{v. 1044 :} \lecon{κροκόεντ'} R (κροκόεν τόδ' Bergk) \espApp ἐνέδυσεν R : \lecon{ἀμφέδυσεν} Hermann ap. $\Sigma$. --- \textbf{v. 1047 :} [ι]ω μοιρ[ας \PsiQuatre{} : ἰώ μοι μοίρας R : \lecon{ἰὼ μοίρας} Hermann \espApp δαίμων R : \lecon{δαῖμον} Burges \espApp ἀνέτικτε R (médivision et épenthèse de $\nu$) : \lecon{ἄτεγκτε} Ellebodius (ἅν ἔτικτε Casaubon). --- \textbf{v. 1051 :} \lecon{βάρβαρον} \PsiQuatre{}, R, Willems. --- \textbf{v. 1054 :} δαιμόνων R, $\Sigma$ (δαιμονῶν Brunck ; δαιμόνι' amicus Fritzschii) : \lecon{διὰ μένων} Alphonse Willems (adopté par L. Hauseux). --- \textbf{v. 1056 :} \lecon{ευρ\auDessus{ι}{π} ηχω} in mg. $\text{R}^{\text{pc}}, \text{M}$. --- \textbf{v. 1065 :} \lecon{ὁ μνησίλοχος ὡς ἀνδρομέδα} in mg. $\text{R}^{\text{pc}}, \text{M}$. --- \textbf{v. 1069 :} \lecon{δι' Ὀλύμπου} R (†δι'† ὀλυμπίας \lecon{Bibi} ap. Hésychius $\text{ο}\ 649$). --- \textbf{v. 1071 :} μέρος ἐξέλαχον; R : \lecon{μόρος ἔξελ' ἄχος.} \lecon{Bibi}. --- \textbf{v. 1072 :} \lecon{θανάτου τλήμων} R (θαλάττης τέμνον \lecon{Bibi} conj. in app.). --- \textbf{v. 1074 :} στωμυλωμένη R : \lecon{στωμυλιομενής} \lecon{Bibi}. --- \textbf{v. 1078 :} (Εὐρ.) παῦσαι R : \lecon{(Εὐρ.) 'παΰσαι} \lecon{Bibi}. --- \textbf{v. 1079 :} βάλλ' ἐς κόρακας R : \lecon{βάλλε σύ κόρακας} \lecon{Bibi}.\par
}
\end{minipage}
\end{adjustbox}
\end{adjustwidth}
\newpage
\enTeteD{Codex Ravennas 429 (Biblioteca Classense) --- Fac-similé du f. 175r (vv. 1044–1081)}
\begin{adjustwidth}{-1.5cm}{-3.0cm}
\begin{center}
\includegraphics[width=16.5cm,height=20.5cm,keepaspectratio]{/workspaces/Thesmophories/dossier_ravennas_thesmophories/folios_individuels_ff162v-177r/26_f175r_classense177r_vv1044-1081.jpg}\par
\vspace{2pt}
{\scriptsize\textbf{Codex Ravennas 429, f. 175r (Classense 177r)} --- ΕΠΕΙΣΟΔΙΟΝ Ε' (Ἠχώ), \textbf{vv. 1044–1081} ($\text{R}^1, \Rpc, \Rdeux$, \PsiQuatre).}
\end{center}
\end{adjustwidth}
\newpage
\clearpage\ifodd\value{page}\hbox{}\newpage\fi
\enTeteG{ΘΕΣΜΟΦΟΡΙΑΖΟΥΣΑΙ --- Texte grec établi (f. 175r, vv. 1044–1081)}
\enTeteD{LES THESMOPHORIES --- Traduction française (f. 175r, vv. 1044–1081)}
\begin{pages}
\begin{Leftside}
\beginnumbering
%% v. 1015-1055
\begin{groupe}%
\typeVers{}\paratragiques{\ligne{1044}ὃς ἐμὲ κροκόεντ᾽ ἐνέδυσεν·\apparat{1044}{\lecon{κροκόεντ᾽} R : κροκόεν τόδ' Bergk \espApp ἐνέδυσεν R : ἀμφέδυσεν Hermann ap Σ}%
%\commentaire{1044}{}{Phéniciennes d'Euripide 1491 : 
%  Ανtigone.                 οὐ προκαλυπτομένα βοτρυχώδεος   (1485)
%ἁβρὰ παρηίδος οὐδ’ ὑπὸ παρθενί-    (1486-1487)
%ας τὸν ὑπὸ βλεφάροις φοίνικ’, ἐρύθημα προσώπου,   (1488)
%αἰδομένα φέρομαι βάκχα νεκύ-
%ων, κράδεμνα δικοῦσα κόμας ἀπ’ ἐ-    (1490)
%μᾶς, στολίδος κροκόεσσαν ἀνεῖσα τρυφάν,
%ἁγεμόνευμα νεκροῖσι πολύστονον. αἰαῖ, ἰώ μοι.
%ὦ Πολύνεικες, ἔφυς ἄρ’ ἐπώνυμος· ὤμοι μοι, Θῆβαι·   (1493-1494)
%σὰ δ’ ἔρις—οὐκ ἔρις, ἀλλὰ φόνῳ φόνος—    (1495)
%Οἰδιπόδα δόμον ὤλεσε κρανθεῖσ’
%αἵματι δεινῷ, αἵματι λυγρῷ.
%τίνα προσῳδὸν
%ἢ τίνα μουσοπόλον στοναχὰν ἐπὶ
%δάκρυσι δάκρυσιν, ὦ δόμος, ὦ δόμος,   (1500)
%ἀγκαλέσωμαι,
%τρισσὰ φέρουσα τάδ’ αἵματα σύγγονα,}
}
\end{groupe}%
%% v. 1015-1055
\begin{groupe}%
\typeVers{dim. troch.}\paratragiques{ἐπὶ δὲ τοῖσδε τόδ᾽ ἀνέπεμψεν}
\end{groupe}%
%% v. 1015-1055
\begin{groupe}%
\typeVers{pherecr.}\paratragiques{ἱερόν, ἔνθα γυναῖκες.}
\end{groupe}%
%% v. 1015-1055
\begin{groupe}%
\typeVers{dochm.}\paratragiques{ἰώ μοι μοίρας\apparat{1047}{[ι]\douteux{ω}μοι\douteux{ρ}[ας \PsiQuatre : ἰώ μοι μοίρας R : ἰώ μοίρας Hermann \espApp δαίμων R : δαῖμον Burges \espApp ἀνέτικτε R : \lecon{ἄτεγκτε} Ellebodius : ἅν ἔτικτε Casaubon}%
%\commentaire{1047}{ἰώ μοι μοίρας.}{\italique{Andromaque}, vers 824 : \citComm{ἰώ μοί μοι.}}
}
\end{groupe}%
%% v. 1015-1055
\begin{groupe}%
\typeVers{ithyph.}\paratragiques{ἄτεγκτε δαίμων·%\apparat{1047}{δαίμων R : δαῖμον Burges \espApp ἀνέτικτε R : \lecon{ἄτεγκτε} Ellebodius : ἅν ἔτικτε Casaubon}
}
\end{groupe}%
%% v. 1015-1055
\begin{groupe}%
\typeVers{hem.}\paratragiques{ὦ κατάρατος ἐγώ·}
\end{groupe}%
%% v. 1015-1055
\begin{groupe}%
\typeVers{lecyth.}\paratragiques{\ligne{1048}τίς ἐμὸν οὐκ ἐπόψεται}
\end{groupe}%
%% v. 1015-1055
\begin{groupe}%
\typeVers{trim. iamb.}\paratragiques{\ligne{1049}πάθος ἀμέγαρτον ἐπὶ κακῶν παρουσίᾳ;\testimonia{1049}{(παρουσίᾳ) Moeris π 14}}
\end{groupe}%
%% v. 1015-1055
\begin{groupe}%
\typeVers{$4$ dactyl.}\paratragiques{εἴθε με πυρφόρος αἰθέρος ἀστὴρ—}
\end{groupe}%
%% v. 1015-1055
\begin{groupe}%
\typeVers{paroem.}\paratragiques{τὸν βάρβαρον ἐξολέσειεν.}
\end{groupe}%
%% v. 1015-1055
\begin{groupe}%
\typeVers{$4$ dactyl.}\paratragiques{οὐ γὰρ ἔτ᾽ ἀθανάταν φλόγα λεύσσειν\apparat{1052}{λεύσειν R : \lecon{λεύσσειν} Biset}}
\end{groupe}%
%% v. 1015-1055
\begin{groupe}%
\typeVers{$+$ $4$ dact.}\paratragiques{ἐστὶν ἐμοὶ φίλον, ὡς ἐκρεμάσθην,}
\end{groupe}%
%% v. 1015-1055
\begin{groupe}%
%\typeVers{hypod. $+$ doch.}
\paratragiques{λαιμότμητ᾽ ἄχη δία μένων αἰόλαν\apparat{1054}{δαιμόνων R, Σ : δαιμονῶν Brunck : δαιμόνι᾽ amicus Fritzschii : \lecon{διὰ μένων} Willems}%
\commentaire{1054-5}{λαιμότμητ᾽ ἄχη δαιμόνων αἰόλαν / νέκυσιν ἐπὶ πορείαν.}{%ὡς ἐκρεμάσθην = depuis que j'ai été attaché ? 
Willems a consacré une note à ce passage\marqueur{1}. Si les vues qu'il y développe sont ingénieuses, elle n'emporte pas pour autant ma conviction. %D'abord, la traduction λαιμότμητ᾽ ἄχη reporter à la  note de Willems, reproduite en annexe --, le texte en est corrompu, certainement en cet étonnant †δαιμόνων†. 
Willems propose de construire ἄχη δία μένων là où le \italique{Ravennas} porte ἄχη †δαιμόνων†. À défaut d'autres propositions plus convaincantes, c'est néanmoins la correction que je retiendrai.%\marqueur{2} ; je préfère y voir le verbe ἀδημονέω (ἀδαιμονέω DGE)
%\eff{Retoucher ! MAsculin adhmonwn ? ὡς marque l'exclamation ?} 
}%
\debutNote{1054-5}%
\note{1}{\nieme{27} note aux \italique{Thesmophories} d'Alphonse Willems, reproduite en annexe.}%
%\note{2}{Si le sens est tout-à-fait satisfait par la conjecture de Willems ; qu'il se justifie habilement de la construction peu habituelle; ; ma réserve vient de ce que je m'explique difficilement le passage du epsilon de μένων au omicron de δαιμόνων. Tandis que la transformation du êta en le digramme ΑΙ est beaucoup plus naturelle en supposant l'étape intermédiaire d'un epsilon : Η $\xrightarrow{}$ Ε  $\xrightarrow{prononciat. ident.}$ ΑΙ.}%
}%
\end{groupe}%
%% v. 1015-1055
\begin{groupe}%
\typeVers{ithyphall.}\paratragiques{νέκυσιν ἐπὶ πορείαν.\\\passeLigne}%
\end{groupe}%
%% v. 1056-1057
\begin{groupe}%
\personnageGrec{Εὐριπίδης}%
\trimIamb{\ligne{1056}χαῖρ᾽ ὦ φίλη παῖ· τὸν δὲ πατέρα Κηφέα\apparat{1056}{ευρ\auDessus{ι}{π} ηχω in mg. \Rpc, \Mon}\\
ὅς σ᾽ ἐξέθηκεν ἀπολέσειαν οἱ θεοί.\\\passeLigne}%
\end{groupe}%
%% v. 1058
\begin{groupe}%
\personnageGrec{Μνησίλοχος}%
\trimIamb{σὺ δ᾽ εἶ τίς ἥτις τοὐμὸν ᾤκτιρας πάθος;\apparat{1058}{ὤκτειρας R : \lecon{ᾤκτιρας} van Herwerden \espApp πάθος R : \lecon{πάθος;} Biset}\\\passeLigne}%
\end{groupe}%
%% v. 1059-1063a
\begin{groupe}%
\personnageGrec{Εὐριπίδης}%
\trimIamb{Ἠχὼ, λόγων ἀντῳδὸς ἐπικοκκάστρια,\testimonia{1059}{Suidas η 686 (ἐπικοκκάστρια codd. GM\expo{a.c.} vel ἐπικοκάστρια cod. I vel ἐπικαγχάστρια cod. A vel ἐπικαχάστρια codd. FV\expo{mg.} vel ἐπικοστράτ textu V) \espTest (ἀντῳδὸς ἐπικοκκάστρια) \Eustathe 1761.26 ἐκτίθεται καινοφώνους λέξεις οἷον [ ... ] καὶ τὸ ἐπικοκκάζειν ὕβρει ἤ κατά τινας ἐπιμηκάζειν, ὃθεν παρὰ Ἀριστοφάνει τὸ `ἀοιδὸς ἐπικοκκάστρια'}%
\commentaire{1059}{ἀντῳδὸς ἐπικοκκάστρια.}{Ceci rappelle la jolie épigramme d'Archias (\italique{Anthol. Palat.}, VII, 191) sur une pie (κίσσα\marqueur{1})%
\citComm{qui souvent, pareille à l'Écho, faisait entendre, d'un bec prompt à répliquer, une équisonnance bruyante et moqueuse.\\
πολλάκι δὲ κρέξασα πολύθροον, οἷά τις ἀχώ,\\
κέρτομον ἀντῳδοῖς χείλεσιν ἁρμονίαν.}%
(Note de Willems.)}%
\debutNote{1059}%
\note{1}{Si la pie est un oiseau capable de parler, selon Descartes, l'oiseau dont il question dans cette épigramme serait à en croire le \LSJ \ un corvidé cousin, de la famille des geais : un \italique{geai des chênes}, \latin{Garrulus glandarius}. (NDÉ.)}%
\commentaire{1059}{Ἠχὼ λόγων ἀντῳδὸς ἐπικοκκάστρια.}{Seule compagnie d'Andromède : la nymphe Écho, cachée dans un antre non loin, et qui répète les paroles. L'occasion d'un détournement comique de ces échos était trop belle pour qu'Aristophane n'en profitât pas !
Euripide – qui joue dans ce passage la nymphe Écho – répond effectivement à l'appel de l'Andromède postiche qu'essaye d'interpréter le peu crédible vieillard Mnésiloque. 
Dans la tragédie éponyme d'Euripide d'antan, l'acteur d'Écho ne se montrait pas sur scène mais renvoyait les répliques depuis un antre à proximité de la falaise où était enchaînée Andromède. %\eff{sûr ? cf. représentation}. 
Ajoutez à cela la μηχανή permettant une entrée aérienne de Persée et vous comprendrez quelque peu l'impression forte et durable que la pièce suscita auprès des spectateurs.\\
Revenons au vers 1059 de notre comédie : Ἠχὼ, λόγων ἀντῳδὸς ἐπικοκκάστρια.\\
On a passé un peu vite sur l'adjectif ἐπικοκκάστρια ; si le sens de λόγων ἀντῳδὸς est clair : « la rediseuse en chantant des paroles », qualificatifs prévisibles de la nymphe Écho (notez tout-de-même que les échos étaient probablement renvoyés en \italique{chantant}, ou au moins en \italique{persiflant}), Aristophane ajoute l'hapax ἐπικοκκάστρια. 
Scrutons ce mot : ἐπι-κοκκάσ-τρια ; le préfixe ἐπι- apporte selon nous sa nuance d'hostilité, ici de causticité %(en fait, les autres mots composés sur ἐπι- jouent avec les différentes nuances du préfixe qui peut marquer la succession temporelle ainsi que l'hostilité) 
; le suffixe -τρια est courant, principalement attesté chez les comiques et en prose\marqueur{1} ; reste κοκκάς, qu'une oreille française ne peut s'empêcher de rapprocher de notre `cocasse', risible mais avec une certaine violence comme le coq dont il tire son origine… 
Si le volatile d'Aristophane n'est pas un coq ; κοκκάς renvoie au κόκκυξ ou \italique{coucou}, les Grecs ne prêtaient pas moins à cet oiseau de malheur une méchanceté bien supérieure. \italique{Confer} le magistrale note de Willems sur le \nieme{507} vers des \italique{Oiseaux} et les propos salés (\latin{petulantiæ sales}) qu'on jetait aux laboureurs en imitant le cri du coucou. (Soit dit en passant, on retrouve les mêmes traits dans feu nos traditions orales : le coucou chantait les événements fâcheux (dont en premier lieu, bien-sûr, le cocufiage)\marqueur{2}.\\
Le scholiaste nous est ici précieux\marqueur{3} : il définit ainsi ἐπικοκκάστρια : εἰωθυῖα γελᾶν, γελάστρια, « ayant l'habitude de rire, moqueuse ». Aristophane nous prévient donc que les échos auront quelque chose de piquant et que si on les trouve d'une « \anglais{annoying loquacity} »\marqueur{4} , c'est que, peut-être, le sel attique nous en a échappé.\\
Une autre scholie nous permet de préciser un peu la teneur des échos :%
\citComm{ἐπεὶ εἰσήγαγε κακοστένακτον τὴν Ἡχὼ ὁ Εὐριπίδης ἐν τῇ Ανδρομέδᾳ, εἰς τοῦτο παίζει [...]}%
La bonne interprétation de cette scholie réside dans le sens à accorder à l'adjectif κακοστένακτος : \italique{à la méchante lamentation} ?\\
Il est étonnant que les modernes éditions n'aient pas cherché à percer ce que les répliques d'Écho-Euripide pouvaient avoir d'amusantes ou de méchantes… Pourtant, elles disposaient déjà d'un exemple de détournement par la « rediseuse goguenarde » : à la fin de la scène, au vers 1094, alors que l'archer Scythe tente d'attraper Écho-Euripide et que ce dernier s'enfuit sans demander son reste, une scholie nous indiquait clairement ce que disait en son jargon l'archer Scythe :%
\citComm{ἀντὶ τοῦ `οὐ χαιρήσεις'\\
À la place de : « tu ne vas pas te réjouir ».}%
On peut déduire de cette scholie que le Scythe – qui ne tient aucun compte dans son parler des aspirations – avait menacé Écho-Euripide d'un : `οὐ καιρήσεις'. Or, le manuscrit de Ravenne comportait effectivement la réponse ἐπικοκκάστρια d'Euripide-Écho\marqueur{5} : `οὐκ᾽ αἰρήσεις', c'est-à-dire en bon attique : `οὐχ αἱρήσεις', « tu ne m'attraperas pas ! ». Le cas était simple mais il aura fallu tout de même des siècles de philologiques hésitations entre les répartitions : Τοξότης (l'Archer) : « οὐ καιρήσεις » — Εὐριπίδης : « οὐ καιρήσεις » et Τοξότης : « οὐκ αἰρήσεις » — Εὐριπίδης : « οὐκ αἰρήσεις » avant que le grand Wilamowitz\marqueur{6} ne s'avise de la solution conciliatrice, la seule sensée et la seule comique : Τοξότης : « οὐ καιρήσεις » — Εὐριπίδης : « οὐχ αἱρήσεις ».\\
Bref ! les répétitions de la nymphe servent à Aristophane pour son jeu de massacre de la tragédie tant appréciée d'Euripide\marqueur{7} : par de légers détournements phonétiques (facilités par les barbarismes de l'archer Scythe), Écho transforme ce qui devait être une scène d'émotion grave en une foire d'empoigne générale.\\
C'est -- reconnaissons-le -- un exercice des plus périlleux que de s'aventurer dans des tentatives de reconstruction ; laisser ainsi libre cours à l'imagination peut exposer à de rudes retours de bâton ; pour Aristophane, nous devons cependant impérieusement nous y livrer puisqu'il vient de nous prévenir que les échos seraient \italique{moqueurs} ! Avec une réserve cependant : lorsque l'ἐπικοκκάστρια d'une réplique nous est demeurée celée, et que malgré nos efforts les seules propositions que nous avions à présenter n'étaient pas même dignes d'intégrer l'apparat critique -- le cas n'est pas isolé --, nous lui avons apposé la \italique{croix du désespoir} †.. Charge à d'autres de compléter notre ébauche ! je pense, pour ma part, ne pas avoir démérité d'Aristophane.}%
\debutNote{1059}%
\note{1}{Cf. Chantraine pour une longue liste des mots formés sur le même modèle, \italique{La formation des noms en grec ancien}, p. 106-7.}%
\note{2}{Pour s'en convaincre, il n'y a qu'à aller feuilleter le \italique{Répertoire} Coirault \italique{des chansons françaises de tradition orale}.}%
\note{3}{%Une fois n'est pas coutume ! ne boudons pas notre plaisir : pour une fois que dans pièces, si pauvre en scholies de qualité par rapport aux autres pièces d'Aristophane, les commentaires relèvent d'une érudition tout alexandrine ; connaissance lexicographique, donc, d'une part ; scholies qui font état parfois état d'un texte mieux conservé que celui qui nous est parvenu d'autre part ; enfin, possibilité pour le commentateur d'aller comparer avec l'original d'Euripide.
Les glosateurs avaient sans doute un texte des \italique{Thesmophories} en meilleur état que le nôtre (ce passage est le moins bien conservé de toute la pièce ; et pour cause : le comique a très tôt échappé aux scribes qui le recopiaient). De plus, ils avaient un avantage indéniable sur nous : ils pouvaient aller comparer la parodie avec l'original d'Euripide.}%
\note{4}{Austin-Olson… et pis encore (p. LXX) : « \anglais{much of the effect of the parodies of Helen and Andromeda depends on the fact that the tragic heroines are played by an individual wearing a prominent stagephallus} » (!!!) Mais pour qui prennent-ils Aristophane ?! le plus gras des vaudevillistes serait pris de scrupules devant les procédés que ces brillants professeurs prêtent à Aristophane, le plus fin des Attiques ! Il faudrait leur expliquer qu'il y a un temps pour tout, et qu'un pastiche d'euripidéenne tragédie n'est pas une scène burlesque de déshabillage ou encore une pièce sur la grève du sexe ; c'est calomnier Aristophane que de croire – et plus grave : faire croire à ses lecteurs, circonstance aggravante lorsque l'on prétend à être l'édition de référence des \italique{Thesmophories} ! – qu'il ne soit capable que de jouer d'un seul registre. Au reste, cette opinion est un expédient bien pratique qui dispense de se pencher sur ce que signifie véritablement le texte… ou de reconnaître honnêtement qu'on ne le comprend pas. Enfin, le véritable Aristophane nous dit lui-même ce qu'il pense du recours abusif à ces « phallus de cuir, épais, rouges par le bout, pour donner à rire aux gamins », c'est au vers 538 des \italique{Nuées} (voir aussi la \italique{Paix}, v. 739-753) ; que les futurs exégètes d'Aristophane ruminent bien ce passage avant de prétendre que notre poète concoctait à toutes les sauces des scènes de sexes sémillants. Et qu'ils lisent aussi l'article que Willems y avait consacré, « Le nu dans la Comédie ancienne », t. III, p. 381-400 (reproduit dans les annexes) : Aristophane recourait à l'\italique{olisbos} avec la plus grande parcimonie, et toujours pour de bonnes raisons.}%
\note{5}{En attribuant, il est vrai, la réplique à l'archer. Mais les indications de personnages sont de peu de valeur : les anciens parchemins ne comportaient pour la plupart que de légers signes (\italique{paragraphos} ou parfois simple indentation) pour indiquer le changement d'interlocuteur, et ce sans préciser quels étaient ces interlocuteurs ; les indications de personnages furent ajoutées par des scribes très postérieurement… et sont parfois assez folkloriques. (Pour celles du \italique{Ravennas}, quelques indications sont de la main d'Euphrosynus Boninus, celui qui compulsa le manuscrit pour la première édition imprimée des \italique{Thesmophories} -- la \italique{Juntina} (Venise, 1516) des apparats critiques. Von Velsen sut identifier les annotations de la main d'Euphrosynus Boninus sur le \italique{Ravennas} et ne les fit pas apparaître dans son apparat critique. Nous l'avons suivi en cela. Cf. Albert Martin, \italique{Les scolies du manuscrit d'Aristophane à Ravenne} (1882), p. V ; A. von Velsen, \italique{Über den Codex Urbinas der Lysistrata und der Thesmophoriazusen des Aristophanes} (1871).)}%
\note{6}{En simple note de ses \italique{Kleine Schriften}, IV (1962), p. 486. Mais ces écrits n'ont été publiés qu'après la mort du grand philologue. Si une partie des travaux modernes rétablissent correctement la menace du Scythe et son écho, il semble que le mérite en revienne plutôt à Sommerstein, qui eut la bonne intuition indépendamment de Wilamowitz.}%
\note{7}{Rappelons, par exemple, que c'est en relisant \italique{Andromède} que le Dionysos des \italique{Grenouilles} est pris d'un soudain désir d'Euripide.}\\
ἥπερ πέρυσιν ἐν τῷδε ταὐτῷ χωρίῳ\apparat{1060}{ταυτῷ R : \lecon{ταὐτῷ} Dindorf (1830)}%
\commentaire{1060}{πέρυσιν.}{L'\italique{Andromède} avait été produite l'année précédente, cf. la scholie :%
\citComm{ἐπεὶ πέρυσιν ἐδιδάχθη <ἡ> Ἀνδρομέδα.}
(Voir l'introduction sur la question de la date des \italique{Thesmophories}.)}\\
Εὐριπίδῃ καὐτὴ ξυνηγωνιζόμην.\apparat{1061}{καὶ αὐτὴ R : \lecon{καὐτὴ} Biset}\\
ἀλλ᾽ ὦ τέκνον σὲ μὲν τὸ σαυτῆς χρὴ ποιεῖν,\apparat{1062}{τοσαύτης R : \lecon{τὸ σαυτῆς} Dindorf (1839)}\\
κλαίειν ἐλεινῶς.\apparat{1063}{\lecon{κλαίειν} et \lecon{ἐπικλαίειν} R : κλάειν et ἐπικλάειν Brunck}\\\passeLigne}%
\end{groupe}%
%% v. 1063b
\begin{groupe}%
\personnageGrec{Μνησίλοχος}%
\typeVers{trim. iamb.}\trimIamb{\ligne{1063}\decalage{κλαίειν ἐλεινῶς.}σὲ δ᾽ ἐπικλαίειν ὕστερον.%
\commentaire{1063}{σὲ δ᾽ ἐπικλάειν ὕστερον.}{Puisqu'Écho répète tout, il est logique qu'elle répète \italique{en retour} (ἐπι-) les pleurs versés par Andromède... il y a là cependant une menace de vengeance clairement exprimée par Mnésiloque pour après sa libération (avec l'emploi à double sens du préfixe ἐπι- : hostilité et succession temporelle). 
%C'est en tout cas en ce sens que le comprit le scholiaste : il nous précise que Mnésiloque répond \italique{à part} (ἠρέμα) -- ce qui n'est d'ailleurs pas le cas puisqu'Euripide lui répond -- :%
%\citComm{ἠρέμα ὁ κηδεστὴς τοῦτο πρὸς τὸν Εὐριπίδην.}%
%Mais le verbe ἐπικλαίειν n'est pas autrement attesté à l'époque classique, au contraire d'ἐπικλάειν\marqueur{1}
}%
%\debutNote{1063}\note{1}{Remarque : la syllabe qui distingue les deux verbes (par sa longueur) tombe dans les deux cas qui se présentent au vers 1063 en position qui peut être indifféremment brève ou longue.}
%\commentaire{1063}{σὲ δ᾽ ἐπικλαίειν ὕστερον.}{\eff{Émouvoir ?}}
\\\passeLigne}%
\end{groupe}%
%% v. 1064
\begin{groupe}%
\personnageGrec{Εὐριπίδης}%
\trimIamb{ἐμοὶ μελήσει ταῦτά γ᾽. ἀλλ᾽ ἄρχου λόγων.\\\passeLigne}%
\end{groupe}%
%% v. 1065-1069a
\begin{groupe}%
\personnageGrec{Μνησίλοχος}%
\typeVers{metr. anap.}\anap{ὦ νὺξ ἱερὰ\apparat{1065}{ὁ μνησίλοχος ὡς ἀνδρομέδα in mg. \Rpc, \Mon}%
\commentaire{1065-72}{ὦ νὺξ ἱερὰ κτλ.}{Nous savons par une scholie que ces vers sont les premiers d'\italique{Andromède} (τοῦ προλόγου Ἀνδρομέδας εἰσβολή). Le mètre est anapestique\marqueur{1} ; les poètes tragiques y recourent effectivement pour certaines scènes de lamentation\marqueur{2}.}%
\debutNote{1065-72}%
\note{1}{De même que dans le prologue d'\italique{Iphigénie en Aulis}.}%
\note{2}{Voir \italique{The Songs of Aristophanes}, p. 57.}\\
ὡς μακρὸν ἵππευμα διώκεις\testimonia{1066-8}{(ὡς μακρὸν -- αἰθέρος) Suidas ι 531 $\sim$ Zonaras 1117}\\
ἀστεροειδέα νῶτα διφρεύουσ᾽\testimonia{1067}{(νῶτα) Moeris ν 6}\testimonia{1067-8}{Suidas αι 122 $\sim$ Zonaras 67, 78 (sed ἀστρο- pro ἀστερο-)}\commentaire{1067}{ἀστεροειδέα νῶτα.}{La voûte solide du firmament incorruptible. (Note de Willems.)}\\
αἰθέρος ἱερᾶς\\
τοῦ σεμνοτάτου δι᾽ Ὀλύμπου;%\commentaire{1068}{δι᾽ Ὀλύμπου.}{Hésychius (ο 654) :\citComm{Ὀλύμπιος· οὐράνιος.}}
\\\passeLigne}%
\end{groupe}%
%% v. 1069b
\begin{groupe}%
\personnageGrec{Εὐριπίδης}%
\anap{\ligne{1069}\decalage{τοῦ σεμνοτάτου δι᾽ Ὀλύμπου;}†δι᾽ Ὀλύμπου;†\apparat{1069}{\lecon{δι᾽ Ὀλύμπου} R : †δι᾽† ὀλυμπίας Bibi ap. Hesych. (ο 649)}%
\commentaire{1069}{τοῦ σεμνοτάτου δι᾽ Ὀλύμπου; --- δι᾽ Ὀλύμπου.}{%Prov : Oλυμπικὴ στοά: ἐπὶ τοῦ λάλου, ἢ ὅτι εἰκὸς ἔχειν τὴν στοὰν ἀντῳδὸn ἠχὼ, ἢ διὰ τὰςτῶνἐκεῖσυνηγμένωνβοάς
En bonne rigueur philologique, nous devons nous demander sur quels éléments probants s'appuyer dans notre aventureuse entreprise de rétablissement de l'\allemand{Urtext} des \italique{Thesmophories}. Car il ne suffit de partir à l'aventure, toutes brides de l'imaginations abattues ; encore faut-il apporter des indices de ce que nous avançons. 
Nous avons déjà évoqué des bribes de scholies ; nous en croiserons d'autres en route. 
Une variante orthographique entre un mot et son écho nous fournira aussi une amorce philologique. 
Cela est bien maigre... est-ce qu'Hésychius, ou les commentateurs qu'il recopie\marqueur{1}, dont nous avons assez montré toute l'importance, avait entre les mains une édition d'Aristophane assez bien conservée pour qu'y apparaisse certains détournements de la nymphe Écho ? Réponse : ce n'est pas impossible ! %\eff{Hum..} 
voyons cela avec première réplique ἐπικοκκάστρια de la nymphe Écho. \\
Voici la glose, parcellaire, d'Hésychius (ο 649) qui éclaire et justifie peut-être l'interruption d'Écho :%
\citComm{ὀλυμπία· δοκός\\
olympia : poutre.}%
La langue d'Aristophane possède la matrice métaphorique : \italique{poutre} $=$ \italique{poutre à laquelle on se pend} $=$ \italique{pendaison}. (Le français dispose d'une matrice semblable avec le mot \italique{corde}.) J'en veux pour preuves : \ \ \ 1° \ ce commentaire\marqueur{2} : [δοκὸς] παρὰ Ἀριστοφάνει ἐν Ταγηνισταῖς σημαίνει τὴν ἀγχόνην, commentaire repris par Hésychius (δ 2129) : δοκόν· τὴν ἀγχόνην.
\ \ \ 2° \ Dans les \italique{Oiseaux}, après avoir poutré l'importun chresmologue Hiéroclès, on l'envoie se faire pendre εἰς Ἐλύμνιον (vers 1126), ville d'Eubée dont le nom fait entendre à l'oreille grecque sa \italique{poutre} : ἐλύμνιαι· δοκοὶ ὀρόφιναι (Hésychius δ 2227).\\
Rappeler à Andromède sa condition de crucifiée \italique{pendue} à la falaise, voilà une incise qui pourrait être qualifiée de κακοστένακτος.\\
Je m'arrête ici, au bord de la reconstruction tentatrice. Car grande est la tentation de retoucher au texte même d'Andromède en y intégrant le iota d'ὀλυμπία et voir dans `δι᾽ Ὀλύμπου' une mention de Ζεύς Ὀλύμπιος\marqueur{3} ; qu'on se souvienne du vers 272 : Euripide jurait sur l'\italique{Éther, demeure de Zeus} (αἰθέρ᾽, οἴκησιν Διός\marqueur{4}). Le vrai Euripide parle ailleurs du \italique{ciel étoilé de Zeus} (\italique{Ion}, v. 1077 sq.) :%
\citComm{καὶ Διὸς ἀστερωπὸς\\
ἀνεχόρευσεν αἰθήρ.}%
La mention de Zeus lorsqu'un poète évoque le ciel constellé est d'ailleurs très naturelle puisque le quatrième astre le plus lumineux, après le Soleil, la lune et Vénus, n'est autre que \italique{Jupiter}, Διὸς ἀστήρ.
%279 crux ἱέρας· σημεῖον
%La première réplique ἐπικοκκάστρια de la nymphe Écho est-elle une consigne de διδάσκαλος, \italique{metteur en scène} ?\\
%En tout cas, le verbe qui répond le mieux au δι᾽ Ὀλύμπου est le verbe διαλάμπειν, \italique{briller}. (On verra que les assonances reposent essentiellement sur les consonnes.) Le verbe λάμπειν peut s'employer en parlant de la voix : παιὰν λάμπει (Soph.), \italique{le péan retentit}. Aristote (\italique{Problemata}, 922a31-39) utilise le composé dans l'un des ses \italique{problèmes} : Pourquoi ceux qui chantent en grand nombre gardent-ils mieux le rythme que ceux qui sont peu nombreux ? Διὰ τί οἱ πολλοὶ ᾄδοντες σώζουσι μᾶλλον τὸν ῥυθμὸν ἢ οἱ ὀλίγοι; 
%διολύγιον: ἐξηχοῦν ἐπὶ πολύ 
%μογ-είω, = μογέω, Dorian acc. to Hsch. (leg. μογίοντι).
% Exemples d'erreurs dans un lemme chez Hésychius où le 'ou' final est transformé en autre chose ('on')
%606 ἀγλαοπίστου (on):λαμπροπίστου
%1450 ἀζείρου:ἀζώστου. ἢ πολυζώστου (ἀζείριος: ἄζωτος Κ)
%3614 ἀποδοῦ (uson):ἀπόλυσον 
%8974 ἀψορρόου (on) :ἀνάπαλινῥέοντος,ἰταχυρρόου
%7255 εὐτυκάζου (on):εὔτυκον ἔχε, ἕτοιμον ()
%7355 "ἐφ᾽ἑαυτοῦ:καθ᾽ἑαυτοῦ (on)
%3693 προσβαλοῦ (on) :προσκατακλίθητι
%110  χαλκοπαρήου (reih):χαλκᾶςπαρειὰςἐχούσης
%En nombre, aucun des chanteurs ne \italique{brillerait}, c'est-à-dire ne \italique{se ferait remarquer} en s'élevant au-dessus de la masse : οὐδεὶς ἂν αὐτῶν διαλάμψειεν ὑπεράρας τὸ πλῆθος. Au contraire, moins nombreux, il leur est plus facile de \italique{briller} (et les chanteurs de rivaliser entre eux plutôt que de se régler sur le chef), ἐν δὲ τοῖς ὀλίγοις μᾶλλον διαλάμπουσιν· διὸ καθ’ αὑτοὺς ἐν αὐτοῖς μᾶλλον ἀγωνίζονται ἢ πρὸς τὸν ἡγεμόνα. Le verbe διαλάμπειν s'emploie donc pour des chanteurs qui \italique{brillent}, notamment par des modulations de rythme. \citComm{τοῦ σεμνοτάτου δι᾽ Ὀλύμπου; --- διαλάμπου\\Plus haut ! }%
}%
\debutNote{1069}%
\note{1}{Cf. le passionnant livre posthume de Pierre Boudreaux, \italique{Le texte d'Aristophane et ses commentateurs}.}%
\note{2}{Fragment d'Aristophane 528 K.-A.(515 Kock). Cf. l'édition Kassel-Austin des \italique{Poetae Comici Graeci} pour la juste distinction entre δοκός oxyton et δόκος paroxyton.% \eff{(Le \LSJ est à corriger sur ce point.)}
}%
\note{3}{Διός Ὀλυμπίου \italique{contra metrum} cependant.}%
\note{4}{Aristophane y revient au vers 100 des \italique{Grenouilles} : αἰθέρα Διὸς δωμάτιον.}%
%\commentaire{1069-72}{δι᾽ Ὀλύμπου — μέρος ἐξέλαχον — θανάτου τλήμων.}{Justifions notre modification du texte ainsi que la traduction. En philologie, tout finit en consultant son Hésychius (ε 2507) :\eff{Jeu du diable boîteux voir Louis Beck de Fouqières}\citComm{Ἔμπουσα· φάσμα δαιμονιῶδες ὑπὸ Ἑκάτης ἐπιπεμπόμενον καί, ὥς τινες, ἑνὶ ποδὶ χρώμενον. Ἀριστοφάνης δὲ τὴν Ἑκάτην ἔφη Ἔμπουσαν.\\Empuse : phantasme démonique envoyé par Hécate et, selon certains, n'ayant qu'une seule jambe. Aristophane appelle Hécate Empuse.}\italique{Nomen omen}, l'Empuse n'a qu'une seule jambe ; ἔμπους (singulier) suggère une entrave à \italique{une seule} jambe quand ἐμποδίζειν consiste à entraver les deux jambes d'\italique{impedementa}. Par dérivation de sens, que l'\italique{empuse} ait pu désigner un \italique{carcan}, cela nous paraît hautement probable au vu de la glose suivante d'Hésychius (ε 1264) :\citComm{ἑκάτη· ξύλον ἐν τοῖς φυλακίοις ᾧ τοὺς κακούργους προσδεσμεύοντες ἐμαστίγουν.\\hécate : morceau de bois dans les corps de garde auquel on attachait les malfaiteurs pour la flagellation.}L'explication qui me paraît la plus simple est bien qu'ἑκάτη ait servi de permutant métaphorique à ἔμπουσα -- Ἔμπουσα et Ἑκάτη étant parfois synonymes. (De toute façon, un glissement métaphorique dans l'autre sens, de ἑκάτη vers ἔμπουσα, suffirait tout autant à notre explication.)\\Tout cela pour justifier notre traduction : Euripide enjoint Mnésiloque de se libérer de ses entraves (qu'il s'aide et du ciel où il est attaché en costume d'Écho, il pourra l'aider ?) :\citComm{δι᾽ Ὀλύμπου — διάλυ(ε) (ἔ)μπουσ(αν).\\Sépare toi de cette empuse.}Voilà un exercice périlleux que de se livrer à des tentatives de reconstruction de texte sur ces trois répliques d'Écho pour faire apparaître d'hypothétiques jeux de mots ; laisser ainsi libre cours à l'imagination peut exposer à de rudes retours de bâton ; pour Aristophane, nous devons cependant impérieusement nous y livrer puisqu'il vient de nous prévenir que les échos seraient \italique{moqueurs}. Il nous semble bien ici qu'Aristophane joue sur le gore pour décrédibiliser la solennité de la pièce d'Euripide ; l'intuition provient de μέρος ἐξέλαχον qu'il est tentant de découper ainsi : μέρος ἐκ σέλαχος, \italique{la part du requin} (τό σέλαχος désignant tout type de poisson cartilagineux, raies, requins.) ; en effet, n'oublions pas qu'Andromède est attachée à une falaise pour servir de pâture à un monstre marin : le κῆτος dont il est question quelques vers plus haut. Aristophane, ce nous semble, s'adonne à un jeu de massacre en décrivant ce qu'est le sort attendu d'Andromède. Le reste se reconstruit selon cette perspective : δι᾽ Ὀλύμπου μέρος ἐξέλαχον θανάτου τλήμων $\longrightarrow$ διό λυποῦ μέρος ἐκ σέλαχος θαλάττης τέμνον « Et pour cette raison (quelle raison ?), tu seras chagrinée ; la part du requin cisailleur sorti de mer ». Pas très satisfaisant... à retoucher ! Pour ce qui est de l'apposition entre μέρος et σέλαχος, fréquente chez Euripide, voir Austin en français}%\debutNote{1069}%\note{1}{Voir, par exemple, le vers 293 des \italique{Grenouilles} et le scholiaste \italique{ad. loc.}.
%Tossi 2202 ou encore les vers 700-705 de \italique{Lysistrata} : Aristophane joue entre l'invocation à Hécate et le proverbe (Érasme n°394) ἀπ' οὐρᾶς τὴν ἔγχελυν ἔχεις, \italique{tu tiens l'anguille par la queue}, expliquant le singulier τοῦ σκέλους (ἐκτραχηλίσῃ était-il le terme propre d'halieutique pour la mise à mort du poisson pêché ?).}%
\\\passeLigne}%
\end{groupe}%
%% v. 1070-1071a
\begin{groupe}%
\personnageGrec{Μνησίλοχος}%
\anap{τί ποτ᾽ Ἀνδρομέδα περίαλλα κακῶν\apparat{1070}{περὶ ἄλλα R : \lecon{περίαλλα} Ellebodius}\\
μέρος ἐξέλαχον;\\\passeLigne}%
\end{groupe}%
%% v. 1071b
\begin{groupe}%
\personnageGrec{Εὐριπίδης}%
\anap{\ligne{1071}\decalage{μέρος ἐξέλαχον;}μόρος ἔξελ᾽ ἄχος.\apparat{1071}{μέρος ἐξέλαχον; R : \lecon{μόρος ἔξελ᾽ ἄχος.} Bibi}%
\commentaire{1071}{μέρος ἐξέλαχον;}{Sans \italique{tordre} le texte, il me semble que la réplique ἐπικοκκάστρια d'Écho s'entend assez bien -- soit dit avec toute la prudence qu'une telle chirurgie du texte impose :%
\citComm{μέρος ἐξέλαχον --- μόρος ἔξελ(ε)\marqueur{1} ἄχος\\
La Mort\marqueur{2} extirpe la douleur.}
Andromède se plaint des malheurs qui l'accablent ; Écho la rassure à sa manière : la mort l'en délivrera !.. Était-ce ce genre de répliques que le scholiaste qualifiait de `κακοστένακτος' ?\\
On trouve un tour analogue chez Sophocle\marqueur{3} :%
\citComm{Ἀλλ᾽ ἔσθ᾽ ὁ θάνατος λοῖσθος ἰατρὸς νόσων.}%
%Ταῖς νόσοις ὁ ϑάνατος ἀλλήλαις ἐπιλαβανούσαις οὐδ᾽ αὐτὸς παρεῖναι ἀναβάλλεται:.Μοσχίωνο
Et dans \italique{la langue des dieux} :%
\citComm{\textbf{Le trépas vient tout guérir} ;\\
Mais ne bougeons d'où nous sommes.\\
Plutôt souffrir que mourir,\\
C'est la devise des hommes.}
%\eff{maux sans guère de répit
%La mort guérit des maux
%la malheur t'a guère de répit
%la mâle-mort à son heure/ oeuvre}
}%
\debutNote{1071}%
\note{1}{Aoriste \italique{gnomique}.}%
\note{2}{Cette acception de μόρος a été relevée par à peu près tous les lexicographes anciens : Hésychius μ 1683 ; Photius μ 538 ; Συναγωγὴ λέξεων χρησίμων μ 267 ; Suidas μ 1257 ; \latin{Etymologicum genuinum} \italique{sub voce} μόρος· ὁ θάνατος ; \&c.}%
\note{3}{$=$ fr. 698 Radt $=$ fr. 636 Dindorf. Stobée (IV, 52, 26) a inséré la citation à son \italique{Florilège} et Michel Apostolius (II, 52c) la reprit pour ses \italique{Parœmiæ}.}\\\passeLigne}%
\end{groupe}%
%% v. 1072a
\begin{groupe}%
\personnageGrec{Μνησίλοχος}%
\anap{θανάτου τλήμων\\\passeLigne}%
\end{groupe}%
%% v. 1072b
\begin{groupe}%
\personnageGrec{Εὐριπίδης}%
\anap{\ligne{1072}\decalage{θανάτου τλήμων}†θανάτου τλήμων†.%\apparat{1072}{θανάτου τλήμων R : θαλάττης τέμνον Bibi}%
\commentaire{1072}{θανάτου τλήμων...}{L'on connaît par une scholie la fin du vers original d'\italique{Andromède} :%
\citComm{λείπει `μέλλουσα τυχεῖν'.}}\\\passeLigne}%
\end{groupe}%
%% v. 1073
\begin{groupe}%
\personnageGrec{Μνησίλοχος}%
\anap{ἀπολεῖς μ᾽ ὦ γραῦ στωμυλλομένη.%
\commentaire{1073-5}{ὦ γραῦ στωμυλλωμένη — στωμυλωμένη — νὴ Δί᾽.}{Le manuscrit de Ravenne (qui, rappelons-le, est le seul témoin complet de notre pièce) comporte deux orthographes différentes pour στωμυλλωμένη : avec un double lambda -λλ- au vers 1073 et un unique -λ- au vers 1074... Étourderie ou fantaisie de notre scribe jonglant avec une orthographe dont il est peu sûr ? Possible... ou bien preuve q'il recopiait scrupuleusement un texte sans parfois – comme ici – le comprendre ? c'est ce que nous allons tenter de montrer : il se trouve que le verbe στωμύλλω, \italique{babiller}, \italique{déblatérer}, s'écrit effectivement avec deux lambdas en grec classique. Cependant, l'adjectif στωμύλος (\italique{qui parle beaucoup}, et souvent avec nuance méliorative par rapport au verbe) qui en dérive n'en comporte qu'un seul ! Si donc, l'on fait le pari que le scribe ne s'y est pas trompé en conservant l'orthographe, savoir : on est en présence d'un participe moyen-passif au vers 1073 et de l'adjectif au vers 1074, il est alors aisé de retrouver la véritable réponse \italique{ἐπικοκκάστρια} d'Écho : στωμυλιο-μενής, « [je suis] de tempérament prolixe ».
Le composé est de la même facture que στωμυλιο-συλλεκτάδης, hapax qu'Aristophane mit dans la bouche d'Eschyle au vers 839 des \italique{Grenouilles} pour clore le bec à Euripide, cet intarissable \italique{collectionneur de commérages}. La seconde partie du composé est forgée avec le suffixe -μενής, \italique{au tempérament}, \italique{à la disposition...} (sur le modèle d'εὐ-μενής, δυσ-μενής, πρευ-μενής \&c).\\
Notre conjecture est renforcée par la concessive de Mnésiloque au vers suivant : il abonde, νὴ Δί᾽, ... néanmoins, γ᾽ : « ça oui par Zeus ! du moins... ».%
% Même scansion qu'aithéros iéras
}%
\apparat{1073}{στωμυλλωμένη R : \lecon{στωμυλλομένη} Grynaeus}\\\passeLigne}%
\end{groupe}%
%% v. 1074
\begin{groupe}%
\personnageGrec{Εὐριπίδης}%
\anap{στωμυλιομενής.\apparat{1074}{στωμυλωμένη R : \lecon{στωμυλιομενής} Bibi}\\\passeLigne}%
\end{groupe}%
%% v. 1075-1076a
\begin{groupe}%
\personnageGrec{Μνησίλοχος}%
\anap{νὴ Δί᾽, ὀχληρά γ᾽ εἰσήρρηκας\\
λίαν.\\\passeLigne}%
\end{groupe}%
%% v. 1076b
\begin{groupe}%
\personnageGrec{Εὐριπίδης}%
\anap{\ligne{1076}\decalage{λίαν.}†λίαν†.\\\passeLigne}%
\end{groupe}%
%% v. 1077-1078a
\begin{groupe}%
\personnageGrec{Μνησίλοχος}%
\anap{ὦγάθ᾽ ἔασόν με μονῳδῆσαι,\\
καὶ χαριεῖ μοι παῦσαι.%\apparat{1078}{Μνη. χαριεῖ μοι παῦσαι --- Εὐρ. παῦσαι R : \lecon{Μνη. χαριεῖ μοι ᾽παΰσαι --- simul Μνη. et Εὐρ. παῦσαι} Bibi}%
\commentaire{1078}{καὶ χαριεῖ μοι παῦσαι --- παῦσαι.}{Mnésiloque-Andromède, qui ne peut chanter sa monodie correctement, demande à Euripide-Écho de cesser de l'interrompre :%
\citComm{καὶ χαριεῖ μοι παῦσαι...}%
«~Tu m'obligeras d'arrêter...~» -- «~De crier~» interrompt facilement Écho. (Et, bien-sûr, en le lui criant à l'oreille :)%
\citComm{ (ἐ)παΰσαι.}%
(Le verbe ἐπαΰω, attesté par Eschyle dans ses \italique{Choéphores} (v. 828), est forgé avec le préfixe ἐπι- à l'instar d'ἐπικοκκάστρια et le verbe αὔω, \italique{crier} (cf. le vers 1019).)\\
Une telle provocation explique la réaction d'emportement de Mnésiloque :%
\citComm{βάλλ᾽ ἐς κόρακας.}%
%Andromède, qui ne peut chanter sa monodie correctement, demande à Écho de cesser de l'interrompre.
%	 par \italique{retour} (ἐπι-) de \italique{cri} (αὔω, cf. le vers 1019). 
%Mnésiloque s'apprêtait donc à dire :%
%\citComm{καὶ χαριεῖ μοι (ἐ)παΰσαι παῦσαι.}%
%Évidemment, cette répétition de deux `παυσαι' est entrecoupée par Écho.\\
%Pour le sens qu'il faut donner à son bien peu compatissant παῦσαι ; l'on le rapprochera du vers 1382 des \italique{Oiseaux} :%
%\citComm{παῦσαι μελῳδῶν, ἀλλ' ὅ τι λέγεις εἰπέ μοι.}%
%Ou encore de ce fragment attribué à Aristophane (962 (adesp. 601) Kassel-Austin = Photius π 519) :%
%\citComm{παῦσαι μελῳδοῦσ', ἀλλὰ πεζῇ μοι φράσον}
%%
%La réplique déclenche pas moins la fureur de Mnésiloque : βάλλ᾽ ἐς κόρακας.%
%παῦσαι peut être à la fois un infinitif et un impératif aoriste ; et c'est la-dessus qu'Aristophane nous semble jouer. La plupart des éditeurs ponctuent avec une pause entre χαριεῖ μοι et παῦσαι – et inclinent donc vers un double impératif avec le sens de « finis ! », c'est-à-dire « tais-toi ! ». Pourtant, juste après, Mnésiloque-Andromède demande, à être libéré de ses entraves, il est donc bien plus indiqué de comprendre un appel au secours – interrompu par Écho – de Mnésiloque, une supplique pour une délivrance morale avant la délivrance physique du vers suivant (ainsi, chez Sophocle, \italique{Électre}, 987, παῦσον ἐκ κακῶν ἐμέ ; ou dans l'\italique{Odyssée}, XXIV, 384, παύσαντο πόνου). On peut donc traduire ainsi : « Et tu m'obligeras en abrégeant...[mon tourment ? mes peines ?] — Abrège ! »
}\\\passeLigne}%
\end{groupe}%
%% v. 1078b
\begin{groupe}%
\personnageGrec{Εὐριπίδης}%
\anap{\ligne{1078}\decalage{καὶ χαριεῖ μοι παῦσαι.}᾽παΰσαι.\apparat{1078}{(Εὐρ.) παῦσαι R : \lecon{(Εὐρ.) ᾽παΰσαι} Bibi}%
\\\passeLigne}%
\end{groupe}%
%% v. 1079a
\begin{groupe}%
\personnageGrec{Μνησίλοχος}%
\anap{βάλλ᾽ ἐς κόρακας.%
\commentaire{1079}{βάλλ᾽ ἐς κόρακας — βάλλ᾽ ἐς κόρακας.}{Dans notre entreprise fort périlleuse de dénicher des jeux de mots tout au long de ce passage, en voici un qui est moins spéculatif et qui va nous permettre de reprendre haleine ! %
%\eff{Nous avons déjà analysé au vers 942 comment Aristophane jouait avec le mot κόραξ et ses deux sens possibles : le corbeau ou bien le crochet qui entrave Mnésiloque. }
Les Grecs usaient d'une formule d'imprécation : ἐς κόρακας, équivalente de notre « [va] au diable ! ». Aristophane l'emploie fréquemment avec le verbe βάλλειν : \italique{Guêpes} v. 835 ; \italique{Ploutos} v. 782 ; \italique{Nuées} v. 133 – digressons d'ailleurs sur ce passage des \italique{Nuées} puisqu'Aristophane recourt au même jeu de mots qu'ici : Strepsiade arrive devant le pensoir de Socrate et toque à la porte, il se fait recevoir par un peu accueillant :%
\citComm{Βάλλ᾽ ἐς κόρακας· / Τίς ἐσθ᾽ ὁ κόψας τήν θύραν;}%
signalons à toutes fins utiles que les κόρακες dont parle le disciple de Socrate sont des marteaux de porte\marqueur{1}. 
Le disciple tint donc à peu près ce langage : %
\citComm{Βάλλε κόρακας·/ Τίς ἐσθ᾽ ὁ κόψας τὴν θύραν;\\
Bas les marteaux ! Qui est ce tambourineur de porte ?}%
(compatible avec le mètre iambique et surtout avec le bon sens : qui – fût-il disciple de Socrate – s'aviserait de demander l'identité de quelqu'un qu'on vient d'envoyer se faire voir chez les corbeaux grecs ?)\\
Il en va de même dans le passage qui nous intéresse présentement : Écho conseille frivolement à Mnésiloque de retirer les crochets qui le retiennent : « Mais va te faire voir ! — Et toi mets bas le fermoir ! »}%
\debutNote{1079}%
\note{1}{Sens attesté par Pollux, X, 22.20.}%
\\\passeLigne}%
\end{groupe}%
%% v. 1079b
\begin{groupe}%
\personnageGrec{Εὐριπίδης}%
\anap{\ligne{1079}\decalage{βάλλ᾽ ἐς κόρακας.}βάλλε σύ κόρακας.\apparat{1079}{βάλλ᾽ ἐς κόρακας R : \lecon{βάλλε σύ κόρακας} Bibi}\\\passeLigne}%
\end{groupe}%
%% v. 1080a
\begin{groupe}%
\personnageGrec{Μνησίλοχος}%
\anap{†τί τὸ κακόν†;\apparat{1080}{\lecon{τί τὸ κακόν} bis R : τί κακόν bis Bentley}%
\commentaire{1080}{τί τὸ κακόν;}{La leçon du \italique{Ravennnas}, «~τί τὸ κακόν; --- τί τὸ κακόν;~», ne respecte pas le mètre anapestique. Gardons-nous pour autant de supprimer avec Bentley l'article `τὸ'. 
Si la formule «~τί τὸ κακόν;~» (ou d'autres très approchantes) se rencontre couramment chez Aristophane\marqueur{1}, sa version lapidaire, «~τί κακόν;~», elle, ne se trouve jamais\marqueur{2}. La bonne solution n'est donc pas celle de Bentley, et pour notre part nous nous résignons à \latin{crux}-désespérer †%
%κάγκελον, ου (τὸ) c. le suiv. Sch.-Ar. Eq. 638, 672.
%καγχάζω
%Hésychius 346 †κακόρας: κατακόψας. παρὰ Εὔκλῳ 
}%
\debutNote{1080}%
\note{1}{\italique{Thesmophories}, v. 610 ; \italique{Paix}, v. 322 ; \italique{Oiseaux}, v. 1213 ; \italique{Guêpes} v. 1136, τουτὶ τὸ κακὸν τί ἐστι ; \italique{Paix}, v. 181.}%
\note{2}{Certains éditeurs prétendent que si : au vers 1085. En vérité, ces faussaires commettent là une double forfaiture en «~corrigeant~» le présent vers \italique{metri gratia} ainsi que le vers 1085 ; le second forfait servant d'alibi au premier : au vers 1085, ils font apparaître «~une seconde fois~» le tour `τί κακόν' alors que le manuscrit porte `σί κακόν', confirmé par le scholiaste (ἀντὶ τοῦ `σοὶ κακόν'). Et quand bien même : c'est dans la bouche d'un Scythe qu'Aristophane la met au vers 1085, l'argument pour la faire dire aussi à Mnésiloque est irrecevable. En dernière analyse, ces éditeurs manquent cruellement de bon sens : le Scythe en son jargon ne tient déjà aucun compte des aspirations, de la grammaire, il ignore sa conjugaison et ses déclinaisons, il iotacise à l'envi... et l'on voudrait encore qu'il sigmatise ses `τ' ?! un tel massacre du grec aurait achevé d'être tout-à-fait inintelligible aux spectateurs.}\\\passeLigne}%
\end{groupe}%
%% v. 1080b
\begin{groupe}%
\personnageGrec{Εὐριπίδης}%
\anap{\ligne{1080}\decalage{†τί τὸ κακόν†;}†τί τὸ κακόν;†%
\commentaire{1080-2}{τί τὸ κακόν; [...] ὀτότυζ᾽.}{Quel est le jeu de scène de tout ce passage, dont les répliques d'Écho demeurent obscures ? Il me semble qu'on peut peut y répondre avec vraisemblance : Euripide-Écho s'est approchée de Mnésiloque (cf. v. 1090, γυνὴ πλησίον αὕτη) ; la raison en est évidente : Euripide tente de libérer Mnésiloque. 
Écho parvient cependant à se cacher dans un premier temps du Scythe (οὖτος τί λαλῖς; --- πῶτε τὸ πωνή;) ; c'est l'iconographie qui nous fournit la réponse : l'acteur jouant Écho se cachait derrière un grand voile\marqueur{1}.
Et lorsque l'archer Scythe veut mettre la main sur Euripide, ce dernier \italique{s'envole} (ποῖ ποῖ πεύγει -- certainement à l'aide de la μηχανή), lançant un narquois `οὐχ αἱρήσεις'.}%
\debutNote{1080-2}%
\note{1}{Sur un vase attique à figure rouge du \nieme{IV} siècle av. J.-C. (n° 7864 au catalogue de Beazley) représentant Dionysos et sa suite, on voit la nymphe Écho, ailée, se voiler le visage derrière son manteau. Ce détail est reproduit en annexe.}\\\passeLigne}%
\end{groupe}%
%% v. 1080c
\begin{groupe}%
\personnageGrec{Μνησίλοχος}%
\anap{\ligne{1080}\decalage{†τί τὸ κακόν†;†τί τὸ κακόν;†}ληρεῖς.\\\passeLigne}%
\end{groupe}%
%% v. 1080d
\begin{groupe}%
\personnageGrec{Εὐριπίδης}%
\anap{\ligne{1080}\decalage{†τί τὸ κακόν†;†τί τὸ κακόν;†ληρεῖς.}†ληρεῖς†.\\\passeLigne}%
\end{groupe}%
%% v. 1081a
\begin{groupe}%
\personnageGrec{Μνησίλοχος}%
\anap{οἴμωζ᾽.\\\passeLigne}%
\end{groupe}%
%% v. 1081b
\begin{groupe}%
\personnageGrec{Εὐριπίδης}%
\anap{\ligne{1081}\decalage{οἴμωζ᾽.}†οἴμωζ᾽†.\\\passeLigne}%
\end{groupe}%
\endnumbering
\end{Leftside}
\begin{Rightside}
\beginnumbering
%% v. 1015-1055
\begin{groupe}%
\paratragiquesFr{me revêtit de cette crocote,}
\end{groupe}%
%% v. 1015-1055
\begin{groupe}%
\paratragiquesFr{puis de surcroît me dépêcha vers ce}
\end{groupe}%
%% v. 1015-1055
\begin{groupe}%
\paratragiquesFr{temple, au milieu de femmes.}
\end{groupe}%
%% v. 1015-1055
\begin{groupe}%
\paratragiquesFr{Hélas!}
\end{groupe}%
%% v. 1015-1055
\begin{groupe}%
\paratragiquesFr{démon implacable de mon destin !}
\end{groupe}%
%% v. 1015-1055
\begin{groupe}%
\paratragiquesFr{Ô moi le maudit !}
\end{groupe}%
%% v. 1015-1055
\begin{groupe}%
\paratragiquesFr{Qui, en présence de mes misères, ne tournerait les yeux}
\end{groupe}%
%% v. 1015-1055
\begin{groupe}%
\paratragiquesFr{sur cette inenviable souffrance ?}
\end{groupe}%
%% v. 1015-1055
\begin{groupe}%
\paratragiquesFr{Puisse l'Étoile dévorante de l'éther}
\end{groupe}%
%% v. 1015-1055
\begin{groupe}%
\paratragiquesFr{anéantir un infortuné !%\commentaire{1050-1}{πυρφόρος αἰθέρος ἀστὴρ ...}{Cf. sur le sens de ce passage et le rôle de l'étoile Sirius : Bull. cité, note 26, ci-après.} !
}
\end{groupe}%
%% v. 1015-1055
\begin{groupe}%
\paratragiquesFr{Car il ne m'est plus doux de voir la lumière immortelle,}
\end{groupe}%
%% v. 1015-1055
\begin{groupe}%
\paratragiquesFr{suspendu que je suis,}
\end{groupe}%
%% v. 1015-1055
\begin{groupe}%
\paratragiquesFr{attendant comme terme d'un supplice jugulant%\commentaire{1054}{}{Allusion au carcan qui lui serre le cou. Lisez : δία μένων. Cf. sur ce passage : Bull. cité, note 27, ci-après.}
}
\end{groupe}%
%% v. 1015-1055
\begin{groupe}%
\paratragiquesFr{un trajet rapide vers les morts.%\commentaire{1015-55}{}{Tout ce chant est d'une incohérence voulue.}.
\\\passeLigne}%
\end{groupe}%
%% v. 1056-1057
\begin{groupe}%
\personnageFr{Euripide}%
\trimIambFr{\ligne{1056}(Dans la coulisse et sans être vu). Salut, chère enfant. Mais ton père Cepheus\\
qui t'exposa, que les dieux le confondent.\\\passeLigne}%
\end{groupe}%
%% v. 1058
\begin{groupe}%
\personnageFr{Mnésiloque}%
\trimIambFr{Qui es-tu, toi qui compatis à mon tourment ?\\\passeLigne}%
\end{groupe}%
%% v. 1059-1063a
\begin{groupe}%
\personnageFr{Euripide}%
\trimIambFr{Je suis Écho, la rediseuse goguenarde,\\
celle qui, l'an passé, en ce même lieu,\\
prêtai mon aide à Euripide.\\
Songe, mon enfant, à tenir bien ta partie\\
et à gémir lamentablement.\\\passeLigne}%
\end{groupe}%
%% v. 1063b
\begin{groupe}%
\personnageFr{Mnésiloque}%
\trimIambFr{\ligne{1063}\espace{2.4cm}Et toi, à gémir de plus belle après moi.\\\passeLigne}%
\end{groupe}%
%% v. 1064
\begin{groupe}%
\personnageFr{Euripide}%
\trimIambFr{Je n'y manquerai pas. Commence.\\\passeLigne}%
\end{groupe}%
%% v. 1065-1069a
\begin{groupe}%
\personnageFr{Mnésiloque}%
\anapFr{Ô Nuit sainte,\\
quelle longue carrière tu fournis,\\
en guidant ton char sur le dos constellé\\
de l'éther sacré,\\
à travers le très auguste Olympe.\\\passeLigne}%
\end{groupe}%
%% v. 1069b
\begin{groupe}%
\personnageFr{Euripide}%
\anapFr{\ligne{1069}\espace{4cm}Olympe.\\\passeLigne}%
\end{groupe}%
%% v. 1070-1071a
\begin{groupe}%
\personnageFr{Mnésiloque}%
\anapFr{Pourquoi seule Andromède reçut-elle tant de maux\\
en partage ?\\\passeLigne}%
\end{groupe}%
%% v. 1071b
\begin{groupe}%
\personnageFr{Euripide}%
\anapFr{\ligne{1071}\espace{2.5cm}En partage.\\\passeLigne}%
\end{groupe}%
%% v. 1072a
\begin{groupe}%
\personnageFr{Mnésiloque}%
\anapFr{Malheureuse que la mort...%\commentaire{1072}{θανάτου τλήμων...}{(Attend.)}
\\\passeLigne}%
\end{groupe}%
%% v. 1072b
\begin{groupe}%
\personnageFr{Euripide}%
\anapFr{\ligne{1072}\espace{2.7cm}Malheureuse que la mort.\\\passeLigne}%
\end{groupe}%
%% v. 1073
\begin{groupe}%
\personnageFr{Mnésiloque}%
\anapFr{Tu me tueras, vieille, avec ton babil.\\\passeLigne}%
\end{groupe}%
%% v. 1074
\begin{groupe}%
\personnageFr{Euripide}%
\anapFr{Ton babil.\\\passeLigne}%
\end{groupe}%
%% v. 1075-1076a
\begin{groupe}%
\personnageFr{Mnésiloque}%
\anapFr{Par Zeus, ton ingérence m'excède\\
au dernier point.\\\passeLigne}%
\end{groupe}%
%% v. 1076b
\begin{groupe}%
\personnageFr{Euripide}%
\anapFr{\ligne{1076}\espace{1cm}Au dernier point.\\\passeLigne}%
\end{groupe}%
%% v. 1077-1078a
\begin{groupe}%
\personnageFr{Mnésiloque}%
\anapFr{Mon brave, laisse-moi chanter ma monodie.\\
Tu m'obligeras : cesse !\\\passeLigne}%
\end{groupe}%
%% v. 1078b
\begin{groupe}%
\personnageFr{Euripide}%
\anapFr{\ligne{1078}\espace{3.5cm}(Criard :) C'est ce... que j'allais dire.\\\passeLigne}%
\end{groupe}%
%% v. 1079a
\begin{groupe}%
\personnageFr{Mnésiloque}%
\anapFr{Va-t'en aux corbeaux.\\\passeLigne}%
\end{groupe}%
%% v. 1079b
\begin{groupe}%
\personnageFr{Euripide}%
\anapFr{\ligne{1079}\espace{2.9cm}Va-t'en aux corbeaux.\\\passeLigne}%
\end{groupe}%
%% v. 1080a
\begin{groupe}%
\personnageFr{Mnésiloque}%
\anapFr{Qu'est-ce qui t'arrive ?\\\passeLigne}%
\end{groupe}%
%% v. 1080b
\begin{groupe}%
\personnageFr{Euripide}%
\anapFr{\ligne{1080}\espace{1.5cm}Qu'est-ce qui t'arrive ?\\\passeLigne}%
\end{groupe}%
%% v. 1080c
\begin{groupe}%
\personnageFr{Mnésiloque}%
\anapFr{\ligne{1081}\espace{3cm}Tu veux rire.\\\passeLigne}%
\end{groupe}%
%% v. 1080d
\begin{groupe}%
\personnageFr{Euripide}%
\anapFr{\ligne{1082}\espace{4.2cm}Tu veux rire.\\\passeLigne}%
\end{groupe}%
%% v. 1081a
\begin{groupe}%
\personnageFr{Mnésiloque}%
\anapFr{Va te promener..\\\passeLigne}%
\end{groupe}%
%% v. 1081b
\begin{groupe}%
\personnageFr{Euripide}%
\anapFr{\ligne{1081}\espace{1.1cm}Va te promener..\\\passeLigne}%
\end{groupe}%
\endnumbering
\end{Rightside}
\end{pages}
\Pages
\clearpage\ifodd\value{page}\hbox{}\newpage\fi
\enTeteG{ΕΠΕΙΣΟΔΙΟΝ Ε' (Περσεύς) --- Codex Ravennas 429, f. 175v (Classense 177v) --- vv. 1082–1127}
\begin{adjustwidth}{-1.5cm}{-3.0cm}
\begin{adjustbox}{max totalheight=21.0cm, max width=16.5cm, keepaspectratio}
\begin{minipage}{16.5cm}
{\fontsize{7.2pt}{8.5pt}\selectfont
\noindent\textbf{Folio 27 --- Codex Ravennas 429, \texttt{f. 175v} (\texttt{Classense 177v}) : ΕΠΕΙΣΟΔΙΟΝ Ε' (Περσεύς), vv. 1082–1127}\hfill \textit{Témoins :} $\text{R}^1, \Rpc, \Rdeux$ (\Mon)\par\vspace{2pt}
\hrule height 0.4pt\vspace{2.5pt}
}
{\fontsize{8.1pt}{9.5pt}\selectfont
\ligneRav{1082}{μ\auDessus{ν}{Η}}{\hspace{1.1em}ὀτότυζ᾽. : †ὀτότυζ᾽†·}{\textbf{1082} \lecon{ὀτότυζ᾽} R : ὀλόφυγ᾽ Bibi}
\ligneRav{1083}{τοξ/}{\hspace{1.1em}οὖτος τί λαλ\inRas{εῖ}{ῖ}ς; \hfill \textbf{[Ἠχώ :]} οὖτος τί λαλ\inRas{εῖ}{ῖ}ς;}{\textbf{1083} οῦτος οὗτος R : \lecon{οὖτος} bis Brunck \espApp λαλεῖς bis \Rac : \lecon{λαλῖς} bis \Rdeux, Σ}
\ligneRav{1084}{τοξ/}{\hspace{1.1em}πρυτάνεις καλέσω· \hfill \textbf{[Ἠχώ :]} πρυτάνεις καλέσω·}{}
\ligneRav{1085}{τοξ/}{\hspace{1.1em}σί κακόν; \hfill \textbf{[Ἠχώ :]} σί κακόν;}{\textbf{1085} σί bis R, sch. : τί Coulon}
\ligneRav{1086}{τοξ/}{\hspace{1.1em}\difference{†}π\inRas{ρ}{}ῶτε\difference{†} τὸ φωνή; \hfill \textbf{[Ἠχώ :]} \difference{†}π\inRas{ρ}{}ῶτε\difference{†} τὸ φωνή;}{\textbf{1086} π\inRas{ρ}{}ωτε bis \Rpc : πωτε bis \Mon}
\ligneRav{1087}{τοξ/}{\hspace{1.1em}σὺ λαλ\auDessus{ῖ}{εῖ}ς; \hfill \textbf{[Ἠχώ :]} σὺ λαλ\auDessus{ῖ}{εῖ}ς;}{\textbf{1087} \lecon{λαλῖς} bis R : λαλεῖς bis \Rpc}
\ligneRav{1088}{τοξ/}{\hspace{1.1em}\difference{κλαύσαι}· \hfill \textbf{[Ἠχώ :]} \difference{κλαύσαι}·}{\textbf{1088} κλαύσαιμι bis R : κλαύσει bis Brunck : \lecon{κλαύσαι} bis Rogers : κλαῦσ' ἔτι bis Austin}
\ligneRav{1089}{τοξ/}{\hspace{1.1em}κακκάσκι μοι; \hfill \textbf{[Ἠχώ :]} κακκάσκι μοι;}{\textbf{1089} κακκασκι bis R : κἀγκάσκις bis van Leeuwen (<ς> add. Fritzsche) : κἀγκάσκι bis Coulon : alii alia}
\ligneRav{1090}{χορ/}{\hspace{1.1em}μὰ Δί᾽ ἀλλὰ γυνὴ πλησίον αὕτη·}{\textbf{1090} (Μνη.) edd. plerique : \lecon{(Χορ.)} Bibi, cf. Σ}
\ligneRav{1091}{\auDessus{εὐ}{ρ}}{\hspace{1.1em}\difference{πλήξασ᾿ αὕτη}·}{\textbf{1091} πλησίον αὕτη R : \lecon{πλήξασ᾿ αὕτη} Bibi}
\ligneRav{1092}{τοξ/}{\hspace{1.1em}ποῦ ᾿στ᾽ ἠ μιαρά; : καὶ δὴ φεύγει·}{\textbf{1092} 'σθ' ἡ R : ᾿στ᾽ ἠ Brunck\espApp (Μνη.) edd. plerique : \lecon{(Χορ.)} Bibi}
\ligneRav{1093}{τοξ/}{\hspace{1.1em}ποῖ ποῖ \difference{πεύγεις}; : ποῖ ποῖ πέτει·}{\textbf{1093} φεύγεις R : \lecon{πεύγεις} Enger\espApp (Εὐρ.) ποῖ ποῖ πέτει om. R : \lecon{suppl.} Bibi}
\ligneRav{1094}{τοξ/}{\hspace{1.1em}οὐ καιρήσεις. : \difference{οὐχ αἱρήσεις}·}{\textbf{1094} (Τοξ.) οὐ καιρήσεις om. R : \lecon{suppl.} Brunck ap. Σ\espApp οὐκ᾽ αἰρήσεις; R : \lecon{οὐχ αἱρήσεις.} Wilamowitz (\italique{Kleine Schriften} IV, p.486n.1)}
\ligneRav{1095}{τοξ/}{\hspace{1.1em}ἔτι γὰρ γρύζεις; : †ἔτι γὰρ γρύζεις;†}{}
\ligneRav{1096}{τοξ/}{\hspace{1.1em}λαβὲ \difference{τὴ} μιαρά. : †λαβὲ τὴ μιαρά.†}{\textbf{1096} τῇ et τῆ R : \lecon{τὴ} bis Bothe}
\ligneRav{}{τοξ/}{\hspace{1.1em}λάλο καὶ κατάρατο γύναικο·}{}
\ligneRav{1098}{ὡς περσευς}{ὦ θεοὶ τίν᾽ εἰς γῆν βαρβάρων ἀφίγμεθα}{\textbf{1098} ευρ\auDessus{ι}{π} ως περσευς in mg. add. \Rpc, \Mon}
\ligneRav{}{}{ταχεῖ πεδίλῳ; διὰ μέσου γὰρ αἰθέρος}{}
\ligneRav{1100}{}{τέμνων κέλευθον πόδα τίθημ᾽ ὑπόπτερον}{}
\ligneRav{1101}{}{Περσεὺς πρὸς Ἄργος \difference{ναυστολῶν} τὸ Γοργόνος}{\textbf{1101} ναυτολῶν R : \lecon{ναυστολῶν} Canini (\italique{navigans} iam Divus)}
\ligneRav{1102}{—}{$[\![\text{in marg. inf. \Rdeux{} : }$κάρα κομίζων : τί λέγι; τὴ \difference{Γόργος} πέρι$]\!]$}{\textbf{1102} Γοργόνος R : \lecon{Γόργος} Σ, Bothe}
\ligneRav{1103}{—}{$[\![\text{in marg. inf. \Rdeux{} : }$τὸ γραμματέο σὺ τὴ \difference{κεπαλή}; : τὴν Γοργόνος$]\!]$}{\textbf{1103} κεφαλή R : \lecon{κεπαλή} Bothe}
\ligneRav{}{—}{ἔγωγε φημί. : Γόργο τοι κἀγὼ λέγι·}{}
\ligneRav{1105}{\auDessus{εὐ}{ρ}}{ἔα· τίν᾽ ὄχθον τόνδ᾽ ὁρῶ καὶ παρθένον}{}
\ligneRav{}{}{θεαῖς ὁμοίαν ναῦν ὅπως ὡρμισμένην;}{}
\ligneRav{1107}{κ\auDessus{η}{δ} ἀνδρ.}{ὦ ξένε \difference{κατοίκτιρόν} με τὴν παναθλίαν,}{\textbf{1107} κ\auDessus{η}{δ} ως ανδρομ\auDessus{ε}{δ} in mg. \Rpc, \Mon \espApp κατοικτειρόν R : \lecon{κατοίκτιρόν} van Herwerden}
\ligneRav{1108}{—}{λῦσόν με δεσμῶν. : οὐκὶ μὶ λαλῆσι σύ;}{\textbf{1108} \lecon{μὶ λαλῆσι} R : μὴ λαλῆσι Brunck}
\ligneRav{1109}{}{κατάρατο τολμᾷς ἀποτανουμένη \difference{λαλῖς};}{\textbf{1109} λαλᾶς R : \lecon{λαλῖς;} Bergk (cf. vv. 1082 et 1087) : λαλᾷς; Hall-Geldart}
\ligneRav{1110}{\auDessus{εὐ}{ρ}}{ὦ παρθέν᾽ \difference{οἰκτίρω} σὲ κρεμαμένην ὁρῶν·}{\textbf{1110} οἰκτείρω R : \lecon{οἰκτίρω} van Herwerden}
\ligneRav{}{τοξ/}{οὐ παρτέν᾽ ἐστίν, ἀλλ᾽ ἀμαρτωλὴ γέρων}{}
\ligneRav{}{—}{καὶ κλέπτο καὶ πανοῦργο. : ληρεῖς ὦ Σκύθα·}{}
\ligneRav{}{}{αὕτη γάρ ἐστιν Ἀνδρομέδα παῖς Κηφέως·}{}
\ligneRav{1114}{τοξ/}{σκέψαι τὸ \difference{κύστο}· \difference{μή τι} \difference{μικκὸν} παίνεται;}{\textbf{1114} σκύτο R : \lecon{κύστο} Scaliger grat. metr. : σῦκο Sommerstein \espApp μῆτι R : \lecon{μή τι} \Mon (\italique{numquid} Divus) \espApp μικτὸν R : \lecon{μικκὸν} Bentley : μικρὸ Enger}
\ligneRav{1115}{\auDessus{εὐ}{ρ}}{φέρε δεῦρό μοι τὴν χεῖρ᾽, ἵν᾽ ἅψωμαι κόρης·}{\textbf{1115} δεῦρο δεῦρο R : \lecon{δεῦρο} Biset}
\ligneRav{}{}{φέρε Σκύθ᾽· ἀνθρώποισι γὰρ νοσήματα}{}
\ligneRav{}{}{ἅπασίν ἐστιν· ἐμὲ δὲ καὐτὸν τῆς κόρης}{}
\ligneRav{}{—}{ταύτης ἔρως εἴληφεν. : οὐ ζηλῶσί σε·}{}
\ligneRav{1119}{}{ἀτὰρ εἰ \difference{τὸ πρωκτὸ} δεῦρο περιεστραμμένον,}{\textbf{1119} τῶ πρωκτῶ R : \lecon{τὸ πρωκτὸ} Biset}
\ligneRav{1120}{}{οὐκ \difference{ἐπτόνησά σ᾽} αὐτὸ πυγίζεις ἄγων·}{\textbf{1120} ἐπτόνησας R : \lecon{ἐπτόνησά σ᾽} Bentley}
\ligneRav{}{\auDessus{εὐ}{ρ}}{τί δ᾽ οὐκ ἐᾷς λύσαντά μ᾽ αὐτὴν ὦ Σκύθα}{}
\ligneRav{1122}{}{πεσεῖν \difference{ἐς} εὐνὴν καὶ γαμήλιον λέχος;}{\textbf{1122} ἐστ᾽ R : ἐς τ᾽ \Mon : \lecon{ἐς} Grynaeus}
\ligneRav{}{τοξ/}{εἰ σπόδρ᾽ ἐπιτυμεῖς τὴ γέροντο πύγισο,}{}
\ligneRav{1124}{}{τὴ σανίδο τρήσας \difference{ἐξόπιστο} πρώκτισον·}{\textbf{1124} ἐξόπισθο R : \lecon{ἐξόπιστο} Daubuz}
\ligneRav{1125}{\auDessus{εὐ}{ρ}}{μὰ Δί᾽ ἀλλὰ λύσω \difference{δεσμά}. : μαστιγῶ σ᾽ ἄρα·}{\textbf{1125} δέμας R : \lecon{δεσμά} Grynaeus}
\ligneRav{1126}{\auDessus{εὐ}{ρ}}{καὶ μὴν ποιήσω τοῦτο : \difference{τὸ κεπαλή σ᾽} (τὸ κ\inRas{αὶ }{ε}παλῆς) ἄρα}{\textbf{1126} τὸ κεπαλῆς (ex καὶ παλῆς) \Rpc, \Mon : \lecon{τὸ κεπαλή σ᾽} Brunck}
\ligneRav{1127}{}{τῇ ξιπομάκαιρᾳ \difference{ἀποκεκόψι} (ἀποκεκόψ\inRas{᾽}{οι}) τουτογί·}{\textbf{1127} ἀποκέκοιψ᾽ \Rac : ἀποκεκόψοι \Rpc : ἀποκεκόψο \Mon : \lecon{ἀποκεκόψι} Brunck\espApp πέπυσται R : \lecon{πέπαισται} Σ}
}
\vspace{3pt}\noindent\rule{4.5cm}{0.4pt}\par\vspace{1.5pt}
{\fontsize{7.0pt}{8.3pt}\selectfont
\noindent\textbf{Loci paralleli (vv. 1082–1127) ---} \textbf{1125} (δεσμά) Moeris δ 28\par
}
\vspace{2.5pt}\noindent\rule{4.5cm}{0.4pt}\par\vspace{1.5pt}
{\fontsize{6.9pt}{8.2pt}\selectfont
\noindent\textbf{Suite (en commentaire) de l'apparat critique \& notes paléographiques (\texttt{f. 175v}) ---} \textbf{Omission des vv. 1102–1103 par $\text{R}^1$ et ajout par \Rdeux{} dans la marge inférieure :} Le copiste $\text{R}^1$ avait sauté les \textbf{vv. 1102–1103} ($\text{Γόργος πέρι λέγις τὸ γραμματέο / σὺ τὴν κεπαλή;}$ — omis également par \Mon{} qui n'a pas regardé le bas de page) ; le réviseur \textbf{\Rdeux{} les a rajoutés tout en bas du folio dans la marge inférieure}, avec un signe de renvoi. --- \textbf{Analyse critique des vv. 1085–1086 par Louis Hauseux (\texttt{vv 642-fin-bis.tex}, ll. 4520, 4588–4612) :} --- \textbf{Deux notes scéniques marginales de \Rpc{} :} Au \textbf{v. 1098}, \lecon{ευρ\auDessus{ι}{π} ως περσευς} ; au \textbf{v. 1107}, \lecon{κ\auDessus{η}{δ} ως ανδρομ\auDessus{ε}{δ}}. --- \textbf{Prothèse de $\sigma$ au v. 1114 :} R écrit \lecon{σκύτο} pour \lecon{κύστο} (Scaliger, = $\text{κύσθον}$ en prononciation scythe), parallèle exact invoqué par Louis Hauseux pour sa conjecture τὸν κύσθον au v. 557. --- \textbf{v. 1082 :} \lecon{ὀτότυζ'} R (ὀλόφυγ' \lecon{Bibi} conj. in app.). --- \textbf{v. 1083 :} οῦτος οὗτος R : \lecon{οὖτος} bis Brunck \espApp λαλεῖς bis \Rac{} : \lecon{λαλῖς} bis $\text{R}^2, \Sigma$. --- \textbf{v. 1085 :} \textbf{σί κακόν; bis} R, $\Sigma^{\text{R}}$ (ἀντὶ τοῦ 'σοὶ κακόν', défendu par \lecon{Bibi} contre τί κακόν; Coulon). --- \textbf{v. 1086 :} π\inRas{ρ}{}ωτε bis \Rpc{} (πωτε bis \Mon{}) : \lecon{†πῶτε†} bis \lecon{Bibi} (rejetant l'interprétation πόθεν de $\Sigma^{\text{R}}$). --- \textbf{vv. 1090 et 1092 :} (Μνη.) edd. plerique : \lecon{(Χορ.)} \lecon{Bibi}, cf. $\Sigma^{\text{R}}$. --- \textbf{v. 1091 :} πλησίον αὕτη R : \lecon{πλήξασ' αὕτη} \lecon{Bibi}. --- \textbf{v. 1093 :} φεύγεις R : \lecon{πεύγεις} Enger \espApp \lecon{<(Εὐρ.) ποῖ ποῖ πέτει>} om. R : \textbf{suppl. \textbf{Bibi}}. --- \textbf{v. 1094 :} <(Τοξ.) οὐχ αἱρήσεις> om. R : \textbf{suppl. Brunck} ap. $\Sigma$ \espApp οὐκ' αἰρήσεις; R : \lecon{οὐχ αἱρήσεις.} Wilamowitz (\italique{Kleine Schriften} IV, p. 486 n. 1). --- \textbf{v. 1098 :} \lecon{ευρ\auDessus{ι}{π} ως περσευς} in mg. add. $\text{R}^{\text{pc}}, \text{M}$. --- \textbf{vv. 1102–1103 :} om. $\text{R}^1, \text{M}$ : \textbf{in mg. inferiore add. \Rdeux{}} \espApp \textbf{v. 1102 :} Γοργόνος R : \lecon{Γόργος} $\Sigma$, Bothe \espApp \textbf{v. 1103 :} κεφαλή R : \lecon{κεπαλή} Bothe. --- \textbf{v. 1107 :} \lecon{κ\auDessus{η}{δ} ως ανδρομ\auDessus{ε}{δ}} in mg. $\text{R}^{\text{pc}}, \text{M}$ \espApp κατοικτειρόν R : \lecon{κατοίκτιρόν} van Herwerden. --- \textbf{v. 1115 :} δεῦρο δεῦρο R (dittographie) : \lecon{δεῦρο} Biset. --- \textbf{v. 1125 :} δέμας R (métathèse de $\sigma$) : \lecon{δεσμά} Grynaeus. --- \textbf{v. 1127 :} ἀποκέκοιψ' \Rac{} (ἀποκεκόψοι \Rpc{}, ἀποκεκόψο \Mon{}) : \lecon{ἀποκεκόψι} Brunck \espApp πέπυσται R (faute récurrente de R, cf. v. 1227) : \lecon{πέπαισται} $\Sigma$.\par
}
\end{minipage}
\end{adjustbox}
\end{adjustwidth}
\newpage
\enTeteD{Codex Ravennas 429 (Biblioteca Classense) --- Fac-similé du f. 175v (vv. 1082–1127)}
\begin{adjustwidth}{-1.5cm}{-3.0cm}
\begin{center}
\includegraphics[width=16.5cm,height=20.5cm,keepaspectratio]{/workspaces/Thesmophories/dossier_ravennas_thesmophories/folios_individuels_ff162v-177r/27_f175v_classense177v_vv1082-1127.jpg}\par
\vspace{2pt}
{\scriptsize\textbf{Codex Ravennas 429, f. 175v (Classense 177v)} --- ΕΠΕΙΣΟΔΙΟΝ Ε' (Περσεύς), \textbf{vv. 1082–1127} ($\text{R}^1, \Rpc, \Rdeux$ (\Mon)).}
\end{center}
\end{adjustwidth}
\newpage
\clearpage\ifodd\value{page}\hbox{}\newpage\fi
\enTeteG{ΘΕΣΜΟΦΟΡΙΑΖΟΥΣΑΙ --- Texte grec établi (f. 175v, vv. 1082–1127)}
\enTeteD{LES THESMOPHORIES --- Traduction française (f. 175v, vv. 1082–1127)}
\begin{pages}
\begin{Leftside}
\beginnumbering
%% v. 1081c
\begin{groupe}%
\personnageGrec{Μνησίλοχος}%
\anap{\ligne{1081}\decalage{οἴμωζ᾽.†οἴμωζ᾽†.}ὀτότυζ᾽.\\\passeLigne}%
\end{groupe}%
%% v. 1081d
\begin{groupe}%
\personnageGrec{Εὐριπίδης}%
\anap{\ligne{1081}\decalage{οἴμωζ᾽.†οἴμωζ᾽†.ὀτότυζ᾽.}†ὀτότυζ᾽†.%\apparat{1082}{\lecon{ὀτότυζ᾽} R : ὀλόφυγ᾽ Bibi}%
\commentaire{1082}{ὀτότυζ᾽ --- ὀτότυζ᾽.}{Signalons à l'attention du lecteur -- que nous espérons exégète à ses heures -- qu'Aristophane invente le gentilé Ὀτοτύξιοι au vers 1043 des \italique{Oiseaux} par jeu de consonance avec Ὀλοφύξιοι (v. 1042), \italique{habitants d'Olophyxos}. Y a-t-il ici un jeu analogue ?}\\\passeLigne}%
\end{groupe}%
%% v. 1083a
\begin{groupe}%
\personnageGrec{Τοξότης}%
\anap{\ligne{1083}οὖτος τί λαλῖς;\apparat{1083}{οῦτος οὗτος R : \lecon{οὖτος} bis Brunck \espApp λαλεῖς bis \Rac : \lecon{λαλῖς} bis \Rdeux, Σ}%
\commentaire{1083}{λαλῖς.}{Scholie à propos de la graphie :%
\citComm{χωρὶς τοῦ ε γράφεται· ὁ γὰρ Σκύθης βαρβαρίζει.}}\\\passeLigne}%
\end{groupe}%
%% v. 1083b
\begin{groupe}%
\personnageGrec{Εὐριπίδης}%
\anap{\ligne{1083}\decalage{οὖτος τί λαλῖς;}†οὖτος τί λαλῖς;†\\\passeLigne}%
\end{groupe}%
%% v. 1084a
\begin{groupe}%
\personnageGrec{Τοξότης}%
\anap{πρυτάνεις καλέσω;\\\passeLigne}%
\end{groupe}%
%% v. 1084b
\begin{groupe}%
\personnageGrec{Εὐριπίδης}%
\anap{\ligne{1084}\decalage{πρυτάνεις καλέσω;}†πρυτάνεις καλέσω;†\\\passeLigne}%
\end{groupe}%
%% v. 1085a
\begin{groupe}%
\personnageGrec{Τοξότης}%
\anap{σί κακόν;%\apparat{1085}{σί bis R, sch. : τί Coulon}
\commentaire{1085}{σί κακόν;}{Commentaire du scholiaste :%
\citComm{ἀντὶ τοῦ `σοὶ κακόν'.}}%
\\\passeLigne}%
\end{groupe}%
%% v. 1085b
\begin{groupe}%
\personnageGrec{Εὐριπίδης}%
\anap{\ligne{1085}\decalage{σί κακόν;}†σί κακόν;†\\\passeLigne}%
\end{groupe}%
%% v. 1086a
\begin{groupe}%
\personnageGrec{Τοξότης}%
\anap{†πῶτε† τὸ πωνή;\apparat{1086}{π\inRas{ρ}{}ωτε bis \Rpc : πωτε bis \Mon}%
\commentaire{1086}{πρωτε το πωνη;}{Le scholiaste interprète ainsi le jargon scythe :%
\citComm{πόθεν ἡ φωνή; θαυμάζει τὴν ἠχώ.}%
On lui objectera que le \italique{Ravennas} comporte `πρωτε'. Et quand bien même on adopterait la correction `πωτε' (qu'une main postérieure a corrigé sur le manuscrit -- peut-être d'après la scholie elle-même), qu'il n'en resterait pas moins qu'à aucun autre endroit le Scythe ne confond un omicron ο avec un oméga ω\marqueur{1}.%
%962 γρωθώνη:σαπρὰγραῦς.οἱ δὲτὴνπαλαιὰνὄνονἢ σῦν
}%
\debutNote{1086}%
\note{1}{Voir la seconde note du commentaire au vers 1080 à propos du nombre déjà bien suffisant de fautes commises par le Scythe. 
%Faut-il aller chercher une explication du côté d'Hésychius ? (γ 962) :%
%\citComm{γρωθώνη· σαπρὰ γραῦς· οἱ δὲ τὴν παλαιὰν ὄνον ἢ σῦν.}
}%
\\\passeLigne}%
\end{groupe}%
%% v. 1086b
\begin{groupe}%
\personnageGrec{Εὐριπίδης}%
\anap{\ligne{1086}\decalage{†πῶτε† τὸ πωνή;}†πῶτε τὸ πωνή;†%\commentaire{1086}{}{Hésychius : πρόπωνα: vεὐχρατῆ. εὔφημα. πρόχειρα, ἕτοιμα, ἀνεμπόδιστα}
\\\passeLigne}%
\end{groupe}%
%% v. 1087a
\begin{groupe}%
\personnageGrec{Τοξότης}%
\anap{σὺ λαλῖς;\apparat{1087}{\lecon{λαλῖς} bis R : λαλεῖς bis \Rpc}\\\passeLigne}%
\end{groupe}%
%% v. 1087b
\begin{groupe}%
\personnageGrec{Εὐριπίδης}%
\anap{\ligne{1087}\decalage{σὺ λαλῖς;}†σὺ λαλῖς;†\\\passeLigne}%
\end{groupe}%
%% v. 1088a
\begin{groupe}%
\personnageGrec{Τοξότης}%
\anap{κλαύσαι.\apparat{1088}{κλαύσαιμι bis R : κλαύσει bis Brunck : \lecon{κλαύσαι} bis Rogers : κλαῦσ' ἔτι bis Austin}\\\passeLigne}%
\end{groupe}%
%% v. 1088b
\begin{groupe}%
\personnageGrec{Εὐριπίδης}%
\anap{\ligne{1088}\decalage{κλαύσαι.}†κλαύσαι.†%\commentaire{1088}{}{Hésychius 2856 κλᾶις:μοχλός 2881 χλαῦσιν:ἀγῶνα,ἄνθοςτι}%
\\\passeLigne}%
\end{groupe}%
%% v. 1089a
\begin{groupe}%
\personnageGrec{Τοξότης}%
\anap{κἀκκάσκι μοι;\apparat{1089}{κακκασκι bis R : κἀγκάσκις bis van Leeuwen (<ς> add. Fritzsche) : κἀγκάσκι bis Coulon : alii alia}%
\commentaire{1089}{κακκασκι μοι;}{$=$ καὶ ἐγχάσκει μοι. Cf. \italique{Guêpes}, vers 721 et le scholiaste : καταγελᾷς μοι.}\\\passeLigne}%
\end{groupe}%
%% v. 1089b
\begin{groupe}%
\personnageGrec{Εὐριπίδης}%
\anap{\ligne{1089}\decalage{κἀκκάσκι μοι;}†κἀκκάσκι μοι;†\\\passeLigne}%
\end{groupe}%
%% v. 1090
\begin{groupe}%
\personnageGrec{Χορός}%
\anap{μὰ Δί᾽ ἀλλὰ γυνὴ πλησίον αὕτη.\apparat{1090}{(Μνη.) edd. plerique : \lecon{(Χορ.)} Bibi, cf. Σ}%
\commentaire{1090}{μὰ Δί᾽ ἀλλὰ γυνὴ πλησίον αὕτη.}{C'est le chœur qui donne cette indication à l'archer désemparé, comme le suggérait timidement le scholiaste :%
\citComm{τοῦτο ὁ χορὸς ἢ ὁ γέρων.}%
Aucun éditeur cependant n'a suivi la première hypothèse du scholiaste, quoique la seule vraisemblable -- aussi agacé soit-il, Mnésiloque ne peut en aucun cas dénoncer son complice, son dernier espoir de survie.}\\\passeLigne}%
\end{groupe}%
%% v. 1091
\begin{groupe}%
\personnageGrec{Εὐριπίδης}%
\anap{πλήξασ᾿ αὕτη.\apparat{1091}{πλησίον αὕτη R : \lecon{πλήξασ᾿ αὕτη} Bibi}%
\commentaire{1091}{πλησίον αὕτη.}{Après un long calvaire parsemé de \latin{cruces desperationis}, il nous semble enfin saisir le sel ἐπικοκκάστρια de la réplique d'Écho, fournissant par la même occasion un indice sur l'action.\\
Sur l'indication du chœur, l'archer Scythe est en train d'inspecter \italique{à proximité} (πλησίον) de Mnésiloque. Il se rapproche effectivement d'Euripide (ce dernier est toujours camouflé derrière son manteau). 
Quand l'archer arrive à portée de bras, Euripide lui assène alors un coup : πλήξασ᾿ αὕτη.\\
Ce coup explique la fureur du Scythe, qui change de langage et passe à l'offensive : ποῦ ᾿στ᾽ ἠ μιαρά;}\\\passeLigne}%
\end{groupe}%
%% v. 1092-1094a
\begin{groupe}%
\personnageGrec{Τοξότης}%
\anap{ποῦ ᾿στ᾽ ἠ μιαρά;\apparat{1092}{'σθ' ἡ R : ᾿στ᾽ ἠ Brunck}\\\passeLigne}%
\end{groupe}%
%% v. 1092-1094a
\begin{groupe}%
\personnageGrec{Χορός}%
\anap{\decalage{ποῦ ᾿στ᾽ ἠ μιαρά;}καὶ δὴ φεύγει.\apparat{1092}{(Μνη.) edd. plerique : \lecon{(Χορ.)} Bibi}\\\passeLigne}%
\end{groupe}%
%% v. 1092-1094a
\begin{groupe}%
\personnageGrec{Τοξότης}%
\anap{ποῖ ποῖ πεύγεις;\apparat{1093}{φεύγεις R : \lecon{πεύγεις} Enger}\\\passeLigne}%
\end{groupe}%
%% v. 1092-1094a
\begin{groupe}%
\personnageGrec{Εὐριπίδης}%
\anap{\ligne{1092}\decalage{ποῖ ποῖ πεύγεις;}ποῖ ποῖ πέτει.\apparat{1093}{(Εὐρ.) ποῖ ποῖ πέτει om. R : \lecon{suppl.} Bibi}%
\commentaire{1093}{ποῖ ποῖ φεύγεις;}{Il y a longtemps que les scribes chargés de recopier cette scène n'en saisissaient plus le comique ; ne nous étonnons donc pas qu'elle soit la plus altérée de toute la comédie : il manque plusieurs répliques (cf. le vers 1094 ; où la lacune est certaine et peut-être comblée grâce à une ancienne scholie). C'est le cas aussi ici.\\
Euripide s'est approché de Mnésiloque pour le délivrer. Sur les indications du chœur (μὰ Δί᾽ ἀλλὰ γυνὴ πλησίον αὕτη), le Scythe a fini par repérer Euripide, et ce malgré son camouflage. C'est ce coup qu'il est bon de partir ; et Euripide de déloger sans trompette (καὶ δὴ φεύγει) mais non sans abandonner sa verve moqueuse. \\
Remettons-nous bien la scène : Euripide incarne la nymphe ailée Écho -- sa fuite est donc aérienne ; et sur scène, l'acteur est tiré dans les airs à l'aide de la μηχανή. La réplique ἐπικοκκάστρια d'Écho fait vraisemblablement allusion à cette \italique{envolée}. Or pour la circonstance, nous ne sommes pas tout-à-fait démunis puisque nous possédons une comédie entière, les \italique{Oiseaux}, avec des jeux théâtraux similaires.\\
En \italique{Oiseaux} 226, lorsque la Huppe se dispose à chanter, elle entonne : ἐποποποῖ ποποῖ, ποποποποῖ ποποῖ.\\
De même, lorsque le chœur au vers 310 répond à l'appel de la Huppe, c'est au cri de : ποποποποποπο ποῦ μ' ὃς ἐκάλεσε;\\
C'est ainsi qu'Aristophane transcrivait les onomatopées d'oiseaux\marqueur{1} ; on comprend mieux pourquoi lorqu'Iris descend du ciel au moyen de la μηχανή, Pisthétaire lui lance ces paroles ailées (vers 1199) : αὕτη σύ, ποῖ ποῖ ποῖ πέτει;\\
Il en va de même dans notre passage, alors que le Scythe poursuit Écho en train de s'envoler, il lui jette à la figure des cris d'oiseaux : ποῖ ποῖ πεύγεις; -- qui fait entendre ποῖ ποῖ πέτει -- qu'on rétablit donc comme réplique ἐπικοκκάστρια manquante d'Écho.}%
\debutNote{1093}%
\note{1}{Pour d'autres exemples de telles onomatopées chez Aristophane, voir l'article de Loredana Di Virgilio «~Some thoughts on humor and music in Aristophanes' plays~», in \italique{Mediterranean World}, 24 (2019), pp. 36 sqq.}\\\passeLigne}%
\end{groupe}%
%% v. 1092-1094a
\begin{groupe}%
\personnageGrec{Τοξότης}%
\anap{οὐ καιρήσεις.\apparat{1094}{(Τοξ.) οὐ καιρήσεις om. R : \lecon{suppl.} Brunck ap. Σ}%
\commentaire{1094}{<οὐ καιρήσεις.> --- οὐκ᾽ αἰρήσεις.}{Le texte manquant a été rétabli à partir de la scholie :%
\citComm{ἀντὶ τοῦ οὐ χαιρήσεις.}%
\italique{Confer} la note \italique{supra} au vers 1056.}\\\passeLigne}%
\end{groupe}%
%% v. 1094b
\begin{groupe}%
\personnageGrec{Εὐριπίδης}%
\anap{\ligne{1094}\decalage{οὐ καιρήσεις.}οὐχ αἱρήσεις.\apparat{1094}{οὐκ᾽ αἰρήσεις; R : \lecon{οὐχ αἱρήσεις.} Wilamowitz (\italique{Kleine Schriften} IV, p.486n.1)}\\\passeLigne}%
\end{groupe}%
%% v. 1095a
\begin{groupe}%
\personnageGrec{Τοξότης}%
\anap{ἔτι γὰρ γρύζεις;\\\passeLigne}%
\end{groupe}%
%% v. 1095b
\begin{groupe}%
\personnageGrec{Εὐριπίδης}%
\anap{\ligne{1095}\decalage{ἔτι γὰρ γρύζεις;}†ἔτι γὰρ γρύζεις;†\\\passeLigne}%
\end{groupe}%
%% v. 1096a
\begin{groupe}%
\personnageGrec{Τοξότης}%
\anap{λαβὲ τὴ μιαρά.\apparat{1096}{τῇ et τῆ R : \lecon{τὴ} bis Bothe}\\\passeLigne}%
\end{groupe}%
%% v. 1096b
\begin{groupe}%
\personnageGrec{Εὐριπίδης}%
\anap{\ligne{1096}\decalage{λαβὲ τὴ μιαρά.}†λαβὲ τὴ μιαρά.†\\\passeLigne}%
\end{groupe}%
%% v. 1097
\begin{groupe}%
\personnageGrec{Τοξότης}%
\anap{λάλο καὶ κατάρατο γύναικο.\\\passeLigne}%
\end{groupe}%
%% v. 1098-1102a
\begin{groupe}%
\personnageGrec{Εὐριπίδης}%
\typeVers{trim. iamb.}\trimIamb{ὦ θεοὶ τίν᾽ εἰς γῆν βαρβάρων ἀφίγμεθα\apparat{1098}{ευρ\auDessus{ι}{π} ως περσευς in mg. add. \Rpc, \Mon}%
\commentaire{1098-102}{ὦ θεοὶ ... κάρα κομίζων.}{La scholie à propos de la tirade de Persée\marqueur{1} :%
\citComm{εἰς Περσέα ἐξ Ἀνδρομέδας τρία τὰ πρῶτα. καὶ λοιπὸν (i.e. vv. 1101-2) ἐπέζευξεν τὰ ἑξῆς.}%
}%
\debutNote{1098-102}%
\note{1}{Fr. 9 de l'\italique{Andromède} dans l'édition Jouan-van Looy $=$ fr. 124 Kannicht.}\\
ταχεῖ πεδίλῳ; διὰ μέσου γὰρ αἰθέρος\\
τέμνων κέλευθον πόδα τίθημ᾽ ὑπόπτερον\\
Περσεὺς πρὸς Ἄργος ναυστολῶν τὸ Γοργόνος\apparat{1101}{ναυτολῶν R : \lecon{ναυστολῶν} Canini (\italique{navigans} iam Divus)}\\
κάρα κομίζων.\\\passeLigne}%
\end{groupe}%
%% v. 1102b-1103a
\begin{groupe}%
\personnageGrec{Τοξότης}%
\trimIamb{\ligne{1102}\decalage{κάρα κομίζων.}τί λέγι; τὴ Γόργος πέρι\apparat{vv. 1102-3}{om R, \Mon : in mg. inferiore add. \Rdeux \ligneApp{1102} Γοργόνος R : \lecon{Γόργος} Σ, Bothe \ligneApp{1103} κεφαλή R : \lecon{κεπαλή} Bothe}%
\commentaire{1102-3}{τί λέγι; τὴ Γόργος πέρι / τὸ γραμματέο σὺ τὴ κεπαλή;}{$=$ τί λέγεις; τοῦ Γόργου φέρεις / τοῦ γραμματέως σὺ τὴν κεφαλήν;\\
Et la scholie (plus que douteuse) :%
\citComm{ὁ δὲ Γόργος γραμματεὺς, ἀλλὰ καὶ βάρβαρος.}%
%\eff{mais la scansion ?..}
Il n'y a en grec qu'une différence d'accent entre Γοργών (la Gorgone) et Γόργων (Gorgon).}\\
τὸ γραμματέο σὺ τὴ κεπαλή;\\\passeLigne}%
\end{groupe}%
%% v. 1103b-1104a
\begin{groupe}%
%
\personnageGrec{Εὐριπίδης}%
\trimIamb{\ligne{1103}\decalage{τὸ γραμματέο σὺ τὴ κεπαλή;}τὴν Γοργόνος\\
ἔγωγε φημί.\\\passeLigne}%
\end{groupe}%
%% v. 1104b
\begin{groupe}%
\personnageGrec{Τοξότης}%
\trimIamb{\ligne{1104}\decalage{ἔγωγε φημί.}Γόργο τοι κἀγὼ λέγι.\\\passeLigne}%
\end{groupe}%
%% v. 1105-1106
\begin{groupe}%
\personnageGrec{Εὐριπίδης}%
\trimIamb{ἔα· τίν᾽ ὄχθον τόνδ᾽ ὁρῶ καὶ παρθένον%
\commentaire{1105}{ἔα· τίν᾽ ὄχθον τόνδ᾽ ὁρῶ καὶ παρθένον.}{On connaît le texte original d'Euripide grâce à une scholie\marqueur{1} :%
\citComm{πάλιν ἐξ Ἀνδρομέδας·\\
ἔα, τίν᾽ ὄχθον τόνδ᾽ ὁρῶ περίρρυτον\\
ἀφρῷ θαλάσσης παρθένου τ᾽ εἰκώ τίνα\\
$[$ ἐξ αὐτομόρφων λαΐνων τυκισμάτων,
σοφῆς ἄγαλμα χειρός; $]$ }%
}%
\debutNote{1105}%
\note{1}{Dixième fragment de l'\italique{Andromède} dans l'édition Jouan-van Looy $=$ fr. 125, 1-2 Kannicht. La scholie rapporte les seuls deux premiers vers ; le fragment a pu être complété en le recoupant avec d'autres témoignages.}\\
θεαῖς ὁμοίαν ναῦν ὅπως ὡρμισμένην;%
\commentaire{1106}{ναῦν ὅπως ὡρμισμένην;}{Quoique le scholiaste ne l'indique pas, le vers 1106 est probablement lui aussi repris de l'\italique{Andromède} d'Euripide. On trouve en tout cas la même image chez l'auteur d'\italique{Hercule Furieux} (vers 1094-1097), Hercule dont la situation était analogue à celle d'Andromède :%
\citComm{τί δεσμοῖς ναῦς ὅπως ὡρμισμένος\\
νεανίαν θώρακα καὶ βραχίονα,\\
πρὸς ἡμιθραύστῳ λαΐνῳ τυκίσματι\\
ἧμαι, νεκροῖσι γείτονας θάκους ἔχων;}%
}\\\passeLigne}%
\end{groupe}%
%% v. 1107-1108a
\begin{groupe}%
\personnageGrec{Μνησίλοχος}%
\trimIamb{ὦ ξένε κατοίκτιρόν με τὴν παναθλίαν,\apparat{1107}{κ\auDessus{η}{δ} ως ανδρομ\auDessus{ε}{δ} in mg. \Rpc, \Mon \espApp κατοικτειρόν R : \lecon{κατοίκτιρόν} van Herwerden}\\
λῦσόν με δεσμῶν.\\\passeLigne}%
\end{groupe}%
%% v. 1108b-1109
\begin{groupe}%
\personnageGrec{Τοξότης}%
\trimIamb{\ligne{1108}\decalage{λῦσόν με δεσμῶν.}οὐκὶ μὶ λαλῆσι σύ;\apparat{1108}{\lecon{μὶ λαλῆσι} R : μὴ λαλῆσι Brunck}\\
κατάρατο τολμᾷς ἀποτανουμένη λαλῖς;\apparat{1109}{λαλᾶς R : \lecon{λαλῖς;} Bergk (cf. vv. 1082 et 1087) : λαλᾷς; Hall-Geldart}\\\passeLigne}%
\end{groupe}%
%% v. 1110
\begin{groupe}%
\personnageGrec{Εὐριπίδης}%
\trimIamb{ὦ παρθέν᾽ οἰκτίρω σὲ κρεμαμένην ὁρῶν.\apparat{1110}{οἰκτείρω R : \lecon{οἰκτίρω} van Herwerden}\\\passeLigne}%
\end{groupe}%
%% v. 1111-1112a
\begin{groupe}%
\personnageGrec{Τοξότης}%
\trimIamb{οὐ παρτέν᾽ ἐστίν, ἀλλ᾽ ἀμαρτωλὴ γέρων\\
καὶ κλέπτο καὶ πανοῦργο.\\\passeLigne}%
\end{groupe}%
%% v. 1112b-1113
\begin{groupe}%
\personnageGrec{Εὐριπίδης}%
\trimIamb{\ligne{1112}\decalage{καὶ κλέπτο καὶ πανοῦργο.}ληρεῖς ὦ Σκύθα.\\
αὕτη γάρ ἐστιν Ἀνδρομέδα παῖς Κηφέως.\\\passeLigne}%
\end{groupe}%
%% v. 1114
\begin{groupe}%
\personnageGrec{Τοξότης}%
\trimIamb{σκέψαι τὸ κύστο· μή τι μικκὸν παίνεται;\apparat{1114}{σκύτο R : \lecon{κύστο} Scaliger grat. metr. : σῦκο Sommerstein \espApp μῆτι R : \lecon{μή τι} \Mon (\italique{numquid} Divus) \espApp μικτὸν R : \lecon{μικκὸν} Bentley : μικρὸ Enger \ligneApp{1115} δεῦρο δεῦρο R : \lecon{δεῦρο} Biset}%
\commentaire{1114}{σκέψαι τὸ σκύτο·}{Fait remarquable : l'erreur du scribe, copiant τὸ σκύτο en place de τὸ κύστο ($=$τὸν κύσθον), est la même qu'au vers 557 : métathèse du `s' (cf. le commentaire à ce vers). 
Le scholiaste devait avoir devant les yeux la bonne leçon qui précise : %
\citComm{δείκνυσιν αὐτῷ τὸ αἰδοῖον.}%
Second fait remarquable : en première rédaction, le \italique{Ravennas} faisait apparaître un `σ' supplémentaire au début d'un mot du vers 1187 : τὸ \textbf{σ}ποστίον à la de τὸ πόστιον ($=$ τὸ πόσθιον)...  on constate qu'il s'agit là-encore d'un organe sexuel -- mâle cette fois.}\\\passeLigne}%
\end{groupe}%
%% v. 1115-1118a
\begin{groupe}%
\personnageGrec{Εὐριπίδης}%
\trimIamb{φέρε δεῦρό μοι τὴν χεῖρ᾽, ἵν᾽ ἅψωμαι κόρης·\\
φέρε Σκύθ᾽· ἀνθρώποισι γὰρ νοσήματα\\
ἅπασίν ἐστιν· ἐμὲ δὲ καὐτὸν τῆς κόρης\\
ταύτης ἔρως εἴληφεν.\\\passeLigne}%
\end{groupe}%
%% v. 1118b-1120
\begin{groupe}%
%
\personnageGrec{Τοξότης}%
\trimIamb{\ligne{1118}\decalage{ταύτης ἔρως εἴληφεν.}οὐ ζηλῶσί σε·\\
ἀτὰρ εἰ τὸ πρωκτὸ δεῦρο περιεστραμμένον,\apparat{1119}{τῶ πρωκτῶ R : \lecon{τὸ πρωκτὸ} Biset}%
%\commentaire{1019}{τῶ πρωκτῶ}{Le prokton peut donc être (mal-)employé au pluriel. \cf la scène où Mnésiloque a le feu aux fesses.}
\\
οὐκ ἐπτόνησά σ᾽ αὐτὸ πυγίζεις ἄγων.\apparat{1120}{ἐπτόνησας R : \lecon{ἐπτόνησά σ᾽} Bentley}\\\passeLigne}%
\end{groupe}%
%% v. 1121-1122
\begin{groupe}%
\personnageGrec{Εὐριπίδης}%
\trimIamb{τί δ᾽ οὐκ ἐᾷς λύσαντά μ᾽ αὐτὴν ὦ Σκύθα\\
πεσεῖν ἐς εὐνὴν καὶ γαμήλιον λέχος;\apparat{1122}{ἐστ᾽ R : ἐς τ᾽ \Mon : \lecon{ἐς} Grynaeus}\\\passeLigne}%
\end{groupe}%
%% v. 1123-1124
\begin{groupe}%
\personnageGrec{Τοξότης}%
\trimIamb{εἰ σπόδρ᾽ ἐπιτυμεῖς τὴ γέροντο πύγισο,\\
τὴ σανίδο τρήσας ἐξόπιστο πρώκτισον.\apparat{1124}{ἐξόπισθο R : \lecon{ἐξόπιστο} Daubuz}\\\passeLigne}%
\end{groupe}%
%% v. 1125a
\begin{groupe}%
\personnageGrec{Εὐριπίδης}%
\trimIamb{μὰ Δί᾽ ἀλλὰ λύσω δεσμά.\testimonia{1125}{(δεσμά) Moeris δ 28}\apparat{1125}{δέμας R : \lecon{δεσμά} Grynaeus}\\\passeLigne}%
\end{groupe}%
%% v. 1125b
\begin{groupe}%
\personnageGrec{Τοξότης}%
\trimIamb{\ligne{1125}\decalage{μὰ Δί᾽ ἀλλὰ λύσω δεσμά.}μαστιγῶ σ᾽ ἄρα.\\\passeLigne}%
\end{groupe}%
%% v. 1126a
\begin{groupe}%
\personnageGrec{Εὐριπίδης}%
\trimIamb{καὶ μὴν ποιήσω τοῦτο.\\\passeLigne}%
\end{groupe}%
%% v. 1126b-1127
\begin{groupe}%
\personnageGrec{Τοξότης}%
\trimIamb{\ligne{1126}\decalage{καὶ μὴν ποιήσω τοῦτο.}τὸ κεπαλή σ᾽ ἄρα\apparat{1126}{τὸ κεπαλῆς (ex καὶ παλῆς) \Rpc, \Mon : \lecon{τὸ κεπαλή σ᾽} Brunck \ligneApp{1127} ἀποκέκοιψ᾽ \Rac : ἀποκεκόψοι \Rpc : ἀποκεκόψο \Mon : \lecon{ἀποκεκόψι} Brunck}\\
τὸ ξιπομάκαιραν ἀποκεκόψι τουτοΐ.%\apparat{1127}{ἀποκέκοιψ᾽ \Rac : ἀποκεκόψοι \Rpc : ἀποκεκόψο \Mon : \lecon{ἀποκεκόψι} Brunck}%
\commentaire{1127}{τὸ ξιπομάκαιραν.}{La harpé de Persée. (Note de Willems.)}\\\passeLigne}%
\end{groupe}%
\endnumbering
\end{Leftside}
\begin{Rightside}
\beginnumbering
%% v. 1081c
\begin{groupe}%
\personnageFr{Mnésiloque}%
\anapFr{\ligne{1081}\espace{2.3cm}Foin de toi.\\\passeLigne}%
\end{groupe}%
%% v. 1081d
\begin{groupe}%
\personnageFr{Euripide}%
\anapFr{\ligne{1081}\espace{3.5cm}Foin de toi.\\\passeLigne}%
\end{groupe}%
%% v. 1083a
\begin{groupe}%
\personnageFr{Archer}%
\anapFr{\ligne{1083}Holà ! qu'as-tu à bavarder ?\\\passeLigne}%
\end{groupe}%
%% v. 1083b
\begin{groupe}%
\personnageFr{Euripide}%
\anapFr{\ligne{1083}\espace{2.6cm}Holà ! qu'as-tu à bavarder ?\\\passeLigne}%
\end{groupe}%
%% v. 1084a
\begin{groupe}%
\personnageFr{Archer}%
\anapFr{J'appellerai les prytanes.\\\passeLigne}%
\end{groupe}%
%% v. 1084b
\begin{groupe}%
\personnageFr{Euripide}%
\anapFr{\ligne{1084}\espace{2.7cm}J'appellerai les prytanes.\\\passeLigne}%
\end{groupe}%
%% v. 1085a
\begin{groupe}%
\personnageFr{Archer}%
\anapFr{Qu'est-ce qui arrive ?\\\passeLigne}%
\end{groupe}%
%% v. 1085b
\begin{groupe}%
\personnageFr{Euripide}%
\anapFr{\ligne{1085}\espace{1.6cm}Qu'est-ce qui arrive ?\\\passeLigne}%
\end{groupe}%
%% v. 1086a
\begin{groupe}%
\personnageFr{Archer}%
\anapFr{D'où part cette voix ?\\\passeLigne}%
\end{groupe}%
%% v. 1086b
\begin{groupe}%
\personnageFr{Euripide}%
\anapFr{\ligne{1086}\espace{2.4cm}D'où part cette voix ?\\\passeLigne}%
\end{groupe}%
%% v. 1087a
\begin{groupe}%
\personnageFr{Archer}%
\anapFr{Est-ce toi qui bavardes?\\\passeLigne}%
\end{groupe}%
%% v. 1087b
\begin{groupe}%
\personnageFr{Euripide}%
\anapFr{\ligne{1087}\espace{1.8cm}Est-ce toi qui bavardes?\\\passeLigne}%
\end{groupe}%
%% v. 1088a
\begin{groupe}%
\personnageFr{Archer}%
\anapFr{Gare à toi.\\\passeLigne}%
\end{groupe}%
%% v. 1088b
\begin{groupe}%
\personnageFr{Euripide}%
\anapFr{\ligne{1088}\espace{1.4cm}Gare à toi.\\\passeLigne}%
\end{groupe}%
%% v. 1089a
\begin{groupe}%
\personnageFr{Archer}%
\anapFr{Tu te gausses de moi.\\\passeLigne}%
\end{groupe}%
%% v. 1089b
\begin{groupe}%
\personnageFr{Euripide}%
\anapFr{\ligne{1089}\espace{2.2cm}Tu te gausses de moi.\\\passeLigne}%
\end{groupe}%
%% v. 1090
\begin{groupe}%
\personnageFr{Chœur}%
\anapFr{ Non, par Zeus : c'est cette femme, là tout près.\\\passeLigne}%
\end{groupe}%
%% v. 1091
\begin{groupe}%
\personnageFr{Euripide}%
\anapFr{Là tout près.\\\passeLigne}%
\end{groupe}%
%% v. 1092-1094a
\begin{groupe}%
\personnageFr{Archer}%
\anapFr{Où est la canaille ?\\\passeLigne}%
\end{groupe}%
%% v. 1092-1094a
\begin{groupe}%
\personnageFr{Chœur}%
\anapFr{\espace{2.2cm}La voilà qui se sauve.\\\passeLigne}%
\end{groupe}%
%% v. 1092-1094a
\begin{groupe}%
\personnageFr{Archer}%
\anapFr{Où, où te sauves-tu ?\\\passeLigne}%
\end{groupe}%
%% v. 1092-1094a
\begin{groupe}%
\personnageFr{Euripide}%
\anapFr{\ligne{1092}\decalage{Où, où te sauves-tu ?}Où, où voles-tu ?\\\passeLigne}%
\end{groupe}%
%% v. 1092-1094a
\begin{groupe}%
\personnageFr{Archer}%
\anapFr{Il t'en cuira.\\\passeLigne}%
\end{groupe}%
%% v. 1094b
\begin{groupe}%
\personnageFr{Euripide}%
\anapFr{\ligne{1094}\espace{2.1cm}Tu ne m'auras.\\\passeLigne}%
\end{groupe}%
%% v. 1095a
\begin{groupe}%
\personnageFr{Archer}%
\anapFr{Tu maronnes encore ?\\\passeLigne}%
\end{groupe}%
%% v. 1095b
\begin{groupe}%
\personnageFr{Euripide}%
\anapFr{\ligne{1095}\espace{2.5cm}Tu maronnes encore ?\\\passeLigne}%
\end{groupe}%
%% v. 1096a
\begin{groupe}%
\personnageFr{Archer}%
\anapFr{Arrête, la canaille.\\\passeLigne}%
\end{groupe}%
%% v. 1096b
\begin{groupe}%
\personnageFr{Euripide}%
\anapFr{\ligne{1096}\espace{2.3cm}Arrête, la canaille.\\\passeLigne}%
\end{groupe}%
%% v. 1097
\begin{groupe}%
\personnageFr{Archer}%
\anapFr{La bavarde, la maudite femme.\\\passeLigne}%
\end{groupe}%
%% v. 1098-1102a
\begin{groupe}%
\personnageFr{Euripide}%
\trimIambFr{(En Persée :) « En quel pays barbare, ô dieux, ces talonnières\\ 
ont-elles amené Persée ? Au milieu de l'éther\\
me frayant un chemin, prestement je fends l'air d'un pied ailé,\\
portant à la cité argienne la tête de\\
Gorgone. »\\\passeLigne}%
\end{groupe}%
%% v. 1102b-1103a
\begin{groupe}%
\personnageFr{Archer}%
\trimIambFr{\ligne{1102}\espace{2.2cm}Que dis-tu? Tu portes la tête de Gorgon\\
le scribe ?\\\passeLigne}%
\end{groupe}%
%% v. 1103b-1104a
\begin{groupe}%
\personnageFr{Euripide}%
\trimIambFr{\ligne{1103}\espace{4.4cm}De Gorgone, ai-je dit.\\
ἔγωγε φημί.\\\passeLigne}%
\end{groupe}%
%% v. 1104b
\begin{groupe}%
\personnageFr{Archer}%
\trimIambFr{\ligne{1104}\espace{2cm}Gorgon, parfaitement, je dis de même.\\\passeLigne}%
\end{groupe}%
%% v. 1105-1106
\begin{groupe}%
\personnageFr{Euripide}%
\trimIambFr{« Ciel! Quelle est cette vierge aux déesses pareille amarrée à ce roc,\\
comme un navire au port ? »\\\passeLigne}%
\end{groupe}%
%% v. 1107-1108a
\begin{groupe}%
\personnageFr{Mnésiloque}%
\trimIambFr{« Étranger, prends pitié de moi, la désolée.\\
Détache mes liens. »\\\passeLigne}%
\end{groupe}%
%% v. 1108b-1109
\begin{groupe}%
\personnageFr{Archer}%
\trimIambFr{\ligne{1108}\espace{2.8cm}Veux-tu ne pas bavarder ?\\
Maudit, oses-tu bien ? Tu vas mourir et tu bavardes.\\\passeLigne}%
\end{groupe}%
%% v. 1110
\begin{groupe}%
\personnageFr{Euripide}%
\trimIambFr{« Ô vierge, une pitié me prend quand je te vois suspendue à ce roc. »\\\passeLigne}%
\end{groupe}%
%% v. 1111-1112a
\begin{groupe}%
\personnageFr{Archer}%
\trimIambFr{Ce n'est pas une vierge, c'est un vieux coquin,\\
et un voleur, et un fripon.\\\passeLigne}%
\end{groupe}%
%% v. 1112b-1113
\begin{groupe}%
\personnageFr{Euripide}%
\trimIambFr{\ligne{1112}\espace{3.9cm}Tu plaisantes, ô Scythe.\\
C'est Andromède, fille de Cépheus.\\\passeLigne}%
\end{groupe}%
%% v. 1114
\begin{groupe}%
\personnageFr{Archer}%
\trimIambFr{Regarde son cas : te paraît-il petit ?\\\passeLigne}%
\end{groupe}%
%% v. 1115-1118a
\begin{groupe}%
\personnageFr{Euripide}%
\trimIambFr{Donne-moi la main, que j'atteigne à la jeune fille.\\
Voyons, Scythe « chaque homme a ses faiblesses » :\\
la mienne est d'être féru\\
d'amour pour cette jeune fille.\\\passeLigne}%
\end{groupe}%
%% v. 1118b-1120
\begin{groupe}%
\personnageFr{Archer}%
\trimIambFr{\ligne{1118}\espace{3.3cm}Je ne suis pas jaloux.\\
S'il tournait les fesses par ici,\\
je ne te refuserais pas d'en passer ta fantaisie.\\\passeLigne}%
\end{groupe}%
%% v. 1121-1122
\begin{groupe}%
\personnageFr{Euripide}%
\trimIambFr{Que ne me permets-tu de la délier, ô Scythe\\
« et de me laisser choir dans le lit nuptial »\\\passeLigne}%
\end{groupe}%
%% v. 1123-1124
\begin{groupe}%
\personnageFr{Archer}%
\trimIambFr{Si tu as tant envie de jouir de ce vieillard,\\
perce-moi l'ais par derrière et enfourne.\\\passeLigne}%
\end{groupe}%
%% v. 1125a
\begin{groupe}%
\personnageFr{Euripide}%
\trimIambFr{Non, par Zeus, je briserai ses liens.\\\passeLigne}%
\end{groupe}%
%% v. 1125b
\begin{groupe}%
\personnageFr{Archer}%
\trimIambFr{\ligne{1125}\espace{3.7cm}Alors je te fouette.\\\passeLigne}%
\end{groupe}%
%% v. 1126a
\begin{groupe}%
\personnageFr{Euripide}%
\trimIambFr{N'importe, je vais le faire.\\\passeLigne}%
\end{groupe}%
%% v. 1126b-1127
\begin{groupe}%
\personnageFr{Archer}%
\trimIambFr{\ligne{1126}\espace{3.2cm}Alors je te coupe la tête\\
avec ce sabre.\\\passeLigne}%
\end{groupe}%
\endnumbering
\end{Rightside}
\end{pages}
\Pages
\clearpage\ifodd\value{page}\hbox{}\newpage\fi
\enTeteG{ΧΟΡΙΚΟΝ / ΕΞΟΔΟΣ --- Codex Ravennas 429, f. 176r (Classense 178r) --- vv. 1128–1168}
\begin{adjustwidth}{-1.5cm}{-3.0cm}
\begin{adjustbox}{max totalheight=21.0cm, max width=16.5cm, keepaspectratio}
\begin{minipage}{16.5cm}
{\fontsize{7.2pt}{8.5pt}\selectfont
\noindent\textbf{Folio 28 --- Codex Ravennas 429, \texttt{f. 176r} (\texttt{Classense 178r}) : ΧΟΡΙΚΟΝ / ΕΞΟΔΟΣ, vv. 1128–1168}\hfill \textit{Témoins :} $\text{R}^1, \Rpc, \Rdeux$ (\Mon)\par\vspace{2pt}
\hrule height 0.4pt\vspace{2.5pt}
}
{\fontsize{8.1pt}{9.5pt}\selectfont
\ligneRav{1128}{\auDessus{εὐ}{ρ}}{αἰαῖ· τί δράσω; πρὸς τίνας στρεφθῶ λόγους;}{}
\ligneRav{1129}{}{ἀλλ᾽ οὐκ \difference{ἂν ἐνδέξαιτο} βάρβαρος φύσις·}{\textbf{1129} ἂν δέξαιτο R : ἂν ἀνδέξαιτο Ellebodius : \lecon{ἂν ἐνδέξαιτο} Küster}
\ligneRav{1130}{}{σκαιοῖσι γάρ τοι καινὰ προσφέρων σοφὰ}{}
\ligneRav{}{}{μάτην ἀναλίσκοις ἄν, ἀλλ᾽ ἄλλην τινὰ}{}
\ligneRav{1132}{}{\difference{τούτῳ} πρέπουσαν μηχανὴν προσοιστέον·}{\textbf{1132} τοῦτο R : \lecon{τούτῳ} Zanetti}
\ligneRav{1133}{τοξ/}{μιαρὸς ἀλώπηξ, \difference{οἶον} \difference{ἐπιτήκιζί} μοι·}{\textbf{1133} οἷον R : \lecon{οἶον} Brunck \espApp ἐπιτηκίζει R : \lecon{ἐπιτήκιζί} anonymus Parisinus}
\ligneRav{}{μ\auDessus{ν}{Η}}{μέμνησο Περσεῦ μ᾽ ὡς καταλείπεις ἀθλίαν·}{}
\ligneRav{1135}{τοξ/}{ἔτι γὰρ σὺ \difference{τῇ μάστιγ᾽ ἂν} (τῆ μαστ\inRas{η}{ι}γᾶν) ἐπεθύμεις λαβεῖν;}{\textbf{1135} τῆ μαστιγᾶν (ex μαστηγαν) \Rpc \espApp λαβεῖν· R : \lecon{λαβεῖν;} Scaliger}
\ligneRav{}{χορ/}{\hspace{1.1em}Παλλάδα τὴν φιλόχορον ἐμοὶ}{}
\ligneRav{1138}{}{\hspace{1.1em}δεῦρο καλεῖν νόμος \difference{εἰς} χορόν,}{\textbf{1138} ἐς R : \lecon{εἰς} von Velsen}
\ligneRav{1139}{}{\hspace{1.1em}παρθένον ἄζυγα \difference{κούρην},}{\textbf{1139} κόρην R : \lecon{κούρην} Bothe}
\ligneRav{1140}{}{\hspace{1.1em}ἣ πόλιν ἡμετέραν ἔχει}{}
\ligneRav{}{}{\hspace{1.1em}καὶ κράτος φανερὸν μόνη}{}
\ligneRav{}{}{\hspace{1.1em}κλῃδοῦχός τε καλεῖται·}{}
\ligneRav{}{}{\hspace{1.1em}φάνηθ᾽ ὦ τυράννους στυγοῦσ᾽ ὥσπερ εἰκός·}{}
\ligneRav{1145}{}{\hspace{1.1em}δῆμός τοί σε καλεῖ γυναι-}{}
\ligneRav{}{}{\hspace{1.1em} κῶν· ἔχουσα δέ μοι μόλοις}{}
\ligneRav{}{}{\hspace{1.1em}εἰρήνην φιλέορτον·}{}
\ligneRav{1148}{}{\hspace{1.1em}\difference{ἥκετέ τ᾽} εὔφρονες ἵλαοι,}{\textbf{1148} ἥκετ᾽ R : \lecon{ἥκετέ τ᾽} Enger : ἥκετε δ᾽ Hermann}
\ligneRav{}{}{\hspace{1.1em}πότνιαι, ἄλσος ἐς ὑμέτερον,}{}
\ligneRav{1150}{}{\hspace{1.1em}οὗ δὴ ἀνδράσιν οὐ \difference{θεμιτεύετον} εἰσορᾶν}{\textbf{1150-1} \lecon{οὗ δὴ ἀνδράσιν} R : οὗ δῆτ' ἀνδράσιν Burges : ἄνδρας ἵν(α) Hermann \espApp θεμιτὸν R : θέμις Burges : θεμιτ' Bothe : \lecon{θεμιτεύετον} Bibi}
\ligneRav{}{}{\hspace{1.1em}ὄργια σεμνὰ θεαῖν ἵνα λαμπάσι·}{}
\ligneRav{}{}{\hspace{1.1em}φαίνετον ἄμβροτον ὄψιν·}{}
\ligneRav{1155}{}{\hspace{1.1em}μόλετον, ἔλθετον, ἀντόμεθ᾽ ὦ}{}
\ligneRav{1156}{}{\hspace{1.1em}Θεσμοφόρω \difference{πολυποτνία},}{\textbf{1156} πολὺ πότνια R : \lecon{πολυποτνία} Hermann}
\ligneRav{}{}{\hspace{1.1em}εἰ καὶ πρότερόν ποτ᾽ ἐπηκόω ἤλθετον}{}
\ligneRav{1159}{}{\hspace{1.1em}νῦν ἀφίκεσθον †ἱκετεύομεν ἐνθάδ᾽ ἡμῖν†·}{\textbf{1159} \lecon{ἐνθάδ᾽ ἡμῖν} R : ἐνθάδε χἠμῖν Hermann}
\ligneRav{1160}{\auDessus{εὐ}{ρ}}{γυναῖκες εἰ βούλεσθε τὸν λοιπόν χρόνον}{}
\ligneRav{}{}{σπονδὰς ποιήσασθαι πρὸς ἐμέ, νυνὶ πάρα,}{}
\ligneRav{}{}{ἐφ᾽ ᾧτ᾽ ἀκοῦσαι μηδὲν ὑπ᾽ ἐμοῦ μηδαμὰ}{}
\ligneRav{}{}{κακὸν τὸ λοιπόν. ταῦτ᾽ ἐπικηρυκεύομαι·}{}
\ligneRav{}{χορ/}{χρείᾳ δὲ ποίᾳ τόνδ᾽ ἐπεσφέρεις λόγον;}{}
\ligneRav{1165}{\auDessus{εὐ}{ρ}}{ὅδ᾽ ἐστὶν \difference{οὑν} τῇ σανίδι κηδεστὴς ἐμός·}{\textbf{1165} οὖν R : \lecon{οὑν} Biset}
\ligneRav{1166}{}{ἢν οὖν κομίσωμαι τοῦτον, οὐδὲν μή ποτε}{\textbf{1166} \lecon{μή} R : μου Bentley}
\ligneRav{1167}{}{κακῶς \difference{ἀκούσητ᾽}· ἢν δὲ μὴ \difference{πίθησθέ} μοι,}{\textbf{1167} ἀκούσαιτ' R : \lecon{ἀκούσητ᾽} Elmsley (ad \italique{Ach.} v.295) \espApp πείθησθέ R : \lecon{πίθησθέ} Hirschig (ad \italique{Vesp} v.155)}
\ligneRav{}{}{ἃ νῦν ὑποικουρεῖτε τοῖσιν ἀνδράσιν}{}
}
\vspace{3pt}\noindent\rule{4.5cm}{0.4pt}\par\vspace{1.5pt}
{\fontsize{7.0pt}{8.3pt}\selectfont
\noindent\textbf{Loci paralleli (vv. 1128–1168) ---} \textbf{1139} (ἄζυγα) \Phryn 100.3\par
}
\vspace{2.5pt}\noindent\rule{4.5cm}{0.4pt}\par\vspace{1.5pt}
{\fontsize{6.9pt}{8.2pt}\selectfont
\noindent\textbf{Suite (en commentaire) de l'apparat critique \& notes paléographiques (\texttt{f. 176r}) ---} \textbf{Dittographie anticipatrice exponctuée après le v. 1143 :} Entre les vv. 1143 et 1144, \Rac{} avait écrit par anticipation déformée du v. 1144 (στυγοῦσ' ὥσπερ τὸν τύραννον) les mots \lecon{στυγνὰς ὡσσε} ; \Rpc{} les a supprimés (\italique{del.}), mais le copiste de \Mon{} les a néanmoins transcrits sous la forme \lecon{στυγνὰς ὡς ἐ}. --- \textbf{Glose interne et conjecture métrique de \textbf{Bibi} au v. 1151 :} R comporte la glose interne ὅπου (expliquant un ἵνα ou οὗ). Là où R porte \lecon{οὗ δὴ ἀνδράσιν οὐ θεμιτὸν} (déficient de deux syllabes dans la responsion lyrique), \textbf{Louis Hauseux (\textbf{Bibi}) conserve intact οὗ δὴ ἀνδράσιν de R et restitue le duel verbal οὐ θεμιτεύετον} (« où toutes deux [les Déesses] interdisent par leur loi sacrée aux hommes de voir leurs orgies... », haplographie de la finale -εύετον en -όν). --- \textbf{vv. 1150–1151 :} \lecon{οὗ δὴ ἀνδράσιν} R (avec glose interne ὅπου ; οὗ δῆτ' ἀνδράσιν Burges ; ἄνδρας ἵν(α) Hermann) \espApp θεμιτὸν R (θέμις Burges ; θεμιτ' Bothe) : \lecon{θεμιτεύετον} \lecon{Bibi}. --- \textbf{v. 1153 :} λαμπάσι R : \lecon{λαμπάσι.} \lecon{Bibi}.\par
}
\end{minipage}
\end{adjustbox}
\end{adjustwidth}
\newpage
\enTeteD{Codex Ravennas 429 (Biblioteca Classense) --- Fac-similé du f. 176r (vv. 1128–1168)}
\begin{adjustwidth}{-1.5cm}{-3.0cm}
\begin{center}
\includegraphics[width=16.5cm,height=20.5cm,keepaspectratio]{/workspaces/Thesmophories/dossier_ravennas_thesmophories/folios_individuels_ff162v-177r/28_f176r_classense178r_vv1128-1168.jpg}\par
\vspace{2pt}
{\scriptsize\textbf{Codex Ravennas 429, f. 176r (Classense 178r)} --- ΧΟΡΙΚΟΝ / ΕΞΟΔΟΣ, \textbf{vv. 1128–1168} ($\text{R}^1, \Rpc, \Rdeux$ (\Mon)).}
\end{center}
\end{adjustwidth}
\newpage
\clearpage\ifodd\value{page}\hbox{}\newpage\fi
\enTeteG{ΘΕΣΜΟΦΟΡΙΑΖΟΥΣΑΙ --- Texte grec établi (f. 176r, vv. 1128–1168)}
\enTeteD{LES THESMOPHORIES --- Traduction française (f. 176r, vv. 1128–1168)}
\begin{pages}
\begin{Leftside}
\beginnumbering
%% v. 1128-1132
\begin{groupe}%
\personnageGrec{Εὐριπίδης}%
\trimIamb{αἰαῖ· τί δράσω; πρὸς τίνας στρεφθῶ λόγους;\\
ἀλλ᾽ οὐκ ἂν ἐνδέξαιτο βάρβαρος φύσις.\apparat{1129}{ἂν δέξαιτο R : ἂν ἀνδέξαιτο Ellebodius : \lecon{ἂν ἐνδέξαιτο} Küster}\\
σκαιοῖσι γάρ τοι καινὰ προσφέρων σοφὰ%
\commentaire{1130}{σκαιοῖσι γάρ τοι καινὰ προσφέρων σοφὰ.}{Cf. le reproche qu'Euripide mit dans la bouche de Médée dans la tragédie de ce nom (vers 298) :%
\citComm{σκαιοῖσι μὲν γὰρ καινὰ προσφέρων σοφὰ.}}\\
μάτην ἀναλίσκοις ἄν, ἀλλ᾽ ἄλλην τινὰ\\
τούτῳ πρέπουσαν μηχανὴν προσοιστέον.\apparat{1132}{τοῦτο R : \lecon{τούτῳ} Zanetti% (\italique{huic} Divus)
}\\\passeLigne}%
\end{groupe}%
%% v. 1133
\begin{groupe}%
\personnageGrec{Τοξότης}%
\trimIamb{μιαρὸς ἀλώπηξ, οἶον ἐπιτήκιζί μοι.\apparat{1133}{οἷον R : \lecon{οἶον} Brunck \espApp ἐπιτηκίζει R : \lecon{ἐπιτήκιζί} anonymus Parisinus}%
\commentaire{1133}{οἷον ἐπιτηκίζει μοι}{$=$ οἷον ἐπιθήκιζέ μοι. À propos des images animales du trompeur : du renard, ἡ ἀλώπηξ, et du singe, ὁ πίθηκος, vus comme les représentant de la ruse et de la fourbe, le grand Taillardat a consacré deux paragraphes auxquels je renvoie le lecteur : §§405-406. Notons que sur le même modèle que πιθηκίζειν (qui se rencontre aussi au vers 1290 des \italique{Guêpes}), il existe chez Aristophane le verbe ἀλωπεκίζειν (\italique{Guêpes}, vers 1241).}\\\passeLigne}%
\end{groupe}%
%% v. 1134
\begin{groupe}%
\personnageGrec{Μνησίλοχος}%
\trimIamb{μέμνησο Περσεῦ μ᾽ ὡς καταλείπεις ἀθλίαν.\\\passeLigne}%
\end{groupe}%
%% v. 1135
\begin{groupe}%
\personnageGrec{Τοξότης}%
\trimIamb{ἔτι γὰρ σὺ τὴ μάστιγαν ἐπιτυμεῖς λαβεῖν;\apparat{1135}{τῆ μαστιγᾶν (ex μαστηγαν) \Rpc \espApp λαβεῖν· R : \lecon{λαβεῖν;} Scaliger %(\italique{consequi?} iam Ellebodius in vers. Lat.)
}\\\passeLigne}%
\end{groupe}%
%% v. 1136-1159
\begin{groupe}%
%
\personnageGrec{Χορός}%
\typeVers{ibyc.}\chantChoeur{Παλλάδα τὴν φιλόχορον ἐμοὶ%
\commentaire{1136-59}{Chant ibycéen du chœur.}{Il a été beaucoup discuté du caractère strophique ou a-strophique de ce chant. Si la balance des commentateurs penche plutôt en faveur d'un chant astrophique\marqueur{1}, il est tentant, à l'instar d'Hermann, von Velsen ou Prato, de décomposer ce chant en deux parties avec pour chacune la division classique du chœur en deux demi-chœurs (strophe -- antistrophe). Voici la structure que je propose :%
\begin{alignat*}{4}
\text{1\expo{ère} partie } & \text{ : invocation d'Athéna, } && \text{ vv. 1140-1142 } &&& \text{ \ } = \text{ vv. 1145-1147 ; }\\
\text{2\expo{nde} partie } & \text{ : invocation des Thesmophores, } && \text{ vv. 1148-1153 } &&& \simeq \text{ vv. 1155-1159.}
\end{alignat*}
(Nous n'avons pas cependant forcé la concordance métrique au point de retoucher le texte lorsqu'il ne s'y prêtait que de mauvaise grâce ; si les vers 1140 ou 1148 peuvent se voir comme des glyconéens à base trissyllabique\marqueur{2}, ainsi, à la rigueur, que le vers 1155\marqueur{3}, il subsiste un problème de correspondance entre le vers 1153, ὄργια σεμνὰ θεαῖν ἵνα λαμπάσι, et le vers 1159, νῦν ἀφίκεσθον ἱκετεύομεν ἐνθάδ᾽ ἡμῖν.)\\
Mis de cette façon en ordre, le chant est entrecoupé entre chaque paire strophique d'un intermède, μεσῳδός, commençant par un impératif de φαίνω : aux vers 1143-4, φάνηθ᾽ ὦ τυράννους στυγοῦσ᾽ ὥσπερ εἰκός, et au vers 1154, φαίνετον ἄμβροτον ὄψιν.}%
\debutNote{1136-59}%
\note{1}{L'analyse la plus poussée de ce chant est celle de Lætitia, \italique{The Songs of Aristophanes}, pp. 446-453 ; on pourra aussi consulter les non moins magistraux ouvrages : \italique{The Lyric Metres of Greek Drama}\expo{2} (1968) pp. 166-167 d'Amy M. Dale ainsi que \italique{The Verse of Greek Comedy} (1912), §387, de White. White ne perçoit de répons antistrophique que pour les $6 = 3 + 3$ premiers vers.}%
\note{2}{Cf. Euripide, \italique{Iphigénie en Aulis}, vers 759.}%
\note{3}{Voir ce qu'en dit Amy, \italique{The Lyric Metres of Greek Drama}\expo{2}, p. 166 n.2.}}
\end{groupe}%
%% v. 1136-1159
\begin{groupe}%
\typeVers{ibyc.}\chantChoeur{δεῦρο καλεῖν νόμος εἰς χορόν,\apparat{1138}{ἐς R : \lecon{εἰς} von Velsen}}
\end{groupe}%
%% v. 1136-1159
\begin{groupe}%
\typeVers{ibyc\cat}\chantChoeur{\ligne{1139}παρθένον ἄζυγα κούρην,\testimonia{1139}{(ἄζυγα) \Phryn 100.3}\apparat{1139}{κόρην R : \lecon{κούρην} Bothe}\\\passeLigne}
\end{groupe}%
%% v. 1136-1159
\begin{groupe}%
\Strophe{(στροφή α)}\typeVers{glyc.}\chantChoeur{ἣ πόλιν ἡμετέραν ἔχει}
\end{groupe}%
%% v. 1136-1159
\begin{groupe}%
\typeVers{glyc.}\chantChoeur{καὶ κράτος φανερὸν μόνη}
\end{groupe}%
%% v. 1136-1159
\begin{groupe}%
\typeVers{pherecr.}\chantChoeur{κλῃδοῦχός τε καλεῖται.\apparat{post v. 1143}{στυγνὰς ὡσσε \Rac (cf. στυγοῦσ᾽ ὥσπερ v. sq. ?) : \lecon{del.} \Rpc : στυγνὰς ὡς ἐ \Mon}\\\passeLigne}
\end{groupe}%
%% v. 1136-1159
\begin{groupe}%
\Strophe{(μεσῳδός α)}\typeVers{$4$ bacch.}\chantChoeur{\ligne{1144}φάνηθ᾽ ὦ τυράννους στυγοῦσ᾽ ὥσπερ εἰκός.\\\passeLigne}
\end{groupe}%
%% v. 1136-1159
\begin{groupe}%
\Strophe{(ἀντιστρ. α)}\typeVers{glyc.}\chantChoeur{δῆμός τοί σε καλεῖ γυναι-}
\end{groupe}%
%% v. 1136-1159
\begin{groupe}%
\typeVers{$+$ glyc.}\chantChoeur{\alinea κῶν· ἔχουσα δέ μοι μόλοις}
\end{groupe}%
%% v. 1136-1159
\begin{groupe}%
\typeVers{pherecr.}\chantChoeur{εἰρήνην φιλέορτον.\\\passeLigne}
\end{groupe}%
%% v. 1136-1159
\begin{groupe}%
%
\Strophe{(στροφή β)}\typeVers{glyc.}\chantChoeur{ἥκετέ τ᾽ εὔφρονες ἵλαοι,\apparat{1148}{ἥκετ᾽ R : \lecon{ἥκετέ τ᾽} Enger : ἥκετε δ᾽ Hermann}}
\end{groupe}%
%% v. 1136-1159
\begin{groupe}%
\typeVers{ibyc.}\chantChoeur{\ligne{1149}πότνιαι, ἄλσος ἐς ὑμέτερον,}
\end{groupe}%
%% v. 1136-1159
\begin{groupe}%
\typeVers{dact. aeol.}\chantChoeur{οὗ δὴ ἀνδράσιν οὐ θεμιτεύετον εἰσορᾶν\apparat{1150-1}{\lecon{οὗ δὴ ἀνδράσιν} R : οὗ δῆτ' ἀνδράσιν Burges : ἄνδρας ἵν(α) Hermann \espApp θεμιτὸν R : θέμις Burges : θεμιτ' Bothe : \lecon{θεμιτεύετον} Bibi \ligneApp{1153} λαμπάσι R : \lecon{λαμπάσι.} Bibi}%
\commentaire{1150-3}{οὗ δὴ ἀνδράσιν οὐ / θεμιτὸν εἰσορᾶν / ὄργια σεμνὰ θεαῖν / ἵνα λαμπάσι φαί-/}{Les mots †οὐ θεμιτὸν† ont été diversement retouchés, jamais de façon satisfaisante. Pour ma part, j'y verrais le verbe θεμιτεύω au duel.
La première partie du chant est une prière adressée à Athéna ; un scribe -- point bien réveillé -- n'a pas tout-de-suite vu que la seconde partie s'adressait aux deux déesses thesmophores.
Au premier vers de cette seconde partie (1148), il aura omis de mettre le verbe au pluriel : ἥκετ<έ τ᾽> εὔφρονες (simple \italique{haplographie}, comme on dit).
\italique{Bis repetita} : non encore averti par la longue énumération des duels aux vers 1154 et suivants, il persévère dans son erreur et le premier duel lui échappe tout aussi haplographiquement :%
\citComm{θεμι<τεύε>τον.}
Le verbe θεμιτεύω, \italique{permettre selon la loi ou le rite}, est utilisé par Euripide avec le même complément, ὄργια, au vers 79 des \italique{Bacchantes} :%
\citComm{τά τε ματρὸς μεγάλας ὄρ-\\
\alinea για Κυβέλας θεμιτεύων\\
Vénérant les orgies de la mère Cybèle la grande}\par 
Cerise sur le gâteau : cette émendation permet de rétablir la correspondance métrique avec le vers 1158 : %
\citComm{οὗ δ\synizese{ὴ ἀ}νδράσιν\marqueur{1} οὐ θεμιτεύετον εἰσορᾶν\\
εἰ καὶ πρότερόν ποτ᾽ ἐπηκόω ἤλθετον\\
$ - \ - \ \cup \ \cup \ - \ \cup \ \cup \ - \ \cup \ \cup \ - \ \cup \ -$\marqueur{2}}%
}%
\debutNote{1150-53}%
\note{1}{Signalons la synalèphe au début du vers : δ\synizese{ὴ ἀ}νδράσιν -- je reprends le vocable de Koster dans son \italique{Traité de métrique grecque}. Il la distingue de la synizèse, terme qu'il applique lorsque la contraction entre deux voyelles a lieu au sein d'un même mot. Subtilité de grammairien ? pas que : la synizèse semble obéir à des lois différentes quant à la longueur des voyelles. Prato, lui, évoque à ce vers la présence d'une \italique{crase}.% \eff{italien ?}.
}%
\note{2}{Séquence interprétée par Lætitia comme un \italique{dactyle éolien}, c'est-à-dire dans son jargon une base (ici unissyllabique $ \times $, permise par l'éolisme du vers) sur laquelle viennent se fixer un nombre arbitraire de dactyles purs ($ - \cup \cup $) et finissant par un \italique{dodrans} ($ - \cup \cup - \cup - $), incontournable en contexte choriambique.}}
\end{groupe}%
%% v. 1136-1159
\begin{groupe}%
\typeVers{tetram. dact.}\chantChoeur{ὄργια σεμνὰ θεαῖν ἵνα λαμπάσι.%\apparat{1152-3}{\lecon{θεαῖν} R : θεοῖν Cobet \espApp λαμπάσι R : \lecon{λαμπάσι.} Bibi}%
\commentaire{1151}{ἵνα,}{Une glose à l'intérieur du texte du manuscrit\marqueur{} précise quel sens donner à (ἵνα : ) ὅπου\marqueur{1}.}%
\debutNote{1151}%
\note{1}{On retrouve une glose identique dans de nombreux lexiques : Συναγωγή ι 82 ; Hésychius ι 648 ; Photius ι 133 ; Suidas ι 360 ; \&c.}\\\passeLigne}
\end{groupe}%
%% v. 1136-1159
\begin{groupe}%
\Strophe{(μεσῳδός β)}\typeVers{trim. dact.}\chantChoeur{\ligne{1154}φαίνετον ἄμβροτον ὄψιν.\\\passeLigne}
\end{groupe}%
%% v. 1136-1159
\begin{groupe}%
\Strophe{(ἀντιστρ. β)}\typeVers{glyc.}\chantChoeur{μόλετον, ἔλθετον, ἀντόμεθ᾽ ὦ}
\end{groupe}%
%% v. 1136-1159
\begin{groupe}%
\typeVers{ibyc.}\chantChoeur{Θεσμοφόρω πολυποτνία,\apparat{1156}{πολὺ πότνια R : \lecon{πολυποτνία} Hermann}}
\end{groupe}%
%% v. 1136-1159
\begin{groupe}%
\typeVers{dact. aeol.}\chantChoeur{εἰ καὶ πρότερόν ποτ᾽ ἐπηκόω ἤλθετον}
\end{groupe}%
%% v. 1136-1159
\begin{groupe}%
\chantChoeur{\ligne{1159}νῦν ἀφίκεσθον †ἱκετεύομεν ἐνθάδ᾽ ἡμῖν†.\apparat{1159}{\lecon{ἐνθάδ᾽ ἡμῖν} R : ἐνθάδε χἠμῖν Hermann}\\\passeLigne}%
\end{groupe}%
%% v. 1160-1163
\begin{groupe}%
%
\personnageGrec{Εὐριπίδης}%
\typeVers{trim. iamb.}\trimIamb{γυναῖκες εἰ βούλεσθε τὸν λοιπόν χρόνον\\
σπονδὰς ποιήσασθαι πρὸς ἐμέ, νυνὶ πάρα,\\
ἐφ᾽ ᾧτ᾽ ἀκοῦσαι μηδὲν ὑπ᾽ ἐμοῦ μηδαμὰ\\
κακὸν τὸ λοιπόν. ταῦτ᾽ ἐπικηρυκεύομαι.\\
\\\passeLigne}%
\end{groupe}%
%% v. 1164
\begin{groupe}%
\personnageGrec{Κορυφαῖος}%
\trimIamb{\ligne{1164}χρείᾳ δὲ ποίᾳ τόνδ᾽ ἐπεσφέρεις λόγον;\\\passeLigne}%
\end{groupe}%
%% v. 1165-1169 (part 1)
\begin{groupe}\personnageGrec{Εὐριπίδης}
\trimIamb{ὅδ᾽ ἐστὶν οὑν τῇ σανίδι κηδεστὴς ἐμός.\apparat{1165}{οὖν R : \lecon{οὑν} Biset}\\
ἢν οὖν κομίσωμαι τοῦτον, οὐδὲν μή ποτε\apparat{1166}{\lecon{μή} R : μου Bentley}\\
κακῶς ἀκούσητ᾽· ἢν δὲ μὴ πίθησθέ μοι,\apparat{1167}{ἀκούσαιτ' R : \lecon{ἀκούσητ᾽} Elmsley (ad \italique{Ach.} v.295) \espApp πείθησθέ R : \lecon{πίθησθέ} Hirschig (ad \italique{Vesp} v.155)}\\
ἃ νῦν ὑποικουρεῖτε τοῖσιν ἀνδράσιν\\\passeLigne}%
\end{groupe}%
\endnumbering
\end{Leftside}
\begin{Rightside}
\beginnumbering
%% v. 1128-1132
\begin{groupe}%
\personnageFr{Euripide}%
\trimIambFr{Aïe, aïe! que faire ? à quelles raisons recourir ?\\
Mais cette nature barbare ne les admettrait pas.\\
« Proposer à des sots des doctrines nouvelles,\\%\commentaire{1130-1}{}{Vers qu'Euripide met dans la bouche de Médée, dans la tragédie de ce nom, v. 298.},\\
c'est perdre sa peine. » Eh bien, dressons quelqu'autre\\
machine, qui soit de sa portée. (Il s'en va.)\\\passeLigne}%
\end{groupe}%
%% v. 1133
\begin{groupe}%
\personnageFr{Archer}%
\trimIambFr{La canaille de renard, comme il pensait m'embabouiner !\\\passeLigne}%
\end{groupe}%
%% v. 1134
\begin{groupe}%
\personnageFr{Mnésiloque}%
\trimIambFr{Souviens-toi, Persée, que tu me laisses en vilaine posture.\\\passeLigne}%
\end{groupe}%
%% v. 1135
\begin{groupe}%
\personnageFr{Archer}%
\trimIambFr{As-tu encore envie d'attraper le fouet? (Il se recouche et s'endort.)\\\passeLigne}%
\end{groupe}%
%% v. 1136-1159
\begin{groupe}%
\personnageFr{Chœur}%
\chantChoeurFr{C'est ma coutume d'appeler parmi nous Pallas,}
\end{groupe}%
%% v. 1136-1159
\begin{groupe}%
\chantChoeurFr{amie du chœur, la jeune fille,}
\end{groupe}%
%% v. 1136-1159
\begin{groupe}%
\chantChoeurFr{\ligne{1138}la vierge libre du joug,\\\passeLigne}
\end{groupe}%
%% v. 1136-1159
\begin{groupe}%
\chantChoeurFr{\ligne{1140}gardienne de notre acropole,}
\end{groupe}%
%% v. 1136-1159
\begin{groupe}%
\chantChoeurFr{seule en possession d'un visible pouvoir,}
\end{groupe}%
%% v. 1136-1159
\begin{groupe}%
\chantChoeurFr{et qui a nom Clédouque.\\\passeLigne}
\end{groupe}%
%% v. 1136-1159
\begin{groupe}%
\chantChoeurFr{Apparais, toi qui détestes les tyrans, comme faire se doit.\\\passeLigne}
\end{groupe}%
%% v. 1136-1159
\begin{groupe}%
\chantChoeurFr{Le peuple des femmes t'appelle.}
\end{groupe}%
%% v. 1136-1159
\begin{groupe}%
\chantChoeurFr{Viens à moi avec}
\end{groupe}%
%% v. 1136-1159
\begin{groupe}%
\chantChoeurFr{la Paix amie des fêtes.\\\passeLigne}
\end{groupe}%
%% v. 1136-1159
\begin{groupe}%
\chantChoeurFr{\ligne{1148}Venez, augustes déesses, bienveillantes et propices,}
\end{groupe}%
%% v. 1136-1159
\begin{groupe}%
\chantChoeurFr{dans votre enceinte sacrée,}
\end{groupe}%
%% v. 1136-1159
\begin{groupe}%
\chantChoeurFr{où vous ne permettez pas aux hommes de contempler}
\end{groupe}%
%% v. 1136-1159
\begin{groupe}%
\chantChoeurFr{vos saintes orgies à la lueur des torches.\\\passeLigne}
\end{groupe}%
%% v. 1136-1159
\begin{groupe}%
\chantChoeurFr{Découvrez vos faces immortelles.\\\passeLigne}
\end{groupe}%
%% v. 1136-1159
\begin{groupe}%
\chantChoeurFr{Venez, approchez, nous vous en conjurons,}
\end{groupe}%
%% v. 1136-1159
\begin{groupe}%
\chantChoeurFr{Thesmophores très vénérées.}
\end{groupe}%
%% v. 1136-1159
\begin{groupe}%
\chantChoeurFr{Si jamais auparavant vous répondîtes à notre appel,}
\end{groupe}%
%% v. 1136-1159
\begin{groupe}%
\chantChoeurFr{\ligne{1159}venez aujourd'hui, sur nos instantes prières, ici au milieu de nous.\\\passeLigne}
\end{groupe}%
%% v. 1160-1163
\begin{groupe}%
\personnageFr{Euripide}%
\trimIambFr{(Déguisé en vieille procureuse et portant une pectis ou harpe lydienne. Il est accompagné d'une danseuse et d'un jeune flûtiste). Femmes, si vous voulez pour le reste du temps
faire la paix avec moi, c'est le moment,\\
à charge que jamais plus à l'avenir je ne dirai de vous\\
le moindre mal. Telles sont mes propositions.\\\passeLigne}%
\end{groupe}%
%% v. 1164
\begin{groupe}%
\personnageFr{Coryphée}%
\trimIambFr{\ligne{1164}Qu'est-ce qui te décide à faire ces ouvertures ?\\\passeLigne}%
\end{groupe}%
%% v. 1165-1169 (part 1)
\begin{groupe}\personnageFr{Euripide}
\trimIambFr{Cet homme attaché à l'ais est mon beau-père.\\
Que je le recouvre, et jamais\\
je ne médirai de vous. Si vous refusez,\\
toutes vos pratiques à huis-clos, je les révélerai\\\passeLigne}%
\end{groupe}%
\endnumbering
\end{Rightside}
\end{pages}
\Pages
\clearpage\ifodd\value{page}\hbox{}\newpage\fi
\enTeteG{ΕΞΟΔΟΣ --- Codex Ravennas 429, f. 176v (Classense 178v) --- vv. 1169–1210}
\begin{adjustwidth}{-1.5cm}{-3.0cm}
\begin{adjustbox}{max totalheight=21.0cm, max width=16.5cm, keepaspectratio}
\begin{minipage}{16.5cm}
{\fontsize{7.2pt}{8.5pt}\selectfont
\noindent\textbf{Folio 29 --- Codex Ravennas 429, \texttt{f. 176v} (\texttt{Classense 178v}) : ΕΞΟΔΟΣ, vv. 1169–1210}\hfill \textit{Témoins :} $\text{R}^1, \Rpc, \Rdeux$, \PsiTrois, \PsiQuatre\par\vspace{2pt}
\hrule height 0.4pt\vspace{2.5pt}
}
{\fontsize{8.1pt}{9.5pt}\selectfont
\ligneRav{1169}{\auDessus{εὐ}{ρ}}{ἀπὸ τῆς στρατιᾶς παροῦσιν ὑμῶν διαβαλῶ·}{}
\ligneRav{1170}{χορ/}{τὰ μὲν παρ᾽ ἡμῶν ἴσθι σοι πεπεισμένα·}{}
\ligneRav{1171}{}{τὸν βάρβαρον δὲ τοῦτον αὐτὸς \difference{πεῖθε} σύ·}{\textbf{1171} πεῖσαι R : \lecon{πεῖθε} Biset}
\ligneRav{1172}{\auDessus{εὐ}{ρ}}{οὐδὲν μὰ Δί᾽ ἐμὸν \inRas{γ᾽}{} ὤν, ὥσπερ ἐχθὲς \difference{ὤμνυες}·}{\textbf{1172} Εὐριπίδης ἐν σχήματι προαγωγοῦ γραός Σ \espApp \lecon{ἐμὸν} \inRas{γ᾽}{} \Rdeux}
\ligneRav{1173}{}{τί δ᾽; οὐκ \expo{ἔφραζον} (ἔφ\expo{ραζ}ον \Rpc) ταῦτ᾽ ἐγώ σοι;}{\textbf{1173} \lecon{ἔφραζον} (ex ἔφρον) \Rpc}
\ligneRav{1174}{}{πρῶτον μὲν οὖν δίελθε \difference{κἀνακάλπασον}·}{\textbf{1174} κἀνακόλπασον R : \lecon{κἀνακάλπασον} Hermann, cf. Σ, Suid. κ 623, Zonar. p.1171}
\ligneRav{1175}{}{σὺ δ᾽ ὦ \difference{Τερηδὸν} ἐπαναφύσα Περσικόν·}{\textbf{1175} Τερηδὼν R, Suid. τ 340 AGFM\expo{ac} : \lecon{Τερηδὸν} Suid. VM\expo{pc}}
\ligneRav{1176}{τοξ/}{τί τὸ βόμβο τοῦτο; κῶμό τις ἀνεγεῖρί μοι;}{\textbf{1176} κωμοτίς ἀν- R : κωμοδρίαν Σ : \lecon{κῶμό τις ἀν-} van Velsen}
\ligneRav{}{\auDessus{εὐ}{ρ}}{ἡ παῖς ἔμελλε προμελετᾶν ὦ τοξότα·}{}
\ligneRav{1178}{}{ὀρχησομένη γὰρ \difference{ἔρχεθ᾽} ὡς ἄνδρας τινάς·}{\textbf{1178} ἔρχετ᾽ R : \lecon{ἔρχεθ᾽} Brunck}
\ligneRav{1179}{τοξ/}{ὀρκῆσι καὶ μελετῆσι, οὐ \difference{κωλῦσ᾽} ἐγώ·}{\textbf{1179} κωλύσ᾽ R : \lecon{κωλῦσ᾽} Blaydes (κωλῦσι $=$ κωλύσω), cf. ad 1003, 1196, 1201, 1225}
\ligneRav{1180}{}{\difference{ὠς} ἐλαπρός, \difference{ὤσπερ} ψύλλο κατὰ τὸ κῴδιο·}{\textbf{1180} ὡς et ὥσπερ R : \lecon{ὠς} et \lecon{ὤσπερ} Brunck}
\ligneRav{}{\auDessus{εὐ}{ρ}}{φέρε θοἰμάτιον ἄνωθεν ὦ τέκνον τοδί·}{}
\ligneRav{1182}{}{καθιζομένη δ᾽ ἐπὶ \difference{τοῖσι} γόνασι τοῦ Σκύθου}{\textbf{1182} τοῖς R : \lecon{τοῖσι} Bentley}
\ligneRav{}{—}{τὼ πόδε πρότεινον, ἵν᾽ ὑπολύσω. : ναῖκι ναὶ}{}
\ligneRav{1184}{}{κάτησο κάτησο, ναίκι ναί, τυγάτριον·}{\textbf{1184} \lecon{ναίκι ναὶ} R : ναίκι ναίκι Bothe}
\ligneRav{1185}{}{ἀνίστασ᾽ ὀρθός· \difference{ὡς στεριπὸ} (ὡ σ\inRas{π}{τ}εριπὸ) τὸ τιττί᾽, ὥσπερ γογγύλη·}{\textbf{1185} ὡ στεριπὸ (ex σπεριπο) \Rpc : \lecon{ὡς στεριπὸ} (sic) \Mon \espApp ὥσπερ R : \lecon{ὤσπερ} Brunck \espApp \lecon{γογγυλί} \PsiTrois, Enger : γογγύλῃ R}
\ligneRav{}{\auDessus{εὐ}{ρ}}{αὔλει σὺ θᾶττον· ἔτι δέδοικας τὸν Σκύθην;}{}
\ligneRav{1187}{τοξ/}{καλόν γε τὸ πυγή· \difference{κλαυσῖ γ᾽} (κλαυσ\inRas{᾽ ε}{}ϊ γ᾽) ἂν μὴ ᾽νδον μένῃ·}{\textbf{1187} κλαυσ\inRas{᾽ ε}{}ϊ γ᾽ \Rdeux : κλαυσεῖ γ᾽ \Mon \espApp ἂν R : \lecon{ἢν} Dindorf (1839) \ligneApp{post 1187} μ$]$\douteux{ε}νο\douteux{ }$[$ \PsiTrois : ἀνακύπτη· καὶ παρακύπτι· ἀπεψωλημένος· R : \lecon{del.} Ellebodius}
\ligneRav{}{}{\hspace{2.2em}\textbf{[παρεπιγραφή R\expo{1}, \PsiTrois{} :]} \lecon{ἀνακύπτη· καὶ παρακύπτι· ἀπεψωλημένος·}}{\textbf{post 1187} \lecon{ἀνακύπτη· καὶ παρακύπτι· ἀπεψωλημένος·} R, μ]ενο[ \PsiTrois}
\ligneRav{1188}{}{εἶἑν· καλὴ τὸ σχῆμα περὶ τὸ \inRas{σ}{}ποστίον·}{\textbf{1188} τὸ \inRas{σ}{}ποστίον \Rpc : τὸ ποστίον \Mon}
\ligneRav{1189}{\auDessus{εὐ}{ρ}}{καλῶς ἔχει. λαβὲ θοἰμάτιον· ὥρα \difference{'στὶ} νῷν}{\textbf{1189} 'στιν R : \lecon{'στὶ} \Mon}
\ligneRav{1190}{—}{ἤδη βαδίζειν. : \difference{οὐκὶ} \difference{πιλῆσι} πρῶτά με;}{\textbf{1190} τί οὐκ ἐπιλήσει R : \lecon{οὐκὶ} (Biset) \lecon{πιλῆσι} Porson (ad \italique{Aristophanica} p.222), cf. Σ `οὐχὶ φιλήσειν'}
\ligneRav{1191}{\auDessus{εὐ}{ρ}}{πάνυ γε· φίλησον αὐτόν. : \difference{ὂ ὂ ὂ} παπαπαπαῖ,}{\textbf{1191} ὂ R : \lecon{ὂ ὂ ὂ} \Mon}
\ligneRav{1192}{}{\difference{ὠς} γλυκερὸ τὸ γλῶσσ᾽, \difference{ὤσπερ} Ἀττικὸς \difference{μέλι}·}{\textbf{1192} ὡς et ὥσπερ R : \lecon{ὠς} et \lecon{ὤσπερ} Brunck \espApp \lecon{μέλι} \PsiTrois, \Mon : μέλις R}
\ligneRav{}{—}{τί οὐ κατεύδει παρ᾽ ἐμέ; : χαῖρε τοξότα,}{}
\ligneRav{1194}{—}{οὐ γὰρ γένοιτ᾽ ἂν τοῦτο. : \difference{ναὶ ναὶ} \difference{γρᾴδιο}·}{\textbf{1194} ναὶ R : ναίκι Bentley : \lecon{ναὶ ναὶ} Brunck \espApp γράδιον R : \lecon{γρᾴδιο} Bothe}
\ligneRav{1195}{—}{ἐμοὶ \inRas{σὺ}{} \difference{κάρισο σὺ} τοῦτο : δώσεις οὖν δραχμή;}{\textbf{1195} \inRas{σὺ}{} κάρισοσ οὐ \Rpc : κάρισο σοῦ \Mon : \lecon{κάρισο σὺ} Brunck}
\ligneRav{}{τοξ/}{ναὶ ναῖκι δῶσι. : τἀργύριον τοίνυν φέρε·}{}
\ligneRav{1197}{τοξ/}{ἀλλ᾽ οὐκ \difference{ἔκὠδέν}· ἀλλὰ τὸ συμβήνην λαβέ·}{\textbf{1197} έκ' ὠδέν R : ἐχω οὐδέν Suid. σ 1273, Photius σ 660 : \lecon{ἔκὠδέν} Bergk \espApp \lecon{συμβήνην} R : συβήνην Σ, Suid., Photius}
\ligneRav{1198}{}{\difference{κομίζι σ᾽ αὐτίκ᾽} (κ\inRas{ω}{ο}μίζεις αὐτοῖς) : ὦ τάλαιν᾽ ἐγώ·}{\textbf{1198} κ\inRas{ω}{ο}μίζεις αὐτοῖς \Rpc : κομιεῖς αὐτό (αὖθις Brunck) Biset : (Τοξότ.) κομίσι σ' αὐτός (αὖτις anonym. Parisinus) Enger : (Τοξότ.) \lecon{κομίζι σ᾽ αὐτίκ᾽} Bibi \espApp ἀκουλούτι R : \lecon{ἀκολούτι} \Mon : ἀκουλούτει Brunck}
\ligneRav{}{}{σὺ δὲ τοῦτο τήρει τὴ γέροντο, γρᾴδιο·}{}
\ligneRav{1200}{—}{ὄνομα δέ σοι τί ἔστιν; : Ἀρτεμισία·}{}
\ligneRav{1201}{τοξ/}{μεμνῆσι τοίνυν τοὔνομ᾽· \difference{Ἀρταμουξία}·}{\textbf{1201} ἀρτομουξία R : \lecon{Ἀρταμουξία} ed. Lipsiensis (1814)}
\ligneRav{}{\auDessus{εὐ}{ρ}}{Ἑρμῆ δόλιε ταυτὶ μὲν ἔτι καλῶς ποιεῖς·}{}
\ligneRav{1203}{}{σὺ μὲν οὖν ἀπότρεχε παιδάριον τουτὶ λαβών·}{\textbf{1203} \lecon{τουτὶ} R : ταυτὶ Dobree, cf. Σ}
\ligneRav{}{}{ἐγὼ δὲ λύσω τόνδε. σὺ δ᾽ ὅπως ἀνδρικῶς}{}
\ligneRav{1205}{}{ὅταν λυθῇς τάχιστα φεύξει καὶ τενεῖς}{}
\ligneRav{}{}{ὡς τὴν γυναῖκα καὶ τὰ παιδί᾽ οἴκαδε·}{}
\ligneRav{}{μ\auDessus{ν}{Η}}{ἐμοὶ μελήσει ταῦτά γ᾽ ἢν ἅπαξ λυθῶ·}{}
\ligneRav{1208}{\auDessus{εὐ}{ρ}}{ὠς τάκιστα σὺ λέλυσο\inRas{ν}{} (cf. \PsiQuatre) καὶ φεῦγ᾽·}{\textbf{1208} ca. 6 litt. $]$· σ\douteux{ο}$[$νεργον \PsiQuatre : \lecon{λέλυσο}\inRas{ν}{} \Rpc : λέλυσαι Bentley}
\ligneRav{}{—}{ἥκοντα καταλαβεῖν. : ἐγὼ δὴ τοῦτο δρῶ·}{}
\ligneRav{1210}{τοξ/}{σὺ δὲ τυγά\inRas{θ}{τ}ριον (\Rpc) κόσμιε καὶ φύλαττε·}{\textbf{1210} τυγάτριον ex τυγάθριον ut videtur \Rpc}
}
\vspace{3pt}\noindent\rule{4.5cm}{0.4pt}\par\vspace{1.5pt}
{\fontsize{7.0pt}{8.3pt}\selectfont
\noindent\textbf{Loci paralleli (vv. 1169–1210) ---} \textbf{1169} (στρατιᾶς) Photius σ 608 στρατιάν \espTest \textbf{1174} (κἀνακάλπασον) Photius α 1514 ἀνακαλπάζει ; cf. Zonaras 1171 καλπάζειν· τὸ ἄκρως βαδίζειν $\simeq$ Scholium $=$ Suidas κ 623 καλπάζειν \espTest \textbf{1175} Suidas τ \espTest (Περσικόν) Pollux IV 100 ; cf. Hesychius π 2005 (?) \espTest \textbf{1183} (τὼ πόδε) Hesychius τ 1758 τὼ πόδε· τοὺς πόδας δυϊκῶς \espTest \textbf{1187} (πυγή) Hesychius π 4279 πυγή· τὸ κάθισμα \espTest \textbf{1197} Photius σ 660 συβήνη $=$ Suidas σ 1273 (sed ἐχω οὐδέν pro έκ' ὠδέν, τὸ omisso et συβήνην pro συμβήνην) \espTest (συμβήνην) Hesychius σ 2137 συμβίνης· καπροβόλον. ἐμβόλιον\par
}
\vspace{2.5pt}\noindent\rule{4.5cm}{0.4pt}\par\vspace{1.5pt}
{\fontsize{6.9pt}{8.2pt}\selectfont
\noindent\textbf{Suite (en commentaire) de l'apparat critique \& notes paléographiques (\texttt{f. 176v}) ---} \textbf{v. 1172 :} $\Sigma^{\text{R}}$ Εὐριπίδης ἐν σχήματι προαγωγοῦ γραός \espApp \textbf{ἐμὸν \inRas{γ'}{}} \Rdeux{}. --- \textbf{v. 1174 :} κἀνακόλπασον R (confusion minuscule $\alpha \to \text{ο}$) : \lecon{κἀνακάλπασον} Hermann, cf. $\Sigma$, Souda $\kappa\ 623$, Zonaras p. 1171 (dans $\Sigma^{\text{R}}$, ἁβρῶς R : \lecon{ἄκρως} \lecon{Bibi} ex Zonaras). --- \textbf{v. 1176 :} κωμοτίς ἀν- R (κωμοδρίαν $\Sigma$) : \lecon{κῶμό τις ἀν-} von Velsen. --- \textbf{v. 1190 :} τί οὐκ ἐπιλήσει R : \lecon{οὐκὶ} (Biset) \lecon{πιλῆσι} Porson (\italique{Aristophanica}, p. 222, cf. $\Sigma$ οὐχὶ φιλήσειν). --- \textbf{v. 1197 :} έκ' ὠδέν R (ἐχω οὐδέν Souda $\sigma\ 1273$, Photius $\sigma\ 660$) : \lecon{ἔκὠδέν} Bergk \espApp \lecon{συμβήνην} R, défendu par \lecon{Bibi} (ap. Hésychius $\sigma\ 2137$) : συβήνην $\Sigma$, Souda, Photius. --- \textbf{v. 1198 :} γϋ. ante v. 1198 (Boninus) et σκ'. ante ἀκολ. in R \espApp κ\inRas{ω}{ο}μίζεις αὐτοῖς \Rpc{} (κομιεῖς αὐτό Biset ; αὖθις Brunck ; (Τοξότ.) κομίσι σ' αὐτός Enger) : \lecon{(Τοξότ.) κομίζι σ' αὐτίκ'} \lecon{Bibi} \espApp ἀκουλούτι R : \lecon{ἀκολούτι} \Mon{} (ἀκουλούτει Brunck). --- \textbf{v. 1203 :} \lecon{τουτὶ} R (ταυτὶ Dobree, cf. $\Sigma$).\par
}
\end{minipage}
\end{adjustbox}
\end{adjustwidth}
\newpage
\enTeteD{Codex Ravennas 429 (Biblioteca Classense) --- Fac-similé du f. 176v (vv. 1169–1210)}
\begin{adjustwidth}{-1.5cm}{-3.0cm}
\begin{center}
\includegraphics[width=16.5cm,height=20.5cm,keepaspectratio]{/workspaces/Thesmophories/dossier_ravennas_thesmophories/folios_individuels_ff162v-177r/29_f176v_classense178v_vv1169-1210.jpg}\par
\vspace{2pt}
{\scriptsize\textbf{Codex Ravennas 429, f. 176v (Classense 178v)} --- ΕΞΟΔΟΣ, \textbf{vv. 1169–1210} ($\text{R}^1, \Rpc, \Rdeux$, \PsiTrois, \PsiQuatre).}
\end{center}
\end{adjustwidth}
\newpage
\clearpage\ifodd\value{page}\hbox{}\newpage\fi
\enTeteG{ΘΕΣΜΟΦΟΡΙΑΖΟΥΣΑΙ --- Texte grec établi (f. 176v, vv. 1169–1210)}
\enTeteD{LES THESMOPHORIES --- Traduction française (f. 176v, vv. 1169–1210)}
\begin{pages}
\begin{Leftside}
\beginnumbering
%% v. 1169+ (part 2)
\begin{groupe}
\trimIamb{\ligne{1169}ἀπὸ τῆς στρατιᾶς παροῦσιν ὑμῶν διαβαλῶ.\testimonia{1169}{(στρατιᾶς) Photius σ 608 στρατιάν}\commentaire{1168-9}{ἃ νῦν ὑποικουρεῖτε τοῖσιν ἀνδράσιν / ἀπὸ τῆς στρατιᾶς παροῦσιν.}{À propos de la légèreté des Athéniennes pendant le service de leur mari, cf. le vers 495 et le commentaire \italique{ad locum}.}\\\passeLigne}%
\end{groupe}%
%% v. 1170-1171
\begin{groupe}%
\personnageGrec{Κορυφαῖος}%
\trimIamb{τὰ μὲν παρ᾽ ἡμῶν ἴσθι σοι πεπεισμένα·\commentaire{1170}{τὰ μὲν παρ᾽ ἡμῶν ἴσθι σοι πεπεισμένα·}{Une scholie intérieure au texte nous précise :%
\citComm{Ἀττικῶς ἀντὶ τοῦ `πεπείσμεθα'.}}\\
τὸν βάρβαρον δὲ τοῦτον αὐτὸς πεῖθε σύ.\apparat{1171}{πεῖσαι R : \lecon{πεῖθε} Biset}\\\passeLigne}%
\end{groupe}%
%% v. 1172-1175
\begin{groupe}%
%
\personnageGrec{Εὐριπίδης}%
\trimIamb{ἐμὸν ἔργον ἐστίν· καὶ σὸν ὦλάφιον ἅ σοι\apparat{1172}{Εὐριπίδης ἐν σχήματι προαγωγοῦ γραός Σ \espApp \lecon{ἐμὸν} \inRas{γ᾽}{} \Rdeux \ligneApp{1173} \lecon{ἔφραζον} (ex ἔφρον) \Rpc \ligneApp{1174} κἀνακόλπασον R : \lecon{κἀνακάλπασον} Hermann, cf. Σ, Suid. κ 623, Zonar. p.1171 %: κᾆτα δ' ὄκλασον Moreau
\ligneApp{1175} Τερηδὼν R, Suid. τ 340 AGFM\expo{ac} : \lecon{Τερηδὸν} Suid. VM\expo{pc}}\\
καθ᾽ ὁδὸν ἔφραζον ταῦτα μεμνῆσθαι ποιεῖν.\\
πρῶτον μὲν οὖν δίελθε κἀνακάλπασον.\testimonia{1174}{(κἀνακάλπασον) Photius α 1514 ἀνακαλπάζει ; cf. Zonaras 1171 καλπάζειν· τὸ ἄκρως βαδίζειν $\simeq$ Scholium $=$ Suidas κ 623 καλπάζειν}%
\commentaire{1174}{πρῶτον μὲν οὖν δίελθε κἀνακόλπασον.}{(Que le lecteur curieux aille faire son miel de la note qu'a consacré Jacques Moreau à ce vers\marqueur{1}.)
Le \italique{Ravennas} porte au vers 1174 la leçon : πρῶτον μὲν οὖν δίελθε κἀνακ\textbf{ό}λπασον.
Le verbe ἀνακολπάζειν est un hapax, certes, mais dont le sens se laisserait entendre -- s'il n'y avait lieu comme nous l'allons voir de le soupçonner : ὁ κόλπος désigne au propre un \italique{pli}, \italique{creux} ; aussi bien le pli d'un vêtement -- notamment au niveau de la ceinture -- que, chez la mère ou la nourrice, le \italique{genou} ou le \italique{sein}. On trouve chez Homère l'expression κόλπον ἀνιεμένη\marqueur{2} lorsqu'Hécube fait tomber le haut de son vêtement et découvre son sein.
Il n'y aurait rien d'étonnant à ce qu'ἀνακολπάζειν ait pu signifier \italique{retrousser un vêtement}\marqueur{3}.
Pour définir les ζειραί, ces vêtements qu'on \italique{resserre au niveau de la ceinture}, le rédacteur anonyme de l'\latin{Etymologicum Magnum} use du verbe ἀνακολπόω :%
\citComm{ζειραί: χιτῶνες ἀνακεκολπωμένοι\marqueur{4}, ἢ ἀνάκωλοι.}%
Si le \italique{Ravennas} ne comportait en sus du texte au vers 1174 une scholie marginale venant renverser tout notre fragile édifice, ἀνακολπάζειν ferait partie de ces nombreux hapax dont le sens n'a été inféré qu'à grand renfort d'imagination sur de maigres indices étymologiques.\\
Cependant un ancien commentateur eut l'heureuse idée de gloser la leçon que son manuscrit donnait :%
\citComm{σ\douteux{η}[ * * * ] τοῦτο. ση[μαίνει] δὲ τὸ `ἁβρῶς βαδίζειν'.}%
Exit donc la leçon ἀνακολπάζειν, sans rapport avec le \italique{déplacement langoureux}\marqueur{5} dont parle le scholiaste, nullement étayée et d'ailleurs inepte : la danseuse court-vêtue n'a aucun vêtement pesant à retrousser ; au moment d'enflammer le Scythe, Euripide lui fera sauter les préliminaires en lui ordonnant à brûle-pourpoint de retirer tout bonnement son léger ἱμάτιον.\\
Hermann\marqueur{6} proposa de lire le verbe ἀνακαλπάζειν\marqueur{7}. 
Son hypothèse a été reprise depuis lors par toutes les éditions ; elle s'appuyait effectivement sur un indice de taille : à l'adresse κ 623 de la \italique{Souda}, l'on retrouve notre glose précédée du lemme καλπάζειν :%
\citComm{καλπάζειν· τὸ ἁβρῶς βαδίζειν.}%
%Ne chicanons pas l'émendation : transposition d'un `α' et d'un `ο', lettres parmi les plus dissemblables -- graphiquement comme phonétiquement -- et dont la confusion reste rare. Regardons plutôt ce que le sens y gagnerait :
Le verbe est formé à partir de ἡ κάλπη, \italique{le trot} (ou le \italique{galop} ?), qui désigne une sorte de course aux jeux olympiques, où le cavalier, en approchant du but, devait prendre le trot, sauter à terre et suivre son cheval.
%On retrouve le simple καλπάζειν chez Hésychius (κ 519) : καλπάζει· ὀξυποδεῖ ; ὀξυ- est à prendre au sens de ταχύ-, \italique{rapide} (cf. le passage de l'\italique{Oreste} d'Euripide, vers 1550, où Ménélas furieux accourt ὀξύπους $=$ ταχύπους). 
Quant au composé, ἀνακαλπάζειν, Photius (α 1514) nous en a conservé une trace en recopiant la glose d'un grammairien versé dans le \italique{Bon usage} attique :%
\citComm{ἀνακαλπάζει· τινὲς μὲν ὡς οὐ δόκιμον ἐφυλάξαντο τὴν φωνήν, Αἰσχύλος δὲ ἐχρήσατο Μυσοῖς \{ὡς δόκιμον\}· λέγει γάρ `εἶδον καλπάζοντας ἐν αἰχμαῖς'. ὁμοίως Σοφοκλῆς καὶ Ἀριστοφάνης καὶ Πλάτων καὶ ἕτεροι.\\
ἀνακαλπάζει· certains défendirent [l'usage de] de ce mot, comme non attesté. Eschyle pourtant l'utilisa dans ses \italique{Mysiens}. Il dit en effet : «~Je les vis trotter au milieu des lances~». De même Sophocle, Aristophane, Platon et d'autres.}%
C'est bien le sens de \italique{trotter}, \italique{galoper}, c'est-à-dire \italique{courir} qu'il faut donner au verbe <ἀνα>καλπάζειν, témoin le fragment d'Eschyle cité par Photius\marqueur{8} : il y fait allusion à une locution proverbiale\marqueur{9} :%
\citComm{`δι᾽ ὀξείας δραμεῖν'· ἐπὶ τῶν διακινδυνευόντων. ὀξεῖαν γὰρ λέγουσι τὴν λόγχην\\
«~Courir en traversant les pointes~» : à propos de ceux qui prennent des risques ; on appelle pointe la lance.}%
Dans le passage d'Eschyle, καλπάζοντας est un permutant métaphorique -- donc synonyme -- de δραμόντας.\\
Willems \& nombre d'exégètes voulurent rattacher <ἀνα>καλπάζειν à un terme technique de danse : \italique{se déplacer} ou \italique{danser sur les pointes}. Le sens leur a-t-il été suggéré par l'image du cheval au trot ? Dans le célèbre roman de Longus (II, 37), lorsque Daphnis imite le dieu Pan avec ses sabots, il se met à courir %
\citComm{ἐπ' ἄκρων τῶν ὀνύχων.}%
Là est le point : car l'on sait le type d'entraînement auquel se soumet la danseuse : au vers 1180, le Scythe qui la voit si \italique{légère} se mettre à \italique{sauter} («~comme une puce~», dixit le Scythe...), il s'exclame : ὡς ἐλαπρός, ὥσπερ ψύλλο κατὰ τὸ κῴδιο.
Moreau avait donc raison de souligner que la glose présentant notre danseuse aguichant l'Archer \italique{mollement} (ἁβρῶς βαδίζειν) ne tenait pas ; mais l'erreur venait de la glose et non de la correction d'Hermann : on eut le tort de vouloir conserver à tout prix cette glose telle que rapportée par le \italique{Ravennas} et Suidas ; elle n'était pourtant en conformité ni avec l'action ni avec le type d'entraînement -- hyperactif -- de la danseuse.
Le lexique de Jean Zonaras\marqueur{10} nous en livre la bonne version (simple substitution de ἄ\textbf{κ}ρως à ἁ\textbf{β}ρῶς) :%
\citComm{καλπάζειν· τὸ ἄκρως βαδίζειν.}%
Willems avait donc vu juste : l'exercice de la danseuse devait évoquer la course de la κάλπη : commencer par un \italique{trot}, c'est-à-dire par un déplacement en \italique{pointes} (τὸ ἄκρως βαδίζειν) qui s'accélérait pour terminer par un saut ; ce qu'Euripide désigne par le verbe ἀνακαλπάζειν.\\
Cette interprétation nous est confirmée par Eustathe (\italique{Commentaire. à l'}Iliade, XX, v. 228 p. 1206, à propos du mot `σκιρτῷεν') :%
\citComm{%τὸ δὲ σκιρτᾶν ἐοικὸς μὲν πρὸς τὸ σκαίρειν, καίριον δὲ ὂν εἰς ἀλλομένους ἵππους δὶς ἔφη, ἐπιμείνας ὡς ἐχρῆν. 
ἔστι δὲ σκιρτᾶν οὐ τὸ ἁπλῶς θέειν, ἀλλὰ τὸ ἀνέτως καὶ δίχα ἐπιταγῆς καὶ πηδητικῶς καὶ οἷον ὀρχηστικῶς \textbf{κατὰ τοὺς καλπάζοντας}.\\
%«~Bondir~» (σκιρτᾶν) est semblable à «~sauter~» (σκαίρειν), utilisé à bon escient par deux fois pour des chevaux bondissants, ἐπιμείνας ὡς ἐχρῆν (?) 
«~Bondir~» ((σκιρτᾶν)) n'est pas le simple fait de courir, mais de le faire librement, sans ordre, fougueusement, ainsi que le font les \textbf{καλπάζοντες} en dansant.}%
Terminons par Hésychius (κ 519)\marqueur{11}, qui ne le comprenait pas autrement :%
\citComm{καλπάζει· ὀξυποδεῖ.}%
ὀξυποδέω signifie \italique{accélérer le pas}\marqueur{12}. Notons aussi l'autre glose d'Hésychius (σ 822) :%
\citComm{σκαλ\{α\}πάζειν· ῥεμβωδῶς βαδίζειν.\\
σκαλπάζειν : se déplacer en tournoyant.}%
%Ce terme technique n'aurait de toute façon pas de rapport avec la définition, ἁβρῶς βαδίζειν\marqueur{6}, qu'en donne le scholiaste.\\
%Moreau propose de voir derrière `κἀνα' une déformation de κᾆτα ; dont la présence serait tout-à-fait naturelle après πρῶτον μέν dans le premier membre de la phrase\marqueur{7}.\\
%Si l'on admet cette émendation, la suite de la restitution paraît couler de source : remplacez le verbe κολπάζειν par μολπάζειν\marqueur{8}. ἡ μολπή est chez Homère un \italique{chant mêlé de danse}, le verbe se rencontre dans les \italique{Grenouilles} (vers 379 ; partie lyrique) ainsi que chez Athénée (XIII, 598f ; un poète qui chante son amour). Le verbe s'inscrirait parfaitement dans l'action et expliquerait l'ordre d'Euripide à la flûtiste de commencer l'accompagnement musical. Ainsi que le réveil en sursaut au vers suivant de l'archer qui, Scythe qu'il est, confond la douce mélopée avec un κῶμος dionysiaque. Toutefois, nous avons assez évoqué la scholie à ce vers pour ne pas reconnaître que μολπάζειν serait étrangement glosé de la sorte. L'explication réside peut-être dans la lacune de cette scholie (le début en est corrompu). Pour confirmer mon hypothèse, il faudrait qu'elle fût peu ou prou rédigée ainsi :%
%\citComm{ση[μαίνει] `ᾄδειν' τοῦτο. ση[μαίνει] δὲ τὸ `ἁβρῶς βαδίζειν'.}%
% Faute de grives on mange des merles : hésychius η δεϊ χελωνης χρέα φαγεϊν, ή μη φαγεϊν.
}%
\debutNote{1174}%
\note{1}{«~Note sur le vers 1174 des \italique{Thesmophories} d'Aristophane~», \italique{L'antiquité classique}, tome 7, fasc. 2 (1938) pp. 227-231. À peine âgé de vingt ans, le jeune étudiant d'alors fait montre d'une érudition et d'un esprit critique peu communs ; il faut dire qu'il avait pour professeur à l'Université de Bruxelles Henri Grégoire -- le successeur d'Alphonse Willems -- et dont l'assistant n'était autre que Roger Goossens ! bouillon intellectuel dont on imagine sans peine l'extraordinaire émulation qui devait en résulter parmi les élèves.}%
\note{2}{\italique{Iliade}, XX, 80. Et l'explication d'Hésychius (κ 3393) : κόλπον ἀνιεμένη· τὸ κατὰ τὸ στῆθος μέρος ἀνέλκουσα, ἵνα τὸν μαστὸν δείξῃ.}%
\note{3}{Ou, à la rigueur, «~tenir haut les seins~» comme le traduit Germaine Prudhommeau.}%
\note{4}{†ἀνακεκολλαμένοι† dans divers lexiques, notamment le lexique platonicien de Timée le Sophiste ou celui de Photius. La correction est-elle du fait du rédacteur de l'\latin{Etymologicum Magnum} ?}%
\note{5}{Ou \italique{lascif}, dira-t-on, au sens que le \nieme{XXI} siècle prête à ce mot ; c'est-à-dire avec \italique{langueur}, un \italique{manque de vigueur} qui trahit un fort penchant pour la sensualité. Est-ce l'influence de \italique{las} qui aura changé radicalement le sens originel de \italique{lascif} (\italique{lubrique}, \italique{folâtre} ; sens que rapportent encore exclusivement nos dictionnaires) ?}%
\note{6}{Dans sa recension des \italique{Thesmophories} éditées par Enger en 1844 (\italique{Aristophanis Comoediae cum scholiis}), in \italique{Zeitschrift für altertumswissenschaft} (1845), coll. 905-920. Article réédité par Fritzsche dans les \italique{Godofredi Hermanni Opuscula}, vol. 8 ; Hermann y discute de notre vers à la page 307.}%
%\note{3}{\eff{Retrouver dans le bouquin de danse attique la réf ; avec aussi la danse des épées.}}%
\note{7}{Au vers 1201, un `α' est lu fautivement de la même manière en `ο'.}%
\note{8}{$=$ fr. 144a Radt.}%
\note{9}{Transmis par les parémiographes Zénobius (III, 13), Suidas (δ 1186 ; ο 414) et Érasme (2014). Il n'est pas impossible que le proverbe soit lui-même tiré de la danse, et d'une éventuelle \italique{danse des épées}. Sur l'utilisation d'épées pour certaines acrobaties périlleuses, voir \italique{La danse grecque antique} de Maurice Emmanuel (1896), fig. 556 et 557 ou l'article «~Histoire anthropologique des danses acrobatiques~» de	Myriam Peignist, \italique{Corps} (2009/2), n° 7, pp. 29-38. La danseuse se livre-t-elle à ce genre d'acrobaties guerrières pour attirer l'attention du Scythe ? cela paraît peu vraisemblable ; il est préférable d'interpréter le fragment d'Eschyle de façon imagée et ne pas rattacher directement le verbe καλπάζειν à une quelconque danse des épées.
Un proverbe analogue et que l'on connaît en français, \italique{être sur le fil du rasoir}, apparaît déjà chez Homère (\italique{Iliade} X, 173).}%
%\note{9}{Et à sa suite \eff{citer les bouquins sur la danse grecque attique}}%
\note{10}{Édité par Tittmann sous le titre de \italique{Iohannis Zonarae Lexicon} (1808), p. 1171. Ce lexique de la fin du \nieme{XII} ou début du \nieme{XIII} siècle, s'il faut en croire Klaus Alpers, serait plutôt à attribuer à un certain Nicéphore. Le Pseudo-Zonaras a puisé abondamment dans la \italique{Souda} (ou des sources très voisines), notamment parmi les adresses de la lettre Α à Κ et les gloses touchant à Aristophane. Notre cas est intéressant, puisqu'il est peu vraisemblable que la transposition de β et κ soit une erreur ou correction du copiste ; il faut donc en conclure que si le Pseudo-Zonaras est allé piocher sa glose dans la \italique{Souda}, il disposait d'un meilleur manuscrit que ceux qui nous sont parvenus. Klaus Alpers a consacré un article très riche au lexique du Pseudo-Zonaras dans la monumentale \italique{Realencyclopädie der classischen Altertumswissenschaft}, 10/1 (1972), à l'entrée \allemand{'Zonarae' Lexikon}, coll. 732-763. Il le résume dans \italique{Der Neue Pauly}, 12/2 (2003), col. 831. Pour une étude plus générale de l'ensemble du corpus des anciens lexiques en langue grecque, l'on consultera la synthèse du même auteur dans l'\italique{Historisches Wörterbuch der Rhetorik}, 5 (2001) à l'entrée \allemand{Lexikographie} B, coll. 194-210.}%
\note{11}{Contrairement à l'usage de notre commentaire, nous n'avons réservé Hésychius que pour la fin de cette note. La raison en est que la glose de l'unicodex -- un \italique{Marcianus} du \nieme{XV} siècle -- est légèrement corrompue et donc sujette à caution :%
\citComm{καλπάγει ὀξυπόδης}%
(Le lexique de Cyrille ajoute : σακκιζει (sic, immo βαδίζει).)}%
\note{12}{Cf. le passage de l'\italique{Oreste} d'Euripide, vers 1550, où Ménélas furieux accourt ὀξύπους ($=$ ταχύπους). Ce sens est celui rapporté avec raison par le \LSJ. Le mot latin \latin{addensator}, \italique{celui qui rend plus serré} (en parlant du foulon et de son pas), est glosé dans certains manuscrits gréco-latin ὀξυποδητής.}%
%\note{6}{Et quoi qu'en pense Suidas (κ 623), qui aura eu la même idée qu'Hermann -- on a suffisamment vu en chemin à quoi s'en tenir avec certaines entrées de la \italique{Souda} :%
%\citComm{καλπάζειν· τὸ ἁβρῶς βαδίζειν.}%
%Mais digressons un tout petit peu, si le lecteur me le permet, et constatons amusés à quel point les modernes n'ont rien inventé : Jean Zonaras en son temps (\nieme{XII} siècle) était déjà gêné par la définition que Suidas donnait à καλπάζειν, incompatible avec le sens de "faire des pointes en danse" qu'il voulait lui plaquer. Il trouva une bonne astuce pour son propre lexique : substituer ἄκρως à ἁβρῶς -- solution si ingénieuse qu'on s'étonne que les modernes n'y aient pas pensé ! On ne sera donc pas surpris de lire à la page  : %
%\citComm{καλπάζειν· τὸ ἄκρως βαδίζειν.}}%
%\note{7}{Cf. par exemple l'arrivée tendue des ambassadeurs laconiens dans \italique{Lysistrata} (vers 1073-1074 ; avec la présence concomitante de πρῶτα μέν / εἶτα et de deux impératifs) :%
%\citComm{Ἄνδρες Λάκωνες, πρῶτα μέν μοι χαίρετε,\\
%εἶτ’ εἴπαθ’ ἡμῖν πῶς ἔχοντες ἥκετε.}}%
%\note{8}{L'on a déjà au vers 570 la forte ressemblance dans certains manuscrits du mu et du kappa.}
\\
σὺ δ᾽ ὦ Τερηδὸν ἐπαναφύσα Περσικόν.\testimonia{1175}{Suidas τ \espTest (Περσικόν) Pollux IV 100 ; cf. Hesychius π 2005 (?)}\\
\commentaire{1175}{Τερηδὼν,}{ἡ αὐλητρίς. (Note du scholiaste ; pourquoi pense-t-il que Térédon soit \italique{une} joueuse de flûte ?)}%
\commentaire{1175}{Περσικόν.}{Euripide demande au flûtiste d'accompagner la danseuse avec un air \italique{persan} -- les oreilles ? Or il existe une danse du même nom\marqueur{1} que Xénophon nous décrit dans son \italique{Anabase} (VI, 1, 10) : %
\citComm{μετὰ τοῦτο Μυσὸς\marqueur{2} εἰσῆλθεν ἐν ἑκατέρᾳ τῇ χειρὶ ἔχων πέλτην, καὶ τοτὲ μὲν ὡς δύο ἀντιταττομένων μιμούμενος ὠρχεῖτο, τοτὲ δὲ ὡς πρὸς ἕνα ἐχρῆτο ταῖς πέλταις, τοτὲ δ᾽ ἐδινεῖτο καὶ ἐξεκυβίστα ἔχων τὰς πέλτας, ὥστε ὄψιν καλὴν φαίνεσθαι. τέλος δὲ \textbf{τὸ περσικὸν ὠρχεῖτο} κρούων τὰς πέλτας καὶ ὤκλαζε καὶ ἐξανίστατο· καὶ ταῦτα πάντα ἐν ῥυθμῷ ἐποίει πρὸς τὸν αὐλόν.\\
Un Mysien\marqueur{2} entra ensuite ; il tenait, dans chacune de ses mains, un bouclier léger ; quelquefois il s'en servait en dansant, comme s'il eût eu à se défendre contre deux adversaires ; quelquefois, comme s'il n'eût eu affaire qu'à un seul. Souvent il tournait et faisait le saut périlleux sans lâcher ses boucliers : c'était un spectacle agréable qu'il donnait ; il finit par \textbf{danser la Persique} ; frappant d'un bouclier sur l'autre ; il s'accroupissait, se relevait, et exécutait tous ces mouvements en mesure et au son de la flûte.}%
L'accroupissement (ὀκλάζω) est le mouvement principal de cette danse et lui a donné son nom ; c'est le scholiaste aux \italique{Thesmophories} qui nous l'apprend : %
\citComm{Περσικὸν τὸ ὄκλασμα καλεῖται.}%
Pollux (IV, 100), compulsant un peu vite les \italique{Thesmophories} et ses commentateurs, croit que le terme ὄκλασμα apparaît dans le texte-même d'Aristophane : il range l'ὄκλασμα auprès du κνισμός parmi les danses guerrières et frénétique :%
\citComm{καὶ ὄκλασμα· οὕτω γὰρ ἐν Θεσμοφοριαζούσαις ὀνομάζεται τὸ ὄρχημα τὸ Περσικὸν καὶ σύντονον. τὴνδ' αὐτὴν καὶ ὑγρὰν ὠνόμαζον.}%
L'on trouve d'ailleurs l'une des scholies des \italique{Thesmophories} les plus savantes à notre vers 1175\marqueur{3} :%
\citComm{Περσικὸν ὄρχημα, περὶ οὗ Ἰόβας μακρὸν πεποίηται λόγον ἐν τοῖς περὶ τῆς θεατρικῆς ἱστορίας. ὥστε λελύσθαι τὴν Σελεύκοὺ πρότασιν. προτείνει γὰρ ἐν τῷ πρὸς Ζήνωνα προτατικῷ τὸ Περσικὸν ὄρχημα.}%
Pour ma part, je doute que l'exercice auquel se livre notre danseuse Ἐλάφιον soit cette danse guerrière illustrée par le Mysien de Xénophon ; l'adjectif Περσικός était assez commun pour qu'il pût désigner à la fois un air de flûte et une forme de danse sans qu'il y eut nécessairement de rapport entre les deux.}%
\debutNote{1175}%
\note{1}{Hésychius (π 2005) : περσικά· [ ... ] καὶ ὄρχησις Περσική.}%
\note{2}{Je cite longuement le passage de Xénophon pour faire apparaître la race du danseur. Dans le cas où la danse persique aurait effectivement un rapport avec le verbe ἀνακαλπάζειν du vers 1174 -- ce que je ne crois pas --, on notera qu'il faisait partie du peuple d'Asie qui servit de titre à la pièce d'Eschyle, citée par Photius à l'adresse ἀνακαλπάζει et dont on a parlé au vers précédent.}%
\note{3}{Elle est discutée par Jacques Moreau dans l'article déjà mentionné : «~Note sur le vers 1174 des \italique{Thesmophories} d'Aristophane~», \italique{L'antiquité classique}, tome 7, fasc. 2 (1938) pp. 227-231.}\\\passeLigne}%
\end{groupe}%
%% v. 1176
\begin{groupe}%
%
\personnageGrec{Τοξότης}%
\trimIamb{\ligne{1176}τί τὸ βόμβο τοῦτο; κῶμό τις ἀνεγεῖρί μοι;\apparat{1176}{κωμοτίς ἀν- R : κωμοδρίαν Σ : \lecon{κῶμό τις ἀν-} van Velsen}%
\commentaire{1176}{κωμοτίς ἀνεγεῖρί μοι;}{Variante de lecture du scholiaste : διχῶς κωμοδρίαν.}%
%\commentaire{1176}{}{Le bourdonnement Hésychius : κόμμωσις:ὁ τοῦ σώματος ὥὡραϊσμός. καὶ ὑπὸ τῶν μελισσουργῶν ἡ τοῦ σμήνους διάχρισις}%
\\\passeLigne}%
\end{groupe}%
%% v. 1177-1178
\begin{groupe}%
\personnageGrec{Εὐριπίδης}%
\trimIamb{\ligne{1177}ἡ παῖς ἔμελλε προμελετᾶν ὦ τοξότα.\\
ὀρχησομένη γὰρ ἔρχεθ᾽ ὡς ἄνδρας τινάς.\apparat{1178}{ἔρχετ᾽ R : \lecon{ἔρχεθ᾽} Brunck}\\\passeLigne}%
\end{groupe}%
%% v. 1179-1180
\begin{groupe}%
\personnageGrec{Τοξότης}%
\trimIamb{ὀρκῆσι καὶ μελετῆσι, οὐ κωλῦσ᾽ ἐγώ.\apparat{1179}{κωλύσ᾽ R : \lecon{κωλῦσ᾽} Blaydes (κωλῦσι $=$ κωλύσω), cf. ad 1003, 1196, 1201, 1225 \ligneApp{1180} ὡς et ὥσπερ R : \lecon{ὠς} et \lecon{ὤσπερ} Brunck}\\
ὠς ἐλαπρός, ὤσπερ ψύλλο κατὰ τὸ κῴδιο.%
\commentaire{1172 et 1180}{ὦλάφιον / ὡς ἐλαπρός, ὥσπερ ψύλλο κατὰ τὸ κῴδιο.}{Les verbes signifiant \italique{sauter}, \italique{danser} (πηδᾶν, ἄλλεσθαι) attirent presque inévitablement une comparaison, et le plus souvent avec un animal.
Les jeunes filles de Sparte dansent en bondissant comme des pouliches, ἇτε πῶλοι ταὶ κόραι / πὰρ τὸν Εὐρώταν / ἀμπαδίοντι (ἀμπαδίοντι Coulon : ἀμπάλλοντι codd.) (\italique{Lysistr.} 1308)\marqueur{1}. 
La Muse laconienne danse et saute comme une biche, ποδοῖν τε πάδη / ἇ τις ἔλαφος (\italique{ib}. 1317)\marqueur{2}.
La danseuse des \italique{Thesmophories} porte d'ailleurs le sobriquet de \italique{Bichette} (Ἐλάφιον, v. 1172).
Bichette danse avec tant de grâce devant l'archer Scythe qu'il s'écrie : « qu'elle est légère ! une vraie puce sur ma peau de mouton » (\italique{Thesm}. 1180). Cette comparaison dénuée de poésie sent le corps de garde et ses peaux de mouton où grouille la vermine (on dormait communément sur des \italique{peaux de mouton}, κῴδια). (Note de Taillardat, §791.)}%
\debutNote{1172 et 1180}%
\note{1}{Comparaison banale, cf. Eur. \italique{Or}. 45 (à propos d'Oreste furieux): πηδᾷ δρομαῖος, πῶλος ὣς ὑπὸ ζυγοῦ ; Eubul. 75, 6 $[=$ Athénée, III, 108ab$]$. Le πωλίον symbolise la vivacité, Ar. \italique{Guêp}. 188. En outre, πῶλος désigne couramment les jeunes femmes : Anacréon 88 \italique{Diehl} [$=$ 417 \italique{Poetae Melici Graeci} Page$]$ ; Eur. \italique{Héc.} 142, \italique{Hipp}. 516, etc. ; Cratin. 87 $[=$ 94 Kassel-Austin$]$ ; Épicrat. 9, 4; Eubul. 84 $[=$ 82 Kassel-Austin$]$.}%
\note{2}{Le cerf (ἔλαφος) est le symbole de l'agilité, \italique{Nuées} 351. (Et aussi, comme chez Homère, de la lâcheté.)}\\\passeLigne}%
\end{groupe}%
%% v. 1181-1183a
\begin{groupe}%
\personnageGrec{Εὐριπίδης}%
\trimIamb{φέρε θοἰμάτιον ἄνωθεν ὦ τέκνον τοδί·%
\commentaire{1181}{φέρε θοἰμάτιον ἄνωθεν.}{Sur le texte et le sens de ce vers 1181, voir la note 28 d'Alphonse Willems, reproduite ci-après.}\\
καθιζομένη δ᾽ ἐπὶ τοῖσι γόνασι τοῦ Σκύθου\apparat{1182}{τοῖς R : \lecon{τοῖσι} Bentley}\\
τὼ πόδε πρότεινον, ἵν᾽ ὑπολύσω.\testimonia{1183}{(τὼ πόδε) Hesychius τ 1758 τὼ πόδε· τοὺς πόδας δυϊκῶς}\\\passeLigne}%
\end{groupe}%
%% v. 1183b-1185
\begin{groupe}%
\personnageGrec{Τοξότης}%
\trimIamb{\ligne{1183}\decalage{τὼ πόδε πρότεινον, ἵν᾽ ὑπολύσω.}ναῖκι ναὶ\\
%
κάτησο κάτησο, ναίκι ναί, τυγάτριον.\apparat{1184}{\lecon{ναίκι ναὶ} R : ναίκι ναίκι Bothe}\\
%
οἴμ᾽ ὠς στέριπο τὸ τιττί᾽, ὤσπερ γογγυλί.%
\apparat{1185}{ὡ στεριπὸ (ex σπεριπο) \Rpc : \lecon{ὡς στεριπὸ} (sic) \Mon \espApp ὥσπερ R : \lecon{ὤσπερ} Brunck \espApp \lecon{γογγυλί} \PsiTrois, Enger : γογγύλῃ R}%
\commentaire{1185}{ὡς στεριπὸ τὸ τιττί᾽}{$=$ ὡς στέριφον τὸ τιτθίον.}%
\commentaire{1185-93}{Papyrus d'Oxyrhynque n° 3840.}{Un mince débris contenant quelques lettres de notre passage fut déniché parmi les papyrus d'Oxyrhynque\marqueur{1}. Il y a néanmoins quelques informations à en tirer : il confirme la conjecture d'Enger au vers 1185 (γογγυλί à la place de γογγύλῃ), ainsi que μέλι au vers 1192. Parsons note la présence d'une ligne supplémentaire après le vers 1187, là où le \italique{Ravennas} possède l'une de ses rares didascalies\marqueur{2} : ἀνακύπτη· καὶ παρακύπτι· ἀπεψωλημένος· (même jeu de scène avec l'olisbos qu'aux vers 643-648).}%
\debutNote{1185-93}%
\note{1}{\italique{The Oxyrhynchus Papyri}, vol. LVI (1989), n°3840. Étudié par Peter J. Parsons, aidé de Colin Austin -- déjà en train de préparer son édition des \italique{Thesmophories}. La transcription qu'en fit Parsons et son analyse sont recopiées en annexe. Voir aussi le papyrus d'Oxyrhynque n° 5132 (\italique{The Oxyrhynchus Papyri} LXXVIII) pour un semblant d'\italique{addendum} : la présence de doubles points `$ : $', comme au vers 1190, pour indiquer un changement de locuteur est confirmée par un autre fragment du même papyrus au vers 1209.}%
\note{2}{Pour notre pièce où les scholies très anciennes et érudites se font rares, on ne relève d'autres didascalies qu'aux vers 100, 129 et 276.}\\\passeLigne}%
\end{groupe}%
%% v. 1186
\begin{groupe}%
%
\personnageGrec{Εὐριπίδης}%
\trimIamb{αὔλει σὺ θᾶττον· ἔτι δέδοικας τὸν Σκύθην;\\\passeLigne}%
\end{groupe}%
%% v. 1187-1189
\begin{groupe}%
\personnageGrec{Τοξότης}%
\trimIamb{καλό γε τὸ πυγή. κλαῦσί γ᾽ ἢν μὴ 'νδον μένῃς.\testimonia{1187}{(πυγή) Hesychius π 4279 πυγή· τὸ κάθισμα}\apparat{1187}{κλαυσ\inRas{᾽ ε}{}ϊ γ᾽ \Rdeux : κλαυσεῖ γ᾽ \Mon \espApp ἂν R : \lecon{ἢν} Dindorf (1839) \ligneApp{post 1187} μ$]$\douteux{ε}νο\douteux{ }$[$ \PsiTrois : ἀνακύπτη· καὶ παρακύπτι· ἀπεψωλημένος· R : \lecon{del.} Ellebodius \ligneApp{1188} τὸ \inRas{σ}{}ποστίον \Rpc : τὸ ποστίον \Mon \ \ligneApp{1189}  'στιν R : \lecon{'στὶ} \Mon}%
\commentaire{1187}{κλαῦσί γ᾽ ἂν μὴ 'νδον μένῃς.}{πρὸς τὸ πέος λέγει ἐπιπλήττων ἐντεινομένῳ. (Note du scholiaste.)\\
L'archer suit d'un œil émérillonné tous les mouvements de la danseuse. Son agitation se trahit d'une manière par trop visible. Le Scythe, qui s'en aperçoit, reprend « gare à toi, etc .. »}\\
\ligne{1188}εἶεν· καλὴ τὸ σκῆμα περὶ τὸ πόστιον.%
%\commentaire{1188}{}{\eff{Cf. fr. ... περι το σπηλαιον / σπειον ... et vers ... περι το τραμιν}}
\\\passeLigne}%
\end{groupe}%
%% v. 1189-1190a
\begin{groupe}%
\personnageGrec{Εὐριπίδης}%
\trimIamb{καλῶς ἔχει. λαβὲ θοἰμάτιον· ὥρα 'στὶ νῷν\\%\apparat{1189}{'στιν R : \lecon{'στὶ} \Mon}\\%
ἤδη βαδίζειν.\\\passeLigne}%
\end{groupe}%
%% v. 1190b
\begin{groupe}%
\personnageGrec{Τοξότης}%
\trimIamb{\ligne{1190}\decalage{ἤδη βαδίζειν.}οὐκὶ πιλῆσι πρῶτά με;\apparat{1190}{τί οὐκ ἐπιλήσει R : \lecon{οὐκὶ} (Biset) \lecon{πιλῆσι} Porson (ad \italique{Aristophanica} p.222), cf. Σ `οὐχὶ φιλήσειν'}\\\passeLigne}%
\end{groupe}%
%% v. 1191a
\begin{groupe}%
\personnageGrec{Εὐριπίδης}%
\trimIamb{πάνυ γε· φίλησον αὐτόν.\\\passeLigne}%
\end{groupe}%
%% v. 1191b-1193a
\begin{groupe}%
\personnageGrec{Τοξότης}%
\trimIamb{\ligne{1191}\decalage{πάνυ γε· φίλησον αὐτόν.}ὂ ὂ ὂ παπαπαπαῖ,\apparat{1191}{ὂ R : \lecon{ὂ ὂ ὂ} \Mon}\\
ὠς γλυκερὸ τὸ γλῶσσ᾽, ὤσπερ Ἀττικὸς μέλι.\apparat{1192}{ὡς et ὥσπερ R : \lecon{ὠς} et \lecon{ὤσπερ} Brunck \espApp \lecon{μέλι} \PsiTrois, \Mon : μέλις R}%
\commentaire{1192}{ὠς γλυκερὸ τὸ γλῶσσ᾽, ὤσπερ Ἀττικὸς μέλις.}{Si l'image des baisers plus doux que le miel\marqueur{1} ou le nectar nous paraît banale, c'est à une longue tradition héritée de ces culs-terreux de latins que nous le devons\marqueur{2} ; pour ce qui est des auteurs de langue grecque, on ne peut guère citer que Alciphron (I, 38, 7), postérieur d'environ six cents ans à Aristophane. 
Les spectateurs d'Aristophane devaient ressentir avec cette comparaison une confusion burlesque : l'épithète poétique μελίγλωσσος existe en grec, mais l'image d'une \italique{langue pleine de miel} signifie tout autre chose. 
Ainsi parlait Taillardat (§§739-743) :\\
«~Le miel étant le symbole de la douceur\marqueur{3}, les poètes se plaisent à dire que leurs chants sont de miel\marqueur{4}. Cette image est chez Aristophane : μελιγλώσσων ἐπέων ... ἀοιδάν (\italique{Ois}. 908) donne un exemple de plus de l'épithète banale μελίγλωσσος\marqueur{5} $[$ ... $]$
On dira donc tout naturellement que la lèvre d'un poète est enduite de miel : c'est le cas de la bouche de Sophocle, Σοφοκλέους τοῦ μέλιτι κεχριµένου... τὸ στόμα (\italique{fr}. Aristophane 581 $[=$ 598 Kassel-Austin$]$). Composant des poèmes doux comme le miel, les poètes aiment à se comparer aux abeilles, ainsi qu'en témoigne Platon (\italique{Ion}, 534b).~»\\
Aristophane fait rire son public aux dépens de ces images trop usées : le Scythe rustique vante la \italique{langue douce comme le miel} de la danseuse, mais c'est une langue toute matérielle -- et qui appartient d'ailleurs à une actrice muette. Il y a des impropriétés de termes qui sont des malpropretés du cœur.}%
\debutNote{1192}%
\note{1}{Et l'Attique est souvent citée pour son miel. Cf. la \italique{Paix}, vers 252-254 ; Antiphane fr. 177(179) Kassel-Austin $=$ Athénée II, 43b \& III, 74d ; le corpus hippocratique mentionne maintes fois du miel Attique.}%
\note{2}{Cf. le \italique{Dictionnaire des sentences latines et grecques} de Renzo Tosi, §1385 \latin{Nectar dulcius} : Catulle 99, 1-2 ; Horace, \italique{Carmina} 1, 13, 14-16 ; \&c. Honoré de Balzac a fameusement repris cette longue tradition : « Les baisers d'une femme sincère ont un miel divin qui semble mettre dans cette caresse une âme, un feu subtil par lequel le cœur est pénétré. Les baisers dénués de cette onction savoureuse sont âpres et secs. »}%
\note{3}{\italique{Guêpes} 878 (voir §361), \italique{Thesm.} 1192, fr. 587, 622 $[=$ 681, 633 K.-A.$]$ ; Alexis 146, 6 $[=$ 150.6 K.-A.$]$ ; Ménand. 708 (?), etc.}%
\note{4}{Pind. \italique{Ol.} 11, 4 : μελιγάρυες ὕμνοι ; id. \italique{Ném.} 11,18 : μελιγδούποισι ... μελιζέμεν ἀοιδαῖς ; id. \italique{Isthm.} 2, 32 : μελικόμπων ἀοιδᾶν ; Alcman 41 $[=$ Athénée XIV, 633a $=$ fr. adesp. 954b \italique{Poetae Melici Graeci}, ed. Page$]$ : μέλεα μελιπτέρωτα ; Simonid. 40 D $[=$ fr 595.3 \italique{Poetae Melici Graeci} Page$]$ : μελιαδέα γᾶρυν, etc. On notera que, déjà chez Homère (\italique{Il.} I, 249), la voix de Nestor est dite $[$ἀπὸ γλώσσης$]$ μέλιτος γλυκίων. D'autre part Pindare, \italique{Ol.} 7, 8, appelle encore la poésie un nectar, νέκταρ.}%
\note{5}{Eschyl. \italique{Prom.} 172 ; Bacchyl. 3, 97 ; fr. 4, 25 Snell (μελιγλώσσων ἀοιδᾶν ἄνθεα).}\\
τί οὐ κατεύδει παρ᾽ ἐμέ;\\\passeLigne}%
\end{groupe}%
%% v. 1193b-1194a
\begin{groupe}%
\personnageGrec{Εὐριπίδης}%
\trimIamb{\ligne{1193}\decalage{τί οὐ κατεύδει παρ᾽ ἐμέ;}χαῖρε τοξότα,\\
%
οὐ γὰρ γένοιτ᾽ ἂν τοῦτο.\\\passeLigne}%
\end{groupe}%
%% v. 1194b-1195a
\begin{groupe}%
\personnageGrec{Τοξότης}%
\trimIamb{\ligne{1194}\decalage{οὐ γὰρ γένοιτ᾽ ἂν τοῦτο.}ναὶ ναὶ γρᾴδιο.\apparat{1194}{ναὶ R : ναίκι Bentley : \lecon{ναὶ ναὶ} Brunck \espApp γράδιον R : \lecon{γρᾴδιο} Bothe \ligneApp{1195} \inRas{σὺ}{} κάρισοσ οὐ \Rpc : κάρισο σοῦ \Mon : \lecon{κάρισο σὺ} Brunck}\\
ἐμοὶ κάρισο σὺ τοῦτο.%
\commentaire{1195}{ἐμοὶ κάρισο σὺ τοῦτο}{$=$ ἐμοὶ χάρισαι σὺ τοῦτο.}\\\passeLigne}%
\end{groupe}%
%% v. 1195b
\begin{groupe}%
\personnageGrec{Εὐριπίδης}%
\trimIamb{\ligne{1195}\decalage{ἐμοὶ κάρισο σὺ τοῦτο.}δώσεις οὖν δραχμήν;%
\\\passeLigne}%
\end{groupe}%
%% v. 1196a
\begin{groupe}%
\personnageGrec{Τοξότης}%
\trimIamb{ναὶ ναῖκι δῶσι.\\\passeLigne}%
\end{groupe}%
%% v. 1196b
\begin{groupe}%
\personnageGrec{Εὐριπίδης}%
\trimIamb{\ligne{1196}\decalage{ναὶ ναῖκι δῶσι.}τἀργύριον τοίνυν φέρε.\\\passeLigne}%
\end{groupe}%
%% v. 1197
\begin{groupe}%
\personnageGrec{Τοξότης}%
\trimIamb{ἀλλ᾽ οὐκ ἔκὠδέν· ἀλλὰ τὸ συμβήνην λαβέ.\testimonia{1197}{Photius σ 660 συβήνη $=$ Suidas σ 1273 (sed ἐχω οὐδέν pro έκ' ὠδέν, τὸ omisso et συβήνην pro συμβήνην) \espTest (συμβήνην) Hesychius σ 2137 συμβίνης· καπροβόλον. ἐμβόλιον}\apparat{1197}{έκ' ὠδέν R : ἐχω οὐδέν Suid. σ 1273, Photius σ 660 : \lecon{ἔκὠδέν} Bergk \espApp \lecon{συμβήνην} R : συβήνην Σ, Suid., Photius \ligneApp{1198} κ\inRas{ω}{ο}μίζεις αὐτοῖς \Rpc : κομιεῖς αὐτό (αὖθις Brunck) Biset : (Τοξότ.) κομίσι σ' αὐτός (αὖτις anonym. Parisinus) Enger : (Τοξότ.) \lecon{κομίζι σ᾽ αὐτίκ᾽} Bibi \espApp ἀκουλούτι R : \lecon{ἀκολούτι} \Mon : ἀκουλούτει Brunck}%
\commentaire{1197-216}{ἀλλὰ τὸ συμβήνην λαβέ $[$ ... $]$ ὀρτῶς δὲ συβήνη 'στί· καταβηνῆσι γάρ.}{Ce passage a été si maltraité par les commentateurs -- anciens comme modernes -- que nous nous sommes cru obligé à commettre cette note rectificatrice. \italique{Recte}, descendons-donc dans l'arène (ὀρτῶς δὲ συμβήνη 'στί· καταβηνῆσι γάρ.)\\
L'origine de tous les malentendus est la glose d'un commentateur ancien qui sera allé un peu vite : le \italique{Ravennas} comporte deux scholies pour définir la συμβήνην (la première scholie se trouve à l'intérieur même du texte) :%
\citComm{τὴν τοξοθήκην\\
συβήνη· αὐλοθήκη. λέγουσι δὲ καὶ τὸν φαρετρεῶνα συβήνην.}%
On retrouve cette scholie chez Hésychius (σ 2136) :%
\citComm{συβήνη· αὐλοθήκη, ἢ τοξοθήκη. ἢ ὁ ναυτικὸς χιτών.}%
Suidas cite notre passage dans une glose qu'on croise à l'identique, ou peu s'en faut, dans d'autres lexiques\marqueur{1}:%
\citComm{Συβήνη: ἡ δερματίνη αὐλοθήκη\marqueur{2}. ἢ ἡ φαρέτρα. Θεσμοφοριαζούσαις Ἀριστοφάνης· ἀλλ' οὐκ ἔχω οὐδέν. ἀλλὰ συβήνην λάβε.}%
Cette unanimité ne doit pas nous tromper : elle prouve seulement l'ancienneté de la première rédaction qui aura diversement infusé dans de nombreux lexiques. Je dirais même plus : ces gloses concordent si bien qu'il est certain qu'elles remontent toutes à une même source -- certes ancienne -- mais unique, et donc faillible.\\
Que le terme συβήνη -- qui n'est d'ailleurs pas celui transmis par notre manuscrit comportant συ\textbf{μ}βήνην --, que συβήνη ait pu désigner un \italique{étui à flûte}, je n'en disconviens pas. À la rigueur, que par dérivation le mot ait aussi pu être appliqué à un \italique{carquois}, pourquoi pas\marqueur{3} ? Mais il serait absurde de lui donner ce sens dans notre passage : croit-on que l'Archer, tout Scythe qu'il fût, pensait un seul instant qu'une vieille maquerelle pût accepter son carquois dont elle n'aurait que faire au lieu d'une drachme sonnante et trébuchante ? ce serait tout bonnement ridicule\marqueur{4} ! Le Scythe, Euripide, Aristophane, le savent bien ; en fait n'importe quel homme, pourvu qu'il ne fût pas versé dans l'exégèse d'Aristophane : %
\citComm{Chez la belle Suzon, pas d'argent, pas de cuisse...\marqueur{5}}%
En réalité, le Scythe demande à Euripide de garder son prisonnier quelques instants, le temps qu'il aille chercher l'argent nécessaire.
Encore une fois, c'est Hésychius (σ 2137) qui détenait la solution :%
\citComm{συμβίνης· καπροβόλον. ἐμβόλιον.}%
Le mot est apparenté à σιβύνη, une \italique{lance} (Hésychius nous précise pour la chasse au \italique{sanglier}, κάπρος), dont il est une légère déformation qui prépare le jeu de mots scabreux du vers 1215.
Le texte est aussi clair qu'un ciel d'azur au milieu de l'été : le Scythe confie sa lance à Euripide pour la garde de Mnésiloque.
Et nous nous garderons bien de biffer le μ de la leçon du manuscrit, συ\textbf{μ}βήνην, qui aura survécu à plus de mille ans de tradition manuscrite mais pas à la griffe implacable des éditeurs d'Aristophane.\\
Le début du vers suivant, altéré, ἔπειτα †κ\inRas{ω}{ο}μίζεις αὐτοῖς†, se laisse lire sans difficulté :%
\citComm{ἔπειτα κομίζι ($=$ κομίζω) σ(οὶ) αὐτίκ(α)\\
Ensuite je te l'apporte $[$la drachme$]$ sur-le-champ.}%
La dernière version du manuscrit attribue, il est vrai, cette réplique à Euripide. Ces indications marginales de personnages ne doivent pas nous arrêter : elles sont récentes et même pour certaines du \nieme{XVI} siècle\marqueur{6} (c'est le cas de cette annotation, «~γϋ.~», que porte le manuscrit (la même qu'au vers 1189) et mis pour γυνή -- c'est-à-dire la \italique{vieille}, γραῦς ; c'est-à-dire Euripide. Von Velsen ne l'a pas intégrée à son apparat critique et avait donc reconnu l'écriture d'Euphrosynus Boninus).\\
Notre émendation appelle quelque justification : l'hypothétique confusion αυτικ / αυτοις peut aisément s'expliquer en constatant la ressemblance entre le digramme `ις' et le `κ'\marqueur{7}. Continuons.\\
Le Scythe part chercher l'argent, accompagné de la jeune beauté, qui se montre si caressante selon les instructions d'Euripide (vers 1172 ; l'attitude de la danseuse est la version muette de celle de la Myrrhine dans \italique{Lysistrata}) qu'il en revient tout émoustillé et en proie à une promesse de félicité qui point\marqueur{8}. Il déchante très vite : Euripide et Mnésiloque ont pris la poudre d'escampette. Élaphion décampe à ce moment sous les imprécations du Scythe (ἀπότρεκ᾽ ὠς τάκιστα σύ).
Le Scythe s'arme donc de sa lance, laissée sur place, et s'apprête à descendre (καταβαίνειν) à leur poursuite :%
\citComm{ὀρτῶς δὲ συμβήνη ᾽στί· καταβηνῆσι γάρ.}%
Je n'explique pas le trait d'esprit -- d'esprit \italique{Scythe} ; si vous ne le comprenez pas, laissez-là Aristophane et allez faire des mathématiques !\\
Nul besoin de retailler le vers : καταβαίνω est le terme consacré lorsqu'un acteur descend de la σκηνή vers l'ὀρχήστρα, d'autant plus adapté quand la scène se situe sur un lieu en hauteur comme la colline de la Pnyx\marqueur{9} ; cf. les \italique{Guêpes} (vers 1514) ou l'\italique{Assemblée des femmes} (vers 1152)\marqueur{10}.
%Ces savants hommes ne manquent jamais de mettre leur érudition en branle dès qu'ils subodorent derrière la plus anodine des répliques le relent du pornographe Aristophane... et passent sans s'arrêter quand le plus grossier de ses personnages paye une maquerelle d'un écot -- à défaut d'un écu -- au nom sonore : écho de *συμβινεῖν. (\italique{Confer} le vers 1215.)
}%
\debutNote{1197-216}%
\note{1}{Photius (σ 660), \latin{Etymologicum magnum}, \latin{genuinum}, Zonaras, \&c. Pollux (VII, 153 ; X, 153) : καὶ τὸ μὲν ἀγγεῖον τὸ τὰς γλώττας ἔχον γλωττοκομεῖον, ἣ δὲ τῶν αὐλῶν θήκη συβήνη. Pour un dossier plus complet, je renvoie le lecteur à la \allemand{Realencyclopädie} : \italique{Paulys Realencyclopädie der classischen Altertumswissenschaft}, IV 1 (1931), coll. 1011-1012.}%
\note{2}{Dans certains manuscrits du \italique{Zonarae Lexicon} (p. 1686), on trouve βελοθήκη en place d'αὐλοθήκη.}%
\note{3}{Encore qu'il me semble plus vraisemblable que ce sens ait été inféré par notre scholiaste, puis scrupuleusement intégré à tous les lexiques.}%
\note{4}{À tous les vulgarisateurs d'Aristophane, tournebouleurs et avilisseurs de ses onze chefs-d'œuvre, rappelons la maxime de bon goût qu'Aristophane s'imposait :%
\citComm{ὡς ἐγὼ φοβοῦμαι περὶ τῶν μελλόντων ῥηθήσεσθαι, οὔ τι μὴ γελοῖα εἴπω — τοῦτο μὲν γὰρ ἂν κέρδος εἴη καὶ τῆς ἡμετέρας μούσης ἐπιχώριον — ἀλλὰ μὴ καταγέλαστα.\\
Ce n'est pas de provoquer le rire que je redoute -- cela serait tout gain, et conforme à ma muse -- mais de dire des choses ridicules.\\
(Trad. P. Boutang.)}%
(Tradition reçue en droite ligne d'Aristodème, témoin des événements, qui en fit le récit à Apollodore, qui lui-même le transmit à Glaucon, frère de Platon, qui -- on peut l'imaginer ainsi -- confia bientôt, après peut-être encore quelques méandres, à son frère Platon le précieux compte-rendu, lequel Platon finit par coucher par écrit cette fameuse nuit de \italique{Banquet}, 189b, quelques quarante années après les faits.)}%
\note{5}{Le proverbe d'origine détourné par Brassens (\italique{Grand-père}) étant «~point d'argent, point de Suisse~». Quel sera l'homme d'assez bonne volonté pour nous offrir le colossal équivalent brassinien du maître-œuvre de Taillardat, quel est le preux qui prendra à bras-le-corps les quelques deux cents chansons de Brassens que ce dernier a bien voulu nous conserver ? Espérons qu'il ne soit pas besoin d'attendre encore vingt-cinq siècles comme pour Aristophane. (C'est quand le vers est troussé qu'il faut l'analyser.) Je tiens à la disposition de cette bonne âme éventuelle un matériau qui poserait de solides fondations à ces \italique{Images de Brassens} en fomentation.}%
%\note{6}{\eff{Retrouver la réf sur l'armement des Scythes ; Hall ?}}%
\note{6}{Cf. \italique{supra}, v. 1059 n. 5.}%
\note{7}{Il faudrait alors supposer l'existence d'une forme intermédiaire αυτιις ou αυτις, bien vite corrigée en αυτοις. Une forme intermédiaire, ou un \italique{chaînon manquant}, pour reprendre le jargon des darwiniens. Ne poussons donc pas plus loin la «~preuve~» afin d'éviter l'écueil de ces demi-savants, capables de tout prouver et même son contraire !}%
\note{8}{Savez-vous ce qu'on peut lire dans le commentaire le plus universitairement reçu, le plus académiquement académisant -- celui d'Austin \& Olson, puisqu'il faut le nommer ? Je vous le donne en mille : le Scythe a consommé la belle Élaphion... il revient donc, précisent leurs auteurs, avec l'\italique{olisbos} abaissé (!!!) Quand on sait -- ou qu'on ne sait pas, comme Austin \& Olson -- combien Aristophane adorait ses chères actrices, qu'il ne manquait jamais de rendre grâce à leur beauté, toujours plein d'attention et de prévenance vis-à-vis de la voracité de la lie athénienne... un archer Scythe... la lie de la lie... suggérer une telle monstruosité ?! «~Paraît qu'elle s'vend même à des flics. / Quelle décadence ! / Y'a plus d'moralité publique / Dans notre France...~» En France peut-être, mais au théâtre d'Aristophane, certainement pas ! (Lisez la note de Willems, t. III pp. 394-395, in « Le nu dans la Comédie ancienne des Grecs », article reproduit en annexe ; purgatif si nécessaire quand on est souillé par de telles salauderies !) A\&O ne nous indiquent pas si le public pouvait se délecter en sus de bruitages émanant des coulisses, quand ça y fornique... Mais pour qui prennent-ils Aristophane ?.. voici donc la littérature qu'on inflige à nos pauvres étudiants... je tiens le cas pour pendable !%
\citComm{et si, par un malheur, j'en avois fait autant,\\
je m'irois, de regret, pendre tout à l'instant.}}%
\note{9}{Le Scythe s'interrompt quelques instants de parler, pause soulignée par l'interjection \italique{extra metrum} οἴμοι qu'il lance une fois descendu : il se demande vers laquelle des deux entrées de l'orchestre se tourner (τί δρᾶσι; ποῖ τὸ γρᾴδι᾽;)}%
\note{10}{\italique{Confer} surtout la note d'Alphonse Willems au vers 1341 des \italique{Guêpes}. Lisez Willems ! aucune excuse : la note est reproduite en annexe.}\\
ἔπειτα κομίζι σ᾽ αὐτίκ᾽. ἀκολούτι τέκνον.%\apparat{1198}{%γϋ. ante v. 1198 et σκ'. ante ἀκολ. in R \espApp 
%κ\inRas{ω}{ο}μίζεις αὐτοῖς \Rpc : κομιεῖς αὐτό (αὖθις Brunck) Biset : (Τοξότ.) κομίσι σ' αὐτός (αὖτις anonym. Parisinus) Enger : (Τοξότ.) \lecon{κομίζι σ᾽ αὐτίκ᾽} Bibi \espApp ἀκουλούτι R : \lecon{ἀκολούτι} \Mon : ἀκουλούτει Brunck}
\\%
%\commentaire{1198}{ἔπειτα κ\inRas{ω}{ο}μίζεις αὐτοῖς.}{Le †αὐτοῖς† du \italique{Ravennas}, en plus d'être \italique{contra metrum}, n'offre aucun sens acceptable. 
%Les diverses tentatives de correction se résument à prolonger les pourparlers. Tantôt Euripide demande son reste ; l'on traduit alors péniblement par "tu le récupéreras une autre fois \auDessus{=}{?} après" :%
%\citComm{(Εὐρ.) ἔπειτα κομιεῖς αὐτό (Biset ; αὐτό \auDessus{=}{?} τὴν συμβήνην)\\
%(Εὐρ.) ἔπειτα κομιεῖς αὖθις (Faber, Brunck, von Velsen, Austin)}%
%Tantôt attribue-t-on au Scythe la réplique qui promet d'apporter à Euripide le reste de la drachme plus tard, ou bien s'enquiert de pouvoir récupérer son carquois en gage... mais le tour paraît bien elliptique :%
%\citComm{(Τοξότ.) ἔπειτα κομίσι σ' αὐτός (... τὸ λοιπόν ?.. τὸ ἀργύριον ?.. ; Enger)\\
%(Τοξότ.) ἔπειτα κομίζεις αὐτόσ᾽ (Fritzsche)\\
%(Τοξότ.) ἔπειτα κομίσι σ' αὖτις (Coulon, Prato)}%
%Ces interprétations me semblent forcer le grec ; elles perdent aussi de vue l'action : Euripide ne cherche assurément pas à être payé rubis sur l'ongle,
%comme dans la chanson, %
%\citComm{Chez la belle Suzon, pas d'argent, pas de cuisse...\marqueur{1}}%
%C'est à cette réplique que se concluent les négociations puisqu'aussitôt le Scythe emmène Élaphion (ἀκολούτι τέκνον).\\
%Euripide, qui ne cherche qu'à éloigner le Cerbère de Mnésiloque, accepte donc l'offre ridicule du Scythe et lui ordonne de quitter la scène -- l'impératif est ici de rigueur : %
%\citComm{ἔπειτα κόμιζ(ε) εἰσαυτίκ(α).\\
%Eh bien alors emmène-la sur-le-champ.}%
%L'émendation appelle quelques justifications : d'abord, εἰσαυτίκα est autrement attesté au vers 367 de la \italique{Paix} ; quant à l'hypothétique confusion αυτικ / αυτοις, elle n'est pas invraisemblable en constatant la ressemblance entre le digramme `ις' et le `κ' ; il faudrait alors supposer l'existence d'une forme intermédiaire\marqueur{2} αυτιις ou αυτις, bien vite corrigée en αυτοις.}%
%\debutNote{1198}%
%\note{1}{}%
%\note{2}{}
σὺ δὲ τοῦτο τήρει τὴ γέροντο, γρᾴδιο.\\
ὄνομα δέ σοι τί ἔστιν;\\\passeLigne}%
\end{groupe}%
%% v. 1200b
\begin{groupe}%
\personnageGrec{Εὐριπίδης}%
\trimIamb{\ligne{1200}\decalage{ὄνομα δέ σοι τί ἔστιν;}Ἀρτεμισία.\\\passeLigne}%
\end{groupe}%
%% v. 1201
\begin{groupe}%
\personnageGrec{Τοξότης}%
\trimIamb{μεμνῆσι τοίνυν τοὔνομ᾽· Ἀρταμουξία.\apparat{1201}{ἀρτομουξία R : \lecon{Ἀρταμουξία} ed. Lipsiensis (1814)}%
\commentaire{1201}{Ἀρταμουξία.}{L'archer scythifie le nom d'Ἀρτεμισία : il l'entend comme un composé de ἀρτα-, multiplie les `α' et ajoute un `ξ'\marqueur{1}. Cf. le parler de Ψευδ\textbf{αρτά}βας dans les \italique{Acharniens} (v. 100) :%
\citComm{ἰαρταμὰν ἐξάρξαν ἀπισσόνα σάτρα.}%
}%
\debutNote{1201}%
\note{1}{Voir le commentaire vv. 91-2 d'Olson aux \italique{Acharniens} ainsi que l'article «~The Archer Scene in Aristophanes' \italique{Thesmophoriazusae}~» d'Édith M. Hall, in \italique{Philologus} 133 (1989), p. 40. On peut citer Ἀρταβάτας, Ξέρξης, Ἀρταξέρξης, Ἀρταφρένης, Ἀρτάβαζος \&c.}\\\passeLigne}%
\end{groupe}%
%% v. 1202-1206
\begin{groupe}%
\personnageGrec{Εὐριπίδης}%
\trimIamb{Ἑρμῆ δόλιε ταυτὶ μὲν ἔτι καλῶς ποιεῖς.\\
σὺ μὲν οὖν ἀπότρεχε παιδάριον τουτὶ λαβών·\apparat{1203}{\lecon{τουτὶ} R : ταυτὶ Dobree, cf. Σ}%
\commentaire{1203-11}{Papyrus d'Oxyrhynque n° 4935 fr. 2 \& n° 5132.}{Deux fragments d'un même papyrus d'Oxyrhynque\marqueur{1} contiennent quelques lettres de ce passage. (Voir en annexe.) Ils n'apportent aucune leçon originale. Au vers 1208, l'espace manque pour être compatible avec la leçon \italique{ante correctionem} λέλυσον du \italique{Ravennas} ; l'état du fragment ne permet pas pour autant de trancher entre les corrections λέλυσο (\Rpc) et λέλυσαι (Bentley).}%
\debutNote{1203-11}%
\note{1}{N° 4935 fr. 2, in \italique{The Oxyrhynchus Papyri}, LXXIII (2009), ed. A. Benaissa. N° 5132, \italique{The Oxyrhynchus Papyri}, LXXVIII (2012), ed. W. B. Henry.}\\
ἐγὼ δὲ λύσω τόνδε. σὺ δ᾽ ὅπως ἀνδρικῶς\\
ὅταν λυθῇς τάχιστα φεύξει καὶ τενεῖς\\
ὡς τὴν γυναῖκα καὶ τὰ παιδί᾽ οἴκαδε.\\\passeLigne}%
\end{groupe}%
%% v. 1207
\begin{groupe}%
\personnageGrec{Μνησίλοχος}%
\trimIamb{ἐμοὶ μελήσει ταῦτά γ᾽ ἢν ἅπαξ λυθῶ.\\\passeLigne}%
\end{groupe}%
%% v. 1208-1209a
\begin{groupe}%
\personnageGrec{Εὐριπίδης}%
\trimIamb{λέλυσο. σὸν ἔργον, φεῦγε πρὶν τὸν τοξότην\apparat{1208}{ca. 6 litt. $]$· σ\douteux{ο}$[$νεργον \PsiQuatre : \lecon{λέλυσο}\inRas{ν}{} \Rpc : λέλυσαι Bentley}\\
ἥκοντα καταλαβεῖν.\\\passeLigne}%
\end{groupe}%
%% v. 1209b
\begin{groupe}%
\personnageGrec{Μνησίλοχος}%
\trimIamb{\ligne{1209}\decalage{ἥκοντα καταλαβεῖν.}ἐγὼ δὴ τοῦτο δρῶ.\\\passeLigne}%
\end{groupe}%
%% v.1210-1216 (part 1)
\begin{groupe}\personnageGrec{Τοξότης}
\trimIamb{ὦ γρᾴδι᾽ ὠς καρίεντό σοι τὸ τυγάτριο,\apparat{1210}{τυγάτριον ex τυγάθριον ut videtur \Rpc}\\\passeLigne}%
\end{groupe}%
\endnumbering
\end{Leftside}
\begin{Rightside}
\beginnumbering
%% v. 1169+ (part 2)
\begin{groupe}
\trimIambFr{\ligne{1169}à vos maris, à leur retour de l'armée.\\\passeLigne}%
\end{groupe}%
%% v. 1170-1171
\begin{groupe}%
\personnageFr{Coryphée}%
\trimIambFr{Pour ce qui est de nous, sache que nous acquiesçons ; \\
mais ce barbare, c'est à toi de le convaincre.\\\passeLigne}%
\end{groupe}%
%% v. 1172-1175
\begin{groupe}%
\personnageFr{Euripide}%
\trimIambFr{Ça, c'est mon affaire. (À la danseuse ) La tienne, Élaphion,\\
est de songer à t'y prendre comme je t'ai dit chemin faisant.\\
Ainsi donc, passe d'abord à l'autre bout, et reviens en courant sur tes pointes.\\% \commentaire{1074}{ἀνακόλπασων}{Hésych. καλπαζεῖ ὀξυποδεῖ, ou, comme aurait dit Longus (II, 37) : ἐπ' ἄκρων τῶν ὀνύχων τρέχει. G. Hermann, à l'aide du scholiaste, a restitué la vraie leçon ἀνακάλπασων (au lieu de ἀνακόλπασων, qui n'est pas grec).}.\\
(Au flûtiste :) Toi, Térédon, joue-nous sur la flûte une Persique.\\
(Térédon joue de la flûte, pendant que la jeune fille danse.)
\\\passeLigne}%
\end{groupe}%
%% v. 1176
\begin{groupe}%
\personnageFr{Archer}%
\trimIambFr{(Se réveillant.) Qu'est-ce que ce ronflement? c'est un comos qui me réveille ?\\\passeLigne}%
\end{groupe}%
%% v. 1177-1178
\begin{groupe}%
\personnageFr{Euripide}%
\trimIambFr{Cette enfant, archer, pensait s'essayer.\\
Elle s'en va danser chez des particuliers.\\\passeLigne}%
\end{groupe}%
%% v. 1179-1180
\begin{groupe}%
\personnageFr{Archer}%
\trimIambFr{Qu'elle danse, qu'elle s'essaie, je ne m'y oppose pas.\\
Qu'elle est vive! On dirait une puce sur une peau de bique.\\\passeLigne}%
\end{groupe}%
%% v. 1181-1183a
\begin{groupe}%
\personnageFr{Euripide}%
\trimIambFr{Tends-moi ce vêtement, ma fille, en le tirant par en haut.\\
Assieds-toi sur les genoux du Scythe,\\
et avance les pieds, que je les déchausse.\\\passeLigne}%
\end{groupe}%
%% v. 1183b-1185
\begin{groupe}%
\personnageFr{Archer}%
\trimIambFr{\ligne{1183}\espace{4.5cm}Si fait, si,
mets-toi, si fait, ma fillette.\\
%
(Pendant qu'Euripide ôte les chaussures à la danseuse.) Ouais! est-il ferme, ce téton ! une vraie rave. (La jeune fille déchaussée et nue danse au son de la flûte.)
%
\\\passeLigne}%
\end{groupe}%
%% v. 1186
\begin{groupe}%
\personnageFr{Euripide}%
\trimIambFr{Joue donc plus vite. As-tu encore peur du Scythe ?\\\passeLigne}%
\end{groupe}%
%% v. 1187-1189
\begin{groupe}%
\personnageFr{Archer}%
\trimIambFr{Ah! que voilà de belles fesses %\commentaire{1187}{}{L'archer suit d'un œil émérillonné tous les mouvements de la danseuse et se penche tantôt à droite, tantôt à gauche (ἀνακύπτει καὶ παρακύπτει ἀπεψωλημένος\eff{...autre sens... cf . v. scène de déshabillage + citer Willems}]). Son agitation se trahit d'une manière par trop visible. Le Scythe, qui s'en aperçoit, reprend « gare à toi, etc .. »}
! — Gare à toi si tu ne restes dans ton coin.\\
(Il s'arrange.) — À la bonne heure ! Elle a repris de la tenue, la broquette.\\\passeLigne}%
\end{groupe}%
%% v. 1189-1190a
\begin{groupe}%
\personnageFr{Euripide}%
\trimIambFr{(À la danseuse). C'est assez. Reprends ton vêtement. Il est temps que\\
nous partions.\\\passeLigne}%
\end{groupe}%
%% v. 1190b
\begin{groupe}%
\personnageFr{Archer}%
\trimIambFr{\ligne{1190}\espace{2.1cm}Sans me donner un baiser?\\\passeLigne}%
\end{groupe}%
%% v. 1191a
\begin{groupe}%
\personnageFr{Euripide}%
\trimIambFr{Très volontiers. (À la danseuse.) Baise-le.\\\passeLigne}%
\end{groupe}%
%% v. 1191b-1193a
\begin{groupe}%
\personnageFr{Archer}%
\trimIambFr{\ligne{1191}\espace{3.9cm}Oh! oh! oh! papapapai!\\
Que sa langue est douce. C'est comme du miel attique.\\
Pourquoi ne couche-t-elle pas avec moi ?\\\passeLigne}%
\end{groupe}%
%% v. 1193b-1194a
\begin{groupe}%
\personnageFr{Euripide}%
\trimIambFr{\ligne{1193}\espace{3.7cm}Adieu, archer,\\
%
cela ne se peut.\\\passeLigne}%
\end{groupe}%
%% v. 1194b-1195a
\begin{groupe}%
\personnageFr{Archer}%
\trimIambFr{\ligne{1194}\espace{3.7cm}Si, si, petite vieille,\\
fais-moi ce plaisir.\\\passeLigne}%
\end{groupe}%
%% v. 1195b
\begin{groupe}%
\personnageFr{Euripide}%
\trimIambFr{\espace{3.3cm}Me donneras-tu une drachme ?\\\passeLigne}%
\end{groupe}%
%% v. 1196a
\begin{groupe}%
\personnageFr{Archer}%
\trimIambFr{Si fait, si, je te la donnerai.\\\passeLigne}%
\end{groupe}%
%% v. 1196b
\begin{groupe}%
\personnageFr{Euripide}%
\trimIambFr{\ligne{1196}\espace{2.4cm}Voyons ton argent.\\\passeLigne}%
\end{groupe}%
%% v. 1197
\begin{groupe}%
\personnageFr{Archer}%
\trimIambFr{Je n'en ai pas. Mais prends ma lance. (Il la lui confie.)\\
Je te l'apporte sur-le-champ. Suis-moi, mon enfant.\\
(À Euripide :) Toi, petite vieille, garde ce bonhomme.\\
Mais comment te nommes-tu ?\\\passeLigne}%
\end{groupe}%
%% v. 1200b
\begin{groupe}%
\personnageFr{Euripide}%
\trimIambFr{\espace{3.5cm}Artémisia.\\\passeLigne}%
\end{groupe}%
%% v. 1201
\begin{groupe}%
\personnageFr{Archer}%
\trimIambFr{Je le retiendrai, ce nom : Artamouxia. (Il s'en va avec la danseuse.)\\\passeLigne}%
\end{groupe}%
%% v. 1202-1206
\begin{groupe}%
\personnageFr{Euripide}%
\trimIambFr{Hermès, dieu de la ruse, jusqu'ici tu régis tout à souhait.\\
Toi, mon petit, enfuis-toi, avec ces instruments. (Il lui donne ses pectis.)\\
Pendant ce temps, je le détacherai.\\
(À Mnésiloque :) Quant à toi, dès que tu seras libre, vois à fuir résolument\\
et cours retrouver chez toi, ta femme et tes enfants.\\\passeLigne}%
\end{groupe}%
%% v. 1207
\begin{groupe}%
\personnageFr{Mnésiloque}%
\trimIambFr{J'en fais mon affaire, une fois délivré.\\\passeLigne}%
\end{groupe}%
%% v. 1208-1209a
\begin{groupe}%
\personnageFr{Euripide}%
\trimIambFr{Sois libre. Le reste te regarde : fuis avant que l'archer\\
ne revienne et ne t'empoigne.\\\passeLigne}%
\end{groupe}%
%% v. 1209b
\begin{groupe}%
\personnageFr{Mnésiloque}%
\trimIambFr{\espace{2.9cm}C'est ce que je fais. (Euripide et Mnésiloque s'enfuient.)% chacun de leur côté.)
\\\passeLigne}%
\end{groupe}%
%% v.1210-1216 (part 1)
\begin{groupe}\personnageFr{Archer}
\trimIambFr{(Revenant en scène). Ah! petite vieille, la charmante fillette que tu as là,\\\passeLigne}%
\end{groupe}%
\endnumbering
\end{Rightside}
\end{pages}
\Pages
\clearpage\ifodd\value{page}\hbox{}\newpage\fi
\enTeteG{ΕΞΟΔΟΣ (Τέλος) --- Codex Ravennas 429, f. 177r (Classense 179r) --- vv. 1211–1231}
\begin{adjustwidth}{-1.5cm}{-3.0cm}
\begin{adjustbox}{max totalheight=21.0cm, max width=16.5cm, keepaspectratio}
\begin{minipage}{16.5cm}
{\fontsize{7.2pt}{8.5pt}\selectfont
\noindent\textbf{Folio 30 --- Codex Ravennas 429, \texttt{f. 177r} (\texttt{Classense 179r}) : ΕΞΟΔΟΣ (Τέλος), vv. 1211–1231}\hfill \textit{Témoins :} $\text{R}^1, \Rpc, \Rdeux$ (\Mon)\par\vspace{2pt}
\hrule height 0.4pt\vspace{2.5pt}
}
{\fontsize{8.1pt}{9.5pt}\selectfont
\ligneRav{1211}{τοξ/}{κοὐ \difference{δύσκολ᾽} ἀλλὰ πρᾶο. ποῦ τὸ γρᾴδιο;}{\textbf{1211} δύσκολλ᾽ R : \lecon{δύσκολ᾽} Grynaeus}
\ligneRav{}{}{οἴμ᾽ ὠς ἀπόλωλο· ποῦ τὸ γέροντ᾽ ἐντευτενί;}{}
\ligneRav{1213}{}{ὦ γρᾴδι᾽, ὦ \difference{γρᾶ᾽}. οὐκ ἐπαινῶ γρᾴδιο·}{\textbf{1213} γρᾶίο R : \lecon{γρᾶ᾽} Bentley}
\ligneRav{}{}{Ἀρταμουξία·}{}
\ligneRav{1214}{}{διέβαλλέ μ᾽ \difference{ἡ γραῦς}. ἀπότρεκ᾽ ὠς τάκιστα σύ·}{\textbf{1214} \lecon{διέβαλλε} R : διέβαλέ Σ, Suid. δ 892, Eust. p. 1406.19 \espApp ὁ γραῦς R : \lecon{ἡ γραῦς} Suid., Eust. \espApp ὡς τάχιστα R : ὠς τάκιστα Brunck}
\ligneRav{1215}{}{ὀρτῶς δὲ \difference{συμβήνη} 'στί· καταβηνῆσι γάρ·}{\textbf{1215} συβήνη R : \lecon{συμβήνη} Bibi, cf. v.1197, Hesych. σ 2137}
\ligneRav{1216}{}{οἴμοι, τί \difference{δρᾶσι}; ποῖ τὸ \difference{γρᾴδιο};}{\textbf{1216} δράσει R : \lecon{δρᾶσι} Blaydes (1842) \espApp γράδιο R : \lecon{γρᾴδιο} \Mon}
\ligneRav{}{}{Ἀρταμουξία·}{}
\ligneRav{1217}{χορ/}{τὴν \difference{γραῦν} ἐρωτᾷς, ἣ 'φερεν τὰς πηκτίδας;}{\textbf{1217} γραῦ R : \lecon{γραῦν} Suid. π 1502}
\ligneRav{}{τοξ/}{ναὶ ναῖκι. εἶδες αὐτό; : ταύτῃ γ᾽ οἴχεται}{}
\ligneRav{1219}{}{αὐτή τ᾽ ἐκείνη καὶ \difference{γέρων} τις εἵπετο·}{\textbf{1219} γέρον R : \lecon{γέρων} \Mon}
\ligneRav{1220}{τοξ/}{κροκῶτ᾽ ἔκοντο τὴ γέροντο; : φήμ᾽ ἐγώ·}{}
\ligneRav{}{}{ἔτ᾽ ἂν καταλάβοις, εἰ διώκοις ταυτῃί·}{}
\ligneRav{1222}{τοξ/}{ὦ μιαρὸ \difference{γρᾷο}· πότερα τρέξι τὴν \difference{ὀδό};}{\textbf{1222} γραῦ et ὁδό R : \lecon{γρᾷο} et \lecon{ὀδό} Brunck}
\ligneRav{}{}{Ἀρταμουξία·}{}
\ligneRav{}{χορ/}{ὀρθὴν ἄνω δίωκε. ποῖ θεῖς; οὐ πάλιν}{}
\ligneRav{1224}{}{τῃδὶ \difference{διώξει}; \difference{᾽ς} τοὔμπαλιν τρέχεις σύ γε·}{\textbf{1224} δ᾽ ὤξεις R : \lecon{διώξει} Elmsley (ad \italique{Ach.} 278) \lecon{᾽ς} Cobet (\italique{Mnemos.} 2 (1853), p.213)}
\ligneRav{1225}{τοξ/}{κακόδαιμον, ἀλλὰ τρέξι. Ἀρταμουξία·}{}
\ligneRav{1226}{χορ/}{\hspace{1.1em}τρέχε νυν \difference{κατ᾽ αὐτοὺς} \difference{ἐς} κόρακας ἐπουρίσας·}{\textbf{1226} κατὰ τοὺς R : \lecon{κατ᾽ αὐτοὺς} Jackson (p.82) \lecon{ἐς} O. Bachmann}
\ligneRav{1227}{}{\hspace{1.1em}ἀλλὰ \difference{πέπαισται} μετρίως ἡμῖν·}{\textbf{1227} πέπυσται R : \lecon{πέπαισται} Σ}
\ligneRav{}{}{\hspace{1.1em}ὥσθ᾽ ὥρα δή 'στι βαδίζειν}{}
\ligneRav{}{}{\hspace{1.1em}οἴκαδ᾽ ἑκάστῃ. τὼ Θεσμοφόρω δ᾽}{}
\ligneRav{1230}{}{\hspace{1.1em}ἡμῖν ἀγαθὴν}{}
\ligneRav{1231}{}{\hspace{1.1em}τούτων χάριν ἀνταποδοίτην·}{\textbf{1231} ἀνταδοῖτον R : \lecon{ἀνταποδοῖτον} Grynaeus}
\vspace{3pt}\hrule height 0.3pt\vspace{2pt}
\centerline{\footnotesize\textbf{$\ast\;\cdot\;\ast$ ἀριστοφάνους θεσμοφοριάζουσαι $\ast\;\cdot\;\ast$}}\vspace{2pt}
\hrule height 0.3pt\vspace{2pt}
{\scriptsize\noindent\textbf{[Moitié inférieure du \texttt{f. 177r} ($\Rdeux$) — \italique{Hypothesis} et \italique{Dramatis personae} des \italique{Ecclesiazusae} :]} Αἱ γυναῖκες συνέθεντο Σκίροις ὥστε εἰς τὴν ἐκκλησίαν ἀνδρείως σκευασάμεναι εἰσελθεῖν... Πραξαγόρα $\cdot$ Βλέπυρος $\cdot$ Γυνή τις $\cdot$ Χρέμης $\cdot$ Ἀνὴρ φειδωλός...}
}
\vspace{3pt}\noindent\rule{4.5cm}{0.4pt}\par\vspace{1.5pt}
{\fontsize{7.0pt}{8.3pt}\selectfont
\noindent\textbf{Loci paralleli (vv. 1211–1231) ---} \textbf{1214} Suidas δ 892 διέβαλέ μ' ἡ γραῦς ; \Eustathe 1406.19 \espTest \textbf{1217} Suidas π 1502 \espTest \textbf{1231} Subscriptio \PsiUn \ fr. 6 ΑΡΙΣΤΟΦΑΝ$]$ΟΥΣ ΘΕΣΜΟΦΟΡΙ$]$ΑΖΥΣΑΙ $]$ΧΗΗΙ ($=$ 1201, versuum numerus)\par
}
\vspace{2.5pt}\noindent\rule{4.5cm}{0.4pt}\par\vspace{1.5pt}
{\fontsize{6.9pt}{8.2pt}\selectfont
\noindent\textbf{Suite (en commentaire) de l'apparat critique \& notes paléographiques (\texttt{f. 177r}) ---} \textbf{Témoin papyrologique :} \textbf{\PsiUn{} (\italique{PSI} XI 1194, fr. 6)} conserve après le \textbf{v. 1231} la \italique{subscriptio} stichométrique antique en chiffres acrophoniques attiques : ΑΡΙΣΤΟΦΑΝ]ΟΥΣ ΘΕΣΜΟΦΟΡΙ]ΑΖΥΣΑΙ ]ΧΗΗΙ (= \textbf{1201 vers}). --- \textbf{Harmonisation de συμβήνη au v. 1215 par \textbf{Bibi} :} Là où R avait conservé συμβήνην avec $-\mu-$ au v. 1197 mais écrit συβήνη sans $-\mu-$ au v. 1215, \textbf{\textbf{Bibi} rétablit συμβήνη} au v. 1215 pour maintenir le jeu de mots avec $\text{συμβινεῖν / καταβηνῆσι}$ (cf. Hésychius $\sigma\ 2137$). --- \textbf{Notes marginales de locuteur au v. 1219 :} R porte Σκ. en marge devant le v. 1219 et Χορ. devant καὶ γέρον à l'intérieur du vers. --- \textbf{Faute récurrente au v. 1227 :} Comme au v. 1127, $\text{R}^1$ écrit πέπυσται au lieu de \lecon{πέπαισται} (« nous avons assez joué », conservé par $\Sigma^{\text{R}}$). --- \textbf{Haplographie syllabique au dernier vers (v. 1231) :} $\text{R}^1$ écrit ἀνταδοῖτον (saut de la syllabe -πο-) au lieu de ἀνταποδοῖτον (Grynaeus) / \lecon{ἀνταποδοίτην} (optatif duel adopté dans le texte par L. Hauseux). --- \textbf{v. 1214 :} \lecon{διέβαλλε} R (διέβαλέ $\Sigma$, Souda $\delta\ 892$, Eustathe p. 1406.19) \espApp ὁ γραῦς R : \lecon{ἡ γραῦς} Souda, Eustathe \espApp ὡς τάχιστα R : \lecon{ὠς τάκιστα} Brunck (adopté par L. Hauseux). --- \textbf{v. 1215 :} συβήνη R : \lecon{συμβήνη} \lecon{Bibi} (cf. v. 1197, Hésychius $\sigma\ 2137$). --- \textbf{v. 1219 :} Σκ. ante versum in mg. et Χορ. ante καὶ γέρον in R \espApp γέρον R : \lecon{γέρων} \Mon{}. --- \textbf{v. 1227 :} πέπυσται R : \lecon{πέπαισται} $\Sigma$.\par
}
\end{minipage}
\end{adjustbox}
\end{adjustwidth}
\newpage
\enTeteD{Codex Ravennas 429 (Biblioteca Classense) --- Fac-similé du f. 177r (vv. 1211–1231)}
\begin{adjustwidth}{-1.5cm}{-3.0cm}
\begin{center}
\includegraphics[width=16.5cm,height=20.5cm,keepaspectratio]{/workspaces/Thesmophories/dossier_ravennas_thesmophories/folios_individuels_ff162v-177r/30_f177r_classense179r_vv1211-1231_fin.jpg}\par
\vspace{2pt}
{\scriptsize\textbf{Codex Ravennas 429, f. 177r (Classense 179r)} --- ΕΞΟΔΟΣ (Τέλος), \textbf{vv. 1211–1231} ($\text{R}^1, \Rpc, \Rdeux$ (\Mon)).}
\end{center}
\end{adjustwidth}
\newpage
\clearpage\ifodd\value{page}\hbox{}\newpage\fi
\enTeteG{ΘΕΣΜΟΦΟΡΙΑΖΟΥΣΑΙ --- Texte grec établi (f. 177r, vv. 1211–1231)}
\enTeteD{LES THESMOPHORIES --- Traduction française (f. 177r, vv. 1211–1231)}
\begin{pages}
\begin{Leftside}
\beginnumbering
%% v. 1211+ (part 2)
\begin{groupe}
\trimIamb{\ligne{1211}κοὐ δύσκολ᾽ ἀλλὰ πρᾶο. ποῦ τὸ γρᾴδιο;\apparat{1211}{δύσκολλ᾽ R : \lecon{δύσκολ᾽} Grynaeus}\\
οἴμ᾽ ὠς ἀπόλωλο· ποῦ τὸ γέροντ᾽ ἐντευτενί;\\
ὦ γρᾴδι᾽, ὦ γρᾶ᾽. οὐκ ἐπαινῶ γρᾴδιο.\apparat{1213}{γρᾶίο R : \lecon{γρᾶ᾽} Bentley}\\
Ἀρταμουξία.\\
\ligne{1214}διέβαλλέ μ᾽ ἡ γραῦς. ἀπότρεκ᾽ ὠς τάκιστα σύ·\testimonia{1214}{Suidas δ 892 διέβαλέ μ' ἡ γραῦς ; \Eustathe 1406.19}\apparat{1214}{\lecon{διέβαλλε} R : διέβαλέ Σ, Suid. δ 892, Eust. p. 1406.19 \espApp ὁ γραῦς R : \lecon{ἡ γραῦς} Suid., Eust. \espApp ὡς τάχιστα R : ὠς τάκιστα Brunck \ligneApp{1215} συβήνη R : \lecon{συμβήνη} Bibi, cf. v.1197, Hesych. σ 2137 \ligneApp{1216} δράσει R : \lecon{δρᾶσι} Blaydes (1842) \espApp γράδιο R : \lecon{γρᾴδιο} \Mon}\\
ὀρτῶς δὲ συμβήνη 'στί· καταβηνῆσι γάρ.\\
οἴμοι, τί δρᾶσι; ποῖ τὸ γρᾴδιο;\commentaire{1216}{οἴμοι, τί δράσει; ποῖ τὸ γράδιο·}{Il manque un pied à ce vers.}\\
Ἀρταμουξία.\\\passeLigne}%
\end{groupe}%
%% v. 1217
\begin{groupe}%
\personnageGrec{Κορυφαῖος}%
\trimIamb{\ligne{1217}τὴν γραῦν ἐρωτᾷς, ἣ 'φερεν τὰς πηκτίδας;\testimonia{1217}{Suidas π 1502}\apparat{1217}{γραῦ R : \lecon{γραῦν} Suid. π 1502}\\\passeLigne}%
\end{groupe}%
%% v. 1218a
\begin{groupe}%
\personnageGrec{Τοξότης}%
\trimIamb{ναὶ ναῖκι. εἶδες αὐτό;\\\passeLigne}%
\end{groupe}%
%% v. 1218b-1219
\begin{groupe}%
\personnageGrec{Κορυφαῖος}%
\trimIamb{\ligne{1218}\decalage{ναὶ ναῖκι. εἶδες αὐτό;}ταύτῃ γ᾽ οἴχεται\\
αὐτή τ᾽ ἐκείνη καὶ γέρων τις εἵπετο.\apparat{1219}{%Σκ. ante versum in mg. et Χορ. ante καὶ γέρον in R \espApp 
γέρον R : \lecon{γέρων} \Mon}\\\passeLigne}%
\end{groupe}%
%% v. 1220a
\begin{groupe}%
\personnageGrec{Τοξότης}%
\trimIamb{κροκῶτ᾽ ἔκοντο τὴ γέροντο;\\\passeLigne}%
\end{groupe}%
%% v. 1220b-1221
\begin{groupe}%
\personnageGrec{Κορυφαῖος}%
\trimIamb{\ligne{1220}\decalage{κροκῶτ᾽ ἔκοντο τὴ γέροντο;}φήμ᾽ ἐγώ.\\
ἔτ᾽ ἂν καταλάβοις, εἰ διώκοις ταυτῃί.\\\passeLigne}%
\end{groupe}%
%% v. 1222a-1222b
\begin{groupe}%
%%
\personnageGrec{Τοξότης}%
\trimIamb{ὦ μιαρὸ γρᾷο· πότερα τρέξι τὴν ὀδό;\apparat{1222}{γραῦ et ὁδό R : \lecon{γρᾷο} et \lecon{ὀδό} Brunck}\\
\ligne{1222}Ἀρταμουξία.\\\passeLigne}%
\end{groupe}%
%% v. 1223-1224
\begin{groupe}%
\personnageGrec{Κορυφαῖος}%
\trimIamb{ὀρθὴν ἄνω δίωκε. ποῖ θεῖς; οὐ πάλιν\\
τῃδὶ διώξει; ᾽ς τοὔμπαλιν τρέχεις σύ γε.\apparat{1224}{δ᾽ ὤξεις R : \lecon{διώξει} Elmsley (ad \italique{Ach.} 278) \lecon{᾽ς} Cobet (\italique{Mnemos.} 2 (1853), p.213)}\\\passeLigne}%
\end{groupe}%
%% v. 1225
\begin{groupe}%
\personnageGrec{Τοξότης}%
\trimIamb{κακόδαιμον, ἀλλὰ τρέξι. Ἀρταμουξία.\\\passeLigne}%
\end{groupe}%
%% v. 1226-1231
\begin{groupe}%
\personnageGrec{Κορυφαῖος}%
\typeVers{anap.}\anap{τρέχε νυν κατ᾽ αὐτοὺς ἐς κόρακας ἐπουρίσας.\apparat{1226}{κατὰ τοὺς R : \lecon{κατ᾽ αὐτοὺς} Jackson (p.82) \lecon{ἐς} O. Bachmann \ligneApp{1127} πέπυσται R : \lecon{πέπαισται} Σ}\\
ἀλλὰ πέπαισται μετρίως ἡμῖν·\\%\apparat{1227}{πέπυσται R : \lecon{πέπαισται} Σ}\\
ὥσθ᾽ ὥρα δή 'στι βαδίζειν\\
οἴκαδ᾽ ἑκάστῃ. τὼ Θεσμοφόρω δ᾽\commentaire{1129-31}{τὼ Θεσμοφόρω δ᾽ / ἡμῖν ἀγαθὴν / τούτων χάριν ἀνταποδοίτην.}{Selon son habitude -- qui était peut-être d'usage dans l'Ancienne comédie --, Artistophane conclue sa pièce par des souhaits de succès (il vise bien-sûr la victoire au concours comique).}\\
ἡμῖν ἀγαθὴν\\
τούτων χάριν ἀνταποδοίτην.\testimonia{1231}{Subscriptio \PsiUn \ fr. 6 ΑΡΙΣΤΟΦΑΝ$]$ΟΥΣ    ΘΕΣΜΟΦΟΡΙ$]$ΑΖΥΣΑΙ   $]$ΧΗΗΙ ($=$ 1201, versuum numerus)}\apparat{1231}{ἀνταδοῖτον R : \lecon{ἀνταποδοῖτον} Grynaeus}\\\passeLigne}%
\end{groupe}%
\endnumbering
\end{Leftside}
\begin{Rightside}
\beginnumbering
%% v. 1211+ (part 2)
\begin{groupe}
\trimIambFr{\ligne{1211}et pas méchante... gentille! Où est la vieille ?\\
Ah ! je suis perdu! Où s'est éclipsé le vieillard ?\\
Ah! petite vieille, petite vieille ! Ce n'est pas bien, petite vieille.\\
(Appelant.) Artamouxia !\\
Elle m'a trompé, la vieille. (La danseuse s'esquive.) Va-t'en au plus vite, toi.\\
(Reprenant sa lance :) La lance est droite, je descends.\\
Ah ! que vais-je faire ? Où est la petite vieille ?\\
Artamouxia !\\\passeLigne}%
\end{groupe}%
%% v. 1217
\begin{groupe}%
\personnageFr{Coryphée}%
\trimIambFr{Tu demandes après la vieille qui portait les pectis ?\\\passeLigne}%
\end{groupe}%
%% v. 1218a
\begin{groupe}%
\personnageFr{Archer}%
\trimIambFr{Si fait, si. L'as-tu vue ?\\\passeLigne}%
\end{groupe}%
%% v. 1218b-1219
\begin{groupe}%
\personnageFr{Coryphée}%
\trimIambFr{\espace{3.54cm}C'est par là qu'elle a disparu,\\
elle et un vieux qui la suivait.\\\passeLigne}%
\end{groupe}%
%% v. 1220a
\begin{groupe}%
\personnageFr{Archer}%
\trimIambFr{Il avait une crocote, le vieux ?\\\passeLigne}%
\end{groupe}%
%% v. 1220b-1221
\begin{groupe}%
\personnageFr{Coryphée}%
\trimIambFr{\espace{3.9cm}Tout juste.\\
Tu peux encore les rattraper en tirant de ce côté. (Indiquant une autre direction.)\\\passeLigne}%
\end{groupe}%
%% v. 1222a-1222b
\begin{groupe}%
\personnageFr{Archer}%
\trimIambFr{Ô la vieille canaille! Par quel chemin prendre ?\\
\ligne{1222}Artamouxia ! (ll court vers l'entrée de l'orchestre.)\\\passeLigne}%
\end{groupe}%
%% v. 1223-1224
\begin{groupe}%
\personnageFr{Coryphée}%
\trimIambFr{Monte tout droit. Où cours-tu ?\\
Veux-tu bien revenir ici ? Tu cours à contresens.\\\passeLigne}%
\end{groupe}%
%% v. 1225
\begin{groupe}%
\personnageFr{Archer}%
\trimIambFr{Malheureux !... Mais je me sauve. Artamouxia !\\\passeLigne}%
\end{groupe}%
%% v. 1226-1231
\begin{groupe}%
\personnageFr{Coryphée}%
\anapFr{Cours maintenant, cours aux corbeaux, et qu'un bon vent te seconde.\\
Mais nous nous sommes assez ébattues.\\
Il est temps pour chacune de rentrer au logis.\\
Puissent, en retour, les Thesmophores\\
nous accorder\\
leur précieuse faveur !\\\passeLigne}%
\end{groupe}%
\endnumbering
\end{Rightside}
\end{pages}
\Pages
\end{document}
